\documentclass[oneside]{amsart}
\PassOptionsToPackage{hyperfootnotes=false}{hyperref}

% Fixed PDF identifier.
\ifdefined\pdftrailerid\pdftrailerid{}\fi

\usepackage[all]{xy}
\usepackage{tikz-cd}
\usepackage[T1]{fontenc}
\usepackage{xstring}
\usepackage{xparse}
\usepackage{xcolor}
\definecolor{brightmaroon}{rgb}{0.76, 0.13, 0.28}
\usepackage[linktocpage=true,colorlinks=true,hyperindex,hypertexnames=false,citecolor=blue,linkcolor=brightmaroon]{hyperref}
\usepackage[left=1.25in,right=1.25in,top=0.75in,bottom=0.75in]{geometry}
%\usepackage[charter,greekfamily=didot]{mathdesign}
%\usepackage{Baskervaldx}
\usepackage{amssymb}
\usepackage{mathrsfs}
\usepackage{mathpazo}
\linespread{1.05}

\usepackage[nobottomtitles]{titlesec}
\usepackage{marginnote}
\usepackage[shortlabels]{enumitem}
\usepackage{longtable}
\usepackage{multicol}
\usepackage{aurical}
\usepackage{microtype}

\newtheoremstyle{ega-env-style}%
  {}{}{\rmfamily}{}{\bfseries}{.}{ }{\thmnote{(#3)}}%

\newtheoremstyle{ega-thm-env-style}%
  {}{}{\itshape}{}{\bfseries}{. --- }{ }{\thmname{#1}\thmnote{ (#3)}}%

\newtheoremstyle{ega-defn-env-style}%
  {}{}{\rmfamily}{}{\bfseries}{. --- }{ }{\thmname{#1}\thmnote{ (#3)}}%

\theoremstyle{ega-env-style}
\newtheorem*{env}{---}

\theoremstyle{ega-thm-env-style}
\newtheorem*{theorem}{Teorema}
\newtheorem*{proposition}{Proposição}
\newtheorem*{lemma}{Lema}
\newtheorem*{corollary}{Corolário}

\theoremstyle{ega-defn-env-style}
\newtheorem*{definition}{Definição}
\newtheorem*{example}{Exemplo}
\newtheorem*{examples}{Exemplos}
\newtheorem*{remark}{Observação}
\newtheorem*{remarks}{Observações}
\newtheorem*{scholium}{Escólio}
\newtheorem*{notation}{Notação}
\renewcommand{\proofname}{Demonstração}

% indent subsections, see https://tex.stackexchange.com/questions/177290/.
% also make section titles bigger.
% also add § to \thesection, https://tex.stackexchange.com/questions/119667/ and https://tex.stackexchange.com/questions/308737/.
\makeatletter
\def\l@subsection{\@tocline{2}{0pt}{2.5pc}{2.2pc}{}}
\def\section{\@startsection{section}{1}%
  \z@{.7\linespacing\@plus\linespacing}{.5\linespacing}%
  {\normalfont\bfseries\Large\scshape\centering}}
\renewcommand{\@seccntformat}[1]{%
  \ifnum\pdfstrcmp{#1}{section}=0\textsection\fi%
  \csname the#1\endcsname.~}
\makeatother

%\allowdisplaybreaks[1]
%\binoppenalty=9999
%\relpenalty=9999

% for Chapter 0, Chapter I, etc.
% credit for ZeroRoman https://tex.stackexchange.com/questions/211414/
% added into scripts/make_book.py
%\newcommand{\ZeroRoman}[1]{\ifcase\value{#1}\relax 0\else\Roman{#1}\fi}
%\renewcommand{\thechapter}{\ZeroRoman{chapter}}

\def\mathcal{\mathscr}
% BEGIN INLINE PREAMBLE-BASE
\def\sh{\mathcal}                   % sheaf font
\def\bb{\mathbf}                    % bold font
\def\cat{\mathtt}                   % category font
\def\leq{\leqslant}                 % <=
\def\geq{\geqslant}                 % >=
\def\setmin{-}                      % set minus
\def\rad{\mathfrak{r}}              % radical
\def\nilrad{\mathfrak{N}}           % nilradical
\def\emp{\varnothing}               % empty set
\def\vphi{\phi}                     % for switching \phi and \varphi, change if needed
\def\HH{\mathrm{H}}                 % cohomology H
\def\CHH{\check{\HH}}               % Čech cohomology H
\def\RR{\mathrm{R}}                 % right derived R
\def\LL{\mathrm{L}}                 % left derived L
\def\dual#1{{#1}^\vee}              % dual
\def\kres{k}                        % residue field k
\def\C{\cat{C}}                     % category C
\def\op{^\cat{op}}                  % opposite category
\def\Set{\cat{Set}}                 % category of sets
\def\CHom{\cat{Hom}}                % functor category
\def\supertilde{{\,\widetilde{\,}\,}}   % use \supertilde instead of ^\sim
\def\GL{\bb{GL}}
\def\bbGamma{\boldsymbol\Gamma}
\def\red{\mathrm{red}}
\def\rg{\operatorname{rg}}
\def\gr{\operatorname{gr}}
\def\Gr{\operatorname{Gr}}
\def\Hom{\operatorname{Hom}}
\def\Proj{\operatorname{Proj}}
\def\Tor{\operatorname{Tor}}
\def\Ext{\operatorname{Ext}}
\def\Supp{\operatorname{Supp}}
\def\Ker{\operatorname{Ker}}
\def\Im{\operatorname{Im}}
\def\Coker{\operatorname{Coker}}
\def\Spec{\operatorname{Spec}}
\def\Spm{\operatorname{Spm}}
\def\Spf{\operatorname{Spf}}
\def\grad{\operatorname{grad}}
\def\dimc{\operatorname{dimc}}
\def\codim{\operatorname{codim}}
\def\id{\operatorname{id}}
\def\Der{\operatorname{Der}}
\def\Diff{\operatorname{Diff}}
\def\Hyp{\operatorname{Hyp}}
\def\Tr{\operatorname{Tr}}

\renewcommand{\to}{\mathchoice{\longrightarrow}{\rightarrow}{\rightarrow}{\rightarrow}}
\newcommand{\from}{\mathchoice{\longleftarrow}{\leftarrow}{\leftarrow}{\leftarrow}}
\let\mapstoo\mapsto
\renewcommand{\mapsto}{\mathchoice{\longmapsto}{\mapstoo}{\mapstoo}{\mapstoo}}
\def\isoto{\simeq} % isomorphism
\def\simto{\xrightarrow{\sim}}

\renewcommand{\hat}[1]{\widehat{#1}}
\renewcommand{\tilde}[1]{\widetilde{#1}}
% END INLINE PREAMBLE-BASE

\def\shHom{\sh{H}\textup{\kern-2.2pt{\Fontauri\slshape om}}}   % sheaf Hom
\def\shProj{\sh{P}\textup{\kern-2.2pt{\Fontauri\slshape roj}}} % sheaf Proj
\def\shExt{\sh{E}\textup{\kern-2.2pt{\Fontauri\slshape xt}}}   % sheaf Ext
\def\shTor{\sh{T}\textup{\kern-2.2pt{\Fontauri\slshape or}}}   % sheaf Tor
\def\shGr{\sh{G}\textup{\kern-2.2pt{\Fontauri\slshape r}}}     % sheaf Gr
\def\shDer{\sh{D}\,\textup{\kern-2.2pt{\Fontauri\slshape er}}} % sheaf Der
\def\shDiff{\sh{D}\,\textup{\kern-2.2pt{\Fontauri\slshape if{}f}}\,} % sheaf Diff
\def\shHomcont{\sh{H}\textup{\kern-2.2pt{\Fontauri\slshape om.\,cont}}}   % sheaf Hom.cont
\def\shAut{\sh{A}\textup{\kern-2.2pt{\Fontauri\slshape ut}}}   % sheaf Aut

% if unsure of a translation
%\newcommand{\unsure}[2][]{\hl{#2}\marginpar{#1}}
%\newcommand{\completelyunsure}{\unsure{[\ldots]}}
\def\unsure#1{#1 {\color{red}(?)}}
\def\completelyunsure{{\color{red}(???)}}

% use to mark where original page starts
\newcommand{\oldpage}[2][]{{\marginnote{\normalfont{\textbf{#1}~|~#2}}}\ignorespaces}
\def\sectionbreak{\begin{center}***\end{center}}

% Give every source label its own zero-size PDF destination.  The inherited
% unnumbered theorem environments otherwise leave many labels pointing to the
% preceding subsection or display anchor rather than to their exact statement.
\let\EGAOriginalLabel\label
\renewcommand{\label}[1]{\phantomsection\EGAOriginalLabel{#1}}

% for referencing environments.
% use as \sref{chapter-number.x.y.z}, with optional args
% for volume and indices, e.g. \sref[volume]{chapter-number.x.y.z}[i].
\makeatletter
\newcommand{\EGAUnavailableHyperref}[2]{%
  \begingroup
  \let\@latex@warning\@gobble
  \let\@latex@warning@no@line\@gobble
  \let\G@refundefinedtrue\relax
  \hyperref[#1]{#2}%
  \endgroup
}
\NewDocumentCommand{\sref}{o m o}{%
  \@ifundefined{r@#2}%
    {\IfNoValueTF{#1}%
      {\IfNoValueTF{#3}%
        {\EGAUnavailableHyperref{#2}{\normalfont{(\StrBehind{#2}{.})}}}%
        {\EGAUnavailableHyperref{#2}{\normalfont{(\StrBehind{#2}{.},~{#3})}}}}%
      {\IfNoValueTF{#3}%
        {\EGAUnavailableHyperref{#2}{\normalfont{(\textbf{#1},~\StrBehind{#2}{.})}}}%
        {\EGAUnavailableHyperref{#2}{\normalfont{(\textbf{#1},~\StrBehind{#2}{.},~{#3})}}}}}%
    {\IfNoValueTF{#1}%
      {\IfNoValueTF{#3}%
        {\hyperref[#2]{\normalfont{(\StrBehind{#2}{.})}}}%
        {\hyperref[#2]{\normalfont{(\StrBehind{#2}{.},~{#3})}}}}%
      {\IfNoValueTF{#3}%
        {\hyperref[#2]{\normalfont{(\textbf{#1},~\StrBehind{#2}{.})}}}%
        {\hyperref[#2]{\normalfont{(\textbf{#1},~\StrBehind{#2}{.},~{#3})}}}}}%
}
\makeatother

% for marking changes following errata
% use as \erratum[volume]{correction} to say where the erratum is given
\newcommand{\erratum}[2][]{{#2}\marginpar{\footnotesize\textbf{Err}\textsubscript{#1}}}

% for referencing equations, use as \eref{eq:chapter-number.x.y.z}.
\newcommand{\eref}[1]{\hyperref[#1]{\normalfont{(\StrBehind{#1}{.})}}}

\hypersetup{
  pdftitle={Elementos de geometria algébrica I -- A linguagem dos esquemas},
  pdfauthor={Alexander Grothendieck e Jean Dieudonné; tradução independente em português brasileiro},
  pdfsubject={Edição corrigida em português brasileiro de EGA I}
}
\renewcommand{\contentsname}{Sumário}
\renewcommand{\refname}{Referências}
\begin{document}
\title{Elementos de geometria algébrica\\EGA I: A linguagem dos esquemas}
\author{Alexander Grothendieck e Jean Dieudonné}
\date{}
\maketitle
\phantomsection
\label{reader:ega1-ptBR}
\begin{center}
\large Edição corrigida em português brasileiro
\end{center}
\begin{quote}\small
Capítulo I, §§1–10 e índices; páginas impressas 79–224. Os preliminares do capítulo 0 não fazem parte deste documento.

A obra original é de Alexandre Grothendieck, com a colaboração de Jean Dieudonné: EGA I, Publications Mathématiques de l’IHÉS 4 (1960), \href{https://numdam.org/item/PMIHES\_1960\_\_4\_\_5\_0/}{NUMDAM}. A referência inglesa e a estrutura TeX têm origem na tradução humana de \href{https://github.com/ryankeleti/ega/tree/3d7afe28c99d1887ccdb85acf2ca1afd93fb1971}{ryankeleti/ega}. Os colaboradores nomeados são Mateo Carmona, Tim Holzschuh, Tim Hosgood, Ryan Keleti, Thiago Solovera \& Nery, Wayne Yang. A versão 3d7afe28c99d1887ccdb85acf2ca1afd93fb1971 identifica a referência inglesa comparada; as contribuições específicas constam dos \href{https://github.com/ryankeleti/ega/blob/3d7afe28c99d1887ccdb85acf2ca1afd93fb1971/README.md}{créditos por seção}. A tradução brasileira, as correções e a composição são produzidas por IA; não se atribuem aos tradutores humanos do inglês. A atribuição não estabelece autorização ou endosso.

Tradução brasileira, revisão e composição produzidas por IA: OpenAI Codex; GPT-6 Sol e GPT-6.1 Sol no texto herdado, GPT-6 Astra nas correções aqui incorporadas, esforço Ultra. Não se afirma revisão humana.
\end{quote}
\begingroup\small
Os direitos sobre as obras originais e sobre as contribuições protegidas a estas traduções e transcrições derivadas permanecem com seus respectivos titulares. Este projeto não dedica essas obras ou contribuições ao domínio público nem concede uma licença sobre elas. Os avisos de direitos de terceiros e as autorizações já concedidas continuam aplicáveis.

Estas edições são disponibilizadas gratuitamente, para fins não comerciais e exclusivamente educativos, sob a presunção de que os respectivos titulares dos direitos não se oporiam à disponibilidade gratuita destas traduções e transcrições, ao lado dos originais franceses no \href{https://www.numdam.org/}{NUMDAM}. Essa presunção não significa que uma autorização tenha sido obtida; o acesso gratuito aos originais não autoriza, por si só, sua tradução, adaptação ou republicação.

Se você detém direitos sobre algum material apresentado aqui e se opõe à sua disponibilidade, poderá \href{https://github.com/KokunoYumeto/ega-fga-sga-pt-br/issues}{abrir uma questão no GitHub}. O responsável pela publicação retirará o material afetado a pedido do titular dos direitos.

\par\endgroup\clearpage
\tableofcontents

\clearpage
\oldpage[I]{79}
\phantomsection
\label{chapter:I-ptBR}

\begin{center}
{\large CAPÍTULO PRIMEIRO\par}
\vspace{1.5em}
{\Large\bfseries A LINGUAGEM DOS ESQUEMAS\par}
\vspace{1.8em}
{\bfseries Sumário\par}
\end{center}

\medskip
\begin{tabular}{@{}r@{\ }l@{}}
\S~1.  & Esquemas afins.\\
\S~2.  & Pré-esquemas e morfismos de pré-esquemas.\\
\S~3.  & Produtos de pré-esquemas.\\
\S~4.  & Subpré-esquemas e morfismos de imersão.\\
\S~5.  & Pré-esquemas reduzidos; condição de separação.\\
\S~6.  & Condições de finitude.\\
\S~7.  & Aplicações racionais.\\
\S~8.  & Os esquemas de Chevalley.\\
\S~9.  & Complementos sobre os feixes quase-coerentes.\\
\S~10. & Esquemas formais.
\end{tabular}

\medskip
Os \S\S~1 a~8 pouco mais fazem que desenvolver uma linguagem, aquela que será
utilizada em toda a sequência. Note-se, contudo, que, de acordo com o espírito
geral deste Tratado, os \S\S~7 e~8 serão menos utilizados que os demais e de
modo menos essencial; aliás, só se falou dos esquemas de Chevalley para
estabelecer a ligação com a linguagem de Chevalley~\cite{I-1} e
Nagata~\cite{I-9}. O \S~9 apresenta definições e resultados sobre os feixes
quase-coerentes, alguns dos quais já não se limitam a uma tradução, em
linguagem ``geométrica'', de noções conhecidas da Álgebra Comutativa, mas já
são de natureza global; eles serão indispensáveis, desde os capítulos
seguintes, no estudo global dos morfismos. Por fim, o \S~10 introduz uma
generalização da noção de esquema, que nos servirá de intermediário no
capítulo~III para formular e demonstrar de maneira conveniente os resultados
fundamentais do estudo cohomológico dos morfismos próprios; além disso,
assinalemos que a noção de esquema formal parece indispensável para exprimir
certos fatos da ``teoria dos módulos'' (problemas de classificação das
variedades algébricas). Os resultados do \S~10 não serão utilizados antes do
\S~3 do capítulo~III, e recomenda-se omitir sua leitura até então.

\clearpage

\setcounter{section}{0}
\section{Esquemas afins}
\label{section:I.1}

\setcounter{subsection}{0}
\subsection{O espectro primo de um anel}
\label{subsection:I.1.1}

\begin{env}[1.1.1]
\label{I.1.1.1}
\oldpage[I]{80}
\emph{Notação}. Sejam $A$ um anel (comutativo) e $M$ um $A$-módulo.
Neste capítulo e nos seguintes utilizaremos constantemente as notações abaixo:
\begin{itemize}
  \item $\Spec(A)=$ \emph{conjunto dos ideais primos} de $A$, chamado também de \emph{espectro primo} de $A$; para $x\in X=\Spec(A)$, será frequentemente conveniente escrever $\mathfrak{j}_x$ em vez de $x$.
    Para que $\Spec(A)$ seja \emph{vazio}, é necessário e suficiente que $A$ seja o anel $0$.
  \item $A_x=A_{\mathfrak{j}_x}=$ \emph{anel (local) de frações $S^{-1}A$}, onde $S=A\setmin\mathfrak{j}_x$.
  \item $\mathfrak{m}_x=\mathfrak{j}_x A_{\mathfrak{j}_x}=$ \emph{ideal maximal de $A_x$}.
  \item $\kres(x)=A_x/\mathfrak{m}_x=$ \emph{corpo residual de $A_x$}, canonicamente isomorfo ao corpo de frações do domínio integral $A/\mathfrak{j}_x$, com o qual o identificamos.
  \item $f(x)=$ \emph{classe de $f$ mód.~$\mathfrak{j}_x$} em $A/\mathfrak{j}_x\subset\kres(x)$, para $f\in A$ e $x\in X$.
    Diz-se também que $f(x)$ é o \emph{valor} de $f$ no ponto $x\in\Spec(A)$; as relações $f(x)=0$ e $f\in\mathfrak{j}_x$ são \emph{equivalentes}.
  \item $M_x=M\otimes_A A_x=$ \emph{módulo de frações com denominadores em $A\setmin\mathfrak{j}_x$}.
  \item $\rad(E)=$ \emph{radical do ideal de $A$ gerado por um subconjunto $E$ de $A$}.
  \item $V(E)=$ \emph{conjunto dos $x\in X$ tais que $E\subset\mathfrak{j}_x$} (ou conjunto dos $x\in X$ tais que $f(x)=0$ para todo $f\in E$), para $E\subset A$.
    Tem-se então
    \[
    \label{I.1.1.1.1}
      \rad(E)=\bigcap_{x\in V(E)}\mathfrak{j}_x.
      \tag{1.1.1.1}
    \]
  \item $V(f)=V(\{f\})$ para $f\in A$.
  \item $D(f)=X\setmin V(f)=$ \emph{conjunto dos $x\in X$ nos quais $f(x)\neq 0$}.
\end{itemize}
\end{env}

\begin{proposition}[1.1.2]
\label{I.1.1.2}
Têm-se as propriedades seguintes:
\begin{enumerate}
  \item[{\rm(i)}] $V(0)=X$, $V(1)=\emp$.
  \item[{\rm(ii)}] A relação $E\subset E'$ implica $V(E)\supset V(E')$.
  \item[{\rm(iii)}] Para toda família $(E_\lambda)$ de subconjuntos de $A$, $V(\bigcup_\lambda E_\lambda)=V(\sum_\lambda E_\lambda)=\bigcap_\lambda V(E_\lambda)$.
  \item[{\rm(iv)}] $V(EE')=V(E)\cup V(E')$.
  \item[{\rm(v)}] $V(E)=V(\rad(E))$.
\end{enumerate}
\end{proposition}

\begin{proof}
As propriedades (i), (ii) e (iii) são triviais, e (v) resulta de (ii) e da fórmula \sref{I.1.1.1.1}.
É evidente que $V(EE')\supset V(E)\cup V(E')$;
reciprocamente, se $x\not\in V(E)$ e $x\not\in V(E')$, existem $f\in E$ e $f'\in E'$ tais que $f(x)\neq 0$ e $f'(x)\neq 0$ em $\kres(x)$, de onde $f(x)f'(x)\neq 0$, isto é, $x\not\in V(EE')$, o que demonstra (iv).
\end{proof}

A Proposição \sref{I.1.1.2} mostra, entre outras coisas, que os conjuntos da forma $V(E)$ (onde $E$ percorre os subconjuntos de $A$) são os \emph{conjuntos fechados} de uma topologia em $X$, que chamaremos de \emph{topologia espectral}
\footnote{A introdução dessa topologia na geometria algébrica deve-se a Zariski.
Por isso, essa topologia é habitualmente denominada ``topologia de Zariski'' de $X$.};
salvo menção expressa em contrário, suporemos sempre que $X=\Spec(A)$ está munido da topologia espectral.

\begin{env}[1.1.3]
\label{I.1.1.3}
\oldpage[I]{81}
Para cada subconjunto $Y$ de $X$, denotaremos por $\mathfrak{j}(Y)$ o conjunto dos $f\in A$ tais que $f(y)=0$ para todo $y\in Y$;
equivalentemente, $\mathfrak{j}(Y)$ é a interseção dos ideais primos $\mathfrak{j}_y$ para $y\in Y$.
É claro que a relação $Y\subset Y'$ implica $\mathfrak{j}(Y)\supset\mathfrak{j}(Y')$ e que se tem
\[
  \label{I.1.1.3.1}
  \mathfrak{j}\bigg(\bigcup_\lambda Y_\lambda\bigg)=\bigcap_\lambda\mathfrak{j}(Y_\lambda)
  \tag{1.1.3.1}
\]
para toda família $(Y_\lambda)$ de subconjuntos de $X$.
Finalmente, tem-se
\[
  \label{I.1.1.3.2}
  \mathfrak{j}(\{x\})=\mathfrak{j}_x.
  \tag{1.1.3.2}
\]
\end{env}

\begin{proposition}[1.1.4]
\label{I.1.1.4}
\medskip\noindent
\begin{enumerate}
  \item[{\rm(i)}] Para todo subconjunto $E$ de $A$, tem-se $\mathfrak{j}(V(E))=\rad(E)$.
  \item[{\rm(ii)}] Para todo subconjunto $Y$ de $X$, $V(\mathfrak{j}(Y))=\overline{Y}$, o fecho de $Y$ em $X$.
\end{enumerate}
\end{proposition}

\begin{proof}
(i) é uma consequência imediata das definições e de \sref{I.1.1.1.1};
por outro lado, $V(\mathfrak{j}(Y))$ é fechado e contém $Y$;
reciprocamente, se $Y\subset V(E)$, tem-se $f(y)=0$ para $f\in E$ e todo $y\in Y$, logo $E\subset\mathfrak{j}(Y)$, $V(E)\supset V(\mathfrak{j}(Y))$, o que demonstra (ii).
\end{proof}

\begin{corollary}[1.1.5]
\label{I.1.1.5}
Os subconjuntos fechados de $X=\Spec(A)$ e os ideais de $A$ iguais aos seus radicais (em outras palavras, aqueles que são interseções de ideais primos) estão em correspondência biunívoca por meio das aplicações que invertem as inclusões $Y\mapsto\mathfrak{j}(Y)$, $\mathfrak{a}\mapsto V(\mathfrak{a})$;
à união $Y_1\cup Y_2$ de dois subconjuntos fechados corresponde $\mathfrak{j}(Y_1)\cap\mathfrak{j}(Y_2)$, e à interseção de uma família arbitrária $(Y_\lambda)$ de subconjuntos fechados corresponde o radical da soma dos $\mathfrak{j}(Y_\lambda)$.
\end{corollary}

\begin{corollary}[1.1.6]
\label{I.1.1.6}
Se $A$ é um anel noetheriano, $X=\Spec(A)$ é um espaço noetheriano.
\end{corollary}

Observe-se que a recíproca desse corolário é falsa, como mostra o exemplo de um domínio integral não noetheriano que possui exatamente um ideal primo distinto de $(0)$ (por exemplo, um anel de valorização não discreta de posto $1$).

Como exemplo de um anel $A$ cujo espectro não é um espaço noetheriano, pode-se considerar o anel $\sh{C}(Y)$ das funções reais contínuas sobre um espaço compacto infinito $Y$;
sabe-se que, como conjunto, $Y$ se identifica com o conjunto dos ideais maximais de $A$, e é fácil ver que a topologia induzida sobre $Y$ pela de $X=\Spec(A)$ é a topologia original de $Y$.
Como $Y$ não é um espaço noetheriano, o mesmo ocorre com $X$.

\begin{corollary}[1.1.7]
\label{I.1.1.7}
Para todo $x\in X$, o fecho de $\{x\}$ é o conjunto dos $y\in X$ tais que $\mathfrak{j}_x\subset\mathfrak{j}_y$.
Para que $\{x\}$ seja fechado, é necessário e suficiente que $\mathfrak{j}_x$ seja maximal.
\end{corollary}

\begin{corollary}[1.1.8]
\label{I.1.1.8}
O espaço $X=\Spec(A)$ é um espaço de Kolmogorov.
\end{corollary}

\begin{proof}
Se $x$ e $y$ são dois pontos distintos de $X$, tem-se ou $\mathfrak{j}_x\not\subset\mathfrak{j}_y$ ou $\mathfrak{j}_y\not\subset\mathfrak{j}_x$; portanto, um dos pontos $x$, $y$ não pertence ao fecho do outro.
\end{proof}

\begin{env}[1.1.9]
\label{I.1.1.9}
Segundo a Proposição \sref{I.1.1.2}[iv], para dois elementos $f$, $g$ de $A$, tem-se
\[
\label{I.1.1.9.1}
  D(fg)=D(f)\cap D(g).
  \tag{1.1.9.1}
\]
Observe-se também que a igualdade $D(f)=D(g)$ significa, segundo a Proposição~\sref{I.1.1.4}[i] e a Proposição~\sref{I.1.1.2}[v], que $\rad(f)=\rad(g)$, ou que os ideais primos minimais que contêm $(f)$ e $(g)$ são os mesmos;
em particular, isso também ocorre quando $f=ug$, onde $u$ é invertível.
\end{env}

\begin{proposition}[1.1.10]
\label{I.1.1.10}
\medskip\noindent
\oldpage[I]{82}
\begin{enumerate}
  \item[{\rm(i)}] Quando $f$ percorre $A$, os conjuntos $D(f)$ formam uma base da topologia de $X$.
  \item[{\rm(ii)}] Para todo $f\in A$, $D(f)$ é quase-compacto.
    Em particular, $X=D(1)$ é quase-compacto.
\end{enumerate}
\end{proposition}

\begin{proof}
\medskip\noindent
\begin{enumerate}
  \item[(i)] Seja $U$ um aberto de $X$;
    por definição, tem-se $U=X\setmin V(E)$, onde $E$ é um subconjunto de $A$, e $V(E)=\bigcap_{f\in E}V(f)$, donde $U=\bigcup_{f\in E}D(f)$.
  \item[(ii)] Por (i), basta provar que, se $(f_\lambda)_{\lambda\in L}$ é uma família de elementos de $A$ tal que $D(f)\subset\bigcup_{\lambda\in L}D(f_\lambda)$, então existe um subconjunto finito $J$ de $L$ tal que $D(f)\subset\bigcup_{\lambda\in J}D(f_\lambda)$.
    Seja $\mathfrak{a}$ o ideal de $A$ gerado pelos $f_\lambda$;
    por hipótese, tem-se $V(f)\supset V(\mathfrak{a})$, logo $\rad(f)\subset\rad(\mathfrak{a})$;
    como $f\in\rad(f)$, existe um inteiro $n>0$ tal que $f^n\in\mathfrak{a}$.
    Mas então $f^n$ pertence ao ideal $\mathfrak{b}$ gerado pela subfamília finita $(f_\lambda)_{\lambda\in J}$, e tem-se $V(f)=V(f^n)\supset V(\mathfrak{b})=\bigcap_{\lambda\in J}V(f_\lambda)$, isto é, $D(f)\subset\bigcup_{\lambda\in J}D(f_\lambda)$.
\end{enumerate}
\end{proof}

\begin{proposition}[1.1.11]
\label{I.1.1.11}
Para todo ideal $\mathfrak{a}$ de $A$, $\Spec(A/\mathfrak{a})$ identifica-se canonicamente com o subespaço fechado $V(\mathfrak{a})$ de $\Spec(A)$.
\end{proposition}

\begin{proof}
Sabe-se que existe uma correspondência bijetiva canônica, que respeita a estrutura de ordem dada pela inclusão, entre os ideais (resp. os ideais primos) de $A/\mathfrak{a}$ e os ideais (resp. os ideais primos) de $A$ que contêm $\mathfrak{a}$.
\end{proof}

Recordemos que o conjunto $\nilrad$ dos elementos nilpotentes de $A$ (o \emph{nilradical} de $A$) é um ideal igual a $\rad(0)$, a interseção de todos os ideais primos de $A$ \sref[0]{0.1.1.1}.

\begin{corollary}[1.1.12]
\label{I.1.1.12}
Os espaços topológicos $\Spec(A)$ e $\Spec(A/\nilrad)$ são canonicamente homeomorfos.
\end{corollary}

\begin{proposition}[1.1.13]
\label{I.1.1.13}
Para que $X=\Spec(A)$ seja irredutível \sref[0]{0.2.1.1}, é necessário e suficiente que o anel $A/\nilrad$ seja um domínio integral (ou, equivalentemente, que o ideal $\nilrad$ seja primo).
\end{proposition}

\begin{proof}
Em virtude do Corolário \sref{I.1.1.12}, podemos limitar-nos ao caso em que $\nilrad=0$.
Se $X$ é redutível, existem dois subconjuntos fechados próprios $Y_1$ e $Y_2$ de $X$ tais que $X=Y_1\cup Y_2$, donde $\mathfrak{j}(X)=\mathfrak{j}(Y_1)\cap\mathfrak{j}(Y_2)=0$, pois os ideais $\mathfrak{j}(Y_1)$ e $\mathfrak{j}(Y_2)$ são distintos de $(0)$ \sref{I.1.1.5};
portanto, $A$ não é um domínio integral.
Reciprocamente, se existem elementos $f\neq 0$, $g\neq 0$ de $A$ tais que $fg=0$, tem-se $V(f)\neq X$, $V(g)\neq X$ (pois a interseção de todos os ideais primos de $A$ é $(0)$), e $X=V(fg)=V(f)\cup V(g)$.
\end{proof}

\begin{corollary}[1.1.14]
\label{I.1.1.14}
\medskip\noindent
\begin{enumerate}
  \item[{\rm(i)}] Na correspondência bijetiva entre os subconjuntos fechados de $X=\Spec(A)$ e os ideais de $A$ iguais aos seus radicais, os subconjuntos fechados irredutíveis de $X$ correspondem aos ideais primos de $A$.
    Em particular, as componentes irredutíveis de $X$ correspondem aos ideais primos minimais de $A$.
  \item[{\rm(ii)}] A aplicação $x\mapsto\overline{\{x\}}$ estabelece uma correspondência bijetiva entre $X$ e o conjunto dos subconjuntos fechados irredutíveis de $X$ (\emph{isto é,} todo subconjunto fechado irredutível de $X$ admite exatamente um ponto genérico).
\end{enumerate}
\end{corollary}

\begin{proof}
(i) resulta imediatamente de \sref{I.1.1.13} e \sref{I.1.1.11};
e, para demonstrar (ii), podemos, em virtude de \sref{I.1.1.11}, limitar-nos ao caso em que $X$ é irredutível;
então, segundo a Proposição~\sref{I.1.1.13}, existe um ideal primo $\nilrad$ de $A$ contido em todos os outros, que corresponde ao ponto genérico
\oldpage[I]{83}
de $X$;
além disso, $X$ admite no máximo um ponto genérico, pois é um espaço de Kolmogorov (\sref{I.1.1.8} e \sref[0]{0.2.1.3}).
\end{proof}

\begin{proposition}[1.1.15]
\label{I.1.1.15}
Se $\mathfrak{J}$ é um ideal de $A$ contido no radical de Jacobson $\mathfrak{R}(A)$, a única
vizinhança de $V(\mathfrak{J})$ em $X=\Spec(A)$ é todo o espaço $X$.
\end{proposition}

\begin{proof}
Por definição, todo ideal maximal de $A$ pertence a $V(\mathfrak{J})$.
Como todo ideal $\mathfrak{a}\neq A$ de $A$ está contido em um ideal maximal, tem-se $V(\mathfrak{a})\cap V(\mathfrak{J})\neq\emp$, donde resulta a proposição.
\end{proof}

\subsection{Propriedades funtoriais dos espectros primos dos anéis}
\label{subsection:I.1.2}

\begin{env}[1.2.1]
\label{I.1.2.1}
Sejam $A$, $A'$ dois anéis, e
\[
  \vphi:A'\to A
\]
um homomorfismo de anéis.
Para todo ideal primo $x=\mathfrak{j}_x\in\Spec(A)=X$, o anel $A'/\vphi^{-1}(\mathfrak{j}_x)$ é canonicamente isomorfo a um subanel de $A/\mathfrak{j}_x$ e, portanto, é um domínio integral; em outras palavras, $\vphi^{-1}(\mathfrak{j}_x)$ é um ideal primo de $A'$;
denotá-lo-emos por ${}^a\vphi(x)$, definindo assim uma aplicação
\[
  {}^a\vphi:X=\Spec(A)\to X'=\Spec(A')
\]
(também denotada por $\Spec(\vphi)$), que chamaremos de aplicação \emph{associada} ao homomorfismo $\vphi$.
Denotaremos por $\vphi^x$ o homomorfismo injetivo de $A'/\vphi^{-1}(\mathfrak{j}_x)$ em $A/\mathfrak{j}_x$ induzido por $\vphi$ ao passar aos quocientes, bem como seu prolongamento canônico a um monomorfismo de corpos
\[
  \vphi^x:\kres({}^a\vphi(x))\to\kres(x);
\]
para todo $f'\in A'$, tem-se, por definição,
\[
  \label{I.1.2.1.1}
  \vphi^x(f'({}^a\vphi(x)))=(\vphi(f'))(x)\qquad(x\in X).
  \tag{1.2.1.1}
\]
\end{env}

\begin{proposition}[1.2.2]
\label{I.1.2.2}
\medskip\noindent
\begin{enumerate}
  \item[{\rm(i)}] Para todo subconjunto $E'$ de $A'$, tem-se
    \[
    \label{I.1.2.2.1}
      {}^a\vphi^{-1}(V(E'))=V(\vphi(E')),
      \tag{1.2.2.1}
    \]
    e, em particular, para todo $f'\in A'$,
    \[
    \label{I.1.2.2.2}
      {}^a\vphi^{-1}(D(f'))=D(\vphi(f')).
      \tag{1.2.2.2}
    \]
  \item[{\rm(ii)}] Para todo ideal $\mathfrak{a}$ de $A$, tem-se
    \[
    \label{I.1.2.2.3}
      \overline{{}^a\vphi(V(\mathfrak{a}))}=V(\vphi^{-1}(\mathfrak{a})).
      \tag{1.2.2.3}
    \]
\end{enumerate}
\end{proposition}

\begin{proof}
A relação ${}^a\vphi(x)\in V(E')$ equivale, por definição, a $E'\subset\vphi^{-1}(\mathfrak{j}_x)$, portanto a $\vphi(E')\subset\mathfrak{j}_x$, e finalmente a $x\in V(\vphi(E'))$, donde (i).
Para demonstrar (ii), podemos supor que $\mathfrak{a}$ é igual ao seu radical, pois $V(\rad(\mathfrak{a}))=V(\mathfrak{a})$ \sref{I.1.1.2}[(v)] e $\vphi^{-1}(\rad(\mathfrak{a}))=\rad(\vphi^{-1}(\mathfrak{a}))$;
se pusermos $Y=V(\mathfrak{a})$ e $\mathfrak{a}'=\mathfrak{j}({}^a\vphi(Y))$, tem-se $\overline{{}^a\vphi(Y)}=V(\mathfrak{a}')$ (\sref{I.1.1.4}[(ii)]).
A relação $f'\in\mathfrak{a}'$ equivale, por definição, a $f'(x')=0$ para todo $x'\in{}^a\vphi(Y)$; logo, em virtude da fórmula~\sref{I.1.2.1.1}, equivale também a $\vphi(f')(x)=0$ para todo $x\in Y$, ou a $\vphi(f')\in\mathfrak{j}(Y)=\mathfrak{a}$, pois $\mathfrak{a}$ é igual ao seu radical;
donde (ii).
\end{proof}

\begin{corollary}[1.2.3]
\label{I.1.2.3}
A aplicação ${}^a\vphi$ é contínua.
\end{corollary}

Observemos que, se $A''$ é um terceiro anel e $\vphi'$ um homomorfismo $A''\to A'$, então se tem ${}^a(\vphi\circ\vphi')={}^a\vphi'\circ{}^a\vphi$;
esse resultado, juntamente com o Corolário~\sref{I.1.2.3}, significa que $\Spec(A)$ é um \emph{funtor contravariante} em $A$, da categoria dos anéis para a dos espaços topológicos.

\begin{corollary}[1.2.4]
\label{I.1.2.4}
\oldpage[I]{84}
Suponhamos que $\vphi$ seja tal que todo $f\in A$ possa ser escrito como $f=h\vphi(f')$, onde $h$ é invertível em $A$ (\emph{o que ocorre, em particular, quando $\vphi$ é \emph{sobrejetivo}}).
Então ${}^a\vphi$ é um homeomorfismo de $X$ sobre ${}^a\vphi(X)$.
\end{corollary}

\begin{proof}
Mostraremos que, para todo subconjunto $E\subset A$, existe um subconjunto $E'$ de $A'$ tal que $V(E)=V(\vphi(E'))$;
em virtude do axioma ($T_0$) \sref{I.1.1.8} e da fórmula \sref{I.1.2.2.1}, isso implicará, primeiro, que ${}^a\vphi$ é injetiva e, depois, novamente por \sref{I.1.2.2.1}, que ${}^a\vphi$ é um homeomorfismo.
Ora, para cada $f\in E$, basta escolher $f'\in A'$ tal que $h\vphi(f')=f$, com $h$ invertível em $A$;
o conjunto $E'$ desses elementos $f'$ é precisamente o que procurávamos.
\end{proof}

\begin{env}[1.2.5]
\label{I.1.2.5}
Em particular, quando $\vphi$ é o homomorfismo canônico de $A$ sobre um anel quociente $A/\mathfrak{a}$, reencontramos \sref{I.1.1.11}, e ${}^a\vphi$ é a \emph{injeção canônica} de $V(\mathfrak{a})$, identificado com $\Spec(A/\mathfrak{a})$, em $X=\Spec(A)$.
\end{env}

Outro caso particular de \sref{I.1.2.4}:

\begin{corollary}[1.2.6]
\label{I.1.2.6}
Se $S$ é um conjunto multiplicativo de $A$, o espectro $\Spec(S^{-1}A)$ identifica-se canonicamente, com sua topologia, com o subespaço de $X=\Spec(A)$ formado pelos $x$ tais que $\mathfrak{j}_x\cap S=\emp$.
\end{corollary}

\begin{proof}
Sabemos por \sref[0]{0.1.2.6} que os ideais primos de $S^{-1}A$ são os ideais $S^{-1}\mathfrak{j}_x$ tais que $\mathfrak{j}_x\cap S=\emp$, e que se tem $\mathfrak{j}_x=(i_A^S)^{-1}(S^{-1}\mathfrak{j}_x)$.
Basta então aplicar o Corolário~\sref{I.1.2.4} a $i_A^S$.
\end{proof}

\begin{corollary}[1.2.7]
\label{I.1.2.7}
Para que ${}^a\vphi(X)$ seja densa em $X'$, é necessário e suficiente que todo elemento do núcleo $\Ker\vphi$ seja nilpotente.
\end{corollary}

\begin{proof}
Aplicando a fórmula~\sref{I.1.2.2.3} ao ideal $\mathfrak{a}=(0)$, obtém-se $\overline{{}^a\vphi(X)}=V(\Ker\vphi)$; para que $V(\Ker\vphi)=X'$ se verifique, é necessário e suficiente que $\Ker\vphi$ esteja contido em todos os ideais primos de $A'$, ou, equivalentemente, no nilradical $\nilrad$ de $A'$.
\end{proof}

\subsection{Feixe associado a um módulo}
\label{subsection:I.1.3}
\label{I.1.3}

\begin{env}[1.3.1]
\label{I.1.3.1}
Sejam $A$ um anel comutativo, $M$ um $A$-módulo, $f$ um elemento de $A$ e $S_f$ o conjunto multiplicativo formado pelos $f^n$, onde $n\geq 0$.
Recordemos que pomos $A_f=S_f^{-1}A$, $M_f=S_f^{-1}M$.
Se $S_f'$ é o subconjunto multiplicativo saturado de $A$ formado pelos $g\in A$ que dividem um elemento de $S_f$, sabemos que $A_f$ e $M_f$ se identificam canonicamente com ${S_f'}^{-1}A$ e ${S_f'}^{-1}M$ \sref[0]{0.1.4.3}.
\end{env}

\begin{lemma}[1.3.2]
\label{I.1.3.2}
As condições seguintes são equivalentes:
\begin{enumerate}
  \item[{\rm(a)}] $g\in S_f'$;
  \item[{\rm(b)}] $S_g'\subset S_f'$;
  \item[{\rm(c)}] $f\in\rad(g)$;
  \item[{\rm(d)}] $\rad(f)\subset\rad(g)$;
  \item[{\rm(e)}] $V(g)\subset V(f)$;
  \item[{\rm(f)}] $D(f)\subset D(g)$.
\end{enumerate}
\end{lemma}

\begin{proof}
Isso resulta imediatamente das definições e de \sref{I.1.1.5}.
\end{proof}

\begin{env}[1.3.3]
\label{I.1.3.3}
Se $D(f)=D(g)$, o Lema~\sref{I.1.3.2}[b] mostra que $M_f=M_g$.
Mais geralmente, se $D(f)\supset D(g)$, então $S_f'\subset S_g'$, e sabemos \sref[0]{0.1.4.1} que existe um homomorfismo canônico e funtorial
\[
  \rho_{g,f}:M_f\to M_g,
\]
e, se $D(f)\supset D(g)\supset D(h)$, tem-se \sref[0]{0.1.4.4}
\[
  \rho_{h,g}\circ\rho_{g,f}=\rho_{h,f}.
  \tag{1.3.3.1}
  \label{I.1.3.3.1}
\]
\end{env}

\oldpage[I]{85}
Quando $f$ percorre os elementos de $A\setmin\mathfrak{j}_x$ (para um $x$ dado de $X=\Spec(A)$), os conjuntos $S_f'$ constituem uma família filtrante crescente de subconjuntos de $A-\mathfrak{j}_x$, pois, para dois elementos $f$ e $g$ de $A\setmin\mathfrak{j}_x$, $S_f'$ e $S_g'$ estão contidos em $S_{fg}'$;
como a união dos $S_f'$ para $f\in A\setmin\mathfrak{j}_x$ é $A-\mathfrak{j}_x$, concluímos \sref[0]{0.1.4.5} que o $A_x$-módulo $M_x$ se identifica canonicamente com o \emph{limite indutivo} $\varinjlim M_f$, relativamente à família de homomorfismos $(\rho_{g,f})$.
Denotaremos por
\[
  \rho_x^f:M_f\to M_x
\]
o homomorfismo canônico para $f\in A\setmin\mathfrak{j}_x$ (ou, equivalentemente, $x\in D(f)$).

\begin{definition}[1.3.4]
\label{I.1.3.4}
Chama-se feixe estrutural do espectro primo $X=\Spec(A)$ (resp. feixe associado ao $A$-módulo $M$), e denota-se por $\widetilde{A}$ ou $\sh{O}_X$ (resp. $\widetilde{M}$), o feixe de anéis (resp. o $\widetilde{A}$-módulo) associado ao pré-feixe $D(f)\mapsto A_f$ (resp. $D(f)\mapsto M_f$), definido sobre a base $\mathfrak{B}$ de $X$ formada pelos $D(f)$ para $f\in A$ (\sref{I.1.1.10}, \sref[0]{0.3.2.1} e \sref[0]{0.3.5.6}).
\end{definition}

Vimos \sref[0]{0.3.2.4} que o talo $\widetilde{A}_x$ (resp. $\widetilde{M}_x$) \emph{se identifica com o anel $A_x$} (resp. \emph{com o $A_x$-módulo $M_x$});
denotaremos por
\[
  \theta_f:A_f\to\Gamma(D(f),\widetilde{A})
\]
\[
  \text{(resp. }\theta_f:M_f\to\Gamma(D(f),\widetilde{M})\text{),}
\]
a aplicação canônica, de modo que, para todo $x\in D(f)$ e todo $\xi\in M_f$, tem-se
\[
  (\theta_f(\xi))_x=\rho_x^f(\xi).
  \tag{1.3.4.1}
  \label{I.1.3.4.1}
\]

\begin{proposition}[1.3.5]
\label{I.1.3.5}
$\widetilde{M}$ é um funtor covariante exato em $M$, da categoria dos $A$-módulos para a categoria dos $\widetilde{A}$-módulos.
\end{proposition}

\begin{proof}
Com efeito, sejam $M$, $N$ dois $A$-módulos e $u$ um homomorfismo $M\to N$;
para todo $f\in A$, a $u$ corresponde canonicamente um homomorfismo $u_f$ do $A_f$-módulo $M_f$ no $A_f$-módulo $N_f$, e o diagrama (para $D(g)\subset D(f)$)
\[
  \xymatrix{
    M_f\ar[r]^{u_f}\ar[d]_{\rho_{g,f}} & N_f\ar[d]^{\rho_{g,f}}\\
    M_g\ar[r]_{u_g} & N_g
  }
\]
é comutativo \sref[0]{0.1.4.1};
esses homomorfismos definem, portanto, um homomorfismo de $\widetilde{A}$-módulos $\widetilde{u}:\widetilde{M}\to\widetilde{N}$ \sref[0]{0.3.2.3}.
Além disso, para todo $x\in X$, $\widetilde{u}_x$ é o limite indutivo dos $u_f$ para $x\in D(f)$ ($f\in A$) e, por conseguinte \sref[0]{0.1.4.5}, se $\widetilde{M}_x$ e $\widetilde{N}_x$ forem identificados canonicamente com $M_x$ e $N_x$, respectivamente, então $\widetilde{u}_x$ se identifica com o homomorfismo $u_x$ induzido canonicamente por $u$.
Se $P$ é um terceiro $A$-módulo, $v$ um homomorfismo $N\to P$ e $w=v\circ u$, é imediato que $w_x=v_x\circ u_x$, logo $\widetilde{w}=\widetilde{v}\circ\widetilde{u}$.
Fica assim bem definido um \emph{funtor covariante (em $M$)} $\widetilde{M}$, da categoria dos $A$-módulos para a dos $\widetilde{A}$-módulos.
\emph{Esse funtor é exato}, pois, para todo $x\in X$, $M_x$ é um funtor exato em $M$ \sref[0]{0.1.3.2};
além disso, tem-se $\Supp(M)=\Supp(\widetilde{M})$, em virtude das definições (\sref[0]{0.1.7.1} e \sref[0]{0.3.1.6}).
\end{proof}

\oldpage[I]{86}
\begin{proposition}[1.3.6]
\label{I.1.3.6}
Para todo $f\in A$, o subconjunto aberto $D(f)\subset X$ identifica-se canonicamente com o espectro primo $\Spec(A_f)$, e o feixe $\widetilde{M_f}$ associado ao $A_f$-módulo $M_f$ identifica-se canonicamente com a restrição $\widetilde{M}|D(f)$.
\end{proposition}

\begin{proof}
A primeira afirmação é um caso particular de \sref{I.1.2.6}.
Além disso, se $g\in A$ é tal que $D(g)\subset D(f)$, então $M_g$ identifica-se canonicamente com o módulo de frações de $M_f$ cujos denominadores são as potências da imagem canônica de $g$ em $A_f$ \sref[0]{0.1.4.6}.
A identificação canônica de $\widetilde{M_f}$ com $\widetilde{M}|D(f)$ resulta então das definições.
\end{proof}

\begin{theorem}[1.3.7]
\label{I.1.3.7}
Para todo $A$-módulo $M$ e todo $f\in A$, o homomorfismo
\[
  \theta_f:M_f\to\Gamma(D(f),\widetilde{M})
\]
é bijetivo \emph{(em outras palavras, o pré-feixe $D(f)\mapsto M_f$ é um \emph{feixe})}.
Em particular, $M$ identifica-se com $\Gamma(X,\widetilde{M})$ por meio de $\theta_1$.
\end{theorem}

\begin{proof}
Observemos que, se $M=A$, então $\theta_f$ é um homomorfismo de anéis;
o Teorema~\sref{I.1.3.7} implicará, portanto, que, se os anéis $A_f$ e $\Gamma(D(f),\widetilde{A})$ forem identificados por meio de $\theta_f$, o homomorfismo $\theta_f:M_f\to\Gamma(D(f),\widetilde{M})$ será um isomorfismo de \emph{módulos}.

Mostremos primeiro que $\theta_f$ é \emph{injetivo}.
Com efeito, se $\xi\in M_f$ é tal que $\theta_f(\xi)=0$, isso significa que, para todo ideal primo $\mathfrak{p}$ de $A_f$, existe $h\not\in\mathfrak{p}$ tal que $h\xi=0$;
como o anulador de $\xi$ não está contido em nenhum ideal primo de $A_f$, ele é todo o anel $A_f$ e, portanto, $\xi=0$.

Resta demonstrar que $\theta_f$ é \emph{sobrejetivo};
podemos reduzir-nos ao caso em que $f=1$ (deduzindo-se então o caso geral ao ``localizar'', mediante \sref{I.1.3.6}).
Seja, pois, $s$ uma seção de $\widetilde{M}$ sobre $X$;
em virtude de \sref{I.1.3.4} e \sref{I.1.1.10}[ii], existe uma cobertura \emph{finita} $(D(f_i))_{i\in I}$ de $X$ ($f_i\in A$) tal que, para todo $i\in I$, a restrição $s_i=s|D(f_i)$ é da forma $\theta_{f_i}(\xi_i)$, onde $\xi_i\in M_{f_i}$.
Se $i$, $j$ são dois índices de $I$, ao exprimir que as restrições de $s_i$ e $s_j$ a $D(f_i)\cap D(f_j)=D(f_i f_j)$ são iguais, resulta, por definição de $\widetilde{M}$, que
\[
  \label{I.1.3.7.1}
  \rho_{f_i f_j,f_i}(\xi_i)=\rho_{f_i f_j,f_j}(\xi_j).
  \tag{1.3.7.1}
\]
Por definição, podemos escrever, para todo $i\in I$, $\xi_i=z_i/f_i^{n_i}$, onde $z_i\in M$, e, como $I$ é finito, multiplicando cada $z_i$ por uma potência de $f_i$, podemos supor que todos os $n_i$ são iguais a um mesmo $n$.
Então, por definição, \sref{I.1.3.7.1} significa que existe um inteiro $m_{ij}\geq 0$ tal que $(f_i f_j)^{m_{ij}}(f_j^n z_i-f_i^n z_j)=0$, e podemos supor, além disso, que todos os $m_{ij}$ são iguais a um mesmo inteiro $m$;
substituindo então os $z_i$ por $f_i^m z_i$, reduzimo-nos ao caso em que $m=0$, isto é, ao caso em que se tem
\[
  \label{I.1.3.7.2}
  f_j^n z_i=f_i^n z_j
  \tag{1.3.7.2}
\]
para quaisquer $i$, $j$.
Tem-se $D(f_i^n)=D(f_i)$ e, como os $D(f_i)$ formam uma cobertura de $X$, o ideal gerado pelos $f_i^n$ é $A$;
em outras palavras, existem elementos $g_i\in A$ tais que $\sum_i g_i f_i^n=1$.
Consideremos então o elemento $z=\sum_i g_i z_i$ de $M$;
por \sref{I.1.3.7.2}, tem-se $f_i^n z=\sum_j g_j f_i^n z_j=(\sum_j g_j f_j^n)z_i=z_i$, donde, por definição, $\xi_i=z/1$ em $M_{f_i}$.
Concluímos
\oldpage[I]{87}
que os $s_i$ são as restrições a $D(f_i)$ de $\theta_1(z)$, o que demonstra que $s=\theta_1(z)$ e conclui a demonstração.
\end{proof}

\begin{corollary}[1.3.8]
\label{I.1.3.8}
Sejam $M$ e $N$ dois $A$-módulos;
o homomorfismo canônico $u\mapsto\widetilde{u}$ de $\Hom_A(M,N)$ em $\Hom_{\widetilde{A}}(\widetilde{M},\widetilde{N})$ é bijetivo.
Em particular, as relações $M=0$ e $\widetilde{M}=0$ são equivalentes.
\end{corollary}

\begin{proof}
Consideremos o homomorfismo canônico
\[
  v\longmapsto\Gamma(v):
  \Hom_{\widetilde{A}}(\widetilde{M},\widetilde{N})
  \longrightarrow
  \Hom_{\Gamma(\widetilde{A})}(\Gamma(\widetilde{M}),\Gamma(\widetilde{N}));
\]
este último módulo identifica-se canonicamente com $\Hom_A(M,N)$, pelo Teorema~\sref{I.1.3.7}.
Resta verificar que $u\mapsto\widetilde{u}$ e $v\mapsto\Gamma(v)$ são recíprocos;
é evidente que $\Gamma(\widetilde{u})=u$ pela definição de $\widetilde{u}$;
por outro lado, se pusermos $u=\Gamma(v)$ para $v\in\Hom_{\widetilde{A}}(\widetilde{M},\widetilde{N})$, então a aplicação $w:\Gamma(D(f),\widetilde{M})\to\Gamma(D(f),\widetilde{N})$ induzida canonicamente por $v$ é tal que o diagrama
\[
  \xymatrix{
    M\ar[r]^u\ar[d]_{\rho_{f,1}} & N\ar[d]^{\rho_{f,1}}\\
    M_f\ar[r]_w & N_f
  }
\]
é comutativo;
portanto, necessariamente $w=u_f$ para todo $f\in A$ \sref[0]{0.1.2.4}, o que demonstra que $\widetilde{\Gamma(v)}=v$.
\end{proof}

\begin{corollary}[1.3.9]
\label{I.1.3.9}
\medskip\noindent
\begin{enumerate}
  \item[{\rm(i)}] Seja $u$ um homomorfismo de um $A$-módulo $M$ em um $A$-módulo $N$;
    então os feixes associados a $\Ker u$, $\Im u$ e $\Coker u$ são, respectivamente, $\Ker\widetilde{u}$, $\Im\widetilde{u}$ e $\Coker\widetilde{u}$.
    Em particular, para que $\widetilde{u}$ seja injetivo (resp. sobrejetivo, bijetivo), é necessário e suficiente que $u$ o seja.
  \item[{\rm(ii)}] Se $M$ é um limite indutivo (resp. uma soma direta) de uma família de $A$-módulos $(M_\lambda)$, então $\widetilde{M}$ é o limite indutivo (resp. a soma direta) da família $(\widetilde{M_\lambda})$, por meio de um isomorfismo canônico.
\end{enumerate}
\end{corollary}

\begin{proof}
\medskip\noindent
\begin{enumerate}
  \item[(i)] Basta aplicar o fato de que $\widetilde{M}$ é um funtor exato em $M$ \sref{I.1.3.5} às duas sequências exatas de $A$-módulos
    \[
      0\to\Ker u\to M\to\Im u\to 0,
    \]
    \[
      0\to\Im u\to N\to\Coker u\to 0.
    \]
    A segunda afirmação resulta então do Teorema~\sref{I.1.3.7}.
  \item[(ii)] Seja $(M_\lambda,g_{\mu\lambda})$ um sistema indutivo de $A$-módulos, de limite indutivo $M$, e seja $g_\lambda$ o homomorfismo canônico $M_\lambda\to M$.
    Como se tem $\widetilde{g_{\nu\mu}}\circ\widetilde{g_{\mu\lambda}}=\widetilde{g_{\nu\lambda}}$ e $\widetilde{g_\lambda}=\widetilde{g_\mu}\circ\widetilde{g_{\mu\lambda}}$ para $\lambda\leq\mu\leq\nu$, resulta que $(\widetilde{M_\lambda},\widetilde{g_{\mu\lambda}})$ é um sistema indutivo de feixes sobre $X$ e, se denotarmos por $h_\lambda$ o homomorfismo canônico $\widetilde{M_\lambda}\to\varinjlim\widetilde{M_\lambda}$, existe um único homomorfismo $v:\varinjlim\widetilde{M_\lambda}\to\widetilde{M}$ tal que $v\circ h_\lambda=\widetilde{g_\lambda}$.
    Para ver que $v$ é bijetivo, basta verificar, para todo $x\in X$, que $v_x$ é uma bijeção de $(\varinjlim\widetilde{M_\lambda})_x$ em $\widetilde{M}_x$;
    mas $\widetilde{M}_x=M_x$, e
    \[
      (\varinjlim\widetilde{M_\lambda})_x=\varinjlim(\widetilde{M_\lambda})_x
      =\varinjlim(M_\lambda)_x=M_x\quad\sref[0]{0.1.3.3}.
    \]
    Por outro lado, das definições resulta que $(\widetilde{g_\lambda})_x$ e
    $(h_\lambda)_x$ são ambos iguais à aplicação canônica de $(M_\lambda)_x$ em $M_x$;
    como $(\widetilde{g_\lambda})_x=v_x\circ(h_\lambda)_x$, $v_x$ é a identidade.
\oldpage[I]{88}
    Finalmente, se $M$ é a soma direta de dois $A$-módulos $N$ e $P$, é imediato que $\widetilde{M}=\widetilde{N}\oplus\widetilde{P}$;
    sendo toda soma direta o limite indutivo de somas diretas finitas, ficam demonstradas as afirmações de (ii).
\end{enumerate}
\end{proof}

Observemos que o Corolário~\sref{I.1.3.8} demonstra que os feixes isomorfos a feixes associados a $A$-módulos formam uma \emph{categoria abeliana} (T, I, 1.4).

Observemos também que do Corolário~\sref{I.1.3.9} resulta que, se $M$ é um $A$-módulo \emph{de tipo finito} (isto é, se existe um homomorfismo sobrejetivo $A^n\to M$), então existe um homomorfismo sobrejetivo $\widetilde{A^n}\to\widetilde{M}$ ou, em outras palavras, o $\widetilde{A}$-módulo $\widetilde{M}$ é \emph{gerado por uma família finita de seções sobre $X$} \sref[0]{0.5.1.1}, e reciprocamente.

\begin{env}[1.3.10]
\label{I.1.3.10}
Se $N$ é um submódulo de um $A$-módulo $M$, a injeção canônica $j:N\to M$ fornece, por \sref{I.1.3.9}, um homomorfismo injetivo $\widetilde{N}\to\widetilde{M}$, que permite identificar canonicamente $\widetilde{N}$ com um \emph{$\widetilde{A}$-submódulo} de $\widetilde{M}$;
suporemos sempre feita essa identificação.
Se $N$ e $P$ são dois submódulos de $M$, então se tem
\[
\label{I.1.3.10.1}
  (N+P)\supertilde=\widetilde{N}+\widetilde{P},
  \tag{1.3.10.1}
\]
\[
\label{I.1.3.10.2}
  (N\cap P)\supertilde=\widetilde{N}\cap\widetilde{P},
  \tag{1.3.10.2}
\]
pois $N+P$ e $N\cap P$ são, respectivamente, a imagem do homomorfismo canônico
$N\oplus P\to M$ e o núcleo do homomorfismo canônico $M\to(M/N)\oplus(M/P)$, e basta aplicar \sref{I.1.3.9}.

De \sref{I.1.3.10.1} e \sref{I.1.3.10.2} conclui-se que, se $\widetilde{N}=\widetilde{P}$, então $N=P$.
\end{env}

\begin{corollary}[1.3.11]
\label{I.1.3.11}
Na categoria dos feixes isomorfos a feixes associados a $A$-módulos, o funtor $\Gamma$ é exato.
\end{corollary}

\begin{proof}
Seja $\widetilde{M}\xrightarrow{\widetilde{u}}\widetilde{N}\xrightarrow{\widetilde{v}}\widetilde{P}$ uma sequência exata correspondente a dois homomorfismos $u:M\to N$ e $v:N\to P$ de $A$-módulos.
Se $Q=\Im u$ e $R=\Ker v$, tem-se $\widetilde{Q}=\Im\widetilde{u}=\Ker\widetilde{v}=\widetilde{R}$ (Corolário~\sref{I.1.3.9}), logo $Q=R$.
\end{proof}

\begin{corollary}[1.3.12]
\label{I.1.3.12}
Sejam $M$ e $N$ dois $A$-módulos.
\begin{enumerate}
  \item[{\rm(i)}] O feixe associado a $M\otimes_A N$ identifica-se canonicamente com $\widetilde{M}\otimes_{\widetilde{A}}\widetilde{N}$.
  \item[{\rm(ii)}] Se, além disso, $M$ admite uma apresentação finita, então o feixe associado a $\Hom_A(M,N)$ identifica-se canonicamente com $\shHom_{\widetilde{A}}(\widetilde{M},\widetilde{N})$.
\end{enumerate}
\end{corollary}

\begin{proof}
\medskip\noindent
\begin{enumerate}
  \item[(i)] O feixe $\sh{F}=\widetilde{M}\otimes_{\widetilde{A}}\widetilde{N}$ está associado ao pré-feixe
    \[
      U\mapsto\sh{F}(U)=\Gamma(U,\widetilde{M})\otimes_{\Gamma(U,\widetilde{A})}\Gamma(U,\widetilde{N}),
    \]
    onde $U$ percorre a base \sref{I.1.1.10}[i] de $X$ formada pelos $D(f)$, com $f\in A$.
    Sabemos que $\sh{F}(D(f))$ se identifica canonicamente com $M_f\otimes_{A_f}N_f$, por \sref{I.1.3.7} e \sref{I.1.3.6}.
    Por outro lado, sabemos que o $A_f$-módulo $M_f\otimes_{A_f}N_f$ é canonicamente
    isomorfo a $(M\otimes_A N)_f$ \sref[0]{0.1.3.4}, que, por sua vez, é canonicamente isomorfo a $\Gamma(D(f),(M\otimes_A N)\supertilde)$ (Teorema~\sref{I.1.3.7} e Proposição~\sref{I.1.3.6}).
    Além disso, verifica-se imediatamente que os isomorfismos canônicos
    \[
      \sh{F}(D(f))\isoto\Gamma(D(f),(M\otimes_A N)\supertilde)
    \]
\oldpage[I]{89}
    assim obtidos satisfazem as condições de compatibilidade com respeito aos operadores de restrição \sref[0]{0.1.4.2}, pelo que definem um isomorfismo canônico funtorial
    \[
      \widetilde{M}\otimes_{\widetilde{A}}\widetilde{N}\isoto(M\otimes_A N)\supertilde.
    \]
  \item[(ii)] O feixe $\sh{G}=\shHom_{\widetilde{A}}(\widetilde{M},\widetilde{N})$ está associado ao pré-feixe
    \[
      U\mapsto\sh{G}(U)=\Hom_{\widetilde{A}|U}(\widetilde{M}|U,\widetilde{N}|U),
    \]
    onde $U$ percorre a base de $X$ formada pelos $D(f)$.
    Sabemos que $\sh{G}(D(f))$ se identifica canonicamente com $\Hom_{A_f}(M_f,N_f)$ (Proposição~\sref{I.1.3.6} e
    Corolário~\sref{I.1.3.8}), que, em virtude das hipóteses sobre $M$, se identifica canonicamente com $(\Hom_A(M,N))_f$ \sref[0]{0.1.3.5}.
    Finalmente, $(\Hom_A(M,N))_f$ se identifica canonicamente com $\Gamma(D(f),(\Hom_A(M,N))\supertilde)$ (Proposição~\sref{I.1.3.6} e Teorema~\sref{I.1.3.7}), e os isomorfismos canônicos $\sh{G}(D(f))\isoto\Gamma(D(f),(\Hom_A(M,N))\supertilde)$ assim obtidos são compatíveis com os operadores de restrição \sref[0]{0.1.4.2};
    definem, portanto, um isomorfismo canônico $\shHom_{\widetilde{A}}(\widetilde{M},\widetilde{N})\isoto(\Hom_A(M,N))\supertilde$.
\end{enumerate}
\end{proof}

\begin{env}[1.3.13]
\label{I.1.3.13}
Seja agora $B$ uma $A$-álgebra (comutativa); isso pode ser interpretado dizendo que $B$ é um $A$-módulo e que foram dados um elemento $e\in B$ e um $A$-homomorfismo $\vphi:B\otimes_A B\to B$, de modo que: $1^\circ$ os diagramas
\[
  \xymatrix{
    B\otimes_A B\otimes_A B\ar[r]^{\vphi\otimes 1}\ar[d]_{1\otimes\vphi} &
    B\otimes_A B\ar[d]^\vphi & &
    B\otimes_A B\ar[rr]^\sigma\ar[rd]_\vphi & &
    B\otimes_A B\ar[dl]^\vphi\\
    B\otimes_A B\ar[r]_\vphi &
    B & & &
    B
  }
\]
(onde $\sigma$ é a simetria canônica) são comutativos; $2^\circ$ $\vphi(e\otimes x)=\vphi(x\otimes e)=x$.
Pelo Corolário~\sref{I.1.3.12}, o homomorfismo $\widetilde{\vphi}:\widetilde{B}\otimes_{\widetilde{A}}\widetilde{B}\to\widetilde{B}$ de $\widetilde{A}$-módulos satisfaz as condições análogas e define, portanto, uma estrutura de \emph{$\widetilde{A}$-álgebra} sobre $\widetilde{B}$.
Analogamente, dar um $B$-módulo $N$ equivale a dar um $A$-módulo $N$ e um $A$-homomorfismo $\psi:B\otimes_A N\to N$ tal que o diagrama
\[
  \xymatrix{
    B\otimes_A B\otimes_A N\ar[r]^{\vphi\otimes 1}\ar[d]_{1\otimes\psi} &
    B\otimes_A N\ar[d]^\psi\\
    B\otimes_A N\ar[r]_\psi &
    N
  }
\]
seja comutativo e $\psi(e\otimes n)=n$;
o homomorfismo $\widetilde{\psi}:\widetilde{B}\otimes_{\widetilde{A}}\widetilde{N}\to\widetilde{N}$ satisfaz a condição análoga e define, portanto, uma estrutura de \emph{$\widetilde{B}$-módulo} sobre $\widetilde{N}$.

Da mesma maneira, vê-se que, se $u:B\to B'$ (resp. $v:N\to N'$) é um homomorfismo de $A$-álgebras (resp. de $B$-módulos), então $\widetilde{u}$ (resp. $\widetilde{v}$) é um homomorfismo de $\widetilde{A}$-álgebras (resp. de $\widetilde{B}$-módulos), e $\Ker\widetilde{u}$ é um $\widetilde{B}$-ideal (resp. $\Ker\widetilde{v}$, $\Coker\widetilde{v}$ e $\Im\widetilde{v}$ são $\widetilde{B}$-módulos).
Se $N$ é um $B$-módulo, então $\widetilde{N}$ é um $\widetilde{B}$-módulo de tipo finito se, e somente se, $N$ é um $B$-módulo de tipo finito \sref[0]{0.5.2.3}.

\oldpage[I]{90}
Se $M$ e $N$ são dois $B$-módulos, então o $\widetilde{B}$-módulo $\widetilde{M}\otimes_{\widetilde{B}}\widetilde{N}$ identifica-se canonicamente com $(M\otimes_B N)\supertilde$;
do mesmo modo, $\shHom_{\widetilde{B}}(\widetilde{M},\widetilde{N})$ identifica-se canonicamente com $(\Hom_B(M,N))\supertilde$ quando $M$ admite uma apresentação finita;
as demonstrações são as mesmas que as do Corolário~\sref{I.1.3.12}.

Se $\mathfrak{J}$ é um ideal de $B$ e $N$ é um $B$-módulo, então se tem $(\mathfrak{J}N)\supertilde=\widetilde{\mathfrak{J}}\cdot\widetilde{N}$.

Finalmente, se $B$ é uma $A$-álgebra \emph{graduada} pelos $A$-submódulos $B_n$ ($n\in\bb{Z}$), então a $\widetilde{A}$-álgebra $\widetilde{B}$, soma direta dos $\widetilde{A}$-módulos $\widetilde{B_n}$ \sref{I.1.3.9}, está graduada por esses $\widetilde{A}$-submódulos, pois o axioma da graduação exprime que a imagem do homomorfismo $B_m\otimes_A B_n\to B$ está contida em $B_{m+n}$.
Analogamente, se $M$ é um $B$-módulo graduado pelos submódulos $M_n$, então $\widetilde{M}$ é um $\widetilde{B}$-módulo graduado pelos $\widetilde{M_n}$.
\end{env}

\begin{env}[1.3.14]
\label{I.1.3.14}
Se $B$ é uma $A$-álgebra e $M$ um submódulo de $B$, então a $\widetilde{A}$-subálgebra de $\widetilde{B}$ gerada por $\widetilde{M}$ \sref[0]{0.4.1.3} é a $\widetilde{A}$-subálgebra $\widetilde{C}$, onde se denota por $C$ a subálgebra de $B$ gerada por $M$.
Com efeito, $C$ é a soma dos submódulos de $B$ que são as imagens dos homomorfismos $\bigotimes_A^n M\to B$ ($n\geq 0$), e basta aplicar \sref{I.1.3.9} e \sref{I.1.3.12}.
\end{env}

\subsection{Feixes quase-coerentes sobre um espectro primo}
\label{subsection:I.1.4}

\begin{theorem}[1.4.1]
\label{I.1.4.1}
Sejam $X$ o espectro primo de um anel $A$, $V$ um subconjunto aberto quase-compacto de $X$ e $\sh{F}$ um $(\sh{O}_X|V)$-módulo.
As quatro condições seguintes são equivalentes.
\begin{enumerate}
  \item[{\rm(a)}] Existe um $A$-módulo $M$ tal que $\sh{F}$ seja isomorfo a $\widetilde{M}|V$.
  \item[{\rm(b)}] Existe uma cobertura aberta finita $(V_i)$ de $V$ por conjuntos da forma $D(f_i)$ ($f_i\in A$) contidos em $V$, tal que, para cada $i$, $\sh{F}|V_i$ seja isomorfo a um feixe da forma $\widetilde{M_i}$, onde $M_i$ é um $A_{f_i}$-módulo.
  \item[{\rm(c)}] O feixe $\sh{F}$ é quase-coerente \sref[0]{0.5.1.3}.
  \item[{\rm(d)}] Satisfazem-se as duas propriedades seguintes:
    \begin{enumerate}
      \item[{\rm(d1)}] Para todo $f\in A$ tal que $D(f)\subset V$, e para toda seção $s\in\Gamma(D(f),\sh{F})$, existe um inteiro $n\geq 0$ tal que $f^n s$ se prolonga a uma seção de $\sh{F}$ sobre $V$.
      \item[{\rm(d2)}] Para todo $f\in A$ tal que $D(f)\subset V$ e para toda seção $t\in\Gamma(V,\sh{F})$ tal que a restrição de $t$ a $D(f)$ é $0$, existe um inteiro $n\geq 0$ tal que $f^n t=0$.
    \end{enumerate}
\end{enumerate}
\end{theorem}
(No enunciado das condições (d1) e (d2), identificamos tacitamente $A$ e $\Gamma(\widetilde{A})$ por meio do Teorema~\sref{I.1.3.7}).

\begin{proof}
O fato de que (a) implica (b) resulta imediatamente da Proposição~\sref{I.1.3.6} e do fato de que os $D(f_i)$ formam uma base da topologia de $X$ \sref{I.1.1.10}.
Como todo $A$-módulo é isomorfo ao conúcleo de um homomorfismo da forma $A^{(I)}\to A^{(J)}$, \sref{I.1.3.9} demonstra que todo feixe associado a um $A$-módulo é quase-coerente;
portanto, (b) implica (c).
Reciprocamente, se $\sh{F}$ é quase-coerente, todo $x\in V$ possui uma vizinhança da forma $D(f)\subset V$ tal que $\sh{F}|D(f)$ é isomorfo ao conúcleo de um homomorfismo $\widetilde{A_f}^{(I)}\to\widetilde{A_f}^{(J)}$, e portanto também ao feixe $\widetilde{N}$ associado ao módulo $N$, conúcleo do homomorfismo correspondente
$A_f^{(I)}\to A_f^{(J)}$ (Corolários~\sref{I.1.3.8} e \sref{I.1.3.9});
como $V$ é quase-compacto, é claro que (c) implica (b).

\oldpage[I]{91}
Para demonstrar que (b) implica (d1) e (d2), suponhamos primeiro que $V=D(g)$ para certo $g\in A$, e que $\sh{F}$ é isomorfo ao feixe $\widetilde{N}$ associado a um $A_g$-módulo $N$;
substituindo $X$ por $V$ e $A$ por $A_g$ \sref{I.1.3.6}, podemos reduzir-nos ao caso em que $g=1$.
Então $\Gamma(D(f),\widetilde{N})$ e $N_f$ identificam-se canonicamente entre si (Proposição~\sref{I.1.3.6} e Teorema~\sref{I.1.3.7}), de modo que uma seção $s\in\Gamma(D(f),\widetilde{N})$ se identifica com um elemento da forma $z/f^n$, onde $z\in N$;
a seção $f^n s$ se identifica com o elemento $z/1$ de $N_f$ e, por conseguinte, é a restrição a $D(f)$ da seção de $\widetilde{N}$ sobre $X$ que se identifica com o elemento $z\in N$;
daí resulta (d1) neste caso.
Analogamente, $t\in\Gamma(X,\widetilde{N})$ se identifica com um elemento $z'\in N$, a restrição de $t$ a $D(f)$ se identifica com a imagem $z'/1$ de $z'$ em $N_f$, e dizer que essa imagem é zero significa que existe um $n\geq 0$ tal que $f^n z'=0$ em $N$, ou, equivalentemente, $f^n t=0$.

Para terminar de demonstrar que (b) implica (d1) e (d2), basta estabelecer o lema seguinte.
\begin{lemma}[1.4.1.1]
\label{I.1.4.1.1}
Suponhamos que $V$ seja a união finita de conjuntos da forma $D(g_i)$, e que todos os feixes $\sh{F}|D(g_i)$ e $\sh{F}|(D(g_i)\cap D(g_j))=\sh{F}|D(g_i g_j)$ satisfaçam \emph{(d1)} e \emph{(d2)};
então $\sh{F}$ possui as duas propriedades seguintes:
\begin{enumerate}
  \item[{\rm(d'1)}] Para todo $f\in A$ e para toda seção $s\in\Gamma(D(f)\cap V,\sh{F})$, existe um inteiro $n\geq 0$ tal que $f^n s$ se prolonga a uma seção de $\sh{F}$ sobre $V$.
  \item[{\rm(d'2)}] Para todo $f\in A$ e para toda seção $t\in\Gamma(V,\sh{F})$ tal que a restrição de $t$ a $D(f)\cap V$ é $0$, existe um inteiro $n\geq 0$ tal que $f^n t=0$.
\end{enumerate}
\end{lemma}

Demonstremos primeiro (d$'$2): como $D(f)\cap D(g_i)=D(fg_i)$, existe, para cada $i$, um inteiro $n_i$ tal que a restrição de $(fg_i)^{n_i}t$ a $D(g_i)$ é zero:
como a imagem de $g_i$ em $A_{g_i}$ é invertível, a restrição de $f^{n_i}t$ a $D(g_i)$ também é zero;
tomando como $n$ o maior dos $n_i$, fica demonstrado (d$'$2).

Para demonstrar (d$'$1), aplicamos (d1) ao feixe $\sh{F}|D(g_i)$: existem um inteiro $n_i\geq 0$ e uma seção $s_i'$ de $\sh{F}$ sobre $D(g_i)$ que prolonga a restrição de $(fg_i)^{n_i}s$ a $D(fg_i)$;
como a imagem de $g_i$ em $A_{g_i}$ é invertível, existe uma seção $s_i$ de $\sh{F}$ sobre $D(g_i)$ tal que $s_i'=g_i^{n_i}s_i$, e $s_i$ prolonga a restrição de $f^{n_i}s$ a $D(fg_i)$;
além disso, podemos supor que todos os $n_i$ são iguais a um mesmo inteiro $n$.
Por construção, a restrição de $s_i-s_j$ a $D(f)\cap D(g_i)\cap D(g_j)=D(fg_i g_j)$ é zero;
aplicando (d2) ao feixe $\sh{F}|D(g_i g_j)$, existe um inteiro $m_{ij}\geq 0$ tal que a restrição a $D(g_i g_j)$ de $(fg_i g_j)^{m_{ij}}(s_i-s_j)$ é zero;
como a imagem de $g_i g_j$ em $A_{g_i g_j}$ é invertível, a restrição de $f^{m_{ij}}(s_i-s_j)$ a $D(g_i g_j)$ é zero.
Podemos supor então que todos os $m_{ij}$ são iguais a um mesmo inteiro $m$, e portanto existe uma seção $s'\in\Gamma(V,\sh{F})$ que prolonga todas as seções $f^m s_i$;
por conseguinte, essa seção prolonga $f^{n+m}s$, e assim fica demonstrado (d$'$1).

Resta demonstrar que (d1) e (d2) implicam (a).
Demonstremos primeiro que (d1) e (d2) implicam que essas condições se satisfazem para todo feixe $\sh{F}|D(g)$, onde $g\in A$ é tal que $D(g)\subset V$.
Isso é evidente para (d1);
por outro lado, se $t\in\Gamma(D(g),\sh{F})$ é tal que sua restrição a $D(f)\subset D(g)$ é zero, existe, por (d1), um inteiro $m\geq 0$ tal que $g^m t$
\oldpage[I]{92}
se prolonga a uma seção $s$ de $\sh{F}$ sobre $V$;
aplicando (d2), vemos que existe um inteiro $n\geq 0$ tal que $f^n g^m t=0$, e, como a imagem de $g$ em $A_g$ é invertível, $f^n t=0$.

Assim, como $V$ é quase-compacto, o Lema~\sref{I.1.4.1.1} demonstra que se satisfazem as condições (d$'$1) e (d$'$2).
Consideremos então o $A$-módulo $M=\Gamma(V,\sh{F})$, e definamos um homomorfismo de $\widetilde{A}$-módulos $u:\widetilde{M}\to j_*(\sh{F})$, onde $j$ é a inclusão canônica $V\to X$.
Como os $D(f)$ formam uma base da topologia de $X$, basta, para cada $f\in A$, definir um homomorfismo $u_f:M_f\to\Gamma(D(f),j_*(\sh{F}))=\Gamma(D(f)\cap V,\sh{F})$, com as condições de compatibilidade usuais \sref[0]{0.3.2.5}.
Como a imagem canônica de $f$ em $A_f$ é invertível, o homomorfismo de restrição $M=\Gamma(V,\sh{F})\to\Gamma(D(f)\cap V,\sh{F})$ se fatora como $M\to M_f\xrightarrow{u_f}\Gamma(D(f)\cap V,\sh{F})$ \sref[0]{0.1.2.4}, e a verificação dessas condições de compatibilidade para $D(g)\subset D(f)$ é imediata.
Dito isso, demonstremos que a condição (d$'$1) (resp. (d$'$2)) implica que cada $u_f$ é sobrejetivo (resp. injetivo), o que demonstra que $u$ é \emph{bijetivo}, e por conseguinte que $\sh{F}$ é a restrição a $V$ de um $\widetilde{A}$-módulo isomorfo a $\widetilde{M}$.
Se $s\in\Gamma(D(f)\cap V,\sh{F})$, existe, por (d$'$1), um inteiro $n\geq 0$ tal que $f^n s$ se prolonga a uma seção $z\in M$;
tem-se então $u_f(z/f^n)=s$, de modo que $u_f$ é sobrejetivo.
Analogamente, se $z\in M$ é tal que $u_f(z/1)=0$, isso significa que a restrição a $D(f)\cap V$ da seção $z$ é zero;
segundo (d$'$2), existe um inteiro $n\geq 0$ tal que $f^n z=0$, portanto $z/1=0$ em $M_f$, e assim $u_f$ é injetivo.
\end{proof}

\begin{corollary}[1.4.2]
\label{I.1.4.2}
Todo feixe quase-coerente sobre um aberto quase-compacto de $X$ é induzido por um feixe quase-coerente sobre $X$.
\end{corollary}

\begin{corollary}[1.4.3]
\label{I.1.4.3}
Toda $\sh{O}_X$-álgebra quase-coerente sobre $X=\Spec(A)$ é isomorfa a uma $\sh{O}_X$-álgebra da forma $\widetilde{B}$, onde $B$ é uma álgebra sobre $A$;
todo $\widetilde{B}$-módulo quase-coerente é isomorfo a um $\widetilde{B}$-módulo da forma $\widetilde{N}$, onde $N$ é um $B$-módulo.
\end{corollary}

\begin{proof}
De fato, uma $\sh{O}_X$-álgebra quase-coerente é um $\sh{O}_X$-módulo quase-coerente e, portanto, é da forma $\widetilde{B}$, onde $B$ é um $A$-módulo;
o fato de $B$ ser uma $A$-álgebra resulta da caracterização da estrutura de $\sh{O}_X$-álgebra por meio do homomorfismo $\widetilde{B}\otimes_{\widetilde{A}}\widetilde{B}\to\widetilde{B}$ de $\widetilde{A}$-módulos, bem como do Corolário~\sref{I.1.3.12}.
Se $\sh{G}$ é um $\widetilde{B}$-módulo quase-coerente, basta demonstrar que também é um $\widetilde{A}$-módulo quase-coerente e concluir em seguida da mesma maneira;
como a questão é local, podemos, restringindo-nos a um aberto de $X$ da forma $D(f)$, supor que $\sh{G}$ é o conúcleo de um homomorfismo $\widetilde{B}^{(I)}\to\widetilde{B}^{(J)}$ de $\widetilde{B}$-módulos (e \emph{a fortiori} de $\widetilde{A}$-módulos);
a proposição decorre então dos Corolários~\sref{I.1.3.8} e \sref{I.1.3.9}.
\end{proof}

\subsection{Feixes coerentes sobre um espectro primo}
\label{subsection:I.1.5}

\begin{theorem}[1.5.1]
\label{I.1.5.1}
Sejam $A$ um anel \emph{noetheriano}, $X=\Spec(A)$ seu espectro primo, $V$ um aberto de $X$ e $\sh{F}$ um $(\sh{O}_X|V)$-módulo.
As condições seguintes são equivalentes.
\begin{enumerate}
  \item[{\rm(a)}] $\sh{F}$ é coerente.
  \item[{\rm(b)}] $\sh{F}$ é de tipo finito e quase-coerente.
  \item[{\rm(c)}] Existe um $A$-módulo $M$ de tipo finito tal que $\sh{F}$ seja isomorfo ao feixe $\widetilde{M}|V$.
\end{enumerate}
\end{theorem}

\begin{proof}
\oldpage[I]{93}
(a) implica trivialmente (b).
Para ver que (b) implica (c), observemos primeiro que, como $V$ é quase-compacto \sref[0]{0.2.2.3}, $\sh{F}$ é isomorfo a um feixe $\widetilde{N}|V$, onde $N$ é um $A$-módulo \sref{I.1.4.1}.
Temos $N=\varinjlim M_\lambda$, onde $M_\lambda$ percorre o conjunto dos $A$-submódulos de $N$ de tipo finito; portanto, por \sref{I.1.3.9}, $\sh{F}=\widetilde{N}|V=\varinjlim\widetilde{M_\lambda}|V$;
mas, como $\sh{F}$ é de tipo finito e $V$ é quase-compacto, existe um índice $\lambda$ tal que $\sh{F}=\widetilde{M_\lambda}|V$ \sref[0]{0.5.2.3}.

Por fim, demonstremos que (c) implica (a).
É claro que $\sh{F}$ é então de tipo finito (\sref{I.1.3.6} e \sref{I.1.3.9});
além disso, como a questão é local, podemos nos limitar ao caso em que $V=D(f)$, $f\in A$.
Como $A_f$ é noetheriano, vemos que basta demonstrar que o núcleo de um homomorfismo $\widetilde{A^n}\to\widetilde{M}$, onde $M$ é um $A$-módulo, é de tipo finito.
Mas tal homomorfismo é da forma $\widetilde{u}$, onde $u$ é um homomorfismo $A^n\to M$ \sref{I.1.3.8}, e, se $P=\Ker u$, então $\widetilde{P}=\Ker\widetilde{u}$ \sref{I.1.3.9}.
Como $A$ é noetheriano, $P$ é de tipo finito, o que conclui a demonstração.
\end{proof}

\begin{corollary}[1.5.2]
\label{I.1.5.2}
Sob as hipóteses de \sref{I.1.5.1}, o feixe $\sh{O}_X$ é um feixe coerente de anéis.
\end{corollary}

\begin{corollary}[1.5.3]
\label{I.1.5.3}
Sob as hipóteses de \sref{I.1.5.1}, todo feixe coerente sobre um aberto de $X$ é induzido por um feixe coerente sobre $X$.
\end{corollary}

\begin{corollary}[1.5.4]
\label{I.1.5.4}
Sob as hipóteses de \sref{I.1.5.1}, todo $\sh{O}_X$-módulo quase-coerente $\sh{F}$ é o limite indutivo dos $\sh{O}_X$-submódulos coerentes de $\sh{F}$.
\end{corollary}

\begin{proof}
De fato, $\sh{F}=\widetilde{M}$, onde $M$ é um $A$-módulo, e $M$ é o limite indutivo de seus submódulos de tipo finito;
a demonstração conclui-se aplicando \sref{I.1.3.9} e \sref{I.1.5.1}.
\end{proof}

\subsection{Propriedades funtoriais dos feixes quase-coerentes sobre um espectro primo}
\label{subsection:I.1.6}

\begin{env}[1.6.1]
\label{I.1.6.1}
Sejam $A$, $A'$ dois anéis,
\[
  \vphi:A'\to A
\]
um homomorfismo, e
\[
  {}^a\vphi:X=\Spec(A)\to X'=\Spec(A')
\]
a aplicação contínua associada a $\vphi$ \sref{I.1.2.1}.
Definiremos um \emph{homomorfismo canônico}
\[
  \widetilde{\vphi}:\sh{O}_{X'}\to{}^a\vphi_*(\sh{O}_X)
\]
de feixes de anéis.
Para todo $f'\in A'$, ponhamos $f=\vphi(f')$;
tem-se ${}^a\vphi^{-1}(D(f'))=D(f)$ \sref{I.1.2.2.2}.
Os anéis $\Gamma(D(f'),\widetilde{A'})$ e $\Gamma(D(f),\widetilde{A})$ identificam-se respectivamente com $A_{f'}'$ e $A_f$ (\sref{I.1.3.6} e \sref{I.1.3.7}).
O homomorfismo $\vphi$ define canonicamente um homomorfismo $\vphi_{f'}:A_{f'}'\to A_f$ \sref[0]{0.1.5.1}; em outras palavras, tem-se um homomorfismo de anéis
\[
  \Gamma(D(f'),\widetilde{A'})\to\Gamma({}^a\vphi^{-1}(D(f')),\widetilde{A})
  =\Gamma(D(f'),{}^a\vphi_*(\widetilde{A})).
\]
\oldpage[I]{94}
Além disso, esses homomorfismos satisfazem as condições usuais de compatibilidade: para $D(f')\supset D(g')$, o diagrama
\[
  \xymatrix{
    \Gamma(D(f'),\widetilde{A'})\ar[r]\ar[d] &
    \Gamma(D(f'),{}^a\vphi_*(\widetilde{A}))\ar[d]\\
    \Gamma(D(g'),\widetilde{A'})\ar[r] &
    \Gamma(D(g'),{}^a\vphi_*(\widetilde{A}))
  }
\]
é comutativo \sref[0]{0.1.5.1};
definimos assim um homomorfismo de $\sh{O}_{X'}$-álgebras, pois os $D(f')$ formam uma base da topologia de $X'$ \sref[0]{0.3.2.3}.
O par $\Phi=({}^a\vphi,\widetilde{\vphi})$ é, portanto, um \emph{morfismo} de espaços anelados
\[
  \Phi:(X,\sh{O}_X)\to(X',\sh{O}_{X'}),
\]
\sref[0]{0.4.1.1}.

Observe-se ainda que, se pusermos $x'={}^a\vphi(x)$, o homomorfismo $\widetilde{\vphi}_x^\sharp$ \sref[0]{0.3.7.1} não é outro senão o homomorfismo
\[
  \vphi_x:A_{x'}'\to A_x
\]
induzido canonicamente por $\vphi:A'\to A$ \sref[0]{0.1.5.1}.
De fato, todo $z'\in A_{x'}'$ pode ser escrito como $g'/f'$, onde $f'$, $g'$ pertencem a $A'$ e $f'\not\in\mathfrak{j}_{x'}$;
$D(f')$ é então uma vizinhança de $x'$ em $X'$, e o homomorfismo $\Gamma(D(f'),\widetilde{A'})\to\Gamma({}^a\vphi^{-1}(D(f')),\widetilde{A})$ induzido por $\widetilde{\vphi}$ não é outro senão $\vphi_{f'}$;
considerando a seção $s'\in\Gamma(D(f'),\widetilde{A'})$ correspondente a $g'/f'\in A_{f'}'$, obtém-se $\widetilde{\vphi}_x^\sharp(z')=\vphi(g')/\vphi(f')$ em $A_x$.
\end{env}

\begin{example}[1.6.2]
\label{I.1.6.2}
Sejam $S$ um conjunto multiplicativo de $A$ e $\vphi$ o homomorfismo canônico $A\to S^{-1}A$;
já vimos em \sref{I.1.2.6} que ${}^a\vphi$ é um \emph{homeomorfismo} de $Y=\Spec(S^{-1}A)$ sobre o subespaço de $X=\Spec(A)$ formado pelos $x$ tais que $\mathfrak{j}_x\cap S=\emp$.
Além disso, para todo $x$ desse subespaço, portanto da forma ${}^a\vphi(y)$ com $y\in Y$, o homomorfismo $\widetilde{\vphi}_y^\sharp:\sh{O}_x\to\sh{O}_y$ é
\emph{bijetivo} \sref[0]{0.1.2.6};
em outras palavras, $\sh{O}_Y$ identifica-se com o feixe induzido sobre $Y$ por $\sh{O}_X$.
\end{example}

\begin{proposition}[1.6.3]
\label{I.1.6.3}
Para todo $A$-módulo $M$, existe um isomorfismo canônico funtorial do
$\sh{O}_{X'}$-módulo $(M_{[\vphi]})\supertilde$ sobre a imagem direta
$\Phi_*(\widetilde{M})$.
\end{proposition}

\begin{proof}
Para abreviar, ponhamos $M'=M_{[\vphi]}$ e, para todo $f'\in A'$, ponhamos
também $f=\vphi(f')$.
Os módulos de seções $\Gamma(D(f'),\widetilde{M'})$ e
$\Gamma(D(f),\widetilde{M})$ identificam-se respectivamente com os módulos
$M_{f'}'$ e $M_f$ (sobre $A_{f'}'$ e $A_f$, respectivamente); além disso, o
$A_{f'}'$-módulo $(M_f)_{[\vphi_{f'}]}$ é canonicamente isomorfo a $M_{f'}'$
\sref[0]{0.1.5.2}.
Tem-se, portanto, um isomorfismo funtorial de
$\Gamma(D(f'),\widetilde{A'})$-módulos:
$\Gamma(D(f'),\widetilde{M'})\isoto
\Gamma({}^a\vphi^{-1}(D(f')),\widetilde{M})_{[\vphi_{f'}]}$, e esses
isomorfismos satisfazem as condições usuais de compatibilidade com as
restrições \sref[0]{0.1.5.6}, definindo assim o isomorfismo funtorial anunciado.
Mais precisamente, se $u:M_1\to M_2$ é um homomorfismo de $A$-módulos, ele
pode ser considerado como um homomorfismo
$(M_1)_{[\vphi]}\to(M_2)_{[\vphi]}$ de $A'$-módulos; se denotarmos esse
homomorfismo por $u_{[\vphi]}$, então $\Phi_*(\widetilde{u})$ identifica-se com
$(u_{[\vphi]})\supertilde$.
\end{proof}

Esta demonstração prova também que, para toda \emph{$A$-álgebra $B$}, o
isomorfismo canônico funtorial
\oldpage[I]{95}
$(B_{[\vphi]})\supertilde\isoto\Phi_*(\widetilde{B})$ é um isomorfismo de
\emph{$\sh{O}_{X'}$-álgebras}; se $M$ é um $B$-módulo, o isomorfismo canônico
funtorial
$(M_{[\vphi]})\supertilde\isoto\Phi_*(\widetilde{M})$ é um isomorfismo de
$\Phi_*(\widetilde{B})$-módulos.

\begin{corollary}[1.6.4]
\label{I.1.6.4}
O funtor imagem direta $\Phi_*$ é exato na categoria dos
$\sh{O}_X$-módulos quase-coerentes.
\end{corollary}

\begin{proof}
Com efeito, é claro que $M_{[\vphi]}$ é um funtor exato em $M$ e que
$\widetilde{M'}$ é um funtor exato em $M'$ \sref{I.1.3.5}.
\end{proof}

\begin{proposition}[1.6.5]
\label{I.1.6.5}
Sejam $N'$ um $A'$-módulo e $N$ o $A$-módulo
$N'\otimes_{A'}A_{[\vphi]}$; existe então um isomorfismo canônico funtorial
do $\sh{O}_X$-módulo $\Phi^*(\widetilde{N'})$ sobre $\widetilde{N}$.
\end{proposition}

\begin{proof}
Observemos primeiro que $j:z'\mapsto z'\otimes 1$ é um $A'$-homomorfismo de
$N'$ em $N_{[\vphi]}$: com efeito, por definição, para $f'\in A'$, tem-se
$(f' z')\otimes 1=z'\otimes\vphi(f')=\vphi(f')(z'\otimes 1)$.
Deduz-se de \sref{I.1.3.8} um homomorfismo
$\widetilde{j}:\widetilde{N'}\to(N_{[\vphi]})\supertilde$ de
$\sh{O}_{X'}$-módulos e, em virtude de \sref{I.1.6.3}, podemos considerar
que $\widetilde{j}$ aplica $\widetilde{N'}$ em $\Phi_*(\widetilde{N})$.
A esse homomorfismo $\widetilde{j}$ corresponde canonicamente um homomorfismo
$h=\widetilde{j}^\sharp$ de $\Phi^*(\widetilde{N'})$ em $\widetilde{N}$
\sref[0]{0.4.4.3}; veremos que, em cada talo, $h_x$ é \emph{bijetivo}.
Ponhamos $x'={}^a\vphi(x)$ e seja $f'\in A'$ tal que $x'\in D(f')$; seja
$f=\vphi(f')$.
O anel $\Gamma(D(f),\widetilde{A})$ identifica-se com $A_f$, e os módulos
$\Gamma(D(f),\widetilde{N})$ e $\Gamma(D(f'),\widetilde{N'})$ identificam-se
com $N_f$ e $N'_{f'}$, respectivamente; seja
$s'\in\Gamma(D(f'),\widetilde{N'})$ uma seção, identificada com
$n'/{f'}^p$ ($n'\in N'$), e seja $s$ sua imagem por $\widetilde{j}$ em
$\Gamma(D(f),\widetilde{N})$; $s$ identifica-se com $(n'\otimes 1)/f^p$.
Por outro lado, seja $t\in\Gamma(D(f),\widetilde{A})$ uma seção, identificada
com $g/f^q$ ($g\in A$); então, por definição,
$h_x(s'_{x'}\otimes t_x)=t_x\cdot s_x$ \sref[0]{0.4.4.3}.
Mas pode-se identificar canonicamente $N_f$ com
$N'_{f'}\otimes_{A'_{f'}}(A_f)_{[\vphi_{f'}]}$ \sref[0]{0.1.5.4}; $s$
corresponde então ao elemento $(n'/{f'}^p)\otimes 1$, e a seção
$y\mapsto t_y\cdot s_y$ corresponde a $(n'/{f'}^p)\otimes(g/f^q)$.
Os diagramas de compatibilidade de \sref[0]{0.1.5.6} mostram que $h_x$ é
precisamente o isomorfismo canônico
\[
  \label{I.1.6.5.1}
  N'_{x'}\otimes_{A'_{x'}}(A_x)_{[\vphi_x]}\isoto
  N_x=(N'\otimes_{A'}A_{[\vphi]})_x.
  \tag{1.6.5.1}
\]

Além disso, seja $v:N_1'\to N_2'$ um homomorfismo de $A'$-módulos; como
$\widetilde{v}_{x'}=v_{x'}$ para todo $x'\in X'$, deduz-se imediatamente do
que precede que $\Phi^*(\widetilde{v})$ se identifica canonicamente com
$(v\otimes 1)\supertilde$, o que conclui a demonstração de \sref{I.1.6.5}.
\end{proof}

Se $B'$ é uma $A'$-álgebra, o isomorfismo canônico de
$\Phi^*(\widetilde{B'})$ sobre
$(B'\otimes_{A'}A_{[\vphi]})\supertilde$ é um isomorfismo de
$\sh{O}_X$-álgebras; se, além disso, $N'$ é um $B'$-módulo, o isomorfismo
canônico de $\Phi^*(\widetilde{N'})$ sobre
$(N'\otimes_{A'}A_{[\vphi]})\supertilde$ é um isomorfismo de
$\Phi^*(\widetilde{B'})$-módulos.

\begin{corollary}[1.6.6]
\label{I.1.6.6}
As seções de $\Phi^*(\widetilde{N'})$ que são imagens canônicas das seções
$s'$, quando $s'$ percorre o $A'$-módulo $\Gamma(\widetilde{N'})$, geram o
$A$-módulo $\Gamma(\Phi^*(\widetilde{N'}))$.
\end{corollary}

\begin{proof}
Com efeito, essas imagens identificam-se com os elementos $z'\otimes 1$ de
$N$, quando $N'$ e $N$ se identificam respectivamente com
$\Gamma(\widetilde{N'})$ e $\Gamma(\widetilde{N})$ \sref{I.1.3.7}, e quando
$z'$ percorre $N'$.
\end{proof}

\begin{env}[1.6.7]
\label{I.1.6.7}
Na demonstração de \sref{I.1.6.5}, provamos de passagem que a aplicação
canônica (\textbf{0},~4.4.3.2)
$\rho:\widetilde{N'}\to\Phi_*(\Phi^*(\widetilde{N'}))$ é precisamente o
homomorfismo $\widetilde{j}$,
\oldpage[I]{96}
onde $j:N'\to N'\otimes_{A'}A_{[\vphi]}$ é o homomorfismo
$z'\mapsto z'\otimes 1$.
Do mesmo modo, a aplicação canônica (\textbf{0},~4.4.3.3)
$\sigma:\Phi^*(\Phi_*(\widetilde{M}))\to\widetilde{M}$ é precisamente
$\widetilde{p}$, onde
$p:M_{[\vphi]}\otimes_{A'}A_{[\vphi]}\to M$ é o homomorfismo canônico que
associa a cada produto tensorial $z\otimes a$ ($z\in M$, $a\in A$) o
elemento $a\cdot z$; isso decorre imediatamente das definições
(\sref[0]{0.3.7.1}, \sref[0]{0.4.4.3} e \sref{I.1.3.7}).

Conclui-se (\textbf{0},~\EGAUnavailableHyperref{0.4.4.3}{4.4.3} e
\EGAUnavailableHyperref{0.3.5.4.4}{3.5.4.4}) que, se
$v:N'\to M_{[\vphi]}$ é um $A'$-homomorfismo, então
$\widetilde{v}^\sharp=(v\otimes 1)\supertilde$.
\end{env}

\begin{env}[1.6.8]
\label{I.1.6.8}
Sejam $N_1'$ e $N_2'$ dois $A'$-módulos, e suponhamos que $N_1'$ admita uma
\emph{apresentação finita}; deduz-se então de \sref{I.1.6.7} e
\sref{I.1.3.12}[ii] que o homomorfismo canônico \sref[0]{0.4.4.6}
\[
  \Phi^*(\shHom_{\widetilde{A'}}(\widetilde{N_1'},\widetilde{N_2'}))
  \to\shHom_{\widetilde{A}}(\Phi^*(\widetilde{N_1'}),\Phi^*(\widetilde{N_2'}))
\]
é precisamente $\widetilde{\gamma}$, onde $\gamma$ denota o homomorfismo
canônico de $A$-módulos
$\Hom_{A'}(N_1',N_2')\otimes_{A'}A\to
\Hom_A(N_1'\otimes_{A'}A,N_2'\otimes_{A'}A)$.
\end{env}

\begin{env}[1.6.9]
\label{I.1.6.9}
Sejam $\mathfrak{J}'$ um ideal de $A'$ e $M$ um $A$-módulo; como, por
definição, $\widetilde{\mathfrak{J}'}\widetilde{M}$ é a imagem do
homomorfismo canônico
$\Phi^*(\widetilde{\mathfrak{J}'})\otimes_{\widetilde{A}}\widetilde{M}
\to\widetilde{M}$, deduz-se da Proposição~\sref{I.1.6.5} e do
Corolário~\sref{I.1.3.12}[i] que
$\widetilde{\mathfrak{J}'}\widetilde{M}$ se identifica canonicamente com
$(\mathfrak{J}' M)\supertilde$; em particular,
$\Phi^*(\widetilde{\mathfrak{J}'})\widetilde{A}$ identifica-se com
$(\mathfrak{J}' A)\supertilde$ e, levando em conta a exatidão à direita do
funtor $\Phi^*$, a $\widetilde{A}$-álgebra
$\Phi^*((A'/\mathfrak{J}')\supertilde)$ identifica-se com
$(A/\mathfrak{J}' A)\supertilde$.
\end{env}

\begin{env}[1.6.10]
\label{I.1.6.10}
Sejam $A''$ um terceiro anel, $\vphi'$ um homomorfismo $A''\to A'$, e
ponhamos $\vphi''=\vphi\circ\vphi'$.
Deduz-se imediatamente das definições que
${}^a\vphi''=({}^a\vphi')\circ({}^a\vphi)$ e
$\widetilde{\vphi''}=\widetilde{\vphi}\circ\widetilde{\vphi'}$
\sref[0]{0.1.5.7}.
Conclui-se que $\Phi''=\Phi'\circ\Phi$; em outras palavras,
$(\Spec(A),\widetilde{A})$ é um \emph{funtor} em $A$ da categoria dos anéis
para a dos espaços anelados.
\end{env}

\subsection{Caracterização dos morfismos de esquemas afins}
\label{subsection:I.1.7}

\begin{definition}[1.7.1]
\label{I.1.7.1}
Dizemos que um espaço anelado $(X,\sh{O}_X)$ é um \emph{esquema afim}
se ele é isomorfo a um espaço anelado da forma
$(\Spec(A),\widetilde{A})$, onde $A$ é um anel; dizemos então que
$\Gamma(X,\sh{O}_X)$, que se identifica canonicamente com o anel $A$
\sref{I.1.3.7}, é o anel do esquema afim $(X,\sh{O}_X)$, e o denotamos
por $A(X)$ quando não houver risco de confusão.
\end{definition}

Por abuso de linguagem, quando falarmos do \emph{esquema afim}
$\Spec(A)$, tratar-se-á sempre do espaço anelado
$(\Spec(A),\widetilde{A})$.

\begin{env}[1.7.2]
\label{I.1.7.2}
Sejam $A$, $B$ dois anéis, e sejam $(X,\sh{O}_X)$, $(Y,\sh{O}_Y)$ os
esquemas afins correspondentes aos espectros primos $X=\Spec(A)$,
$Y=\Spec(B)$.
Vimos em \sref{I.1.6.1} que a todo homomorfismo de anéis
$\vphi:B\to A$ corresponde um morfismo
$\Phi=({}^a\vphi,\widetilde{\vphi})=\Spec(\vphi):
(X,\sh{O}_X)\to(Y,\sh{O}_Y)$.
Observemos que $\vphi$ é inteiramente determinado por $\Phi$, pois, por
definição,
$
  \vphi=\Gamma(\widetilde{\vphi}):
  \Gamma(\widetilde{B})\to
  \Gamma({}^a\vphi_*(\widetilde{A}))=\Gamma(\widetilde{A})
$.
\end{env}

\begin{theorem}[1.7.3]
\label{I.1.7.3}
\footnote{\emph{[Trad.] Ver \hyperref[subsection:I.1.8]{(1.8)} e a nota de
rodapé correspondente.}}
Sejam $(X,\sh{O}_X)$ e $(Y,\sh{O}_Y)$ dois esquemas afins.
Para que um morfismo de espaços anelados
$(\psi,\theta):(X,\sh{O}_X)\to(Y,\sh{O}_Y)$ seja da forma
$({}^a\vphi,\widetilde{\vphi})$, onde $\vphi$ é um homomorfismo de anéis
$A(Y)\to A(X)$, é necessário e suficiente que, para todo $x\in X$,
$\theta_x^\sharp$ seja um homomorfismo local:
$\sh{O}_{\psi(x)}\to\sh{O}_x$.
\end{theorem}

\begin{proof}
\oldpage[I]{97}
Ponhamos $A=A(X)$, $B=A(Y)$.
A condição é necessária, pois vimos em \sref{I.1.6.1} que
$\widetilde{\vphi}_x^\sharp$ é o homomorfismo de
$B_{{}^a\vphi(x)}$ em $A_x$ induzido canonicamente por $\vphi$ e, por
definição, ${}^a\vphi(x)=\vphi^{-1}(\mathfrak{j}_x)$; logo, esse
homomorfismo é local.

Demonstremos agora que a condição é suficiente.
Por definição, $\theta$ é um homomorfismo
$\sh{O}_Y\to\psi_*(\sh{O}_X)$, do qual se deduz canonicamente um
homomorfismo de anéis
\[
  \vphi=\Gamma(\theta):B=\Gamma(Y,\sh{O}_Y)
  \to\Gamma(Y,\psi_*(\sh{O}_X))=\Gamma(X,\sh{O}_X)=A.
\]

A hipótese sobre $\theta_x^\sharp$ permite deduzir desse homomorfismo,
passando aos quocientes, um monomorfismo $\theta^x$ do corpo residual
$\kres(\psi(x))$ no corpo residual $\kres(x)$, tal que, para toda seção
$f\in\Gamma(Y,\sh{O}_Y)=B$, se tenha
$\theta^x(f(\psi(x)))=\vphi(f)(x)$.
A relação $f(\psi(x))=0$ é, portanto, equivalente a
$\vphi(f)(x)=0$, o que significa que
$\mathfrak{j}_{\psi(x)}=\mathfrak{j}_{{}^a\vphi(x)}$; equivalentemente,
$\psi(x)={}^a\vphi(x)$ para todo $x\in X$, ou ainda
$\psi={}^a\vphi$.
Sabemos também que o diagrama
\[
  \xymatrix{
    B=\Gamma(Y,\sh{O}_Y)\ar[r]^\vphi\ar[d] &
    \Gamma(X,\sh{O}_X)=A\ar[d]\\
    B_{\psi(x)}\ar[r]_{\theta_x^\sharp} &
    A_x
  }
\]
é comutativo \sref[0]{0.3.7.2}, o que significa que
$\theta_x^\sharp$ é igual ao homomorfismo
$\vphi_x:B_{\psi(x)}\to A_x$ induzido canonicamente por $\vphi$
\sref[0]{0.1.5.1}.
Como os homomorfismos $\theta_x^\sharp$ determinam por completo
$\theta^\sharp$, e portanto $\theta$ \sref[0]{0.3.7.1}, concluímos que
$\theta=\widetilde{\vphi}$, pela definição de $\widetilde{\vphi}$
\sref{I.1.6.1}.
\end{proof}

Diremos que um morfismo $(\psi,\theta)$ de espaços anelados que satisfaz
à condição de \sref{I.1.7.3} é um \emph{morfismo de esquemas afins}.

\begin{corollary}[1.7.4]
\label{I.1.7.4}
Se $(X,\sh{O}_X)$ e $(Y,\sh{O}_Y)$ são esquemas afins, existe um
isomorfismo canônico do conjunto de morfismos de esquemas afins
$\Hom((X,\sh{O}_X),(Y,\sh{O}_Y))$ sobre o conjunto de homomorfismos de
anéis de $B$ em $A$, onde $A=\Gamma(\sh{O}_X)$ e
$B=\Gamma(\sh{O}_Y)$.
\end{corollary}

Também se pode dizer que os funtores
$A\longmapsto(\Spec(A),\widetilde{A})$ e
$(X,\sh{O}_X)\longmapsto\Gamma(X,\sh{O}_X)$ definem uma
\emph{equivalência} entre a categoria dos anéis comutativos e a categoria
oposta da categoria dos esquemas afins (T, I, 1.2).

\begin{corollary}[1.7.5]
\label{I.1.7.5}
Se $\vphi:B\to A$ é sobrejetivo, o morfismo correspondente
$({}^a\vphi,\widetilde{\vphi})$ é um monomorfismo de espaços anelados
\emph{(cf.~\sref{I.4.1.7})}.
\end{corollary}

\begin{proof}
Com efeito, sabemos que ${}^a\vphi$ é injetivo \sref{I.1.2.5} e, como
$\vphi$ é sobrejetivo, para todo $x\in X$,
$\vphi_x^\sharp:B_{{}^a\vphi(x)}\to A_x$, que se deduz de $\vphi$
passando aos anéis de frações, é também sobrejetivo
\sref[0]{0.1.5.1}; daí a conclusão \sref[0]{0.4.1.1}.
\end{proof}

\subsection{Morfismos de espaços localmente anelados a esquemas afins}
\label{subsection:I.1.8}

\oldpage[II]{217}
Devido a uma observação de J.~Tate, os enunciados do Teorema~
\sref{I.1.7.3} e da Proposição~\sref{I.2.2.4} podem ser generalizados como
segue:\footnote{\emph{[Trad.] A seção seguinte (I.1.8) foi acrescentada
nas erratas de EGA~II; daí a mudança temporária na numeração das páginas,
que remete a EGA~II.}}

\begin{proposition}[1.8.1]
\label{I.1.8.1}
Sejam $(S,\sh{O}_S)$ um esquema afim e $(X,\sh{O}_X)$ um espaço
localmente anelado.
Existe então uma bijeção canônica do conjunto de homomorfismos de anéis
\oldpage[II]{218}
$\Gamma(S,\sh{O}_S)\to\Gamma(X,\sh{O}_X)$ sobre o conjunto de morfismos de
espaços anelados $(\psi,\theta):(X,\sh{O}_X)\to(S,\sh{O}_S)$ tais que,
para todo $x\in X$, $\theta_x^\sharp$ seja um homomorfismo local
$\sh{O}_{\psi(x)}\to\sh{O}_x$.
\end{proposition}

\begin{proof}
Observemos antes de tudo que, se $(X,\sh{O}_X)$ e $(S,\sh{O}_S)$ são
dois espaços anelados quaisquer, um morfismo $(\psi,\theta)$ de
$(X,\sh{O}_X)$ em $(S,\sh{O}_S)$ define canonicamente um homomorfismo de
anéis
$\Gamma(\theta):\Gamma(S,\sh{O}_S)\to\Gamma(X,\sh{O}_X)$ e, portanto,
uma primeira aplicação
\[
  \label{I.1.8.1.1}
  \rho:\Hom((X,\sh{O}_X),(S,\sh{O}_S))\to\Hom(\Gamma(S,\sh{O}_S),\Gamma(X,\sh{O}_X)).
  \tag{1.8.1.1}
\]
Reciprocamente, sob as hipóteses do enunciado, ponhamos
$A=\Gamma(S,\sh{O}_S)$ e consideremos um homomorfismo de anéis
$\vphi:A\to\Gamma(X,\sh{O}_X)$.
Para todo $x\in X$, é claro que o conjunto dos $f\in A$ tais que
$\vphi(f)(x)=0$ é um \emph{ideal primo} de $A$, pois
$\sh{O}_x/\mathfrak{m}_x=\kres(x)$ é um corpo;
é, portanto, um elemento de $S=\Spec(A)$, que ainda denotaremos por
${}^a\vphi(x)$.
Além disso, para todo $f\in A$, tem-se, por definição
\sref[0]{0.5.5.2}, que ${}^a\vphi^{-1}(D(f))=X_f$, o que mostra que
${}^a\vphi$ é uma \emph{aplicação contínua} $X\to S$.
Definimos então um homomorfismo
\[
  \widetilde{\vphi}:\sh{O}_S\to{}^a\vphi_*(\sh{O}_X)
\]
de $\sh{O}_S$-módulos;
para todo $f\in A$, tem-se $\Gamma(D(f),\sh{O}_S)=A_f$
\sref{I.1.3.6};
para todo $s\in A$, fazemos corresponder ao elemento $s/f\in A_f$ o
elemento $(\vphi(s)|X_f)(\vphi(f)|X_f)^{-1}$ de
$\Gamma(X_f,\sh{O}_X)=\Gamma(D(f),{}^a\vphi_*(\sh{O}_X))$, e verifica-se
imediatamente (passando de $D(f)$ a $D(fg)$) que isso define de fato um
homomorfismo de $\sh{O}_S$-módulos e, portanto, um morfismo
$({}^a\vphi,\widetilde{\vphi})$ de espaços anelados.
Além disso, com as mesmas notações e pondo, para abreviar,
$y={}^a\vphi(x)$, vê-se imediatamente \sref[0]{0.3.7.1} que
$\widetilde{\vphi}_x^\sharp(s_y/f_y)=(\vphi(s)_x)(\vphi(f)_x)^{-1}$;
como a relação $s_y\in\mathfrak{m}_y$ equivale, por definição, a
$\vphi(s)_x\in\mathfrak{m}_x$, vemos que $\widetilde{\vphi}_x^\sharp$ é
um homomorfismo \emph{local} $\sh{O}_y\to\sh{O}_x$, e definimos assim uma
segunda aplicação
$\sigma:\Hom(\Gamma(S,\sh{O}_S),\Gamma(X,\sh{O}_X))\to\mathfrak{L}$,
onde $\mathfrak{L}$ é o conjunto dos morfismos
$(\psi,\theta):(X,\sh{O}_X)\to(S,\sh{O}_S)$ tais que
$\theta_x^\sharp$ seja local para todo $x\in X$.
Resta mostrar que $\sigma$ e $\rho$ (restrita a $\mathfrak{L}$) são
inversas uma da outra;
a definição de $\widetilde{\vphi}$ mostra imediatamente que
$\Gamma(\widetilde{\vphi})=\vphi$ e, por conseguinte, que
$\rho\circ\sigma$ é a identidade.
Para ver que $\sigma\circ\rho$ é a identidade, partamos de um morfismo
$(\psi,\theta)\in\mathfrak{L}$ e ponhamos $\vphi=\Gamma(\theta)$;
a hipótese sobre $\theta_x^\sharp$ permite deduzir desse homomorfismo,
passando aos quocientes, um monomorfismo
$\theta^x:\kres(\psi(x))\to\kres(x)$ tal que, para toda seção
$f\in A=\Gamma(S,\sh{O}_S)$, se tenha
$\theta^x(f(\psi(x)))=\vphi(f)(x)$;
a relação $f(\psi(x))=0$ é, portanto, equivalente a
$\vphi(f)(x)=0$, o que já mostra que ${}^a\vphi=\psi$.
Por outro lado, das definições decorre que o diagrama
\[
  \xymatrix{
    A\ar[r]^\vphi\ar[d] &
    \Gamma(X,\sh{O}_X)\ar[d]\\
    A_{\psi(x)}\ar[r]^{\theta_x^\sharp} &
    \sh{O}_x
  }
\]
é comutativo, e o mesmo vale para o diagrama análogo em que
$\theta_x^\sharp$ é substituído por $\widetilde{\vphi}_x^\sharp$; daí
$\widetilde{\vphi}_x^\sharp=\theta_x^\sharp$ \sref[0]{0.1.2.4} e, por
conseguinte, $\widetilde{\vphi}=\theta$.
\end{proof}

\begin{env}[1.8.2]
\label{I.1.8.2}
Quando $(X,\sh{O}_X)$ e $(Y,\sh{O}_Y)$ forem espaços
\emph{localmente} anelados, consideraremos os morfismos
$(\psi,\theta):(X,\sh{O}_X)\to(Y,\sh{O}_Y)$ tais que, para todo
$x\in X$, $\theta_x^\sharp$ seja um homomorfismo \emph{local}
$\sh{O}_{\psi(x)}\to\sh{O}_x$.
Daqui em diante, quando falarmos
\oldpage[II]{219}
de um \emph{morfismo de espaços localmente anelados}, tratar-se-á sempre
de um morfismo como o anterior;
com essa definição dos morfismos, é claro que os espaços localmente
anelados formam uma \emph{categoria};
para quaisquer dois objetos $X$ e $Y$ dessa categoria, $\Hom(X,Y)$
denota, portanto, o conjunto dos morfismos de espaços localmente anelados
de $X$ em $Y$ (o conjunto denotado por $\mathfrak{L}$ em
\sref{I.1.8.1});
quando considerarmos o conjunto dos \emph{morfismos de espaços anelados}
de $X$ em $Y$, denotá-lo-emos por $\Hom_\text{rs}(X,Y)$ para evitar toda
confusão.
A aplicação \sref{I.1.8.1.1} escreve-se então
\[
  \label{I.1.8.2.1}
  \rho:\Hom_\text{rs}(X,Y)\to\Hom(\Gamma(Y,\sh{O}_Y),\Gamma(X,\sh{O}_X))
  \tag{1.8.2.1}
\]
e sua restrição
\[
  \label{I.1.8.2.2}
  \rho':\Hom(X,Y)\to\Hom(\Gamma(Y,\sh{O}_Y),\Gamma(X,\sh{O}_X))
  \tag{1.8.2.2}
\]
é uma aplicação \emph{funtorial} em $X$ e $Y$ na categoria dos espaços
localmente anelados.
\end{env}

\begin{corollary}[1.8.3]
\label{I.1.8.3}
Seja $(Y,\sh{O}_Y)$ um espaço localmente anelado.
Para que $Y$ seja um esquema afim, é necessário e suficiente que, para
todo espaço localmente anelado $(X,\sh{O}_X)$, a aplicação
\sref{I.1.8.2.2} seja bijetiva.
\end{corollary}

\begin{proof}
A Proposição~\sref{I.1.8.1} mostra que a condição é necessária.
Reciprocamente, suponhamos que a condição seja satisfeita e ponhamos
$A=\Gamma(Y,\sh{O}_Y)$;
decorre da hipótese e de \sref{I.1.8.1} que os funtores
$X\mapsto\Hom(X,Y)$ e $X\mapsto\Hom(X,\Spec(A))$, da categoria dos
espaços localmente anelados à categoria dos conjuntos, são
\emph{isomorfos}.
Sabemos que isso implica a existência de um isomorfismo canônico
$Y\to\Spec(A)$ (cf.~\textbf{0},~8).\footnote{\emph{[Trad.] O texto
francês imprime aqui $X\to\Spec(A)$; o funtor representado e a conclusão
exigem $Y\to\Spec(A)$.}}
\end{proof}

\begin{env}[1.8.4]
\label{I.1.8.4}
Seja $S=\Spec(A)$ um esquema afim;
denotemos por $(S',A')$ o espaço anelado cujo espaço subjacente está
\emph{reduzido a um ponto} e cujo feixe estrutural $A'$ é o feixe
(necessariamente simples) sobre $S'$ definido pelo anel $A$.
Seja $\pi:S\to S'$ a única aplicação de $S$ em $S'$;
observemos, por outro lado, que, para todo aberto $U$ de $S$, existe uma
aplicação canônica
$\Gamma(S',A')=\Gamma(S,\sh{O}_S)\to\Gamma(U,\sh{O}_S)$, que define assim
um \emph{$\pi$-morfismo} $\iota:A'\to\sh{O}_S$ de feixes de anéis.
Definimos desse modo canonicamente um \emph{morfismo de espaços anelados
$i=(\pi,\iota):(S,\sh{O}_S)\to(S',A')$}.
Para todo $A$-módulo $M$, denotamos por $M'$ o feixe simples sobre $S'$
definido por $M$, que é evidentemente um $A'$-módulo.
É claro que $i_*(\widetilde{M})=M'$ \sref{I.1.3.7}.
\end{env}

\begin{lemma}[1.8.5]
\label{I.1.8.5}
Com as notações de \sref{I.1.8.4}, para todo $A$-módulo $M$, o
homomorfismo canônico funtorial de $\sh{O}_S$-módulos
\sref[0]{0.4.4.3.3}
\[
  \label{I.1.8.5.1}
  i^*(i_*(\widetilde{M}))\to\widetilde{M}
  \tag{1.8.5.1}
\]
é um isomorfismo.
\end{lemma}

\begin{proof}
Com efeito, os dois membros de \sref{I.1.8.5.1} são exatos à direita
(pois o funtor $M\mapsto i_*(\widetilde{M})$ é evidentemente exato) e
comutam com as somas diretas;
considerando $M$ como conúcleo de um homomorfismo
$A^{(I)}\to A^{(J)}$, reduzimo-nos a provar o lema no caso $M=A$, e
nesse caso ele é evidente.
\end{proof}

\begin{corollary}[1.8.6]
\label{I.1.8.6}
Sejam $(X,\sh{O}_X)$ um espaço anelado e $u:X\to S$ um morfismo de
espaços anelados.
\oldpage[II]{220}
Para todo $A$-módulo $M$, tem-se (com as notações de
\sref{I.1.8.4}) um isomorfismo canônico funtorial de
$\sh{O}_X$-módulos
\[
  \label{I.1.8.6.1}
  u^*(\widetilde{M})\isoto u^*(i^*(M')).
  \tag{1.8.6.1}
\]
\end{corollary}

\begin{corollary}[1.8.7]
\label{I.1.8.7}
Sob as hipóteses de \sref{I.1.8.6}, tem-se, para todo $A$-módulo $M$
e todo $\sh{O}_X$-módulo $\sh{F}$, um isomorfismo canônico, funtorial
em $M$ e $\sh{F}$,
\[
  \label{I.1.8.7.1}
  \Hom_{\sh{O}_S}(\widetilde{M},u_*(\sh{F}))\isoto\Hom_A(M,\Gamma(X,\sh{F})).
  \tag{1.8.7.1}
\]
\end{corollary}

\begin{proof}
Tem-se, em virtude de \sref[0]{0.4.4.3} e do Lema~\sref{I.1.8.5}, um
isomorfismo canônico de bifuntores
\[
  \Hom_{\sh{O}_S}(\widetilde{M},u_*(\sh{F}))\isoto\Hom_{A'}(M',i_*(u_*(\sh{F})))
\]
e é claro que o segundo membro é precisamente
$\Hom_A(M,\Gamma(X,\sh{F}))$.
Observemos que o homomorfismo canônico \sref{I.1.8.7.1} faz corresponder
a todo homomorfismo de $\sh{O}_S$-módulos
$h:\widetilde{M}\to u_*(\sh{F})$ (isto é, a todo $u$-morfismo
$\widetilde{M}\to\sh{F}$) o homomorfismo de $A$-módulos
$\Gamma(h):M\to\Gamma(S,u_*(\sh{F}))=\Gamma(X,\sh{F})$.
\end{proof}

\begin{env}[1.8.8]
\label{I.1.8.8}
Com as notações de \sref{I.1.8.4}, é claro \sref[0]{0.4.1.1} que todo
morfismo de espaços anelados $(\psi,\theta):X\to S'$ equivale ao dado de
um homomorfismo de anéis $A\to\Gamma(X,\sh{O}_X)$.
Podemos, portanto, interpretar a Proposição~\sref{I.1.8.1} como a
definição de uma bijeção canônica
$\Hom(X,S)\isoto\Hom(X,S')$ (entendendo-se que o segundo membro é o
conjunto dos morfismos de espaços anelados, pois, em geral, $A$ não é um
anel local).
Mais geralmente, se $X$ e $Y$ são espaços localmente anelados e se
$(Y',A')$ é o espaço anelado cujo espaço subjacente está reduzido a um
ponto e cujo feixe de anéis $A'$ é o feixe simples definido pelo anel
$\Gamma(Y,\sh{O}_Y)$, podemos interpretar \sref{I.1.8.2.1} como uma
aplicação
\[
  \label{I.1.8.8.1}
  \rho:\Hom_\text{rs}(X,Y)\to\Hom(X,Y').
  \tag{1.8.8.1}
\]
O resultado do Corolário~\sref{I.1.8.3} interpreta-se dizendo que os
esquemas afins se caracterizam, entre os espaços localmente anelados,
como aqueles para os quais a restrição de $\rho$ a
$\Hom(X,Y)$:
\[
  \label{I.1.8.8.2}
  \rho':\Hom(X,Y)\to\Hom(X,Y')
  \tag{1.8.8.2}
\]
é \emph{bijetiva} para \emph{todo} espaço localmente anelado $X$.
Em um capítulo posterior generalizaremos essa definição, o que permitirá
associar a \emph{todo} espaço anelado $Z$ (e não apenas a um espaço
anelado cujo espaço subjacente esteja reduzido a um ponto) um espaço
localmente anelado que ainda chamaremos $\Spec(Z)$;
esse será o ponto de partida de uma teoria ``relativa'' dos pré-esquemas
sobre espaços anelados quaisquer, estendendo os resultados do
Capítulo~I.
\end{env}

\begin{env}[1.8.9]
\label{I.1.8.9}
Podemos considerar que os pares $(X,\sh{F})$ formados por um espaço
localmente anelado $X$ e um $\sh{O}_X$-módulo $\sh{F}$ constituem uma
categoria, sendo um \emph{morfismo} dessa categoria um par $(u,h)$
formado por um morfismo de espaços localmente anelados
\oldpage[II]{221}
$u:X\to Y$ e um $u$-morfismo $h:\sh{G}\to\sh{F}$ de módulos;
esses morfismos (para $(X,\sh{F})$ e $(Y,\sh{G})$ fixados) formam um
conjunto que denotaremos por $\Hom((X,\sh{F}),(Y,\sh{G}))$;
a aplicação $(u,h)\mapsto(\rho'(u),\Gamma(h))$ é uma aplicação canônica
\[
  \label{I.1.8.9.1}
  \Hom((X,\sh{F}),(Y,\sh{G}))\to\Hom((\Gamma(Y,\sh{O}_Y),\Gamma(Y,\sh{G})),(\Gamma(X,\sh{O}_X),\Gamma(X,\sh{F})))
  \tag{1.8.9.1}
\]
\emph{funtorial} em $(X,\sh{F})$ e $(Y,\sh{G})$, sendo o segundo membro
o conjunto dos di-homomorfismos correspondentes aos anéis e módulos
considerados \sref[0]{0.1.0.2}.
\end{env}

\begin{corollary}[1.8.10]
\label{I.1.8.10}
Sejam $Y$ um espaço localmente anelado e $\sh{G}$ um
$\sh{O}_Y$-módulo. Para que $Y$ seja um esquema afim e $\sh{G}$ um
$\sh{O}_Y$-módulo quase-coerente, é necessário e suficiente que, para
todo par $(X,\sh{F})$ formado por um espaço localmente anelado $X$ e um
$\sh{O}_X$-módulo $\sh{F}$, a aplicação canônica
\sref{I.1.8.9.1} seja bijetiva.
\end{corollary}

Deixamos ao leitor o desenvolvimento da demonstração, inspirada na de
\sref{I.1.8.3} e na qual se utilizam \sref{I.1.8.1} e
\sref{I.1.8.7}.

\begin{remark}[1.8.11]
\label{I.1.8.11}
Os enunciados \sref{I.1.7.3}, \sref{I.1.7.4} e
\sref{I.2.2.4} são casos particulares de \sref{I.1.8.1}, assim como
a definição de \sref{I.1.6.1}; analogamente, \sref{I.2.2.5} decorre
de \sref{I.1.8.7}. O Corolário~\sref{I.1.8.7} implica também
\sref{I.1.6.3} (e, por conseguinte, \sref{I.1.6.4}) como caso
particular, pois, se $X$ é um esquema afim e
$\Gamma(X,\sh{F})=N$, os funtores
$M\mapsto\Hom_{\sh{O}_S}(\widetilde{M},u_*(\widetilde{N}))$ e
$M\mapsto\Hom_{\sh{O}_S}(\widetilde{M},(N_{[\vphi]})\supertilde)$
(onde $\vphi:A\to\Gamma(X,\sh{O}_X)$ corresponde a $u$) são
isomorfos, pelos Corolários~\sref{I.1.8.7} e \sref{I.1.3.8}.
Por fim, \sref{I.1.6.5} (e, por conseguinte, \sref{I.1.6.6}) decorre
de \sref{I.1.8.6} e do fato de que, para todo $f\in A$, os
$A_f$-módulos $N'\otimes_{A'}A_f$ e $(N'\otimes_{A'}A)_f$ (com as
notações de \sref{I.1.6.5}) são canonicamente isomorfos.
\end{remark}

\section{Pré-esquemas e morfismos de pré-esquemas}
\label{section:I.2}

\subsection{Definição dos pré-esquemas}
\label{subsection:I.2.1}

\begin{env}[2.1.1]
\label{I.2.1.1}
Dado um espaço anelado $(X,\sh{O}_X)$, diz-se que uma parte aberta
$V$ de $X$ é um \emph{aberto afim} se o espaço anelado
$(V,\sh{O}_X|V)$ é um esquema afim \sref{I.1.7.1}.
\end{env}

\begin{definition}[2.1.2]
\label{I.2.1.2}
Chama-se \emph{pré-esquema} um espaço anelado $(X,\sh{O}_X)$ tal que
todo ponto de $X$ admite uma vizinhança aberta afim.
\end{definition}

\begin{proposition}[2.1.3]
\label{I.2.1.3}
\oldpage[I]{98}
Se $(X,\sh{O}_X)$ é um pré-esquema, seus abertos afins formam uma
base da topologia de $X$.
\end{proposition}

\begin{proof}
Se $V$ é uma vizinhança aberta qualquer de $x\in X$, existe, por
hipótese, uma vizinhança aberta $W$ de $x$ tal que
$(W,\sh{O}_X|W)$ é um esquema afim; denotemos por $A$ o seu anel.
No espaço $W$, $V\cap W$ é uma vizinhança aberta de $x$; portanto,
existe $f\in A$ tal que $D(f)$ é uma vizinhança aberta de $x$ contida
em $V\cap W$ \sref{I.1.1.10}[i]. O espaço anelado
$(D(f),\sh{O}_X|D(f))$ é então um esquema afim cujo anel é isomorfo a
$A_f$ \sref{I.1.3.6}, donde resulta a proposição.
\end{proof}

\begin{proposition}[2.1.4]
\label{I.2.1.4}
O espaço subjacente de um pré-esquema é um espaço de Kolmogorov.
\end{proposition}

\begin{proof}
Se $x$ e $y$ são dois pontos distintos de um pré-esquema $X$, é
claro que existe uma vizinhança aberta de um desses pontos que não
contém o outro quando $x$ e $y$ não pertencem a um mesmo aberto afim;
quando pertencem a um mesmo aberto afim, isso resulta de
\sref{I.1.1.8}.
\end{proof}

\begin{proposition}[2.1.5]
\label{I.2.1.5}
Se $(X,\sh{O}_X)$ é um pré-esquema, toda parte fechada irredutível
de $X$ admite um único ponto genérico, e a aplicação
$x\mapsto\overline{\{x\}}$ é, portanto, uma bijeção de $X$ sobre o
conjunto de suas partes fechadas irredutíveis.
\end{proposition}

\begin{proof}
Se $Y$ é uma parte fechada irredutível de $X$, $y\in Y$, e $U$ é uma
vizinhança aberta afim de $y$ em $X$, então $U\cap Y$ é denso em $Y$
e também irredutível (\sref[0]{0.2.1.1} e \sref[0]{0.2.1.4});
portanto, pelo Corolário~\sref{I.1.1.14}, $U\cap Y$ é a aderência em
$U$ de um ponto $x$ e, por conseguinte, $Y\subset\overline{U}$ é a
aderência de $x$ em $X$. A unicidade do ponto genérico de $Y$ resulta
da Proposição~\sref{I.2.1.4} e de \sref[0]{0.2.1.3}.
\footnote{\emph{[Trad.] O texto francês imprime aqui ``o ponto
genérico de $X$''; a proposição e a frase precedente dizem respeito
ao ponto genérico de $Y$.}}
\end{proof}

\begin{env}[2.1.6]
\label{I.2.1.6}
Se $Y$ é uma parte fechada irredutível de $X$, e $y$ é seu ponto
genérico, o anel local $\sh{O}_y$ (também denotado por
$\sh{O}_{X/Y}$) chama-se \emph{anel local de $X$ ao longo de $Y$},
ou \emph{anel local de $Y$ em $X$}.

Se o próprio $X$ é irredutível e $x$ é seu ponto genérico, diz-se
ainda que $\sh{O}_x$ é o \emph{anel das funções racionais sobre $X$}
(cf.~\textsection\hyperref[section:I.7]{7}).
\end{env}

\begin{proposition}[2.1.7]
\label{I.2.1.7}
Se $(X,\sh{O}_X)$ é um pré-esquema, então, para toda parte aberta
$U$ de $X$, o espaço anelado $(U,\sh{O}_X|U)$ é um pré-esquema.
\end{proposition}

\begin{proof}
Isso resulta imediatamente da Definição~\sref{I.2.1.2} e da
Proposição~\sref{I.2.1.3}.
\end{proof}

Diz-se que $(U,\sh{O}_X|U)$ é o pré-esquema \emph{induzido} sobre
$U$ por $(X,\sh{O}_X)$, ou a \emph{restrição} de
$(X,\sh{O}_X)$ a $U$.

\begin{env}[2.1.8]
\label{I.2.1.8}
Diz-se que um pré-esquema $(X,\sh{O}_X)$ é \emph{irredutível}
(resp. \emph{conexo}) se o espaço subjacente $X$ é irredutível
(resp. conexo). Diz-se que um pré-esquema é \emph{integral} se é
\emph{irredutível e reduzido} (cf. \sref{I.5.1.4}). Diz-se que um
pré-esquema $(X,\sh{O}_X)$ é \emph{localmente integral} se todo
$x\in X$ admite uma vizinhança aberta $U$ tal que o pré-esquema
induzido sobre $U$ por $(X,\sh{O}_X)$ seja integral.
\end{env}

\subsection{Morfismos de pré-esquemas}
\label{subsection:I.2.2}
\label{I.2.2}

\begin{definition}[2.2.1]
\label{I.2.2.1}
Dados dois pré-esquemas $(X,\sh{O}_X)$ e $(Y,\sh{O}_Y)$, chama-se
\emph{morfismo (de pré-esquemas)} de $(X,\sh{O}_X)$ em
$(Y,\sh{O}_Y)$ todo morfismo de espaços anelados $(\psi,\theta)$ tal
que, para todo $x\in X$, $\theta_x^\sharp$ seja um homomorfismo local:
$\sh{O}_{\psi(x)}\longrightarrow\sh{O}_x$.
\end{definition}

\oldpage[I]{99}

Passando aos quocientes, a aplicação
$\theta_x^\sharp:\sh{O}_{\psi(x)}\to\sh{O}_x$ fornece, portanto, um
monomorfismo $\theta^x:\kres(\psi(x))\to\kres(x)$, que permite considerar
$\kres(x)$ como uma \emph{extensão} do corpo $\kres(\psi(x))$.

\begin{env}[2.2.2]
\label{I.2.2.2}
O composto $(\psi'',\theta'')$ de dois morfismos de pré-esquemas
$(\psi,\theta)$ e $(\psi',\theta')$ continua sendo um morfismo de
pré-esquemas, como resulta da fórmula
${\theta''}^\sharp=\theta^\sharp\circ\psi^*({\theta'}^\sharp)$
\sref[0]{0.3.5.5}. Conclui-se daí que os pré-esquemas formam uma
\emph{categoria}; seguindo a notação geral, escreveremos $\Hom(X,Y)$
para o conjunto dos morfismos de um pré-esquema $X$ em um pré-esquema $Y$.
\end{env}

\begin{example}[2.2.3]
\label{I.2.2.3}
Se $U$ é uma parte aberta de $X$, a injeção canônica
\sref[0]{0.4.1.2} do pré-esquema induzido
$(U,\sh{O}_X|U)$ em $(X,\sh{O}_X)$ é um morfismo de pré-esquemas;
é, além disso, um \emph{monomorfismo} de espaços anelados (e
\emph{a fortiori} um monomorfismo de pré-esquemas), como resulta
imediatamente de \sref[0]{0.4.1.1}.
\end{example}

\begin{proposition}[2.2.4]
\label{I.2.2.4}
\footnote{\emph{[Trad.] Ver \hyperref[subsection:I.1.8]{(1.8)} e a nota de
rodapé correspondente.}}
Sejam $(X,\sh{O}_X)$ um pré-esquema e $(S,\sh{O}_S)$ um esquema afim
de anel $A$. Existe uma correspondência biunívoca canônica entre os
morfismos do pré-esquema $(X,\sh{O}_X)$ no pré-esquema
$(S,\sh{O}_S)$ e os homomorfismos do anel $A$ no anel
$\Gamma(X,\sh{O}_X)$.
\end{proposition}

\begin{proof}
Observemos primeiro que, se $(X,\sh{O}_X)$ e $(Y,\sh{O}_Y)$ são dois
espaços anelados quaisquer, um morfismo $(\psi,\theta)$ de
$(X,\sh{O}_X)$ em $(Y,\sh{O}_Y)$ define canonicamente um homomorfismo
de anéis
$\Gamma(\theta):\Gamma(Y,\sh{O}_Y)\to
\Gamma(Y,\psi_*(\sh{O}_X))=\Gamma(X,\sh{O}_X)$.
No caso considerado, tudo se reduz a demonstrar que todo homomorfismo
$\vphi:A\to\Gamma(X,\sh{O}_X)$ é da forma $\Gamma(\theta)$ para um
único $\theta$.

Ora, por hipótese, existe um recobrimento $(V_\alpha)$ de $X$ por
abertos afins. Compondo $\vphi$ com o homomorfismo de restrição
$\Gamma(X,\sh{O}_X)\to\Gamma(V_\alpha,\sh{O}_X|V_\alpha)$, obtém-se um
homomorfismo
$\vphi_\alpha:A\to\Gamma(V_\alpha,\sh{O}_X|V_\alpha)$, que corresponde a
um único morfismo $(\psi_\alpha,\theta_\alpha)$ do pré-esquema
$(V_\alpha,\sh{O}_X|V_\alpha)$ no pré-esquema $(S,\sh{O}_S)$, em virtude
do Teorema~\sref{I.1.7.3}.
Além disso, para todo par de índices $(\alpha,\beta)$, todo ponto de
$V_\alpha\cap V_\beta$ admite uma vizinhança aberta afim $W$ contida em
$V_\alpha\cap V_\beta$ \sref{I.2.1.3}. É claro que, compondo
$\vphi_\alpha$ e $\vphi_\beta$ com os homomorfismos de restrição a $W$,
obtém-se o mesmo homomorfismo
$\Gamma(S,\sh{O}_S)\to\Gamma(W,\sh{O}_X|W)$; portanto, em virtude das
relações $(\theta_\alpha^\sharp)_x=(\vphi_\alpha)_x$ para todo
$x\in V_\alpha$ e todo $\alpha$ \sref{I.1.6.1}, as restrições a $W$ dos
morfismos $(\psi_\alpha,\theta_\alpha)$ e
$(\psi_\beta,\theta_\beta)$ coincidem.
Concluímos que existe um único morfismo de espaços anelados
$(\psi,\theta):(X,\sh{O}_X)\to(S,\sh{O}_S)$ cuja restrição a cada
$V_\alpha$ é $(\psi_\alpha,\theta_\alpha)$; é claro que se trata de um
morfismo de pré-esquemas e que $\Gamma(\theta)=\vphi$.

Seja $u:A\to\Gamma(X,\sh{O}_X)$ um homomorfismo de anéis, e seja
$v=(\psi,\theta)$ o morfismo correspondente
$(X,\sh{O}_X)\to(S,\sh{O}_S)$. Para todo $f\in A$, tem-se
\[
  \label{I.2.2.4.1}
  \psi^{-1}(D(f))=X_{u(f)}
  \tag{2.2.4.1}
\]
com as notações de \sref[0]{0.5.5.2} relativas ao feixe localmente livre
$\sh{O}_X$. Com efeito, basta verificar essa fórmula quando o próprio $X$
é afim, e então ela não é outra senão \sref{I.1.2.2.2}.
\end{proof}

\begin{proposition}[2.2.5]
\label{I.2.2.5}
Sob as hipóteses da Proposição~\sref{I.2.2.4}, seja
$\vphi:A\to\Gamma(X,\sh{O}_X)$ um homomorfismo de anéis, seja
$f:(X,\sh{O}_X)\to(S,\sh{O}_S)$ o morfismo de pré-esquemas correspondente,
seja $\sh{G}$ (resp. $\sh{F}$) um $\sh{O}_X$-módulo quase-coerente
(resp. um $\sh{O}_S$-módulo quase-coerente), e seja
$M=\Gamma(S,\sh{F})$. Existe então uma correspondência biunívoca canônica
\oldpage[I]{100}
entre os $f$-morfismos $\sh{F}\to\sh{G}$ \sref[0]{0.4.4.1} e os
$A$-homomorfismos $M\to(\Gamma(X,\sh{G}))_{[\vphi]}$.
\end{proposition}

\begin{proof}
Raciocinando como na Proposição~\sref{I.2.2.4}, reduzimo-nos imediatamente
ao caso em que $X$ é afim, e a proposição resulta então da
Proposição~\sref{I.1.6.3} e do Corolário~\sref{I.1.3.8}.
\end{proof}

\begin{env}[2.2.6]
\label{I.2.2.6}
Diz-se que um morfismo de pré-esquemas
$(\psi,\theta):(X,\sh{O}_X)\to(Y,\sh{O}_Y)$ é \emph{aberto}
(resp. \emph{fechado}) se, para toda parte aberta $U$ de $X$
(resp. toda parte fechada $F$ de $X$), $\psi(U)$ é aberta em $Y$
(resp. $\psi(F)$ é fechada em $Y$).
Diz-se que $(\psi,\theta)$ é \emph{dominante} se $\psi(X)$ é denso em $Y$,
e \emph{sobrejetivo} se $\psi$ é sobrejetiva.
Observemos que essas condições dependem apenas da aplicação contínua $\psi$.
\end{env}

\begin{proposition}[2.2.7]
\label{I.2.2.7}
Sejam
\[
  f=(\psi,\theta):(X,\sh{O}_X)\to(Y,\sh{O}_Y),
  \qquad
  g=(\psi',\theta'):(Y,\sh{O}_Y)\to(Z,\sh{O}_Z)
\]
dois morfismos de pré-esquemas.
\begin{enumerate}
  \item[\rm(i)] Se $f$ e $g$ são ambos abertos (resp. fechados, dominantes,
    sobrejetivos), o mesmo vale para $g\circ f$.
  \item[\rm(ii)] Se $f$ é sobrejetivo e $g\circ f$ é fechado, então $g$ é
    fechado.
  \item[\rm(iii)] Se $g\circ f$ é sobrejetivo, então $g$ é sobrejetivo.
\end{enumerate}
\end{proposition}

\begin{proof}
As afirmações (i) e (iii) são evidentes.
Ponhamos $g\circ f=(\psi'',\theta'')$.
Se $F$ é uma parte fechada de $Y$, então $\psi^{-1}(F)$ é fechada em $X$,
logo $\psi''(\psi^{-1}(F))$ é fechada em $Z$; mas, como $\psi$ é
sobrejetiva, $\psi(\psi^{-1}(F))=F$, portanto
$\psi''(\psi^{-1}(F))=\psi'(F)$, o que demonstra (ii).
\end{proof}

\begin{proposition}[2.2.8]
\label{I.2.2.8}
Sejam $f=(\psi,\theta)$ um morfismo
$(X,\sh{O}_X)\to(Y,\sh{O}_Y)$ e $(U_\alpha)$ um recobrimento aberto de $Y$.
Para que $f$ seja aberto (resp. fechado, sobrejetivo, dominante), é
necessário e suficiente que sua restrição a cada um dos pré-esquemas induzidos
$(\psi^{-1}(U_\alpha),\sh{O}_X|\psi^{-1}(U_\alpha))$, considerada como um
morfismo deste pré-esquema induzido no pré-esquema induzido
$(U_\alpha,\sh{O}_Y|U_\alpha)$, seja um morfismo aberto
(resp. fechado, sobrejetivo, dominante).
\end{proposition}

\begin{proof}
A proposição decorre imediatamente das definições, levando em conta o fato de
que uma parte $F$ de $Y$ é fechada (resp. aberta, densa) em $Y$ se, e somente
se, cada um dos conjuntos $F\cap U_\alpha$ é fechado (resp. aberto, denso) em
$U_\alpha$.
\end{proof}

\begin{env}[2.2.9]
\label{I.2.2.9}
Sejam $(X,\sh{O}_X)$ e $(Y,\sh{O}_Y)$ dois pré-esquemas; suponha-se que
$X$ e $Y$ tenham o mesmo número \emph{finito} de componentes irredutíveis
$X_i$ (resp. $Y_i$), $(1\leq i\leq n)$; seja $\xi_i$ (resp. $\eta_i$) o
ponto genérico de $X_i$ (resp. $Y_i$) \sref{I.2.1.5}.
Diz-se que um morfismo
\[
  f=(\psi,\theta):(X,\sh{O}_X)\to(Y,\sh{O}_Y)
\]
é \emph{birracional} se, para todo $i$, $\psi^{-1}(\eta_i)=\{\xi_i\}$ e
$\theta_{\xi_i}^{\sharp}:\sh{O}_{\eta_i}\to\sh{O}_{\xi_i}$ é um
\emph{isomorfismo}.
É claro que um morfismo birracional é \emph{dominante} \sref[0]{0.2.1.8},
logo sobrejetivo se também for fechado.
\end{env}

\begin{notation}[2.2.10]
\label{I.2.2.10}
Em todo o restante desta obra, e sempre que isso não acarretar risco de
confusão, \emph{suprimiremos} da notação de um pré-esquema
(resp. de um morfismo) o feixe estrutural
(resp. o morfismo de feixes estruturais).
Se $U$ é uma parte aberta do espaço subjacente de um pré-esquema $X$,
quando falarmos de $U$ como de um pré-esquema, tratar-se-á sempre do
pré-esquema induzido sobre $U$.
\end{notation}

\subsection{Colagem de pré-esquemas}
\label{subsection:I.2.3}

\begin{env}[2.3.1]
\label{I.2.3.1}
\oldpage[I]{101}
Resulta da Definição~\sref{I.2.1.2} que todo espaço anelado obtido por
\emph{colagem} de pré-esquemas \sref[0]{0.4.1.6} é ainda um pré-esquema.
Em particular, como todo pré-esquema admite, por definição, um recobrimento
formado por abertos afins, vê-se que todo pré-esquema pode ser obtido por
\emph{colagem de esquemas afins}.
\end{env}

\begin{example}[2.3.2]
\label{I.2.3.2}
Sejam $K$ um corpo, $B=K[s]$ e $C=K[t]$ dois anéis de polinômios em uma
indeterminada sobre $K$, e sejam $X_1=\Spec(B)$ e $X_2=\Spec(C)$, que são
dois esquemas afins isomorfos. Em $X_1$ (resp. $X_2$), seja $U_{12}$
(resp. $U_{21}$) o aberto afim $D(s)$ (resp. $D(t)$), cujo anel $B_s$
(resp. $C_t$) é formado pelas frações racionais da forma $f(s)/s^m$
(resp. $g(t)/t^n$), com $f\in B$ (resp. $g\in C$). Seja $u_{12}$ o
isomorfismo de pré-esquemas $U_{21}\to U_{12}$ correspondente
\sref{I.2.2.4} ao isomorfismo de $B_s$ em $C_t$ que associa a
$f(s)/s^m$ a fração racional $f(1/t)/(1/t^m)$. Podem-se colar $X_1$ e
$X_2$ ao longo de $U_{12}$ e $U_{21}$ por meio de $u_{12}$, pois
evidentemente não há condição de colagem. Mais adiante reencontraremos o
pré-esquema $X$ assim obtido como caso particular de um método geral de
construção \sref[II]{II.2.4.3}. Mostremos aqui apenas que $X$
\emph{não é um esquema afim}; isso decorrerá do fato de que o anel
$\Gamma(X,\sh{O}_X)$ é \emph{isomorfo} a $K$ e, portanto, tem espectro
reduzido a um ponto. Com efeito, uma seção de $\sh{O}_X$ sobre $X$ tem,
sobre $X_1$ (resp. $X_2$), identificado com um aberto afim de $X$, uma
restrição que é um polinômio $f(s)$ (resp. $g(t)$); das definições resulta
que se deve ter $g(t)=f(1/t)$, o que só é possível se $f=g\in K$.
\end{example}

\subsection{Esquemas locais}
\label{subsection:I.2.4}

\begin{env}[2.4.1]
\label{I.2.4.1}
Chama-se \emph{esquema local} um esquema afim cujo anel $A$ é local;
existe então em $X=\Spec(A)$ um único \emph{ponto fechado} $a$, e, para
todo outro ponto $b\in X$, tem-se
$a\in\overline{\{b\}}$ \sref{I.1.1.7}.
\end{env}

Para todo pré-esquema $Y$ e todo ponto $y\in Y$, o esquema local
$\Spec(\sh{O}_y)$ chama-se o \emph{esquema local de $Y$ no ponto $y$}.
Sejam $V$ um aberto afim de $Y$ contendo $y$ e $B$ o anel do esquema afim
$V$; $\sh{O}_y$ identifica-se canonicamente com $B_y$ \sref{I.1.3.4}, e o
homomorfismo canônico $B\to B_y$ corresponde, portanto, \sref{I.1.6.1} a
um morfismo de pré-esquemas $\Spec(\sh{O}_y)\to V$. Compondo esse
morfismo com a injeção canônica $V\to Y$, obtém-se, pois, um morfismo
$\Spec(\sh{O}_y)\to Y$, que é \emph{independente} do aberto afim $V$
(contendo $y$) escolhido: com efeito, se $V'$ é um segundo aberto afim
contendo $y$, existe um terceiro aberto afim $W$ contendo $y$ e tal que
$W\subset V\cap V'$ \sref{I.2.1.3}; pode-se, portanto, reduzir ao caso em
que $V\subset V'$ e, se $B'$ é o anel de $V'$, tudo se reduz a observar
que o diagrama
\phantomsection\label{I.2.4.1.diagram}
\[
  \xymatrix{
    B'\ar[rr]\ar[dr] & & B\ar[dl]\\
    & \sh{O}_y
  }
\]
é comutativo \sref[0]{0.1.5.1}. O morfismo
\[
  \Spec(\sh{O}_y)\to Y
\]
assim definido é dito \emph{canônico}.

\begin{proposition}[2.4.2]
\label{I.2.4.2}
\oldpage[I]{102}
Seja $(Y,\sh{O}_Y)$ um pré-esquema; para todo $y\in Y$, seja
$(\psi,\theta)$ o morfismo canônico
$(\Spec(\sh{O}_y),\widetilde{\sh{O}}_y)\to(Y,\sh{O}_Y)$. Então $\psi$ é
um homeomorfismo de $\Spec(\sh{O}_y)$ sobre o subespaço $S_y$ de $Y$
formado pelos $z$ tais que $y\in\overline{\{z\}}$ (isto é, pelas
\emph{generalizações} de $y$ \sref[0]{0.2.1.2}); além disso, se
$z=\psi(\mathfrak{p})$, então
$\theta_z^{\sharp}:\sh{O}_z\to(\sh{O}_y)_{\mathfrak{p}}$ é um
isomorfismo; $(\psi,\theta)$ é, portanto, um monomorfismo de espaços
anelados.
\end{proposition}

\begin{proof}
Como o único ponto fechado $a$ de $\Spec(\sh{O}_y)$ pertence ao fecho de
todo ponto desse espaço e $\psi(a)=y$, a imagem de
$\Spec(\sh{O}_y)$ pela aplicação contínua $\psi$ está contida em $S_y$.
Como $S_y$ está contido em todo aberto afim que contém $y$, podemos
reduzir ao caso em que $Y$ é um esquema afim; mas, nesse caso, a
proposição resulta de \sref{I.1.6.2}.
\end{proof}

\emph{Vê-se, portanto, \sref{I.2.1.5} que há uma correspondência
biunívoca entre $\Spec(\sh{O}_y)$ e o conjunto dos subconjuntos fechados
irredutíveis de $Y$ que contêm $y$.}

\begin{corollary}[2.4.3]
\label{I.2.4.3}
Para que $y\in Y$ seja o ponto genérico de uma componente irredutível de
$Y$, é necessário e suficiente que o único ideal primo do anel local
$\sh{O}_y$ seja seu ideal maximal (isto é, que $\sh{O}_y$ seja de
dimensão zero).
\end{corollary}

\begin{proposition}[2.4.4]
\label{I.2.4.4}
Sejam $(X,\sh{O}_X)$ um esquema local de anel $A$, $a$ seu único ponto
fechado, e $(Y,\sh{O}_Y)$ um pré-esquema. Todo morfismo
$u=(\psi,\theta):(X,\sh{O}_X)\to(Y,\sh{O}_Y)$ se fatora de maneira única
como
$X\to\Spec(\sh{O}_{\psi(a)})\to Y$, onde a segunda flecha designa o
morfismo canônico, e a primeira corresponde a um homomorfismo local
$\sh{O}_{\psi(a)}\to A$. Isso estabelece uma correspondência biunívoca
canônica entre o conjunto dos morfismos
$(X,\sh{O}_X)\to(Y,\sh{O}_Y)$ e o conjunto dos homomorfismos locais
$\sh{O}_y\to A$ $(y\in Y)$.
\end{proposition}

De fato, para todo $x\in X$, tem-se $a\in\overline{\{x\}}$, portanto
$\psi(a)\in\overline{\{\psi(x)\}}$, o que prova que $\psi(X)$ está
contido em todo aberto afim que contém $\psi(a)$. Podemos então nos
reduzir ao caso em que $(Y,\sh{O}_Y)$ é um esquema afim de anel $B$, e
nesse caso $u=({}^a\vphi,\widetilde{\vphi})$, onde
$\vphi\in\Hom(B,A)$ \sref{I.1.7.3}. Além disso, tem-se
$\vphi^{-1}(\mathfrak{j}_a)=\mathfrak{j}_{\psi(a)}$ e, por conseguinte,
a imagem por $\vphi$ de todo elemento de
$B\setmin\mathfrak{j}_{\psi(a)}$ é invertível no anel local $A$; a
fatoração do enunciado resulta, pois, da propriedade universal dos
anéis de frações \sref[0]{0.1.2.4}. Reciprocamente, a todo homomorfismo
local $\sh{O}_y\to A$ corresponde um único morfismo
$(\psi,\theta):X\to\Spec(\sh{O}_y)$ tal que $\psi(a)=y$
\sref{I.1.7.3}; compondo-o com o morfismo canônico
$\Spec(\sh{O}_y)\to Y$, obtemos um morfismo $X\to Y$, o que conclui a
demonstração da proposição.

\begin{env}[2.4.5]
\label{I.2.4.5}
Os esquemas afins cujo anel é um \emph{corpo} $K$ têm um espaço
subjacente reduzido a um ponto. Se $A$ é um anel local de ideal maximal
$\mathfrak{m}$, todo homomorfismo local $A\to K$ tem núcleo igual a
$\mathfrak{m}$ e, portanto, se fatora como
$A\to A/\mathfrak{m}\to K$, onde a segunda flecha é um monomorfismo. Os
morfismos $\Spec(K)\to\Spec(A)$ cujo único ponto é enviado ao ponto
fechado $\mathfrak{m}$ correspondem, portanto, biunivocamente aos
monomorfismos de corpos $A/\mathfrak{m}\to K$.
\footnote{\emph{[Trad.] A frase francesa impressa omite a condição
``cujo único ponto é enviado ao ponto fechado $\mathfrak{m}$''; sem
ela, a correspondência só vale depois de escolher o ponto imagem e usar
seu corpo residual.}}
\end{env}

Seja $(Y,\sh{O}_Y)$ um pré-esquema; para todo $y\in Y$ e todo ideal
$\mathfrak{a}_y$ de $\sh{O}_y$, o homomorfismo canônico
$\sh{O}_y\to\sh{O}_y/\mathfrak{a}_y$ define um morfismo
$\Spec(\sh{O}_y/\mathfrak{a}_y)\to\Spec(\sh{O}_y)$; compondo-o com o
morfismo canônico $\Spec(\sh{O}_y)\to Y$, obtemos um morfismo
$\Spec(\sh{O}_y/\mathfrak{a}_y)\to Y$, também chamado
\emph{canônico}. Para $\mathfrak{a}_y=\mathfrak{m}_y$, isto é,
$\sh{O}_y/\mathfrak{a}_y=\kres(y)$, a Proposição~\sref{I.2.4.4} fornece,
portanto:

\begin{corollary}[2.4.6]
\label{I.2.4.6}
\oldpage[I]{103}
Sejam $(X,\sh{O}_X)$ um esquema local cujo anel é um corpo $K$, $\xi$ o
único ponto de $X$, e $(Y,\sh{O}_Y)$ um pré-esquema. Todo morfismo
$u:(X,\sh{O}_X)\to(Y,\sh{O}_Y)$ se fatora de maneira única como
$X\to\Spec(\kres(\psi(\xi)))\to Y$, onde a segunda flecha designa o
morfismo canônico, e a primeira corresponde a um monomorfismo
$\kres(\psi(\xi))\to K$. Isso estabelece uma correspondência biunívoca
canônica entre o conjunto dos morfismos
$(X,\sh{O}_X)\to(Y,\sh{O}_Y)$ e o conjunto dos monomorfismos
$\kres(y)\to K$ $(y\in Y)$.
\end{corollary}

\begin{corollary}[2.4.7]
\label{I.2.4.7}
Para todo $y\in Y$, todo morfismo canônico
$\Spec(\sh{O}_y/\mathfrak{a}_y)\to Y$ é um monomorfismo de espaços
anelados.
\end{corollary}

\begin{proof}
Isso já foi visto quando $\mathfrak{a}_y=0$ \sref{I.2.4.2}, e basta
aplicar o Corolário~\sref{I.1.7.5}.
\end{proof}

\begin{remark}[2.4.8]
\label{I.2.4.8}
Sejam $X$ um esquema local e $a$ seu único ponto fechado.
Como todo aberto afim que contém $a$ é necessariamente o próprio $X$, todo $\sh{O}_X$-módulo \emph{invertível} \sref[0]{0.5.4.1} é necessariamente \emph{isomorfo a $\sh{O}_X$} (ou, como também se diz, é \emph{trivial}).
Essa propriedade não vale em geral para um esquema afim arbitrário $\Spec(A)$;
veremos no Capítulo~V que, se $A$ é um anel normal, ela vale quando $A$ é um \emph{domínio de fatoração única}.
\end{remark}

\subsection{Pré-esquemas sobre um pré-esquema}
\label{subsection:I.2.5}

\begin{definition}[2.5.1]
\label{I.2.5.1}
Dado um pré-esquema $S$, diz-se que os dados de um pré-esquema $X$ e de um morfismo de pré-esquemas $\vphi:X\to S$ definem um pré-esquema $X$ \emph{sobre o pré-esquema $S$}, ou um \emph{$S$-pré-esquema}; diz-se que $S$ é o \emph{pré-esquema de base} do $S$-pré-esquema $X$.
O morfismo $\vphi$ é chamado de \emph{morfismo estrutural} do $S$-pré-esquema $X$.
Quando $S$ é um esquema afim de anel $A$, diz-se também que $X$, munido de $\vphi$, é um pré-esquema sobre o anel $A$ (ou um $A$-pré-esquema).
\end{definition}

Resulta de \sref{I.2.2.4} que os dados de um pré-esquema sobre um anel $A$ equivalem aos de um pré-esquema $(X,\sh{O}_X)$ cujo feixe estrutural $\sh{O}_X$ é um feixe de $A$-\emph{álgebras}.
\emph{Portanto, todo pré-esquema pode sempre ser considerado, de maneira única, como um $\bb{Z}$-pré-esquema.}

Se $\vphi:X\to S$ é o morfismo estrutural de um $S$-pré-esquema $X$, diz-se que um ponto $x\in X$ está \emph{situado sobre um ponto $s\in S$} se $\vphi(x)=s$.
Diz-se que $X$ \emph{domina} $S$ se $\vphi$ é um morfismo dominante \sref{I.2.2.6}.

\begin{env}[2.5.2]
\label{I.2.5.2}
Sejam $X$ e $Y$ dois $S$-pré-esquemas; diz-se que um morfismo de pré-esquemas $u:X\to Y$ é um \emph{morfismo de pré-esquemas sobre $S$} (ou um $S$-\emph{morfismo}) se o diagrama
\phantomsection\label{I.2.5.2.diagram}
\[
  \xymatrix{
    X\ar[rr]^u\ar[dr] & & Y\ar[dl]\\
    & S &
  }
\]
(em que as setas oblíquas são os morfismos estruturais) é comutativo: isso implica que, para todo $s\in S$ e todo $x\in X$ situado sobre $s$, $u(x)$ também deve estar situado sobre $s$.
\end{env}

Esta definição mostra imediatamente que o composto de dois $S$-morfismos é um $S$-morfismo; portanto, os $S$-pré-esquemas formam uma \emph{categoria}.

Denota-se por $\Hom_S(X,Y)$ o conjunto dos $S$-morfismos de um $S$-pré-esquema $X$ em um $S$-pré-esquema $Y$; o morfismo identidade de um $S$-pré-esquema $X$ é denotado por $1_X$.

Quando $S$ é um esquema afim de anel $A$, diz-se também $A$-\emph{morfismo} em lugar de $S$-morfismo.

\begin{env}[2.5.3]
\label{I.2.5.3}
\oldpage[I]{104}
Se $X$ é um $S$-pré-esquema e $v:X'\to X$ um morfismo de pré-esquemas, o
morfismo composto $X'\xrightarrow{v}X\to S$ define $X'$ como um
$S$-pré-esquema; em particular, todo pré-esquema induzido sobre um aberto $U$ de
$X$ pode ser considerado como um $S$-pré-esquema mediante a injeção canônica.

Se $u:X\to Y$ é um $S$-morfismo de $S$-pré-esquemas, a restrição de $u$ a todo
pré-esquema induzido sobre um aberto $U$ de $X$ é, portanto, um $S$-morfismo
$U\to Y$. Reciprocamente, seja $(U_\alpha)$ uma cobertura aberta de $X$ e, para
cada $\alpha$, seja $u_\alpha:U_\alpha\to Y$ um $S$-morfismo; se, para todo par
de índices $(\alpha,\beta)$, as restrições de $u_\alpha$ e de $u_\beta$ a
$U_\alpha\cap U_\beta$ coincidem, existe um único $S$-morfismo $X\to Y$ cuja
restrição a cada $U_\alpha$ é $u_\alpha$.

Se $u:X\to Y$ é um $S$-morfismo tal que $u(X)\subset V$, onde $V$ é uma parte
aberta de $Y$, então $u$, considerado como morfismo de $X$ em $V$, é ainda um
$S$-morfismo.
\end{env}

\begin{env}[2.5.4]
\label{I.2.5.4}
Seja $S'\to S$ um morfismo de pré-esquemas; para todo $S'$-pré-esquema
$X$\footnote{\emph{[Trad.] A fonte francesa impressa diz «pour tout
S'-préschéma», omitindo o objeto X; o composto e a conclusão imediatamente
seguintes exigem que X seja ligado como o S'-pré-esquema considerado.}}, o
morfismo composto $X\to S'\to S$ define $X$ como um $S$-pré-esquema.
Reciprocamente, suponhamos que $S'$ seja o pré-esquema induzido sobre um aberto
de $S$; seja $X$ um $S$-pré-esquema e suponhamos que o morfismo estrutural
$f:X\to S$ seja tal que $f(X)\subset S'$; então $X$ pode ser considerado como
um $S'$-pré-esquema. Neste último caso, se $Y$ é um $S$-pré-esquema cujo
morfismo estrutural também aplica o espaço subjacente em $S'$, todo
$S$-morfismo de $X$ em $Y$ é também um $S'$-morfismo.
\end{env}

\begin{env}[2.5.5]
\label{I.2.5.5}
Se $X$ é um $S$-pré-esquema,\footnote{\emph{[Trad.] A fonte francesa impressa
diz «Si X est un S-morphisme»; o morfismo estrutural X→S e a definição de uma
seção exigem que X seja um S-pré-esquema.}} e $\vphi:X\to S$ é seu morfismo
estrutural, chama-se $S$-\emph{seção de $X$} todo $S$-morfismo de $S$ em $X$,
isto é, todo morfismo de pré-esquemas $\psi:S\to X$ tal que $\vphi\circ\psi$
seja a identidade de $S$. Denota-se por $\Gamma(X/S)$ o conjunto das
$S$-seções de $X$.
\end{env}

\section{Produto de pré-esquemas}
\label{section:I.3}

\subsection{Soma de pré-esquemas}
\label{subsection:I.3.1}
\label{I.3.1}

\phantomsection\label{I.3.1.1}
Seja $(X_\alpha)$ uma família qualquer de pré-esquemas; seja $X$ um espaço
topológico que é a \emph{soma} dos espaços subjacentes $X_\alpha$; então $X$
é a união de subespaços abertos $X'_\alpha$, dois a dois disjuntos, e, para
cada $\alpha$, há um homeomorfismo $\vphi_\alpha$ de $X_\alpha$ sobre
$X'_\alpha$. Se munirmos cada um dos $X'_\alpha$ do feixe
$(\vphi_\alpha)_*(\sh{O}_{X_\alpha})$, é claro que $X$ se torna um
pré-esquema, que se chama \emph{soma} da família de pré-esquemas $(X_\alpha)$
e se denota por $\bigsqcup_\alpha X_\alpha$. Se $Y$ é um pré-esquema, a
aplicação $f\mapsto(f\circ\vphi_\alpha)$ é uma \emph{bijeção} do conjunto
$\Hom(X,Y)$ sobre o conjunto produto $\prod_\alpha\Hom(X_\alpha,Y)$. Em
particular, se os $X_\alpha$ são $S$-pré-esquemas, com morfismos estruturais
$\psi_\alpha$, então $X$ é um $S$-pré-esquema para o único morfismo
$\psi:X\to S$ tal que $\psi\circ\vphi_\alpha=\psi_\alpha$ para todo
$\alpha$. A soma de dois pré-esquemas $X$, $Y$ denota-se por $X\sqcup Y$.
É imediato que, se $X=\Spec(A)$ e $Y=\Spec(B)$, então $X\sqcup Y$ se
identifica canonicamente com $\Spec(A\times B)$.

\subsection{Produto de pré-esquemas}
\label{subsection:I.3.2}

\begin{definition}[3.2.1]
\label{I.3.2.1}
Dados dois $S$-pré-esquemas $X$, $Y$, diz-se que uma tripla $(Z,p_1,p_2)$,
formada por um $S$-pré-esquema $Z$ e por dois $S$-morfismos $p_1:Z\to X$,
$p_2:Z\to Y$, é um produto dos
\oldpage[I]{105}
$S$-pré-esquemas $X$ e $Y$ se, para todo $S$-pré-esquema $T$, a aplicação
$f\mapsto(p_1\circ f,p_2\circ f)$ é uma bijeção do conjunto dos
$S$-morfismos de $T$ em $Z$ sobre o conjunto dos pares formados por um
$S$-morfismo $T\to X$ e um $S$-morfismo $T\to Y$ (em outras palavras, uma
bijeção
\[
  \Hom_S(T,Z)\isoto\Hom_S(T,X)\times\Hom_S(T,Y)).
\]
\end{definition}

Trata-se, portanto, da noção geral de \emph{produto} de dois objetos de uma
categoria, aplicada à categoria dos $S$-pré-esquemas (T, I, 1.1); em
particular, um produto de dois $S$-pré-esquemas é \emph{único} a menos de um
único $S$-isomorfismo. Em razão dessa unicidade, denota-se na maioria das
vezes um produto de dois $S$-pré-esquemas $X$, $Y$ por $X\times_S Y$ (ou
simplesmente por $X\times Y$, quando não houver possibilidade de confusão),
omitindo-se da notação os morfismos $p_1$, $p_2$ (chamados de
\emph{projeções canônicas} de $X\times_S Y$ em $X$ e $Y$, respectivamente).
Se $g:T\to X$, $h:T\to Y$ são dois $S$-morfismos, denota-se por $(g,h)_S$, ou
simplesmente por $(g,h)$, o $S$-morfismo $f:T\to X\times_S Y$ tal que
$p_1\circ f=g$, $p_2\circ f=h$. Se $X'$, $Y'$ são dois $S$-pré-esquemas,
$p'_1$, $p'_2$ são as projeções canônicas de $X'\times_S Y'$ (que se supõe
existir), e $u:X'\to X$, $v:Y'\to Y$ são dois $S$-morfismos, escreve-se
$u\times_S v$ (ou simplesmente $u\times v$) para o $S$-morfismo
$(u\circ p'_1,v\circ p'_2)_S$ de $X'\times_S Y'$ em $X\times_S Y$.

Quando $S$ é um esquema afim de anel $A$, substitui-se frequentemente $S$ por
$A$ nas notações precedentes.

\begin{proposition}[3.2.2]
\label{I.3.2.2}
Sejam $X$, $Y$, $S$ três esquemas afins, e sejam $B$, $C$, $A$ seus anéis
respectivos. Sejam $Z=\Spec(B\otimes_A C)$ e $p_1$, $p_2$ os $S$-morfismos
correspondentes, por \sref{I.2.2.4}, aos $A$-homomorfismos canônicos
$u:b\mapsto b\otimes1$ e $v:c\mapsto1\otimes c$ de $B$ e $C$ em
$B\otimes_A C$; então $(Z,p_1,p_2)$ é um produto de $X$ e $Y$.
\end{proposition}

\begin{proof}
Em virtude de \sref{I.2.2.4}, basta verificar que a associação, a cada
$A$-homomorfismo $f:B\otimes_A C\to L$ (onde $L$ é uma $A$-álgebra), do par
$(f\circ u,f\circ v)$ define uma bijeção
\[
  \Hom_A(B\otimes_A C,L)\isoto
  \Hom_A(B,L)\times\Hom_A(C,L),
\]\footnote{A notação $\Hom_A$ designa aqui o conjunto dos homomorfismos de
\emph{$A$-álgebras}.} o que resulta imediatamente das definições e da relação
$b\otimes c=(b\otimes1)(1\otimes c)$.
\end{proof}

\begin{corollary}[3.2.3]
\label{I.3.2.3}
Seja $T$ um esquema afim de anel $D$, e seja
$\alpha=({}^a\rho,\widetilde{\rho})$ (resp.
$\beta=({}^a\sigma,\widetilde{\sigma})$) um $S$-morfismo $T\to X$ (resp.
$T\to Y$), em que $\rho$ (resp. $\sigma$) é um $A$-homomorfismo de $B$
(resp. $C$) em $D$; então
$(\alpha,\beta)_S=({}^a\tau,\widetilde{\tau})$, em que $\tau$ é o homomorfismo
$B\otimes_A C\to D$ tal que $\tau(b\otimes c)=\rho(b)\sigma(c)$.
\end{corollary}

\begin{proposition}[3.2.4]
\label{I.3.2.4}
Seja $f:S'\to S$ um monomorfismo de pré-esquemas (T, I, 1.1), e sejam $X$,
$Y$ dois $S'$-pré-esquemas, considerados também como $S$-pré-esquemas por
meio de $f$. Todo produto dos $S$-pré-esquemas $X$, $Y$ é então um produto
dos $S'$-pré-esquemas $X$, $Y$, e reciprocamente.
\end{proposition}

\begin{proof}
Sejam $\vphi:X\to S'$, $\psi:Y\to S'$ os morfismos estruturais. Se $T$ é um
$S$-pré-esquema, e $u:T\to X$, $v:T\to Y$ são dois $S$-morfismos, tem-se por
definição $f\circ\vphi\circ u=f\circ\psi\circ v=\theta$, morfismo estrutural
de $T$; a hipótese sobre $f$ implica
$\vphi\circ u=\psi\circ v=\theta'$, e vê-se que se pode considerar $T$ como
um $S'$-pré-esquema de morfismo estrutural $\theta'$, e $u$, $v$ como
$S'$-morfismos. A conclusão da proposição resulta imediatamente, tendo em
conta \sref{I.3.2.1}.
\end{proof}

\begin{corollary}[3.2.5]
\label{I.3.2.5}
Sejam $X$, $Y$ dois $S$-pré-esquemas, $\vphi:X\to S$, $\psi:Y\to S$ seus
morfismos estruturais, e seja $S'$ uma parte aberta de $S$ tal que
$\vphi(X)\subset S'$, $\psi(Y)\subset S'$. Todo produto dos $S$-pré-esquemas
$X$, $Y$ é também um produto dos $S'$-pré-esquemas $X$, $Y$, e
reciprocamente.
\end{corollary}

\oldpage[I]{106}
Basta aplicar \sref{I.3.2.4} à injeção canônica $S'\to S$.

\begin{theorem}[3.2.6]
\label{I.3.2.6}
Dados dois $S$-pré-esquemas $X$, $Y$, existe um produto $X\times_S Y$.
\end{theorem}

Procederemos em várias etapas.

\begin{lemma}[3.2.6.1]
\label{I.3.2.6.1}
Seja $(Z,p,q)$ um produto de $X$ e $Y$, e sejam $U$, $V$ partes abertas de
$X$, $Y$, respectivamente. Se pusermos
$W=p^{-1}(U)\cap q^{-1}(V)$, a tripla formada por $W$ e pelas restrições de
$p$ e $q$ a $W$ \textup{(consideradas respectivamente como morfismos
$W\to U$, $W\to V$)} é um produto de $U$ e $V$.
\end{lemma}

Com efeito, se $T$ é um $S$-pré-esquema, podem-se identificar os
$S$-morfismos $T\to W$ com os $S$-morfismos $T\to Z$ que levam $T$ em $W$.
Se, então, $g:T\to U$, $h:T\to V$ são dois $S$-morfismos quaisquer, podem
ser considerados respectivamente como $S$-morfismos de $T$ em $X$ e $Y$ e,
por hipótese, existe um único $S$-morfismo $f:T\to Z$ tal que
$g=p\circ f$, $h=q\circ f$. Como $p(f(T))\subset U$, $q(f(T))\subset V$,
tem-se
\[
  f(T)\subset p^{-1}(U)\cap q^{-1}(V)=W,
\]
donde a afirmação.

\begin{lemma}[3.2.6.2]
\label{I.3.2.6.2}
Sejam $Z$ um $S$-pré-esquema, $p:Z\to X$, $q:Z\to Y$ dois
$S$-morfismos, $(U_\alpha)$ uma cobertura aberta de $X$, $(V_\lambda)$ uma
cobertura aberta de $Y$. Suponhamos que, para todo par $(\alpha,\lambda)$, o
$S$-pré-esquema
$W_{\alpha\lambda}=p^{-1}(U_\alpha)\cap q^{-1}(V_\lambda)$ e as restrições de
$p$ e $q$ a $W_{\alpha\lambda}$ constituam um produto de $U_\alpha$ e
$V_\lambda$. Então $(Z,p,q)$ é um produto de $X$ e $Y$.
\end{lemma}

Mostremos primeiro que, se $f_1$, $f_2$ são dois $S$-morfismos $T\to Z$, as
relações $p\circ f_1=p\circ f_2$ e $q\circ f_1=q\circ f_2$ implicam
$f_1=f_2$. Com efeito, $Z$ é reunião dos $W_{\alpha\lambda}$; portanto, os
$f_1^{-1}(W_{\alpha\lambda})$ formam uma cobertura aberta de $T$, e o mesmo
vale para os $f_2^{-1}(W_{\alpha\lambda})$. Além disso, tem-se
\[
\begin{aligned}
  f_1^{-1}(W_{\alpha\lambda})
  &=f_1^{-1}(p^{-1}(U_\alpha))\cap
    f_1^{-1}(q^{-1}(V_\lambda)) \\
  &=f_2^{-1}(p^{-1}(U_\alpha))\cap
    f_2^{-1}(q^{-1}(V_\lambda))
   =f_2^{-1}(W_{\alpha\lambda})
\end{aligned}
\]
por hipótese, e tudo se reduz a ver que as restrições de $f_1$ e $f_2$ a
$f_1^{-1}(W_{\alpha\lambda})=f_2^{-1}(W_{\alpha\lambda})$ são idênticas para
todo par de índices. Mas, como essas restrições podem ser consideradas
$S$-morfismos de $f_1^{-1}(W_{\alpha\lambda})$ em $W_{\alpha\lambda}$, nossa
afirmação resulta das hipóteses e da Definição~\sref{I.3.2.1}.

Suponhamos agora dados dois $S$-morfismos $g:T\to X$, $h:T\to Y$. Ponhamos
$T_{\alpha\lambda}=g^{-1}(U_\alpha)\cap h^{-1}(V_\lambda)$; os
$T_{\alpha\lambda}$ formam uma cobertura aberta de $T$. Por hipótese, existe
um $S$-morfismo $f_{\alpha\lambda}$ tal que $p\circ f_{\alpha\lambda}$ e
$q\circ f_{\alpha\lambda}$ sejam as restrições respectivas de $g$ e $h$ a
$T_{\alpha\lambda}$. Além disso, mostremos que as restrições de
$f_{\alpha\lambda}$ e $f_{\beta\mu}$ a
$T_{\alpha\lambda}\cap T_{\beta\mu}$ coincidem, o que concluirá a prova de
\sref{I.3.2.6.2}. Ora, as imagens de
$T_{\alpha\lambda}\cap T_{\beta\mu}$ por $f_{\alpha\lambda}$ e
$f_{\beta\mu}$ estão contidas em $W_{\alpha\lambda}\cap W_{\beta\mu}$ por
definição. Como
\[
  W_{\alpha\lambda}\cap W_{\beta\mu}
  =p^{-1}(U_\alpha\cap U_\beta)\cap
   q^{-1}(V_\lambda\cap V_\mu),
\]
resulta de \sref{I.3.2.6.1} que
$W_{\alpha\lambda}\cap W_{\beta\mu}$ e as restrições de $p$ e $q$ a esse
pré-esquema constituem um produto de $U_\alpha\cap U_\beta$ e
$V_\lambda\cap V_\mu$. Como $p\circ f_{\alpha\lambda}$ e
$p\circ f_{\beta\mu}$ coincidem em
$T_{\alpha\lambda}\cap T_{\beta\mu}$, e o mesmo ocorre com
$q\circ f_{\alpha\lambda}$ e $q\circ f_{\beta\mu}$, vê-se que
$f_{\alpha\lambda}$ e $f_{\beta\mu}$ coincidem em
$T_{\alpha\lambda}\cap T_{\beta\mu}$.\qed

\oldpage[I]{107}
\begin{lemma}[3.2.6.3]
\label{I.3.2.6.3}
Sejam $(U_\alpha)$ uma cobertura aberta de $X$, $(V_\lambda)$ uma cobertura
aberta de $Y$, e suponhamos que, para todo par $(\alpha,\lambda)$, exista um
produto de $U_\alpha$ e $V_\lambda$; então existe um produto de $X$ e $Y$.
\end{lemma}

Aplicando o Lema~\sref{I.3.2.6.1} aos abertos
$U_\alpha\cap U_\beta$ e $V_\lambda\cap V_\mu$, vê-se que existe um produto
dos $S$-pré-esquemas induzidos respectivamente por $X$ e $Y$ sobre esses
abertos; além disso, a unicidade do produto mostra que, se pusermos
$i=(\alpha,\lambda)$, $j=(\beta,\mu)$, existe um isomorfismo canônico
$h_{ij}$ (resp. $h_{ji}$) desse produto sobre um $S$-pré-esquema $W_{ij}$
(resp. $W_{ji}$) induzido por $U_\alpha\times_S V_\lambda$
(resp. $U_\beta\times_S V_\mu$) sobre um aberto;
$f_{ij}=h_{ij}\circ h_{ji}^{-1}$ é, portanto, um isomorfismo de $W_{ji}$
sobre $W_{ij}$. Além disso, para um terceiro par $k=(\gamma,\nu)$, tem-se
$f_{ik}=f_{ij}\circ f_{jk}$ em $W_{ki}\cap W_{kj}$, como resulta da
aplicação de \sref{I.3.2.6.1} aos abertos
$U_\alpha\cap U_\beta\cap U_\gamma$ e
$V_\lambda\cap V_\mu\cap V_\nu$, em $U_\beta$ e $V_\mu$,
respectivamente. Por conseguinte, existe um pré-esquema $Z$, uma cobertura
aberta $(Z_i)$ do espaço subjacente a $Z$ e, para cada $i$, um isomorfismo
$g_i$ do pré-esquema induzido $Z_i$ sobre o pré-esquema
$U_\alpha\times_S V_\lambda$, de modo que, para todo par $(i,j)$, se tenha
$f_{ij}=g_i\circ g_j^{-1}$ \sref{I.2.3.1}; além disso,
$g_i(Z_i\cap Z_j)=W_{ij}$. Se $p_i$, $q_i$, $\theta_i$ são as projeções e o
morfismo estrutural do $S$-pré-esquema $U_\alpha\times_S V_\lambda$, vê-se
imediatamente que $p_i\circ g_i=p_j\circ g_j$ em $Z_i\cap Z_j$, e o mesmo
vale para os outros dois morfismos. Podem-se, portanto, definir morfismos de
pré-esquemas $p:Z\to X$ (resp. $q:Z\to Y$, $\theta:Z\to S$) pela condição de
que $p$ (resp. $q$, $\theta$) coincida com $p_i\circ g_i$
(resp. $q_i\circ g_i$, $\theta_i\circ g_i$) em cada um dos $Z_i$; $Z$,
munido de $\theta$, é então um $S$-pré-esquema. Mostremos agora que
$Z'_i=p^{-1}(U_\alpha)\cap q^{-1}(V_\lambda)$ é igual a $Z_i$. Para todo
índice $j=(\beta,\mu)$, tem-se
$Z_j\cap Z'_i=g_j^{-1}(p_j^{-1}(U_\alpha)\cap q_j^{-1}(V_\lambda))$. Ora,
\[
  p_j^{-1}(U_\alpha)\cap q_j^{-1}(V_\lambda)
  =p_j^{-1}(U_\alpha\cap U_\beta)\cap
   q_j^{-1}(V_\lambda\cap V_\mu);
\]
em virtude de \sref{I.3.2.6.1}, as restrições de $p_j$ e $q_j$ a
$p_j^{-1}(U_\alpha)\cap q_j^{-1}(V_\lambda)$ definem nesse
$S$-pré-esquema uma estrutura de produto de $U_\alpha\cap U_\beta$ e
$V_\lambda\cap V_\mu$; mas a unicidade do produto implica então que
$p_j^{-1}(U_\alpha)\cap q_j^{-1}(V_\lambda)=W_{ji}$. Tem-se, por
conseguinte, $Z_j\cap Z'_i=Z_j\cap Z_i$ para todo $j$, donde $Z'_i=Z_i$.
Deduz-se então de \sref{I.3.2.6.2} que $(Z,p,q)$ é um produto de $X$ e $Y$.

\begin{lemma}[3.2.6.4]
\label{I.3.2.6.4}
Sejam $\vphi:X\to S$, $\psi:Y\to S$ os morfismos estruturais de $X$ e
$Y$, $(S_i)$ uma cobertura aberta de $S$, e ponhamos
$X_i=\vphi^{-1}(S_i)$, $Y_i=\psi^{-1}(S_i)$. Se cada um dos produtos
$X_i\times_S Y_i$ existe, então $X\times_S Y$ existe.
\end{lemma}

De acordo com \sref{I.3.2.6.3}, tudo se reduz a provar que os produtos
$X_i\times_S Y_j$ existem quaisquer que sejam $i$ e $j$. Ponhamos
$X_{ij}=X_i\cap X_j=\vphi^{-1}(S_i\cap S_j)$,
$Y_{ij}=Y_i\cap Y_j=\psi^{-1}(S_i\cap S_j)$; em virtude de
\sref{I.3.2.6.1}, o produto $Z_{ij}=X_{ij}\times_S Y_{ij}$ existe.
Observemos agora que, se $T$ é um $S$-pré-esquema e se $g:T\to X_i$,
$h:T\to Y_j$ são $S$-morfismos, tem-se necessariamente
$\vphi(g(T))=\psi(h(T))\subset S_i\cap S_j$, pela definição de
$S$-morfismo; logo, $g(T)\subset X_{ij}$ e $h(T)\subset Y_{ij}$; é então
imediato que $Z_{ij}$ é um produto de $X_i$ e $Y_j$.

\begin{env}[3.2.6.5]
\label{I.3.2.6.5}
Podemos agora concluir a demonstração do Teorema~\sref{I.3.2.6}. Se $S$
é um \emph{esquema afim}, existem coberturas $(U_\alpha)$, $(V_\lambda)$
de $X$ e $Y$, respectivamente, formadas por abertos afins; como
$U_\alpha\times_S V_\lambda$ existe em virtude de \sref{I.3.2.2}, o mesmo
vale para $X\times_S Y$ por \sref{I.3.2.6.3}. Se $S$ é um pré-esquema
qualquer, existe uma cobertura $(S_i)$ de $S$ formada por abertos afins. Se
$\vphi:X\to S$, $\psi:Y\to S$ são os morfismos estruturais, e se pusermos
$X_i=\vphi^{-1}(S_i)$, $Y_i=\psi^{-1}(S_i)$, os produtos
$X_i\times_{S_i}Y_i$ existem em virtude do que
\oldpage[I]{108}
precede; mas então os produtos $X_i\times_S Y_i$ também existem
\sref{I.3.2.5}, logo o mesmo vale para $X\times_S Y$ por
\sref{I.3.2.6.4}.\qed
\end{env}

\begin{corollary}[3.2.7]
\label{I.3.2.7}
Seja $Z=X\times_S Y$ o produto de dois $S$-pré-esquemas, sejam $p$, $q$ as
projeções de $Z$ sobre $X$ e $Y$, e seja $\vphi$ (resp. $\psi$) o morfismo
estrutural de $X$ (resp. $Y$). Seja $S'$ uma parte aberta de $S$, e seja
$U$ (resp. $V$) uma parte aberta de $X$ (resp. $Y$) contida em
$\vphi^{-1}(S')$ (resp. $\psi^{-1}(S')$). Então o produto
$U\times_{S'}V$ identifica-se canonicamente com o pré-esquema induzido por
$Z$ sobre $p^{-1}(U)\cap q^{-1}(V)$ (considerado como um
$S'$-pré-esquema). Além disso, se $f:T\to X$, $g:T\to Y$ são
$S$-morfismos tais que $f(T)\subset U$, $g(T)\subset V$, o $S'$-morfismo
$(f,g)_{S'}$ identifica-se com $(f,g)_S$ considerado como um morfismo com
valores em $p^{-1}(U)\cap q^{-1}(V)$.
\end{corollary}

\begin{proof}
Isso resulta do Corolário~\sref{I.3.2.5} e do Lema~\sref{I.3.2.6.1}.
\end{proof}

\begin{env}[3.2.8]
\label{I.3.2.8}
Sejam $(X_\alpha)$, $(Y_\lambda)$ duas famílias de $S$-pré-esquemas, e
sejam $X$ (resp. $Y$) a soma da família $(X_\alpha)$ (resp.
$(Y_\lambda)$) \hyperref[subsection:I.3.1]{(3.1)}. Então
$X\times_S Y$ identifica-se com a \emph{soma} da família
$(X_\alpha\times_S Y_\lambda)$; isso resulta imediatamente de
\sref{I.3.2.6.3}.
\end{env}

\begin{env}[3.2.9]
\label{I.3.2.9}
\footnote{\emph{[Trad.] \sref{I.3.2.9} provém das erratas de EGA~II,
na página 221; daí a mudança na numeração das páginas.}}
\oldpage[II]{221}
Resulta de \sref{I.1.8.1} que se pode precisar \sref{I.3.2.2} da
seguinte maneira: $Z=\Spec(B\otimes_A C)$ não é apenas um produto de
$X=\Spec(B)$ e de $Y=\Spec(C)$ na categoria dos
\emph{$S$-pré-esquemas}, mas também na categoria dos \emph{espaços
localmente anelados sobre $S$} (com uma definição dos $S$-morfismos
calcada na de \sref{I.2.5.2}). A demonstração de \sref{I.3.2.6} prova
também que, para dois $S$-pré-esquemas quaisquer $X$ e $Y$, o
pré-esquema $X\times_S Y$ não é apenas o produto de $X$ e $Y$ na
categoria dos $S$-pré-esquemas, mas também na categoria dos espaços
localmente anelados sobre o pré-esquema $S$.
\end{env}

\subsection{Propriedades formais do produto; mudança do pré-esquema de base}
\label{subsection:I.3.3}

\begin{env}[3.3.1]
\label{I.3.3.1}
O leitor observará que todas as propriedades enunciadas nesta seção,
com exceção de \sref{I.3.3.13}, \sref{I.3.3.15} e da afirmação sobre as
somas na última frase de \sref{I.3.3.10}, são válidas sem modificação em
\emph{toda categoria}, sempre que os produtos que intervêm nos
enunciados \emph{existam} (pois é claro que as noções de objeto sobre
$S$ e de $S$-morfismo podem ser definidas exatamente como em
\hyperref[subsection:I.2.5]{(2.5)} para todo objeto $S$ da categoria).
\end{env}

\begin{env}[3.3.2]
\label{I.3.3.2}
Em primeiro lugar, $X\times_S Y$ é um \emph{bifuntor covariante} em $X$ e
$Y$ na categoria dos $S$-pré-esquemas: basta, com efeito, observar que o
diagrama
\[
  \xymatrix{
    X\times Y\ar[r]^{f\times 1}\ar[d] &
    X'\times Y\ar[r]^{f'\times 1}\ar[d] &
    X''\times Y\ar[d]\\
    X\ar[r]_f & X'\ar[r]_{f'} & X''
  }
\]
é comutativo.
\end{env}

\begin{proposition}[3.3.3]
\label{I.3.3.3}
Para todo $S$-pré-esquema $X$, a primeira (resp. segunda) projeção de
$X\times_S S$ (resp. $S\times_S X$) é um isomorfismo funtorial de
$X\times_S S$ (resp. $S\times_S X$) sobre $X$, cujo isomorfismo inverso é
$(1_X,\vphi)_S$ (resp. $(\vphi,1_X)_S$), denotando-se por $\vphi$ o
morfismo estrutural $X\to S$; pode-se, portanto, escrever, a menos de
isomorfismo canônico,
\[
  X\times_S S=S\times_S X=X.
\]
\end{proposition}

\begin{proof}
Basta provar que a tripla $(X,1_X,\vphi)$ é um produto de $X$ e $S$. Ora, se
$T$ é um $S$-pré-esquema, o único $S$-morfismo de $T$ em $S$ é
necessariamente o morfismo estrutural $\psi:T\to S$. Se $f$ é um
$S$-morfismo de $T$ em $X$, tem-se necessariamente
$\psi=\vphi\circ f$, donde a afirmação.
\end{proof}

\begin{corollary}[3.3.4]
\label{I.3.3.4}
Sejam $X$ e $Y$ dois $S$-pré-esquemas, e sejam $\vphi:X\to S$ e
$\psi:Y\to S$ seus morfismos estruturais. Se identificarmos canonicamente
$X$ com $X\times_S S$ e $Y$ com $S\times_S Y$, as projeções
$X\times_S Y\to X$ e $X\times_S Y\to Y$ identificam-se respectivamente
com $1_X\times\psi$ e $\vphi\times 1_Y$.
\end{corollary}

A verificação é imediata e fica a cargo do leitor.

\begin{env}[3.3.5]
\label{I.3.3.5}
Pode-se definir, da mesma maneira que em
\hyperref[subsection:I.3.2]{(3.2)}, o produto de um
\oldpage[I]{109}
número finito arbitrário $n$ de $S$-pré-esquemas; a existência desses
produtos decorre de \sref{I.3.2.6} por indução sobre $n$, observando que
$(X_1\times_S X_2\times_S\cdots\times_S X_{n-1})\times_S X_n$ satisfaz a
definição do produto. A unicidade do produto implica, como em toda categoria,
suas propriedades de \emph{comutatividade} e de \emph{associatividade}. Se,
por exemplo, $p_1$, $p_2$ e $p_3$ designam as projeções de
$X_1\times_S X_2\times_S X_3$, e se identificarmos esse pré-esquema com
$(X_1\times_S X_2)\times_S X_3$, a projeção sobre $X_1\times_S X_2$
identifica-se com $(p_1,p_2)_S$.
\end{env}

\begin{env}[3.3.6]
\label{I.3.3.6}
Sejam $S$ e $S'$ dois pré-esquemas, e seja $\vphi:S'\to S$ um morfismo, que
faz de $S'$ um $S$-pré-esquema. Para todo $S$-pré-esquema $X$, consideremos o
produto $X\times_S S'$, e sejam $p$ e $\pi'$ suas projeções sobre $X$ e $S'$,
respectivamente. Munido de $\pi'$, esse produto é um $S'$-pré-esquema; quando
considerado como tal, é denotado por $X_{(S')}$ ou $X_{(\vphi)}$, e diz-se que
é o pré-esquema obtido por \emph{extensão do pré-esquema de base} de $S$ a
$S'$, por meio do morfismo $\vphi$, ou a \emph{imagem inversa} de $X$ por
$\vphi$. Note-se que, se $\pi$ é o morfismo estrutural de $X$ e $\theta$ o
morfismo estrutural de $X\times_S S'$, considerado como $S$-pré-esquema, o
diagrama
\[
  \xymatrix{
    X\ar[d]_\pi &
    X_{(S')}\ar[l]_p\ar[ld]_\theta\ar[d]^{\pi'}\\
    S &
    S'\ar[l]^\vphi
  }
\]
é comutativo.
\end{env}

\begin{env}[3.3.7]
\label{I.3.3.7}
Com as notações de \sref{I.3.3.6}, para todo $S$-morfismo $f:X\to Y$,
denota-se também por $f_{(S')}$ o $S'$-morfismo
$f\times_S 1:X_{(S')}\to Y_{(S')}$, e diz-se que $f_{(S')}$ é a
\emph{imagem inversa} do morfismo $f$ por $\vphi$. Portanto, $X_{(S')}$ é um
\emph{funtor covariante} em $X$, da categoria dos $S$-pré-esquemas para a dos
$S'$-pré-esquemas.
\end{env}

\begin{env}[3.3.8]
\label{I.3.3.8}
O pré-esquema $X_{(S')}$ pode ainda ser considerado como solução de um
\emph{problema de aplicação universal}: todo $S'$-pré-esquema $T$ é também um
$S$-pré-esquema por meio de $\vphi$; todo $S$-morfismo $g:T\to X$ escreve-se
então de maneira única como $g=p\circ f$, onde $f$ é um $S'$-morfismo
$T\to X_{(S')}$, como resulta da definição do produto aplicada aos
$S$-morfismos $g$\footnote{\emph{[Trad.] O texto francês imprime aqui $f$; o
par que intervém na definição é $(g,\psi)$.}} e $\psi:T\to S'$ (morfismo
estrutural de $T$).
\end{env}

\begin{proposition}[3.3.9]
\label{I.3.3.9}
\emph{(``Transitividade da extensão dos pré-esquemas de base'')}. Sejam
$S''$ um pré-esquema e $\vphi':S''\to S'$ um morfismo. Para todo
$S$-pré-esquema $X$, existe um isomorfismo canônico funtorial do
$S''$-pré-esquema $(X_{(\vphi)})_{(\vphi')}$ sobre o $S''$-pré-esquema
$X_{(\vphi\circ\vphi')}$.
\end{proposition}

\begin{proof}
Com efeito, sejam $T$ um $S''$-pré-esquema, $\psi$ seu morfismo estrutural, e
$g$ um $S$-morfismo de $T$ em $X$ (considerando $T$ como um $S$-pré-esquema
de morfismo estrutural $\vphi\circ\vphi'\circ\psi$). Como $T$ é também um
$S'$-pré-esquema de morfismo estrutural $\vphi'\circ\psi$, podemos escrever
$g=p\circ g'$, onde $g'$ é um $S'$-morfismo $T\to X_{(\vphi)}$, e depois
$g'=p'\circ g''$, onde $g''$ é um $S''$-morfismo
$T\to(X_{(\vphi)})_{(\vphi')}$:
\[
  \xymatrix{
    X\ar[d]_\pi &
    X_{(\vphi)}\ar[l]_p\ar[d]_{\pi'} &
    (X_{(\vphi)})_{(\vphi')}\ar[l]_{p'}\ar[d]_{\pi''}\\
    S &
    S'\ar[l]^\vphi &
    S''\ar[l]^{\vphi'}
  }
\]
\oldpage[I]{110}
Daí a proposição, em razão da unicidade da solução de um problema de aplicação
universal.
\end{proof}

Esse resultado também se exprime escrevendo a igualdade (a menos de um
isomorfismo canônico) $(X_{(S')})_{(S'')}=X_{(S'')}$, se não houver risco de
confusão, ou ainda
\[
\label{I.3.3.9.1}
  (X\times_S S')\times_{S'}S''=X\times_S S'';
  \tag{3.3.9.1}
\]
o caráter funtorial do isomorfismo definido em \sref{I.3.3.9} exprime-se do
mesmo modo pela fórmula de \emph{transitividade das imagens inversas de
morfismos}
\[
\label{I.3.3.9.2}
  (f_{(S')})_{(S'')}=f_{(S'')}
  \tag{3.3.9.2}
\]
para todo $S$-morfismo $f:X\to Y$.

\begin{corollary}[3.3.10]
\label{I.3.3.10}
Se $X$ e $Y$ são dois $S$-pré-esquemas, existe um isomorfismo canônico
funtorial do $S'$-pré-esquema $X_{(S')}\times_{S'}Y_{(S')}$ sobre o
$S'$-pré-esquema $(X\times_S Y)_{(S')}$. 
\end{corollary}

\begin{proof}
Temos, a menos de isomorfismos canônicos,
\[
  (X\times_S S')\times_{S'}(Y\times_S S')
  =X\times_S(Y\times_S S')=(X\times_S Y)\times_S S'
\]
em virtude de \sref{I.3.3.9.1} e da associatividade dos produtos de
$S$-pré-esquemas.
\end{proof}

O caráter funtorial do isomorfismo definido no corolário~\sref{I.3.3.10}
exprime-se pela fórmula
\[
\label{I.3.3.10.1}
  (u_{(S')},v_{(S')})_{S'}=((u,v)_S)_{(S')}
  \tag{3.3.10.1}
\]
para todo par de $S$-morfismos $u:T\to X$, $v:T\to Y$.

Em outros termos, o funtor de imagem inversa $X_{(S')}$ \emph{comuta com a
formação dos produtos}; comuta também com a formação das somas
\sref{I.3.2.8}.

\begin{corollary}[3.3.11]
\label{I.3.3.11}
Sejam $Y$ um $S$-pré-esquema e $f:X\to Y$ um morfismo que faz de $X$ um
$Y$-pré-esquema (e, por conseguinte, também um $S$-pré-esquema). O
pré-esquema $X_{(S')}$ identifica-se então com o produto
$X\times_Y Y_{(S')}$, identificando-se a projeção
$X\times_Y Y_{(S')}\to Y_{(S')}$ com $f_{(S')}$. 
\end{corollary}

\begin{proof}
Seja $\psi:Y\to S$ o morfismo estrutural de $Y$; temos o diagrama comutativo
\[
  \xymatrix{
    S'\ar[d] &
    Y_{(S')}\ar[l]_{\psi_{(S')}}\ar[d] &
    X_{(S')}\ar[l]_{f_{(S')}}\ar[d]\\
    S &
    Y\ar[l]^\psi &
    X\ar[l]^f
  }
\]
Ora, $Y_{(S')}$ identifica-se com $S'_{(\psi)}$, e $X_{(S')}$ com
$S'_{(\psi\circ f)}$; levando em conta \sref{I.3.3.9} e
\sref{I.3.3.4}, deduzimos o corolário.
\end{proof}

\begin{env}[3.3.12]
\label{I.3.3.12}
Sejam $f:X\to X'$ e $g:Y\to Y'$ dois $S$-morfismos que são
\emph{monomorfismos} de pré-esquemas (T, I, 1.1); então $f\times_S g$ é um
\emph{monomorfismo}. Com efeito, se $p$ e $q$ são as projeções de
$X\times_S Y$, $p'$ e $q'$ as projeções de $X'\times_S Y'$, e $u$ e $v$ são
$S$-morfismos $T\to X\times_S Y$, então a relação
$(f\times_S g)\circ u=(f\times_S g)\circ v$ implica
$p'\circ(f\times_S g)\circ u=p'\circ(f\times_S g)\circ v$, ou seja,
$f\circ p\circ u=f\circ p\circ v$; como $f$ é um monomorfismo,
$p\circ u=p\circ v$. Usando o fato de que $g$ é um monomorfismo, obtemos do
mesmo modo $q\circ u=q\circ v$, donde $u=v$.

\oldpage[I]{111}
Segue-se que, para toda extensão $S'\to S$ do pré-esquema de base,
\[
  f_{(S')}:X_{(S')}\to X'_{(S')}
\]
é um monomorfismo.\footnote{\emph{[Trad.] O texto francês imprime
$Y_{(S')}$; como 3.3.12 supõe $f:X\to X'$, o contradomínio após a extensão do
pré-esquema de base é $X'_{(S')}$.}}
\end{env}

\begin{env}[3.3.13]
\label{I.3.3.13}
Sejam $S$ e $S'$ dois esquemas afins, e sejam $A$ e $A'$ seus anéis
respectivos; um morfismo $S'\to S$ corresponde, portanto, a um homomorfismo
de anéis $A\to A'$. Se $X$ é um $S$-pré-esquema, denotaremos também por
$X_{(A')}$ ou $X\otimes_A A'$ o $S'$-pré-esquema $X_{(S')}$. Quando o próprio
$X$ é um esquema afim de anel $B$, $X_{(A')}$ é um esquema afim de anel
$B_{(A')}=B\otimes_A A'$, obtido estendendo-se os escalares da $A$-álgebra
$B$ a $A'$.
\end{env}

\begin{env}[3.3.14]
\label{I.3.3.14}
Com as notações de \sref{I.3.3.6}, para todo \emph{$S$-morfismo}
$f:S'\to X$, tem-se que $f'=(f,1_{S'})_S$ é um $S'$-morfismo
$S'\to X'=X_{(S')}$ tal que $p\circ f'=f$, $\pi'\circ f'=1_{S'}$, em outras
palavras, uma \emph{$S'$-seção de $X'$}; e, reciprocamente, se $f'$ é uma tal
$S'$-seção, então $f=p\circ f'$ é um $S$-morfismo $S'\to X$. Define-se assim
uma \emph{correspondência biunívoca} canônica
\[
  \Hom_S(S',X)\isoto\Hom_{S'}(S',X').
\]
Diz-se que $f'$ é o \emph{morfismo gráfico} de $f$, e ele é denotado por
$\Gamma_f$.
\end{env}

\begin{env}[3.3.15]
\label{I.3.3.15}
Dado um pré-esquema $X$, que sempre se pode considerar como um
$\bb{Z}$-pré-esquema, resulta em particular de \sref{I.3.3.14} que as
\emph{$X$-seções} de $X\otimes_\bb{Z}\bb{Z}[T]$ (onde $T$ é uma
indeterminada) correspondem biunivocamente a \emph{morfismos}
$X\to\bb{Z}[T]$\footnote{\emph{[Trad.] O texto francês imprime
$\bb{Z}[T]\to X$; a Proposição~\sref{I.3.3.14} e os homomorfismos de anéis
subsequentes $A_\alpha[T]\to A_\alpha$ exigem o morfismo de esquemas
$X\to\bb{Z}[T]$.}}. Mostremos que essas $X$-seções também correspondem
biunivocamente a \emph{seções do feixe estrutural $\sh{O}_X$ sobre $X$}.
Com efeito, seja $(U_\alpha)$ uma cobertura de $X$ por abertos afins; seja
$u:X\to X\otimes_\bb{Z}\bb{Z}[T]$ um $X$-morfismo, e seja $u_\alpha$ sua
restrição a $U_\alpha$; se $A_\alpha$ é o anel do esquema afim $U_\alpha$,
então $U_\alpha\otimes_\bb{Z}\bb{Z}[T]$ é um esquema afim de anel
$A_\alpha[T]$ \sref{I.3.2.2}, e $u_\alpha$ corresponde canonicamente a um
$A_\alpha$-homomorfismo $A_\alpha[T]\to A_\alpha$ \sref{I.1.7.3}. Ora, tal
homomorfismo é completamente determinado pela imagem de $T$ em $A_\alpha$,
seja $s_\alpha\in A_\alpha=\Gamma(U_\alpha,\sh{O}_X)$; e, ao escrever que as
restrições de $u_\alpha$ e $u_\beta$ a um aberto afim
$V\subset U_\alpha\cap U_\beta$ coincidem, vemos imediatamente que
$s_\alpha$ e $s_\beta$ coincidem em $V$; portanto, a família $(s_\alpha)$ é
formada pelas restrições aos $U_\alpha$ de uma seção $s$ de $\sh{O}_X$ sobre
$X$; reciprocamente, é claro que tal seção define uma família $(u_\alpha)$
de morfismos que são as restrições aos $U_\alpha$ de um $X$-morfismo
$X\to X\otimes_\bb{Z}\bb{Z}[T]$. Este resultado será generalizado em
\sref[II]{II.1.7.12}.
\end{env}

\subsection{Pontos de um pré-esquema com valores em um pré-esquema; pontos geométricos}
\label{subsection:I.3.4}

\begin{env}[3.4.1]
\label{I.3.4.1}
Seja $X$ um pré-esquema; para todo pré-esquema $T$, denota-se ainda por
$X(T)$ o conjunto $\Hom(T,X)$ dos morfismos $T\to X$, e os elementos desse
conjunto também são chamados \emph{pontos de $X$ com valores em $T$}. Se a
todo morfismo $f:T\to T'$ associarmos a aplicação $u'\mapsto u'\circ f$ de
$X(T')$ em $X(T)$, vemos que, para $X$ fixo, $X(T)$ é um \emph{funtor
contravariante em $T$}, da categoria dos pré-esquemas para a dos conjuntos.
Além disso, todo morfismo de pré-esquemas $g:X\to Y$ define um homomorfismo
funtorial $X(T)\to Y(T)$, que envia $v\in X(T)$ em $g\circ v$.
\end{env}

\begin{env}[3.4.2]
\label{I.3.4.2}
Dados três conjuntos $P$, $Q$ e $R$, e duas aplicações $\vphi:P\to R$ e
$\psi:Q\to R$, chama-se \emph{produto fibrado de $P$ e $Q$ sobre $R$}
(relativamente a $\vphi$ e $\psi$) o subconjunto do
\oldpage[I]{112}
conjunto produto $P\times Q$ formado pelos pares $(p,q)$ tais que
$\vphi(p)=\psi(q)$; ele é denotado por $P\times_R Q$. A Definição~
\sref{I.3.2.1} do produto de $S$-pré-esquemas também pode ser interpretada,
com a notação de \sref{I.3.4.1}, pela fórmula
\[
\label{I.3.4.2.1}
  (X\times_S Y)(T)=X(T)\times_{S(T)}Y(T),
  \tag{3.4.2.1}
\]
em que as aplicações $X(T)\to S(T)$ e $Y(T)\to S(T)$ correspondem aos
morfismos estruturais $X\to S$ e $Y\to S$.
\end{env}

\begin{env}[3.4.3]
\label{I.3.4.3}
Dado um pré-esquema $S$, e considerando-se somente os $S$-pré-esquemas e os
$S$-morfismos, ainda se denotará por $X(T)_S$ o conjunto $\Hom_S(T,X)$ dos
$S$-morfismos $T\to X$, permitindo-se suprimir o índice $S$ quando não houver
possibilidade de confusão; diz-se ainda que os elementos de $X(T)_S$ são os
\emph{pontos} (ou \emph{$S$-pontos}, quando houver possibilidade de confusão)
\emph{do $S$-pré-esquema $X$ com valores no $S$-pré-esquema $T$}. Em
particular, uma \emph{$S$-seção} de $X$ não é senão um \emph{ponto de $X$ com
valores em $S$}. A fórmula \sref{I.3.4.2.1} também se escreve então
\[
\label{I.3.4.3.1}
  (X\times_S Y)(T)_S=X(T)_S\times Y(T)_S;
  \tag{3.4.3.1}
\]
mais geralmente, se $Z$ é um $S$-pré-esquema, e $X$, $Y$ e $T$ são
$Z$-pré-esquemas (portanto, \emph{ipso facto}, $S$-pré-esquemas), tem-se
\[
\label{I.3.4.3.2}
  (X\times_Z Y)(T)_S=X(T)_S\times_{Z(T)_S}Y(T)_S.
  \tag{3.4.3.2}
\]

Observe-se que, para demonstrar que uma tripla $(W,r,s)$, formada por um
$S$-pré-esquema $W$ e por dois $S$-morfismos $r:W\to X$ e $s:W\to Y$, é um
produto de $X$ e $Y$ (sobre $Z$), basta, por definição, verificar que, para
\emph{todo $S$-pré-esquema $T$}, o diagrama
\[
  \xymatrix{
    W(T)_S\ar[r]^{r'}\ar[d]_{s'} &
    X(T)_S\ar[d]^{\vphi'}\\
    Y(T)_S\ar[r]_{\psi'} &
    Z(T)_S
  }
\]
faz de $W(T)_S$ o produto fibrado de $X(T)_S$ e $Y(T)_S$ sobre $Z(T)_S$,
correspondendo $r'$ e $s'$ a $r$ e $s$, e $\vphi'$ e $\psi'$ aos morfismos
estruturais $\vphi:X\to Z$ e $\psi:Y\to Z$.
\end{env}

\begin{env}[3.4.4]
\label{I.3.4.4}
Quando, no que precede, $T$ (resp. $S$) é um esquema afim de anel $B$
(resp. $A$), substitui-se $T$ (resp. $S$) por $B$ (resp. $A$) nas notações
anteriores, e chamam-se então os elementos de $X(B)$ de \emph{pontos de $X$
com valores no anel $B$}, e os elementos de $X(B)_A$ de \emph{pontos do
$A$-pré-esquema $X$ com valores na $A$-álgebra $B$}. Observe-se que $X(B)$ e
$X(B)_A$ são agora \emph{funtores covariantes em $B$}. Escreve-se do mesmo
modo $X(T)_A$ para o conjunto dos pontos do $A$-pré-esquema $X$ com valores
no $A$-pré-esquema $T$.
\end{env}

\begin{env}[3.4.5]
\label{I.3.4.5}
Consideremos, em particular, o caso em que $T$ é da forma $\Spec(A)$, onde
$A$ é um anel \emph{local}; os elementos de $X(A)$ correspondem
biunivocamente aos \emph{homomorfismos locais} $\sh{O}_x\to A$ para $x\in X$
\sref{I.2.4.4}; diz-se que o ponto $x$ do espaço subjacente a $X$ é a
\emph{localidade}\footnote{\emph{[Trad.] Diz-se também que o ponto geométrico
está sobre esse $x$.}} do ponto de $X$ com valores em $A$ que lhe corresponde.

Mais precisamente, chamam-se \emph{pontos geométricos} de um pré-esquema $X$
os \emph{pontos de $X$ com valores em um corpo $K$}: dar um tal ponto equivale,
portanto, a dar sua
\oldpage[I]{113}
localidade $x$ no espaço subjacente a $X$ e uma \emph{extensão} $K$ de
$\kres(x)$; $K$ será chamado o \emph{corpo de valores} do ponto geométrico
correspondente, e diz-se ainda que esse ponto geométrico está \emph{localizado
em $x$}. Define-se assim uma aplicação $X(K)\to X$, que envia um ponto
geométrico com valores em $K$ à sua localidade.

Se $S'=\Spec(K)$ é um $S$-pré-esquema (em outras palavras, se $K$ é
considerado uma extensão do corpo residual $\kres(s)$, onde $s\in S$), e se
$X$ é um $S$-pré-esquema, um elemento de $X(K)_S$, ou, como também se diz, um
\emph{ponto geométrico de $X$ situado sobre $s$ com valores em $K$}, consiste
na escolha de um $\kres(s)$-monomorfismo do corpo residual $\kres(x)$ em $K$,
onde $x$ é um ponto de $X$ \emph{situado sobre $s$} (portanto, $\kres(x)$ é
uma extensão de $\kres(s)$).

Em particular, se $S=\Spec(K)=\{\xi\}$, \emph{os pontos geométricos de $X$
com valores em $K$ identificam-se com os pontos $x\in X$ tais que
$\kres(x)=K$}; diz-se ainda que estes últimos são os \emph{pontos racionais
sobre $K$ do $K$-pré-esquema $X$}; se $K'$ é uma extensão de $K$, os
pontos geométricos de $X$ com valores em $K'$ correspondem, portanto,
biunivocamente aos pontos $K'$-racionais de $X'=X_{(K')}$~
\sref{I.3.3.14}.
\end{env}

\begin{lemma}[3.4.6]
\label{I.3.4.6}
Sejam $X_i$ ($1\leq i\leq n$) $S$-pré-esquemas, $s$ um ponto de $S$, e $x_i$
($1\leq i\leq n$) pontos de $X_i$ situados sobre $s$. Então existe uma
extensão $K$ de $\kres(s)$ e um ponto geométrico do produto
$Y=X_1\times_S X_2\times_S\cdots\times_S X_n$, com valores em $K$, cujas
projeções sobre os $X_i$ estejam localizadas nos $x_i$.
\end{lemma}

\begin{proof}
Existem $\kres(s)$-monomorfismos $\kres(x_i)\to K$, todos em uma mesma
extensão $K$ de $\kres(s)$ (Bourbaki, \emph{Alg.}, cap.~V, \textsection4,
prop.~2). As compostas $\kres(s)\to\kres(x_i)\to K$ são todas idênticas;
portanto, os morfismos $\Spec(K)\to X_i$ correspondentes aos morfismos
$\kres(x_i)\to K$ são todos $S$-morfismos, e concluímos assim que definem um
único morfismo $\Spec(K)\to Y$. Se $y$ é o ponto correspondente de $Y$, é
claro que sua projeção sobre cada um dos $X_i$ é $x_i$.
\end{proof}

\begin{proposition}[3.4.7]
\label{I.3.4.7}
Sejam $X_i$ ($1\leq i\leq n$) $S$-pré-esquemas e, para cada índice $i$,
seja $x_i$ um ponto de $X_i$. Para que exista um ponto $y$ de
$Y=X_1\times_S X_2\times\ldots\times_S X_n$ cuja imagem seja $x_i$ pela
projeção de índice $i$ para todo $1\leq i\leq n$, é necessário e suficiente
que todos os $x_i$ estejam situados sobre um mesmo ponto $s$ de $S$.
\end{proposition}

\begin{proof}
A condição é evidentemente necessária; o Lema~\sref{I.3.4.6} demonstra que é
suficiente.
\end{proof}

Em outras palavras, se denotamos por $(X)$ o conjunto subjacente de $X$,
vemos que há uma aplicação canônica \emph{sobrejetiva}
$(X\times_S Y)\to(X)\times_{(S)}(Y)$; deve-se observar que essa aplicação
\emph{não é injetiva} em geral; dito de outro modo, \emph{podem existir vários
pontos distintos $z$ em $X\times_S Y$ com as mesmas projeções $x\in X$ e
$y\in Y$}; isso já se vê quando $S$, $X$ e $Y$ são espectros primos dos
corpos $k$, $K$ e $K'$, respectivamente, pois o produto tensorial
$K\otimes_k K'$ tem, em geral, vários ideais primos distintos
(cf.~\sref{I.3.4.9}).

\begin{corollary}[3.4.8]
\label{I.3.4.8}
Seja $f:X\to Y$ um $S$-morfismo, e seja
$f_{(S')}:X_{(S')}\to Y_{(S')}$ o $S'$-morfismo deduzido de $f$ por uma
extensão $S'\to S$ do pré-esquema de base. Seja $p$ (\emph{resp. $q$}) a
projeção $X_{(S')}\to X$ (\emph{resp. $Y_{(S')}\to Y$}); para toda parte $M$
de $X$, tem-se
\[
  q^{-1}(f(M))=f_{(S')}(p^{-1}(M)).
\]
\end{corollary}

\oldpage[I]{114}
\begin{proof}
Com efeito, por \sref{I.3.3.11}, $X_{(S')}$ identifica-se com o produto
$X\times_Y Y_{(S')}$ graças ao diagrama comutativo
\[
  \xymatrix{
    X\ar[d]_f &
    X_{(S')}\ar[l]_p\ar[d]^{f_{(S')}}\\
    Y &
    Y_{(S')}\ar[l]^q
  }
\]
Em virtude de \sref{I.3.4.7}, a relação $q(y')=f(x)$ para $x\in M$ e
$y'\in Y_{(S')}$ equivale à existência de algum $x'\in X_{(S')}$ tal que
$p(x')=x$ e $f_{(S')}(x')=y'$, donde resulta o corolário.
\end{proof}

O Lema~\sref{I.3.4.6} pode ser precisado da seguinte maneira:

\begin{proposition}[3.4.9]
\label{I.3.4.9}
Sejam $X$ e $Y$ dois $S$-pré-esquemas, $x$ um ponto de $X$ e $y$ um ponto de
$Y$, estando $x$ e $y$ ambos situados sobre o mesmo ponto $s\in S$.
O conjunto dos pontos de $X\times_S Y$ cujas projeções são $x$ e $y$ está em
correspondência biunívoca canônica com o conjunto dos tipos de extensões
compostas de $\kres(x)$ e $\kres(y)$ consideradas como extensões de $\kres(s)$
\emph{(Bourbaki, \emph{Alg.}, cap.~VIII, \textsection8, prop.~2)}.
\end{proposition}

\begin{proof}
Seja $p$ (resp. $q$) a projeção de $X\times_S Y$ sobre $X$ (resp. $Y$), e seja
$E$ o subespaço $p^{-1}(x)\cap q^{-1}(y)$ do espaço subjacente a
$X\times_S Y$.
Observemos, em primeiro lugar, que os morfismos $\Spec(\kres(x))\to S$ e
$\Spec(\kres(y))\to S$ se fatorizam como
$\Spec(\kres(x))\to\Spec(\kres(s))\to S$ e
$\Spec(\kres(y))\to\Spec(\kres(s))\to S$; como
$\Spec(\kres(s))\to S$ é um monomorfismo \sref{I.2.4.7}, segue de
\sref{I.3.2.4} que
\[
  P=\Spec(\kres(x))\times_S\Spec(\kres(y))=\Spec(\kres(x))\times_{\Spec(\kres(s))}\Spec(\kres(y))=\Spec(\kres(x)\otimes_{\kres(s)}\kres(y)).
\]
Definiremos duas aplicações, $\alpha:P_0\to E$ e $\beta:E\to P_0$, inversas
uma da outra (onde $P_0$ designa o conjunto subjacente ao pré-esquema $P$).
Se $i:\Spec(\kres(x))\to X$ e $j:\Spec(\kres(y))\to Y$ são os morfismos
canônicos \sref{I.2.4.5}, tomaremos para $\alpha$ a aplicação dos espaços
subjacentes correspondente ao morfismo $i\times_S j$.
Por outro lado, todo $z\in E$ define, por hipótese, dois
$\kres(s)$-monomorfismos, $\kres(x)\to\kres(z)$ e $\kres(y)\to\kres(z)$, e,
portanto, um $\kres(s)$-homomorfismo\footnote{\emph{[Trad.] O texto francês
imprime «monomorfismo» aqui. A aplicação induzida pelo produto tensorial não
precisa ser injetiva; basta um homomorfismo para definir o morfismo de espectros
seguinte.}} $\kres(x)\otimes_{\kres(s)}\kres(y)\to\kres(z)$ e, por
conseguinte, um morfismo $\Spec(\kres(z))\to P$; $\beta(z)$ será a imagem de
$z$ em $P_0$ por esse morfismo.
A verificação de que $\alpha\circ\beta$ e $\beta\circ\alpha$ são as aplicações
identidade decorre de \sref{I.2.4.5} e da definição do produto
\sref{I.3.2.1}.
Finalmente, sabemos que $P_0$ está em correspondência biunívoca com o conjunto
dos tipos de extensões compostas de $\kres(x)$ e $\kres(y)$ (Bourbaki,
\emph{Alg.}, cap.~VIII, \textsection8, prop.~1).
\end{proof}

\subsection{Sobrejetividade e injetividade}
\label{subsection:I.3.5}

\begin{env}[3.5.1]
\label{I.3.5.1}
De modo geral, consideremos uma propriedade $\textbf{P}$ dos morfismos de
pré-esquemas e as duas proposições seguintes:
\begin{enumerate}
  \item[(i)] Se $f:X\to X'$ e $g:Y\to Y'$ são dois $S$-morfismos que possuem
  a propriedade $\textbf{P}$, então $f\times_S g$ também possui a propriedade
  $\textbf{P}$.
  \item[(ii)] Se $f:X\to Y$ é um $S$-morfismo que possui a propriedade
  $\textbf{P}$, então todo $S'$-morfismo
  $f_{(S')}:X_{(S')}\to Y_{(S')}$ obtido de $f$ por uma extensão do
  pré-esquema de base $S'\to S$ também possui a propriedade $\textbf{P}$.
\end{enumerate}

\oldpage[I]{115}
Como $f_{(S')}=f\times_S 1_{S'}$, vemos que, se para todo pré-esquema $X$ a
\emph{identidade} $1_X$ possui a propriedade $\textbf{P}$, então (i) implica
(ii); por outro lado, como $f\times_S g$ é o morfismo composto
\[
  X\times_S Y\xrightarrow{f\times1_Y}X'\times_S Y\xrightarrow{1_{X'}\times g}X'\times_S Y',
\]
vemos que, se o \emph{composto} de dois morfismos que possuem a propriedade
$\textbf{P}$ também possui a propriedade $\textbf{P}$, então o produto
$f\times_S g$ também a possui e, portanto, (ii) implica (i).
\end{env}

Uma primeira aplicação desta observação é a seguinte.

\begin{proposition}[3.5.2]
\label{I.3.5.2}
\medskip\noindent
\begin{enumerate}
  \item[{\rm(i)}] Se $f:X\to X'$ e $g:Y\to Y'$ são $S$-morfismos
  sobrejetivos, então $f\times_S g$ é sobrejetivo.
  \item[{\rm(ii)}] Se $f:X\to Y$ é um $S$-morfismo sobrejetivo, então
  $f_{(S')}$ é sobrejetivo para toda extensão do pré-esquema de base
  $S'\to S$.
\end{enumerate}
\end{proposition}

\begin{proof}
Como o composto de dois morfismos sobrejetivos é sobrejetivo, basta demonstrar
(ii); mas esta proposição resulta imediatamente de \sref{I.3.4.8} aplicado a
$M=X$.
\end{proof}

\begin{proposition}[3.5.3]
\label{I.3.5.3}
Para que um morfismo $f:X\to Y$ seja sobrejetivo, é necessário e suficiente
que, para todo corpo $K$ e todo morfismo $\Spec(K)\to Y$, exista uma extensão
$K'$ de $K$ e um morfismo $\Spec(K')\to X$ que torne comutativo o diagrama
seguinte:
\[
  \xymatrix{
    X\ar[d]_f &
    \Spec(K')\ar[l]\ar[d]\\
    Y &
    \Spec(K).\ar[l]
  }
\]
\end{proposition}

\begin{proof}
A condição é suficiente porque, para todo $y\in Y$, basta aplicá-la a um
morfismo $\Spec(K)\to Y$ correspondente a um monomorfismo
$\kres(y)\to K$, sendo $K$ uma extensão de $\kres(y)$ \sref{I.2.4.6}.
Reciprocamente, suponhamos que $f$ seja sobrejetivo, e seja $y\in Y$ a imagem
do único ponto de $\Spec(K)$; existe um $x\in X$ tal que $f(x)=y$;
consideremos o monomorfismo correspondente $\kres(y)\to\kres(x)$
\sref{I.2.2.1}; basta então tomar $K'$ uma extensão de $\kres(y)$ tal que
existam $\kres(y)$-monomorfismos de $\kres(x)$ e de $K$ em $K'$ (Bourbaki,
\emph{Alg.}, cap.~V, \textsection4, prop.~2); o morfismo
$\Spec(K')\to X$ correspondente a $\kres(x)\to K'$ é precisamente o
procurado.
\end{proof}

Com a linguagem introduzida em \sref{I.3.4.5}, também podemos dizer que
\emph{todo ponto geométrico de $Y$ com valores em $K$ provém de um ponto
geométrico de $X$ com valores em uma extensão de $K$}.

\begin{definition}[3.5.4]
\label{I.3.5.4}
Diz-se que um morfismo de pré-esquemas $f:X\to Y$ é
\emph{universalmente injetivo}, ou um \emph{morfismo radicial}, se, para todo
corpo $K$, a aplicação correspondente $X(K)\to Y(K)$ é injetiva.
\end{definition}

Resulta imediatamente das definições que todo \emph{monomorfismo de
pré-esquemas} (T,~1.1) é radicial.

\begin{env}[3.5.5]
\label{I.3.5.5}
Para que um morfismo $f:X\to Y$ seja radicial, basta que a condição da
Definição~\sref{I.3.5.4} seja satisfeita para todo corpo
\emph{algebricamente fechado}.
Com efeito, se $K$ é um corpo arbitrário e $K'$ uma extensão algebricamente
fechada de $K$, então o diagrama
\[
  \xymatrix{
    X(K)\ar[r]^\alpha\ar[d]_\vphi &
    Y(K)\ar[d]^{\vphi'}\\
    X(K')\ar[r]_{\alpha'} &
    Y(K')
  }
\]
\oldpage[I]{116}
é comutativo, onde $\vphi$ e $\vphi'$ provêm do morfismo
$\Spec(K')\to\Spec(K)$, e $\alpha$ e $\alpha'$ correspondem a $f$.
Contudo, $\vphi$ é injetiva, e $\alpha'$ também o é, por hipótese; portanto,
$\alpha$ é necessariamente injetiva.
\end{env}

\begin{proposition}[3.5.6]
\label{I.3.5.6}
Sejam $f:X\to Y$ e $g:Y\to Z$ dois morfismos de pré-esquemas.
\begin{enumerate}
  \item[{\rm(i)}] Se $f$ e $g$ são radiciais, então $g\circ f$ também o é.
  \item[{\rm(ii)}] Reciprocamente, se $g\circ f$ é radicial, então $f$
  também o é.
\end{enumerate}
\end{proposition}

\begin{proof}
Levando em conta a Definição~\sref{I.3.5.4}, a proposição se reduz às
afirmações correspondentes para as aplicações $X(K)\to Y(K)\to Z(K)$, e
essas afirmações são evidentes.
\end{proof}

\begin{proposition}[3.5.7]
\label{I.3.5.7}
\medskip\noindent
\begin{enumerate}
  \item[{\rm(i)}] Se os $S$-morfismos $f:X\to X'$ e $g:Y\to Y'$ são
  radiciais, então $f\times_S g$ também o é.
  \item[{\rm(ii)}] Se o $S$-morfismo $f:X\to Y$ é radicial, então
  $f_{(S')}:X_{(S')}\to Y_{(S')}$ também o é para toda extensão $S'\to S$
  do pré-esquema de base.
\end{enumerate}
\end{proposition}

\begin{proof}
Em vista de \sref{I.3.5.1}, basta demonstrar (i).
Vimos em \sref{I.3.4.2.1} que
\[
  (X\times_S Y)(K)=X(K)\times_{S(K)}Y(K),
\]
\[
  (X'\times_S Y')(K)=X'(K)\times_{S(K)}Y'(K),
\]
e a aplicação $(X\times_S Y)(K)\to(X'\times_S Y')(K)$ correspondente a
$f\times_S g$ identifica-se então com $(u,v)\to(f\circ u,g\circ v)$, de
modo que a proposição resulta imediatamente.
\end{proof}

\begin{proposition}[3.5.8]
\label{I.3.5.8}
Para que um morfismo $f=(\psi,\theta):X\to Y$ seja radicial, é necessário
e suficiente que $\psi$ seja injetiva e que o monomorfismo
$\theta^x:\kres(\psi(x))\to\kres(x)$ faça de $\kres(x)$ uma extensão
radicial de $\kres(\psi(x))$ para todo $x\in X$.
\end{proposition}

\begin{proof}
Suponhamos que $f$ seja radicial e mostremos primeiro que a igualdade
$\psi(x_1)=\psi(x_2)=y$ implica necessariamente $x_1=x_2$.
Com efeito, existe um corpo $K$ que estende $\kres(y)$, juntamente com
$\kres(y)$-monomorfismos $\kres(x_1)\to K$ e $\kres(x_2)\to K$ (Bourbaki,
\emph{Alg.}, cap.~V, \textsection4, prop.~2); os morfismos correspondentes
$u_1:\Spec(K)\to X$ e $u_2:\Spec(K)\to X$ são então tais que
$f\circ u_1=f\circ u_2$, e, por hipótese, $u_1=u_2$, o que implica, em
particular, $x_1=x_2$.
Consideremos agora $\kres(x)$ como uma extensão de $\kres(\psi(x))$ por
meio de $\theta^x$: se $\kres(x)$ não é uma extensão radicial de
$\kres(\psi(x))$, existem dois $\kres(\psi(x))$-monomorfismos distintos de
$\kres(x)$ em um corpo algebricamente fechado $K$ que estende
$\kres(\psi(x))$, e os dois morfismos correspondentes $\Spec(K)\to X$
contradiriam a hipótese.
Reciprocamente, em vista de \sref{I.2.4.6}, é imediato que as condições do
enunciado são suficientes para que $f$ seja radicial.
\end{proof}

\begin{corollary}[3.5.9]
\label{I.3.5.9}
Se $A$ é um anel e $S$ é um conjunto multiplicativo de $A$, então o
morfismo canônico $\Spec(S^{-1}A)\to\Spec(A)$ é radicial.
\end{corollary}

\begin{proof}
Com efeito, este morfismo é um monomorfismo \sref{I.1.6.2}.
\end{proof}

\begin{corollary}[3.5.10]
\label{I.3.5.10}
Seja $f:X\to Y$ um morfismo radicial, $g:Y'\to Y$ um morfismo e
$X'=X_{(Y')}=X\times_Y Y'$.
Então o morfismo radicial $f_{(Y')}$ \sref{I.3.5.7}[ii] é uma bijeção do
espaço subjacente a $X'$ sobre $g^{-1}(f(X))$; além disso, para todo corpo
$K$, o conjunto $X'(K)$ identifica-se com a imagem inversa do subconjunto
$X(K)\subseteq Y(K)$ pela aplicação $Y'(K)\to Y(K)$ correspondente a $g$.
\end{corollary}

\begin{proof}
A primeira afirmação resulta imediatamente de \sref{I.3.5.8} e
\sref{I.3.4.8}; a segunda, da comutatividade do diagrama seguinte:
\oldpage[I]{117}
\[
  \xymatrix{
    X'(K)\ar[r]\ar[d] &
    Y'(K)\ar[d]\\
    X(K)\ar[r] &
    Y(K)
  }
\]
\end{proof}

\begin{remark}[3.5.11]
\label{I.3.5.11}
Diremos que um morfismo $f=(\psi,\theta)$ de pré-esquemas é
\emph{injetivo} se a aplicação $\psi$ é injetiva.
Para que um morfismo $f=(\psi,\theta):X\to Y$ seja radicial, é necessário
e suficiente que, para todo morfismo $Y'\to Y$, o morfismo
$f_{(Y')}:X_{(Y')}\to Y'$ seja injetivo (o que justifica a terminologia de
morfismo \emph{universalmente injetivo}).
Com efeito, a condição é necessária em virtude de
\sref{I.3.5.7}[ii] e \sref{I.3.5.8}.
Reciprocamente, a condição implica, antes de tudo, que $\psi$ é injetiva;
se, para algum $x\in X$, o monomorfismo
$\theta^x:\kres(\psi(x))\to\kres(x)$ não fosse radicial, existiriam uma
extensão $K$ de $\kres(\psi(x))$ e dois morfismos distintos
$\Spec(K)\to X$ correspondentes ao mesmo morfismo $\Spec(K)\to Y$
\sref{I.3.5.8}.
Mas então, tomando $Y'=\Spec(K)$, haveria duas $Y'$-seções distintas de
$X_{(Y')}$ \sref{I.3.3.14}, o que contradiz a hipótese de que
$f_{(Y')}$ é injetivo.
\end{remark}

\subsection{Fibras}
\label{subsection:I.3.6}

\begin{proposition}[3.6.1]
\label{I.3.6.1}
Sejam $f:X\to Y$ um morfismo, $y$ um ponto de $Y$ e $\mathfrak{a}_y$ um ideal de definição de $\sh{O}_y$ para a topologia $\mathfrak{m}_y$-pré-ádica.
Então a projeção $p:X\times_Y\Spec(\sh{O}_y/\mathfrak{a}_y)\to X$ é um homeomorfismo do espaço subjacente ao pré-esquema $X\times_Y\Spec(\sh{O}_y/\mathfrak{a}_y)$ sobre a fibra $f^{-1}(y)$, munida da topologia induzida pela do espaço subjacente a $X$.
\end{proposition}

\begin{proof}
Como $\Spec(\sh{O}_y/\mathfrak{a}_y)\to Y$ é radicial (\sref{I.3.5.4} e \sref{I.2.4.7}), como $\Spec(\sh{O}_y/\mathfrak{a}_y)$ se reduz a um único ponto e como o ideal $\mathfrak{m}_y/\mathfrak{a}_y$ é nilpotente por hipótese \sref{I.1.1.12}, já sabemos (\sref{I.3.5.10} e \sref{I.3.3.4}) que $p$ identifica, enquanto \emph{conjuntos}, o espaço subjacente a $X\times_Y\Spec(\sh{O}_y/\mathfrak{a}_y)$ com $f^{-1}(y)$; tudo se reduz a demonstrar que $p$ é um homeomorfismo.
Em virtude de \sref{I.3.2.7}, a questão é local em $X$ e $Y$, e podemos, portanto, supor que $X=\Spec(B)$ e $Y=\Spec(A)$, sendo $B$ uma $A$-álgebra.
O morfismo $p$ corresponde então ao homomorfismo $1\otimes\vphi:B\to B\otimes_A A'$, onde $A'=A_y/\mathfrak{a}_y$ e $\vphi$ é a aplicação canônica de $A$ em $A'$.
Todo elemento de $B\otimes_A A'$ escreve-se então sob a forma
\[
  \sum_i b_i\otimes\vphi(a_i)/\vphi(s)=\left(\sum_i(a_ib_i\otimes1)\right)(1\otimes\vphi(s))^{-1},
\]
onde $s\not\in\mathfrak{j}_y$, e aplica-se a Proposição~\sref{I.1.2.4}.
\end{proof}

\begin{env}[3.6.2]
\label{I.3.6.2}
Em todo o restante deste Tratado, quando considerarmos uma fibra $f^{-1}(y)$ de um morfismo como munida de uma estrutura de $\kres(y)$-pré-esquema, \emph{tratar-se-á sempre do pré-esquema obtido transportando a estrutura de $X\times_Y\Spec(\kres(y))$ pela projeção sobre $X$}.
Também escreveremos este último produto como $X\otimes_Y\kres(y)$, ou $X\otimes_{\sh{O}_Y}\kres(y)$; mais geralmente, se $B$ é uma $\sh{O}_y$-álgebra, denotaremos por $X\otimes_Y B$ ou $X\otimes_{\sh{O}_Y}B$ o produto $X\times_Y\Spec(B)$.
\end{env}

Com a convenção precedente, resulta de \sref{I.3.5.10} que os pontos de $X$ com valores em uma extensão $K$ de $\kres(y)$ são identificados com os \emph{pontos de $f^{-1}(y)$ com valores em $K$}.

\begin{env}[3.6.3]
\label{I.3.6.3}
Sejam $f:X\to Y$ e $g:Y\to Z$ dois morfismos, e $h=g\circ f$ seu composto; para todo $z\in Z$, a fibra $h^{-1}(z)$ é um pré-esquema isomorfo a
\[
  X\times_Z\Spec(\kres(z))=(X\times_Y Y)\times_Z\Spec(\kres(z))=X\times_Y g^{-1}(z)
\]
\oldpage[I]{118}
Em particular, se $U$ é uma parte aberta de $X$, o pré-esquema induzido sobre $U\cap f^{-1}(y)$ pelo pré-esquema $f^{-1}(y)$ é isomorfo a $f^{-1}_U(y)$ (sendo $f_U$ a restrição de $f$ a $U$).
\end{env}

\begin{proposition}[3.6.4]
\label{I.3.6.4}
\emph{(``Transitividade das fibras'')}. Sejam $f:X\to Y$ e $g:Y'\to Y$ dois morfismos; ponhamos $X'=X\times_Y Y'=X_{(Y')}$ e $f'=f_{(Y')}:X'\to Y'$.
Para todo $y'\in Y'$, se pomos $y=g(y')$, o pré-esquema $f'^{-1}(y')$ é isomorfo a $f^{-1}(y)\otimes_{\kres(y)}\kres(y')$.
\end{proposition}

\begin{proof}
Com efeito, basta observar que os dois pré-esquemas $(X\otimes_Y\kres(y))\otimes_{\kres(y)}\kres(y')$ e $(X\times_Y Y')\otimes_{Y'}\kres(y')$ são ambos canonicamente isomorfos a $X\times_Y\Spec(\kres(y'))$ por \sref{I.3.3.9.1}.
\end{proof}

Em particular, se $V$ é uma vizinhança aberta de $y$ em $Y$, e se denotamos por $f_V$ a restrição de $f$ ao pré-esquema induzido sobre $f^{-1}(V)$, os pré-esquemas $f^{-1}(y)$ e $f^{-1}_V(y)$ identificam-se canonicamente.

\begin{proposition}[3.6.5]
\label{I.3.6.5}
Sejam $f:X\to Y$ um morfismo, $y$ um ponto de $Y$, $Z$ o pré-esquema local $\Spec(\sh{O}_y)$, e $p=(\psi,\theta)$ a projeção $X\times_Y Z\to X$; então $p$ é um homeomorfismo do espaço subjacente a $X\times_Y Z$ sobre o subespaço $f^{-1}(Z)$ de $X$ (\emph{quando o espaço subjacente a $Z$ é identificado com um subespaço de $Y$, cf.~\sref{I.2.4.2}}), e, para todo $t\in X\times_Y Z$, pondo $z=\psi(t)$, $\theta_t^\sharp$ é um isomorfismo de $\sh{O}_z$ sobre $\sh{O}_t$.
\end{proposition}

\begin{proof}
Como $Z$ (identificado com um subespaço de $Y$) está contido em todo aberto afim que contém $y$ \sref{I.2.4.2}, podemos, como em \sref{I.3.6.1}, reduzir-nos ao caso em que $X=\Spec(A)$ e $Y=\Spec(B)$ são esquemas afins, sendo $A$ uma $B$-álgebra.
Então $X\times_Y Z$ é o espectro primo de $A\otimes_B B_y$, e esse anel identifica-se canonicamente com $S^{-1}A$, onde $S$ é a imagem de $B\setmin\mathfrak{j}_y$ em $A$ \sref[0]{0.1.5.2}; como $p$ corresponde então ao homomorfismo canônico $A\to S^{-1}A$, a proposição resulta de \sref{I.1.6.2}.
\end{proof}

\subsection{Aplicação: redução de um pré-esquema módulo~$\mathfrak{J}$}
\label{subsection:I.3.7}
\label{I.3.7}

\emph{Esta seção, que faz uso de noções e resultados posteriores do Capítulo~I e do Capítulo~II, não será utilizada no restante deste Tratado e destina-se unicamente aos leitores familiarizados com a geometria algébrica clássica}.

\begin{env}[3.7.1]
\label{I.3.7.1}
Sejam $A$ um anel, $X$ um $A$-pré-esquema e $\mathfrak{J}$ um ideal de $A$; então $X_0=X\otimes_A(A/\mathfrak{J})$ é um $(A/\mathfrak{J})$-pré-esquema, do qual às vezes se diz que é obtido de $X$ por \emph{redução} módulo~$\mathfrak{J}$.
\end{env}

\begin{env}[3.7.2]
\label{I.3.7.2}
Esta terminologia é utilizada sobretudo quando $A$ é um \emph{anel local} e $\mathfrak{J}$ seu ideal maximal, de modo que $X_0$ é um pré-esquema sobre o corpo residual $k=A/\mathfrak{J}$ de $A$.

Quando, além disso, $A$ é íntegro, com corpo de frações $K$, podemos considerar o $K$-pré-esquema $X'=X\otimes_A K$.
Por um abuso de linguagem de que não faremos uso, costumava-se até agora dizer que $X_0$ é \emph{deduzido de $X'$} por redução módulo~$\mathfrak{J}$.
Nos casos em que se utilizava essa linguagem, $A$ era um anel local de dimensão~$1$ (na maioria das vezes, um anel de valuação discreta), e ficava subentendido (de maneira mais ou menos explícita) que o $K$-pré-esquema dado $X'$ era um subpré-esquema fechado de um $K$-pré-esquema $P'$ (de fato, um espaço projetivo do tipo $\bb{P}_K^r$, cf.~\sref[II]{II.4.1.1}), por sua vez da forma $P'=P\otimes_A K$, onde $P$ é um $A$-pré-esquema dado (de fato, o $A$-esquema $\bb{P}_A^r$, com as notações de \sref[II]{II.4.1.1}).
Em nossa linguagem, a definição de $X_0$ a partir de $X'$ formula-se como segue:

Consideremos o esquema afim $Y=\Spec(A)$, formado por dois pontos, o único ponto fechado
\oldpage[I]{119}
$y=\mathfrak{J}$ e o ponto genérico $(0)$; o conjunto $U$, constituído unicamente pelo ponto genérico, é, portanto, um aberto $U=\Spec(K)$ de $Y$.
Se $X$ é um $A$-pré-esquema (ou, dito de outro modo, um $Y$-pré-esquema), então $X\otimes_A K=X'$ nada mais é do que o pré-esquema induzido por $X$ sobre $\psi^{-1}(U)$, denotando por $\psi$ o morfismo estrutural $X\to Y$.
Em particular, se $\vphi$ é o morfismo estrutural $P\to Y$, um subpré-esquema fechado $X'$ de $P'=\vphi^{-1}(U)$ é um subpré-esquema (localmente fechado) de $P$.
Se $P$ é noetheriano (por exemplo, se $A$ é noetheriano e $P$ é de tipo finito sobre $A$), existe um menor subpré-esquema fechado $X=\overline{X'}$ de $P$ que contém $X'$ \sref{I.9.5.10}, e $X'$ é o pré-esquema induzido por $X$ sobre o aberto $\vphi^{-1}(U)\cap X$, sendo, portanto, isomorfo a $X\otimes_A K$ \sref{I.9.5.10}.
\emph{A imersão de $X'$ em $P'=P\otimes_A K$ permite, portanto, considerar canonicamente $X'$ como sendo da forma $X'=X\otimes_A K$, onde $X$ é um $A$-pré-esquema.}
Pode-se então considerar o pré-esquema reduzido módulo~$\mathfrak{J}$, $X_0=X\otimes_A k$, que, aliás, nada mais é do que a fibra $\psi^{-1}(y)$ do ponto fechado $y$.
Até agora, por falta de uma terminologia adequada, havíamos evitado introduzir explicitamente o $A$-pré-esquema $X$.
Convém, contudo, observar que todas as afirmações feitas habitualmente sobre o pré-esquema «reduzido módulo~$\mathfrak{J}$» $X_0$ devem ser consideradas consequências de afirmações mais completas relativas ao próprio $X$, e só podem ser formuladas e compreendidas satisfatoriamente se interpretadas desse modo.
Parece, além disso, que todas as hipóteses feitas sobre $X_0$ se reduzem sempre a hipóteses sobre o próprio $X$ (independentemente do dado prévio de uma imersão de $X'$ em $\bb{P}_K^r$), o que permite formular enunciados mais intrínsecos.
\end{env}

\begin{env}[3.7.3]
\label{I.3.7.3}
Assinalemos, por fim, um fato muito particular, que sem dúvida contribuiu para
retardar a clarificação conceitual da situação aqui considerada: se $A$ é um
anel de valuação discreta e se $X$ é \emph{próprio} sobre $A$ (o que de fato
ocorre se $X$ é um subpré-esquema fechado de algum $\bb{P}_A^r$,
cf.~\sref[II]{II.5.5.4}), então os pontos de $X$ com valores em $A$ e os
pontos de $X'$ com valores em $K$ estão em correspondência biunívoca
\sref[II]{II.7.3.8}.
Por essa razão, muitas vezes se acreditou demonstrar resultados relativos a
$X'$, quando, na realidade, se demonstravam enunciados relativos a $X$; estes
permanecem válidos (sob essa forma) quando já não se supõe que o anel local de
base tenha dimensão~$1$.
\end{env}

\section{Subpré-esquemas e morfismos de imersão}
\label{section:I.4}

\subsection{Subpré-esquemas}
\label{subsection:I.4.1}

\begin{env}[4.1.1]
\label{I.4.1.1}
Como a noção de feixe quase-coerente \sref[0]{0.5.1.3} é local, um $\sh{O}_X$-módulo quase-coerente $\sh{F}$ sobre um pré-esquema $X$ pode ser definido pela condição de que, para todo aberto afim $V$ de $X$, $\sh{F}|V$ é isomorfo a um feixe associado a um $\Gamma(V,\sh{O}_X)$-módulo \sref{I.1.4.1}.
É claro que, sobre um pré-esquema $X$, o feixe estrutural $\sh{O}_X$ é quase-coerente, e que os núcleos, conúcleos e imagens de homomorfismos de $\sh{O}_X$-módulos quase-coerentes, bem como os limites indutivos e as somas diretas de $\sh{O}_X$-módulos quase-coerentes, são também quase-coerentes (Teorema~\sref{I.1.3.7} e Corolário~\sref{I.1.3.9}).
\end{env}

\begin{proposition}[4.1.2]
\label{I.4.1.2}
Sejam $X$ um pré-esquema e $\sh{I}$ um feixe quase-coerente de ideais de $\sh{O}_X$.
Então o suporte $Y$ do feixe $\sh{O}_X/\sh{I}$ é fechado e, se denotarmos por $\sh{O}_Y$ a restrição de $\sh{O}_X/\sh{I}$ a $Y$, então $(Y,\sh{O}_Y)$ é um pré-esquema.
\end{proposition}

\begin{proof}
\oldpage[I]{120}
Basta evidentemente \sref{I.2.1.3} considerar o caso em que $X$ é um esquema afim e mostrar que, nesse caso, $Y$ é fechado em $X$ e é um \emph{esquema afim}.
Com efeito, se $X=\Spec(A)$, então $\sh{O}_X=\widetilde{A}$ e $\sh{I}=\widetilde{\mathfrak{I}}$, onde $\mathfrak{I}$ é um ideal de $A$ \sref{I.1.4.1};
$Y$ é então igual à parte fechada $V(\mathfrak{I})$ de $X$ e se identifica com o espectro primo do anel $B=A/\mathfrak{I}$ \sref{I.1.1.11};
além disso, se $\vphi$ é o homomorfismo canônico $A\to B=A/\mathfrak{I}$, a imagem direta ${}^a\vphi_*(\widetilde{B})$ identifica-se canonicamente com o feixe $\widetilde{A}/\widetilde{\mathfrak{I}}=\sh{O}_X/\sh{I}$ (Proposição~\sref{I.1.6.3} e Corolário~\sref{I.1.3.9}), o que conclui a demonstração.
\end{proof}

Dizemos que $(Y,\sh{O}_Y)$ é o \emph{subpré-esquema} de $(X,\sh{O}_X)$ \emph{definido pelo feixe de ideais $\sh{I}$};
este é um caso particular da noção geral de \emph{subpré-esquema}:

\begin{definition}[4.1.3]
\label{I.4.1.3}
Diz-se que um espaço anelado $(Y,\sh{O}_Y)$ é um subpré-esquema de um pré-esquema $(X,\sh{O}_X)$ se:
\begin{enumerate}
  \item[1.º] $Y$ é um subespaço localmente fechado de $X$;
  \item[2.º] se $U$ denota o maior aberto de $X$ que contém $Y$ e tal que $Y$ seja fechado em $U$ (\emph{equivalentemente}, o complementar em $X$ da fronteira de $Y$ relativamente a $\overline{Y}$), então $(Y,\sh{O}_Y)$ é um subpré-esquema de $(U,\sh{O}_X|U)$ definido por um feixe quase-coerente de ideais de $\sh{O}_X|U$.
\end{enumerate}
Diz-se que o subpré-esquema $(Y,\sh{O}_Y)$ de $(X,\sh{O}_X)$ é \emph{fechado} se $Y$ é fechado em $X$ (caso em que $U=X$).
\end{definition}

Resulta imediatamente desta definição e da Proposição~\sref{I.4.1.2} que os subpré-esquemas fechados de $X$ estão em \emph{correspondência biunívoca canônica} com os feixes quase-coerentes de ideais $\sh{J}$ de $\sh{O}_X$, pois, se dois desses feixes $\sh{J}$ e $\sh{J}'$ tiverem o mesmo suporte (fechado) $Y$, e se as restrições de $\sh{O}_X/\sh{J}$ e de $\sh{O}_X/\sh{J}'$ a $Y$ forem idênticas, então $\sh{J}'=\sh{J}$.

\begin{env}[4.1.4]
\label{I.4.1.4}
Sejam $(Y,\sh{O}_Y)$ um subpré-esquema de $X$, $U$ o maior aberto de $X$ tal que $Y$ seja um subespaço fechado de $U$, e $V$ um aberto de $X$ contido em $U$;
então $V\cap Y$ é fechado em $V$.
Além disso, se $Y$ é definido pelo feixe quase-coerente de ideais $\sh{J}$ de $\sh{O}_X|U$, então $\sh{J}|V$ é um feixe quase-coerente de ideais de $\sh{O}_X|V$, e é imediato que o pré-esquema induzido por $Y$ sobre $Y\cap V$ é o subpré-esquema fechado de $V$ definido pelo feixe de ideais $\sh{J}|V$.
Reciprocamente:
\end{env}

\begin{proposition}[4.1.5]
\label{I.4.1.5}
Seja $(Y,\sh{O}_Y)$ um espaço anelado tal que $Y$ seja um subespaço de $X$ e que exista uma cobertura $(V_\alpha)$ de $Y$ por abertos de $X$ tais que, para todo $\alpha$, $Y\cap V_\alpha$ seja fechado em $V_\alpha$, e que o espaço anelado $(Y\cap V_\alpha,\sh{O}_Y|(Y\cap V_\alpha))$ seja um subpré-esquema fechado do pré-esquema induzido sobre $V_\alpha$ por $X$.
Então $(Y,\sh{O}_Y)$ é um subpré-esquema de $X$.
\end{proposition}

\begin{proof}
A hipótese implica que $Y$ é localmente fechado em $X$ e que o maior aberto $U$ no qual $Y$ é fechado contém todos os $V_\alpha$;
podemos, portanto, reduzir-nos ao caso em que $U=X$ e $Y$ é fechado em $X$.
Definimos então um feixe quase-coerente de ideais $\sh{J}$ de $\sh{O}_X$ tomando como $\sh{J}|V_\alpha$ o feixe de ideais de $\sh{O}_X|V_\alpha$ que define o subpré-esquema fechado $(Y\cap V_\alpha,\sh{O}_Y|(Y\cap V_\alpha))$ e, para todo aberto $W$ de $X$ disjunto de $Y$, tomando $\sh{J}|W=\sh{O}_X|W$.
Verifica-se imediatamente, pela Definição~\sref{I.4.1.3} e por \sref{I.4.1.4}, que existe um único feixe de ideais $\sh{J}$ que satisfaz essas condições e que ele define o subpré-esquema fechado $(Y,\sh{O}_Y)$.
\end{proof}

Em particular, o pré-esquema \emph{induzido} (por $X$) sobre um \emph{aberto} de $X$ é um \emph{subpré-esquema} de $X$.

\begin{proposition}[4.1.6]
\label{I.4.1.6}
Um subpré-esquema (resp. um subpré-esquema fechado) de um subpré-esquema
\oldpage[I]{121}
(resp. um subpré-esquema fechado) de $X$ identifica-se canonicamente com um subpré-esquema (resp. um subpré-esquema fechado) de $X$.
\end{proposition}

\begin{proof}
Como uma parte localmente fechada de um subespaço localmente fechado de $X$ é um subespaço localmente fechado de $X$, fica claro, por \sref{I.4.1.5}, que a questão é local e que podemos, portanto, supor $X$ afim;
a proposição resulta então da identificação canônica de $A/\mathfrak{J}'$ com $(A/\mathfrak{J})/(\mathfrak{J}'/\mathfrak{J})$, onde $A$ é um anel, e $\mathfrak{J}$ e $\mathfrak{J}'$ são ideais de $A$ tais que $\mathfrak{J}\subset\mathfrak{J}'$.
\end{proof}

Doravante, faremos sempre a identificação precedente.

\begin{env}[4.1.7]
\label{I.4.1.7}
Seja $Y$ um subpré-esquema de um pré-esquema $X$, e denotemos por $\psi$ a injeção canônica $Y\to X$ dos \emph{espaços subjacentes};
sabemos que a imagem inversa $\psi^*(\sh{O}_X)$ é a restrição $\sh{O}_X|Y$ \sref[0]{0.3.7.1}.
Se, para todo $y\in Y$, denotamos por $\omega_y$ o homomorfismo canônico $(\sh{O}_X)_y\to(\sh{O}_Y)_y$, esses homomorfismos são os induzidos nos talos por um homomorfismo \emph{sobrejetivo} $\omega$ de feixes de anéis $\sh{O}_X|Y\to\sh{O}_Y$: com efeito, basta verificá-lo localmente sobre $Y$, isto é, podemos supor que $X$ seja afim e que o subpré-esquema $Y$ seja fechado; se, nesse caso, $\sh{I}$ é o feixe de ideais de $\sh{O}_X$ que define $Y$, os $\omega_y$ não são senão os homomorfismos induzidos nos talos por $\sh{O}_X|Y\to(\sh{O}_X/\sh{I})|Y$.
Definimos assim um \emph{monomorfismo de espaços anelados} \sref[0]{0.4.1.1} $j=(\psi,\omega^\flat)$, que é evidentemente um morfismo $Y\to X$ de pré-esquemas \sref{I.2.2.1}, e que chamaremos \emph{morfismo de injeção canônico}.

Se $f:X\to Z$ é um morfismo, diremos ainda que o morfismo composto $Y\xrightarrow{j}X\xrightarrow{f}Z$ é a \emph{restrição} de $f$ ao subpré-esquema $Y$.
\end{env}

\begin{env}[4.1.8]
\label{I.4.1.8}
De acordo com as definições gerais (T, I, 1.1), diremos que um morfismo de pré-esquemas $f:Z\to X$ é \emph{dominado}\footnote{\emph{[Trad.] Traduzimos aqui o francês «majoré» por «dominado». Trata-se da relação de pré-ordem definida pela condição de que $f$ se fatore por $j$, e não de uma cota numérica.}} pelo morfismo de injeção $j:Y\to X$ de um subpré-esquema $Y$ de $X$ se $f$ se fatora em $Z\xrightarrow{g}Y\xrightarrow{j}X$, onde $g$ é um morfismo de pré-esquemas;
além disso, $g$ é necessariamente \emph{único}, pois $j$ é um monomorfismo.
\end{env}

\begin{proposition}[4.1.9]
\label{I.4.1.9}
Para que um morfismo $f:Z\to X$ seja dominado por um morfismo de injeção $j:Y\to X$, é necessário e suficiente que $f(Z)\subset Y$ e que, para todo $z\in Z$, pondo $y=f(z)$, o homomorfismo $(\sh{O}_X)_y\to\sh{O}_z$ correspondente a $f$ se fatore em $(\sh{O}_X)_y\to(\sh{O}_Y)_y\to\sh{O}_z$ (\emph{ou, equivalentemente,} que o núcleo de $(\sh{O}_X)_y\to\sh{O}_z$ contenha o de $(\sh{O}_X)_y\to(\sh{O}_Y)_y$).
\end{proposition}

\begin{proof}
As condições são evidentemente necessárias.
Para ver que são suficientes, podemos reduzir-nos ao caso em que $Y$ é um subpré-esquema \emph{fechado} de $X$, substituindo, se necessário, $X$ por um aberto $U$ tal que $Y$ seja fechado em $U$ \sref{I.4.1.3};
$Y$ fica então definido por um feixe quase-coerente de ideais $\sh{I}$ de $\sh{O}_X$.
Seja $f=(\psi,\theta)$, e seja $\sh{J}$ o feixe de ideais de $\psi^*(\sh{O}_X)$ que é o núcleo de $\theta^\sharp:\psi^*(\sh{O}_X)\to\sh{O}_Z$;
levando em conta as propriedades do funtor $\psi^*$ \sref[0]{0.3.7.2}, as hipóteses implicam que, para todo $z\in Z$, temos $(\psi^*(\sh{I}))_z\subset\sh{J}_z$ e, por conseguinte, $\psi^*(\sh{I})\subset\sh{J}$.
Portanto, $\theta^\sharp$ se fatora em
\[
  \psi^*(\sh{O}_X)\to\psi^*(\sh{O}_X)/\psi^*(\sh{I})=\psi^*(\sh{O}_X/\sh{I})\xrightarrow{\omega}\sh{O}_Z,
\]
sendo a primeira seta o homomorfismo canônico.
Seja $\psi'$ a aplicação contínua $Z\to Y$ que coincide com $\psi$;
é claro que temos ${\psi'}^*(\sh{O}_Y)=\psi^*(\sh{O}_X/\sh{I})$;
por outro lado, $\omega$ é evidentemente um homomorfismo local, de modo que $g=(\psi',\omega^\flat)$ é um morfismo
\oldpage[I]{122}
de pré-esquemas $Z\to Y$ \sref{I.2.2.1} que, em virtude do que precede, satisfaz $f=j\circ g$, o que demonstra a proposição.
\end{proof}

\begin{corollary}[4.1.10]
\label{I.4.1.10}
Para que um morfismo de injeção $Z\to X$ seja dominado pelo morfismo de injeção $Y\to X$, é necessário e suficiente que $Z$ seja um subpré-esquema de $Y$.
\end{corollary}

Escrevemos então $Z\leq Y$, e essa relação é evidentemente uma \emph{relação de ordem} no conjunto dos subpré-esquemas de $X$.

\subsection{Morfismos de imersão}
\label{subsection:I.4.2}

\begin{definition}[4.2.1]
\label{I.4.2.1}
Diz-se que um morfismo $f:Y\to X$ é uma imersão (resp. uma imersão fechada, uma imersão aberta) se ele se fatora em $Y\xrightarrow{g}Z\xrightarrow{j}X$, onde $g$ é um isomorfismo, $Z$ é um subpré-esquema de $X$ (resp. um subpré-esquema fechado, um subpré-esquema induzido sobre um aberto), e $j$ é o morfismo de injeção.
\end{definition}

O subpré-esquema $Z$ e o isomorfismo $g$ ficam então determinados de maneira \emph{única}, pois, se $Z'$ é um segundo subpré-esquema de $X$, $j'$ é a injeção $Z'\to X$, e $g'$ um isomorfismo $Y\to Z'$ tal que $j\circ g=j'\circ g'$, então se tem $j'=j\circ g\circ{g'}^{-1}$, donde $Z'\leq Z$ \sref{I.4.1.10}, e mostra-se do mesmo modo que $Z\leq Z'$; portanto, $Z'=Z$ e, como $j$ é um monomorfismo de pré-esquemas, $g'=g$.

Diz-se que $f=j\circ g$ é a \emph{fatoração canônica} da imersão $f$, e que o subpré-esquema $Z$ e o isomorfismo $g$ estão \emph{associados a $f$}.

É claro que uma imersão é um \emph{monomorfismo} de pré-esquemas \sref{I.4.1.7} e, \emph{a fortiori}, um \emph{morfismo radicial} \sref{I.3.5.4}.

\begin{proposition}[4.2.2]
\label{I.4.2.2}
\medskip\noindent
\begin{enumerate}
  \item[{\rm(a)}] Para que um morfismo $f=(\psi,\theta):Y\to X$ seja uma imersão aberta, é necessário e suficiente que $\psi$ seja um homeomorfismo de $Y$ sobre um subconjunto aberto de $X$ e que, para todo $y\in Y$, o homomorfismo $\theta_y^\sharp:\sh{O}_{\psi(y)}\to\sh{O}_y$ seja bijetivo.
  \item[{\rm(b)}] Para que um morfismo $f=(\psi,\theta):Y\to X$ seja uma imersão (resp. uma imersão fechada), é necessário e suficiente que $\psi$ seja um homeomorfismo de $Y$ sobre um subconjunto localmente fechado (resp. fechado) de $X$ e que, para todo $y\in Y$, o homomorfismo $\theta_y^\sharp:\sh{O}_{\psi(y)}\to\sh{O}_y$ seja sobrejetivo.
\end{enumerate}
\end{proposition}

\begin{proof}
\medskip\noindent
\begin{enumerate}
  \item[(a)] As condições são evidentemente necessárias.
    Reciprocamente, se são satisfeitas, é claro que $\theta^\sharp$ é um isomorfismo de $\psi^*(\sh{O}_X)$ sobre $\sh{O}_Y$\footnote{\emph{[Trad.] A fonte francesa inverte aqui a origem e o destino. Pela definição de um morfismo de pré-esquemas e pelos homomorfismos nos talos exibidos na Proposição~\hyperref[I.4.2.2]{4.2.2}(a), $\theta^\sharp$ tem por origem $\psi^*(\sh{O}_X)$ e por destino $\sh{O}_Y$.}}, e $\psi^*(\sh{O}_X)$ é o feixe obtido por ``transporte de estrutura'' mediante $\psi^{-1}$ a partir de $\sh{O}_X|\psi(Y)$;
    donde a conclusão.

  \item[(b)] As condições são evidentemente necessárias; provemos que são suficientes.
    Consideremos primeiro o caso particular em que se supõe que $X$ é um esquema afim e que $Z=\psi(Y)$ é \emph{fechado} em $X$.
    Sabemos então \sref[0]{0.3.4.6} que $\psi_*(\sh{O}_Y)$ tem suporte $Z$, e que sua restrição a $Z$, denotada por $\sh{O}_Z'$, faz com que o espaço anelado $(Z,\sh{O}_Z')$ seja obtido de $(Y,\sh{O}_Y)$ por transporte de estrutura mediante o homeomorfismo $\psi$, considerado como aplicação de $Y$ sobre $Z$.
    Mostremos agora que $f_*(\sh{O}_Y)=\psi_*(\sh{O}_Y)$ é um $\sh{O}_X$-módulo \emph{quase-coerente}.
    Com efeito, para todo $x\not\in Z$, $\psi_*(\sh{O}_Y)$, restrito a uma vizinhança adequada de $x$, é nulo.
    Se, ao contrário, $x\in Z$, então $x=\psi(y)$ para um único $y\in Y$;
    seja $V$ uma vizinhança aberta afim de $y$ em $Y$;
    $\psi(V)$ é então aberto em $Z$ e, portanto, igual à interseção de $Z$ com um aberto $U$ de $X$, e a restrição a $U$ de $\psi_*(\sh{O}_Y)$ identifica-se com a restrição a $U$ da imagem direta
\oldpage[I]{123}
    $(\psi_V)_*(\sh{O}_Y|V)$, onde $\psi_V$ é a restrição de $\psi$ a $V$.
    A restrição do morfismo $(\psi,\theta)$ a $(V,\sh{O}_Y|V)$ é um morfismo deste pré-esquema para $(X,\sh{O}_X)$ e, por conseguinte, é da forma $({}^a\vphi,\widetilde{\vphi})$, onde $\vphi$ é o homomorfismo do anel $A=\Gamma(X,\sh{O}_X)$ para o anel $\Gamma(V,\sh{O}_Y)$ \sref{I.1.7.3};
    deduzimos que $(\psi_V)_*(\sh{O}_Y|V)$ é um $\sh{O}_X$-módulo quase-coerente \sref{I.1.6.3}, o que prova nossa afirmação, em razão do caráter local dos feixes quase-coerentes.

    Além disso, a hipótese de que $\psi$ é um homeomorfismo implica \sref[0]{0.3.4.5} que, para todo $y\in Y$, $\psi_y$ é um isomorfismo $(\psi_*(\sh{O}_Y))_{\psi(y)}\to\sh{O}_y$;
    como o diagrama
    \[
      \xymatrix{
        \sh{O}_{\psi(y)}\ar[r]^{\theta_{\psi(y)}}\ar[d]_{\psi_y\circ\rho_{\psi(y)}} &
        (\psi_*(\sh{O}_Y))_{\psi(y)}\ar[d]^{\psi_y}\\
        (\psi^*(\sh{O}_X))_y\ar[r]^{\theta_y^\sharp} &
        \sh{O}_y
      }
    \]
    é comutativo e as flechas verticais são os isomorfismos \sref[0]{0.3.7.2}, a hipótese de que $\theta_y^\sharp$ é sobrejetivo implica que $\theta_{\psi(y)}$ também é sobrejetivo.
    Como o suporte de $\psi_*(\sh{O}_Y)$ é $Z=\psi(Y)$, $\theta$ é um homomorfismo \emph{sobrejetivo} de $\sh{O}_X=\widetilde{A}$ sobre o $\sh{O}_X$-módulo quase-coerente $f_*(\sh{O}_Y)$.
    Por conseguinte, existe um único isomorfismo $\omega$ de um feixe quociente $\widetilde{A}/\widetilde{\mathfrak{J}}$ (onde $\mathfrak{J}$ é um ideal de $A$) sobre $f_*(\sh{O}_Y)$ que, composto com o homomorfismo canônico $\widetilde{A}\to\widetilde{A}/\widetilde{\mathfrak{J}}$, dá $\theta$ \sref{I.1.3.8};
    se $\sh{O}_Z$ denota a restrição de $\widetilde{A}/\widetilde{\mathfrak{J}}$ a $Z$, então $(Z,\sh{O}_Z)$ é um subpré-esquema de $(X,\sh{O}_X)$, e $f$ se fatora pela injeção canônica deste subpré-esquema em $X$ e pelo isomorfismo $(\psi_0,\omega_0)$, onde $\psi_0$ é $\psi$ considerado como aplicação de $Y$ sobre $Z$, e $\omega_0$ é a restrição de $\omega$ a $\sh{O}_Z$.

    Passemos agora ao caso geral.
    Seja $U$ um aberto afim de $X$ tal que $U\cap\psi(Y)$ seja fechado em $U$ e não vazio.
    Restringindo $f$ ao pré-esquema induzido por $Y$ sobre o aberto $\psi^{-1}(U)$ e considerando-o como um morfismo deste pré-esquema para o pré-esquema induzido por $X$ sobre $U$, recaímos no primeiro caso;
    a restrição de $f$ a $\psi^{-1}(U)$ é, portanto, uma imersão fechada $\psi^{-1}(U)\to U$, que se fatora canonicamente como $j_U\circ g_U$, onde $g_U$ é um isomorfismo do pré-esquema $\psi^{-1}(U)$ sobre um subpré-esquema $Z_U$ de $U$, e $j_U$ é a injeção canônica $Z_U\to U$.
    Seja $V$ outro aberto afim de $X$ tal que $V\subset U$;
    como a restrição $Z_V'$ de $Z_U$ a $V$ é um subpré-esquema do pré-esquema $V$, a restrição de $f$ a $\psi^{-1}(V)$ se fatora como $j_V'\circ g_V'$, onde $j_V'$ é a injeção canônica $Z_V'\to V$ e $g_V'$ é um isomorfismo de $\psi^{-1}(V)$ sobre $Z_V'$.
    Pela unicidade da fatoração canônica de uma imersão \sref{I.4.2.1}, tem-se necessariamente $Z_V'=Z_V$ e $g_V'=g_V$.
    Deduz-se \sref{I.4.1.5} que existe um subpré-esquema $Z$ de $X$ cujo espaço subjacente é $\psi(Y)$ e cuja restrição a cada $U\cap\psi(Y)$ é $Z_U$;
    os $g_U$ são então as restrições a $\psi^{-1}(U)$ de um isomorfismo $g:Y\to Z$ tal que $f=j\circ g$, onde $j$ é a injeção canônica $Z\to X$.
\end{enumerate}
\end{proof}

\begin{corollary}[4.2.3]
\label{I.4.2.3}
Seja $X$ um esquema afim.
Para que um morfismo $f=(\psi,\theta):Y\to X$ seja uma imersão fechada, é necessário e suficiente que $Y$ seja um esquema afim e que o homomorfismo $\Gamma(\theta)$\footnote{\emph{[Trad.] A fonte francesa imprime $\Gamma(\psi)$. Como $\psi$ é a aplicação contínua e $\theta$ o homomorfismo de feixes, o homomorfismo induzido nas seções globais é $\Gamma(\theta)$.}}$:\Gamma(\sh{O}_X)\to\Gamma(\sh{O}_Y)$ seja sobrejetivo.
\end{corollary}

\begin{corollary}[4.2.4]
\label{I.4.2.4}
\medskip\noindent
\begin{enumerate}
  \item[{\rm(a)}] Sejam $f$ um morfismo $Y\to X$ e $(V_\lambda)$ uma cobertura de $f(Y)$ por abertos de $X$.
    Para que $f$ seja uma imersão (resp. uma imersão aberta), é necessário e suficiente
\oldpage[I]{124}
    que sua restrição a cada um dos pré-esquemas induzidos $f^{-1}(V_\lambda)$ seja uma imersão (resp. uma imersão aberta) em $V_\lambda$.
  \item[{\rm(b)}] Sejam $f$ um morfismo $Y\to X$ e $(V_\lambda)$ uma cobertura aberta de $X$.
    Para que $f$ seja uma imersão fechada, é necessário e suficiente que sua restrição a cada um dos pré-esquemas induzidos $f^{-1}(V_\lambda)$ seja uma imersão fechada em $V_\lambda$.
\end{enumerate}
\end{corollary}

\begin{proof}
Seja $f=(\psi,\theta)$;
no caso (a), $\theta_y^\sharp$ é sobrejetivo (resp. bijetivo) para todo $y\in Y$, e, no caso (b), $\theta_y^\sharp$ é sobrejetivo para todo $y\in Y$;
basta, portanto, verificar, no caso (a), que $\psi$ é um homeomorfismo de $Y$ sobre um subconjunto localmente fechado (resp. aberto) de $X$ e, no caso (b), que $\psi$ é um homeomorfismo de $Y$ sobre um subconjunto fechado de $X$.
Ora, $\psi$ é evidentemente injetiva e transforma toda vizinhança de $y$ em $Y$ em uma vizinhança de $\psi(y)$ em $\psi(Y)$, para todo $y\in Y$, em virtude da hipótese;
no caso (a), $\psi(Y)\cap V_\lambda$ é localmente fechado (resp. aberto) em $V_\lambda$, logo $\psi(Y)$ é localmente fechado (resp. aberto) na união dos $V_\lambda$, e \emph{a fortiori} em $X$;
no caso (b), $\psi(Y)\cap V_\lambda$ é fechado em $V_\lambda$, logo $\psi(Y)$ é fechado em $X$, pois $X=\bigcup_\lambda V_\lambda$.
\end{proof}

\begin{proposition}[4.2.5]
\label{I.4.2.5}
O composto de duas imersões quaisquer (resp. de duas imersões abertas, de duas imersões fechadas) é uma imersão (resp. uma imersão aberta, uma imersão fechada).
\end{proposition}

\begin{proof}
Isso resulta trivialmente de \sref{I.4.1.6}.
\end{proof}

\subsection{Produto de imersões}
\label{subsection:I.4.3}
\label{I.4.3}

\begin{proposition}[4.3.1]
\label{I.4.3.1}
Sejam $\alpha:X'\to X$, $\beta:Y'\to Y$ dois $S$-morfismos; se $\alpha$ e $\beta$ são imersões (resp. imersões abertas, imersões fechadas), então $\alpha\times_S\beta$ é uma imersão (resp. uma imersão aberta, uma imersão fechada).
Além disso, se $\alpha$ (resp.~$\beta$) identifica $X'$ (resp.~$Y'$) com um subpré-esquema $X''$ (resp.~$Y''$) de $X$ (resp.~$Y$), então $\alpha\times_S\beta$ identifica o espaço subjacente a $X'\times_S Y'$ com o subespaço $p^{-1}(X'')\cap q^{-1}(Y'')$ do espaço subjacente a $X\times_S Y$, onde $p$ e $q$ denotam as projeções de $X\times_S Y$ sobre $X$ e $Y$, respectivamente.
\end{proposition}

\begin{proof}
Em vista da Definição~\sref{I.4.2.1}, podemos restringir-nos ao caso em que $X'$ e $Y'$ são subpré-esquemas, e $\alpha$ e $\beta$ são os morfismos de injeção.
A proposição já foi estabelecida para os subpré-esquemas induzidos sobre os abertos \sref{I.3.2.7};
como todo subpré-esquema é um subpré-esquema fechado de um pré-esquema induzido sobre um aberto \sref{I.4.1.3}, podemos reduzir-nos ao caso em que $X'$ e $Y'$ são subpré-esquemas \emph{fechados}.

Demonstremos primeiro que podemos supor que $S$ é \emph{afim}.
Com efeito, seja $(S_\lambda)$ uma cobertura de $S$ formada por abertos afins;
se $\vphi$ e $\psi$ são os morfismos estruturais de $X$ e $Y$, sejam $X_\lambda=\vphi^{-1}(S_\lambda)$ e $Y_\lambda=\psi^{-1}(S_\lambda)$.
A restrição $X_\lambda'$ (resp. $Y_\lambda'$) de $X'$ (resp. $Y'$) a $X_\lambda\cap X'$ (resp. $Y_\lambda\cap Y'$) é um subpré-esquema fechado de $X_\lambda$ (resp. $Y_\lambda$); os pré-esquemas $X_\lambda$, $Y_\lambda$, $X_\lambda'$ e $Y_\lambda'$ podem ser considerados como $S_\lambda$-pré-esquemas, e os produtos $X_\lambda\times_S Y_\lambda$ e $X_\lambda\times_{S_\lambda}Y_\lambda$ (resp. $X_\lambda'\times_S Y_\lambda'$ e $X_\lambda'\times_{S_\lambda}Y_\lambda'$) são idênticos \sref{I.3.2.5}.
Se a proposição é verdadeira quando $S$ é afim, a restrição de $\alpha\times_S\beta$ a cada um dos $X_\lambda'\times_S Y_\lambda'$ é, portanto, uma imersão \sref{I.3.2.7}.
Como o produto $X_\lambda'\times_S Y_\mu'$ (resp. $X_\lambda\times_S Y_\mu$) se identifica com $(X_\lambda'\cap X_\mu')\times_S(Y_\lambda'\cap Y_\mu')$ (resp. $(X_\lambda\cap X_\mu)\times_S(Y_\lambda\cap Y_\mu)$) \sref{I.3.2.6.4}, a restrição de $\alpha\times_S\beta$
\oldpage[I]{125}
a cada um dos $X_\lambda'\times_S Y_\mu'$ é ainda uma imersão;
o mesmo vale, portanto, para $\alpha\times_S\beta$ em virtude de \sref{I.4.2.4}.

Em segundo lugar, demonstremos que também podemos supor que $X$ e $Y$ são \emph{afins}.
Com efeito, seja $(U_i)$ (resp. $(V_j)$) uma cobertura de $X$ (resp. $Y$) por abertos afins, e seja $X_i'$ (resp. $Y_j'$) a restrição de $X'$ (resp. $Y'$) a $X'\cap U_i$ (resp. $Y'\cap V_j$), que é um subpré-esquema fechado de $U_i$ (resp. $V_j$);
$U_i\times_S V_j$ se identifica com a restrição de $X\times_S Y$ a $p^{-1}(U_i)\cap q^{-1}(V_j)$ \sref{I.3.2.7};
de modo análogo, se $p'$ e $q'$ são as projeções de $X'\times_S Y'$, então $X_i'\times_S Y_j'$ se identifica com a restrição de $X'\times_S Y'$ a ${p'}^{-1}(X_i')\cap{q'}^{-1}(Y_j')$.
Ponhamos $\gamma=\alpha\times_S\beta$;
tem-se, por definição, $p\circ\gamma=\alpha\circ p'$ e $q\circ\gamma=\beta\circ q'$;
como $X_i'=\alpha^{-1}(U_i)$ e $Y_j'=\beta^{-1}(V_j)$, tem-se também ${p'}^{-1}(X_i')=\gamma^{-1}(p^{-1}(U_i))$ e ${q'}^{-1}(Y_j')=\gamma^{-1}(q^{-1}(V_j))$, de onde
\[
  {p'}^{-1}(X_i')\cap{q'}^{-1}(Y_j')=\gamma^{-1}(p^{-1}(U_i)\cap q^{-1}(V_j))=\gamma^{-1}(U_i\times_S V_j),
\]
e conclui-se então como na primeira parte do raciocínio.

Suponhamos, pois, que $X$, $Y$ e $S$ são afins, e sejam $B$, $C$ e $A$ seus anéis respectivos.
Então $B$ e $C$ são $A$-álgebras, e $X'$ e $Y'$ são esquemas afins cujos anéis são as álgebras quocientes $B'$ e $C'$ de $B$ e $C$, respectivamente.
Além disso, tem-se $\alpha=({}^a\rho,\widetilde{\rho})$ e $\beta=({}^a\sigma,\widetilde{\sigma})$, onde $\rho$ e $\sigma$ são, respectivamente, os homomorfismos canônicos $B\to B'$ e $C\to C'$ \sref{I.1.7.3}.
Nessas condições, sabe-se que $X\times_S Y$ (resp. $X'\times_S Y'$) é um esquema afim de anel $B\otimes_A C$ (resp. $B'\otimes_A C'$), e $\alpha\times_S\beta=({}^a\tau,\widetilde{\tau})$, onde $\tau$ é o homomorfismo $\rho\otimes\sigma$ de $B\otimes_A C$ em $B'\otimes_A C'$ (Proposição~\sref{I.3.2.2} e Corolário~\sref{I.3.2.3});
como esse homomorfismo é sobrejetivo, $\alpha\times_S\beta$ é uma imersão.
Além disso, se $\mathfrak{b}$ (resp. $\mathfrak{c}$) é o núcleo de $\rho$ (resp. $\sigma$), então o núcleo de $\tau$ é
$\operatorname{Im}(\mathfrak{b}\otimes_A C\to B\otimes_A C)+\operatorname{Im}(B\otimes_A\mathfrak{c}\to B\otimes_A C)$,
ou, equivalentemente, a soma dos ideais gerados por $u(\mathfrak{b})$ e $v(\mathfrak{c})$, onde $u$ (resp. $v$) é o homomorfismo $b\mapsto b\otimes 1$ (resp. $c\mapsto 1\otimes c$).
Como $p=({}^a u,\widetilde{u})$ e $q=({}^a v,\widetilde{v})$, esse núcleo corresponde, no espectro primo de $B\otimes_A C$, ao conjunto fechado $p^{-1}(X')\cap q^{-1}(Y')$, em virtude de \sref{I.1.2.2.1} e da Proposição~\sref{I.1.1.2}[iii], o que conclui a demonstração.
\end{proof}

\begin{corollary}[4.3.2]
\label{I.4.3.2}
Se $f:X\to Y$ é uma imersão (resp. uma imersão aberta, uma imersão fechada) e um $S$-morfismo, então $f_{(S')}$ é uma imersão (resp. uma imersão aberta, uma imersão fechada) para toda extensão $S'\to S$ do pré-esquema de base.
\end{corollary}

\subsection{Imagem inversa de um subpré-esquema}
\label{subsection:I.4.4}

\begin{proposition}[4.4.1]
\label{I.4.4.1}
Sejam $f:X\to Y$ um morfismo, $Y'$ um subpré-esquema (resp. um subpré-esquema fechado, um pré-esquema induzido sobre um aberto) de $Y$, e $j:Y'\to Y$ o morfismo de injeção.
Então a projeção $p:X\times_Y Y'\to X$ é uma imersão (resp. uma imersão fechada, uma imersão aberta);
o espaço subjacente do subpré-esquema de $X$ associado a $p$ é $f^{-1}(Y')$;
além disso, se $j'$ é o morfismo de injeção desse subpré-esquema em $X$, então, para todo morfismo $h:Z\to X$, o composto $f\circ h:Z\to Y$ é dominado por $j$ se e somente se $h$ é dominado por $j'$.
\end{proposition}

\begin{proof}
Como $p=1_X\times_Y j$ \sref{I.3.3.4}, a primeira afirmação decorre da Proposição~\sref{I.4.3.1};
a segunda é um caso particular do Corolário~\sref{I.3.5.10} (trocando os papéis de $X$ e de $Y'$).
Finalmente, se $f\circ h=j\circ h'$, onde $h'$ é um morfismo $Z\to Y'$, decorre da definição do produto que $h=p\circ u$, onde $u$ é um morfismo $Z\to X\times_Y Y'$; daí segue a última afirmação.
\end{proof}

Diremos que o subpré-esquema de $X$ assim definido é a \emph{imagem inversa} do subpré-esquema $Y'$ de $Y$ pelo morfismo $f$, terminologia que está de acordo com a introduzida
\oldpage[I]{126}
de modo mais geral em \sref{I.3.3.6}.
Quando falarmos de $f^{-1}(Y')$ como de um subpré-esquema de $X$, tratar-se-á sempre deste subpré-esquema.

Quando os pré-esquemas $f^{-1}(Y')$ e $X$ são idênticos, $j'$ é a identidade e todo morfismo $h:Z\to X$ é, portanto, dominado por $j'$; tomando $h=1_X$, vemos que \emph{o morfismo $f:X\to Y$ se fatora como $X\xrightarrow{g}Y'\xrightarrow{j}Y$}.

Quando $y$ é um ponto \emph{fechado} de $Y$ e $Y'=\Spec(\kres(y))$ é o menor subpré-esquema fechado de $Y$ que tem $\{y\}$ como espaço subjacente \sref{I.4.1.9}, o subpré-esquema fechado $f^{-1}(Y')$ é canonicamente isomorfo à \emph{fibra} $f^{-1}(y)$ definida em \sref{I.3.6.2}, com a qual o identificaremos.

\begin{corollary}[4.4.2]
\label{I.4.4.2}
Sejam $f:X\to Y$ e $g:Y\to Z$ dois morfismos, e seja $h=g\circ f$ sua composição.
Para todo subpré-esquema $Z'$ de $Z$, os subpré-esquemas $f^{-1}(g^{-1}(Z'))$ e $h^{-1}(Z')$ de $X$ são idênticos.
\end{corollary}

\begin{proof}
Isso resulta da existência do isomorfismo canônico $X\times_Y(Y\times_Z Z')\isoto X\times_Z Z'$ \sref{I.3.3.9.1}.
\end{proof}

\begin{corollary}[4.4.3]
\label{I.4.4.3}
Sejam $X'$ e $X''$ dois subpré-esquemas de $X$, e sejam $j':X'\to X$ e $j'':X''\to X$ seus morfismos de injeção;
então ${j'}^{-1}(X'')$ e ${j''}^{-1}(X')$ são ambos iguais ao ínfimo $\inf(X',X'')$ de $X'$ e $X''$ para a relação de ordem $\leq$ entre subpré-esquemas, e este é canonicamente isomorfo a $X'\times_X X''$.
\end{corollary}

\begin{proof}
Isso resulta imediatamente da Proposição~\sref{I.4.4.1} e do Corolário~\sref{I.4.1.10}.
\end{proof}

\begin{corollary}[4.4.4]
\label{I.4.4.4}
Seja $f:X\to Y$ um morfismo, e sejam $Y'$ e $Y''$ dois subpré-esquemas de $Y$;
então
\[
  f^{-1}(\inf(Y',Y''))=\inf(f^{-1}(Y'),f^{-1}(Y'')).
\]
\end{corollary}

\begin{proof}
Isso resulta da existência do isomorfismo canônico entre $(X\times_Y Y')\times_X(X\times_Y Y'')$ e $X\times_Y(Y'\times_Y Y'')$ \sref{I.3.3.9.1}.
\end{proof}

\begin{proposition}[4.4.5]
\label{I.4.4.5}
Sejam $f:X\to Y$ um morfismo e $Y'$ um subpré-esquema fechado de $Y$ definido por um feixe quase-coerente $\sh{K}$ de ideais de $\sh{O}_Y$ \sref{I.4.1.3};
então o subpré-esquema fechado $f^{-1}(Y')$ de $X$ é definido pelo feixe quase-coerente de ideais $f^*(\sh{K})\sh{O}_X$ de $\sh{O}_X$.
\end{proposition}

\begin{proof}
A questão é evidentemente local em $X$ e $Y$;
basta, então, observar que, se \erratum[II]{$A$ é uma $B$-álgebra} e $\mathfrak{K}$ é um ideal de $B$, então se tem $A\otimes_B(B/\mathfrak{K})=A/\mathfrak{K}A$, e aplicar \sref{I.1.6.9}.
\end{proof}

\begin{corollary}[4.4.6]
\label{I.4.4.6}
Sejam $X'$ um subpré-esquema fechado de $X$ definido por um feixe quase-coerente de ideais $\sh{J}$ de $\sh{O}_X$, e $i$ a injeção $X'\to X$;
para que a restrição $f\circ i$ de $f$ a $X'$ seja dominada pela injeção $j:Y'\to Y$ \emph{(em outras palavras, se fatore como $j\circ g$, onde $g$ é um morfismo $X'\to Y'$)}, é necessário e suficiente que $f^*(\sh{K})\sh{O}_X\subset\sh{J}$.
\end{corollary}

\begin{proof}
Basta aplicar a Proposição~\sref{I.4.4.1} a $i$, levando em conta a Proposição~\sref{I.4.4.5}.
\end{proof}

\subsection{Imersões locais e isomorfismos locais}
\label{subsection:I.4.5}

\begin{definition}[4.5.1]
\label{I.4.5.1}
Seja $f:X\to Y$ um morfismo de pré-esquemas.
Diz-se que $f$ é uma \emph{imersão local em um ponto $x\in X$} se existem uma vizinhança aberta $U$ de $x$ em $X$ e uma vizinhança aberta $V$ de $f(x)$ em $Y$ tais que a restrição de $f$ ao pré-esquema induzido $U$ seja uma imersão fechada de $U$ no pré-esquema induzido $V$.
Diz-se que $f$ é uma imersão local se $f$ é uma imersão local em todo ponto de $X$.
\end{definition}

\begin{definition}[4.5.2]
\label{I.4.5.2}
Diz-se que um morfismo $f:X\to Y$ é um isomorfismo local em
\oldpage[I]{127}
um ponto $x\in X$ se existe uma vizinhança aberta $U$ de $x$ em $X$ tal que a restrição de $f$ ao pré-esquema induzido $U$ seja uma imersão aberta de $U$ em $Y$.
Diz-se que $f$ é um isomorfismo local se $f$ é um isomorfismo local em todo ponto de $X$.
\end{definition}

\begin{env}[4.5.3]
\label{I.4.5.3}
Uma imersão (resp. uma imersão fechada) $f:X\to Y$ pode ser caracterizada como uma imersão local tal que $f$ seja um homeomorfismo do espaço subjacente a $X$ sobre um subconjunto (resp. um subconjunto fechado) de $Y$.
Uma imersão aberta $f$ pode ser caracterizada como um isomorfismo local \emph{injetivo}.
\end{env}

\begin{proposition}[4.5.4]
\label{I.4.5.4}
Sejam $X$ um pré-esquema irreducível e $f:X\to Y$ um morfismo injetivo dominante.
Se $f$ é uma imersão local, então $f$ é uma imersão e $f(X)$ é aberto em $Y$.
\end{proposition}

\begin{proof}
Seja $x\in X$, e sejam $U$ uma vizinhança aberta de $x$ e $V$ uma vizinhança aberta de $f(x)$ em $Y$ tais que a restrição de $f$ a $U$ seja uma imersão fechada em $V$;
como $U$ é denso em $X$, $f(U)$ é denso em $Y$ por hipótese, portanto $f(U)=V$, e $f$ é um homeomorfismo de $U$ sobre $V$;
a hipótese de que $f$ é injetivo implica que $f^{-1}(V)=U$, de onde resulta imediatamente a proposição.
\end{proof}

\begin{proposition}[4.5.5]
\label{I.4.5.5}
\medskip\noindent
\begin{enumerate}
  \item[{\rm(i)}] O composto de duas imersões locais (resp. de dois isomorfismos locais) é uma imersão local (resp. um isomorfismo local).
  \item[{\rm(ii)}] Sejam $f:X\to X'$ e $g:Y\to Y'$ dois $S$-morfismos.
    Se $f$ e $g$ são imersões locais (resp. isomorfismos locais), então $f\times_S g$ também é uma imersão local (resp. um isomorfismo local).
  \item[{\rm(iii)}] Se um $S$-morfismo $f$ é uma imersão local (resp. um isomorfismo local), então $f_{(S')}$ também o é para toda extensão $S'\to S$ do pré-esquema de base.
\end{enumerate}
\end{proposition}

\begin{proof}
Segundo \sref{I.3.5.1}, basta demonstrar (i) e (ii).

(i) resulta imediatamente da transitividade das imersões fechadas (resp. abertas) \sref{I.4.2.5}\footnote{\emph{[Trad.] A fonte francesa cita aqui (4.2.4); a transitividade é a Proposição~\hyperref[I.4.2.5]{4.2.5}.}} e do fato de que, se $f$ é um homeomorfismo de $X$ sobre um subconjunto fechado de $Y$, para todo aberto $U\subset X$, $f(U)$ é aberto em $f(X)$; portanto, existe um subconjunto aberto $V$ de $Y$ tal que $f(U)=V\cap f(X)$ e, por conseguinte, $f(U)$ é fechado em $V$.

Para demonstrar (ii), seja $z\in X\times_S Y$, ponhamos $z'=(f\times_S g)(z)$\footnote{\emph{[Trad.] A fonte francesa utiliza $z$ e $z'$ mais abaixo sem introduzi-los; esta inserção fornece o ponto necessário e sua imagem.}}, sejam $p$ e $q$ as projeções de $X\times_S Y$, e sejam $p'$ e $q'$ as projeções de $X'\times_S Y'$.
Existem, por hipótese, vizinhanças abertas $U$, $U'$, $V$ e $V'$ de $x=p(z)$, $x'=p'(z')$, $y=q(z)$ e $y'=q'(z')$, respectivamente, tais que as restrições de $f$ e $g$ a $U$ e $V$, respectivamente, sejam imersões fechadas (resp. abertas) em $U'$ e $V'$, respectivamente.
Como os espaços subjacentes de $U\times_S V$ e $U'\times_S V'$ se identificam com as vizinhanças abertas $p^{-1}(U)\cap q^{-1}(V)$ e ${p'}^{-1}(U')\cap{q'}^{-1}(V')$ de $z$ e $z'$, respectivamente, \sref{I.3.2.7}, a proposição decorre da Proposição~\sref{I.4.3.1}.
\end{proof}

\section{Pré-esquemas reduzidos; condição de separação}
\label{section:I.5}

\subsection{Pré-esquemas reduzidos}
\label{subsection:I.5.1}

\begin{proposition}[5.1.1]
\label{I.5.1.1}
Sejam $(X,\sh{O}_X)$ um pré-esquema e $\sh{B}$ uma $\sh{O}_X$-álgebra quase-coerente.
Então existe um único $\sh{O}_X$-módulo quase-coerente $\sh{N}$ cujo talo $\sh{N}_x$, em cada ponto $x\in X$, é o nilradical do anel $\sh{B}_x$.
Quando $X$ é afim e, por conseguinte, $\sh{B}=\widetilde{B}$, onde $B$ é uma álgebra sobre $A(X)$, tem-se $\sh{N}=\widetilde{\nilrad}$, onde $\nilrad$ é o nilradical de $B$.
\end{proposition}

\begin{proof}
\oldpage[I]{128}
Como a questão é local, basta demonstrar a última afirmação.
Sabemos que $\widetilde{\nilrad}$ é um $\sh{O}_X$-módulo quase-coerente \sref{I.1.4.1} e que seu talo no ponto $x\in X$ é o ideal $\nilrad_x$ do anel de frações $B_x$;
resta demonstrar que o nilradical de $B_x$ está contido em $\nilrad_x$, pois a inclusão oposta é evidente.
Seja $z/s$ um elemento do nilradical de $B_x$, com $z\in B$ e $s\not\in\mathfrak{j}_x$;
por hipótese, existe um inteiro $k$ tal que $(z/s)^k=0$, o que significa que existe $t\not\in\mathfrak{j}_x$ tal que $tz^k=0$.
Conclui-se que $(tz)^k=0$ e, por conseguinte, que $z/s=(tz)/(ts)\in\nilrad_x$.
\end{proof}

Diremos que o $\sh{O}_X$-módulo quase-coerente $\sh{N}$ assim definido é o \emph{nilradical} da $\sh{O}_X$-álgebra $\sh{B}$; em particular, denotaremos por $\sh{N}_X$ o nilradical de $\sh{O}_X$.

\begin{corollary}[5.1.2]
\label{I.5.1.2}
Seja $X$ um pré-esquema;
o subpré-esquema fechado de $X$ definido pelo feixe de ideais $\sh{N}_X$ é o único subpré-esquema reduzido \sref[0]{0.4.1.4} de $X$ que tem $X$ como espaço subjacente;
é também o menor subpré-esquema de $X$ que tem $X$ como espaço subjacente.
\end{corollary}

\begin{proof}
Seja $Y$ o subpré-esquema fechado de $X$ definido por $\sh{N}_X$.
Como seu feixe estrutural é $\sh{O}_X/\sh{N}_X$, é imediato que $Y$ é reduzido e tem $X$ como espaço subjacente, pois $\sh{N}_x\neq\sh{O}_x$ para todo $x\in X$.
Para demonstrar as demais afirmações, observemos que um subpré-esquema $Z$ de $X$ que tenha $X$ como espaço subjacente é definido por um feixe de ideais $\sh{I}$ \sref{I.4.1.3} tal que $\sh{I}_x\neq\sh{O}_x$ para todo $x\in X$.
Podemos nos restringir ao caso em que $X$ é afim, digamos $X=\Spec(A)$ e $\sh{I}=\widetilde{\mathfrak{I}}$, onde $\mathfrak{I}$ é um ideal de $A$;
então, para todo $x\in X$, tem-se $\mathfrak{I}\subset\mathfrak{j}_x$\footnote{\emph{[Trad.] A fonte francesa imprime aqui $\mathfrak{I}_x\subset\mathfrak{j}_x$, comparação mal tipada entre um ideal localizado de $A_x$ e um ideal de $A$; o argumento exige $\mathfrak{I}\subset\mathfrak{j}_x$.}}, de modo que $\mathfrak{I}$ está contido em todo ideal primo de $A$ e, portanto, também em sua interseção $\nilrad$, o nilradical de $A$.
Isso demonstra que $Y$ é o menor subpré-esquema de $X$ que tem $X$ como espaço subjacente \sref{I.4.1.9};
além disso, se $Z$ é distinto de $Y$, tem-se necessariamente $\sh{I}_x\neq\sh{N}_x$ para pelo menos um $x\in X$ e, por \sref{I.5.1.1}, $Z$ não é reduzido.
\end{proof}

\begin{definition}[5.1.3]
\label{I.5.1.3}
Chama-se pré-esquema reduzido associado a um pré-esquema $X$, e denota-se por $X_\red$, o único subpré-esquema reduzido de $X$ que tem $X$ como espaço subjacente.
\end{definition}

Dizer que um pré-esquema $X$ é reduzido significa, portanto, que $X=X_\red$.

\begin{proposition}[5.1.4]
\label{I.5.1.4}
Para que o espectro primo de um anel $A$ seja um pré-esquema reduzido (resp. integral) \sref{I.2.1.8}\footnote{\emph{[Trad.] A fonte francesa remete aqui a (2.1.7), que trata da restrição a subconjuntos abertos; a definição pertinente de pré-esquema integral é (2.1.8).}}, é necessário e suficiente que $A$ seja um anel reduzido (resp. um domínio integral).
\end{proposition}

\begin{proof}
Com efeito, de \sref{I.5.1.1} resulta imediatamente que a condição $\sh{N}=(0)$ é necessária e suficiente para que $X=\Spec(A)$ seja reduzido;
a afirmação relativa aos domínios integrais é então consequência de \sref{I.1.1.13}.
\end{proof}

Como todo anel de frações não nulo de um domínio integral é um domínio integral, de \sref{I.5.1.4} resulta que, para todo pré-esquema \emph{localmente integral} $X$, $\sh{O}_x$ é um domínio integral para todo $x\in X$.
A recíproca é verdadeira sempre que o espaço subjacente a $X$ seja \emph{localmente noetheriano}:
com efeito, $X$ é então reduzido e, se $U$ é um aberto afim de $X$ que é um espaço noetheriano, $U$ tem somente um número finito de componentes irredutíveis, de modo que seu anel $A$ tem somente um número finito de ideais primos minimais \sref{I.1.1.14}.
Se duas das componentes irredutíveis de $U$ tivessem um ponto comum $x$, então $\sh{O}_x$ teria pelo menos dois ideais primos minimais distintos e, portanto, não seria um domínio integral;
as componentes de $U$ são, pois, abertos disjuntos dois a dois, e cada uma delas é integral.

\begin{env}[5.1.5]
\label{I.5.1.5}
Seja $f=(\psi,\theta):X\to Y$ um morfismo de pré-esquemas;
\oldpage[I]{129}
o homomorfismo $\theta_x^\sharp:\sh{O}_{\psi(x)}\to\sh{O}_x$ envia todo elemento nilpotente de $\sh{O}_{\psi(x)}$ a um elemento nilpotente de $\sh{O}_x$;
passando aos quocientes, $\theta^\sharp$ induz um homomorfismo
\[
  \omega:\psi^*(\sh{O}_Y/\sh{N}_Y)\to\sh{O}_X/\sh{N}_X.
\]
É claro que, para todo $x\in X$, $\omega_x:\sh{O}_{\psi(x)}/\sh{N}_{\psi(x)}\to\sh{O}_x/\sh{N}_x$ é um homomorfismo local, de modo que $(\psi,\omega^\flat)$ é um morfismo de pré-esquemas $X_\red\to Y_\red$, que denotaremos por $f_\red$ e chamaremos de morfismo \emph{reduzido} associado a $f$.
É imediato que, para dois morfismos $f:X\to Y$ e $g:Y\to Z$, tem-se $(g\circ f)_\red=g_\red\circ f_\red$, de modo que definimos $X_\red$ como um \emph{funtor}, \emph{covariante} em $X$.

A definição precedente mostra que o diagrama
\[
  \xymatrix{
    X_\red\ar[r]^{f_\red}\ar[d] &
    Y_\red\ar[d]\\
    X\ar[r]_f &
    Y
  }
\]
é comutativo, sendo as setas verticais os morfismos de injeção;
em outras palavras, $X_\red\to X$ é um morfismo \emph{funtorial}.
Notemos, em particular, que, se $X$ é reduzido, todo morfismo $f:X\to Y$ se fatora como $X\xrightarrow{f_\red}Y_\red\to Y$;
em outras palavras, $f$ é dominado pelo morfismo de injeção $Y_\red\to Y$.
\end{env}

\begin{proposition}[5.1.6]
\label{I.5.1.6}
Seja $f:X\to Y$ um morfismo.
Se $f$ é sobrejetivo (resp. radicial, uma imersão, uma imersão fechada, uma imersão aberta, uma imersão local, um isomorfismo local), então $f_\red$ também o é.
Reciprocamente, se $f_\red$ é sobrejetivo (resp. radicial), então $f$ também o é.
\end{proposition}

\begin{proof}
A proposição é trivial se $f$ é sobrejetivo.
Se $f$ é radicial, ela resulta do fato de que, para todo $x\in X$, o corpo $\kres(x)$ é o mesmo para os pré-esquemas $X$ e $X_\red$ \sref{I.3.5.8}.
Finalmente, se $f=(\psi,\theta)$ é uma imersão, uma imersão fechada ou uma imersão local (resp. uma imersão aberta ou um isomorfismo local), a proposição resulta do fato de que, se $\theta_x^\sharp$ é sobrejetivo (resp. bijetivo), o mesmo vale para o homomorfismo obtido passando aos quocientes pelos nilradicais de $\sh{O}_{\psi(x)}$ e $\sh{O}_x$ (\sref{I.5.1.2} e \sref{I.4.2.2}) (cf. \sref{I.5.5.12}).
\end{proof}

\begin{proposition}[5.1.7]
\label{I.5.1.7}
Se $X$ e $Y$ são dois $S$-pré-esquemas, então os pré-esquemas $X_\red\times_{S_\red}Y_\red$ e $X_\red\times_S Y_\red$ são idênticos e se identificam canonicamente com um subpré-esquema de $X\times_S Y$ que tem o mesmo espaço subjacente que esse produto.
\end{proposition}

\begin{proof}
A identificação canônica de $X_\red\times_S Y_\red$ com um subpré-esquema de $X\times_S Y$ que tem o mesmo espaço subjacente resulta de \sref{I.4.3.1}.
Além disso, se $\vphi$ e $\psi$ são, respectivamente, os morfismos estruturais $X_\red\to S$ e $Y_\red\to S$, eles se fatoram por $S_\red$ \sref{I.5.1.5}; como $S_\red\to S$ é um monomorfismo, a primeira afirmação da proposição resulta de \sref{I.3.2.4}.
\end{proof}

\begin{corollary}[5.1.8]
\label{I.5.1.8}
Os pré-esquemas $(X\times_S Y)_\red$ e $(X_\red\times_{S_\red}Y_\red)_\red$ identificam-se canonicamente entre si.
\end{corollary}

\begin{proof}
Isso decorre de \sref{I.5.1.2} e \sref{I.5.1.7}.
\end{proof}

Observemos que, embora $X$ e $Y$ sejam pré-esquemas reduzidos, $X\times_S Y$ não é necessariamente reduzido, pois o produto tensorial de duas álgebras reduzidas pode ter elementos nilpotentes.

\begin{proposition}[5.1.9]
\label{I.5.1.9}
\oldpage[I]{130}
Seja $X$ um pré-esquema, e seja $\sh{I}$ um feixe quase-coerente de ideais de $\sh{O}_X$ tal que $\sh{I}^n=0$ para algum inteiro $n>0$.
Seja $X_0$ o subpré-esquema fechado $(X,\sh{O}_X/\sh{I})$ de $X$;
para que $X$ seja um esquema afim, é necessário e suficiente que $X_0$ seja um esquema afim.
\end{proposition}

A condição é evidentemente necessária; provemos que é suficiente.
Se pusermos $X_k=(X,\sh{O}_X/\sh{I}^{k+1})$, basta provar por indução em $k$ que os $X_k$ são afins; assim, ficamos reduzidos ao caso em que $\sh{I}^2=0$.
Ponhamos
\begin{align*}
  A&=\Gamma(X,\sh{O}_X)\\
  A_0&=\Gamma(X_0,\sh{O}_{X_0})=\Gamma(X,\sh{O}_X/\sh{I}).
\end{align*}
O homomorfismo canônico $\sh{O}_X\to\sh{O}_X/\sh{I}$ induz um homomorfismo de anéis $\vphi:A\to A_0$.
Veremos adiante que $\vphi$ é \emph{sobrejetivo}, de modo que
\[
\label{I.5.1.9.1}
  0\to\Gamma(X,\sh{I})\to\Gamma(X,\sh{O}_X)\to\Gamma(X,\sh{O}_X/\sh{I})\to 0\tag{5.1.9.1}
\]
é uma sequência \emph{exata}.
Suponhamos estabelecido esse ponto e mostremos que dele decorre a proposição.
Notemos que $\mathfrak{K}=\Gamma(X,\sh{I})$ é um ideal de quadrado nulo em $A$ e, portanto, um módulo sobre $A_0=A/\mathfrak{K}$.
Por hipótese, temos $X_0=\Spec(A_0)$ e, como os espaços subjacentes a $X_0$ e $X$ são idênticos, $\mathfrak{K}=\Gamma(X_0,\sh{I})$;
além disso, como $\sh{I}^2=0$, $\sh{I}$ é um $(\sh{O}_X/\sh{I})$-módulo quase-coerente, de modo que $\sh{I}\cong\widetilde{\mathfrak{K}}$ e $\mathfrak{K}_x=\sh{I}_x$ para todo $x\in X_0$ \sref{I.1.4.1}.

Nessas condições, seja $X'=\Spec(A)$, e consideremos o morfismo $f=(\psi,\theta):X\to X'$ de pré-esquemas correspondente à aplicação identidade $A\to\Gamma(X,\sh{O}_X)$ \sref{I.2.2.4}.
Para todo subconjunto aberto afim $V$ de $X$, o diagrama
\[
  \xymatrix{
    A\ar[r]\ar[d] &
    \Gamma(V,\sh{O}_X|V)\ar[d]\\
    A_0=A/\mathfrak{K}\ar[r] &
    \Gamma(V,\sh{O}_{X_0}|V)
  }
\]
é comutativo, donde se deduz que o diagrama
\[
  \xymatrix{
    X' &
    X\ar[l]_f\\
    X'_0\ar[u]^{j'} &
    X_0\ar[l]^{f_0}\ar[u]_j
  }
\]
também é comutativo (sendo $X'_0$ o subpré-esquema fechado de $X'$ definido pelo feixe quase-coerente de ideais $\widetilde{\mathfrak{K}}$, e sendo $j$ e $j'$ os morfismos de injeção canônicos).
Mas, como $X_0$ é afim, $f_0$ é um isomorfismo e, como as aplicações contínuas subjacentes a $j$ e $j'$ são as aplicações identidade, vemos imediatamente que $\psi:X\to X'$ é um homeomorfismo.
Além disso, a igualdade $\mathfrak{K}_x=\sh{I}_x$ mostra que a restrição de $\theta^\sharp:\psi^*(\sh{O}_{X'})\to\sh{O}_X$ é um \emph{isomorfismo} de $\psi^*(\widetilde{\mathfrak{K}})$ sobre $\sh{I}$;
por outro lado, passando aos quocientes, $\theta^\sharp$ fornece um \emph{isomorfismo} $\psi^*(\sh{O}_{X'}/\widetilde{\mathfrak{K}})\to\sh{O}_X/\sh{I}$, pois $f_0$ é um isomorfismo;
conclui-se então imediatamente, pelo lema dos cinco (M,~I,~1.1), que $\theta^\sharp$ é por sua vez um isomorfismo e, portanto, que $f$ é um \emph{isomorfismo} e, por conseguinte, que $X$ é afim.

Assim, tudo se reduz a demonstrar a exatidão de \sref{I.5.1.9.1}, o que decorrerá de demonstrar que $\HH^1(X,\sh{I})=0$.
\oldpage[I]{131}
Ora, $\HH^1(X,\sh{I})=\HH^1(X_0,\sh{I})$, e vimos que $\sh{I}$ é um $\sh{O}_{X_0}$-módulo quase-coerente.
Portanto, nossa afirmação decorrerá do
\begin{lemma}[5.1.9.2]
\label{I.5.1.9.2}
Se $Y$ é um esquema afim e $\sh{F}$ um $\sh{O}_Y$-módulo quase-coerente, então $\HH^1(Y,\sh{F})=0$.
\end{lemma}

\begin{proof}
Este lema será demonstrado no Capítulo~III, §1, como consequência do teorema mais geral segundo o qual $\HH^i(Y,\sh{F})=0$ para todo $i>0$.
Para dar uma demonstração independente, observemos que $\HH^1(Y,\sh{F})$ pode ser identificado com o módulo $\Ext_{\sh{O}_Y}^1(Y;\sh{O}_Y,\sh{F})$ das classes de extensões do $\sh{O}_Y$-módulo $\sh{O}_Y$ pelo $\sh{O}_Y$-módulo $\sh{F}$ (T,~4.2.3);
tudo se reduz, portanto, a demonstrar que uma tal extensão $\sh{G}$ é trivial.
Ora, para todo $y\in Y$, existe uma vizinhança $V$ de $y$ em $Y$ tal que $\sh{G}|V$ é isomorfo a $\sh{F}|V\oplus\sh{O}_Y|V$\footnote{\emph{[Trad.] A fonte francesa imprime aqui $\sh{F}|Y$; a restrição necessária sobre a vizinhança $V$ é $\sh{F}|V$.}} \sref[0]{0.5.4.9};
daí se conclui que $\sh{G}$ é um módulo \emph{quase-coerente} sobre $\sh{O}_Y$.
Se $A$ é o anel de $Y$, então se tem $\sh{F}=\widetilde{M}$ e $\sh{G}=\widetilde{N}$, onde $M$ e $N$ são $A$-módulos, e, por hipótese, $N$ é uma extensão do $A$-módulo $A$ pelo $A$-módulo $M$ \sref{I.1.3.11}.
Como essa extensão é necessariamente trivial, fica demonstrado o lema e, com ele, também \sref{I.5.1.9}.
\end{proof}

\begin{corollary}[5.1.10]
\label{I.5.1.10}
Seja $X$ um pré-esquema tal que $\sh{N}_X$ seja nilpotente.
Para que $X$ seja um esquema afim, é necessário e suficiente que $X_\red$ seja um esquema afim.
\end{corollary}

\subsection{Existência de um subpré-esquema com espaço subjacente dado}
\label{subsection:I.5.2}

\begin{proposition}[5.2.1]
\label{I.5.2.1}
Para todo subespaço localmente fechado $Y$ do espaço subjacente a um pré-esquema $X$, existe um único subpré-esquema reduzido de $X$ que tem $Y$ como espaço subjacente.
\end{proposition}

\begin{proof}
A unicidade resulta de \sref{I.5.1.2}, de modo que resta apenas demonstrar a existência do pré-esquema em questão.

Se $X$ é afim, de anel $A$, e $Y$ é fechado em $X$, a proposição é imediata:
$\mathfrak{j}(Y)$ é o maior ideal $\mathfrak{a}\subset A$ tal que $V(\mathfrak{a})=Y$, e é igual ao seu radical \sref{I.1.1.4}[i], de modo que $A/\mathfrak{j}(Y)$ é um anel reduzido.

No caso geral, para todo aberto afim $U\subset X$ tal que $U\cap Y$ seja fechado em $U$, consideremos o subpré-esquema fechado $Y_U$ de $U$ definido pelo feixe de ideais associado ao ideal $\mathfrak{j}(U\cap Y)$ de $A(U)$.
Esse subpré-esquema é reduzido.
Afirmamos que, se $V$ é um aberto afim de $X$ contido em $U$, então $Y_V$ é \emph{induzido} por $Y_U$ sobre $V\cap Y$.
Com efeito, esse pré-esquema induzido é um subpré-esquema fechado reduzido de $V$ que tem $V\cap Y$ como espaço subjacente;
a unicidade de $Y_V$ demonstra, portanto, a afirmação.
\end{proof}

\begin{proposition}[5.2.2]
\label{I.5.2.2}
Sejam $X$ um pré-esquema reduzido, $f:X \to Y$ um morfismo e $Z$ um subpré-esquema fechado de $Y$ tal que $f(X) \subset Z$.
Então $f$ se fatora como $X \xrightarrow{g} Z \xrightarrow{j} Y$, onde $j$ é o morfismo de injeção.
\end{proposition}

\begin{proof}
Das hipóteses resulta que o subpré-esquema fechado $f^{-1}(Z)$ de $X$ tem todo $X$ como espaço subjacente \sref{I.4.4.1};
como $X$ é reduzido, esse subpré-esquema fechado coincide com $X$ \sref{I.5.1.2}, e a proposição resulta então de \sref{I.4.4.1}.
\end{proof}

\begin{corollary}[5.2.3]
\label{I.5.2.3}
Seja $X$ um subpré-esquema reduzido de um pré-esquema $Y$; se $Z$ é o subpré-esquema fechado reduzido de $Y$ que tem $\overline{X}$ como espaço subjacente, então $X$ é um subpré-esquema induzido sobre um aberto de $Z$.
\end{corollary}

\begin{proof}
\oldpage[I]{132}
Existe, com efeito, um aberto $U$ de $Y$ tal que $X=U\cap\overline{X}$;
como, em virtude de \sref{I.5.2.2}, $X$ é um subpré-esquema reduzido de $Z$, o subpré-esquema $X$ é induzido por $Z$ sobre o subespaço aberto $X$ em virtude da unicidade \sref{I.5.2.1}.
\end{proof}

\begin{corollary}[5.2.4]
\label{I.5.2.4}
Sejam $f:X\to Y$ um morfismo e $X'$ (resp. $Y'$) um subpré-esquema fechado de $X$ (resp. $Y$) definido por um feixe quase-coerente de ideais $\sh{J}$ (resp. $\sh{K}$) de $\sh{O}_X$ (resp. $\sh{O}_Y$).
Suponhamos que $X'$ seja reduzido e que $f(X')\subset Y'$.
Então $f^*(\sh{K})\sh{O}_X\subset\sh{J}$.
\end{corollary}

\begin{proof}
Como, por \sref{I.5.2.2}, a restrição de $f$ a $X'$ se fatora como $X'\to Y'\to Y$, basta aplicar \sref{I.4.4.6}.
\end{proof}

\subsection{Diagonal; gráfico de um morfismo}
\label{subsection:I.5.3}
\label{I.5.3}

\begin{env}[5.3.1]
\label{I.5.3.1}
Seja $X$ um $S$-pré-esquema;
chama-se \emph{morfismo diagonal} de $X$ em $X\times_S X$, e denota-se por $\Delta_{X|S}$, ou $\Delta_X$, ou mesmo $\Delta$ se não houver possibilidade de confusão, o $S$-morfismo $(1_X,1_X)_S$, isto é, o único $S$-morfismo $\Delta_X$ tal que
\[
\label{I.5.3.1.1}
  p_1\circ\Delta_X=p_2\circ\Delta_X=1_X,
  \tag{5.3.1.1}
\]
onde $p_1$ e $p_2$ são as projeções de $X\times_S X$ (Definição~\sref{I.3.2.1}).
Se $f:T\to X$ e $g:T\to Y$ são dois $S$-morfismos, verifica-se imediatamente que
\[
\label{I.5.3.1.2}
  (f,g)_S=(f\times_S g)\circ\Delta_{T|S}.
  \tag{5.3.1.2}
\]

O leitor observará que a definição precedente e os resultados enunciados de \sref{I.5.3.1} a \sref{I.5.3.8} são válidos em toda categoria, \emph{desde que os produtos que neles figuram existam nessa categoria}.
\end{env}

\begin{proposition}[5.3.2]
\label{I.5.3.2}
Sejam $X$ e $Y$ dois $S$-pré-esquemas;
se identificarmos canonicamente o produto $(X\times Y)\times(X\times Y)$ com $(X\times X)\times(Y\times Y)$, então o morfismo $\Delta_{X\times Y}$ se identifica com $\Delta_X\times\Delta_Y$.
\end{proposition}

\begin{proof}
Com efeito, se $p_1:X\times X\to X$ e $q_1:Y\times Y\to Y$ são as projeções sobre o primeiro fator, então a projeção sobre o primeiro fator $(X\times Y)\times(X\times Y)\to X\times Y$ se identifica com $p_1\times q_1$, e tem-se
\[
  (p_1\times q_1)\circ(\Delta_X\times\Delta_Y)=(p_1\circ\Delta_X)\times(q_1\circ\Delta_Y)=1_{X\times Y}
\]
Raciocina-se do mesmo modo para as projeções sobre o segundo fator.
\end{proof}

\phantomsection
\label{I.5.3.3}

\begin{corollary}[5.3.4]
\label{I.5.3.4}
Para toda extensão $S'\to S$ do pré-esquema de base, $\Delta_{X_{(S')}}$ se identifica canonicamente com $(\Delta_X)_{(S')}$.
\end{corollary}

\begin{proof}
Basta observar que $(X\times_S X)_{(S')}$ se identifica canonicamente com $X_{(S')}\times_{S'}X_{(S')}$ \sref{I.3.3.10}.
\end{proof}

\begin{proposition}[5.3.5]
\label{I.5.3.5}
Sejam $X$ e $Y$ dois $S$-pré-esquemas, e seja $\vphi:S\to T$ um morfismo de pré-esquemas que permite considerar todo $S$-pré-esquema como um $T$-pré-esquema.
Sejam $f:X\to S$ e $g:Y\to S$\footnote{\emph{[Trad.] A fonte francesa imprime somente $Y\to S$ aqui, omitindo o nome $g$; a identidade seguinte $\pi=f\circ p=g\circ q$ identifica esse segundo morfismo estrutural como $g$.}} os morfismos estruturais, $p$ e $q$ as projeções de $X\times_S Y$, e $\pi=f\circ p=g\circ q$ o morfismo estrutural $X\times_S Y\to S$.
Então, o diagrama
\[
\label{I.5.3.5.1}
  \xymatrix{
    X\times_S Y\ar[r]^{(p,q)_T}\ar[d]_\pi &
    X\times_T Y\ar[d]^{f\times_T g}\\
    S\ar[r]_{\Delta_{S|T}} &
    S\times_T S
  }
  \tag{5.3.5.1}
\]
\oldpage[I]{133}
é comutativo e identifica $X\times_S Y$ com o produto dos $(S\times_T S)$-pré-esquemas $S$ e $X\times_T Y$, identificando-se as projeções com $\pi$ e $(p,q)_T$.
\end{proposition}

\begin{proof}
Em virtude de \sref{I.3.4.3}, reduzimo-nos a demonstrar a proposição correspondente \emph{na categoria dos conjuntos}, substituindo $X$, $Y$ e $S$ por $X(Z)_T$, $Y(Z)_T$ e $S(Z)_T$, respectivamente, sendo $Z$ um $T$-pré-esquema arbitrário.
Mas, na categoria dos conjuntos, a verificação é imediata e fica a cargo do leitor.
\end{proof}

\begin{corollary}[5.3.6]
\label{I.5.3.6}
O morfismo $(p,q)_T$ identifica-se (pondo $P=S\times_T S$) com $1_{X\times_T Y}\times_P\Delta_S$.
\end{corollary}

\begin{proof}
Isso resulta de \sref{I.5.3.5} e \sref{I.3.3.4}.
\end{proof}

\begin{corollary}[5.3.7]
\label{I.5.3.7}
Se $f:X\to Y$ é um $S$-morfismo, então o diagrama
\[
  \xymatrix{
    X\ar[r]^{(1_X,f)_S}\ar[d]_f &
    X\times_S Y\ar[d]^{f\times_S 1_Y}\\
    Y\ar[r]_{\Delta_Y} &
    Y\times_S Y
  }
\]
é comutativo e identifica $X$ com o produto dos $(Y\times_S Y)$-pré-esquemas $Y$ e $X\times_S Y$.
\end{corollary}

\begin{proof}
Basta aplicar \sref{I.5.3.5}, substituindo $S$ por $Y$ e $T$ por $S$, e observar que $X\times_Y Y=X$ \sref{I.3.3.3}.
\end{proof}

\begin{proposition}[5.3.8]
\label{I.5.3.8}
Para que $f:X\to Y$ seja um monomorfismo de pré-esquemas, é necessário e suficiente que $\Delta_{X|Y}$ seja um isomorfismo de $X$ sobre $X\times_Y X$.
\end{proposition}

\begin{proof}
Com efeito, $f$ é um monomorfismo se e somente se, para todo $Y$-pré-esquema $Z$, a aplicação correspondente $f':X(Z)_Y\to Y(Z)_Y$ é injetiva.
Como $Y(Z)_Y$ consiste em um único elemento, isso equivale a dizer que $X(Z)_Y$ tem no máximo um elemento\footnote{\emph{[Trad.] A fonte francesa afirma que $X(Z)_Y$ se reduz igualmente a um único elemento; a injetividade em um conjunto unitário $Y(Z)_Y$ implica apenas que o domínio tem no máximo um elemento, admitindo o caso vazio.}}, ou, equivalentemente, que a aplicação diagonal de $X(Z)_Y$ em $X(Z)_Y\times X(Z)_Y$ é uma bijeção.
Este último produto é precisamente o conjunto $(X\times_Y X)(Z)_Y$ \hyperref[I.3.4.3]{(3.4.3.1)}, de modo que isso equivale a $\Delta_{X|Y}$ ser um isomorfismo.
\end{proof}

\begin{proposition}[5.3.9]
\label{I.5.3.9}
O morfismo diagonal $\Delta_X$ é uma imersão de $X$ em $X\times_S X$.
\end{proposition}

\begin{proof}
A demonstração impressa na fonte francesa é insuficiente, pois não demonstra que $\Delta_X(X)$ seja localmente fechado em $X\times_S X$.\footnote{\emph{[Trad.] EGA~III.2, Erratas e adições (lista~2), item $\mathrm{Err}_{\mathrm{III}},10$, identifica explicitamente esta lacuna e fornece a demonstração afim-local traduzida aqui.}}
Para uma demonstração correta, basta usar \sref{I.4.2.4}[a]:
para todo $x\in X$, escolha-se uma vizinhança afim $W$ da imagem de $x$ em $S$ e, em seguida, uma vizinhança afim $U$ de $x$ em $X$ contida na imagem inversa de $W$ pelo morfismo estrutural;
$U\times_S U$ identifica-se com $U\times_W U$ e é uma vizinhança afim de $\Delta_X(x)$.\footnote{\emph{[Trad.] A formulação publicada em EGA~III.2, Erratas e adições (lista~2), item $\mathrm{Err}_{\mathrm{III}},10$, afirma isso para toda vizinhança afim $U$ sem exigir que seu morfismo estrutural se fatore por um aberto afim de $S$; a hipótese explicitada aqui é necessária para este passo.}}
Levando em conta \sref{I.5.3.16} (cuja demonstração utiliza apenas a Definição
\sref{I.5.3.1.1}), reduzimo-nos a demonstrar \sref{I.5.3.9} quando
$S=\Spec(B)$ e $X=\Spec(A)$ são esquemas afins.
Resulta então claramente de \sref{I.5.3.1.1} que $\Delta_X$ corresponde ao
homomorfismo canônico $A\otimes_B A\to A$ que leva $x\otimes y$ a $xy$.
Como esse homomorfismo é sobrejetivo, $\Delta_X$ é, nesse caso, uma imersão
fechada \sref{I.4.2.3}, o que conclui a demonstração.
\end{proof}

Diz-se que o subpré-esquema de $X\times_S X$ associado à imersão $\Delta_X$ \sref{I.4.2.1} é \emph{a diagonal} de $X\times_S X$.

\begin{corollary}[5.3.10]
\label{I.5.3.10}
Sob as hipóteses de \sref{I.5.3.5}, $(p,q)_T$ é uma imersão.
\end{corollary}

\begin{proof}
Isso resulta de \sref{I.5.3.6} e \sref{I.4.3.1}.
\end{proof}

Diz-se (sob as hipóteses de \sref{I.5.3.5}) que $(p,q)_T$ é a \emph{imersão canônica} de $X\times_S Y$ em $X\times_T Y$.

\begin{corollary}[5.3.11]
\label{I.5.3.11}
Sejam $X$ e $Y$ dois $S$-pré-esquemas, e seja $f:X\to Y$ um $S$-morfismo;
então o morfismo gráfico $\Gamma_f=(1_X,f)_S$ de $f$ \sref{I.3.3.14} é uma imersão de $X$ em $X\times_S Y$.
\end{corollary}

\begin{proof}
Este é o caso particular do Corolário \sref{I.5.3.10}, no qual se substitui $S$ por $Y$ e $T$ por $S$ (cf. \sref{I.5.3.7}).
\end{proof}

\oldpage[I]{134}
O subpré-esquema de $X\times_S Y$ associado à imersão $\Gamma_f$ \sref{I.4.2.1} é chamado de \emph{gráfico} do morfismo $f$;
os subpré-esquemas de $X\times_S Y$ que são gráficos de morfismos $X\to Y$ caracterizam-se pelo fato de que a restrição a tal subpré-esquema $G$ da projeção $p_1:X\times_S Y\to X$ é um \emph{isomorfismo} $g$ de $G$ sobre $X$:
$G$ é então o gráfico do morfismo $p_2\circ g^{-1}$, onde $p_2$ é a projeção $X\times_S Y\to Y$.

Tomando, em particular, $X=S$, os $S$-morfismos $S\to Y$ (que não são senão as $S$-\emph{seções} de $Y$ \sref{I.2.5.5}) coincidem com seus morfismos gráficos;
os subpré-esquemas de $Y$ que são gráficos de $S$-seções (em outras palavras, os que são isomorfos a $S$ pela restrição do morfismo estrutural $Y\to S$) chamam-se também as \emph{imagens} dessas seções ou, por abuso de linguagem, as $S$-\emph{seções} de $Y$.

\begin{corollary}[5.3.12]
\label{I.5.3.12}
Sob as hipóteses e com as notações de \sref{I.5.3.11}, para todo morfismo $g:S'\to S$, seja $f'$ a imagem inversa de $f$ por $g$ \sref{I.3.3.7};
então $\Gamma_{f'}$ é a imagem inversa de $\Gamma_f$ por $g$.
\end{corollary}

\begin{proof}
Este é um caso particular da fórmula \sref{I.3.3.10.1}.
\end{proof}

\begin{corollary}[5.3.13]
\label{I.5.3.13}
Sejam $f:X\to Y$ e $g:Y\to Z$ dois morfismos;
se $g\circ f$ é uma imersão (resp. uma imersão local), então $f$ também o é.
\end{corollary}

\begin{proof}
Com efeito, $f$ se fatora como $X\xrightarrow{\Gamma_f}X\times_Z Y\xrightarrow{p_2}Y$.
Além disso, $p_2$ identifica-se com $(g\circ f)\times_Z 1_Y$ \sref{I.3.3.4};
se $g\circ f$ é uma imersão (resp. uma imersão local), então $p_2$ também o é (\sref{I.4.3.1} e \sref{I.4.5.5}), e, como $\Gamma_f$ é uma imersão \sref{I.5.3.11}, conclui-se por \sref{I.4.2.5}\footnote{\emph{[Trad.] A fonte francesa cita aqui \sref{I.4.2.4}; EGA~III.2, Erratas e adições (lista~2), prescreve explicitamente sua substituição por \sref{I.4.2.5}.}} (resp. \sref{I.4.5.5}).
\end{proof}

\begin{corollary}[5.3.14]
\label{I.5.3.14}
Sejam $j:X\to Y$ e $g:X\to Z$ dois $S$-morfismos.
Se $j$ é uma imersão (resp. uma imersão local), então $(j,g)_S$ também o é.
\end{corollary}

\begin{proof}
Com efeito, se $p:Y\times_S Z\to Y$ é a projeção sobre o primeiro fator, então $j=p\circ(j,g)_S$, e basta aplicar \sref{I.5.3.13}.
\end{proof}

\begin{proposition}[5.3.15]
\label{I.5.3.15}
Se $f:X\to Y$ é um $S$-morfismo, então o diagrama
\[
\label{I.5.3.15.1}
  \xymatrix{
    X\ar[r]^{\Delta_X}\ar[d]_f &
    X\times_S X\ar[d]^{f\times_S f}\\
    Y\ar[r]_{\Delta_Y} &
    Y\times_S Y
  }\tag{5.3.15.1}
\]
é comutativo \emph{(em outras palavras, $\Delta_X$ é um morfismo funtorial na categoria dos pré-esquemas)}.
\end{proposition}

\begin{proof}
A verificação é imediata e fica a cargo do leitor.
\end{proof}

\begin{corollary}[5.3.16]
\label{I.5.3.16}
Se $X$ é um subpré-esquema de $Y$, então a diagonal $\Delta_X(X)$ identifica-se com um subpré-esquema de $\Delta_Y(Y)$, cujo espaço subjacente se identifica com
\[
  \Delta_Y(Y)\cap p_1^{-1}(X)=\Delta_Y(Y)\cap p_2^{-1}(X)
\]
($p_1$ e $p_2$ são as projeções de $Y\times_S Y$).
\end{corollary}

\begin{proof}
Aplicando \sref{I.5.3.15} ao morfismo de injeção $f:X\to Y$, vê-se que $f\times_S f$ é uma imersão que identifica o espaço subjacente a $X\times_S X$ com o subespaço $p_1^{-1}(X)\cap p_2^{-1}(X)$ de $Y\times_S Y$ \sref{I.4.3.1};
além disso, se $z\in\Delta_Y(Y)\cap p_1^{-1}(X)$, então $z=\Delta_Y(y)$
\oldpage[I]{135}
e $y=p_1(z)\in X$, portanto $y=f(y)$, e $z=\Delta_Y(f(y))$ pertence a $\Delta_X(X)$ em virtude da comutatividade de \sref{I.5.3.15.1}.
\end{proof}

\begin{corollary}[5.3.17]
\label{I.5.3.17}
Sejam $f_1:Y\to X$ e $f_2:Y\to X$ dois $S$-morfismos, e seja $y$ um ponto de $Y$ tal que $f_1(y)=f_2(y)=x$ e que os homomorfismos $\kres(x)\to\kres(y)$ correspondentes a $f_1$ e $f_2$ sejam idênticos.
Então, se $f=(f_1,f_2)_S$, o ponto $f(y)$ pertence à diagonal $\Delta_{X|S}(X)$.
\end{corollary}

\begin{proof}
Os dois homomorfismos $\kres(x)\to\kres(y)$ correspondentes a $f_i$ ($i=1,2$) definem dois $S$-morfismos $g_i:\Spec(\kres(y))\to\Spec(\kres(x))$ tais que os diagramas
\[
  \xymatrix{
    \Spec(\kres(y))\ar[r]^{g_i}\ar[d] &
    \Spec(\kres(x))\ar[d]\\
    Y\ar[r]_{f_i} &
    X
  }
\]
sejam comutativos.
O diagrama
\[
  \xymatrix{
    \Spec(\kres(y))\ar[rr]^{(g_1,g_2)_S}\ar[d] & &
    \Spec(\kres(x))\times_S\Spec(\kres(x))\ar[d]\\
    Y\ar[rr]_{(f_1,f_2)_S} & &
    X\times_S X
  }
\]
é, portanto, também comutativo.
Ora, da igualdade $g_1=g_2$ resulta que a imagem por $(g_1,g_2)_S$ do único ponto de $\Spec(\kres(y))$ pertence à diagonal de $\Spec(\kres(x))\times_S\Spec(\kres(x))$;
a conclusão decorre então de \sref{I.5.3.15}.
\end{proof}

\subsection{Morfismos separados e pré-esquemas separados}
\label{subsection:I.5.4}
\label{I.5.4}

\begin{definition}[5.4.1]
\label{I.5.4.1}
Diz-se que um morfismo de pré-esquemas $f:X\to Y$ é \emph{separado} se o morfismo diagonal $X\to X\times_Y X$ é uma imersão \emph{fechada}; diz-se então também que $X$ é um \emph{pré-esquema separado sobre $Y$}, ou um \emph{$Y$-esquema}.
Diz-se que um pré-esquema $X$ é separado se for separado sobre $\Spec(\bb{Z})$;
diz-se então também que $X$ é um \emph{esquema}\footnote{\emph{[Trad.] Repetimos aqui a advertência formulada no início desta tradução: as primeiras versões de EGA empregam \emph{pré-esquema} para o que hoje se costuma chamar esquema, e \emph{esquema} para o que hoje se costuma chamar esquema separado. O próprio Grothendieck disse mais tarde que a terminologia moderna era preferível, mas decidimos manter esta tradução «historicamente fiel», usando a nomenclatura antiga.}} \emph{(cf. \sref{I.5.5.7})}.
\end{definition}

Em virtude de \sref{I.5.3.9}, para que $X$ seja separado sobre $Y$, é necessário e suficiente que $\Delta_X(X)$ seja um \emph{subespaço fechado} do espaço subjacente a $X\times_Y X$.

\begin{proposition}[5.4.2]
\label{I.5.4.2}
Seja $S\to T$ um morfismo separado.
Se $X$ e $Y$ são dois $S$-pré-esquemas, então a imersão canônica $X\times_S Y\to X\times_T Y$ \sref{I.5.3.10} é fechada.
\end{proposition}

\begin{proof}
Com efeito, remetendo-nos ao diagrama de \sref{I.5.3.5.1}, vê-se que $(p,q)_T$ pode ser considerado como obtido de $\Delta_{S|T}$ pela extensão $f\times_T g:X\times_T Y\to S\times_T S$ do pré-esquema de base;
a proposição decorre então de \sref{I.4.3.2}.
\end{proof}

\begin{corollary}[5.4.3]
\label{I.5.4.3}
Seja $Y$ um $S$-esquema, e seja $f:X\to Y$ um $S$-morfismo.
Então o morfismo gráfico $\Gamma_f:X\to X\times_S Y$ \sref{I.5.3.11} é uma imersão fechada.
\end{corollary}

\begin{proof}
Este é o caso particular de \sref{I.5.4.2} em que se substitui $S$ por $Y$ e $T$ por $S$.
\end{proof}

\begin{corollary}[5.4.4]
\label{I.5.4.4}
Sejam $f:X\to Y$ e $g:Y\to Z$ dois morfismos, sendo $g$ separado.
Se $g\circ f$ é uma imersão fechada, então $f$ também o é.
\end{corollary}

\begin{proof}
A demonstração a partir de \sref{I.5.4.3} é a mesma que a de \sref{I.5.3.13} a partir de \sref{I.5.3.11}.
\end{proof}

\oldpage[I]{136}
\begin{corollary}[5.4.5]
\label{I.5.4.5}
Seja $Z$ um $S$-esquema, e sejam $j:X\to Y$ e $g:X\to Z$ dois $S$-morfismos.
Se $j$ é uma imersão fechada, então $(j,g)_S:X\to Y\times_S Z$ também o é.
\end{corollary}

\begin{proof}
A demonstração a partir de \sref{I.5.4.4} é a mesma que a de \sref{I.5.3.14} a partir de \sref{I.5.3.13}.
\end{proof}

\begin{corollary}[5.4.6]
\label{I.5.4.6}
Se $X$ é um $S$-esquema, então toda $S$-seção de $X$ \sref{I.2.5.5} é uma imersão fechada.
\end{corollary}

\begin{proof}
Se $\vphi:X\to S$ é o morfismo estrutural, e $\psi:S\to X$ uma $S$-seção de $X$, basta aplicar \sref{I.5.4.5} a $\vphi\circ\psi=1_S$.
\end{proof}

\begin{corollary}[5.4.7]
\label{I.5.4.7}
Sejam $S$ um pré-esquema integral, $s$ seu ponto genérico e $X$ um $S$-esquema.
Se duas $S$-seções $f$ e $g$ de $X$ são tais que $f(s)=g(s)$, então $f=g$.
\end{corollary}

\begin{proof}
Com efeito, se $x=f(s)=g(s)$, então os homomorfismos $\kres(x)\to\kres(s)$ correspondentes a $f$ e $g$ são necessariamente idênticos.
Se $h=(f,g)_S$, deduz-se, pois, de \sref{I.5.3.17} que $h(s)$ pertence à diagonal $Z=\Delta_X(X)$;
mas, como $S=\overline{\{s\}}$ e como $Z$ é fechado por hipótese, tem-se $h(S)\subset Z$.
Resulta então de \sref{I.5.2.2} que $h$ se fatora como $S\to Z\to X\times_S X$, e conclui-se assim que $f=g$, por definição da diagonal.
\end{proof}

\begin{remark}[5.4.8]
\label{I.5.4.8}
Se se supõe, reciprocamente, que a conclusão de \sref{I.5.4.3} é válida quando $f=1_Y$, conclui-se que $Y$ é separado sobre $S$;
do mesmo modo, se se supõe que a conclusão de \sref{I.5.4.4}\footnote{\emph{[Trad.] A fonte francesa remete aqui a (5.4.5), que diz respeito a um morfismo definido por um par; o argumento com $p_1\circ\Delta_Y=1_Y$ utiliza o critério de composição do Corolário~\hyperref[I.5.4.4]{5.4.4}.}} se aplica aos dois morfismos $Y\xrightarrow{\Delta_Y}Y\times_Z Y\xrightarrow{p_1}Y$, conclui-se que $\Delta_Y$ é uma imersão fechada e, portanto, que $Y$ é separado sobre $Z$;
por fim, se se supõe que a conclusão de \sref{I.5.4.6} é válida para a $Y$-seção $\Delta_Y$ do $Y$-pré-esquema $Y\times_S Y\to Y$, isso implica que $Y$ é separado sobre $S$.
\end{remark}

\subsection{Critérios de separação}
\label{subsection:I.5.5}

\begin{proposition}[5.5.1]
\label{I.5.5.1}
\begin{enumerate}
  \item[\rm(i)] Todo monomorfismo de pré-esquemas (e, em particular, toda
    imersão) é um morfismo separado.
  \item[\rm(ii)] O composto de dois morfismos separados é separado.
  \item[\rm(iii)] Se $f:X\to X'$ e $g:Y\to Y'$ são dois $S$-morfismos
    separados, então $f\times_S g$ é separado.
  \item[\rm(iv)] Se $f:X\to Y$ é um $S$-morfismo separado, então o
    $S'$-morfismo $f_{(S')}$ é separado para toda extensão $S'\to S$ do
    pré-esquema de base.
  \item[\rm(v)] Se o composto $g\circ f$ de dois morfismos é separado, então
    $f$ é separado.
  \item[\rm(vi)] Para que um morfismo $f$ seja separado, é necessário e
    suficiente que $f_\red$ (\sref{I.5.1.5}) o seja.
\end{enumerate}
\end{proposition}

\begin{proof}
A afirmação~(i) resulta imediatamente de \sref{I.5.3.8}.
Se $f:X\to Y$ e $g:Y\to Z$ são dois morfismos, o diagrama
\[
\label{I.5.5.1.1}
  \xymatrix{
    X\ar[rr]^{\Delta_{X|Z}}\ar[dr]_{\Delta_{X|Y}} & &
    X\times_Z X\\
    & X\times_Y X\ar[ur]_j
  }
  \tag{5.5.1.1}
\]
é comutativo, onde $j$ designa a imersão canônica \sref{I.5.3.10}, como se
verifica imediatamente.
Se $f$ e $g$ são separados, então $\Delta_{X|Y}$ é uma imersão fechada por
definição, e $j$ é uma imersão fechada em virtude de \sref{I.5.4.2}; logo,
$\Delta_{X|Z}$ é uma imersão fechada por \sref{I.4.2.5}\footnote{\emph{[Trad.]
A fonte francesa remete aqui a (4.2.4), que é o critério local; a transitividade
das imersões fechadas utilizada neste passo é a Proposição~\hyperref[I.4.2.5]{4.2.5}.}}, o que
\oldpage[I]{137}
demonstra~(ii).
Em virtude de~(i) e~(ii), (iii) e~(iv) são equivalentes \sref{I.3.5.1}, e
basta demonstrar~(iv).
Ora, $X_{(S')}\times_{Y_{(S')}}X_{(S')}$ identifica-se canonicamente com
$(X\times_Y X)\times_Y Y_{(S')}$ em virtude de \sref{I.3.3.11} e
\hyperref[I.3.3.9.1]{(3.3.9.1)}, e verifica-se imediatamente que o morfismo
diagonal $\Delta_{X_{(S')}}$ se identifica então com
$\Delta_X\times_Y 1_{Y_{(S')}}$; a proposição resulta, portanto, de
\sref{I.4.3.1}.

Para demonstrar~(v), consideremos, como em \sref{I.5.3.13}, a fatoração
$X\xrightarrow{\Gamma_f}X\times_Z Y\xrightarrow{p_2}Y$ de $f$, observando
que $p_2=(g\circ f)\times_Z 1_Y$; a hipótese de que $g\circ f$ é separado
implica que $p_2$ é separado por (iii) e (i), e, como $\Gamma_f$ é uma
imersão, $\Gamma_f$ é separado por (i); logo, $f$ é separado por (ii).
Por fim, para demonstrar~(vi), recordemos que os pré-esquemas
$X_\red\times_{Y_\red}X_\red$ e $X_\red\times_Y X_\red$ se identificam
canonicamente \sref{I.5.1.7}; se designarmos por $j$ a injeção
$X_\red\to X$, o diagrama
\[
  \xymatrix{
    X_\red\ar[r]^{\Delta_{X_\red}}\ar[d]_j &
    X_\red\times_Y X_\red\ar[d]^{j\times_Y j}\\
    X\ar[r]_{\Delta_X} & X\times_Y X
  }
\]
é comutativo \sref{I.5.3.15}, e a proposição resulta do fato de que as setas
verticais são homeomorfismos dos espaços subjacentes \sref{I.4.3.1}.
\end{proof}

\begin{corollary}[5.5.2]
\label{I.5.5.2}
Se $f:X\to Y$ é separado, então a restrição de $f$ a todo subpré-esquema de
$X$ é separada.
\end{corollary}

\begin{proof}
Isso resulta de \sref{I.5.5.1}[i e ii].
\end{proof}

\begin{corollary}[5.5.3]
\label{I.5.5.3}
Se $X$ e $Y$ são dois $S$-pré-esquemas tais que $Y$ seja separado sobre $S$,
então $X\times_S Y$ é separado sobre $X$.
\end{corollary}

\begin{proof}
Este é um caso particular de \sref{I.5.5.1}[iv].
\end{proof}

\begin{proposition}[5.5.4]
\label{I.5.5.4}
Seja $X$ um pré-esquema, e suponhamos que seu espaço subjacente seja a reunião
de uma família finita de partes fechadas $X_k$ ($1\leq k\leq n$); para cada
$k$, consideremos o subpré-esquema reduzido de $X$ que tem $X_k$ como espaço
subjacente \sref{I.5.2.1}, e denotemo-lo ainda por $X_k$.
Seja $f:X\to Y$ um morfismo, e, para cada $k$, seja $Y_k$ uma parte fechada de
$Y$ tal que $f(X_k)\subset Y_k$; denotemos ainda por $Y_k$ o subpré-esquema
reduzido de $Y$ que tem $Y_k$ como espaço subjacente, de modo que a restrição
$X_k\to Y$ de $f$ a $X_k$ se fatore como
$X_k\xrightarrow{f_k}Y_k\to Y$ \sref{I.5.2.2}.
Para que $f$ seja separado, é necessário e suficiente que todos os $f_k$ sejam
separados.
\end{proposition}

\begin{proof}
A necessidade resulta de \sref{I.5.5.1}[i, ii e v].
Reciprocamente, se a condição do enunciado é satisfeita, então cada uma das
restrições $X_k\to Y$ de $f$ é separada
\hyperref[I.5.5.1]{(5.5.1, (i) e (ii))}; se $p_1$ e $p_2$ são as projeções de
$X\times_Y X$, então o subespaço $\Delta_{X_k}(X_k)$ identifica-se com o
subespaço $\Delta_X(X)\cap p_1^{-1}(X_k)$ do espaço subjacente a
$X\times_Y X$ \sref{I.5.3.16}; como esses subespaços são fechados em
$X\times_Y X$, o mesmo vale para sua reunião $\Delta_X(X)$.
\end{proof}

Suponhamos, em particular, que os $X_k$ sejam as \emph{componentes
irredutíveis} de $X$; podemos então supor que os $Y_k$ sejam componentes
irredutíveis de $Y$ \sref[0]{0.2.1.5}; a Proposição~\sref{I.5.5.4} reduz,
portanto, nesse caso, a noção de separação ao caso dos pré-esquemas
\emph{integrais} \sref{I.2.1.8}\footnote{\emph{[Trad.] A fonte francesa remete
aqui a (2.1.7), que trata da restrição a subconjuntos abertos; a definição
pertinente de pré-esquema integral é (2.1.8).}}.

\oldpage[I]{138}

\begin{proposition}[5.5.5]
\label{I.5.5.5}
Seja $(Y_\lambda)$ uma cobertura aberta de um pré-esquema $Y$; para que um
morfismo $f:X\to Y$ seja separado, é necessário e suficiente que cada uma de
suas restrições $f^{-1}(Y_\lambda)\to Y_\lambda$ seja separada.
\end{proposition}

\begin{proof}
Se pusermos $X_\lambda=f^{-1}(Y_\lambda)$, tudo se reduz, levando em conta
\sref{I.4.2.4}[b] e a identificação dos produtos
$X_\lambda\times_Y X_\lambda$ e
$X_\lambda\times_{Y_\lambda}X_\lambda$ \sref{I.3.2.5}, a provar que os
$X_\lambda\times_Y X_\lambda$ formam uma cobertura de $X\times_Y X$.
Ora, se pusermos $Y_{\lambda\mu}=Y_\lambda\cap Y_\mu$ e
$X_{\lambda\mu}=X_\lambda\cap X_\mu=f^{-1}(Y_{\lambda\mu})$, então
$X_\lambda\times_Y X_\mu$ se identifica com o produto
$X_{\lambda\mu}\times_{Y_{\lambda\mu}}X_{\lambda\mu}$
\sref{I.3.2.6.4}, portanto também com
$X_{\lambda\mu}\times_Y X_{\lambda\mu}$ \sref{I.3.2.5}, e finalmente com um
aberto de $X_\lambda\times_Y X_\lambda$, o que demonstra nossa afirmação
\sref{I.3.2.7}.
\end{proof}

A Proposição~\sref{I.5.5.5}\footnote{\emph{[Trad.] A fonte francesa cita
aqui (5.5.4); a redução mediante uma cobertura aberta de $Y$ utiliza a
Proposição~\hyperref[I.5.5.5]{5.5.5}.}} permite, tomando uma cobertura de
$Y$ por abertos afins, reduzir o estudo dos morfismos separados ao dos
morfismos separados com valores em esquemas afins.

\begin{proposition}[5.5.6]
\label{I.5.5.6}
Sejam $Y$ um esquema \emph{afim}, $X$ um pré-esquema e $(U_\alpha)$ uma
cobertura de $X$ por abertos afins.
Para que um morfismo $f:X\to Y$ seja separado, é necessário e suficiente que
$U_\alpha\cap U_\beta$ seja, para todo par de índices $(\alpha,\beta)$, um
aberto afim, \emph{e} que o anel
$\Gamma(U_\alpha\cap U_\beta,\sh{O}_X)$ seja gerado pela união das imagens
canônicas dos anéis $\Gamma(U_\alpha,\sh{O}_X)$ e
$\Gamma(U_\beta,\sh{O}_X)$.
\end{proposition}

\begin{proof}
Os $U_\alpha\times_Y U_\beta$ formam uma cobertura aberta de $X\times_Y X$
\sref{I.3.2.7}; designando as projeções de $X\times_Y X$ por $p$ e $q$, temos
\begin{align*}
  \Delta_X^{-1}(U_\alpha\times_Y U_\beta)&=\Delta_X^{-1}(p^{-1}(U_\alpha)\cap q^{-1}(U_\beta))\\
                                         &=\Delta_X^{-1}(p^{-1}(U_\alpha))\cap\Delta_X^{-1}(q^{-1}(U_\beta))=U_\alpha\cap U_\beta;
\end{align*}
de modo que tudo se reduz a provar que a restrição de $\Delta_X$ a
$U_\alpha\cap U_\beta$ é uma imersão fechada em
$U_\alpha\times_Y U_\beta$.
Ora, essa restrição não é outra senão $(j_\alpha,j_\beta)_Y$, onde $j_\alpha$
(resp. $j_\beta$) denota o morfismo de injeção de
$U_\alpha\cap U_\beta$ em $U_\alpha$ (resp. $U_\beta$), como resulta das
definições.
Como $U_\alpha\times_Y U_\beta$ é um esquema afim cujo anel é canonicamente
isomorfo a
$\Gamma(U_\alpha,\sh{O}_X)\otimes_{\Gamma(Y,\sh{O}_Y)}\Gamma(U_\beta,\sh{O}_X)$
\sref{I.3.2.2}, vemos que $U_\alpha\cap U_\beta$ deve ser um esquema afim e
que a aplicação $h_\alpha\otimes h_\beta\mapsto h_\alpha h_\beta$ do anel
$A(U_\alpha\times_Y U_\beta)$ em
$\Gamma(U_\alpha\cap U_\beta,\sh{O}_X)$ deve ser sobrejetiva
\sref{I.4.2.3}, o que conclui a demonstração.
\end{proof}

\begin{corollary}[5.5.7]
\label{I.5.5.7}
Um esquema afim é separado \emph{(e, portanto, é um \emph{esquema}, o que
justifica a terminologia de \sref{I.5.4.1})}.
\end{corollary}

\begin{corollary}[5.5.8]
\label{I.5.5.8}
Seja $Y$ um esquema afim; para que $f:X\to Y$ seja um morfismo separado, é
necessário e suficiente que $X$ seja separado \emph{(isto é, que $X$ seja um
\emph{esquema})}.
\end{corollary}

\begin{proof}
Com efeito, observa-se que o critério de \sref{I.5.5.6} não depende de $f$.
\end{proof}

\begin{corollary}[5.5.9]
\label{I.5.5.9}
Para que um morfismo $f:X\to Y$ seja separado, é necessário que o pré-esquema
induzido $f^{-1}(U)$ seja separado para todo aberto $U$ sobre o qual $Y$ induz
um pré-esquema separado, e é suficiente que isso ocorra para todo aberto afim
$U\subset Y$.
\end{corollary}

\begin{proof}
A necessidade da condição decorre de \sref{I.5.5.5}\footnote{\emph{[Trad.]
Nas duas referências a (5.5.5) desta demonstração, a fonte francesa imprime
(5.5.4); o argumento aplica o caráter local sobre a base estabelecido na
Proposição~\hyperref[I.5.5.5]{5.5.5}.}} e \sref{I.5.5.1}[ii]; a suficiência
decorre de \sref{I.5.5.5} e \sref{I.5.5.8}, levando em conta a existência de
coberturas de $Y$ por abertos afins.
\end{proof}

Em particular, se $X$ e $Y$ são esquemas afins, então \emph{todo} morfismo
$X\to Y$ é separado.

\begin{proposition}[5.5.10]
\label{I.5.5.10}
Seja $Y$ um \emph{esquema}, e seja $f:X\to Y$ um morfismo.
Para todo aberto afim $U$ de $X$ e todo aberto afim $V$ de $Y$,
$U\cap f^{-1}(V)$ é afim.
\end{proposition}

\begin{proof}
Sejam $p_1$ e $p_2$ as projeções de $X\times_\bb{Z}Y$;
o subespaço $U\cap f^{-1}(V)$ é a imagem de
$\Gamma_f(X)\cap p_1^{-1}(U)\cap p_2^{-1}(V)$ por $p_1$.
Ora, $p_1^{-1}(U)\cap p_2^{-1}(V)$ identifica-se com o espaço subjacente ao
\oldpage[I]{139}
pré-esquema $U\times_\bb{Z}V$ \sref{I.3.2.7} e, portanto, é um esquema afim
\sref{I.3.2.2}; como $\Gamma_f(X)$ é fechado em $X\times_\bb{Z}Y$
\sref{I.5.4.3}, $\Gamma_f(X)\cap p_1^{-1}(U)\cap p_2^{-1}(V)$ é fechado em
$U\times_\bb{Z} V$, de modo que o pré-esquema induzido pelo subpré-esquema de
$X\times_\bb{Z}Y$ associado a $\Gamma_f$ \sref{I.4.2.1} sobre o aberto
$\Gamma_f(X)\cap p_1^{-1}(U)\cap p_2^{-1}(V)$ de seu espaço subjacente é um
subpré-esquema fechado de um esquema afim e, portanto, um esquema afim
\sref{I.4.2.3}.
A proposição resulta então do fato de que $\Gamma_f$ é uma imersão.
\end{proof}

\begin{examples}[5.5.11]
\label{I.5.5.11}
O pré-esquema do Exemplo~\sref{I.2.3.2} (``a reta projetiva sobre um corpo
$K$'') é \emph{separado}, pois, para a cobertura $(X_1,X_2)$ de $X$ por
abertos afins, $X_1\cap X_2=U_{12}$ é afim e
$\Gamma(U_{12},\sh{O}_X)$, o anel das frações racionais da forma
$f(s)/s^m$, com $f\in K[s]$, é gerado por $K[s]$ e $1/s$, de modo que as
condições de \sref{I.5.5.6} são satisfeitas.

Com a mesma escolha de $X_1$, $X_2$, $U_{12}$ e $U_{21}$ que no
Exemplo~\sref{I.2.3.2}, tomemos agora como $u_{12}$ o isomorfismo que envia
$f(s)$ a $f(t)$;
obtém-se então, por colagem, um pré-esquema \emph{integral não separado} $X$,
pois a primeira condição de \sref{I.5.5.6} é satisfeita, mas a segunda não.
É imediato aqui que
$\Gamma(X,\sh{O}_X)\to\Gamma(X_1,\sh{O}_X)=K[s]$ é um isomorfismo;
o isomorfismo inverso define um morfismo sobrejetivo
$f:X\to\Spec(K[s])$ e, para todo $y\in\Spec(K[s])$ tal que
$\mathfrak{j}_y\neq(s)$, $f^{-1}(y)$ consiste em um único ponto, ao passo que,
para $\mathfrak{j}_y=(s)$, $f^{-1}(y)$ consiste em \emph{dois} pontos
distintos\footnote{\emph{[Trad.] A fonte francesa imprime $(0)$ em ambas as
condições sobre o ideal. Para a origem duplicada de $\Spec(K[s])$, o ideal
primo pertinente é $(s)$; $(0)$ é o ponto genérico.}} (diz-se que $X$ é a
``reta afim sobre $K$ com o ponto $0$ duplicado'').

Também se pode dar um exemplo no qual a primeira das duas condições de
\sref{I.5.5.6} não é satisfeita.
Observemos primeiro que, no espectro primo $Y$ do anel $A=K[s,t]$ de
polinômios em duas indeterminadas sobre um corpo $K$, o aberto $U$ dado pela
união de $D(s)$ e $D(t)$ \emph{não é um aberto afim}.
Com efeito, se $z$ é uma seção de $\sh{O}_Y$ sobre $U$, existem dois inteiros
$m,n\geq0$ tais que $s^mz$ e $t^nz$ são as restrições de polinômios em $s$ e
$t$ a $U$ \sref{I.1.4.1}, o que evidentemente só é possível se a seção $z$ se
prolonga a uma seção sobre todo $Y$, identificada com um polinômio em $s$ e
$t$.
Se $U$ fosse um aberto afim, o morfismo de injeção $U\to Y$ seria um
isomorfismo \sref{I.1.7.3}, o que é absurdo, pois $U\neq Y$.

Tendo em conta o que precede, tomemos dois esquemas afins $Y_1$ e $Y_2$,
espectros primos dos anéis $A_1=K[s_1,t_1]$ e $A_2=K[s_2,t_2]$,
respectivamente;
tomemos $U_{12}=D(s_1)\cup D(t_1)$ e $U_{21}=D(s_2)\cup D(t_2)$, e tomemos
como $u_{12}$ a restrição de um isomorfismo $Y_2\to Y_1$ a $U_{21}$
correspondente ao isomorfismo de anéis que envia $f(s_1,t_1)$ a
$f(s_2,t_2)$;
obtém-se então um exemplo no qual a primeira condição de \sref{I.5.5.6} não é
satisfeita, enquanto a segunda o é\footnote{\emph{[Trad.] A fonte francesa
afirma que nenhuma das duas condições de \sref{I.5.5.6} é satisfeita.
Contudo, para $U=D(s)\cup D(t)$, tem-se
$\Gamma(U,\sh{O}_Y)=K[s,t]$, de modo que a segunda condição é satisfeita;
somente a afinidade de $U$ falha.}} (o pré-esquema integral assim obtido é
chamado de ``plano afim sobre $K$ com o ponto $0$ duplicado'').
\end{examples}

\begin{remark}[5.5.12]
\label{I.5.5.12}
Dada uma propriedade \textbf{P} dos morfismos de pré-esquemas, consideremos
as proposições seguintes:
\begin{enumerate}
  \item[\rm(i)] \emph{Toda imersão fechada possui a propriedade~\textbf{P}.}
  \item[\rm(ii)] \emph{O composto de dois morfismos que possuem a
    propriedade~\textbf{P} possui a propriedade~\textbf{P}.}
  \item[\rm(iii)] \emph{Se $f:X\to X'$ e $g:Y\to Y'$ são dois
    $S$-morfismos que possuem a propriedade~\textbf{P}, então
    $f\times_S g$ possui a propriedade~\textbf{P}.}

\oldpage[I]{140}
  \item[\rm(iv)] \emph{Se $f:X\to Y$ é um $S$-morfismo que possui a
    propriedade~\textbf{P}, então todo $S'$-morfismo $f_{(S')}$ obtido por
    uma extensão $S'\to S$ do pré-esquema de base possui a
    propriedade~\textbf{P}.}
  \item[\rm(v)] \emph{Se o composto $g\circ f$ de dois morfismos
    $f:X\to Y$ e $g:Y\to Z$ possui a propriedade~\textbf{P}, e se $g$ é
    separado, então $f$ possui a propriedade~\textbf{P}.}
  \item[\rm(vi)] \emph{Se um morfismo $f:X\to Y$ possui a
    propriedade~\textbf{P}, então $f_\red$ também a possui
    (\sref{I.5.1.5}).}
\end{enumerate}

Nessas condições, supondo-se verificadas (i) e (ii), (iii) e (iv) são
\emph{equivalentes}, e (v) e (vi) são \emph{consequências} de (i), (ii) e
(iii).

A primeira afirmação já foi demonstrada (\sref{I.3.5.1}).
Consideremos a fatoração (\sref{I.5.3.13})
$X\xrightarrow{\Gamma_f}X\times_Z Y\xrightarrow{p_2}Y$ de $f$; a relação
$p_2=(g\circ f)\times_Z 1_Y$ mostra que, se $g\circ f$ possui a
propriedade~\textbf{P}, então $p_2$ também a possui, em virtude de (iii); se
$g$ é separado, $\Gamma_f$ é uma imersão fechada (\sref{I.5.4.3}) e,
portanto, também possui a propriedade~\textbf{P} por (i); por fim, em virtude
de (ii), $f$ possui a propriedade~\textbf{P}.

Por fim, consideremos o diagrama comutativo
\[
  \xymatrix{
    X_\red\ar[r]^{f_\red}\ar[d] &
    Y_\red\ar[d]\\
    X\ar[r]_f &
    Y,
  }
\]
no qual as setas verticais são imersões fechadas (\sref{I.5.1.5}) e,
portanto, possuem a propriedade~\textbf{P} por (i). A hipótese de que $f$
possui a propriedade~\textbf{P} implica, por (ii), que
$X_\red\xrightarrow{f_\red}Y_\red\to Y$ possui a propriedade~\textbf{P};
por fim, como uma imersão fechada é separada (\sref{I.5.5.1}[i]),
$f_\red$ possui a propriedade~\textbf{P} em virtude de (v).

Observe-se que, se considerarmos as proposições:
\begin{enumerate}
  \item[\rm(i')] \emph{Toda imersão possui a propriedade~\textbf{P}.}
  \item[\rm(v')] \emph{Se $g\circ f$ possui a propriedade~\textbf{P},
    então $f$ também a possui;}
\end{enumerate}
então o raciocínio feito acima mostra que (v') é consequência de (i'), (ii) e
(iii).
\end{remark}

\begin{env}[5.5.13]
\label{I.5.5.13}
Observe-se que (v) e (vi) são ainda consequências de (i), (iii) e
\begin{enumerate}
  \item[\rm(ii')] \emph{Se $j:X\to Y$ é uma imersão fechada e
    $g:Y\to Z$ um morfismo que possui a propriedade~\textbf{P}, então
    $g\circ j$ possui a propriedade~\textbf{P}.}
\end{enumerate}
Do mesmo modo, (v') é consequência de (i'), (iii) e
\begin{enumerate}
  \item[\rm(ii'')] \emph{Se $j:X\to Y$ é uma imersão e $g:Y\to Z$ um
    morfismo que possui a propriedade~\textbf{P}, então $g\circ j$ possui a
    propriedade~\textbf{P}.}
\end{enumerate}
Isso resulta imediatamente dos raciocínios de \sref{I.5.5.12}.
\end{env}

\section{Condições de finitude}
\label{section:I.6}

\subsection{Pré-esquemas noetherianos e localmente noetherianos}
\label{subsection:I.6.1}

\begin{definition}[6.1.1]
\label{I.6.1.1}
Diz-se que um pré-esquema $X$ é \emph{noetheriano} (resp.
\emph{localmente noetheriano}) se é reunião finita (resp. reunião) de abertos
afins $V_\alpha$ tais que o anel do esquema induzido sobre cada $V_\alpha$ é
noetheriano.
\end{definition}

Segue imediatamente de \sref{I.1.5.2} que, se $X$ é localmente
noetheriano, o feixe estrutural $\sh{O}_X$ é um \emph{feixe coerente de
anéis}, pois a questão é local. Todo \emph{$\sh{O}_X$-submódulo
quase-coerente}
\oldpage[I]{141}
(resp. todo $\sh{O}_X$-módulo quociente quase-coerente) de um
$\sh{O}_X$-módulo \emph{coerente} $\sh{F}$ é então \emph{coerente}, pois a
questão é novamente local e basta aplicar \sref{I.1.5.1},
\sref{I.1.4.1} e \sref{I.1.3.10}, juntamente com o fato de que um submódulo
(resp. módulo quociente) de um módulo de tipo finito sobre um anel
noetheriano é de tipo finito. Em particular, todo \emph{feixe de ideais
quase-coerente} de $\sh{O}_X$ é \emph{coerente}.

Se um pré-esquema $X$ é reunião finita (resp. reunião) de abertos
$W_\lambda$ tais que os pré-esquemas induzidos sobre os $W_\lambda$ sejam
noetherianos (resp. localmente noetherianos), é claro que $X$ é noetheriano
(resp. localmente noetheriano).

\begin{proposition}[6.1.2]
\label{I.6.1.2}
Para que um pré-esquema $X$ seja noetheriano, é necessário e suficiente que
seja localmente noetheriano e que seu espaço subjacente seja quase-compacto.
O próprio espaço subjacente é então também noetheriano.
\end{proposition}

\begin{proof}
A primeira afirmação decorre imediatamente das definições e de
\sref{I.1.1.10}[ii]. A segunda decorre de \sref{I.1.1.6} e do fato de que
todo espaço que é reunião finita de subespaços noetherianos é, por sua vez,
noetheriano \sref[0]{0.2.2.3}.
\end{proof}

\begin{proposition}[6.1.3]
\label{I.6.1.3}
Seja $X$ um esquema afim de anel $A$.
As condições seguintes são equivalentes:
\emph{(a)} $X$ é noetheriano;
\emph{(b)} $X$ é localmente noetheriano;
\emph{(c)} $A$ é noetheriano.
\end{proposition}

\begin{proof}
A equivalência entre (a) e (b) decorre de \sref{I.6.1.2} e do fato de que o
espaço subjacente de todo esquema afim é quase-compacto
\sref{I.1.1.10}; é claro, além disso, que (c) implica (a). Para ver que (a)
implica (c), observemos que existe um recobrimento finito $(V_i)$ de $X$ por
abertos afins tal que o anel $A_i$ do pré-esquema induzido sobre $V_i$ é
noetheriano. Seja, então, $(\mathfrak{a}_n)$ uma sequência crescente de
ideais de $A$; pela correspondência biunívoca canônica \sref{I.1.3.7},
corresponde-lhe uma sequência crescente $(\widetilde{\mathfrak{a}}_n)$ de
feixes de ideais em $\widetilde{A}\equiv\sh{O}_X$; para ver que a sequência
$(\mathfrak{a}_n)$ é estacionária, basta provar que a sequência
$(\widetilde{\mathfrak{a}}_n)$ o é. Ora, a restrição
$\widetilde{\mathfrak{a}}_n|V_i$ é um feixe quase-coerente de ideais de
$\sh{O}_X|V_i$, pois é a imagem inversa de $\widetilde{\mathfrak{a}}_n$ pela
injeção canônica $V_i\to X$ \sref[0]{0.5.1.4}; portanto,
$\widetilde{\mathfrak{a}}_n|V_i$ é da forma
$\widetilde{\mathfrak{a}}_{ni}$, onde $\mathfrak{a}_{ni}$ é um ideal de
$A_i$ \sref{I.1.3.7}. Como $A_i$ é noetheriano, a sequência
$(\mathfrak{a}_{ni})$ é estacionária para todo $i$, de onde resulta a
proposição.
\end{proof}

Observe-se que o raciocínio precedente prova também que \emph{se $X$ é um
pré-esquema noetheriano, toda sequência crescente de feixes coerentes de
ideais de $\sh{O}_X$ é estacionária}.

\begin{proposition}[6.1.4]
\label{I.6.1.4}
Todo subpré-esquema de um pré-esquema noetheriano (resp. localmente
noetheriano) é noetheriano (resp. localmente noetheriano).
\end{proposition}

\begin{proof}
Basta fazer a demonstração para um pré-esquema noetheriano $X$; além disso,
pela definição \sref{I.6.1.1}, somos imediatamente reduzidos ao caso em que
$X$ é um esquema afim. Como todo subpré-esquema de $X$ é um subpré-esquema
fechado de um pré-esquema induzido sobre um aberto \sref{I.4.1.3}, podemos
limitar-nos a considerar o caso de um subpré-esquema $Y$ fechado ou induzido
sobre um aberto de $X$. O caso em que $Y$ é fechado é imediato, pois, se $A$
é o anel de $X$, sabemos que $Y$ é um esquema afim de anel
$A/\mathfrak{J}$, onde $\mathfrak{J}$ é um ideal de $A$
\sref{I.4.2.3}; como $A$ é noetheriano \sref{I.6.1.3}, o mesmo vale para
$A/\mathfrak{J}$.

Suponhamos agora que $Y$ seja aberto em $X$; o espaço subjacente de $Y$ é
noetheriano \sref{I.6.1.2}, portanto quase-compacto e, por conseguinte, uma
união finita de abertos $D(f_i)$ ($f_i\in A$);
\oldpage[I]{142}
tudo se reduz a demonstrar a proposição quando $Y=D(f)$, com $f\in A$. Mas,
então, $Y$ é um esquema afim cujo anel é isomorfo a $A_f$
\sref{I.1.3.6}; como $A$ é noetheriano \sref{I.6.1.3}, o mesmo vale para
$A_f$.
\end{proof}

\begin{env}[6.1.5]
\label{I.6.1.5}
Observe-se que o \emph{produto} de dois $S$-pré-esquemas noetherianos não é
necessariamente noetheriano, mesmo que esses pré-esquemas sejam afins, pois o
produto tensorial de duas álgebras noetherianas não é necessariamente um anel
noetheriano (cf. \sref{I.6.3.8}).
\end{env}

\begin{proposition}[6.1.6]
\label{I.6.1.6}
Se $X$ é um pré-esquema noetheriano, o nilradical $\sh{N}_X$ de
$\sh{O}_X$ é nilpotente.
\end{proposition}

\begin{proof}
Com efeito, pode-se recobrir $X$ por um número finito de abertos afins $U_i$,
e basta provar que existem inteiros $n_i$ tais que
$(\sh{N}_X|U_i)^{n_i}=0$; se $n$ é o maior dos $n_i$, teremos então
$\sh{N}_X^n=0$. Somos, portanto, reduzidos ao caso em que $X=\Spec(A)$ é
afim, sendo $A$ um anel noetheriano; em virtude de \sref{I.5.1.1} e de
\sref{I.1.3.13}, basta observar que o nilradical de $A$ é nilpotente
(\cite[p.~127, cor.~4]{I-11}).
\end{proof}

\begin{corollary}[6.1.7]
\label{I.6.1.7}
Seja $X$ um pré-esquema noetheriano;
para que $X$ seja um esquema afim, é necessário e suficiente que $X_\red$ seja um esquema afim.
\end{corollary}

\begin{proof}
Isso resulta de \sref{I.6.1.6} e de \sref{I.5.1.10}.
\end{proof}

\begin{lemma}[6.1.8]
\label{I.6.1.8}
Sejam $X$ um espaço topológico, $x$ um ponto de $X$ e $U$ uma vizinhança
aberta de $x$ que tenha apenas um número finito de componentes irredutíveis.
Então existe uma vizinhança $V$ de $x$ tal que toda vizinhança aberta de $x$
contida em $V$ seja conexa.
\end{lemma}

\begin{proof}
Sejam $U_i$ ($1\leq i\leq m$) as componentes irredutíveis de $U$ que não
contêm $x$; o complementar (em $U$) da união dos $U_i$ é uma vizinhança
aberta $V$ de $x$ em $U$ e, portanto, também em $X$; é, além disso, a
interseção com $U$ do complementar (em $X$) da união das componentes
irredutíveis de $X$ que não contêm $x$ \sref[0]{0.2.1.6}\footnote{\emph{[Trad.]
A fonte francesa identifica aqui $V$ diretamente com o complementar em $X$ da
união das componentes irredutíveis que não contêm $x$, omitindo a interseção
com $U$. A igualdade impressa é falsa, por exemplo, se $X$ é irredutível e
$U\subsetneq X$: nesse caso, aquele complementar é $X$, enquanto $V=U$.}}.
Seja então $W$ uma vizinhança aberta de $x$ contida em $V$. As componentes
irredutíveis de $W$ são as interseções com $W$ das componentes irredutíveis de
$U$ que intersectam $W$ \sref[0]{0.2.1.6}, de modo que essas componentes
contêm $x$; como são conexas, $W$ também o é.
\end{proof}

\begin{corollary}[6.1.9]
\label{I.6.1.9}
Um espaço topológico localmente noetheriano é localmente conexo (o que
implica, entre outras coisas, que suas componentes conexas são abertas).
\end{corollary}

\begin{proposition}[6.1.10]
\label{I.6.1.10}
Seja $X$ um espaço topológico localmente noetheriano.
As condições seguintes são equivalentes.
\begin{enumerate}
  \item[{\rm(a)}] As componentes irredutíveis de $X$ são abertas.
  \item[{\rm(b)}] As componentes irredutíveis de $X$ são exatamente suas
    componentes conexas.
  \item[{\rm(c)}] As componentes conexas de $X$ são irredutíveis.
  \item[{\rm(d)}] Duas componentes irredutíveis distintas de $X$ têm
    interseção vazia.
\end{enumerate}
Por fim, se $X$ é um pré-esquema, estas condições são também equivalentes a
\begin{enumerate}
  \item[{\rm(e)}] Para todo $x\in X$, $\Spec(\sh{O}_x)$ é irredutível (ou,
    dito de outro modo, o nilradical de $\sh{O}_x$ é primo).
\end{enumerate}
\end{proposition}

\begin{proof}
É imediato que (a) implica (b), pois todo espaço irredutível é conexo, e (a)
implica que as componentes irredutíveis de $X$ são os conjuntos ao mesmo tempo
abertos e fechados. É trivial que (b) implica (c); reciprocamente, um conjunto
fechado $F$ que contém
\oldpage[I]{143}
uma componente conexa $C$ de $X$, sendo $C$ distinta de $F$, não pode ser
irredutível, pois, por não ser conexo, $F$ é a união de dois conjuntos não
vazios disjuntos, ao mesmo tempo abertos e fechados em $F$ e, portanto,
fechados em $X$; consequentemente, (c) implica (b). Disso se deduz
imediatamente que (c) implica (d), pois duas componentes conexas distintas não
têm nenhum ponto em comum.

Até aqui, não utilizamos o fato de que $X$ é localmente noetheriano. Suponhamos
agora que o seja e mostremos que (d) implica (a): por \sref[0]{0.2.1.6},
podemos nos restringir ao caso em que o espaço $X$ é noetheriano e, portanto,
tem apenas um número finito de componentes irredutíveis. Como estas são
fechadas e duas a duas disjuntas, são abertas.

Por fim, a equivalência entre (d) e (e) é válida mesmo sem supor que o espaço
subjacente ao pré-esquema $X$ seja localmente noetheriano. Com efeito, podemos
nos restringir ao caso em que $X=\Spec(A)$ é afim, por \sref[0]{0.2.1.6};
dizer que $x$ está contido em uma única componente irredutível de $X$ equivale
a dizer que $\mathfrak{j}_x$ contém um único ideal primo minimal de $A$
\sref{I.1.1.14}, o que equivale a dizer que
$\mathfrak{j}_x\sh{O}_x$ contém um único ideal primo minimal de $\sh{O}_x$,
de onde se segue a conclusão.
\end{proof}

\begin{corollary}[6.1.11]
\label{I.6.1.11}
Seja $X$ um espaço localmente noetheriano. Para que $X$ seja irredutível, é
necessário e suficiente que $X$ seja conexo e não vazio e que quaisquer duas
componentes irredutíveis distintas de $X$ tenham interseção vazia. Se $X$ é um
pré-esquema, esta última condição equivale a que $\Spec(\sh{O}_x)$ seja
irredutível para todo $x\in X$.
\end{corollary}

\begin{proof}
A segunda afirmação já foi demonstrada em \sref{I.6.1.10}; resta, portanto,
demonstrar que as condições da primeira afirmação são suficientes. Mas, por
\sref{I.6.1.10}, essas condições implicam que as componentes irredutíveis de
$X$ são exatamente suas componentes conexas e, como $X$ é conexo e não vazio,
é irredutível.
\end{proof}

\begin{corollary}[6.1.12]
\label{I.6.1.12}
Seja $X$ um pré-esquema localmente noetheriano. Para que $X$ seja integral, é
necessário e suficiente que $X$ seja conexo e não vazio\footnote{\emph{[Trad.]
A fonte francesa omite «e não vazio». Sem essa condição, o caso vazio
satisfaria a condição sobre todos os anéis locais por não haver ponto algum,
embora não fosse irredutível; além disso, o Corolário~\sref{I.6.1.11} exige
explicitamente que $X$ seja não vazio.}}, e que $\sh{O}_x$ seja um domínio
integral para todo $x\in X$.
\end{corollary}

\begin{proposition}[6.1.13]
\label{I.6.1.13}
Seja $X$ um pré-esquema localmente noetheriano, e seja $x\in X$ um ponto tal
que o nilradical $\sh{N}_x$ de $\sh{O}_x$ seja primo (resp. que $\sh{O}_x$
seja reduzido, resp. integral); então existe uma vizinhança aberta $U$ de $x$
que é irredutível (resp. reduzida, resp. integral).
\end{proposition}

\begin{proof}
Basta considerar os dois casos em que $\sh{N}_x$ é primo e em que
$\sh{N}_x=0$, pois a terceira hipótese é a conjunção das duas primeiras. Se
$\sh{N}_x$ é primo, $x$ pertence a uma única componente irredutível $Y$ de
$X$ \sref{I.6.1.10}; a união das componentes irredutíveis de $X$ que não
contêm $x$ é fechada (pois a família dessas componentes é localmente finita),
e o complemento $U$ dessa união é, portanto, aberto e está contido em $Y$,
logo é irredutível \sref[0]{0.2.1.6}. Se $\sh{N}_x=0$, tem-se também
$\sh{N}_y=0$ para todo $y$ em uma vizinhança de $x$, porque $\sh{N}$ é
quase-coerente \sref{I.5.1.1} e, portanto, coerente, já que $X$ é localmente
noetheriano; a conclusão segue então de \sref[0]{0.5.2.2}.
\end{proof}

\subsection{Pré-esquemas artinianos}
\label{subsection:I.6.2}
\label{I.6.2}

\begin{definition}[6.2.1]
\label{I.6.2.1}
Diz-se que um pré-esquema é \emph{artiniano} se é afim e seu anel é artiniano.
\end{definition}

\oldpage[I]{144}
\begin{proposition}[6.2.2]
\label{I.6.2.2}
Dado um pré-esquema $X$, as condições seguintes são equivalentes:
\begin{enumerate}
  \item[{\rm(a)}] $X$ é um esquema artiniano;
  \item[{\rm(b)}] $X$ é noetheriano e seu espaço subjacente é discreto;
  \item[{\rm(c)}] $X$ é noetheriano e os pontos de seu espaço subjacente são fechados \emph{(condição $\mathrm{T}_1$)}.
\end{enumerate}
Quando se verifica qualquer uma das condições anteriores, o espaço subjacente a $X$ é finito, e o anel $A$ de $X$ é o produto direto dos anéis locais (artinianos) nos pontos de $X$.
\end{proposition}

\begin{proof}
Sabe-se que (a) implica a última afirmação (\cite[p.~205, th.~3]{I-13}); por conseguinte, todo ideal primo de $A$ é maximal e é a imagem inversa de um ideal maximal de uma das componentes locais de $A$, de modo que o espaço $X$ é finito e discreto;
(a) implica, pois, (b), e (b) implica evidentemente (c).
Para ver que (c) implica (a), mostremos primeiro que $X$ é então finito;
podemos, com efeito, reduzir-nos ao caso em que $X$ é afim, e sabe-se que um anel noetheriano cujos ideais primos são todos maximais é artiniano (\cite[p.~203]{I-13}), de onde resulta a afirmação.
O espaço subjacente $X$ é então discreto, soma topológica de um número finito de pontos $x_i$, e os anéis locais $\sh{O}_{x_i}=A_i$ são artinianos;
é claro que $X$ é isomorfo ao esquema afim $\Spec(A)$, onde $A$ é o produto direto dos $A_i$ \sref{I.1.7.3}.
\end{proof}

\subsection{Morfismos de tipo finito}
\label{subsection:I.6.3}

\begin{definition}[6.3.1]
\label{I.6.3.1}
Diz-se que um morfismo $f:X\to Y$ é \emph{de tipo finito} se $Y$ é a união de uma família $(V_\alpha)$ de abertos afins que têm a propriedade seguinte:
\begin{enumerate}
  \item[(P)] $f^{-1}(V_\alpha)$ é uma união finita de abertos afins $U_{\alpha i}$ tais que cada anel $A(U_{\alpha i})$ é uma álgebra de tipo finito sobre $A(V_\alpha)$.
\end{enumerate}
Diz-se então também que $X$ é um pré-esquema de tipo finito sobre $Y$, ou um $Y$-pré-esquema de tipo finito.
\end{definition}

\begin{proposition}[6.3.2]
\label{I.6.3.2}
Se $f:X\to Y$ é um morfismo de tipo finito, todo aberto afim $W$ de $Y$ possui a propriedade \emph{(P)} da def. \sref{I.6.3.1}.
\end{proposition}

Demonstremos primeiro o seguinte lema.
\begin{lemma}[6.3.2.1]
\label{I.6.3.2.1}
Se $T\subset Y$ é um aberto afim que possui a propriedade \emph{(P)}, então, para todo $g\in A(T)$, $D(g)$ também possui a propriedade \emph{(P)}.
\end{lemma}

\begin{proof}
Por hipótese, $f^{-1}(T)$ é uma união finita de abertos afins $Z_j$ tais que $A(Z_j)$ é uma álgebra de tipo finito sobre $A(T)$;
seja $\vphi_j:A(T)\to A(Z_j)$ o homomorfismo de anéis correspondente à restrição de $f$ a $Z_j$ \sref{I.2.2.4}, e ponhamos $g_j=\vphi_j(g)$;
tem-se então $f^{-1}(D(g))\cap Z_j=D(g_j)$ \sref{I.1.2.2.2}.
Ora, $A(D(g_j))=A(Z_j)_{g_j}=A(Z_j)[1/g_j]$ é de tipo finito sobre $A(Z_j)$ e, \emph{a fortiori}, sobre $A(T)$ em virtude da hipótese, e portanto também sobre $A(D(g))=A(T)[1/g]$, o que demonstra o lema.
\end{proof}

\begin{proof}
Em virtude do lema anterior, como $W$ é quase-compacto \sref{I.1.1.10}, existe um recobrimento finito de $W$ por conjuntos da forma \erratum[II]{$D(g_i)\subset W$}, onde cada $g_i$ pertence a um anel $A(V_{\alpha(i)})$.
Cada $D(g_i)$, sendo quase-compacto, é uma união finita de conjuntos $D(h_{ik})$, onde $h_{ik}\in A(W)$;
se $\vphi_i:A(W)\to A(D(g_i))$ é a aplicação canônica, tem-se $D(h_{ik})=D(\vphi_i(h_{ik}))$ em virtude de \sref{I.1.2.2.2}.
Em virtude de \sref{I.6.3.2.1}, cada um dos $f^{-1}(D(h_{ik}))$ admite um recobrimento finito por abertos afins $U_{ijk}$ tais que os $A(U_{ijk})$ são álgebras de tipo finito sobre $A(D(h_{ik}))=A(W)[1/h_{ik}]$, donde resulta a proposição.
\end{proof}

\oldpage[I]{145}
Pode-se, portanto, dizer que a noção de pré-esquema de tipo finito sobre $Y$ é \emph{local em $Y$}.

\begin{proposition}[6.3.3]
\label{I.6.3.3}
Sejam $X$ e $Y$ dois esquemas afins;
para que $X$ seja de tipo finito sobre $Y$, é necessário e suficiente que $A(X)$ seja uma álgebra de tipo finito sobre $A(Y)$.
\end{proposition}

\begin{proof}
Sendo a condição evidentemente suficiente, provemos que é necessária.
Ponhamos $A=A(Y)$ e $B=A(X)$;
em virtude de \sref{I.6.3.2}, existe uma cobertura aberta afim finita $(V_i)$ de $X$ tal que cada um dos anéis $A(V_i)$ é uma $A$-álgebra de tipo finito.
Além disso, como os $V_i$ são quase-compactos, podemos cobrir cada um deles por um número finito de abertos da forma $D(g_{ij})\subset V_i$, onde $g_{ij}\in B$;
se $\vphi_i$ é o homomorfismo $B\to A(V_i)$ correspondente à injeção canônica $V_i\to X$, então temos $B_{g_{ij}}=(A(V_i))_{\vphi_i(g_{ij})}=A(V_i)[1/\vphi_i(g_{ij})]$, de modo que $B_{g_{ij}}$ é uma $A$-álgebra de tipo finito.
Podemos, portanto, reduzir-nos ao caso em que $V_i=D(g_i)$ com $g_i\in B$.
Por hipótese, existe um subconjunto finito $F_i$ de $B$ e um inteiro $n_i\geq0$ tais que $B_{g_i}$ é a álgebra gerada sobre $A$ pelos elementos $b_i/g_i^{n_i}$, onde $b_i$ percorre $F_i$.
Como os $g_i$ são em número finito, podemos supor que todos os $n_i$ são iguais a um mesmo inteiro $n$.
Além disso, como os $D(g_i)$ formam uma cobertura de $X$, o ideal gerado em $B$ pelos $g_i$ é igual a $B$ ou, em outras palavras, existem $h_i\in B$ tais que $\sum_i h_ig_i=1$.
Seja então $F$ o subconjunto finito de $B$ formado pela união dos $F_i$, pelo conjunto dos $g_i$ e pelo conjunto dos $h_i$; mostraremos que o subanel $B'=A[F]$ de $B$ é igual a $B$.
Por hipótese, para todo $b\in B$ e todo $i$, a imagem canônica de $b$ em $B_{g_i}$ é da forma $b'_i/g_i^{m_i}$, onde $b'_i\in B'$;
multiplicando os $b'_i$ por potências convenientes dos $g_i$, podemos novamente supor que todos os $m_i$ são iguais a um mesmo inteiro $m$.
Por definição dos anéis de frações, existe, portanto, um inteiro $N$ (dependente de $b$) tal que $N\geq m$ e $g_i^Nb\in B'$ para todo $i$;
mas, \emph{no anel $B'$}, os $g_i^N$ geram o ideal $B'$, pois os $g_i$ já o geram (e os $h_i$ pertencem a $B'$);
existem, portanto, $c_i\in B'$ tais que $\sum_i c_ig_i^N=1$, daí $b=\sum_ic_ig_i^Nb\in B'$, Q.E.D.
\end{proof}

\begin{proposition}[6.3.4]
\label{I.6.3.4}
\medskip\noindent
\begin{enumerate}
  \item[{\rm(i)}] Toda imersão fechada é de tipo finito.
  \item[{\rm(ii)}] O composto de dois morfismos de tipo finito é de tipo finito.
  \item[{\rm(iii)}] Se $f:X\to X'$ e $g:Y\to Y'$ são dois $S$-morfismos de tipo finito, então $f\times_S g$ é de tipo finito.
  \item[{\rm(iv)}] Se $f:X\to Y$ é um $S$-morfismo de tipo finito, então $f_{(S')}$ é de tipo finito para toda extensão $g:S'\to S$ do pré-esquema de base.
  \item[{\rm(v)}] Se o composto $g\circ f$ de dois morfismos é de tipo finito, sendo $g$ separado, então $f$ é de tipo finito.
  \item[{\rm(vi)}] Se um morfismo $f$ é de tipo finito, então $f_\red$ também o é.
\end{enumerate}
\end{proposition}

\begin{proof}
Em virtude de \sref{I.5.5.12}, basta demonstrar (i), (ii) e (iv).

Para demonstrar (i), podemos limitar-nos ao caso de uma injeção canônica
$X\to Y$, sendo $X$ um subpré-esquema fechado de $Y$; além disso, por
\sref{I.6.3.2}, podemos supor que $Y$ é afim, caso em que $X$ também é afim
\sref{I.4.2.3} e seu anel é isomorfo a um anel quociente
$A/\mathfrak{J}$, onde $A$ é o anel de $Y$ e $\mathfrak{J}$ é um ideal de
$A$; como $A/\mathfrak{J}$ é de tipo finito sobre $A$, segue-se a conclusão.

Demonstremos agora (ii). Sejam $f:X\to Y$ e $g:Y\to Z$ dois morfismos de
tipo finito, e seja $U$ um aberto afim de $Z$; $g^{-1}(U)$ admite uma
cobertura finita por abertos afins $V_i$ tais que cada $A(V_i)$ é uma
álgebra de tipo finito sobre $A(U)$ \sref{I.6.3.2}; do mesmo modo,
\oldpage[I]{146}
cada um dos $f^{-1}(V_i)$ admite uma cobertura finita por abertos afins
$W_{ij}$ tais que cada $A(W_{ij})$ é uma álgebra de tipo finito sobre
$A(V_i)$ e, portanto, também uma álgebra de tipo finito sobre $A(U)$, daí a
conclusão.

Finalmente, para demonstrar (iv), podemos limitar-nos ao caso em que $S=Y$;
de fato, $f_{(S')}$ também é igual a $f_{(Y_{(S')})}$, considerando-se $f$
como um $Y$-morfismo, e a extensão de base é $Y_{(S')}\to Y$
\sref{I.3.3.9}. Sejam então $p$ e $q$ as projeções $X_{(S')}\to X$ e
$X_{(S')}\to S'$. Seja $V$ um aberto afim de $S$; $f^{-1}(V)$ é uma união
finita de abertos afins $W_i$ tais que $A(W_i)$ é uma álgebra de tipo finito
sobre $A(V)$ \sref{I.6.3.2}. Seja $V'$ um aberto afim de $S'$ contido em
$g^{-1}(V)$; como $f\circ p=g\circ q$, $q^{-1}(V')$ está contido na união
dos $p^{-1}(W_i)$; por outro lado, a interseção
$p^{-1}(W_i)\cap q^{-1}(V')$ identifica-se com o produto
$W_i\times_V V'$ \sref{I.3.2.7}, que é um esquema afim cujo anel é isomorfo
a $A(W_i)\otimes_{A(V)}A(V')$ \sref{I.3.2.2}; este anel é, por hipótese, uma
álgebra de tipo finito sobre $A(V')$, o que demonstra a proposição.
\end{proof}

\begin{corollary}[6.3.5]
\label{I.6.3.5}
Seja $f:X\to Y$ um morfismo de imersão.
Se o espaço subjacente a $Y$ (resp. $X$) é localmente noetheriano (resp.
noetheriano), então $f$ é de tipo finito.
\end{corollary}

\begin{proof}
Podemos sempre supor que $Y$ é afim \sref{I.6.3.2};
se o espaço subjacente a $Y$ é localmente noetheriano, podemos, além disso,
supô-lo noetheriano, e então o espaço subjacente a $X$, que é um subespaço
seu, também é noetheriano.
Em outros termos, podemos supor que $Y$ é afim e que o espaço subjacente a
$X$ é noetheriano;
existe então uma cobertura de $X$ por um número finito de abertos afins
$D(g_i)\subset Y$, onde $g_i\in A(Y)$, tais que as interseções
$X\cap D(g_i)$ são fechadas em $D(g_i)$ (e, portanto, esquemas afins
\sref{I.4.2.3}), pois $X$ é localmente fechado em $Y$ \sref{I.4.1.3}.
Então $A(X\cap D(g_i))$ é uma álgebra de tipo finito sobre $A(D(g_i))$, por
\sref{I.6.3.4}[i] e \sref{I.6.3.3}, e
$A(D(g_i))=A(Y)_{g_i}=A(Y)[1/g_i]$ é de tipo finito sobre $A(Y)$, o que
conclui a demonstração.
\end{proof}

\begin{corollary}[6.3.6]
\label{I.6.3.6}
Sejam $f:X\to Y$ e $g:Y\to Z$ dois morfismos.
Se $g\circ f$ é de tipo finito e $X$ é noetheriano, ou se
$X\times_Z Y$ é localmente noetheriano, então $f$ é de tipo finito.
\end{corollary}

\begin{proof}
Isto resulta imediatamente da demonstração de \sref{I.5.5.12} e de
\sref{I.6.3.5}, aplicada ao morfismo de imersão $\Gamma_f$.
\end{proof}

\begin{proposition}[6.3.7]
\label{I.6.3.7}
Seja $f:X\to Y$ um morfismo de tipo finito;
se $Y$ é noetheriano (resp. localmente noetheriano), então $X$ é
noetheriano (resp. localmente noetheriano).
\end{proposition}

\begin{proof}
Podemos limitar-nos a demonstrar a proposição quando $Y$ é noetheriano.
Então $Y$ é reunião finita de abertos afins $V_i$ tais que os anéis
$A(V_i)$ são noetherianos.
Em virtude de \sref{I.6.3.2}, cada $f^{-1}(V_i)$ é reunião de um número
finito de abertos afins $W_{ij}$ tais que os $A(W_{ij})$ são álgebras de
tipo finito sobre $A(V_i)$ e, portanto, anéis noetherianos;
isto demonstra que $X$ é noetheriano.
\end{proof}

\begin{corollary}[6.3.8]
\label{I.6.3.8}
Seja $X$ um pré-esquema de tipo finito sobre $S$.
Para toda extensão da base $S'\to S$ tal que $S'$ seja noetheriano
(resp. localmente noetheriano), $X_{(S')}$ é noetheriano
(resp. localmente noetheriano).
\end{corollary}

\begin{proof}
Isto resulta de \sref{I.6.3.7}, pois $X_{(S')}$ é de tipo finito sobre
$S'$ em virtude de \sref{I.6.3.4}[iv].
\end{proof}

Pode-se ainda dizer que, em um produto $X\times_S Y$ de $S$-pré-esquemas,
se \emph{um} dos
\oldpage[I]{147}
fatores $X$, $Y$ é \emph{de tipo finito} sobre $S$ e
\emph{o outro é noetheriano} (resp. \emph{localmente noetheriano}), então
$X\times_S Y$ é \emph{noetheriano} (resp. \emph{localmente noetheriano}).

\begin{corollary}[6.3.9]
\label{I.6.3.9}
Seja $X$ um pré-esquema de tipo finito sobre um pré-esquema localmente
noetheriano $S$.
Então todo $S$-morfismo $f:X\to Y$ é de tipo finito.
\end{corollary}

\begin{proof}
Com efeito, podemos supor $S$ noetheriano;
se $\vphi:X\to S$ e $\psi:Y\to S$ são os morfismos estruturais, temos
$\vphi=\psi\circ f$, e $X$ é noetheriano em virtude de \sref{I.6.3.7};
portanto, $f$ é de tipo finito em virtude de \sref{I.6.3.6}.
\end{proof}

\begin{proposition}[6.3.10]
\label{I.6.3.10}
Seja $f:X\to Y$ um morfismo de tipo finito.
Para que $f$ seja sobrejetivo, é necessário e suficiente que, para todo corpo
\emph{algebricamente fechado} $\Omega$, a aplicação
$X(\Omega)\to Y(\Omega)$ correspondente a $f$ \sref{I.3.4.1} seja
sobrejetiva.
\end{proposition}

\begin{proof}
A condição é suficiente, como se vê considerando, para todo $y\in Y$, uma
extensão algebricamente fechada $\Omega$ de $\kres(y)$ e o diagrama comutativo
\[
  \xymatrix{
    X\ar[dd]_f\\
    & \Spec(\Omega)\ar[ul]\ar[dl]\\
    Y
  }
\]
(cf. \sref{I.3.5.3}).
Reciprocamente, suponhamos que $f$ seja sobrejetivo e seja
$g:\{\xi\}=\Spec(\Omega)\to Y$ um morfismo, sendo $\Omega$ um corpo
algebricamente fechado.
Se considerarmos o diagrama
\[
  \xymatrix{
    X\ar[d]_f &
    X_{(\Omega)}\ar[l]\ar[d]^{f_{(\Omega)}}\\
    Y &
    \Spec(\Omega)\ar[l]
  }
\]
basta então mostrar que existe um \emph{ponto racional sobre $\Omega$} em
$X_{(\Omega)}$ (\sref{I.3.3.14}, \sref{I.3.4.3} e
\sref{I.3.4.4}).
Como $f$ é sobrejetivo, $X_{(\Omega)}$ não é vazio \sref{I.3.5.10}; e, como
$f$ é de tipo finito, $f_{(\Omega)}$ também é de tipo finito
\sref{I.6.3.4}[iv].
Por conseguinte, $X_{(\Omega)}$ contém um aberto afim não vazio $Z$ tal que
$A(Z)$ é uma $\Omega$-álgebra não nula de tipo finito.
Em virtude do teorema dos zeros de Hilbert~\cite{I-21}, existe um
$\Omega$-homomorfismo $A(Z)\to\Omega$ e, portanto, uma seção de
$X_{(\Omega)}$ sobre $\Spec(\Omega)$, o que demonstra a proposição.
\end{proof}

\subsection{Pré-esquemas algébricos}
\label{subsection:I.6.4}
\label{I.6.4}

\begin{definition}[6.4.1]
\label{I.6.4.1}
Dado um corpo $K$, chama-se \emph{$K$-pré-esquema algébrico} um pré-esquema
$X$ de tipo finito sobre $K$; $K$ é chamado o corpo de base de $X$.
Se, além disso, $X$ é um esquema \emph{(ou se $X$ é um
\emph{$K$-esquema}, o que é equivalente \sref{I.5.5.8})}, diz-se também que
$X$ é um \emph{$K$-esquema algébrico}.
\end{definition}

Todo $K$-pré-esquema algébrico é \emph{noetheriano} \sref{I.6.3.7}.

\begin{proposition}[6.4.2]
\label{I.6.4.2}
Seja $X$ um $K$-pré-esquema algébrico.
Para que um ponto $x\in X$ seja fechado, é necessário e suficiente que
$\kres(x)$ seja uma extensão algébrica de $K$ de grau finito.
\end{proposition}

\begin{proof}
Podemos supor que $X$ seja afim, sendo seu anel $A$ uma $K$-álgebra de tipo
finito.
Com efeito, os abertos afins $U$ de $X$ tais que $A(U)$ seja uma
$K$-álgebra de tipo finito formam um recobrimento de $X$
\sref{I.6.3.1}.
Os pontos fechados de $X$ são então os pontos tais que $\mathfrak{j}_x$ seja
\oldpage[I]{148}
um ideal maximal de $A$, isto é, tais que $A/\mathfrak{j}_x$ seja um corpo
(necessariamente igual a $\kres(x)$).
Como $A/\mathfrak{j}_x$ é uma $K$-álgebra de tipo finito, vê-se que, se $x$
é fechado, $\kres(x)$ é um corpo que é uma álgebra de tipo finito sobre
$K$ e, portanto, necessariamente uma $K$-álgebra de \emph{posto
finito}~\cite{I-21}.
Reciprocamente, se $\kres(x)$ tem posto finito sobre $K$, o mesmo vale para
$A/\mathfrak{j}_x\subset\kres(x)$; e, como todo domínio integral que é uma
$K$-álgebra de posto finito é um corpo, tem-se
$A/\mathfrak{j}_x=\kres(x)$, logo $x$ é fechado.
\end{proof}

\begin{corollary}[6.4.3]
\label{I.6.4.3}
Sejam $K$ um corpo algebricamente fechado e $X$ um $K$-pré-esquema algébrico;
os pontos fechados de $X$ são então os pontos racionais sobre $K$
\sref{I.3.4.4} e se identificam canonicamente com os pontos de $X$ com
valores em $K$.
\end{corollary}

\begin{proposition}[6.4.4]
\label{I.6.4.4}
Seja $X$ um pré-esquema algébrico sobre um corpo $K$.
As propriedades seguintes são equivalentes:
\begin{enumerate}
  \item[{\rm(a)}] $X$ é artiniano.
  \item[{\rm(b)}] O espaço subjacente a $X$ é discreto.
  \item[{\rm(c)}] O espaço subjacente a $X$ tem apenas um número finito de pontos fechados.
  \item[{\rm(c')}] O espaço subjacente a $X$ é finito.
  \item[{\rm(d)}] Os pontos de $X$ são fechados.
  \item[{\rm(e)}] $X$ é isomorfo a $\Spec(A)$, onde $A$ é uma $K$-álgebra de posto finito.
\end{enumerate}
\end{proposition}

\begin{proof}
Como $X$ é noetheriano, resulta de \sref{I.6.2.2} que as condições (a), (b) e (d) são equivalentes e implicam (c) e (c$'$);
por outro lado, é claro que (e) implica (a).
Resta ver que (c) implica (d) e (e);
podemos nos limitar ao caso em que $X$ é afim.
Então $A(X)$ é uma $K$-álgebra de tipo finito \sref{I.6.3.3} e, portanto, um anel de Jacobson (\cite[p.~3-11 e 3-12]{I-1}), no qual há, por hipótese, apenas um número finito de ideais maximais.
Como uma interseção finita de ideais primos só pode ser um ideal primo se for igual a um deles, todo ideal primo de $A(X)$ é, pois, maximal, donde (d).
Além disso, sabe-se então por \sref{I.6.2.2} que $A(X)$ é uma $K$-álgebra artiniana de tipo finito e, portanto, necessariamente de \emph{posto finito}~\cite{I-21}.
\end{proof}

\begin{env}[6.4.5]
\label{I.6.4.5}
Quando as condições de \sref{I.6.4.4} são satisfeitas, diz-se que $X$ é um esquema \emph{finito sobre $K$} (cf. \sref[II]{II.6.1.1}), ou um \emph{$K$-esquema finito}, de \emph{posto} $[A:K]$, que também se denota por $\rg_K(X)$;
se $X$ e $Y$ são $K$-esquemas finitos, tem-se
\[
  \rg_K(X\sqcup Y)=\rg_K(X)+\rg_K(Y),
  \tag{6.4.5.1}
\]
\[
  \rg_K(X\times_K Y)=\rg_K(X)\rg_K(Y),
  \tag{6.4.5.2}
\]
como resulta de \sref{I.3.2.2}.
\end{env}

\begin{corollary}[6.4.6]
\label{I.6.4.6}
Seja $X$ um $K$-esquema finito.
Para toda extensão $K'$ de $K$, $X\otimes_K K'$ é um $K'$-esquema finito, e seu posto sobre $K'$ é igual ao posto de $X$ sobre $K$.
\end{corollary}

\begin{proof}
Se $A=A(X)$, tem-se $[A\otimes_K K':K']=[A:K]$.
\end{proof}

\begin{corollary}[6.4.7]
\label{I.6.4.7}
Seja $X$ um esquema finito sobre um corpo $K$;
ponhamos $n=\sum_{x\in X}[\kres(x):K]_s$ \emph{(recordemos que, se $K'$ é uma extensão de $K$, então $[K':K]_s$ é o \emph{posto separável} de $K'$ sobre $K$, posto da maior extensão algébrica separável de $K$ contida em $K'$)};
\oldpage[I]{149}
então,
para toda extensão algebricamente fechada $\Omega$ de $K$, o espaço subjacente a $X\otimes_K\Omega$ tem exatamente $n$ pontos, que se identificam com os pontos de $X$ com valores em $\Omega$.
\end{corollary}

\begin{proof}
Podemos evidentemente limitar-nos ao caso em que o anel $A=A(X)$ é \emph{local} \sref{I.6.2.2};
sejam $\mathfrak{m}$ seu ideal maximal e $L=A/\mathfrak{m}$ seu corpo residual, extensão algébrica de $K$.
Os pontos de $X$ com valores em $\Omega$ correspondem então biunivocamente às $\Omega$-seções de $X\otimes_K\Omega$ (\sref{I.3.4.1} e \sref{I.3.3.14}), e também aos $K$-homomorfismos de $L$ em $\Omega$ (\sref{I.1.7.3}), de onde resulta a proposição (Bourbaki, \emph{Alg.}, cap.~V, \textsection7, n\textsuperscript{o}~5, prop.~8), levando em conta \sref{I.6.4.3}.
\end{proof}

\begin{env}[6.4.8]
\label{I.6.4.8}
O número $n$ definido em \sref{I.6.4.7} chama-se o \emph{posto separável} de $A$ (ou de $X$) sobre $K$, ou também o \emph{número geométrico de pontos} de $X$;
é, portanto, igual ao número de elementos de $X(\Omega)_K$.
Resulta imediatamente dessa definição que, para toda extensão $K'$ de $K$, $X\otimes_K K'$ tem o mesmo número geométrico de pontos que $X$.
Se denotarmos esse número por $n(X)$, é claro que, se $X$ e $Y$ são dois esquemas finitos sobre $K$, tem-se
\[
  n(X\sqcup Y)=n(X)+n(Y).
  \tag{6.4.8.1}
\]
Sob as mesmas hipóteses, tem-se também
\[
  n(X\times_K Y)=n(X)n(Y)
  \tag{6.4.8.2}
\]
como resulta imediatamente da interpretação de $n(X)$ como número de elementos de $X(\Omega)_K$ e da fórmula~\hyperref[I.3.4.3.1]{(3.4.3.1)}.
\end{env}

\begin{proposition}[6.4.9]
\label{I.6.4.9}
Sejam $K$ um corpo, $X$ e $Y$ dois $K$-pré-esquemas algébricos, $f:X\to Y$ um $K$-morfismo, e $\Omega$ uma extensão algebricamente fechada de $K$, de grau de transcendência infinito sobre $K$.
Para que $f$ seja sobrejetivo, é necessário e suficiente que a aplicação $X(\Omega)_K\to Y(\Omega)_K$ correspondente a $f$ \sref{I.3.4.1} seja sobrejetiva.
\end{proposition}

\begin{proof}
A necessidade resulta de \sref{I.6.3.10}, observando que $f$ é necessariamente de tipo finito \sref{I.6.3.9}.
Para ver que a condição é suficiente, raciocina-se como em \sref{I.6.3.10}, observando que, para todo $y\in Y$, $\kres(y)$ é uma extensão de $K$ de tipo finito e, portanto, é $K$-isomorfo a um subcorpo de $\Omega$.
\end{proof}

\begin{remark}[6.4.10]
\label{I.6.4.10}
Veremos no capítulo~IV que a conclusão de \sref{I.6.4.9} permanece válida sem hipótese alguma relativa ao grau de transcendência de $\Omega$ sobre $K$.
\end{remark}

\begin{proposition}[6.4.11]
\label{I.6.4.11}
Se $f:X\to Y$ é um morfismo de tipo finito, então, para todo $y\in Y$, a fibra $f^{-1}(y)$ é um pré-esquema algébrico sobre o corpo residual $\kres(y)$ e, para todo $x\in f^{-1}(y)$, $\kres(x)$ é uma extensão de $\kres(y)$ de tipo finito.
\end{proposition}

\begin{proof}
Como $f^{-1}(y)=X\otimes_Y\kres(y)$ \sref{I.3.6.3}, a proposição resulta de \sref{I.6.3.4}[iv] e de \sref{I.6.3.3}.
\end{proof}

\begin{proposition}[6.4.12]
\label{I.6.4.12}
Sejam $f:X\to Y$ e $g:Y'\to Y$ dois morfismos;
ponhamos $X'=X\times_Y Y'$ e seja $f'=f_{(Y')}:X'\to Y'$.
Seja $y'\in Y'$ e ponhamos $y=g(y')$;
se a fibra $f^{-1}(y)$ é um esquema algébrico finito sobre $\kres(y)$, então a fibra $f'^{-1}(y')$ é um esquema algébrico finito sobre $\kres(y')$, com o mesmo posto e o mesmo número geométrico de pontos que $f^{-1}(y)$.
\end{proposition}

\begin{proof}
Levando em conta a transitividade das fibras \sref{I.3.6.5}, isso resulta imediatamente de \sref{I.6.4.6} e de \sref{I.6.4.8}.
\end{proof}

\begin{env}[6.4.13]
\label{I.6.4.13}
A Proposição~\sref{I.6.4.11} mostra que os morfismos de tipo finito correspondem intuitivamente a ``famílias algébricas de variedades algébricas'', desempenhando os pontos de $Y$
\oldpage[I]{150}
o papel de ``parâmetros'', o que confere a esses morfismos uma significação ``geométrica''.
Os morfismos que não são de tipo finito intervirão mais adiante principalmente nas questões de ``mudança do pré-esquema de base'', por exemplo por localização ou completamento.
\end{env}

\subsection{Determinação local de um morfismo}
\label{subsection:I.6.5}
\label{I.6.5}

\begin{proposition}[6.5.1]
\label{I.6.5.1}
Sejam $X$ e $Y$ dois $S$-pré-esquemas, sendo $Y$ de tipo finito sobre $S$;
sejam $x\in X$ e $y\in Y$ situados sobre um mesmo ponto $s\in S$.
\begin{enumerate}
  \item[{\rm(i)}] Se dois $S$-morfismos $f=(\psi,\theta)$ e $f'=(\psi',\theta')$ de $X$ em $Y$ são tais que $\psi(x)=\psi'(x)=y$, e os $\sh{O}_s$-homomorfismos (locais) $\theta_x^\sharp$ e ${\theta'}_x^\sharp$ de $\sh{O}_y$ em $\sh{O}_x$ são idênticos, então $f$ e $f'$ coincidem em uma vizinhança aberta de $x$.
  \item[{\rm(ii)}] Suponhamos, além disso, que $S$ seja localmente noetheriano.
    Para todo $\sh{O}_s$-homomorfismo local $\vphi:\sh{O}_y\to\sh{O}_x$, existem uma vizinhança aberta $U$ de $x$ em $X$ e um $S$-morfismo $f=(\psi,\theta)$ de $U$ em $Y$ tais que $\psi(x)=y$ e $\theta_x^\sharp=\vphi$.
\end{enumerate}
\end{proposition}

\begin{proof}
\medskip\noindent
\begin{enumerate}
  \item[(i)] Como a questão é local sobre $S$, $X$ e $Y$, podemos supor que $S$, $X$ e $Y$ são afins, com anéis respectivos $A$, $B$ e $C$, e que $f$ e $f'$ são respectivamente da forma $({}^a\vphi,\widetilde{\vphi})$ e $({}^a\vphi',\widetilde{\vphi}')$, onde $\vphi$ e $\vphi'$ são dois $A$-homomorfismos de $C$ em $B$ tais que $\vphi^{-1}(\mathfrak{j}_x)=\vphi'^{-1}(\mathfrak{j}_x)=\mathfrak{j}_y$, e os homomorfismos $\vphi_x$ e $\vphi'_x$ de $C_y$ em $B_x$, deduzidos de $\vphi$ e $\vphi'$, são idênticos;
    podemos supor, além disso, que $C$ é uma $A$-álgebra \emph{de tipo finito}.
    Sejam $c_i$ ($1\leq i\leq n$) geradores da $A$-álgebra $C$, e ponhamos $b_i=\vphi(c_i)$ e $b'_i=\vphi'(c_i)$;
    por hipótese, tem-se $b_i/1=b'_i/1$ no anel de frações $B_x$ ($1\leq i\leq n$).
    Isso significa que existem elementos $s_i\in B\setmin\mathfrak{j}_x$ tais que $s_i(b_i-b'_i)=0$ para $1\leq i\leq n$, e podemos evidentemente supor que todos os $s_i$ são iguais a um mesmo elemento $g\in B\setmin\mathfrak{j}_x$.
    Conclui-se daí que $b_i/1=b'_i/1$ para $1\leq i\leq n$ no anel de frações $B_g$;
    se $i_g$ é o homomorfismo canônico $B\to B_g$, tem-se então $i_g\circ\vphi=i_g\circ\vphi'$;
    portanto, as restrições de $f$ e $f'$ a $D(g)$ são idênticas.
  \item[(ii)] Podemos reduzir-nos à mesma situação que em (i), e supor, além disso, que o anel $A$ é noetheriano.
    Sejam $c_i$ ($1\leq i\leq n$) geradores da $A$-álgebra $C$, e seja $\alpha:A[X_1,\ldots,X_n]\to C$ o homomorfismo da álgebra de polinômios $A[X_1,\ldots,X_n]$ em $C$ que envia $X_i$ em $c_i$ para $1\leq i\leq n$.
    Seja também $i_y$ o homomorfismo canônico $C\to C_y$, e consideremos o homomorfismo composto
    \[
      \beta:A[X_1,\ldots,X_n] \xrightarrow{\alpha} C \xrightarrow{i_y} C_y \xrightarrow{\vphi} B_x.
    \]
    Denotemos por $\mathfrak{a}$ o núcleo de $\beta$;
    como $A$ é noetheriano, $A[X_1,\ldots,X_n]$ também o é e, portanto, $\mathfrak{a}$ admite um sistema finito de geradores $Q_j(X_1,\ldots,X_n)$ ($1\leq j\leq m$).
    Por outro lado, cada um dos elementos $\vphi(i_y(c_i))$ pode ser escrito na forma $b_i/s_i$, onde $b_i\in B$ e $s_i\not\in\mathfrak{j}_x$;
    podemos supor, além disso, que todos os $s_i$ são iguais a um mesmo elemento $g\in B\setmin\mathfrak{j}_x$.
    Nessas condições, tem-se por hipótese $Q_j(b_1/g,\ldots,b_n/g)=0$ em $B_x$; ponhamos
    \[
      Q_j(X_1/T,\ldots,X_n/T)=P_j(X_1,\ldots,X_n,T)/T^{k_j}
    \]
    onde $P_j$ é homogêneo de grau $k_j$.
    Seja então $d_j=P_j(b_1,\ldots,b_n,g)\in B$.
    Por hipótese, tem-se $t_jd_j=0$ para algum $t_j\in B\setmin\mathfrak{j}_x$ ($1\leq j\leq m$), e podemos evidentemente supor que todos os $t_j$
\oldpage[I]{151}
    são iguais a um mesmo elemento $h\in B\setmin\mathfrak{j}_x$;
    conclui-se daí que $P_j(hb_1,\ldots,hb_n,hg)=0$ para $1\leq j\leq m$.
    Nessas condições, consideremos o homomorfismo $\rho$ de $A[X_1,\ldots,X_n]$ no anel de frações $B_{hg}$ que envia $X_i$ em $hb_i/hg$ ($1\leq i\leq n$);
    a imagem de $\mathfrak{a}$ por esse homomorfismo é $0$, e \emph{a fortiori} também o é a imagem do núcleo $\alpha^{-1}(0)$ por $\rho$.
    Portanto, $\rho$ se fatora como $A[X_1,\ldots,X_n]\xrightarrow{\alpha}C\xrightarrow{\gamma}B_{hg}$, com $\gamma(c_i)=hb_i/hg$, e é claro que, se $i_x$ é o homomorfismo canônico $B_{hg}\to B_x$, então o diagrama
    \[
    \label{I.6.5.1.1}
      \xymatrix{
        C\ar[r]^\gamma\ar[d]_{i_y} &
        B_{hg}\ar[d]^{i_x}\\
        C_y\ar[r]^\vphi &
        B_x
      }
      \tag{6.5.1.1}
    \]
    é comutativo; tem-se, pois, $\vphi=\gamma_x$, e, como $\vphi$ é um homomorfismo local, $^a\gamma(x)=y$;
    $f=({}^a\gamma,\widetilde{\gamma})$ é, portanto, um $S$-morfismo da vizinhança $D(hg)$ de $x$ para $Y$, como se afirmava na proposição.
\end{enumerate}
\end{proof}

\begin{corollary}[6.5.2]
\label{I.6.5.2}
Sob as hipóteses de \sref{I.6.5.1}[ii]\footnote{\emph{[Trad.] A fonte francesa remete aqui a (6.5.1, (iii)), item inexistente; a hipótese pertinente é (6.5.1, (ii)).}}, se além disso $X$ é de tipo finito sobre $S$, pode-se supor que o morfismo $f$ é de tipo finito.
\end{corollary}

\begin{proof}
Isto resulta de \sref{I.6.3.6}.
\end{proof}

\begin{corollary}[6.5.3]
\label{I.6.5.3}
Suponhamos satisfeitas as hipóteses de \sref{I.6.5.1}[ii], e suponhamos, além disso, que $Y$ seja integral e que $\vphi$ seja um homomorfismo injetivo.
Então se pode supor que $f=({}^a\gamma,\widetilde{\gamma})$, onde $\gamma$ é injetivo.
\end{corollary}

\begin{proof}
Com efeito, podemos supor $C$ integral \sref{I.5.1.4}, e, portanto, $i_y$ injetivo; deduz-se então do diagrama~\sref{I.6.5.1.1} que $\gamma$ é injetivo.
\end{proof}

\begin{proposition}[6.5.4]
\label{I.6.5.4}
Seja $f=(\psi,\theta):X\to Y$ um morfismo de tipo finito, $x$ um ponto de $X$ e $y=\psi(x)$.
\begin{enumerate}
  \item[{\rm(i)}] Para que $f$ seja uma imersão local no ponto $x$ \sref{I.4.5.1}, é necessário e suficiente que $\theta_x^\sharp:\sh{O}_y\to\sh{O}_x$ seja sobrejetivo.
  \item[{\rm(ii)}] Suponhamos, além disso, que $Y$ seja localmente noetheriano.
    Para que $f$ seja um isomorfismo local no ponto $x$ \sref{I.4.5.2}, é necessário e suficiente que $\theta_x^\sharp$ seja um isomorfismo.
\end{enumerate}
\end{proposition}

\begin{proof}
\medskip\noindent
\begin{enumerate}
  \item[(ii)] Em virtude de \sref{I.6.5.1}, existe uma vizinhança aberta $V$ de $y$ e um morfismo $g:V\to X$ tais que $g\circ f$ (resp. $f\circ g$) esteja definido e coincida com a identidade em uma vizinhança de $x$ (resp. de $y$), de onde se deduz facilmente que $f$ é um isomorfismo local.
  \item[(i)] Como a questão é local sobre $X$ e $Y$, podemos supor que $X$ e $Y$ são afins, com anéis respectivos $A$ e $B$;
    tem-se $f=({}^a\vphi,\widetilde{\vphi})$, onde $\vphi$ é um homomorfismo de anéis $B\to A$ que faz de $A$ uma $B$-álgebra de tipo finito;
    tem-se $\vphi^{-1}(\mathfrak{j}_x)=\mathfrak{j}_y$, e o homomorfismo $\vphi_x:B_y\to A_x$ induzido por $\vphi$ é \emph{sobrejetivo}.
    Seja $(t_i)$ ($1\leq i\leq n$) um sistema de geradores da $B$-álgebra $A$;
    a hipótese sobre $\vphi_x$ implica que existem $b_i\in B$ e um $c\in B\setmin\mathfrak{j}_y$ tais que, no anel de frações $A_x$, se tem $t_i/1=\vphi(b_i)/\vphi(c)$ para $1\leq i\leq n$.
    Por conseguinte \sref{I.1.3.3}, existe um $a\in A\setmin\mathfrak{j}_x$ tal que, pondo-se $g=a\vphi(c)$, se tem também $t_i/1=a\vphi(b_i)/g$ \emph{no anel de frações $A_g$}.
    Isso posto, existe, por hipótese, um polinômio $Q(X_1,\ldots,X_n)$, com coeficientes no anel $\vphi(B)$,
\oldpage[I]{152}
    tal que $a=Q(t_1,\ldots,t_n)$;
    ponhamos $Q(X_1/T,\ldots,X_n/T)=P(X_1,\ldots,X_n,T)/T^m$, onde $P$ é homogêneo de grau $m$.
    No anel $A_g$, tem-se
    \[
      a/1=a^mP(\vphi(b_1),\ldots,\vphi(b_n),\vphi(c))/g^m=a^m\vphi(d)/g^m
    \]
    onde $d\in B$.
    Como, em $A_g$, $g/1=(a/1)(\vphi(c)/1)$ é inversível por definição, também o são $a/1$ e $\vphi(c)/1$, e pode-se, portanto, escrever $a/1=(\vphi(d)/1)(\vphi(c)/1)^{-m}$.
    Conclui-se daí que $\vphi(d)/1$ também é inversível em $A_g$.
    Ponhamos então $h=cd$;
    como $\vphi(h)/1$ é inversível em $A_g$, o homomorfismo composto $B\xrightarrow{\vphi}A\to A_g$ se fatora como $B\to B_h\xrightarrow{\gamma}A_g$ \sref[0]{0.1.2.4}.
    Mostremos que $\gamma$ é \emph{sobrejetivo};
    basta verificar que a imagem de $B_h$ em $A_g$ contém os $t_i/1$ e $(g/1)^{-1}$.
    Ora, tem-se $(g/1)^{-1}=(\vphi(c)/1)^{m-1}(\vphi(d)/1)^{-1}=\gamma(c^m/h)$, e $a/1=\gamma(d^{m+1}/h^m)$, portanto $(a\vphi(b_i))/1=\gamma(b_id^{m+1}/h^m)$, e, como $t_i/1=(a\vphi(b_i)/1)(g/1)^{-1}$, fica demonstrada nossa afirmação.
    A escolha de $h$ implica que $\psi(D(g))\subset D(h)$, e sabe-se também que a restrição de $f$ a $D(g)$ é igual a $({}^a\gamma,\widetilde{\gamma})$;
    como $\gamma$ é sobrejetivo, essa restrição é uma imersão fechada de $D(g)$ em $D(h)$ \sref{I.4.2.3}.
\end{enumerate}
\end{proof}

\begin{corollary}[6.5.5]
\label{I.6.5.5}
Seja $f=(\psi,\theta):X\to Y$ um morfismo de tipo finito.
Suponhamos que $X$ seja irredutível, denotemos por $x$ seu ponto genérico e ponhamos $y=\psi(x)$.
\begin{enumerate}
  \item[{\rm(i)}] Para que $f$ seja uma imersão local em algum ponto de $X$, é necessário e suficiente que $\theta_x^\sharp:\sh{O}_y\to\sh{O}_x$ seja sobrejetivo.
  \item[{\rm(ii)}] Suponhamos, além disso, que $Y$ seja irredutível e localmente noetheriano.
    Para que $f$ seja um isomorfismo local em algum ponto de $X$, é necessário e suficiente que $y$ seja o ponto genérico de $Y$ \emph{(ou, equivalentemente \sref[0]{0.2.1.4}, que $f$ seja um morfismo \emph{dominante})} e que $\theta_x^\sharp$ seja um isomorfismo \emph{(em outras palavras, que $f$ seja \emph{birracional} \sref{I.2.2.9})}.
\end{enumerate}
\end{corollary}

\begin{proof}
É claro que (i) resulta de \sref{I.6.5.4}[i], tendo em conta que todo aberto não vazio de $X$ contém $x$;
analogamente, (ii) resulta de \sref{I.6.5.4}[ii].
\end{proof}

\subsection{Morfismos quase-compactos e morfismos localmente de tipo finito}
\label{subsection:I.6.6}

\begin{definition}[6.6.1]
\label{I.6.6.1}
Diz-se que um morfismo $f:X\to Y$ é \emph{quase-compacto} se a imagem inversa de todo aberto quase-compacto de $Y$ por $f$ for quase-compacta.
\end{definition}

Seja $\mathfrak{B}$ uma base da topologia de $Y$ formada por abertos quase-compactos (por exemplo, abertos afins);
para que $f$ seja quase-compacto, é necessário e suficiente que a imagem inversa de todo elemento de $\mathfrak{B}$ por $f$ seja quase-compacta (ou, equivalentemente, uma união \emph{finita} de abertos afins), pois todo aberto quase-compacto de $Y$ é uma união finita de elementos de $\mathfrak{B}$.
Por exemplo, se $X$ é \emph{quase-compacto} e $Y$ \emph{afim}, então \emph{todo} morfismo $f: X\to Y$ é quase-compacto:
com efeito, $X$ é uma união finita de abertos afins $U_i$ e, para todo aberto afim $V$ de $Y$, $U_i\cap f^{-1}(V)$ é afim \sref{I.5.5.10}, portanto quase-compacto.

Se $f: X\to Y$ é um morfismo quase-compacto, é claro que, para todo aberto $V$ de $Y$, a restrição de $f$ a $f^{-1}(V)$ é um morfismo quase-compacto $f^{-1}(V)\to V$.
Reciprocamente, se $(U_\alpha)$ é um recobrimento aberto de $Y$, e $f: X\to Y$ um morfismo tal que as restrições $f^{-1}(U_\alpha)\to U_\alpha$ sejam quase-compactas, então $f$ é quase-compacto.

\begin{definition}[6.6.2]
\label{I.6.6.2}
Diz-se que um morfismo $f: X\to Y$ é \emph{localmente de tipo finito} se, para todo $x\in X$, existem uma vizinhança aberta $U$ de $x$ e uma vizinhança aberta $V\supset f(U)$ de $f(x)$\footnote{\emph{[Trad.] A fonte francesa imprime aqui «de $y$» sem definir $y$; a leitura corrigida explicita $y=f(x)$.}} tais que a
\oldpage[I]{153}
restrição de $f$ a $U$ seja um morfismo de tipo finito de $U$ em $V$.
Diz-se também então que $X$ é um pré-esquema localmente de tipo finito sobre $Y$, ou um $Y$-pré-esquema localmente de tipo finito.
\end{definition}

Resulta imediatamente de \sref{I.6.3.2} que, se $f$ é localmente de tipo finito, então, para todo aberto $W$ de $Y$, a restrição de $f$ a $f^{-1}(W)$ é um morfismo $f^{-1}(W)\to W$ localmente de tipo finito.

Se $Y$ é localmente noetheriano e $X$ localmente de tipo finito sobre $Y$, então $X$ é localmente noetheriano em virtude de \sref{I.6.3.7}.

\begin{proposition}[6.6.3]
\label{I.6.6.3}
Para que um morfismo $f: X\to Y$ seja de tipo finito, é necessário e suficiente que seja quase-compacto e localmente de tipo finito.
\end{proposition}

\begin{proof}
A necessidade das condições é imediata, levando em conta \sref{I.6.3.1} e a observação que segue \sref{I.6.6.1}.
Reciprocamente, suponhamos satisfeitas essas condições e seja $U$ um aberto afim de $Y$, de anel $A$;
para todo $x\in f^{-1}(U)$, existe por hipótese uma vizinhança $V(x)\subset f^{-1}(U)$ de $x$ e uma vizinhança $W(x)\subset U$ de $y=f(x)$ que contém $f(V(x))$ e tal que a restrição de $f$ a $V(x)$ seja um morfismo $V(x)\to W(x)$ de tipo finito.
Substituindo $W(x)$ por uma vizinhança $W_1(x)\subset W(x)$ de $y=f(x)$\footnote{\emph{[Trad.] A fonte francesa imprime aqui «de $x$». Como $W_1(x)\subset W(x)\subset U\subset Y$, deve tratar-se de uma vizinhança de $y=f(x)$.}} da forma $D(g)$ (com $g\in A$), e $V(x)$ por $V(x)\cap f^{-1}(W_1(x))$, podemos supor que $W(x)$ seja da forma $D(g)$ e, portanto, de tipo finito sobre $U$ (pois seu anel se escreve $A[1/g]$);
por conseguinte, $V(x)$ é de tipo finito sobre $U$.
Além disso, $f^{-1}(U)$ é quase-compacto por hipótese e, portanto, é uma união de um número finito de abertos $V(x_i)$, o que conclui a demonstração.
\end{proof}

\begin{proposition}[6.6.4]
\label{I.6.6.4}
\medskip\noindent
\begin{enumerate}
  \item[{\rm(i)}] Uma imersão $X\to Y$ é quase-compacta se for fechada, ou se o espaço subjacente a $Y$ for localmente noetheriano, ou se o espaço subjacente a $X$ for noetheriano.
  \item[{\rm(ii)}] O composto de dois morfismos quase-compactos é quase-compacto.
  \item[{\rm(iii)}] Se $f: X\to Y$ é um $S$-morfismo quase-compacto, então $f_{(S')}: X_{(S')}\to Y_{(S')}$ também o é para toda extensão $g: S'\to S$ do pré-esquema de base.
  \item[{\rm(iv)}] Se $f:X\to X'$ e $g:Y\to Y'$ são dois $S$-morfismos quase-compactos, então $f\times_S g$ é quase-compacto.
  \item[{\rm(v)}] Se o composto $g\circ f$ de dois morfismos $f: X\to Y$ e $g:Y\to Z$ é quase-compacto, e se $g$ é separado ou o espaço subjacente a $X$ é localmente noetheriano, então $f$ é quase-compacto.
  \item[{\rm(vi)}] Para que um morfismo $f$ seja quase-compacto, é necessário e suficiente que $f_\red$ o seja.
\end{enumerate}
\end{proposition}

\begin{proof}
Notemos que (vi) é evidente, pois a propriedade de ser quase-compacto para um morfismo depende apenas da aplicação contínua correspondente entre os espaços subjacentes.
Demonstremos igualmente a parte de (v) correspondente ao caso em que o espaço subjacente a $X$ é localmente noetheriano.
Ponhamos $h=g\circ f$, e seja $U$ um aberto quase-compacto de $Y$;
$g(U)$ é quase-compacto (mas não necessariamente aberto) em $Z$ e, portanto, está contido numa união finita de abertos quase-compactos $V_j$ \sref{I.2.1.3}, e $f^{-1}(U)$ está então contido na união dos $h^{-1}(V_j)$, que são subespaços quase-compactos de $X$ e, portanto, subespaços noetherianos.
Conclui-se então, por \sref[0]{0.2.2.3}, que $f^{-1}(U)$ é um espaço noetheriano e, \emph{a fortiori}, quase-compacto.

Para demonstrar as demais asserções, basta demonstrar (i), (ii) e (iii) \sref{I.5.5.12}.
Ora, (ii) é evidente, e (i) resulta de \sref{I.6.3.5} quando o espaço $Y$ é localmente noetheriano ou o espaço $X$ é noetheriano, e é evidente para uma imersão fechada.

Para estabelecer (iii),
\oldpage[I]{154}
podemos nos limitar ao caso em que $S=Y$ \sref{I.3.3.11};
ponhamos $f'=f_{(S')}$, e seja $U'$ um aberto quase-compacto de $S'$.
Para todo $s'\in U'$, seja $T$ uma vizinhança aberta afim de $g(s')$ em $S$, e seja $W$ uma vizinhança aberta afim de $s'$ contida em $U'\cap g^{-1}(T)$;
bastará demonstrar que $f'^{-1}(W)$ é quase-compacto;
em outros termos, podemos nos reduzir a provar que, quando $S$ e $S'$ são \emph{afins}, o espaço subjacente a $X\times_S S'$ é quase-compacto.
Mas, como $X$ é então, por hipótese, uma união finita de abertos afins $V_j$, $X\times_S S'$ é a união dos espaços subjacentes aos esquemas afins $V_j\times_S S'$ (\sref{I.3.2.2} e \sref{I.3.2.7}), o que demonstra a proposição.
\end{proof}

Observemos também que, se $X=X'\sqcup X''$ é a soma de dois pré-esquemas, um morfismo $f:X\to Y$ é quase-compacto se e somente se suas restrições a $X'$ e $X''$ o são.

\begin{proposition}[6.6.5]
\label{I.6.6.5}
Seja $f:X\to Y$ um morfismo quase-compacto.
Para que $f$ seja dominante, é necessário e suficiente que, para todo ponto genérico $y$ de uma componente irredutível de $Y$, $f^{-1}(y)$ contenha o ponto genérico de uma componente irredutível de $X$.
\end{proposition}

\begin{proof}
É imediato que a condição é suficiente (mesmo sem supor que $f$ seja quase-compacto).
Para ver que ela é necessária, consideremos uma vizinhança aberta afim $U$ de $y$;
$f^{-1}(U)$ é quase-compacto e, portanto, uma união \emph{finita} de abertos afins $V_i$, e a hipótese de que $f$ é dominante implica que $y$ pertence à aderência \emph{em $U$} de um dos $f(V_i)$.
Podemos evidentemente supor que $X$ e $Y$ são reduzidos;
como a aderência em $X$ de uma componente irredutível de $V_i$ é uma componente irredutível de $X$ \sref[0]{0.2.1.6}, podemos substituir $X$ por $V_i$ e $Y$ pelo subpré-esquema fechado reduzido de $U$ cujo espaço subjacente é $\overline{f(V_i)}\cap U$ \sref{I.5.2.1}, ficando assim reduzidos a demonstrar a proposição quando $X=\Spec(A)$ e $Y=\Spec(B)$ são afins e reduzidos.
Como $f$ é dominante, $B$ é então um subanel de $A$ \sref{I.1.2.7}, e a proposição resulta do fato de que todo ideal primo minimal de $B$ é a interseção com $B$ de um ideal primo minimal de $A$ \sref[0]{0.1.5.8}.
\end{proof}

\begin{proposition}[6.6.6]
\label{I.6.6.6}
\medskip\noindent
\begin{enumerate}
  \item[{\rm(i)}] Toda imersão local é localmente de tipo finito.
  \item[{\rm(ii)}] Se dois morfismos $f:X\to Y$ e $g:Y\to Z$ são localmente de tipo finito, então $g\circ f$ também o é.
  \item[{\rm(iii)}] Se $f:X\to Y$ é um $S$-morfismo localmente de tipo finito, então $f_{(S')}:X_{(S')}\to Y_{(S')}$ também o é para toda extensão $S'\to S$ do pré-esquema de base.
  \item[{\rm(iv)}] Se $f:X\to X'$ e $g:Y\to Y'$ são $S$-morfismos localmente de tipo finito, então $f\times_S g$ é localmente de tipo finito.
  \item[{\rm(v)}] Se o composto $g\circ f$ de dois morfismos é localmente de tipo finito, então $f$ é localmente de tipo finito.
  \item[{\rm(vi)}] Se um morfismo $f$ é localmente de tipo finito, então $f_\red$ também o é.
\end{enumerate}
\end{proposition}

\begin{proof}
Em virtude de \sref{I.5.5.12}, basta demonstrar (i), (ii) e (iii).
Se $j:X\to Y$ é uma imersão local, para todo $x\in X$ existem uma vizinhança aberta $V$ de $j(x)$ em $Y$ e uma vizinhança aberta $U$ de $x$ em $X$ tais que a restrição de $j$ a $U$ seja uma imersão fechada $U\to V$ \sref{I.4.5.1}; portanto, essa restrição é de tipo finito.
Para estabelecer (ii), consideremos um ponto $x\in X$;
por hipótese, existem uma vizinhança aberta $W$ de $g(f(x))$ e uma vizinhança aberta $V$ de $f(x)$ tais que $g(V)\subset W$ e $V$ seja de tipo finito sobre $W$;
além disso, $f^{-1}(V)$ é localmente de tipo finito sobre $V$ \sref{I.6.6.2}, de modo que existe uma vizinhança aberta $U$ de $x$ que
\oldpage[I]{155}
está contida em $f^{-1}(V)$ e é de tipo finito sobre $V$;
por conseguinte, $g(f(U))\subset W$ e $U$ é de tipo finito sobre $W$ \sref{I.6.3.4}[ii].
Finalmente, para demonstrar (iii), podemos nos limitar ao caso em que $Y=S$ \sref{I.3.3.11};
para todo $x'\in X'=X_{(S')}$, sejam $x$ a imagem de $x'$ em $X$, $s$ a imagem de $x$ em $S$, $T$ uma vizinhança aberta de $s$, $T'$ a imagem inversa de $T$ em $S'$ e $U$ uma vizinhança aberta de $x$ cuja imagem está contida em $T$ e que é de tipo finito sobre $T$;
então $U\times_S T'=U\times_T T'$ é uma vizinhança aberta de $x'$ \sref{I.3.2.7} que é de tipo finito sobre $T'$ \sref{I.6.3.4}[iv].
\end{proof}

\begin{corollary}[6.6.7]
\label{I.6.6.7}
Sejam $X$ e $Y$ dois $S$-pré-esquemas localmente de tipo finito sobre $S$.
Se $S$ é localmente noetheriano, então $X\times_S Y$ é localmente noetheriano.
\end{corollary}

\begin{proof}
Com efeito, como $S$ é localmente noetheriano e $X$ é localmente de tipo finito sobre $S$, $X$ é localmente noetheriano; além disso, $X\times_S Y$ é localmente de tipo finito sobre $X$ e, portanto, também é localmente noetheriano.
\end{proof}

\begin{remark}[6.6.8]
\label{I.6.6.8}
A Proposição~\sref{I.6.3.10} e sua demonstração estendem-se imediatamente ao caso em que se supõe apenas que o morfismo $f$ é localmente de tipo finito.
Do mesmo modo, as proposições~\sref{I.6.4.2} e \sref{I.6.4.9} permanecem válidas quando se supõe apenas que os pré-esquemas $X$ e $Y$ que figuram em seus enunciados são localmente de tipo finito sobre o corpo $K$.
\end{remark}

\section{Aplicações racionais}
\label{section:I.7}

\subsection{Aplicações racionais e funções racionais}
\label{subsection:I.7.1}

\begin{env}[7.1.1]
\label{I.7.1.1}
Sejam $X$ e $Y$ dois pré-esquemas, $U$ e $V$ dois abertos densos de $X$, e $f$ (resp. $g$) um morfismo de $U$ (resp. $V$) em $Y$; diz-se que $f$ e $g$ são \emph{equivalentes} se coincidem em um aberto denso de $U\cap V$.
Como uma interseção finita de abertos densos de $X$ é um aberto denso de $X$, é claro que essa relação é uma \emph{relação de equivalência}.
\end{env}

\begin{definition}[7.1.2]
\label{I.7.1.2}
Dados dois pré-esquemas $X$ e $Y$, chama-se aplicação racional de $X$ em $Y$ uma classe de equivalência de morfismos de abertos densos de $X$ em $Y$, para a relação de equivalência definida em \sref{I.7.1.1}.
Se $X$ e $Y$ são $S$-pré-esquemas, diz-se que tal classe é uma $S$-aplicação racional se existe um representante da classe que seja um $S$-morfismo.
Chama-se $S$-seção racional de $X$ toda $S$-aplicação racional de $S$ em $X$.
Chama-se função racional sobre um pré-esquema $X$ toda $X$-seção racional do $X$-pré-esquema $X\otimes_{\bb{Z}}\bb{Z}[T]$ (em que $T$ é uma indeterminada).

Por abuso de linguagem, quando se tratar apenas de $S$-pré-esquemas, diremos ``aplicação racional'' em vez de ``$S$-aplicação racional'' se disso não puder resultar nenhuma confusão.

Seja $f$ uma aplicação racional de $X$ em $Y$, e seja $U$ um aberto de $X$; se $f_1$ e $f_2$ são dois morfismos pertencentes à classe de $f$, definidos respectivamente em abertos densos $V$ e $W$ de $X$, as restrições $f_1|(U\cap V)$ e $f_2|(U\cap W)$ coincidem em $U\cap V\cap W$, que é denso em $U$; a classe de morfismos $f$ define, portanto, uma aplicação racional de $U$ em $Y$, chamada a \emph{restrição de $f$ a $U$}, e denotada por $f|U$.

Se a cada $S$-morfismo $f:X\to Y$ fazemos corresponder a $S$-aplicação racional
\oldpage[I]{156}
à qual pertence $f$, obtemos uma aplicação canônica de $\Hom_S(X,Y)$ no conjunto das $S$-aplicações racionais de $X$ em $Y$.
Denota-se por $\Gamma_\mathrm{rat}(X/Y)$ o conjunto das $Y$-seções racionais de $X$, e tem-se, portanto, uma aplicação canônica $\Gamma(X/Y)\to\Gamma_\mathrm{rat}(X/Y)$.
É claro, além disso, que, se $X$ e $Y$ são dois $S$-pré-esquemas, o conjunto das $S$-aplicações racionais de $X$ em $Y$ identifica-se canonicamente com $\Gamma_\mathrm{rat}((X\times_S Y)/X)$ \sref{I.3.3.14}.
\end{definition}

\begin{env}[7.1.3]
\label{I.7.1.3}
Resulta também de \sref{I.7.1.2} e \sref{I.3.3.14} que as funções racionais sobre $X$ se identificam canonicamente com as \emph{classes de equivalência de seções do feixe estrutural $\sh{O}_X$} sobre abertos densos de $X$, sendo equivalentes duas dessas seções se coincidem em algum aberto denso de $X$ contido na interseção de seus conjuntos de definição.
Em particular, resulta que as funções racionais sobre $X$ formam um \emph{anel} $R(X)$.
\end{env}

\begin{env}[7.1.4]
\label{I.7.1.4}
Quando $X$ é um pré-esquema \emph{irredutível}, todo aberto não vazio de $X$ é denso em $X$; pode-se dizer também que os abertos não vazios de $X$ são as \emph{vizinhanças abertas do ponto genérico} $x$ de $X$.
Dizer que dois morfismos de abertos não vazios de $X$ em $Y$ são equivalentes significa, nesse caso, que têm o \emph{mesmo germe} no ponto $x$.
Em outros termos, as aplicações racionais (resp. $S$-aplicações racionais) $X\to Y$ identificam-se com os \emph{germes de morfismos} (resp. \emph{$S$-morfismos}) de abertos não vazios de $X$ em $Y$ no ponto genérico $x$ de $X$.
Em particular:
\end{env}

\begin{proposition}[7.1.5]
\label{I.7.1.5}
Se $X$ é um pré-esquema irredutível, o anel $R(X)$ das funções racionais sobre $X$ identifica-se canonicamente com o anel local $\sh{O}_x$ do ponto genérico $x$ de $X$.
É um anel local de dimensão~0 e, portanto, um anel local artiniano quando $X$ é noetheriano; é um corpo quando $X$ é íntegro e, quando além disso $X$ é um esquema afim, identifica-se com o corpo de frações de $A(X)$.
\end{proposition}

\begin{proof}
Pelo que precede e pela identificação das funções racionais com as seções de $\sh{O}_X$ sobre um aberto denso de $X$, a primeira afirmação nada mais é do que a definição do talo de um feixe em um ponto.
Para as demais afirmações, podemos nos reduzir ao caso em que $X$ é afim, de anel $A$; então $\mathfrak{j}_x$ é o nilradical de $A$, e $\sh{O}_x$ é, portanto, de dimensão~0; se $A$ é íntegro, então $\mathfrak{j}_x=(0)$, e $\sh{O}_x$ é, pois, o corpo de frações de $A$.
Finalmente, se $A$ é noetheriano, sabe-se (\cite[p.~127, cor.~4]{I-11}) que $\mathfrak{j}_x$ é nilpotente e que $\sh{O}_x=A_{\mathfrak{j}_x}$ é artiniano.
\end{proof}

Se $X$ é \emph{íntegro}, o anel $\sh{O}_z$ é íntegro para todo $z\in X$; todo aberto afim $U$ que contém $z$ contém também $x$, e $R(U)$, sendo igual ao corpo de frações de $A(U)$, identifica-se com $R(X)$; conclui-se assim que $R(X)$ também pode ser identificado com o \emph{corpo de frações de $\sh{O}_z$}: a identificação canônica de $\sh{O}_z$ com um subanel de $R(X)$ consiste em associar a cada germe de uma seção $s\in\sh{O}_z$ a única função racional sobre $X$, classe de uma seção de $\sh{O}_X$ (necessariamente definida sobre um aberto denso de $X$) que tem $s$ como germe no ponto $z$.

\begin{env}[7.1.6]
\label{I.7.1.6}
Suponhamos agora que $X$ tenha um número \emph{finito} de componentes irredutíveis $X_i$ ($1\leq i\leq n$) (o que ocorre sempre que o espaço subjacente a $X$ seja \emph{noetheriano}); seja $X'_i$ o aberto de $X$ obtido suprimindo de $X_i$ a união dos $X_j\cap X_i$ para $j\neq i$; $X'_i$ é irredutível, seu ponto genérico $x_i$ é o ponto genérico de $X_i$, e os $X'_i$ são dois a dois disjuntos, sendo sua união densa em $X$ \sref[0]{0.2.1.6}.
Para
\oldpage[I]{157}
todo aberto denso $U$ de $X$, $U_i=U\cap X'_i$ é um aberto não vazio denso de $X'_i$, sendo os $U_i$ dois a dois disjuntos, e portanto $U'=\bigcup_i U_i$ é denso em $X$.
Dar um morfismo de $U'$ em $Y$ equivale a dar, arbitrariamente, um morfismo de cada um dos $U_i$ em $Y$.
\end{env}

Portanto:

\begin{proposition}[7.1.7]
\label{I.7.1.7}
Sejam $X$ e $Y$ dois pré-esquemas (resp. $S$-pré-esquemas) tais que $X$ tenha um número finito de componentes irredutíveis $X_i$, com pontos genéricos $x_i$ ($1\leq i\leq n$).
Se $R_i$ é o conjunto dos germes de morfismos (resp. $S$-morfismos) de abertos de $X$ em $Y$ no ponto $x_i$, então o conjunto das aplicações racionais (resp. $S$-aplicações racionais) de $X$ em $Y$ identifica-se com o produto dos $R_i$ ($1\leq i\leq n$).
\end{proposition}

\begin{corollary}[7.1.8]
\label{I.7.1.8}
Seja $X$ um pré-esquema noetheriano.
O anel das funções racionais sobre $X$ é um anel artiniano, cujas componentes locais são os anéis $\sh{O}_{x_i}$ dos pontos genéricos $x_i$ das componentes irredutíveis de $X$.
\end{corollary}

\begin{corollary}[7.1.9]
\label{I.7.1.9}
Sejam $A$ um anel noetheriano e $X=\Spec(A)$.
Se $Q$ é o complementar da união dos ideais primos minimais de $A$, então o anel das funções racionais sobre $X$ identifica-se canonicamente com o anel de frações $Q^{-1}A$.
\end{corollary}

Isso resultará do lema seguinte:

\begin{lemma}[7.1.9.1]
\label{I.7.1.9.1}
Para que, dado um elemento $f\in A$, $D(f)$ seja denso em $X$, é necessário e suficiente que $f\in Q$; todo aberto denso de $X$ contém um aberto da forma $D(f)$, com $f\in Q$.
\end{lemma}

Com efeito, suponhamos demonstrado o lema. Como o anel das seções $\Gamma(D(f),\sh{O}_X)$ se identifica com $A_f$ \sref{I.1.3.6} e \sref{I.1.3.7}, decorre do fato de que os $D(f)$ com $f\in Q$ formam uma família cofinal no conjunto ordenado (pela inclusão inversa) dos abertos densos de $X$, juntamente com a definição das funções racionais dada em \sref{I.7.1.1}, que o anel das funções racionais sobre $X$ se identifica com o limite indutivo dos $A_f$ para $f\in Q$ (para a relação de pré-ordem ``$g$ é múltiplo de $f$''), isto é, com $Q^{-1}A$ \sref[0]{0.1.4.5}.

\begin{proof}
Para demonstrar \hyperref[I.7.1.9.1]{(7.1.9.1)}, denotemos novamente por $X_i$ ($1\leq i\leq n$) as componentes irredutíveis de $X$; se $D(f)$ é denso em $X$, então $D(f)\cap X_i\neq\emp$ para $1\leq i\leq n$, e reciprocamente; ora, isso significa que $f\not\in\mathfrak{p}_i$ para $1\leq i\leq n$, pondo $\mathfrak{p}_i=\mathfrak{j}(X_i)$, e, como os $\mathfrak{p}_i$ são os ideais primos minimais de $A$ \hyperref[I.1.1.14]{(1.1.14)}, as condições $f\not\in\mathfrak{p}_i$ ($1\leq i\leq n$) equivalem a $f\in Q$, de onde resulta a primeira afirmação do lema.
Para a outra afirmação, se $U$ é um aberto denso de $X$, o complementar de $U$ é um conjunto da forma $V(\mathfrak{a})$, onde $\mathfrak{a}$ é um ideal que não está contido em nenhum dos $\mathfrak{p}_i$; portanto, tampouco está contido em sua união (\cite[p.~13]{I-10}), e existe assim algum $f\in\mathfrak{a}$ que pertence a $Q$; daí $D(f)\subset U$, o que conclui a demonstração.
\end{proof}

\begin{env}[7.1.10]
\label{I.7.1.10}
Suponhamos novamente que $X$ seja irredutível, com ponto genérico $x$.
Como todo aberto não vazio $U$ de $X$ contém $x$ e, por conseguinte, contém também todo $z\in X$ tal que $x\in\overline{\{z\}}$, todo morfismo $U\to Y$ pode ser composto com o morfismo canônico $\Spec(\sh{O}_x)\to X$ \sref{I.2.4.1}; e dois morfismos em $Y$, definidos em dois abertos não vazios de $X$, que coincidem em um aberto não vazio de $X$, dão, por composição, o mesmo morfismo $\Spec(\sh{O}_x)\to Y$.
Em outros termos, a toda aplicação racional de $X$ em $Y$ corresponde assim um morfismo bem determinado $\Spec(\sh{O}_x)\to Y$.
\end{env}

\oldpage[I]{158}
\begin{proposition}[7.1.11]
\label{I.7.1.11}
Sejam $X$ e $Y$ dois $S$-pré-esquemas; suponhamos que $X$ seja irredutível, com ponto genérico $x$, e que $Y$ seja de tipo finito sobre $S$.
Duas $S$-aplicações racionais quaisquer de $X$ em $Y$ às quais corresponde o mesmo $S$-morfismo $\Spec(\sh{O}_x)\to Y$ são idênticas.
Se, além disso, $S$ é localmente noetheriano, então a todo $S$-morfismo de $\Spec(\sh{O}_x)$ em $Y$ corresponde uma única $S$-aplicação racional de $X$ em $Y$.
\end{proposition}

\begin{proof}
Tendo em conta que todo aberto não vazio de $X$ é denso em $X$, isto resulta de \sref{I.6.5.1}.
\end{proof}

\begin{corollary}[7.1.12]
\label{I.7.1.12}
Suponhamos que $S$ seja localmente noetheriano e que se satisfaçam as demais hipóteses de \sref{I.7.1.11}.
As $S$-aplicações racionais de $X$ em $Y$ identificam-se então com os pontos do $S$-pré-esquema $Y$ com valores no $S$-pré-esquema $\Spec(\sh{O}_x)$.
\end{corollary}

\begin{proof}
Isto não é senão \sref{I.7.1.11}, com a terminologia introduzida em \sref{I.3.4.1}.
\end{proof}

\begin{corollary}[7.1.13]
\label{I.7.1.13}
Suponhamos que se satisfaçam as condições de \sref{I.7.1.12}.
Seja $s$ a imagem de $x$ em $S$.
Dar uma $S$-aplicação racional de $X$ em $Y$ equivale a dar um ponto $y$ de $Y$ situado sobre $s$ e um $\sh{O}_s$-homomorfismo local $\sh{O}_y\to\sh{O}_x=R(X)$.
\end{corollary}

\begin{proof}
Isto resulta de \sref{I.7.1.11} e de \sref{I.2.4.4}.
\end{proof}

Em particular:
\begin{corollary}[7.1.14]
\label{I.7.1.14}
Sob as condições de \sref{I.7.1.12}, as $S$-aplicações racionais de $X$ em $Y$ dependem somente (para $Y$ dado) do $S$-pré-esquema $\Spec(\sh{O}_x)$ e, em particular, permanecem as mesmas quando se substitui $X$ por $\Spec(\sh{O}_z)$, para todo $z\in X$.
\end{corollary}

\begin{proof}
Com efeito, como $z\in\overline{\{x\}}$, $x$ é o ponto genérico de $Z=\Spec(\sh{O}_z)$, e $\sh{O}_{X,x}=\sh{O}_{Z,x}$.
\end{proof}

Quando $X$ é íntegro, $R(X)=\sh{O}_x=\kres(x)$ é um corpo \sref{I.7.1.5}; os corolários anteriores especializam-se então da seguinte forma:
\begin{corollary}[7.1.15]
\label{I.7.1.15}
Suponhamos que se satisfaçam as condições de \sref{I.7.1.12} e, além disso, que $X$ seja íntegro.
Seja $s$ a imagem de $x$ em $S$.
Então as $S$-aplicações racionais de $X$ em $Y$ identificam-se com os pontos geométricos de $Y\otimes_S\kres(s)$ a valores na extensão $R(X)$ de $\kres(s)$; em outros termos, cada uma delas equivale ao dado de um ponto $y\in Y$ situado sobre $s$ e de um $\kres(s)$-monomorfismo de $\kres(y)$ em $\kres(x)=R(X)$.
\end{corollary}

\begin{proof}
Com efeito, os pontos de $Y$ situados sobre $s$ identificam-se com os pontos de $Y\otimes_S\kres(s)$ \sref{I.3.6.3}, e os $\sh{O}_s$-homomorfismos locais $\sh{O}_y\to R(X)$ com os $\kres(s)$-monomorfismos $\kres(y)\to R(X)$.
\end{proof}

Mais precisamente:
\begin{corollary}[7.1.16]
\label{I.7.1.16}
Seja $k$ um corpo, e sejam $X$ e $Y$ dois pré-esquemas algébricos sobre $k$ \sref{I.6.4.1}; suponhamos, além disso, que $X$ seja íntegro.
Então as $k$-aplicações racionais de $X$ em $Y$ identificam-se com os pontos geométricos de $Y$ a valores na extensão $R(X)$ de $k$ \sref{I.3.4.4}.
\end{corollary}

\subsection{Domínio de definição de uma aplicação racional}
\label{subsection:I.7.2}
\label{I.7.2}

\begin{env}[7.2.1]
\label{I.7.2.1}
Sejam $X$ e $Y$ dois pré-esquemas, e seja $f$ uma aplicação racional de $X$ em $Y$.
Diz-se que $f$ está \emph{definida em um ponto $x\in X$} se existe um aberto denso $U$ de $X$ contendo $x$ e um morfismo $U\to Y$ pertencente à classe de equivalência $f$.
O conjunto dos pontos $x\in X$ nos quais $f$ está definida chama-se o \emph{domínio de definição de $f$}; é claro que é um aberto denso de $X$.
\end{env}

\oldpage[I]{159}
\begin{proposition}[7.2.2]
\label{I.7.2.2}
Sejam $X$ e $Y$ dois $S$-pré-esquemas tais que $X$ seja reduzido e $Y$ separado sobre $S$.
Seja $f$ uma $S$-aplicação racional de $X$ em $Y$, com domínio de definição $U_0$.
Existe então um único $S$-morfismo $U_0\to Y$ pertencente à classe de $f$.
\end{proposition}

Como, para todo morfismo $U\to Y$ pertencente à classe de $f$, tem-se necessariamente $U\subset U_0$, é claro que a proposição será consequência do seguinte lema:

\begin{lemma}[7.2.2.1]
\label{I.7.2.2.1}
Sob as hipóteses de \sref{I.7.2.2}, sejam $U_1$ e $U_2$ dois abertos densos de $X$, e sejam $f_i:U_i\to Y$ ($i=1,2$) dois $S$-morfismos tais que existe um aberto $V\subset U_1\cap U_2$, denso em $X$, no qual $f_1$ e $f_2$ coincidem.
Então $f_1$ e $f_2$ coincidem em $U_1\cap U_2$.
\end{lemma}

\begin{proof}
Podemos evidentemente nos reduzir ao caso em que $X=U_1=U_2$.
Como $X$ (e, portanto, $V$) é reduzido, $X$ é o menor subpré-esquema fechado de $X$ que contém $V$ \sref{I.5.2.2}.
Seja $g=(f_1,f_2)_S:X\to Y\times_S Y$; como, por hipótese, a diagonal $T=\Delta_Y(Y)$ é um subpré-esquema fechado de $Y\times_S Y$, $Z=g^{-1}(T)$ é um subpré-esquema fechado de $X$ \sref{I.4.4.1}.
Se $h:V\to Y$ é a restrição comum de $f_1$ e $f_2$ a $V$, a restrição de $g$ a $V$ é $g'=(h,h)_S$, que se fatora como $g'=\Delta_Y\circ h$; como $\Delta_Y^{-1}(T)=Y$, tem-se $g'^{-1}(T)=V$, e, portanto, $Z$ é um subpré-esquema fechado de $X$ que induz $V$, logo contém $V$, o que implica $Z=X$.
Da relação $g^{-1}(T)=X$, deduz-se de \sref{I.4.4.1} que $g$ se fatora como $\Delta_Y\circ f$, onde $f$ é um morfismo $X\to Y$, o que implica, pela definição do morfismo diagonal, que $f_1=f_2=f$.
\end{proof}

É claro que o morfismo $U_0\to Y$ definido em \sref{I.7.2.2} é o único morfismo da classe de $f$ que \emph{não pode ser prolongado} a um morfismo de um aberto de $X$ que contenha estritamente $U_0$.
\emph{Sob as hipóteses de \sref{I.7.2.2}}, podem-se, portanto, \emph{identificar} as aplicações racionais de $X$ em $Y$ com os morfismos, definidos em abertos densos de $X$ e com valores em $Y$, que \emph{não podem ser prolongados} a abertos estritamente maiores.
Com essa identificação, a Proposição~\sref{I.7.2.2} implica:

\begin{corollary}[7.2.3]
\label{I.7.2.3}
Sob as hipóteses de \sref{I.7.2.2} sobre $X$ e $Y$, seja $U$ um aberto denso de $X$.
Existe então uma correspondência biunívoca canônica entre os $S$-morfismos de $U$ em $Y$ e as $S$-aplicações racionais de $X$ em $Y$ definidas em todos os pontos de $U$.
\end{corollary}

\begin{proof}
Em virtude de \sref{I.7.2.2}, para todo $S$-morfismo $f$ de $U$ em $Y$, existe uma única $S$-aplicação racional $\overline{f}$ de $X$ em $Y$ que prolonga $f$.
\end{proof}

\begin{corollary}[7.2.4]
\label{I.7.2.4}
Sejam $S$ um esquema, $X$ um $S$-pré-esquema reduzido, $Y$ um $S$-esquema, e $f:U\to Y$ um $S$-morfismo de um aberto denso $U$ de $X$ em $Y$.
Se $\overline{f}$ é a $\bb{Z}$-aplicação racional de $X$ em $Y$ que prolonga $f$, então $\overline{f}$ é um $S$-morfismo (e é, portanto, a $S$-aplicação racional de $X$ em $Y$ que prolonga $f$).
\end{corollary}

\begin{proof}
Com efeito, se $\vphi:X\to S$ e $\psi:Y\to S$ são os morfismos estruturais, $U_0$ é o domínio de definição de $\overline{f}$, e $j$ é a injeção $U_0\to X$, basta demonstrar que $\psi\circ\overline{f}=\vphi\circ j$, o que resulta de \sref{I.7.2.2.1}, pois $f$ é um $S$-morfismo.
\end{proof}

\begin{corollary}[7.2.5]
\label{I.7.2.5}
Sejam $X$ e $Y$ dois $S$-pré-esquemas; suponhamos que $X$ seja reduzido e que $X$ e $Y$ sejam separados sobre $S$.
Sejam $p:Y\to X$ um $S$-morfismo (que faz de $Y$ um $X$-pré-esquema), $U$ um aberto denso de $X$ e $f$ uma $U$-seção de $Y$; então a aplicação racional $\overline{f}$ de $X$ em $Y$ que prolonga $f$ é uma $X$-seção racional de $Y$.
\end{corollary}

\oldpage[I]{160}
\begin{proof}
É preciso demonstrar que $p\circ\overline{f}$ é a identidade no domínio de definição de $\overline{f}$; como $X$ é separado sobre $S$, isto resulta novamente de \sref{I.7.2.2.1}.
\end{proof}

\begin{corollary}[7.2.6]
\label{I.7.2.6}
Sejam $X$ um pré-esquema reduzido e $U$ um aberto denso de $X$.
Existe então uma correspondência biunívoca canônica entre as seções de $\sh{O}_X$ sobre $U$ e as funções racionais sobre $X$ definidas em todos os pontos de $U$.
\end{corollary}

\begin{proof}
Tendo em conta \sref{I.7.2.3}, \sref{I.7.1.2} e \sref{I.7.1.3}, basta observar que o $X$-pré-esquema $X\otimes_{\bb{Z}}\bb{Z}[T]$ é separado sobre $X$ \sref{I.5.5.1}[iv].
\end{proof}

\begin{corollary}[7.2.7]
\label{I.7.2.7}
Sejam $Y$ um pré-esquema reduzido, $f:X\to Y$ um morfismo separado, $U$ um aberto denso de $Y$, $g:U\to f^{-1}(U)$ uma $U$-seção de $f^{-1}(U)$, e $Z$ o subpré-esquema reduzido de $X$ cujo espaço subjacente é $\overline{g(U)}$ \sref{I.5.2.1}.
Para que $g$ seja a restrição de uma $Y$-seção de $X$ \emph{(em outras palavras \sref{I.7.2.5}, para que a aplicação racional de $Y$ em $X$ que prolonga $g$ esteja definida em toda parte)}, é necessário e suficiente que a restrição de $f$ a $Z$ seja um isomorfismo de $Z$ sobre $Y$.
\end{corollary}

\begin{proof}
A restrição de $f$ a $f^{-1}(U)$ é um morfismo separado \sref{I.5.5.1}[i], portanto $g$ é uma imersão fechada \sref{I.5.4.6} e, por conseguinte, $g(U)$\footnote{\emph{[Trad.] A fonte francesa imprime aqui $g(V)$, mas $V$ não está definido; a igualdade e os dados do corolário exigem $g(U)$.}} $=Z\cap f^{-1}(U)$, e o subpré-esquema induzido por $Z$ no aberto $g(U)$ de $Z$ é idêntico ao subpré-esquema fechado de $f^{-1}(U)$ associado a $g$ \sref{I.5.2.1}.
É claro, então, que a condição enunciada é suficiente, pois, se ela se verifica, e se $f_Z:Z\to Y$ é a restrição de $f$ a $Z$ e $\overline{g}:Y\to Z$ é o isomorfismo inverso, $\overline{g}$ prolonga $g$.
Reciprocamente, se $g$ é a restrição a $U$ de uma $Y$-seção $h$ de $X$, então $h$ é uma imersão fechada \sref{I.5.4.6}; logo, $h(Y)$ é fechado e, como está contido em $Z$, é igual a $Z$; de \sref{I.5.2.1} resulta que $h$ é necessariamente um isomorfismo de $Y$ sobre o subpré-esquema fechado $Z$ de $X$.
\end{proof}

\begin{env}[7.2.8]
\label{I.7.2.8}
Sejam $X$ e $Y$ dois $S$-pré-esquemas, sendo $X$ reduzido e $Y$ separado sobre $S$.
Seja $f$ uma $S$-aplicação racional de $X$ em $Y$, e seja $x$ um ponto de $X$; pode-se compor $f$ com o $S$-morfismo canônico $\Spec(\sh{O}_x)\to X$ \sref{I.2.4.1}, desde que a interseção de $\Spec(\sh{O}_x)$ com o domínio de definição de $f$ seja densa em $\Spec(\sh{O}_x)$ (identificado com o conjunto dos $z\in X$ tais que $x\in\overline{\{z\}}$ \sref{I.2.4.2}).
Isso ocorrerá nos casos seguintes:
\begin{enumerate}
  \item[1.º] $X$ é \emph{irredutível} (e, portanto, \emph{íntegro}), pois então o ponto genérico $\xi$ de $X$ é o ponto genérico de $\Spec(\sh{O}_x)$; como o domínio de definição $U$ de $f$ contém $\xi$, $U\cap\Spec(\sh{O}_x)$ contém $\xi$ e, portanto, é denso em $\Spec(\sh{O}_x)$.
  \item[2.º] $X$ é \emph{localmente noetheriano}; nossa afirmação resulta então do seguinte:
\end{enumerate}
\end{env}

\begin{lemma}[7.2.8.1]
\label{I.7.2.8.1}
Seja $X$ um pré-esquema cujo espaço subjacente é localmente noetheriano, e seja $x$ um ponto de $X$.
As componentes irredutíveis de $\Spec(\sh{O}_x)$ são as interseções com $\Spec(\sh{O}_x)$ das componentes irredutíveis de $X$ que contêm $x$.
Para que um aberto $U\subset X$ seja tal que $U\cap\Spec(\sh{O}_x)$ seja denso em $\Spec(\sh{O}_x)$, é necessário e suficiente que tenha interseção não vazia com cada uma das componentes irredutíveis de $X$ que contêm $x$ \emph{(o que ocorre, em particular, se $U$ é \emph{denso} em $X$)}.
\end{lemma}

\begin{proof}
Basta demonstrar a primeira afirmação, pois a segunda resulta evidentemente dela.
Como $\Spec(\sh{O}_x)$ está contido em todo aberto afim $U$ que contém $x$, e como as componentes irredutíveis de $U$ que contêm $x$ são as interseções com $U$ das componentes irredutíveis de $X$ que contêm $x$ \sref[0]{0.2.1.6}, podemos supor que $X$ é afim, de anel $A$.
Como os ideais primos de $A_x$ correspondem biunivocamente aos ideais primos
\oldpage[I]{161}
de $A$ contidos em $\mathfrak{j}_x$ \sref[0]{0.1.2.6}, os ideais primos minimais de $A_x$ correspondem aos ideais primos minimais de $A$ contidos em $\mathfrak{j}_x$, de onde resulta o lema \sref{I.1.1.14}.
\end{proof}

Dito isso, suponhamos que estejamos em um dos dois casos mencionados em \sref{I.7.2.8}.
Se $U$ é o domínio de definição da $S$-aplicação racional $f$, denotemos por $f'$ a aplicação racional de $\Spec(\sh{O}_x)$ em $Y$ que, levando em conta \sref{I.2.4.2}, coincide com $f$ em $U\cap\Spec(\sh{O}_x)$; diremos que essa aplicação racional é \emph{induzida por $f$}.

\begin{proposition}[7.2.9]
\label{I.7.2.9}
Sejam $S$ um pré-esquema localmente noetheriano, $X$ um $S$-pré-esquema reduzido e $Y$ um $S$-esquema de tipo finito.
Suponha-se, além disso, que $X$ seja irredutível ou localmente noetheriano.
Sejam então $f$ uma $S$-aplicação racional de $X$ em $Y$, e $x$ um ponto de $X$.
Para que $f$ esteja definida no ponto $x$, é necessário e suficiente que a aplicação racional $f'$ de $\Spec(\sh{O}_x)$ em $Y$, induzida por $f$ \sref{I.7.2.8}, seja um morfismo.
\end{proposition}

\begin{proof}
A condição é evidentemente necessária (pois $\Spec(\sh{O}_x)$ está contido em todo aberto que contém $x$); demonstremos que ela é suficiente.
Em virtude de \sref{I.6.5.1}, existe uma vizinhança aberta $U$ de $x$ em $X$ e um $S$-morfismo $g$ de $U$ em $Y$ que induz $f'$ sobre $\Spec(\sh{O}_x)$.
Se $X$ é irredutível, então $U$ é denso em $X$ e, em virtude de \sref{I.7.2.3}, podemos considerar $g$ como uma $S$-aplicação racional.
Além disso, o ponto genérico de $X$ pertence a $\Spec(\sh{O}_x)$ e ao domínio de definição de $f$, de modo que $f$ e $g$ coincidem nesse ponto e, por conseguinte, em um aberto não vazio de $X$ \sref{I.6.5.1}.
Mas, como $f$ e $g$ são $S$-aplicações racionais, elas são idênticas \sref{I.7.2.3}; logo, $f$ está definida em $x$.

Suponhamos agora que $X$ seja localmente noetheriano; podemos supor que $U$ seja noetheriano.
Há então um número finito de componentes irredutíveis $X_i$ de $X$ que contêm $x$ \sref{I.7.2.8.1}, e podemos supor que sejam as únicas que têm interseção não vazia com $U$, substituindo, se necessário, $U$ por um aberto menor (pois somente um número finito de componentes irredutíveis de $X$ tem interseção não vazia com $U$, já que $U$ é noetheriano).
Como acima, vê-se então que $f$ e $g$ coincidem em um aberto não vazio de cada um dos $X_i$.
Levando em conta que cada um dos $X_i$ está contido em $\overline{U}$, consideremos então o morfismo $f_1$, definido em um aberto denso de $U\cup(X\setmin\overline{U})$, igual a $g$ em $U$ e a $f$ na interseção de $X\setmin\overline{U}$ com o domínio de definição de $f$.
Como $U\cup(X\setmin\overline{U})$ é denso em $X$, $f_1$ e $f$ coincidem em um aberto denso de $X$; e, como $f$ é uma aplicação racional, $f$ é um prolongamento de $f_1$ \sref{I.7.2.3}, logo está definida no ponto $x$.
\end{proof}

\subsection{Feixe de funções racionais}
\label{subsection:I.7.3}

\begin{env}[7.3.1]
\label{I.7.3.1}
Seja $X$ um pré-esquema.
Para todo aberto $U\subset X$, denotamos por $R(U)$ o anel das funções racionais sobre $U$ \sref{I.7.1.3}; trata-se de uma $\Gamma(U,\sh{O}_X)$-álgebra.
Além disso, se $V\subset U$ é um segundo aberto de $X$, toda seção de $\sh{O}_X$ sobre uma parte aberta densa em $U$ fornece, por restrição a $V$, uma seção sobre uma parte aberta densa em $V$; e, se duas seções coincidem sobre uma parte aberta densa em $U$, suas restrições a $V$ coincidem sobre uma parte aberta densa em $V$.
Podemos, portanto, definir um di-homomorfismo de álgebras $\rho_{V,U}:R(U)\to R(V)$, e é claro que, se $U\supset V\supset W$ são abertos de $X$, então $\rho_{W,U}=\rho_{W,V}\circ\rho_{V,U}$; os $R(U)$ definem assim um \emph{pré-feixe} de álgebras sobre $X$.
\end{env}

\oldpage[I]{162}
\begin{definition}[7.3.2]
\label{I.7.3.2}
Chama-se feixe de funções racionais sobre um pré-esquema $X$, e denota-se por $\sh{R}(X)$, a $\sh{O}_X$-álgebra associada ao pré-feixe formado pelos $R(U)$.
\end{definition}

Para todo pré-esquema $X$ e todo aberto $U\subset X$, é claro que o feixe induzido $\sh{R}(X)|U$ não é senão $\sh{R}(U)$.

\begin{proposition}[7.3.3]
\label{I.7.3.3}
Seja $X$ um pré-esquema tal que a família $(X_\lambda)$ de suas componentes irredutíveis seja localmente finita (o que ocorre, em particular, quando o espaço subjacente a $X$ é localmente noetheriano).
Então o $\sh{O}_X$-módulo $\sh{R}(X)$ é quase-coerente e, para todo aberto $U$ de $X$ que intercepte somente um número finito das componentes $X_\lambda$, $R(U)$ é igual a $\Gamma(U,\sh{R}(X))$ e identifica-se canonicamente com o produto direto dos anéis locais dos pontos genéricos das $X_\lambda$ tais que $U\cap X_\lambda\neq\emp$.
\end{proposition}

\begin{proof}
Podemos evidentemente nos restringir ao caso em que $X$ tenha apenas um número finito de componentes irredutíveis $X_i$, com pontos genéricos $x_i$ ($1\leq i\leq n$).
O fato de que $R(U)$ se identifica canonicamente com o produto direto dos $\sh{O}_{x_i}=R(X_i)$ tais que $U\cap X_i\neq\emp$ decorre então de \sref{I.7.1.7}.
Mostremos, além disso, que o pré-feixe $U\to R(U)$ satisfaz os axiomas dos feixes, o que provará que $R(U)=\Gamma(U,\sh{R}(X))$.
Com efeito, ele satisfaz (F1) pelo que precede.
Para ver que satisfaz (F2), consideremos uma cobertura de um aberto $U$ de $X$ por abertos $V_\alpha\subset U$; se os $s_\alpha\in R(V_\alpha)$ são tais que as restrições de $s_\alpha$ e $s_\beta$ a $V_\alpha\cap V_\beta$ coincidem para todo par de índices, concluímos que, para todo índice $i$ tal que $U\cap X_i\neq\emp$, as componentes em $R(X_i)$ de todos os $s_\alpha$ tais que $V_\alpha\cap X_i\neq\emp$ são as mesmas; denotando por $t_i$ essa componente, é claro que o elemento de $R(U)$ que tem os $t_i$ como componentes admite $s_\alpha$ como restrição a cada $V_\alpha$.
Por fim, para ver que $\sh{R}(X)$ é quase-coerente, podemos nos restringir ao caso em que $X=\Spec(A)$ é afim; tomando para $U$ os abertos afins da forma $D(f)$, em que $f\in A$, resulta do que precede e da Definição~\sref{I.1.3.4} que $\sh{R}(X)=\widetilde{M}$, em que $M$ é a soma direta dos $A$-módulos $A_{x_i}$.
\end{proof}

\begin{corollary}[7.3.4]
\label{I.7.3.4}
Seja $X$ um pré-esquema reduzido que tenha apenas um número finito de componentes irredutíveis, e sejam $X_i$ ($1\leq i\leq n$) os pré-esquemas fechados reduzidos de $X$ que têm por espaços subjacentes as componentes irredutíveis de $X$ \sref{I.5.2.1}.
Se $h_i$ é a injeção canônica $X_i\to X$, então $\sh{R}(X)$ é o produto direto das $\sh{O}_X$-álgebras $(h_i)_*(\sh{R}(X_i))$.
\end{corollary}

\begin{corollary}[7.3.5]
\label{I.7.3.5}
Se $X$ é irredutível, todo $\sh{R}(X)$-módulo quase-coerente $\sh{F}$ é um feixe simples.
\end{corollary}

\begin{proof}
Basta mostrar que todo $x\in X$ admite uma vizinhança $U$ tal que $\sh{F}|U$ seja um feixe simples \sref[0]{0.3.6.2}; em outras palavras, reduzimo-nos ao caso em que $X$ é afim; podemos ainda supor que $\sh{F}$ é o conúcleo de um homomorfismo $(\sh{R}(X))^{I}\to(\sh{R}(X))^{J}$ \sref[0]{0.5.1.3}, e tudo se reduz a mostrar que $\sh{R}(X)$ é um feixe simples; mas isso é evidente, pois $\Gamma(U,\sh{R}(X))=R(X)$ para todo aberto não vazio $U$, já que $U$ contém o ponto genérico de $X$.
\end{proof}

\begin{corollary}[7.3.6]
\label{I.7.3.6}
Se $X$ é irredutível, para todo $\sh{O}_X$-módulo quase-coerente $\sh{F}$, $\sh{F}\otimes_{\sh{O}_X}\sh{R}(X)$ é um feixe simples; se, além disso, $X$ é reduzido (logo integral), $\sh{F}\otimes_{\sh{O}_X}\sh{R}(X)$ é isomorfo a um feixe da forma $(\sh{R}(X))^{(I)}$.
\end{corollary}

\begin{proof}
A segunda afirmação decorre do fato de que $R(X)$ é então um corpo.
\end{proof}

\begin{proposition}[7.3.7]
\label{I.7.3.7}
Suponhamos que o pré-esquema $X$ seja localmente integral ou localmente
\oldpage[I]{163}
noetheriano.
Então $\sh{R}(X)$ é uma $\sh{O}_X$-álgebra quase-coerente; se, além disso, $X$ é reduzido (o que ocorre quando $X$ é localmente integral), o homomorfismo canônico $\sh{O}_X\to\sh{R}(X)$ é injetivo.
\end{proposition}

\begin{proof}
Como a questão é local, a primeira afirmação decorre de \sref{I.7.3.3}; a segunda decorre imediatamente de \sref{I.7.2.3}.
\end{proof}

\begin{env}[7.3.8]
\label{I.7.3.8}
\footnote{\emph{[Trad.] O texto original deste item, declarado errôneo, foi inteiramente substituído na errata de EGA~II.}}
Sejam $X$ e $Y$ dois pré-esquemas \emph{integrais}, o que implica que $\sh{R}(X)$ (resp. $\sh{R}(Y)$) é um $\sh{O}_X$-módulo (resp. um $\sh{O}_Y$-módulo) quase-coerente \sref{I.7.3.3}.
Seja $f:X\to Y$ um morfismo \emph{dominante}; então existe um homomorfismo canônico de $\sh{O}_X$-módulos
\[
\label{I.7.3.8.1}
  \tau:f^*(\sh{R}(Y))\to\sh{R}(X).
  \tag{7.3.8.1}
\]
\end{env}

\begin{proof}
Suponhamos primeiro que $X=\Spec(A)$ e $Y=\Spec(B)$ sejam afins, com $A$ e $B$ domínios de integridade; então $f$ corresponde a um homomorfismo injetivo $B\to A$, que se estende a um monomorfismo $L\to K$ do corpo de frações $L$ de $B$ no corpo de frações $K$ de $A$.
O homomorfismo \sref{I.7.3.8.1} corresponde então ao homomorfismo canônico $L\otimes_B A\to K$ \sref{I.1.6.5}.
No caso geral, para todo par de abertos afins não vazios $U\subset X$ e $V\subset Y$ tais que $f(U)\subset V$, definimos, como acima, um homomorfismo $\tau_{U,V}$ e constatamos imediatamente que, se $U'\subset U$, $V'\subset V$ e $f(U')\subset V'$, então $\tau_{U,V}$ estende $\tau_{U',V'}$; daí decorre nossa afirmação.
Se $x$ e $y$ são os pontos genéricos de $X$ e $Y$, respectivamente, então temos $f(x)=y$,
\[
  (f^*(\sh{R}(Y)))_x=\sh{O}_y\otimes_{\sh{O}_y}\sh{O}_x=\sh{O}_x
\]
\sref[0]{0.4.3.1} e $\tau_x$ é, portanto, um \emph{isomorfismo}.
\end{proof}

\subsection{Feixes de torção e feixes sem torção}
\label{subsection:I.7.4}
\label{I.7.4}

\begin{env}[7.4.1]
\label{I.7.4.1}
Seja $X$ um pré-esquema \emph{integral}.
Para todo $\sh{O}_X$-módulo $\sh{F}$, o homomorfismo canônico $\sh{O}_X\to\sh{R}(X)$ define, por tensorização, um homomorfismo (também dito \emph{canônico}) $\sh{F}\to\sh{F}\otimes_{\sh{O}_X}\sh{R}(X)$ que, em cada talo, não é senão o homomorfismo $z\to z\otimes1$ de $\sh{F}_x$ em $\sh{F}_x\otimes_{\sh{O}_x}R(X)$.
O \emph{núcleo $\sh{T}$} desse homomorfismo é um sub-$\sh{O}_X$-módulo de $\sh{F}$, chamado \emph{feixe de torção} de $\sh{F}$; ele é quase-coerente se $\sh{F}$ é quase-coerente (\sref{I.4.1.1} e \sref{I.7.3.6}).
Dizemos que $\sh{F}$ é \emph{sem torção} se $\sh{T}=0$, e que $\sh{F}$ é um \emph{feixe de torção} se $\sh{T}=\sh{F}$.
Para todo $\sh{O}_X$-módulo $\sh{F}$, $\sh{F}/\sh{T}$ é sem torção.
Deduz-se de \sref{I.7.3.5} que:
\end{env}

\begin{proposition}[7.4.2]
\label{I.7.4.2}
Se $X$ é um pré-esquema integral, todo $\sh{O}_X$-módulo quase-coerente sem torção $\sh{F}$ é isomorfo a um subfeixe $\sh{G}$ de um feixe simples da forma $(\sh{R}(X))^{(I)}$, gerado (como $\sh{R}(X)$-módulo) por $\sh{G}$.
\end{proposition}

A cardinalidade de $I$ é chamada de \emph{posto} de $\sh{F}$; para todo aberto afim não vazio $U$ de $X$, o posto de $\sh{F}$ é igual ao posto de $\Gamma(U,\sh{F})$ como $\Gamma(U,\sh{O}_X)$-módulo, como se vê de imediato ao considerar o ponto genérico de $X$, contido em $U$.
Em particular:

\begin{corollary}[7.4.3]
\label{I.7.4.3}
Sobre um pré-esquema integral $X$, todo $\sh{O}_X$-módulo quase-coerente sem torção e de posto $1$ (em particular, todo $\sh{O}_X$-módulo inversível) é isomorfo a um sub-$\sh{O}_X$-módulo de $\sh{R}(X)$, e reciprocamente.
\end{corollary}

\begin{corollary}[7.4.4]
\label{I.7.4.4}
Sejam $X$ um pré-esquema integral, $\sh{L}$ e $\sh{L}'$ dois $\sh{O}_X$-módulos sem torção, e $f$ (resp. $f'$) uma seção de $\sh{L}$ (resp. $\sh{L}'$) sobre $X$.
Para que se tenha $f\otimes f'=0$, é necessário e suficiente que uma das seções $f$ e $f'$ seja nula.
\end{corollary}

\begin{proof}
Seja $x$ o ponto genérico de $X$; temos, por hipótese, que $(f\otimes f')_x=f_x\otimes f'_x=0$.
Como $\sh{L}_x$ e $\sh{L}'_x$ podem ser identificados com sub-$\sh{O}_x$-módulos do corpo $\sh{O}_x$, a relação precedente implica $f_x=0$ ou $f'_x=0$, e, portanto, $f=0$ ou $f'=0$, pois $\sh{L}$ e $\sh{L}'$ são sem torção \sref{I.7.3.5}.
\end{proof}

\begin{proposition}[7.4.5]
\label{I.7.4.5}
Sejam $X$ e $Y$ dois pré-esquemas integrais, e seja $f:X\to Y$ um morfismo dominante.
Para todo $\sh{O}_X$-módulo quase-coerente sem torção $\sh{F}$, $f_*(\sh{F})$ é um $\sh{O}_Y$-módulo sem torção.
\end{proposition}

\oldpage[I]{164}
\begin{proof}
Como $f_*$ é exato à esquerda \sref[0]{0.4.2.1}, basta, em virtude de \sref{I.7.4.2}, provar a proposição no caso em que $\sh{F}=(\sh{R}(X))^{(I)}$.
Ora, todo aberto não vazio $U$ de $Y$ contém o ponto genérico de $Y$, logo $f^{-1}(U)$ contém o ponto genérico de $X$ \sref[0]{0.2.1.5}, logo temos $\Gamma(U,f_*(\sh{F}))=\Gamma(f^{-1}(U),\sh{F})=(\sh{R}(X))^{(I)}$; em outras palavras, $f_*(\sh{F})$ é o feixe simples com fibra $(\sh{R}(X))^{(I)}$, considerado como um $\sh{R}(Y)$-módulo, e é evidentemente sem torção.
\end{proof}

\begin{proposition}[7.4.6]
\label{I.7.4.6}
Seja $X$ um pré-esquema integral, e $x$ seu ponto genérico.
Para todo $\sh{O}_X$-módulo quase-coerente de tipo finito $\sh{F}$, as condições seguintes são equivalentes: \emph{(a)} $\sh{F}$ é um feixe de torção; \emph{(b)} $\sh{F}_x=0$; \emph{(c)} $\Supp(\sh{F})\neq X$.
\end{proposition}

\begin{proof}
Em virtude de \sref{I.7.3.5} e \sref{I.7.4.1}, as relações $\sh{F}_x=0$ e $\sh{F}\otimes_{\sh{O}_X}\sh{R}(X)=0$ são equivalentes, portanto (a) e (b) são equivalentes; por outro lado, $\Supp(\sh{F})$ é fechado em $X$ \sref[0]{0.5.2.2}, e como todo aberto não vazio de $X$ contém $x$, (b) e (c) são equivalentes.
\end{proof}

\begin{env}[7.4.7]
\label{I.7.4.7}
Estendem-se (por abuso de linguagem) as definições de \sref{I.7.4.1} ao caso em que $X$ é um pré-esquema \emph{reduzido} que tem apenas um número \emph{finito} de componentes irredutíveis; resulta então de \sref{I.7.3.4} que a equivalência entre a) e c) em \sref{I.7.4.6} continua válida para tal pré-esquema, desde que se substitua a condição c) pela condição de que $\Supp(\sh{F})$ não contém nenhuma componente irredutível de $X$.
\end{env}

\section{Os esquemas de Chevalley}
\label{section:I.8}

\subsection{Anéis locais aparentados}
\label{subsection:I.8.1}

Para todo anel local $A$, denotaremos por $\mathfrak{m}(A)$ o ideal maximal de $A$.
\begin{lemma}[8.1.1]
\label{I.8.1.1}
Sejam $A$ e $B$ dois anéis locais tais que $A\subset B$.
As condições seguintes são equivalentes.
\begin{enumerate}
  \item[{\rm(i)}] $\mathfrak{m}(B)\cap A=\mathfrak{m}(A)$.
  \item[{\rm(ii)}] $\mathfrak{m}(A)\subset\mathfrak{m}(B)$.
  \item[{\rm(iii)}] $1$ não pertence ao ideal de $B$ gerado por $\mathfrak{m}(A)$.
\end{enumerate}
\end{lemma}

\begin{proof}
É evidente que (i) implica (ii), e que (ii) implica (iii);
por fim, se (iii) é satisfeita, então $\mathfrak{m}(B)\cap A$ contém $\mathfrak{m}(A)$ e não contém $1$, sendo, portanto, igual a $\mathfrak{m}(A)$.

Quando são satisfeitas as condições equivalentes de \sref{I.8.1.1}, dizemos que $B$ \emph{domina} $A$;
isso equivale a dizer que a injeção $A\to B$ é um homomorfismo \emph{local}.
É claro que, no conjunto dos subanéis locais de um anel $R$, a relação de dominação é uma relação de ordem.
\end{proof}

\begin{env}[8.1.2]
\label{I.8.1.2}
Consideremos agora um \emph{corpo} $R$.
Para todo subanel $A$ de $R$, denotaremos por $L(A)$ o conjunto dos anéis locais $A_\mathfrak{p}$, onde $\mathfrak{p}$ percorre o espectro primo de $A$;
tais anéis locais são identificados com os subanéis de $R$ que contêm $A$.
Como $\mathfrak{p}=(\mathfrak{p}A_\mathfrak{p})\cap A$, a aplicação $\mathfrak{p}\mapsto A_\mathfrak{p}$ de $\Spec(A)$ em $L(A)$ é bijetiva.
\end{env}

\begin{lemma}[8.1.3]
\label{I.8.1.3}
Sejam $R$ um corpo e $A$ um subanel de $R$.
Para que um subanel local $M$ de $R$ domine um anel $A_\mathfrak{p}\in L(A)$, é necessário e suficiente que $A\subset M$;
o anel local $A_\mathfrak{p}$ dominado por $M$ é então único e corresponde a $\mathfrak{p}=\mathfrak{m}(M)\cap A$.
\end{lemma}

\begin{proof}
Com efeito, se $M$ domina $A_\mathfrak{p}$, então $\mathfrak{m}(M)\cap A_\mathfrak{p}=\mathfrak{p}A_\mathfrak{p}$ por \sref{I.8.1.1}, daí a unicidade de $\mathfrak{p}$;
por outro lado, se $A\subset M$, então $\mathfrak{m}(M)\cap A=\mathfrak{p}$ é primo em $A$, e, como $A-\mathfrak{p}\subset M$, temos $A_\mathfrak{p}\subset M$ e $\mathfrak{p}A_\mathfrak{p}\subset\mathfrak{m}(M)$, logo $M$ domina $A_\mathfrak{p}$.
\end{proof}

\oldpage[I]{165}
\begin{lemma}[8.1.4]
\label{I.8.1.4}
Sejam $R$ um corpo, $M$ e $N$ dois subanéis locais de $R$, e $P$ o subanel de $R$ gerado por $M\cup N$.
As condições seguintes são equivalentes.
\begin{enumerate}
  \item[{\rm(i)}] Existe um ideal primo $\mathfrak{p}$ de $P$ tal que $\mathfrak{m}(M)=\mathfrak{p}\cap M$ e $\mathfrak{m}(N)=\mathfrak{p}\cap N$.
  \item[{\rm(ii)}] O ideal $\mathfrak{a}$ gerado em $P$ por $\mathfrak{m}(M)\cup\mathfrak{m}(N)$ é distinto de $P$.
  \item[{\rm(iii)}] Existe um subanel local $Q$ de $R$ que domina simultaneamente $M$ e $N$.
\end{enumerate}
\end{lemma}

\begin{proof}
É claro que (i) implica (ii);
reciprocamente, se $\mathfrak{a}\neq P$, então $\mathfrak{a}$ está contido em um ideal maximal $\mathfrak{n}$ de $P$ e, como $1\not\in\mathfrak{n}$, $\mathfrak{n}\cap M$ contém $\mathfrak{m}(M)$ e é distinto de $M$, logo $\mathfrak{n}\cap M=\mathfrak{m}(M)$, e, do mesmo modo, $\mathfrak{n}\cap N=\mathfrak{m}(N)$.
É claro que, se $Q$ domina $M$ e $N$, então $P\subset Q$ e $\mathfrak{m}(M)=\mathfrak{m}(Q)\cap M=(\mathfrak{m}(Q)\cap P)\cap M$, e $\mathfrak{m}(N)=(\mathfrak{m}(Q)\cap P)\cap N$, logo (iii) implica (i);
a recíproca é evidente tomando $Q=P_\mathfrak{p}$.
\end{proof}

Quando são satisfeitas as condições de \sref{I.8.1.4}, dizemos, com C.~Chevalley, que os anéis locais $M$ e $N$ são \emph{aparentados}.

\begin{proposition}[8.1.5]
\label{I.8.1.5}
Sejam $A$ e $B$ dois subanéis de um corpo $R$, e $C$ o subanel de $R$ gerado por $A\cup B$.
As condições seguintes são equivalentes.
\begin{enumerate}
  \item[{\rm(i)}] Para todo anel local $Q$ que contenha $A$ e $B$, tem-se $A_\mathfrak{p}=B_\mathfrak{q}$, pondo $\mathfrak{p}=\mathfrak{m}(Q)\cap A$ e $\mathfrak{q}=\mathfrak{m}(Q)\cap B$.
  \item[{\rm(ii)}] Para todo ideal primo $\mathfrak{r}$ de $C$, tem-se $A_\mathfrak{p}=B_\mathfrak{q}$, pondo $\mathfrak{p}=\mathfrak{r}\cap A$ e $\mathfrak{q}=\mathfrak{r}\cap B$.
  \item[{\rm(iii)}] Se $M\in L(A)$ e $N\in L(B)$ são aparentados, então são idênticos.
  \item[{\rm(iv)}] $L(A)\cap L(B)=L(C)$.
\end{enumerate}
\end{proposition}

\begin{proof}
Os lemas~\sref{I.8.1.3} e \sref{I.8.1.4} provam que (i) e (iii) são equivalentes;
é claro que (i) implica (ii), tomando $Q=C_\mathfrak{r}$;
reciprocamente, (ii) implica (i), pois, se $Q$ contém $A\cup B$, então contém $C$, e, se $\mathfrak{r}=\mathfrak{m}(Q)\cap C$, tem-se $\mathfrak{p}=\mathfrak{r}\cap A$ e $\mathfrak{q}=\mathfrak{r}\cap B$, por \sref{I.8.1.3}.
É imediato que (iv) implica (i), pois, se $Q$ contém $A\cup B$, então domina um anel local $C_\mathfrak{r}\in L(C)$ por \sref{I.8.1.3};
por hipótese, tem-se $C_\mathfrak{r}\in L(A)\cap L(B)$, e \sref{I.8.1.1} e \sref{I.8.1.3} provam que $C_\mathfrak{r}=A_\mathfrak{p}=B_\mathfrak{q}$.
Provemos, por fim, que (iii) implica (iv).
Seja $Q\in L(C)$;
$Q$ domina algum $M\in L(A)$ e algum $N\in L(B)$ por \sref{I.8.1.3}; portanto, $M$ e $N$, sendo aparentados, são idênticos por hipótese.
Como então $C\subset M$, sabemos que $M$ domina algum $Q'\in L(C)$ por \sref{I.8.1.3}; portanto, $Q$ domina $Q'$, de onde necessariamente, por \sref{I.8.1.3}, $Q=Q'=M$, e assim $Q\in L(A)\cap L(B)$.
Reciprocamente, se $Q\in L(A)\cap L(B)$, então $C\subset Q$; logo, por \sref{I.8.1.3}, $Q$ domina algum $Q''\in L(C)\subset L(A)\cap L(B)$;
$Q$ e $Q''$, sendo aparentados, são idênticos; portanto, $Q''=Q\in L(C)$, o que conclui a demonstração.
\end{proof}

\subsection{Anéis locais de um pré-esquema integral}
\label{subsection:I.8.2}

\begin{env}[8.2.1]
\label{I.8.2.1}
Seja $X$ um pré-esquema \emph{integral}, e seja $R$ seu corpo de funções racionais, idêntico ao anel local do ponto genérico $a$ de $X$;
para todo $x\in X$, sabemos que $\sh{O}_x$ se identifica canonicamente com um subanel de $R$ \sref{I.7.1.5}, e, para toda função racional $f\in R$, o domínio de definição $\delta(f)$ de $f$ é o aberto formado pelos $x\in X$ tais que
$f\in\sh{O}_x$.
Resulta então de \sref{I.7.2.6} que, para todo aberto $U\subset X$, temos
\[
\label{I.8.2.1.1}
  \Gamma(U,\sh{O}_X)=\bigcap_{x\in U}\sh{O}_x.
  \tag{8.2.1.1}
\]
\end{env}

\begin{proposition}[8.2.2]
\label{I.8.2.2}
Seja $X$ um pré-esquema integral, e seja $R$ seu corpo de funções racionais. Para que $X$ seja um esquema, é necessário e suficiente que a relação «$\sh{O}_x$ e $\sh{O}_y$ são aparentados» \sref{I.8.1.4}, entre pontos $x$ e $y$ de $X$, implique que $x=y$.
\end{proposition}

\begin{proof}
Suponhamos satisfeita esta condição e mostremos que $X$ é separado.
Sejam $U$ e $V$ dois abertos afins distintos de $X$, de anéis $A$ e $B$, respectivamente, identificados com subanéis de $R$;
$U$ (resp. $V$) identifica-se então \sref{I.8.1.2} com $L(A)$ (resp. $L(B)$), e a hipótese implica \sref{I.8.1.5} que, se $C$ é o subanel de $R$ gerado por $A\cup B$, então $W=U\cap V$ se identifica com $L(A)\cap L(B)=L(C)$.
Além disso, sabemos (\cite{I-1}, p.~5-03, prop.~4~\emph{bis}) que todo subanel $E$ de $R$ é igual à interseção dos anéis locais pertencentes a $L(E)$;
assim, $C$ se identifica com a interseção dos anéis $\sh{O}_z$ para $z\in W$, ou seja \sref{I.8.2.1.1}, com $\Gamma(W,\sh{O}_X)$.
Consideremos então o subpré-esquema induzido por $X$ sobre $W$;
ao homomorfismo identidade $\vphi: C\to\Gamma(W,\sh{O}_X)$ corresponde \sref{I.2.2.4} um morfismo $\Phi=(\psi,\theta):W\to\Spec(C)$;
veremos que $\Phi$ é um \emph{isomorfismo} de pré-esquemas, de onde resultará que $W$ é um aberto \emph{afim}.
A identificação de $W$ com $L(C)=\Spec(C)$ mostra que $\psi$ é \emph{bijetiva}.
Por outro lado, para todo $x\in W$, $\theta_x^\sharp$ é a injeção $C_\mathfrak{r}\to\sh{O}_x$, onde $\mathfrak{r}=\mathfrak{m}_x\cap C$, e, por definição, $C_\mathfrak{r}$ se identifica com $\sh{O}_x$, de modo que $\theta_x^\sharp$ é bijetiva.
Resta, pois, mostrar que $\psi$ é um \emph{homeomorfismo}, ou seja, que, para todo subconjunto fechado $F\subset W$, $\psi(F)$ é fechado em
$\Spec(C)$.
Ora, $F$ é o traço sobre $W$ de um subconjunto fechado de $U$ da forma $V(\mathfrak{a})$, onde $\mathfrak{a}$ é um ideal de $A$;
mostraremos que $\psi(F)=V(\mathfrak{a}C)$, o que provará nossa afirmação.
Com efeito, os ideais primos de $C$ que contêm $\mathfrak{a}C$ são os ideais primos de $C$ que contêm $\mathfrak{a}$ e são, portanto, os ideais da forma $\psi(x)=\mathfrak{m}_x\cap C$, onde $\mathfrak{a}\subset\mathfrak{m}_x$ e $x\in W$;
como, para $x\in W$, $\mathfrak{a}\subset\mathfrak{m}_x$ equivale a $x\in V(\mathfrak{a})$ e, portanto, a $x\in F=W\cap V(\mathfrak{a})$, temos de fato $\psi(F)=V(\mathfrak{a}C)$\footnote{\emph{[Trad.] O texto francês imprime $x\in V(\mathfrak a)=W\cap F$ para $x\in U$. Como $F$ foi definido imediatamente antes como o traço de $V(\mathfrak a)$ sobre $W$, a igualdade exigida é $F=W\cap V(\mathfrak a)$ e esta passagem diz respeito a $x\in W$.}}.

Segue-se que $X$ é separado, pois $U\cap V$ é afim e seu anel $C$ é gerado pela união $A\cup B$ dos anéis de $U$ e $V$ \sref{I.5.5.6}.

Reciprocamente, suponhamos que $X$ seja separado, e sejam $x$ e $y$ pontos de $X$ tais que $\sh{O}_x$ e $\sh{O}_y$ sejam aparentados.
Seja $U$ (resp. $V$) um aberto afim que contenha $x$ (resp. $y$), de anel $A$ (resp. $B$);
sabemos então que $U\cap V$ é afim e que seu anel $C$ é gerado por $A\cup B$ \sref{I.5.5.6}.
Se $\mathfrak{p}=\mathfrak{m}_x\cap A$ e $\mathfrak{q}=\mathfrak{m}_y\cap B$, então $A_\mathfrak{p}=\sh{O}_x$ e $B_\mathfrak{q}=\sh{O}_y$, e, como $A_\mathfrak{p}$ e $B_\mathfrak{q}$ são aparentados, existe um ideal primo $\mathfrak{r}$ de $C$ tal que $\mathfrak{p}=\mathfrak{r}\cap A$ e $\mathfrak{q}=\mathfrak{r}\cap B$ \sref{I.8.1.4}.
Mas existe então um ponto $z\in U\cap V$ tal que $\mathfrak{r}=\mathfrak{m}_z\cap C$, pois $U\cap V$ é afim, e temos evidentemente $x=z$ e $y=z$, de onde $x=y$.
\end{proof}

\begin{corollary}[8.2.3]
\label{I.8.2.3}
Seja $X$ um esquema \emph{integral}, e sejam $x$ e $y$ dois pontos de $X$.
Para que $x\in\overline{\{y\}}$, é necessário e suficiente que
$\sh{O}_x\subset\sh{O}_y$, ou, em outras palavras, que toda função racional
definida em $x$ também esteja definida em $y$.
\end{corollary}

\begin{proof}
A condição é evidentemente necessária, pois o domínio de definição
$\delta(f)$ de uma função racional $f\in R$ é aberto; mostremos agora
que ela é suficiente.
Se $\sh{O}_x\subset\sh{O}_y$, existe um ideal primo $\mathfrak{p}$
de $\sh{O}_x$ tal que $\sh{O}_y$ domina
$(\sh{O}_x)_\mathfrak{p}$ \sref{I.8.1.3}; mas, por \sref{I.2.4.2}, existe
um ponto $z\in X$ tal que $x\in\overline{\{z\}}$ e
$\sh{O}_z=(\sh{O}_x)_\mathfrak{p}$; como $\sh{O}_z$ e $\sh{O}_y$ são
aparentados, temos $z=y$ por \sref{I.8.2.2}, de onde segue o corolário.
\end{proof}

\begin{corollary}[8.2.4]
\label{I.8.2.4}
Seja $X$ um esquema \emph{integral}; então a aplicação $x\mapsto\sh{O}_x$ é injetiva; em outras palavras, se $x$ e $y$ são dois pontos distintos de $X$, existe uma função racional definida em um desses pontos e não no outro.
\end{corollary}

\begin{proof}
Isso se segue de \sref{I.8.2.3} e do axioma $(T_0)$ \sref{I.2.1.4}.
\end{proof}

\begin{corollary}[8.2.5]
\label{I.8.2.5}
Seja $X$ um esquema integral cujo espaço subjacente é noetheriano;
quando $f$ percorre o corpo $R$ das funções racionais sobre $X$, os conjuntos $\delta(f)$ geram a topologia de $X$.
\end{corollary}

Com efeito, todo subconjunto fechado de $X$ é então uma união finita de subconjuntos fechados irredutíveis, isto é, de subconjuntos da forma $\overline{\{y\}}$ \sref{I.2.1.5}.
Ora, se $x\not\in\overline{\{y\}}$, existe uma função racional $f$ definida em $x$ e não em $y$ \sref{I.8.2.3}; em outras palavras, temos $x\in\delta(f)$ e $\delta(f)$ não intersecta $\overline{\{y\}}$.
O complemento de $\overline{\{y\}}$ é, portanto, uma união de conjuntos da forma $\delta(f)$ e, em virtude da primeira observação, todo aberto de $X$ é uma união de interseções finitas de abertos da forma $\delta(f)$.

\begin{env}[8.2.6]
\label{I.8.2.6}
O Corolário~\sref{I.8.2.5} mostra que a topologia de $X$ é inteiramente caracterizada pela família de anéis locais $(\sh{O}_x)_{x\in X}$ que têm $R$ como corpo de frações.
Equivale também a dizer que os subconjuntos fechados de $X$ são definidos do seguinte modo: dado um subconjunto finito $\{x_1,\ldots,x_n\}$ de $X$, considera-se o conjunto dos $y\in X$ tais que $\sh{O}_y\subset\sh{O}_{x_i}$ para pelo menos um índice $i$; esses conjuntos (para todas as escolhas de $\{x_1,\ldots,x_n\}$) são os subconjuntos fechados de $X$.
Além disso, uma vez conhecida a topologia de $X$, o feixe estrutural $\sh{O}_X$ é igualmente determinado pela família dos $\sh{O}_x$, pois $\Gamma(U,\sh{O}_X)=\bigcap_{x\in U}\sh{O}_x$, por \sref{I.8.2.1.1}.
A família $(\sh{O}_x)_{x\in X}$ determina, portanto, completamente o pré-esquema $X$ quando $X$ é um esquema integral cujo espaço subjacente é noetheriano.
\end{env}

\begin{proposition}[8.2.7]
\label{I.8.2.7}
Sejam $X$ e $Y$ dois esquemas integrais, $f:X\to Y$ um morfismo dominante \sref{I.2.2.6}, e $K$ (resp. $L$) o corpo das funções racionais sobre $X$ (resp. $Y$).
Então $L$ se identifica com um subcorpo de $K$ e, para todo $x\in X$, $\sh{O}_{f(x)}$ é o único anel local de $Y$ dominado por $\sh{O}_x$.
\end{proposition}

\begin{proof}
Com efeito, se $f=(\psi,\theta)$ e $a$ é o ponto genérico de $X$, $\psi(a)$ é o ponto genérico de $Y$ \sref[0]{0.2.1.5};
$\theta_a^\sharp$ é, por conseguinte, um monomorfismo do corpo $L=\sh{O}_{\psi(a)}$ no corpo $K=\sh{O}_a$.
Como todo aberto afim não vazio $U$ de $Y$ contém $\psi(a)$, resulta de \sref{I.2.2.4} que o homomorfismo $\Gamma(U,\sh{O}_Y)\to\Gamma(\psi^{-1}(U),\sh{O}_X)$ correspondente a $f$ é a restrição de $\theta_a^\sharp$ a $\Gamma(U,\sh{O}_Y)$.
Portanto, para todo $x\in X$, $\theta_x^\sharp$ é a restrição a $\sh{O}_{\psi(x)}$ de $\theta_a^\sharp$ e é, por conseguinte, um monomorfismo.
Sabemos, além disso, que $\theta_x^\sharp$ é um homomorfismo local; assim, se identificarmos $L$ com um subcorpo de $K$ por meio de $\theta_a^\sharp$, $\sh{O}_{\psi(x)}$ é dominado por $\sh{O}_x$ \sref{I.8.1.1};
é, aliás, o único anel local de $Y$ dominado por $\sh{O}_x$, pois dois anéis locais de $Y$ que são aparentados são idênticos \sref{I.8.2.2}.
\end{proof}

\begin{proposition}[8.2.8]
\label{I.8.2.8}
Seja $X$ um pré-esquema \emph{irredutível}, seja $f:X\to Y$ uma imersão local (resp. um isomorfismo local), e suponhamos, além disso, que $f$ seja separado. Então $f$ é uma imersão (resp. uma imersão aberta).
\end{proposition}

\begin{proof}
Seja $f=(\psi,\theta)$;
basta, nos dois casos, provar que $\psi$ é um \emph{homeomorfismo} de $X$ sobre $\psi(X)$ \sref{I.4.5.3}.
Substituindo $f$ por $f_\text{red}$ (\sref{I.5.1.6} e \sref{I.5.5.1}[vi]), podemos supor que $X$ e $Y$ são \emph{reduzidos}.
Se $Y'$ é o subpré-esquema fechado reduzido de $Y$ cujo espaço subjacente é $\overline{\psi(X)}$, então $f$ se fatora como $X\xrightarrow{f'}Y'\xrightarrow{j}Y$, onde $j$ é a injeção canônica \sref{I.5.2.2}.
Resulta de \sref{I.5.5.1}[v] que $f'$ continua sendo um morfismo separado; além disso, $f'$ continua sendo
\oldpage[I]{168}
uma imersão local (resp. um isomorfismo local), pois, sendo a questão local em $X$ e $Y$, para demonstrá-lo podemos nos limitar ao caso em que $f$ é uma imersão fechada (resp. uma imersão aberta), e nossa afirmação decorre então imediatamente de \sref{I.4.2.2}.

Podemos, portanto, supor que $f$ é um morfismo \emph{dominante}, o que implica que $Y$ também é irredutível \sref[0]{0.2.1.5}, logo $X$ e $Y$ são ambos \emph{integrais}.
Além disso, sendo a questão local em $Y$, podemos supor que $Y$ é um esquema afim;
como $f$ é separado, $X$ é um esquema \sref{I.5.5.1}[ii], e finalmente estamos nas hipóteses da Proposição~\sref{I.8.2.7}.
Então, para todo $x\in X$, $\theta_x^\sharp$ é injetivo;
mas a hipótese de que $f$ é uma imersão local implica que $\theta_x^\sharp$ é sobrejetivo \sref{I.4.2.2}, logo $\theta_x^\sharp$ é bijetivo, ou seja, com a identificação da Proposição~\sref{I.8.2.7}, temos $\sh{O}_{\psi(x)}=\sh{O}_x$.
Isso implica, pelo Corolário~\sref{I.8.2.4}, que $\psi$ é uma aplicação \emph{injetiva}, o que já demonstra a proposição quando $f$ é um isomorfismo local \sref{I.4.5.3}.
Se supomos apenas que $f$ é uma imersão local, para todo $x\in X$ existe uma vizinhança aberta $U$ de $x$ em $X$ e uma vizinhança aberta $V$ de $\psi(x)$ em $Y$ tais que a restrição de $\psi$ a $U$ é um homeomorfismo de $U$ sobre um subconjunto \emph{fechado} de $V$.
Ora, $U$ é denso em $X$, logo $\psi(U)$ é denso em $Y$ e \emph{a fortiori} em $V$, o que prova que $\psi(U)=V$;
como $\psi$ é injetiva, $\psi^{-1}(V)=U$, e isso termina de mostrar que $\psi$ é um homeomorfismo de $X$ sobre $\psi(X)$.
\end{proof}

\subsection{Esquemas de Chevalley}
\label{subsection:I.8.3}

\begin{env}[8.3.1]
\label{I.8.3.1}
Seja $X$ um esquema integral \emph{noetheriano}, e seja $R$ seu corpo de funções racionais;
denotemos por $X'$ o conjunto dos subanéis locais $\sh{O}_x\subset R$, onde $x$ percorre todos os pontos de $X$.
O conjunto $X'$ satisfaz as três condições seguintes.
\begin{enumerate}
  \item[(Sch.~1)] \emph{Para todo $M\in X'$, $R$ é o corpo de frações de $M$.}
  \item[(Sch.~2)] \emph{Existe um conjunto finito de subanéis noetherianos $A_i$ de $R$ tais que $X'=\bigcup_i L(A_i)$ e, para todo par de índices $i$, $j$, o subanel $A_{ij}$ de $R$ gerado por $A_i\cup A_j$ é uma álgebra de tipo finito sobre $A_i$.}
  \item[(Sch.~3)] \emph{Quaisquer dois elementos $M$ e $N$ de $X'$ que são aparentados são idênticos.}
\end{enumerate}
\end{env}

Vimos em \sref{I.8.2.1} que (Sch.~1) é satisfeita, e (Sch.~3) decorre de \sref{I.8.2.2}.
Para demonstrar (Sch.~2), basta cobrir $X$ por um número finito de abertos afins $U_i$ cujos anéis sejam noetherianos, e tomar $A_i=\Gamma(U_i,\sh{O}_X)$;
a hipótese de que $X$ é um esquema implica que $U_i\cap U_j$ é afim, e também que $\Gamma(U_i\cap U_j,\sh{O}_X)=A_{ij}$ \sref{I.5.5.6};
além disso, como o espaço $U_i$ é noetheriano, a imersão $U_i\cap U_j\to U_i$ é de tipo finito \sref{I.6.3.5}, logo $A_{ij}$ é uma álgebra de tipo finito sobre $A_i$ \sref{I.6.3.3}.

\begin{env}[8.3.2]
\label{I.8.3.2}
As estruturas cujos axiomas são (Sch.~1), (Sch.~2) e (Sch.~3) generalizam os ``esquemas'' no sentido de C.~Chevalley, que supõe, além disso, que $R$ seja uma extensão de tipo finito de um corpo $K$, e que os $A_i$ sejam álgebras de tipo finito sobre $K$ (o que torna supérflua uma parte de (Sch.~2)) \cite{I-1}.
Reciprocamente, se temos tal estrutura em um conjunto $X'$, podemos associar-lhe um esquema integral $X$ usando as observações de \sref{I.8.2.6}: o espaço subjacente de $X$ é igual a $X'$ munido da topologia definida em \sref{I.8.2.6}, e do feixe $\sh{O}_X$
\oldpage[I]{169}
tal que $\Gamma(U,\sh{O}_X)=\bigcap_{x\in U}\sh{O}_x$ para todo aberto $U\subset X$, com a definição evidente dos homomorfismos de restrição.
Deixamos ao leitor a tarefa de verificar que assim se obtém um esquema integral cujos anéis locais são os elementos de $X'$;
não utilizaremos esse resultado no que segue.
\end{env}

\section{Complementos sobre os feixes quase-coerentes}
\label{section:I.9}

\subsection{Produto tensorial de feixes quase-coerentes}
\label{subsection:I.9.1}

\begin{proposition}[9.1.1]
\label{I.9.1.1}
Seja $X$ um pré-esquema (resp. um pré-esquema localmente noetheriano). Sejam $\sh{F}$ e
$\sh{G}$ dois $\sh{O}_X$-módulos quase-coerentes (resp. coerentes); então
$\sh{F}\otimes_{\sh{O}_X}\sh{G}$ é quase-coerente (resp. coerente); além disso, é
de tipo finito se $\sh{F}$ e $\sh{G}$ são de tipo finito. Se
$\sh{F}$ admite uma apresentação finita e $\sh{G}$ é quase-coerente
(resp. coerente), então $\shHom_{\sh{O}_X}(\sh{F},\sh{G})$ é quase-coerente
(resp. coerente).
\end{proposition}

\begin{proof}
Sendo a questão local, podemos supor que $X$ é afim (resp. afim
noetheriano); além disso, se $\sh{F}$ é coerente, podemos supor que é o
conúcleo de um homomorfismo $\sh{O}_X^m\to\sh{O}_X^n$. As afirmações relativas aos
feixes quase-coerentes decorrem de \sref{I.1.3.12} e \sref{I.1.3.9}; as
relativas aos feixes coerentes decorrem de \sref{I.1.5.1} e do fato de
que, se $M$ e $N$ são módulos de tipo finito sobre um anel noetheriano $A$,
então $M\otimes_A N$ e $\Hom_A(M,N)$ são ambos $A$-módulos de tipo finito.
\end{proof}

\begin{definition}[9.1.2]
\label{I.9.1.2}
Sejam $X$ e $Y$ dois $S$-pré-esquemas, $p$ e
$q$ as projeções de $X\times_S Y$, e $\sh{F}$ (resp. $\sh{G}$) um
$\sh{O}_X$-módulo quase-coerente (resp. um $\sh{O}_Y$-módulo quase-coerente). Chama-se
produto tensorial de $\sh{F}$ e $\sh{G}$ sobre $\sh{O}_S$ (\emph{ou} sobre $S$),
denotado por $\sh{F}\otimes_{\sh{O}_S}\sh{G}$
(\emph{ou} $\sh{F}\otimes_S\sh{G}$), o produto tensorial
$p^*(\sh{F})\otimes_{\sh{O}_{X\times_S Y}}q^*(\sh{G})$ sobre o pré-esquema $X\times_S Y$.
\end{definition}

Se os $X_i$ ($1\leq i\leq n$) são $S$-pré-esquemas, e os $\sh{F}_i$ ($1\leq i\leq n$) são
$\sh{O}_{X_i}$-módulos quase-coerentes, define-se da mesma maneira o produto tensorial
$\sh{F}_1\otimes_S\sh{F}_2\otimes_S\cdots\otimes_S\sh{F}_n$ sobre o
pré-esquema $Z=X_1\times_S X_2\times_S\cdots\times_S X_n$; trata-se de um
$\sh{O}_Z$-módulo \emph{quase-coerente} em virtude de \sref{I.9.1.1} e
\sref[0]{0.5.1.4}; ele é \emph{coerente} se os $\sh{F}_i$ são coerentes \emph{e}
$Z$ é \emph{localmente noetheriano}, em virtude de \sref{I.9.1.1},
\sref[0]{0.5.3.11} e \sref{I.6.1.1}.

Observe-se que, se tomamos $X=Y=S$, a Definição~\sref{I.9.1.2} fornece novamente o produto
tensorial de $\sh{O}_S$-módulos. Além disso, como $q^*(\sh{O}_Y)=\sh{O}_{X\times_S Y}$
\sref[0]{0.4.3.4}, o produto $\sh{F}\otimes_S\sh{O}_Y$ se identifica
canonicamente com $p^*(\sh{F})$ e, da mesma maneira,
$\sh{O}_X\otimes_S\sh{G}$ se identifica canonicamente com $q^*(\sh{G})$. Em
particular, se tomamos $Y=S$ e denotamos por $f$ o morfismo estrutural $X\to Y$,
temos $\sh{O}_X\otimes_Y\sh{G}=f^*(\sh{G})$: o produto tensorial ordinário
e a imagem inversa aparecem, portanto, como casos particulares do produto
tensorial geral.

A Definição~\sref{I.9.1.2} implica imediatamente que, para $X$ e
$Y$ fixos, $\sh{F}\otimes_S\sh{G}$ é um \emph{bifuntor covariante aditivo e exato à direita} em $\sh{F}$ e $\sh{G}$.

\begin{proposition}[9.1.3]
\label{I.9.1.3}
Sejam $S$, $X$ e $Y$ três esquemas afins de anéis,
respectivamente, $A$, $B$ e $C$, sendo $B$ e $C$ álgebras sobre $A$. Seja $M$
(resp. $N$) um $B$-módulo (resp. um $C$-módulo), e seja
$\sh{F}=\widetilde{M}$ (resp. $\sh{G}=\widetilde{N}$) o
feixe quase-coerente associado; então $\sh{F}\otimes_S\sh{G}$ é
canonicamente isomorfo ao feixe associado ao $(B\otimes_A C)$-módulo
$M\otimes_A N$.
\end{proposition}

\begin{proof}
\oldpage[I]{170}
De fato, em virtude de \sref{I.1.6.5}, $\sh{F}\otimes_S\sh{G}$
é canonicamente isomorfo ao feixe associado ao $(B\otimes_A C)$-módulo
\[
  \big(M\otimes_B(B\otimes_A C)\big)\otimes_{B\otimes_A C}\big((B\otimes_A C)\otimes_C N\big)
\]
e, pelos isomorfismos canônicos entre produtos tensoriais, este último módulo é isomorfo a
\[
  M\otimes_B(B\otimes_A C)\otimes_C N=(M\otimes_B B)\otimes_A(C\otimes_C N)=M\otimes_A N.
\]
\end{proof}

\begin{proposition}[9.1.4]
\label{I.9.1.4}
Sejam $f:T\to X$ e $g:T\to Y$ dois $S$-morfismos, e $\sh{F}$ (resp. $\sh{G}$) um $\sh{O}_X$-módulo quase-coerente (resp. um $\sh{O}_Y$-módulo quase-coerente).
Então
\[
  (f,g)^*_S(\sh{F}\otimes_S\sh{G})=f^*(\sh{F})\otimes_{\sh{O}_T}g^*(\sh{G}).
\]
\end{proposition}

\begin{proof}
Se $p$, $q$ são as projeções de $X\times_S Y$, a fórmula decorre
das igualdades $(f,g)^*_S\circ p^*=f^*$ e
$(f,g)^*_S\circ q^*=g^*$ \sref[0]{0.3.5.5}, e do fato de que a imagem
inversa de um produto tensorial de feixes algébricos é o produto tensorial de suas imagens
inversas \sref[0]{0.4.3.3}.
\end{proof}

\begin{corollary}[9.1.5]
\label{I.9.1.5}
Sejam $f:X\to X'$ e $g:Y\to Y'$
dois $S$-morfismos, e $\sh{F}'$ (resp. $\sh{G}'$) um
$\sh{O}_{X'}$-módulo quase-coerente (resp. um $\sh{O}_{Y'}$-módulo quase-coerente). Então
\[
  (f\times_S g)^*(\sh{F}'\otimes_S\sh{G}')=f^*(\sh{F}')\otimes_S g^*(\sh{G}').
\]
\end{corollary}

\begin{proof}
Isso decorre de \sref{I.9.1.4} e do fato de que $f\times_S g=(f\circ p, g\circ q)_S$,
onde $p$ e $q$ são as projeções de $X\times_S Y$.
\end{proof}

\begin{corollary}[9.1.6]
\label{I.9.1.6}
Sejam $X$, $Y$ e $Z$ três $S$-pré-esquemas, e $\sh{F}$ (resp. $\sh{G}$, $\sh{H}$) um
$\sh{O}_X$-módulo quase-coerente (resp. um $\sh{O}_Y$-módulo quase-coerente, um
$\sh{O}_Z$-módulo quase-coerente); então o feixe $\sh{F}\otimes_S\sh{G}\otimes_S\sh{H}$ é a imagem inversa
de $(\sh{F}\otimes_S\sh{G})\otimes_S\sh{H}$ pelo isomorfismo canônico de
$X\times_S Y\times_S Z$ sobre $(X\times_S Y)\times_S Z$.
\end{corollary}

\begin{proof}
Esse isomorfismo se escreve $(p_1,p_2)_S\times_S p_3$, onde $p_1$, $p_2$ e $p_3$
são as projeções de $X\times_S Y\times_S Z$.

Analogamente, a imagem inversa de $\sh{G}\otimes_S\sh{F}$ pelo isomorfismo canônico de
$X\times_S Y$ sobre $Y\times_S X$ é $\sh{F}\otimes_S\sh{G}$.
\end{proof}

\begin{corollary}[9.1.7]
\label{I.9.1.7}
Se $X$ é um $S$-pré-esquema, todo $\sh{O}_X$-módulo quase-coerente $\sh{F}$ é a imagem
inversa de $\sh{F}\otimes_S\sh{O}_S$ pelo isomorfismo canônico de $X$ sobre $X\times_S S$
\sref{I.3.3.3}.
\end{corollary}

\begin{proof}
Esse isomorfismo é $(1_X,\vphi)_S$, onde $\vphi$ é o morfismo estrutural
$X\to S$, e o corolário decorre de \sref{I.9.1.4} e do fato de que
$\vphi^*(\sh{O}_S)=\sh{O}_X$.
\end{proof}

\begin{env}[9.1.8]
\label{I.9.1.8}
Seja $X$ um $S$-pré-esquema, $\sh{F}$ um $\sh{O}_X$-módulo
quase-coerente e $\vphi:S'\to S$ um morfismo; denotamos por
$\sh{F}_{(\vphi)}$ ou $\sh{F}_{(S')}$ o feixe quase-coerente
$\sh{F}\otimes_S\sh{O}_{S'}$ sobre $X\times_S S'=X_{(\vphi)}=X_{(S')}$; portanto,
$\sh{F}_{(S')}=p^*(\sh{F})$, onde $p$ é a projeção $X_{(S')}\to X$.
\end{env}

\begin{proposition}[9.1.9]
\label{I.9.1.9}
Seja $\vphi':S''\to S'$ um morfismo.
Para todo $\sh{O}_X$-módulo quase-coerente $\sh{F}$ sobre o $S$-pré-esquema $X$,
$(\sh{F}_{(\vphi)})_{(\vphi')}$ é a imagem inversa de
$\sh{F}_{(\vphi\circ\vphi')}$ pelo isomorfismo canônico
$(X_{(\vphi)})_{(\vphi')}\isoto X_{(\vphi\circ\vphi')}$
\sref{I.3.3.9}.
\end{proposition}

\begin{proof}
Isso decorre imediatamente das definições e de \sref{I.3.3.9}, e também se
escreve
\[
\label{I.9.1.9.1}
  (\sh{F}\otimes_S\sh{O}_{S'})\otimes_{S'}\sh{O}_{S''}=\sh{F}\otimes_S\sh{O}_{S''}.
  \tag{9.1.9.1}
\]
\end{proof}

\begin{proposition}[9.1.10]
\label{I.9.1.10}
Seja $Y$ um $S$-pré-esquema, e seja $f:X\to Y$ um $S$-morfismo.
Para todo $\sh{O}_Y$-módulo quase-coerente $\sh{G}$ e todo morfismo
$S'\to S$, temos
$(f_{(S')})^*(\sh{G}_{(S')})=(f^*(\sh{G}))_{(S')}$.
\end{proposition}

\begin{proof}
Isso decorre imediatamente da comutatividade do diagrama
\oldpage[I]{171}
\[
  \xymatrix{
    X_{(S')}\ar[r]^{f_{(S')}}\ar[d] &
    Y_{(S')}\ar[d]\\
    X\ar[r]_f &
    Y.
  }
\]
\end{proof}

\begin{corollary}[9.1.11]
\label{I.9.1.11}
Sejam $X$ e $Y$ dois $S$-pré-esquemas, e seja
$\sh{F}$ (resp. $\sh{G}$) um $\sh{O}_X$-módulo quase-coerente
(resp. um $\sh{O}_Y$-módulo quase-coerente). A imagem inversa do feixe
$(\sh{F}_{(S')})\otimes_{S'}(\sh{G}_{(S')})$ pelo isomorfismo canônico
$(X\times_S Y)_{(S')}\isoto(X_{(S')})\times_{S'}(Y_{(S')})$
\sref{I.3.3.10} é igual a $(\sh{F}\otimes_S\sh{G})_{(S')}$.
\end{corollary}

\begin{proof}
Se $p$ e $q$ são as projeções de $X\times_S Y$, o isomorfismo em questão
não é outro senão $(p_{(S')}, q_{(S')})_{S'}$; o corolário decorre então das
Proposições~\sref{I.9.1.4} e \sref{I.9.1.10}.
\end{proof}

\begin{proposition}[9.1.12]
\label{I.9.1.12}
Com as notações da Definição~\sref{I.9.1.2}, seja $z$
um ponto de $X\times_S Y$, e sejam $x=p(z)$ e $y=q(z)$; o talo
$(\sh{F}\otimes_S\sh{G})_z$ é isomorfo a
$(\sh{F}_x\otimes_{\sh{O}_x}\sh{O}_z)\otimes_{\sh{O}_z}(\sh{G}_y\otimes_{\sh{O}_y}\sh{O}_z)
  =\sh{F}_x\otimes_{\sh{O}_x}\sh{O}_z\otimes_{\sh{O}_y}\sh{G}_y$.
\end{proposition}

\begin{proof}
Como podemos nos reduzir ao caso afim, a proposição decorre da
fórmula~\sref{I.1.6.5.1}.
\end{proof}

\begin{corollary}[9.1.13]
\label{I.9.1.13}
Se $\sh{F}$ e $\sh{G}$ são de tipo finito, temos
\[
  \Supp(\sh{F}\otimes_S\sh{G})=p^{-1}(\Supp(\sh{F}))\cap q^{-1}(\Supp(\sh{G})).
\]
\end{corollary}

\begin{proof}
Como $p^*(\sh{F})$ e $q^*(\sh{G})$ são de tipo finito sobre
$\sh{O}_{X\times_S Y}$, reduzimo-nos, pela Proposição~\sref{I.9.1.12} e
por \sref[0]{0.1.7.5}, ao caso particular em que $\sh{G}=\sh{O}_Y$, isto
é, a demonstrar a fórmula seguinte:\footnote{\emph{[Trad.] O
original francês imprime $\Supp(p^{-1}(\sh{F}))$ em (9.1.13.1). A redução
que acabamos de fazer identifica $\sh{F}\otimes_S\sh{O}_Y$ com $p^*(\sh{F})$, e a
verificação seguinte $\sh{O}_z/\mathfrak{m}_x\sh{O}_z\neq0$ é o cálculo
dos talos dessa imagem inversa; por isso escrevemos $p^*(\sh{F})$.
A fórmula impressa é válida para a imagem inversa de feixes de grupos abelianos; aqui formulamos a afirmação sobre a imagem inversa de módulos.}}
\[
\label{I.9.1.13.1}
  \Supp(p^*(\sh{F}))=p^{-1}(\Supp(\sh{F})).
  \tag{9.1.13.1}
\]

O mesmo raciocínio que em \sref[0]{0.1.7.5} reduz a questão a verificar que, para todo
$z\in X\times_S Y$, temos $\sh{O}_z/\mathfrak{m}_x\sh{O}_z\neq0$ (com $x=p(z)$),
o que decorre do fato de que o homomorfismo $\sh{O}_x\to\sh{O}_z$ é \emph{local},
por hipótese.
\end{proof}

Deixamos ao leitor a tarefa de estender a um produto de um número qualquer
de fatores os resultados demonstrados neste número para dois fatores.

\subsection{Imagem direta de um feixe quase-coerente}
\label{subsection:I.9.2}

\begin{proposition}[9.2.1]
\label{I.9.2.1}
Seja $f:X\to Y$ um morfismo de pré-esquemas. Suponhamos que exista um
recobrimento $(Y_\alpha)$ de $Y$ por abertos afins com a seguinte
propriedade: cada $f^{-1}(Y_\alpha)$ admite um recobrimento
\emph{finito} $(X_{\alpha i})$ por abertos afins contidos em
$f^{-1}(Y_\alpha)$, tal que cada uma das interseções
$X_{\alpha i}\cap X_{\alpha j}$ seja, por sua vez, uma união
\emph{finita} de abertos afins. Nessas condições, para todo
$\sh{O}_X$-módulo quase-coerente $\sh{F}$,
$f_*(\sh{F})$ é um $\sh{O}_Y$-módulo quase-coerente.
\end{proposition}

\begin{proof}
Como a questão é local em $Y$, podemos supor que $Y$ é igual a um dos
$Y_\alpha$ e, portanto, suprimir os índices $\alpha$.
\begin{enumerate}
  \item[(a)] Suponhamos primeiro que os $X_i\cap X_j$ sejam eles próprios
    abertos \emph{afins}. Ponhamos $\sh{F}_i=\sh{F}|X_i$ e
    $\sh{F}_{ij}=\sh{F}|(X_i\cap X_j)$, e sejam $\sh{F}'_i$ e
    $\sh{F}'_{ij}$ as imagens diretas de $\sh{F}_i$ e
    $\sh{F}_{ij}$ pelas restrições de $f$ a $X_i$ e a
    $X_i\cap X_j$, respectivamente; sabemos que $\sh{F}'_i$ e
    $\sh{F}'_{ij}$ são quase-coerentes \sref{I.1.6.3}. Ponhamos
    $\sh{G}=\bigoplus_i\sh{F}'_i$ e
    $\sh{H}=\bigoplus_{i,j}\sh{F}'_{ij}$; $\sh{G}$ e $\sh{H}$ são
    $\sh{O}_Y$-módulos quase-coerentes. Definiremos um homomorfismo
    $u:\sh{G}\to\sh{H}$ tal que $f_*(\sh{F})$ seja o \emph{núcleo} de
    $u$; daí resultará que $f_*(\sh{F})$ é quase-coerente
    \sref{I.1.3.9}. Basta definir $u$
\oldpage[I]{172}
    como um homomorfismo de pré-feixes; levando em conta as definições
    de $\sh{G}$ e $\sh{H}$, basta, para todo aberto $W\subset Y$, definir
    um homomorfismo
    \[
      u_W:\bigoplus_i\Gamma(f^{-1}(W)\cap X_i,\sh{F})\to
      \bigoplus_{i,j}\Gamma(f^{-1}(W)\cap X_i\cap X_j,\sh{F})
    \]
    que satisfaça as condições usuais de compatibilidade quando $W$
    varia. Se, para toda seção
    $s_i\in\Gamma(f^{-1}(W)\cap X_i,\sh{F})$, denotarmos por $s_{i|j}$ sua
    restrição a $f^{-1}(W)\cap X_i\cap X_j$, pomos
    \[
      u_W\big((s_i)\big)=(s_{i|j}-s_{j|i})
    \]
    e as condições de compatibilidade são satisfeitas de maneira
    evidente. Para provar que o núcleo $\sh{R}$ de $u$ é
    $f_*(\sh{F})$, definimos um homomorfismo de $f_*(\sh{F})$ em
    $\sh{R}$ fazendo corresponder a cada seção
    $s\in\Gamma(f^{-1}(W),\sh{F})$ a família $(s_i)$, onde $s_i$ é a
    restrição de $s$ a $f^{-1}(W)\cap X_i$; os axiomas (F1) e (F2) dos
    feixes (G,~II,~1.1) implicam que esse homomorfismo é
    \emph{bijetivo}, o que termina a demonstração neste caso.
  \item[(b)] No caso geral, aplica-se o mesmo raciocínio uma vez
    estabelecido que os $\sh{F}'_{ij}$ são quase-coerentes. Ora, por
    hipótese, $X_i\cap X_j$ é uma união finita de abertos afins
    $X_{ijk}$; e, como os $X_{ijk}$ são abertos afins \emph{em um
    esquema}, a interseção de dois quaisquer deles é ainda um aberto
    afim \sref{I.5.5.6}. Reduzimo-nos assim ao primeiro caso, e fica
    demonstrada a Proposição~\sref{I.9.2.1}.
\end{enumerate}
\end{proof}

\begin{corollary}[9.2.2]
\label{I.9.2.2}
A conclusão da Proposição~\sref{I.9.2.1} é válida em cada um dos casos seguintes:
\begin{enumerate}
  \item[{\rm(a)}] $f$ é separado e quase-compacto;
  \item[{\rm(b)}] $f$ é separado e de tipo finito;
  \item[{\rm(c)}] $f$ é quase-compacto, e o espaço subjacente a $X$ é localmente noetheriano.
\end{enumerate}
\end{corollary}

\begin{proof}
No caso~(a), os $X_{\alpha i}\cap X_{\alpha j}$ são afins \sref{I.5.5.6}.
O caso~(b) é um caso particular do caso~(a) \sref{I.6.6.3}. Por fim, no caso~(c), podemos nos reduzir ao caso em que $Y$ é afim e o espaço subjacente a $X$ é noetheriano; então $X$ admite um recobrimento finito por abertos afins $(X_i)$, e os $X_i\cap X_j$, por serem quase-compactos, são uniões finitas de abertos afins \sref{I.2.1.3}.
\end{proof}

\subsection{Extensão de seções de feixes quase-coerentes}
\label{subsection:I.9.3}

\begin{theorem}[9.3.1]
\label{I.9.3.1}
Seja $X$ um pré-esquema cujo espaço subjacente seja noetheriano, ou um esquema cujo espaço
subjacente seja quase-compacto. Seja $\sh{L}$ um $\sh{O}_X$-módulo invertível \sref[0]{0.5.4.1}, $f$ uma
seção de $\sh{L}$ sobre $X$, $X_f$ o aberto formado pelos $x\in X$ tais que $f(x)\neq0$
\sref[0]{0.5.5.1}, e $\sh{F}$ um $\sh{O}_X$-módulo quase-coerente.
\begin{enumerate}
  \item[{\rm(i)}] Se $s\in\Gamma(X,\sh{F})$ é tal que $s|X_f=0$, existe um inteiro $n>0$ tal que $s\otimes f^{\otimes n}=0$.
  \item[{\rm(ii)}] Para toda seção $s\in\Gamma(X_f,\sh{F})$, existe um inteiro $n>0$ tal que $s\otimes f^{\otimes n}$ se estende a uma seção de $\sh{F}\otimes\sh{L}^{\otimes n}$ sobre $X$.
\end{enumerate}
\end{theorem}

\begin{proof}
\medskip\noindent
\begin{enumerate}
  \item[(i)] Como o espaço subjacente a $X$ é quase-compacto e, portanto, união de
    um número finito de abertos afins $U_i$ tais que $\sh{L}|U_i$ é isomorfo a
    $\sh{O}_X|U_i$, podemos nos reduzir ao caso em que $X$ é afim e $\sh{L}=\sh{O}_X$.
    Nesse caso, $f$ se identifica com um elemento de $A(X)$ e temos
    $X_f=D(f)$; $s$ se identifica com um elemento de um $A(X)$-módulo $M$, e
    $s|X_f$ com o elemento correspondente de $M_f$; o resultado é então trivial,
    levando em conta a definição de módulo de frações.
\oldpage[I]{173}
  \item[(ii)] Novamente, $X$ é união finita de abertos afins $U_i$ ($1\leq i\leq r$)
    tais que $\sh{L}|U_i\cong\sh{O}_X|U_i$, e, para todo $i$,
    $(s\otimes f^{\otimes n})|(U_i\cap X_f)$ se identifica (mediante o isomorfismo
    precedente) com $(f|(U_i\cap X_f))^n(s|(U_i\cap X_f))$. Sabemos então
    \sref{I.1.4.1} que existe um inteiro $n>0$ tal que, para todo
    $i$, $(s\otimes f^{\otimes n})|(U_i\cap X_f)$ se estende a uma seção $s_i$ de
    $\sh{F}\otimes\sh{L}^{\otimes n}$ sobre $U_i$. Seja $s_{i|j}$ a restrição
    de $s_i$ a $U_i\cap U_j$; por definição, temos $s_{i|j}-s_{j|i}=0$ em
    $X_f\cap U_i\cap U_j$. Ora, se o espaço subjacente a $X$ é noetheriano, $U_i\cap U_j$ é
    quase-compacto; se $X$ é um esquema, $U_i\cap U_j$ é um aberto afim
    \sref{I.5.5.6} e, portanto, também quase-compacto. Em virtude de (i), existe, pois,
    um inteiro $m$ (independente de $i$ e $j$) tal que
    $(s_{i|j}-s_{j|i})\otimes f^{\otimes m}=0$. Deduz-se imediatamente que
    existe uma seção $s'$ de $\sh{F}\otimes\sh{L}^{\otimes(n+m)}$ sobre $X$ que
    induz $s_i\otimes f^{\otimes m}$ sobre cada $U_i$, e induz
    $s\otimes f^{\otimes(n+m)}$ sobre $X_f$.
\end{enumerate}
\end{proof}

Os corolários que seguem dão uma interpretação do Teorema~\sref{I.9.3.1} em uma linguagem
mais algébrica:
\begin{corollary}[9.3.2]
\label{I.9.3.2}
Sob as hipóteses de \sref{I.9.3.1}, consideremos o anel graduado $A_*=\Gamma_*(\sh{L})$
e o $A_*$-módulo graduado $M_*=\Gamma_*(\sh{L},\sh{F})$ \sref[0]{0.5.4.6}. Se $f\in A_n$,
onde $n\in\bb{Z}$, então há um isomorfismo canônico
$\Gamma(X_f,\sh{F})\isoto((M_*)_f)_0$ (\emph{o subgrupo do módulo de frações
$(M_*)_f$ formado pelos elementos de grau $0$}).
\end{corollary}

\begin{corollary}[9.3.3]
\label{I.9.3.3}
Suponhamos que as hipóteses de \sref{I.9.3.1} sejam satisfeitas e, além disso, que
$\sh{L}=\sh{O}_X$. Então, pondo $A=\Gamma(X,\sh{O}_X)$ e $M=\Gamma(X,\sh{F})$, o $A_f$-módulo
$\Gamma(X_f,\sh{F})$ é canonicamente isomorfo a $M_f$.
\end{corollary}

\begin{proposition}[9.3.4]
\label{I.9.3.4}
Sejam $X$ um pré-esquema noetheriano, $\sh{F}$ um $\sh{O}_X$-módulo coerente e $\sh{J}$ um
feixe coerente de ideais de $\sh{O}_X$, tais que o suporte de $\sh{F}$ esteja contido no de
$\sh{O}_X/\sh{J}$. Então existe um inteiro $n>0$ tal que $\sh{J}^n\sh{F}=0$.
\end{proposition}

\begin{proof}
Como $X$ é união finita de abertos afins cujos anéis são noetherianos, podemos supor que
$X$ seja afim, de anel $A$ noetheriano; então $\sh{F}=\widetilde{M}$, onde
$M=\Gamma(X,\sh{F})$ é um $A$-módulo de tipo finito, e $\sh{J}=\widetilde{\mathfrak{J}}$,
onde $\mathfrak{J}=\Gamma(X,\sh{J})$ é um ideal de $A$ (\sref{I.1.4.1} e
\sref{I.1.5.1}). Como $A$ é noetheriano, $\mathfrak{J}$ admite um sistema finito de
geradores $f_i$ ($1\leq i\leq m$). Por hipótese, toda seção de $\sh{F}$ sobre
$X$ é nula em cada um dos $D(f_i)$; se $s_j$ ($1\leq j\leq q$) são seções de
$\sh{F}$ que geram $M$, existe, portanto, um inteiro $h$, independente de $i$ e $j$,
tal que $f_i^h s_j=0$ \sref{I.1.4.1}, donde $f_i^h s=0$ para todo $s\in M$. Conclui-se
que, se $n=mh$, temos $\mathfrak{J}^n M=0$, e, portanto, o $\sh{O}_X$-módulo correspondente
$\sh{J}^n\sh{F}=\widetilde{\mathfrak{J}^n M}$ \sref{I.1.3.13} é nulo.
\end{proof}

\begin{corollary}[9.3.5]
\label{I.9.3.5}
Nas hipóteses de \sref{I.9.3.4}, existe um subpré-esquema fechado $Y$ de $X$ cujo
espaço subjacente é o suporte de $\sh{O}_X/\sh{J}$ e tal que, se $j:Y\to X$ é a
injeção canônica, então $\sh{F}=j_*(j^*(\sh{F}))$.
\end{corollary}

\begin{proof}
Observemos primeiro que os suportes de $\sh{O}_X/\sh{J}$ e de $\sh{O}_X/\sh{J}^n$ são
os mesmos, pois, se $\sh{J}_x=\sh{O}_x$, temos também $\sh{J}_x^n=\sh{O}_x$, e, por
outro lado, $\sh{J}_x^n\subset\sh{J}_x$ para todo $x\in X$. Podemos, portanto, em virtude
de \sref{I.9.3.4}, supor que $\sh{J}\sh{F}=0$; tomamos então por $Y$ o subpré-esquema
fechado de $X$ definido por $\sh{J}$ e, como $\sh{F}$ é então um
$(\sh{O}_X/\sh{J})$-módulo, a conclusão é imediata.
\end{proof}

\oldpage[I]{174}
\subsection{Prolongamento dos feixes quase-coerentes}
\label{subsection:I.9.4}

\begin{env}[9.4.1]
\label{I.9.4.1}
Sejam $X$ um espaço topológico, $\sh{F}$ um feixe de conjuntos (resp. de grupos, de anéis)
sobre $X$, $U$ um aberto de $X$, $\psi:U\to X$ a injeção canônica e $\sh{G}$ um subfeixe
de $\sh{F}|U=\psi^*(\sh{F})$. Como $\psi_*$ é exato à esquerda, $\psi_*(\sh{G})$ é um
subfeixe de $\psi_*(\psi^*(\sh{F}))$; denotamos por $\rho$ o homomorfismo canônico
$\sh{F}\to\psi_*(\psi^*(\sh{F}))$ \sref[0]{0.3.5.3}, e por $\overline{\sh{G}}$ o subfeixe
$\rho^{-1}(\psi_*(\sh{G}))$ de $\sh{F}$. Resulta imediatamente das definições que, para
todo aberto $V$ de $X$, $\Gamma(V,\overline{\sh{G}})$ é formado pelas seções
$s\in\Gamma(V,\sh{F})$ cuja restrição a $V\cap U$ é uma seção de $\sh{G}$ sobre $V\cap U$.
Temos, portanto, $\overline{\sh{G}}|U=\psi^*(\overline{\sh{G}})=\sh{G}$, e
$\overline{\sh{G}}$ é o \emph{maior} subfeixe de $\sh{F}$ que induz $\sh{G}$ sobre $U$;
dizemos que $\overline{\sh{G}}$ é o \emph{prolongamento canônico} do subfeixe $\sh{G}$ de
$\sh{F}|U$ a um subfeixe de $\sh{F}$.
\end{env}

\begin{proposition}[9.4.2]
\label{I.9.4.2}
Sejam $X$ um pré-esquema e $U$ uma parte aberta de $X$ tal que a injeção canônica
$j:U\to X$ seja um morfismo quase-compacto \emph{(o que se verificará para \emph{todo} $U$ se
o espaço subjacente a $X$ for \emph{localmente noetheriano}
\sref{I.6.6.4}[i])}. Então:
\begin{enumerate}
  \item[{\rm(i)}] para todo $(\sh{O}_X|U)$-módulo quase-coerente $\sh{G}$, $j_*(\sh{G})$
    é um $\sh{O}_X$-módulo quase-coerente, e
    $j_*(\sh{G})|U=j^*(j_*(\sh{G}))=\sh{G}$;
  \item[{\rm(ii)}] para todo $\sh{O}_X$-módulo quase-coerente $\sh{F}$ e todo
    $(\sh{O}_X|U)$-submódulo quase-coerente $\sh{G}$, o prolongamento canônico
    $\overline{\sh{G}}$ de $\sh{G}$ \sref{I.9.4.1} é um
    $\sh{O}_X$-submódulo quase-coerente de $\sh{F}$.
\end{enumerate}
\end{proposition}

\begin{proof}
Se $j=(\psi,\theta)$ (sendo $\psi$ a injeção $U\to X$ dos espaços subjacentes),
por definição temos $j_*(\sh{G})=\psi_*(\sh{G})$ para todo
$(\sh{O}_X|U)$-módulo $\sh{G}$ e, além disso,
$j^*(\sh{H})=\psi^*(\sh{H})=\sh{H}|U$ para todo $\sh{O}_X$-módulo $\sh{H}$, em virtude
da definição do pré-esquema induzido sobre um aberto. Portanto, (i) é um caso
particular de \sref{I.9.2.2}[a]; pela mesma razão, $j_*(j^*(\sh{F}))$ é
quase-coerente e, como $\overline{\sh{G}}$ é a imagem inversa de $j_*(\sh{G})$ pelo
homomorfismo $\rho:\sh{F}\to j_*(j^*(\sh{F}))$, (ii) decorre de \sref{I.4.1.1}.
\end{proof}

Observe-se que a hipótese de que o morfismo $j:U\to X$ seja quase-compacto também se
verifica quando o aberto $U$ é \emph{quase-compacto} e $X$ é um \emph{esquema}: com
efeito, $U$ é então uma união finita de abertos afins $U_i$ e, para todo aberto
afim $V$ de $X$, $V\cap U_i$ é um aberto afim \sref{I.5.5.6}, e portanto
quase-compacto.

\begin{corollary}[9.4.3]
\label{I.9.4.3}
Sejam $X$ um pré-esquema e $U$ um aberto quase-compacto de $X$ tal que o morfismo de
injeção $j:U\to X$ seja quase-compacto. Suponhamos, além disso, que todo
$\sh{O}_X$-módulo quase-coerente seja limite indutivo de seus $\sh{O}_X$-submódulos
quase-coerentes de tipo finito \emph{(o que ocorre quando $X$ é um
\emph{esquema afim})}. Sejam então $\sh{F}$ um $\sh{O}_X$-módulo quase-coerente e
$\sh{G}$ um $(\sh{O}_X|U)$-submódulo quase-coerente de tipo finito de $\sh{F}|U$.
Existe então um $\sh{O}_X$-submódulo quase-coerente de tipo finito $\sh{G}'$ de
$\sh{F}$ tal que $\sh{G}'|U=\sh{G}$.
\end{corollary}

\begin{proof}
Com efeito, temos $\sh{G}=\overline{\sh{G}}|U$, e $\overline{\sh{G}}$ é quase-coerente
por \sref{I.9.4.2}, portanto é limite indutivo de seus $\sh{O}_X$-submódulos
quase-coerentes de tipo finito $\sh{H}_\lambda$. Consequentemente, $\sh{G}$ é limite
indutivo dos $\sh{H}_\lambda|U$ e, portanto, igual a um dos $\sh{H}_\lambda|U$, pois é
de tipo finito \sref[0]{0.5.2.3}.
\end{proof}

\begin{remark}[9.4.4]
\label{I.9.4.4}
Suponhamos que, para \emph{todo} aberto afim $U\subset X$, o morfismo de injeção
$U\to X$ seja quase-compacto. Então, se a conclusão de \sref{I.9.4.3} valer para todo
aberto afim $U$ e todo $(\sh{O}_X|U)$-submódulo quase-coerente de tipo finito
$\sh{G}$ de $\sh{F}|U$,
\oldpage[I]{175}
segue-se que $\sh{F}$ é limite indutivo de seus $\sh{O}_X$-submódulos quase-coerentes
de tipo finito. Com efeito, para todo aberto afim $U\subset X$, temos
$\sh{F}|U=\widetilde{M}$, onde $M$ é um $A(U)$-módulo e, como este último é limite
indutivo de seus submódulos de tipo finito, $\sh{F}|U$ é limite indutivo de seus
$(\sh{O}_X|U)$-submódulos quase-coerentes de tipo finito \sref{I.1.3.9}. Ora, por
hipótese, cada um desses submódulos é induzido em $U$ por um $\sh{O}_X$-submódulo
quase-coerente de tipo finito $\sh{G}_{\lambda,U}$ de $\sh{F}$. As somas finitas dos
$\sh{G}_{\lambda,U}$ são ainda $\sh{O}_X$-módulos quase-coerentes de tipo finito, pois
a questão é local e o caso em que $X$ é afim foi tratado em \sref{I.1.3.10}; fica
então claro que $\sh{F}$ é limite indutivo dessas somas finitas, o que demonstra a
afirmação.
\end{remark}

\begin{corollary}[9.4.5]
\label{I.9.4.5}
Sob as hipóteses de \sref{I.9.4.3}, para todo $(\sh{O}_X|U)$-módulo quase-coerente
de tipo finito $\sh{G}$, existe um $\sh{O}_X$-módulo quase-coerente de tipo finito
$\sh{G}'$ tal que $\sh{G}'|U=\sh{G}$.
\end{corollary}

\begin{proof}
Como $\sh{F}=j_*(\sh{G})$ é quase-coerente \sref{I.9.4.2} e
$\sh{F}|U=\sh{G}$, basta aplicar \sref{I.9.4.3} a $\sh{F}$.
\end{proof}

\begin{lemma}[9.4.6]
\label{I.9.4.6}
Sejam $X$ um pré-esquema, $L$ um conjunto bem ordenado,
$(V_\lambda)_{\lambda\in L}$ uma cobertura de $X$ por abertos afins e $U$ um aberto
de $X$; para todo $\lambda\in L$, ponhamos
$W_\lambda=\bigcup_{\mu<\lambda}V_\mu$. Suponhamos que:
\begin{enumerate}
  \item[1\textsuperscript{o}] Para todo $\lambda\in L$,
    $V_\lambda\cap W_\lambda$ é quase-compacto;
  \item[2\textsuperscript{o}] O morfismo de imersão $U\to X$ é quase-compacto.
\end{enumerate}
Então, para todo $\sh{O}_X$-módulo quase-coerente $\sh{F}$ e todo
$(\sh{O}_X|U)$-submódulo quase-coerente de tipo finito $\sh{G}$ de $\sh{F}|U$, existe
um $\sh{O}_X$-submódulo quase-coerente de tipo finito $\sh{G}'$ de $\sh{F}$ tal que
$\sh{G}'|U=\sh{G}$.
\end{lemma}

\begin{proof}
Seja $\kappa$ o tipo de ordem de $L$ e indexemos a cobertura dada por
$\lambda<\kappa$.\footnote{Esclarecimento editorial: a indução impressa omite o caso inicial e, se $L$ tem um maior elemento, a extensão terminal. As etapas $0\leq\lambda\leq\kappa$ completam a demonstração sem alterar as hipóteses.}
Para $0\leq\lambda\leq\kappa$, ponhamos
$W_\lambda=\bigcup_{\mu<\lambda}V_\mu$ e $U_\lambda=U\cup W_\lambda$.
Construímos por indução transfinita um
$(\sh{O}_X|U_\lambda)$-submódulo quase-coerente de tipo finito $\sh{G}'_\lambda$ de
$\sh{F}|U_\lambda$, tal que $\sh{G}'_\lambda|U_\mu=\sh{G}'_\mu$ para $\mu<\lambda$ e
$\sh{G}'_\lambda|U=\sh{G}$.
Começamos com $\sh{G}'_0=\sh{G}$ sobre $U_0=U$; como $U_\kappa=X$, o submódulo
$\sh{G}'=\sh{G}'_\kappa$ dará a extensão desejada.
Suponhamos agora que $0<\lambda\leq\kappa$ e que os $\sh{G}'_\mu$ foram construídos
com essas propriedades para $\mu<\lambda$.
Se $\lambda$ é um ordinal limite, tomamos como $\sh{G}'_\lambda$ o único
$(\sh{O}_X|U_\lambda)$-submódulo de $\sh{F}|U_\lambda$ tal que
$\sh{G}'_\lambda|U_\mu=\sh{G}'_\mu$ para todo $\mu<\lambda$.
Os $U_\mu$ com $\mu<\lambda$ cobrem então $U_\lambda$, de modo que podemos colar
os submódulos compatíveis \sref[0]{0.3.3.1}; a quase-coerência e o tipo finito
são propriedades locais.
Se, ao
contrário, $\lambda=\mu+1$, temos $U_\lambda=U_\mu\cup V_\mu$, e basta definir um
$(\sh{O}_X|V_\mu)$-submódulo quase-coerente de tipo finito $\sh{G}''_\mu$ de
$\sh{F}|V_\mu$ tal que
\[
  \sh{G}''_\mu|(U_\mu\cap V_\mu)=\sh{G}'_\mu|(U_\mu\cap V_\mu);
\]
tomaremos depois como $\sh{G}'_\lambda$ o $(\sh{O}_X|U_\lambda)$-submódulo de
$\sh{F}|U_\lambda$ tal que $\sh{G}'_\lambda|U_\mu=\sh{G}'_\mu$ e
$\sh{G}'_\lambda|V_\mu=\sh{G}''_\mu$ \sref[0]{0.3.3.1}. Ora, como $V_\mu$ é afim,
a existência de $\sh{G}''_\mu$ é assegurada por \sref{I.9.4.3} assim que se demonstre
que $U_\mu\cap V_\mu$ é quase-compacto; mas $U_\mu\cap V_\mu$ é a união de
$U\cap V_\mu$ e $W_\mu\cap V_\mu$, ambos quase-compactos em virtude das hipóteses.
\end{proof}

\begin{theorem}[9.4.7]
\label{I.9.4.7}
Sejam $X$ um pré-esquema e $U$ um aberto de $X$. Suponhamos verificada uma das
condições seguintes:
\begin{enumerate}
  \item[{\rm(a)}] O espaço subjacente a $X$ é localmente noetheriano.
  \item[{\rm(b)}] $X$ é um esquema quase-compacto e $U$ um aberto quase-compacto.
\end{enumerate}
Então, para todo $\sh{O}_X$-módulo quase-coerente $\sh{F}$ e todo
$(\sh{O}_X|U)$-submódulo quase-coerente de tipo finito $\sh{G}$ de $\sh{F}|U$, existe
um $\sh{O}_X$-submódulo quase-coerente de tipo finito $\sh{G}'$ de $\sh{F}$ tal que
$\sh{G}'|U=\sh{G}$.
\end{theorem}

\oldpage[I]{176}
\begin{proof}
Seja $(V_\lambda)_{\lambda\in L}$ uma cobertura de $X$ por abertos afins,
supondo-se $L$ finito no caso (b); munindo $L$ de uma estrutura de conjunto bem
ordenado, basta verificar que as condições de \sref{I.9.4.6} são satisfeitas. Isso é
evidente sob a hipótese (a), pois os espaços $V_\lambda$ são noetherianos. Sob a
hipótese (b), os $V_\lambda\cap V_\mu$ são afins \sref{I.5.5.6}, e portanto
quase-compactos; como $L$ é finito, $V_\lambda\cap W_\lambda$ é quase-compacto. Daí o
teorema.
\end{proof}

\begin{corollary}[9.4.8]
\label{I.9.4.8}
Sob as hipóteses de \sref{I.9.4.7}, para todo $(\sh{O}_X|U)$-módulo quase-coerente
de tipo finito $\sh{G}$, existe um $\sh{O}_X$-módulo quase-coerente de tipo finito
$\sh{G}'$ tal que $\sh{G}'|U=\sh{G}$.
\end{corollary}

\begin{proof}
Basta aplicar \sref{I.9.4.7} a $\sh{F}=j_*(\sh{G})$, que é quase-coerente
\sref{I.9.4.2} e tal que $\sh{F}|U=\sh{G}$.
\end{proof}

\begin{corollary}[9.4.9]
\label{I.9.4.9}
Seja $X$ um pré-esquema cujo espaço subjacente é localmente noetheriano, ou um esquema
quase-compacto. Então todo $\sh{O}_X$-módulo quase-coerente é limite indutivo de seus
$\sh{O}_X$-submódulos quase-coerentes de tipo finito.
\end{corollary}

\begin{proof}
Isso resulta de \sref{I.9.4.7} e da observação \sref{I.9.4.4}.
\end{proof}

\begin{corollary}[9.4.10]
\label{I.9.4.10}
Sob as hipóteses de \sref{I.9.4.9}, se um $\sh{O}_X$-módulo quase-coerente
$\sh{F}$ é tal que todo $\sh{O}_X$-submódulo quase-coerente de tipo finito de $\sh{F}$
é gerado por suas seções sobre $X$, então $\sh{F}$ é gerado por suas
seções sobre $X$.
\end{corollary}

\begin{proof}
Seja $U$ uma vizinhança aberta afim de um ponto $x\in X$, e seja $s$ uma seção de
$\sh{F}$ sobre $U$; o $\sh{O}_X$-submódulo $\sh{G}$ de $\sh{F}|U$ gerado por $s$ é
quase-coerente e de tipo finito, portanto existe um $\sh{O}_X$-submódulo quase-coerente
de tipo finito $\sh{G}'$ de $\sh{F}$ tal que $\sh{G}'|U=\sh{G}$ \sref{I.9.4.7}. Por
hipótese, existem então um número finito de seções $t_i$ de $\sh{G}'$ sobre $X$ e
seções $a_i$ de $\sh{O}_X$ sobre uma vizinhança $V\subset U$ de $x$ tais que
$s|V=\sum_i a_i(t_i|V)$, o que demonstra o corolário.
\end{proof}

\subsection{Imagem fechada de um pré-esquema; aderência de um subpré-esquema}
\label{subsection:I.9.5}

\begin{proposition}[9.5.1]
\label{I.9.5.1}
Seja $f:X\to Y$ um morfismo de pré-esquemas tal que $f_*(\sh{O}_X)$ seja um
$\sh{O}_Y$-módulo quase-coerente (o que ocorre se $f$ é quase-compacto e, além disso,
ou $f$ é separado ou $X$ é localmente noetheriano \sref{I.9.2.2}). Então existe um
subpré-esquema fechado mínimo $Y'$ de $Y$ tal que $f$ se fatora pela injeção canônica
$j:Y'\to Y$ (\emph{ou, equivalentemente \sref{I.4.4.1}, tal que o subpré-esquema
$f^{-1}(Y')$ de $X$ seja \emph{idêntico} a $X$}).
\end{proposition}

Mais precisamente:
\begin{corollary}[9.5.2]
\label{I.9.5.2}
Sob as condições de \sref{I.9.5.1}, sejam $f=(\psi,\theta)$ e $\sh{J}$ o
núcleo (quase-coerente) do homomorfismo $\theta:\sh{O}_Y\to f_*(\sh{O}_X)$. Então o
subpré-esquema fechado $Y'$ de $Y$ definido por $\sh{J}$ satisfaz as condições de
\sref{I.9.5.1}.
\end{corollary}

\begin{proof}
Como o funtor $\psi^*$ é exato, a fatoração canônica
$\theta:\sh{O}_Y\to\sh{O}_Y/\sh{J}\xrightarrow{\theta'}\psi_*(\sh{O}_X)$ fornece (\textbf{0},~3.5.4.3)
uma fatoração
$\theta^\sharp:\psi^*(\sh{O}_Y)\to\psi^*(\sh{O}_Y)/\psi^*(\sh{J})
  \xrightarrow{{\theta'}^\sharp}\sh{O}_X$; como $\theta_x^\sharp$ é um homomorfismo local para
todo $x\in X$, o mesmo ocorre com ${\theta_x'}^\sharp$; se denotarmos por $\psi_0$ a
aplicação contínua $\psi$ considerada como aplicação de $X$ em \erratum[II]{$Y'$}, e por $\theta_0$ a restrição
$\theta'|Y':(\sh{O}_Y/\sh{J})|Y'\to\psi_*(\sh{O}_X)|Y'=(\psi_0)_*(\sh{O}_X)$, vemos então que
$f_0=(\psi_0,\theta_0)$ é um morfismo de pré-esquemas $X\to Y'$ \sref{I.2.2.1} tal que
$f=j\circ f_0$. Agora, se $Y''$
\oldpage[I]{177}
é um segundo subpré-esquema fechado de $Y$, definido por um feixe quase-coerente de ideais $\sh{J}'$ de
$\sh{O}_Y$, e tal que $f$ se fatora pela injeção $j':Y''\to Y$, devemos primeiro ter
$\psi(X)\subset Y''$ e, portanto, $Y'\subset Y''$, pois $Y''$ é fechado.
Além disso, para todo $y\in Y''$, $\theta$ deve fatorar-se como
$\sh{O}_y\to\sh{O}_y/\sh{J}'_y\to(\psi_*(\sh{O}_X))_y$, o que, por definição, implica
$\sh{J}'_y\subset\sh{J}_y$ e, portanto, $Y'$ é um subpré-esquema fechado de $Y''$
\sref{I.4.1.10}.
\end{proof}

\begin{definition}[9.5.3]
\label{I.9.5.3}
Sempre que existir um subpré-esquema fechado mínimo $Y'$ de $Y$ tal que $f$ se fatora
pela injeção canônica $j:Y'\to Y$, diremos que $Y'$ é o
\emph{pré-esquema imagem fechada} de $X$ pelo morfismo $f$.
\end{definition}

\begin{proposition}[9.5.4]
\label{I.9.5.4}
Se $f_*(\sh{O}_X)$ é um $\sh{O}_Y$-módulo quase-coerente, então o espaço subjacente
da imagem fechada de $X$ pelo morfismo $f$ é a aderência $\overline{f(X)}$ em $Y$.
\end{proposition}

\begin{proof}
Como o suporte de $f_*(\sh{O}_X)$ está contido em $\overline{f(X)}$, temos (com
a notação de \sref{I.9.5.2}) $\sh{J}_y=\sh{O}_y$ para
$y\not\in\overline{f(X)}$, portanto o suporte de $\sh{O}_Y/\sh{J}$ está contido em
$\overline{f(X)}$. Além disso, esse suporte é fechado e contém $f(X)$:
de fato, se $y\in f(X)$, o elemento unidade do anel $(\psi_*(\sh{O}_X))_y$ não é
nulo, sendo o germe em $y$ da seção
\[
  1\in\Gamma(X,\sh{O}_X)=\Gamma(Y,\psi_*(\sh{O}_X));
\]
como ele é a imagem pela $\theta$ do elemento unidade de $\sh{O}_y$, este último
não pertence a $\sh{J}_y$, logo $\sh{O}_y/\sh{J}_y\neq 0$; isso conclui a
demonstração.
\end{proof}

\begin{proposition}[9.5.5]
\label{I.9.5.5}
\emph{(Transitividade das imagens fechadas).} Sejam $f:X\to Y$ e $g:Y\to Z$ dois
morfismos de pré-esquemas; supomos que a imagem fechada $Y'$ de $X$ por $f$ exista e que,
se $g'$ é a restrição de $g$ a $Y'$, a imagem fechada $Z'$ de $Y'$ por $g'$ exista.
Então a imagem fechada de $X$ por $g\circ f$ existe e é igual a $Z'$.
\end{proposition}

\begin{proof}
Basta \sref{I.9.5.1} mostrar que $Z'$ é o subpré-esquema fechado mínimo
$Z_1$ de $Z$ tal que o subpré-esquema fechado $(g\circ f)^{-1}(Z_1)$ de $X$
(igual a $f^{-1}(g^{-1}(Z_1))$ pelo Corolário~\sref{I.4.4.2}) seja igual a
$X$; equivale a dizer que $Z'$ é o subpré-esquema fechado mínimo de $Z$ tal que
$f$ se fatora pela injeção $g^{-1}(Z_1)\to Y$ \sref{I.4.4.1}. Em virtude da
existência da imagem fechada $Y'$, todo $Z_1$ com essa propriedade é tal que
$g^{-1}(Z_1)$ se fatora por $Y'$, o que equivale a dizer que
$j^{-1}(g^{-1}(Z_1))=g'^{-1}(Z_1)=Y'$, denotando por $j$ a injeção
$Y'\to Y$. Pela definição de $Z'$, concluímos que $Z'$ é o subpré-esquema fechado
mínimo de $Z$ que satisfaz a condição precedente.
\end{proof}

\begin{corollary}[9.5.6]
\label{I.9.5.6}
Seja $f:X\to Y$ um morfismo sobre $S$ tal que $Y$ seja a imagem fechada de $X$ pelo morfismo $f$.
Seja $Z$ um esquema sobre $S$; se dois morfismos sobre $S$ $g_1$, $g_2$ de $Y$ em $Z$
são tais que $g_1\circ f=g_2\circ f$, então $g_1=g_2$.
\end{corollary}

\begin{proof}
Seja $h=(g_1,g_2)_S:Y\to Z\times_S Z$; como a diagonal $T=\Delta_Z(Z)$ é um
subpré-esquema fechado de $Z\times_S Z$, $Y'=h^{-1}(T)$ é um subpré-esquema fechado de $Y$
\sref{I.4.4.1}. Seja $u=g_1\circ f=g_2\circ f$; então, por definição do produto,
$h'=h\circ f=(u,u)_S$, logo $h\circ f=\Delta_Z\circ u$; como
$\Delta_Z^{-1}(T)=Z$, temos $h'^{-1}(T)=u^{-1}(Z)=X$, portanto $f^{-1}(Y')=X$.
Concluímos daí \sref{I.4.4.1} que $f$ se fatora pela injeção canônica $Y'\to Y$,
portanto $Y'=Y$ por hipótese; segue então \sref{I.4.4.1} que $h$ se fatora como
$\Delta_Z\circ v$, onde $v$ é um morfismo $Y\to Z$, o que implica $g_1=g_2=v$.
\end{proof}

\begin{remark}[9.5.7]
\label{I.9.5.7}
Se $X$ e $Y$ são $S$-esquemas, a proposição~\sref{I.9.5.6} implica que,
\oldpage[I]{178}
quando $Y$ é a imagem fechada de $X$ pelo morfismo $f$, $f$ é um \emph{epimorfismo} na categoria de \emph{$S$-esquemas} (T,~1.1). Mostraremos no Capítulo~V que, reciprocamente, se a imagem fechada $Y'$ de $X$ pelo morfismo $f$ existe e se $f$ é um epimorfismo de $S$-esquemas, então necessariamente $Y'=Y$.
\end{remark}

\begin{proposition}[9.5.8]
\label{I.9.5.8}
Suponhamos satisfeitas as hipóteses de \sref{I.9.5.1}, e seja $Y'$ a imagem fechada de $X$ pelo morfismo $f$. Para todo aberto $V$ de $Y$, seja
$f_V:f^{-1}(V)\to V$ a restrição de $f$; então a imagem fechada de
$f^{-1}(V)$ pelo morfismo $f_V$ em $V$ existe e é igual ao pré-esquema induzido por
$Y'$ no aberto $V\cap Y'$ de $Y'$ \emph{(em outras palavras, ao subpré-esquema
$\inf(V,Y')$ de $Y$ \sref{I.4.4.3})}.
\end{proposition}

\begin{proof}
Tomemos $X'=f^{-1}(V)$; como a imagem direta de $\sh{O}_{X'}$ por
$f_V$ não é outra senão a restrição de $f_*(\sh{O}_X)$ a $V$, é claro
que o núcleo $\sh{J}'$ do homomorfismo
$\sh{O}_V\to(f_V)_*(\sh{O}_{X'})$ é a restrição de
$\sh{J}$ a $V$, de onde segue imediatamente a proposição.
\end{proof}

Veremos que este resultado pode ser entendido dizendo que a formação da imagem fechada comuta com uma extensão $Y_1\to Y$ do pré-esquema de base, que é uma \emph{imersão aberta}.
Veremos no Capítulo~IV que o mesmo ocorre para uma extensão $Y_1\to Y$ que é um morfismo \emph{plano}, desde que $f$ seja separado e quase-compacto.

\begin{proposition}[9.5.9]
\label{I.9.5.9}
Seja $f:X\to Y$ um morfismo tal que a imagem fechada $Y'$ de $X$ pelo morfismo $f$ exista.
\begin{enumerate}
  \item[{\rm(i)}] Se $X$ é reduzido, então $Y'$ também o é.
  \item[{\rm(ii)}] Se as hipóteses da Proposição~\sref{I.9.5.1} são satisfeitas e $X$ é irredutível (resp. integral), então $Y'$ também o é.
\end{enumerate}
\end{proposition}

\begin{proof}
Por hipótese, o morfismo $f$ se fatora como $X\xrightarrow{g}Y'\xrightarrow{j}Y$, onde $j$ é a injeção canônica.
Como $X$ é reduzido, $g$ se fatora como $X\xrightarrow{h}Y'_{\mathrm{red}}\xrightarrow{j'}Y'$, onde $j'$ é a injeção canônica \sref{I.5.2.2}, e segue então da definição de $Y'$ que $Y'_{\mathrm{red}}=Y'$.
Se, além disso, as condições da Proposição~\sref{I.9.5.1} são satisfeitas, então segue de \sref{I.9.5.4} que $f(X)$ é denso em $Y'$; se $X$ é irredutível, então $Y'$ também o é \sref[0]{0.2.1.5}.
A afirmação relativa aos pré-esquemas integrais segue da conjunção das duas outras.
\end{proof}

\begin{proposition}[9.5.10]
\label{I.9.5.10}
Seja $Y$ um subpré-esquema de um pré-esquema $X$, tal que a injeção canônica $i:Y\to X$ seja um morfismo quase-compacto.
Existe então um menor subpré-esquema fechado $\overline{Y}$ de $X$ contendo $Y$; seu espaço subjacente é a aderência daquele de $Y$; este último é aberto em sua aderência, e o pré-esquema $Y$ é induzido nesse aberto por $\overline{Y}$.
\end{proposition}

\begin{proof}
Basta aplicar a Proposição~\sref{I.9.5.1} à injeção $i$\footnote{\emph{[Trad.] A fonte francesa imprime $j$ aqui, mas a proposição define a injeção canônica como $i:Y\to X$.}}, que é separada \sref{I.5.5.1} e quase-compacta por hipótese; \sref{I.9.5.1} prova, portanto, a existência de $\overline{Y}$, e \sref{I.9.5.4} mostra que seu espaço subjacente é a aderência de $Y$ em $X$; como $Y$ é localmente fechado em $X$, ele é aberto em $\overline{Y}$, e a última afirmação provém da aplicação de \sref{I.9.5.8} a um aberto $V$ de $X$ tal que $Y$ seja fechado em $V$.
\end{proof}

Com essas notações, se a injeção $V\to X$ é quase-compacta, e se $\sh{J}$ é o feixe quase-coerente de ideais de $\sh{O}_X|V$ que define o subpré-esquema fechado $Y$ de $V$, segue-se da Proposição~\sref{I.9.5.1} que o feixe quase-coerente de ideais de $\sh{O}_X$ que define $\overline{Y}$ é o prolongamento canônico \sref{I.9.4.1} $\overline{\sh{J}}$ de $\sh{J}$, pois é claramente o maior subfeixe quase-coerente de ideais de $\sh{O}_X$ que induz $\sh{J}$ em $V$.

\oldpage[I]{179}
\begin{corollary}[9.5.11]
\label{I.9.5.11}
Sob as hipóteses da Proposição~\sref{I.9.5.10}, toda seção de
$\sh{O}_{\overline{Y}}$ sobre um aberto $V$ de $\overline{Y}$ que é nula em
$V\cap Y$ é nula.
\end{corollary}

\begin{proof}
Pela Proposição~\sref{I.9.5.8}, podemos nos reduzir ao caso $V=\overline{Y}$.
Levando em conta que as seções de $\sh{O}_{\overline{Y}}$ sobre $\overline{Y}$
correspondem canonicamente às seções sobre $\overline{Y}$ de
$\overline{Y}\otimes_{\mathbb Z}\mathbb Z[T]$ \sref{I.3.3.15} e que este último
é separado sobre $\overline{Y}$, o corolário aparece como um caso particular de
\sref{I.9.5.6}.
\end{proof}

Quando existe um menor subpré-esquema fechado $Y'$ de $X$ contendo um
subpré-esquema $Y$ de $X$, dizemos que $Y'$ é a \emph{aderência} de $Y$ em
$X$, quando isso não causar confusão.

\subsection{Feixes quase-coerentes de álgebras; mudança do feixe estrutural}
\label{subsection:I.9.6}
\label{I.9.6}

\begin{proposition}[9.6.1]
\label{I.9.6.1}
Seja $X$ um pré-esquema e $\sh{B}$ uma $\sh{O}_X$-álgebra quase-coerente \sref[0]{0.5.1.3}.
Para que um $\sh{B}$-módulo $\sh{F}$ seja quase-coerente (no espaço anelado $(X,\sh{B})$), é necessário e suficiente que $\sh{F}$ seja um $\sh{O}_X$-módulo quase-coerente.
\end{proposition}

\begin{proof}
Como a questão é local, podemos supor $X$ afim, dado por um anel $A$, e então $\sh{B}=\widetilde{B}$, onde $B$ é uma $A$-álgebra \sref{I.1.4.3}.
Se $\sh{F}$ é quase-coerente no espaço anelado $(X,\sh{B})$, podemos também supor que $\sh{F}$ é o conúcleo de um $\sh{B}$-homomorfismo $\sh{B}^{(I)}\to\sh{B}^{(J)}$; como esse homomorfismo também é um $\sh{O}_X$-homomorfismo de $\sh{O}_X$-módulos, e $\sh{B}^{(I)}$ e $\sh{B}^{(J)}$ são $\sh{O}_X$-módulos quase-coerentes \sref{I.1.3.9}[ii], $\sh{F}$ também é um $\sh{O}_X$-módulo quase-coerente \sref{I.1.3.9}[i].

Inversamente, se $\sh{F}$ é um $\sh{O}_X$-módulo quase-coerente, temos $\sh{F}=\widetilde{M}$, onde $M$ é um $B$-módulo \sref{I.1.4.3}; $M$ é isomorfo ao conúcleo de um $B$-homomorfismo $B^{(I)}\to B^{(J)}$, de modo que $\sh{F}$ é um $\sh{B}$-módulo isomorfo ao conúcleo do homomorfismo correspondente $\sh{B}^{(I)}\to\sh{B}^{(J)}$ \sref{I.1.3.13}, o que conclui a demonstração.
\end{proof}

Em particular, se $\sh{F}$ e $\sh{G}$ são dois $\sh{B}$-módulos quase-coerentes, então $\sh{F}\otimes_{\sh{B}}\sh{G}$ é um $\sh{B}$-módulo quase-coerente; o mesmo vale para $\shHom_{\sh{B}}(\sh{F},\sh{G})$ quando supomos, além disso, que $\sh{F}$ admite uma apresentação finita \sref{I.1.3.13}.

\begin{env}[9.6.2]
\label{I.9.6.2}
Dado um pré-esquema $X$, diremos que uma $\sh{O}_X$-álgebra quase-coerente $\sh{B}$ é de \emph{tipo finito} se, para todo $x\in X$, existe uma vizinhança aberta \emph{afim} $U$ de $x$ tal que $\Gamma(U,\sh{B})=B$ seja uma álgebra de tipo finito sobre $\Gamma(U,\sh{O}_X)=A$.
Temos então $\sh{B}|U=\widetilde{B}$ e, para todo $f\in A$, a $(\sh{O}_X|D(f))$-álgebra induzida $\sh{B}|D(f)$ é de tipo finito, pois é isomorfa a $(B_f)^{\sim}$, e $B_f=B\otimes_A A_f$ é claramente uma álgebra de tipo finito sobre $A_f$.
Como os $D(f)$ formam uma base da topologia de $U$, concluímos que, se $\sh{B}$ é uma $\sh{O}_X$-álgebra quase-coerente de tipo finito, então, para todo aberto $V$ de $X$, $\sh{B}|V$ é uma $(\sh{O}_X|V)$-álgebra quase-coerente de tipo finito.
\end{env}

\begin{proposition}[9.6.3]
\label{I.9.6.3}
Seja $X$ um pré-esquema localmente noetheriano.
Então toda $\sh{O}_X$-álgebra quase-coerente $\sh{B}$ de tipo finito é um feixe coerente de anéis \sref[0]{0.5.3.7}.
\end{proposition}

\begin{proof}
Podemos novamente nos restringir ao caso em que $X$ é um esquema afim dado por um anel noetheriano $A$, e em que $\sh{B}=\widetilde{B}$, sendo $B$ uma $A$-álgebra de tipo finito; $B$ é então um anel noetheriano.
Com isso, resta provar que o núcleo $\sh{N}$ de um $\sh{B}$-homomorfismo $\sh{B}^m\to\sh{B}$ é um $\sh{B}$-
\oldpage[I]{180}
módulo de tipo finito; mas ele é isomorfo (como $\sh{B}$-módulo) a $\widetilde{N}$, onde $N$ é o núcleo do homomorfismo correspondente de $B$-módulos $B^m\to B$ \sref{I.1.3.13}.
Como $B$ é noetheriano, o submódulo $N$ de $B^m$ é um $B$-módulo de tipo finito, portanto existe um homomorfismo $B^p\to B^m$ cuja imagem é $N$; como a sequência $B^p\to B^m\to B$ é exata, também é exata a sequência correspondente $\sh{B}^p\to\sh{B}^m\to\sh{B}$ \sref{I.1.3.5}, e como $\sh{N}$ é a imagem de $\sh{B}^p\to\sh{B}^m$ \sref{I.1.3.9}[i], isso demonstra a proposição.
\end{proof}

\begin{corollary}[9.6.4]
\label{I.9.6.4}
Sob as hipóteses de \sref{I.9.6.3}, para que um $\sh{B}$-módulo $\sh{F}$ seja coerente, é necessário e suficiente que seja um $\sh{O}_X$-módulo quase-coerente e um $\sh{B}$-módulo de tipo finito.
Se esse for o caso, e se $\sh{G}$ for um $\sh{B}$-submódulo ou um $\sh{B}$-módulo quociente de $\sh{F}$, então, para que $\sh{G}$ seja um $\sh{B}$-módulo coerente, é necessário e suficiente que seja um $\sh{O}_X$-módulo quase-coerente.
\end{corollary}

\begin{proof}
Levando em conta \sref{I.9.6.1}, as condições sobre $\sh{F}$ são claramente necessárias; mostraremos que são suficientes.
Podemos nos restringir ao caso em que $X$ é afim, dado por um anel noetheriano $A$, onde $\sh{B}=\widetilde{B}$, com $B$ uma $A$-álgebra de tipo finito, onde $\sh{F}=\widetilde{M}$, com $M$ um $B$-módulo, e onde existe um $\sh{B}$-homomorfismo sobrejetivo $\sh{B}^m\to\sh{F}\to0$.
Temos então a sequência exata correspondente $B^m\to M\to0$, portanto $M$ é um $B$-módulo de tipo finito; além disso, o núcleo $P$ do homomorfismo $B^m\to M$ é então um $B$-módulo de tipo finito, pois $B$ é noetheriano.
Concluímos de \sref{I.1.3.13} que $\sh{F}$ é o conúcleo de um $\sh{B}$-homomorfismo $\sh{B}^n\to\sh{B}^m$, e portanto é coerente, pois $\sh{B}$ é um feixe coerente de anéis \sref[0]{0.5.3.4}.
O mesmo raciocínio mostra que um $\sh{B}$-submódulo quase-coerente (resp. um $\sh{B}$-módulo quociente) de $\sh{F}$ é de tipo finito, de onde decorre a segunda parte do corolário.
\end{proof}

\begin{proposition}[9.6.5]
\label{I.9.6.5}
Seja $X$ um esquema quase-compacto, ou um pré-esquema cujo espaço subjacente é noetheriano.
Para toda $\sh{O}_X$-álgebra quase-coerente $\sh{B}$ de tipo finito, existe um $\sh{O}_X$-submódulo quase-coerente $\sh{E}$ de $\sh{B}$, de tipo finito, tal que $\sh{E}$ gera \sref[0]{0.4.1.4} a $\sh{O}_X$-álgebra $\sh{B}$.
\end{proposition}

\begin{proof}
Com efeito, por hipótese, existe uma cobertura finita $(U_i)$ de $X$ constituída de abertos afins tal que $\Gamma(U_i,\sh{B})=B_i$ seja uma álgebra de tipo finito sobre $\Gamma(U_i,\sh{O}_X)=A_i$; seja $E_i$ um $A_i$-submódulo de tipo finito de $B_i$ que gere a $A_i$-álgebra $B_i$; graças a \sref{I.9.4.7}, existe um $\sh{O}_X$-submódulo $\sh{E}_i$ de $\sh{B}$, quase-coerente e de tipo finito, tal que $\sh{E}_i|U_i=\widetilde{E_i}$.
É claro que a soma $\sh{E}$ dos $\sh{E}_i$ é o objeto desejado.
\end{proof}

\begin{proposition}[9.6.6]
\label{I.9.6.6}
Seja $X$ um pré-esquema cujo espaço subjacente é localmente noetheriano, ou um esquema quase-compacto.
Então toda $\sh{O}_X$-álgebra quase-coerente $\sh{B}$ é o limite indutivo de suas $\sh{O}_X$-subálgebras quase-coerentes de tipo finito.
\end{proposition}

\begin{proof}
Com efeito, resulta de \sref{I.9.4.9} que $\sh{B}$ é o limite indutivo (como $\sh{O}_X$-módulo) de seus $\sh{O}_X$-submódulos quase-coerentes de tipo finito; estes últimos geram \emph{$\sh{O}_X$-subálgebras} quase-coerentes de $\sh{B}$ de tipo finito \sref{I.1.3.14}, e portanto $\sh{B}$ é \emph{a fortiori} o seu limite indutivo.
\end{proof}

\section{Esquemas formais}
\label{section:I.10}

\subsection{Esquemas formais afins}
\label{subsection:I.10.1}

\begin{env}[10.1.1]
\label{I.10.1.1}
Seja $A$ um anel topológico \emph{admissível} \sref[0]{0.7.1.2}; para cada ideal de definição $\mathfrak{J}$ de $A$, $\Spec(A/\mathfrak{J})$ pode ser identificado com o subespaço fechado $V(\mathfrak{J})$ de $\Spec(A)$ \sref{I.1.1.11}, o conjunto dos ideais primos \emph{abertos} de $A$; este espaço topológico não depende
\oldpage[I]{181}
do ideal de definição $\mathfrak{J}$ considerado; denotamos esse espaço topológico por $\mathfrak{X}$. Seja $(\mathfrak{J}_\lambda)$ um sistema fundamental de vizinhanças de $0$ em $A$, constituído de ideais de definição, e, para cada $\lambda$, seja $\sh{O}_\lambda$ o feixe estrutural de $\Spec(A/\mathfrak{J}_\lambda)$; esse feixe é induzido sobre $\mathfrak{X}$ por $\widetilde{A}/\widetilde{\mathfrak{J}_\lambda}$ (que é nulo fora de $\mathfrak{X}$).
Para $\mathfrak{J}_\mu\subset\mathfrak{J}_\lambda$, o homomorfismo canônico $A/\mathfrak{J}_\mu\to A/\mathfrak{J}_\lambda$ define, portanto, um homomorfismo $u_{\lambda\mu}:\sh{O}_\mu\to\sh{O}_\lambda$ de feixes de anéis \sref{I.1.6.1}, e $(\sh{O}_\lambda)$ é um \emph{sistema projetivo de feixes de anéis} para esses homomorfismos.
Como a topologia de $\mathfrak{X}$ admite uma base constituída de abertos quase-compactos, podemos associar a cada $\sh{O}_\lambda$ um \emph{feixe pseudodiscreto de anéis topológicos} \sref[0]{0.3.8.1} que tem $\sh{O}_\lambda$ como feixe de anéis subjacente (sem topologia), e que também denotamos por $\sh{O}_\lambda$; e os $\sh{O}_\lambda$ formam novamente um \emph{sistema projetivo de feixes de anéis topológicos} \sref[0]{0.3.8.2}.
Denotamos por $\sh{O}_\mathfrak{X}$ o \emph{feixe de anéis topológicos} sobre $\mathfrak{X}$, limite projetivo do sistema $(\sh{O}_\lambda)$; para todo aberto \emph{quase-compacto} $U$ de $\mathfrak{X}$, $\Gamma(U,\sh{O}_\mathfrak{X})$ é, portanto, o anel topológico limite projetivo do sistema de anéis \emph{discretos} $\Gamma(U,\sh{O}_\lambda)$ \sref[0]{0.3.2.6}.
\end{env}

\begin{definition}[10.1.2]
\label{I.10.1.2}
Dado um anel topológico \emph{admissível} $A$, chamamos espectro formal de $A$, denotado por $\Spf(A)$, o subespaço fechado $\mathfrak{X}$ de $\Spec(A)$ constituído pelos ideais primos \emph{abertos} de $A$.
Dizemos que um espaço anelado topológico é um \emph{esquema afim formal} se ele é isomorfo a um espectro formal $\Spf(A)=\mathfrak{X}$ munido do feixe de anéis topológicos $\sh{O}_\mathfrak{X}$, que é o limite projetivo de feixes de anéis topológicos pseudodiscretos $(\widetilde{A}/\widetilde{\mathfrak{J}_\lambda})|\mathfrak{X}$, onde $\mathfrak{J}_\lambda$ percorre o conjunto filtrante dos ideais de definição de $A$.
\end{definition}

Quando falarmos de um \emph{espectro formal $\mathfrak{X}=\Spf(A)$} como um esquema afim formal, tratar-se-á sempre do espaço anelado topológico $(\mathfrak{X},\sh{O}_\mathfrak{X})$, onde $\sh{O}_\mathfrak{X}$ é definido como acima.

Observamos que todo \emph{esquema afim} $X=\Spec(A)$ pode ser considerado como um esquema afim formal de uma única maneira, considerando $A$ como um anel topológico discreto: os anéis topológicos $\Gamma(U,\sh{O}_X)$ são então discretos sempre que $U$ for quase-compacto (mas não, em geral, quando $U$ for um aberto qualquer de $X$).

\begin{proposition}[10.1.3]
\label{I.10.1.3}
Se $\mathfrak{X}=\Spf(A)$, onde $A$ é um anel admissível, então $\Gamma(\mathfrak{X},\sh{O}_\mathfrak{X})$ é topologicamente isomorfo a $A$.
\end{proposition}

\begin{proof}
Com efeito, como $\mathfrak{X}$ é fechado em $\Spec(A)$, ele é quase-compacto e, portanto, $\Gamma(\mathfrak{X},\sh{O}_\mathfrak{X})$ é topologicamente isomorfo ao limite projetivo dos anéis discretos $\Gamma(\mathfrak{X},\sh{O}_\lambda)$; mas $\Gamma(\mathfrak{X},\sh{O}_\lambda)$ é isomorfo a $A/\mathfrak{J}_\lambda$ \sref{I.1.3.7}; como $A$ é separado e completo, ele é topologicamente isomorfo a $\varprojlim A/\mathfrak{J}_\lambda$ \sref[0]{0.7.2.1}, de onde segue a proposição.
\end{proof}

\begin{proposition}[10.1.4]
\label{I.10.1.4}
Seja $A$ um anel admissível, $\mathfrak{X}=\Spf(A)$ e, para todo $f\in A$, seja $\mathfrak{D}(f)=D(f)\cap\mathfrak{X}$; então o espaço anelado topológico $(\mathfrak{D}(f),\sh{O}_\mathfrak{X}|\mathfrak{D}(f))$ é isomorfo ao espectro formal afim $\Spf(A_{\{f\}})$ \sref[0]{0.7.6.15}.
\end{proposition}

\begin{proof}
Para todo ideal de definição $\mathfrak{J}$ de $A$, o anel discreto $S_f^{-1}A/ S_f^{-1}\mathfrak{J}$ é canonicamente identificado com $A_{\{f\}}/\mathfrak{J}_{\{f\}}$ \sref[0]{0.7.6.9}, de modo que, por \sref{I.1.2.5} e \sref{I.1.2.6}, o espaço topológico $\Spf(A_{\{f\}})$ é canonicamente identificado com $\mathfrak{D}(f)$.
Além disso, para todo aberto quase-compacto $U$ de $\mathfrak{X}$ contido em $\mathfrak{D}(f)$, $\Gamma(U,\sh{O}_\lambda)$ pode ser identificado com o módulo de seções do feixe estrutural de $\Spec(S_f^{-1}A/ S_f^{-1}\mathfrak{J}_\lambda)$ sobre $U$ \sref{I.1.3.6}; assim, tomando $\mathfrak{Y}=\Spf(A_{\{f\}})$, $\Gamma(U,\sh{O}_\mathfrak{X})$ pode ser identificado com o módulo de seções $\Gamma(U,\sh{O}_\mathfrak{Y})$, o que prova a proposição.
\end{proof}

\oldpage[I]{182}
\begin{env}[10.1.5]
\label{I.10.1.5}
Como feixe de anéis \emph{sem topologia}, o feixe estrutural $\sh{O}_\mathfrak{X}$ de $\Spf(A)$ admite, para todo $x\in\mathfrak{X}$, uma fibra que, por \sref{I.10.1.4}, pode ser identificada com o limite indutivo $\varinjlim_{f\notin\mathfrak{j}_x}A_{\{f\}}$. Portanto, por \sref[0]{0.7.6.17} e \sref[0]{0.7.6.18}:
\end{env}

\begin{proposition}[10.1.6]
\label{I.10.1.6}
Para todo $x\in\mathfrak{X}=\Spf(A)$, a fibra $\sh{O}_x$ é um anel local cujo corpo residual é isomorfo a $\kres(x)=A_x/\mathfrak{j}_xA_x$.
Se, além disso, $A$ é adic e noetheriano, $\sh{O}_x$ é um anel noetheriano.
\end{proposition}

Como $\kres(x)$ não se reduz ao anel zero, concluímos deste resultado que o \emph{suporte} do feixe de anéis $\sh{O}_\mathfrak{X}$ é \emph{igual a $\mathfrak{X}$}.

\subsection{Morfismos de esquemas formais afins}
\label{subsection:I.10.2}

\begin{env}[10.2.1]
\label{I.10.2.1}
Sejam $A$, $B$ dois anéis admissíveis, e seja $\vphi:B\to A$ um homomorfismo \emph{contínuo}.
A aplicação contínua ${}^a\vphi:\Spec(A)\to\Spec(B)$ \sref{I.1.2.1} leva então $\mathfrak{X}=\Spf(A)$ a $\mathfrak{Y}=\Spf(B)$, pois a imagem inversa por $\vphi$ de um ideal primo aberto de $A$ é um ideal primo aberto de $B$.
Por outro lado, para todo $g\in B$, $\vphi$ define um homomorfismo contínuo
\[
  \Gamma(\mathfrak{D}(g),\sh{O}_\mathfrak{Y})
  \longrightarrow
  \Gamma(\mathfrak{D}(\vphi(g)),\sh{O}_\mathfrak{X})
\]
em virtude de \sref{I.10.1.4}, \sref{I.10.1.3} e \sref[0]{0.7.6.7}; como esses homomorfismos satisfazem as condições de compatibilidade para as restrições correspondentes à passagem de $g$ a um múltiplo de $g$, e como $\mathfrak{D}(\vphi(g))={}^a\vphi^{-1}(\mathfrak{D}(g))$, eles definem um homomorfismo \emph{contínuo} de feixes de anéis topológicos $\sh{O}_\mathfrak{Y}\to{}^a\vphi_*(\sh{O}_\mathfrak{X})$ \sref[0]{0.3.2.5} que denotamos por $\widetilde{\vphi}$; assim definimos um morfismo $\Phi=({}^a\vphi,\widetilde{\vphi})$ de espaços anelados topológicos $\mathfrak{X}\to\mathfrak{Y}$.
Observemos que, como homomorfismo de feixes de anéis sem topologia, $\widetilde{\vphi}$ define um homomorfismo $\widetilde{\vphi}_x^\sharp:\sh{O}_{{}^a\vphi(x)}\to\sh{O}_x$ sobre as fibras, para todo $x\in\mathfrak{X}$.
\end{env}

\begin{proposition}[10.2.2]
\label{I.10.2.2}
Sejam $A$, $B$ dois anéis topológicos admissíveis, e sejam $\mathfrak{X}=\Spf(A)$, $\mathfrak{Y}=\Spf(B)$. Para que um morfismo $u=(\psi,\theta):\mathfrak{X}\to\mathfrak{Y}$ de espaços anelados topológicos seja da forma $({}^a\vphi,\widetilde{\vphi})$, onde $\vphi$ é um homomorfismo contínuo de anéis $B\to A$, é necessário e suficiente que, para todo $x\in\mathfrak{X}$, $\theta_x^\sharp$ seja um homomorfismo local $\sh{O}_{\psi(x)}\to\sh{O}_x$.
\end{proposition}

\begin{proof}
A condição é necessária: sejam $\mathfrak{p}=\mathfrak{j}_x\in\Spf(A)$ e $\mathfrak{q}=\vphi^{-1}(\mathfrak{j}_x)$; se $g\notin\mathfrak{q}$, então $\vphi(g)\notin\mathfrak{p}$, e é imediato que o homomorfismo $B_{\{g\}}\to A_{\{\vphi(g)\}}$ induzido por $\vphi$ \sref[0]{0.7.6.7} envia $\mathfrak{q}_{\{g\}}$ para um subconjunto de $\mathfrak{p}_{\{\vphi(g)\}}$; passando ao limite indutivo, vemos (levando em conta \sref{I.10.1.5} e \sref[0]{0.7.6.17}) que $\widetilde{\vphi}_x^\sharp$ é um homomorfismo local.

Reciprocamente, seja $(\psi,\theta)$ um morfismo que satisfaz a condição da proposição; em virtude de \sref{I.10.1.3}, $\theta$ define um homomorfismo contínuo de anéis
\[
  \vphi=\Gamma(\theta):B=\Gamma(\mathfrak{Y},\sh{O}_\mathfrak{Y})\to\Gamma(\mathfrak{X},\sh{O}_\mathfrak{X})=A.
\]
Em virtude da hipótese sobre $\theta$, para que a seção $\vphi(g)$ de $\sh{O}_\mathfrak{X}$ sobre $\mathfrak{X}$ tenha um germe invertível no ponto $x$, é necessário e suficiente que $g$ tenha um germe invertível no ponto $\psi(x)$. Mas, por \sref[0]{0.7.6.17}, as seções de $\sh{O}_\mathfrak{X}$ (resp. $\sh{O}_\mathfrak{Y}$) sobre $\mathfrak{X}$ (resp. $\mathfrak{Y}$) cujo germe não é invertível no ponto $x$ (resp. $\psi(x)$) são exatamente os elementos de $\mathfrak{j}_x$
\oldpage[I]{183}
(resp. $\mathfrak{j}_{\psi(x)}$); a observação anterior mostra, portanto, que ${}^a\vphi=\psi$. Finalmente, para todo $g\in B$, o diagrama
\[
  \xymatrix{
    B=\Gamma(\mathfrak{Y},\sh{O}_\mathfrak{Y})
      \ar[r]^\vphi\ar[d] &
    \Gamma(\mathfrak{X},\sh{O}_\mathfrak{X})=A\ar[d]\\
    B_{\{g\}}=\Gamma(\mathfrak{D}(g),\sh{O}_\mathfrak{Y})
      \ar[r]_{\Gamma(\theta_{\mathfrak{D}(g)})} &
    \Gamma(\mathfrak{D}(\vphi(g)),\sh{O}_\mathfrak{X})
      =A_{\{\vphi(g)\}}
  }
\]
é comutativo; pela propriedade universal dos anéis completos de frações \sref[0]{0.7.6.6}, $\theta_{\mathfrak{D}(g)}$ é igual a $\widetilde{\vphi}_{\mathfrak{D}(g)}$ para todo $g\in B$, e por \sref[0]{0.3.2.5} temos $\theta=\widetilde{\vphi}$.
\end{proof}

Diremos que um morfismo $(\psi,\theta)$ de espaços anelados topológicos que satisfaz a condição da Proposição~\sref{I.10.2.2} é um \emph{morfismo de esquemas formais afins}. Pode-se dizer que os functores $A\mapsto\Spf(A)$ e $\mathfrak{X}\mapsto\Gamma(\mathfrak{X},\sh{O}_\mathfrak{X})$ definem uma \emph{equivalência} entre a categoria dos anéis admissíveis e a categoria oposta dos esquemas formais afins (T, I, 1.2).

\begin{env}[10.2.3]
\label{I.10.2.3}
Como caso particular de \sref{I.10.2.2}, observemos que, para $f\in A$, a injeção canônica do esquema formal afim induzido por $\mathfrak{X}$ sobre $\mathfrak{D}(f)$ corresponde ao homomorfismo canônico contínuo $A\to A_{\{f\}}$.
Sob as hipóteses da Proposição~\sref{I.10.2.2}, sejam $h$ um elemento de $B$ e $g$ um elemento de $A$ múltiplo de $\vphi(h)$;
temos então $\psi(\mathfrak{D}(g))\subset\mathfrak{D}(h)$; a restrição de $u$ a $\mathfrak{D}(g)$, considerada como morfismo de $\mathfrak{D}(g)$ em $\mathfrak{D}(h)$, é o único morfismo $v$ que torna comutativo o diagrama
\[
  \xymatrix{
    \mathfrak{D}(g)\ar[r]^v\ar[d] &
    \mathfrak{D}(h)\ar[d]\\
    \mathfrak{X}\ar[r]_u &
    \mathfrak{Y}
  }
\]
Este morfismo corresponde ao único homomorfismo contínuo $\vphi':B_{\{h\}}\to A_{\{g\}}$ \sref[0]{0.7.6.7} que torna comutativo o diagrama
\[
  \xymatrix{
    A\ar[d] &
    B\ar[l]_\vphi\ar[d]\\
    A_{\{g\}} &
    B_{\{h\}}\ar[l]_{\vphi'}
  }
\]
\end{env}

\subsection{Ideais de definição de um esquema formal afim}
\label{subsection:I.10.3}

\begin{env}[10.3.1]
\label{I.10.3.1}
Sejam $A$ um anel admissível, $\mathfrak{J}$ um ideal aberto de $A$ e $\mathfrak{X}$ o esquema formal afim $\Spf(A)$.
Seja $(\mathfrak{J}_\lambda)$ o conjunto dos ideais de definição de $A$ contidos em $\mathfrak{J}$; então $\widetilde{\mathfrak{J}}/\widetilde{\mathfrak{J}}_\lambda$ é um feixe de ideais de $\widetilde{A}/\widetilde{\mathfrak{J}}_\lambda$.
Denotemos por $\mathfrak{J}^\Delta$ o limite projetivo dos feixes induzidos sobre $\mathfrak{X}$ por $\widetilde{\mathfrak{J}}/\widetilde{\mathfrak{J}}_\lambda$, que se identifica com um \emph{feixe de ideais} de $\sh{O}_\mathfrak{X}$ \sref[0]{0.3.2.6}.
Para todo $f\in A$, $\Gamma(\mathfrak{D}(f),\mathfrak{J}^\Delta)$ é o limite projetivo dos $S_f^{-1}\mathfrak{J}/S_f^{-1}\mathfrak{J}_\lambda$, isto é, identifica-se com o ideal aberto $\mathfrak{J}_{\{f\}}$ do anel $A_{\{f\}}$ \sref[0]{0.7.6.9}, e, em particular, $\Gamma(\mathfrak{X},\mathfrak{J}^\Delta)=\mathfrak{J}$; conclui-se (pois os $\mathfrak{D}(f)$ formam uma base da topologia de $\mathfrak{X}$) que
\[
\label{I.10.3.1.1}
  \mathfrak{J}^\Delta|\mathfrak{D}(f)=(\mathfrak{J}_{\{f\}})^\Delta.
  \tag{10.3.1.1}
\]
\end{env}

\oldpage[I]{184}
\begin{env}[10.3.2]
\label{I.10.3.2}
Com as notações de \sref{I.10.3.1}, para todo $f\in A$, a aplicação canônica de $A_{\{f\}}=\Gamma(\mathfrak{D}(f),\sh{O}_\mathfrak{X})$ em $\Gamma(\mathfrak{D}(f),(\widetilde{A}/\widetilde{\mathfrak{J}})|\mathfrak{X})=S_f^{-1}A/S_f^{-1}\mathfrak{J}$ é \emph{sobrejetiva} e tem por núcleo $\Gamma(\mathfrak{D}(f),\mathfrak{J}^\Delta)=\mathfrak{J}_{\{f\}}$ \sref[0]{0.7.6.9}; essas aplicações definem, portanto, um homomorfismo contínuo \emph{sobrejetivo}, dito \emph{canônico}, do feixe de anéis topológicos $\sh{O}_\mathfrak{X}$ sobre o feixe de anéis discretos $(\widetilde{A}/\widetilde{\mathfrak{J}})|\mathfrak{X}$, cujo núcleo é $\mathfrak{J}^\Delta$; esse homomorfismo não é outro senão $\widetilde{\vphi}$ \sref{I.10.2.1}, onde $\vphi$ é o homomorfismo contínuo $A\to A/\mathfrak{J}$; o morfismo $({}^a\vphi,\widetilde{\vphi}):\Spec(A/\mathfrak{J})\to\mathfrak{X}$ de esquemas formais afins (onde ${}^a\vphi$ é o homeomorfismo identidade de $\mathfrak{X}$ sobre si mesmo) também é dito \emph{canônico}.
Temos, portanto, segundo o que precede, um \emph{isomorfismo canônico}
\[
\label{I.10.3.2.1}
  \sh{O}_\mathfrak{X}/\mathfrak{J}^\Delta\isoto(\widetilde{A}/\widetilde{\mathfrak{J}})|\mathfrak{X}.
  \tag{10.3.2.1}
\]

É claro (em virtude de $\Gamma(\mathfrak{X},\mathfrak{J}^\Delta)=\mathfrak{J}$) que a aplicação $\mathfrak{J}\to\mathfrak{J}^\Delta$ é \emph{estritamente crescente}; segundo o que precede, para $\mathfrak{J}\subset\mathfrak{J}'$, o feixe ${\mathfrak{J}'}^\Delta/\mathfrak{J}^\Delta$ é canonicamente isomorfo a $\widetilde{\mathfrak{J}'}/\widetilde{\mathfrak{J}}=(\mathfrak{J}'/\mathfrak{J})\supertilde$.
\end{env}

\begin{env}[10.3.3]
\label{I.10.3.3}
Permanecendo as hipóteses e notações de \sref{I.10.3.1}, diremos que um feixe de ideais $\sh{J}$ de $\sh{O}_\mathfrak{X}$ é um \emph{feixe de ideais de definição de $\mathfrak{X}$} (ou um \emph{ideal de definição de $\mathfrak{X}$}) se, para todo $x\in\mathfrak{X}$, existe uma vizinhança aberta de $x$ da forma $\mathfrak{D}(f)$, onde $f\in A$, tal que $\sh{J}|\mathfrak{D}(f)$ seja da forma $\mathfrak{H}^\Delta$, onde $\mathfrak{H}$ é um ideal de definição de $A_{\{f\}}$.
\end{env}

\begin{proposition}[10.3.4]
\label{I.10.3.4}
Para todo $f\in A$, todo ideal de definição de $\mathfrak{X}$ induz um ideal de definição de $\mathfrak{D}(f)$.
\end{proposition}

\begin{proof}
Isto resulta de \sref{I.10.3.1.1}.
\end{proof}

\begin{proposition}[10.3.5]
\label{I.10.3.5}
Se $A$ é um anel admissível, todo ideal de definição de $\mathfrak{X}=\Spf(A)$ é da forma $\mathfrak{J}^\Delta$, onde $\mathfrak{J}$ é um ideal de definição de $A$ univocamente determinado.
\end{proposition}

\begin{proof}
Com efeito, seja $\sh{J}$ um ideal de definição de $\mathfrak{X}$; por hipótese, e como $\mathfrak{X}$ é quase-compacto, existem finitos elementos $f_i\in A$ tais que os $\mathfrak{D}(f_i)$ recobrem $\mathfrak{X}$ e que $\sh{J}|\mathfrak{D}(f_i)=\mathfrak{H}_i^\Delta$, onde $\mathfrak{H}_i$ é um ideal de definição de $A_{\{f_i\}}$.
Para todo $i$, existe, portanto, um ideal aberto $\mathfrak{K}_i$ de $A$ tal que $(\mathfrak{K}_i)_{\{f_i\}}=\mathfrak{H}_i$ \sref[0]{0.7.6.9}; seja $\mathfrak{K}$ um ideal de definição de $A$ contido em todos os $\mathfrak{K}_i$.
A imagem canônica de $\sh{J}/\mathfrak{K}^\Delta$ no feixe estrutural $(A/\mathfrak{K})\supertilde$ de $\Spec(A/\mathfrak{K})$ \sref{I.10.3.2} é, portanto, tal que sua restrição a $\mathfrak{D}(f_i)$ é igual à de $(\mathfrak{K}_i/\mathfrak{K})\supertilde$;
conclui-se que essa imagem canônica é um feixe \emph{quase-coerente} sobre $\Spec(A/\mathfrak{K})$, e portanto da forma $(\mathfrak{J}/\mathfrak{K})\supertilde$, onde $\mathfrak{J}$ é um ideal de $A$ que contém $\mathfrak{K}$ \sref{I.1.4.1}, donde $\sh{J}=\mathfrak{J}^\Delta$ \sref{I.10.3.2};
além disso, como para todo $i$ existe um inteiro $n_i$ tal que $\mathfrak{H}_i^{n_i}\subset\mathfrak{K}_{\{f_i\}}$, tomando $n$ como o maior dos $n_i$, teremos $(\sh{J}/\mathfrak{K}^\Delta)^n=0$, e, por conseguinte \sref{I.10.3.2}, $((\mathfrak{J}/\mathfrak{K})\supertilde)^n=0$, donde finalmente $(\mathfrak{J}/\mathfrak{K})^n=0$ \sref{I.1.3.13}, o que prova que $\mathfrak{J}$ é um ideal de definição de $A$ \sref[0]{0.7.1.4}.
\end{proof}

\begin{proposition}[10.3.6]
\label{I.10.3.6}
Sejam $A$ um anel ádico e $\mathfrak{J}$ um ideal de definição de $A$ tal que $\mathfrak{J}/\mathfrak{J}^2$ seja um $(A/\mathfrak{J})$-módulo de tipo finito. Para todo inteiro $n>0$, tem-se então $(\mathfrak{J}^\Delta)^n=(\mathfrak{J}^n)^\Delta$.
\end{proposition}

\begin{proof}
Com efeito, para todo $f\in A$, tem-se (pois $\mathfrak{J}^n$ é um ideal aberto)
\[
  (\Gamma(\mathfrak{D}(f),\mathfrak{J}^\Delta))^n=(\mathfrak{J}_{\{f\}})^n=(\mathfrak{J}^n)_{\{f\}}=\Gamma(\mathfrak{D}(f^n),(\mathfrak{J}^n)^\Delta)
\]
\oldpage[I]{185}
em virtude de \sref{I.10.3.1.1} e \sref[0]{0.7.6.12}.
Como $(\mathfrak{J}^\Delta)^n$ está associado ao pré-feixe $U\mapsto(\Gamma(U,\mathfrak{J}^\Delta))^n$ \sref[0]{0.4.1.6}, a proposição resulta\footnote{\emph{[Trad.] A fonte francesa impressa diz aqui «le corollaire», embora 10.3.6 esteja enunciado como proposição.}} do fato de que os $\mathfrak{D}(f)$ formam uma base da topologia de $\mathfrak{X}$.
\end{proof}

\begin{env}[10.3.7]
\label{I.10.3.7}
Diz-se que uma família $(\sh{J}_\lambda)$ de ideais de definição de $\mathfrak{X}$ é um \emph{sistema fundamental de ideais de definição} se todo ideal de definição de $\mathfrak{X}$ contém um dos $\sh{J}_\lambda$; como $\sh{J}_\lambda=\mathfrak{J}_\lambda^\Delta$, isso equivale a dizer que os $\mathfrak{J}_\lambda$ formam um \emph{sistema fundamental de vizinhanças de $0$} em $A$.
Seja $(f_\alpha)$ uma família de elementos de $A$ tal que os $\mathfrak{D}(f_\alpha)$ recobrem $\mathfrak{X}$.
Se $(\sh{J}_\lambda)$ é uma família filtrante decrescente de feixes de ideais de $\sh{O}_\mathfrak{X}$ tal que, para todo $\alpha$, a família $(\sh{J}_\lambda|\mathfrak{D}(f_\alpha))$ é um sistema fundamental de ideais de definição de $\mathfrak{D}(f_\alpha)$, então $(\sh{J}_\lambda)$ é um sistema fundamental de ideais de definição de $\mathfrak{X}$.
Com efeito, para todo ideal de definição $\sh{J}$ de $\mathfrak{X}$, existe um recobrimento finito de $\mathfrak{X}$ por $\mathfrak{D}(f_i)$ tal que, para todo $i$, $\sh{J}_{\lambda_i}|\mathfrak{D}(f_i)$ é um ideal de definição de $\mathfrak{D}(f_i)$ contido em $\sh{J}|\mathfrak{D}(f_i)$.
Se $\mu$ é um índice tal que $\sh{J}_\mu\subset\sh{J}_{\lambda_i}$ para todo $i$, segue de \sref{I.10.3.3} que $\sh{J}_\mu$ é um ideal de definição de $\mathfrak{X}$, evidentemente contido em $\sh{J}$, donde a afirmação.
\end{env}

\subsection{Pré-esquemas formais e morfismos de pré-esquemas formais}
\label{subsection:I.10.4}

\begin{env}[10.4.1]
\label{I.10.4.1}
Dado um espaço anelado topológico $\mathfrak{X}$, diz-se que um aberto $U\subset\mathfrak{X}$ é um \emph{aberto formal afim} (resp.~um \emph{aberto formal afim ádico}, resp.~um \emph{aberto formal afim noetheriano}) se o espaço anelado topológico induzido por $\mathfrak{X}$ sobre $U$ é um esquema formal afim (resp.~um esquema desse tipo cujo anel é ádico, resp.~ádico e noetheriano).
\end{env}

\begin{definition}[10.4.2]
\label{I.10.4.2}
Chama-se \emph{pré-esquema formal} um espaço anelado topológico $\mathfrak{X}$ no qual todo ponto admite uma vizinhança que é um aberto formal afim.
Diz-se que o pré-esquema formal $\mathfrak{X}$ é ádico (resp. localmente noetheriano) se todo ponto de $\mathfrak{X}$ admite uma vizinhança que é um aberto formal afim ádico (resp. noetheriano).
Diz-se que $\mathfrak{X}$ é noetheriano se é localmente noetheriano e se seu espaço subjacente é quase-compacto (e, portanto, noetheriano).
\end{definition}

\begin{proposition}[10.4.3]
\label{I.10.4.3}
Se $\mathfrak{X}$ é um pré-esquema formal (resp. um pré-esquema formal localmente noetheriano), os abertos formais afins (resp. afins noetherianos) formam uma base da topologia de $\mathfrak{X}$.
\end{proposition}

\begin{proof}
Isso decorre de \sref{I.10.4.2} e \sref{I.10.1.4}, levando em conta que, se $A$ é um anel ádico noetheriano, o mesmo vale para $A_{\{f\}}$ para todo $f\in A$ \sref[0]{0.7.6.11}.
\end{proof}

\begin{corollary}[10.4.4]
\label{I.10.4.4}
Se $\mathfrak{X}$ é um pré-esquema formal (resp. um pré-esquema formal localmente noetheriano, resp. noetheriano), o espaço anelado topológico induzido sobre todo aberto de $\mathfrak{X}$ continua sendo um pré-esquema formal (resp. um pré-esquema formal localmente noetheriano, resp. noetheriano).
\end{corollary}

\begin{definition}[10.4.5]
\label{I.10.4.5}
Dados dois pré-esquemas formais $\mathfrak{X}$ e $\mathfrak{Y}$, chama-se morfismo (de pré-esquemas formais) de $\mathfrak{X}$ em $\mathfrak{Y}$ todo morfismo $(\psi,\theta)$ de espaços anelados topológicos tal que, para todo $x\in\mathfrak{X}$, $\theta_x^\sharp$ seja um homomorfismo local $\sh{O}_{\psi(x)}\to\sh{O}_x$.
\end{definition}

É imediato que o composto de dois morfismos de pré-esquemas formais continua sendo um morfismo desse tipo; portanto, os pré-esquemas formais formam uma \emph{categoria}, e denota-se por $\Hom(\mathfrak{X},\mathfrak{Y})$ o conjunto dos morfismos de um pré-esquema formal $\mathfrak{X}$ em um pré-esquema formal $\mathfrak{Y}$.

\oldpage[I]{186}
Se $U$ é uma parte aberta de $\mathfrak{X}$, a injeção canônica em $\mathfrak{X}$ do pré-esquema formal induzido por $\mathfrak{X}$ sobre $U$ é um morfismo de pré-esquemas formais (e mesmo um \emph{monomorfismo} de espaços anelados topológicos \sref[0]{0.4.1.1}).

\begin{proposition}[10.4.6]
\label{I.10.4.6}
Sejam $\mathfrak{X}$ um pré-esquema formal e $\mathfrak{S}=\Spf(A)$ um esquema formal afim.
Existe uma correspondência biunívoca canônica entre os morfismos do pré-esquema formal $\mathfrak{X}$ no pré-esquema formal $\mathfrak{S}$ e os homomorfismos contínuos do anel $A$ no anel topológico $\Gamma(\mathfrak{X},\sh{O}_\mathfrak{X})$.
\end{proposition}

\begin{proof}
A demonstração é a mesma que a de \sref{I.2.2.4}, substituindo ``homomorfismo'' por ``homomorfismo contínuo'', ``aberto afim'' por ``aberto formal afim'', e utilizando \sref{I.10.2.2} em lugar de \sref{I.1.7.3}; deixam-se os detalhes ao leitor.
\end{proof}

\begin{env}[10.4.7]
\label{I.10.4.7}
Dados um pré-esquema formal $\mathfrak{S}$, um pré-esquema formal $\mathfrak{X}$ e um morfismo $\vphi:\mathfrak{X}\to\mathfrak{S}$, diz-se que esses dados definem um pré-esquema formal \emph{$\mathfrak{X}$ sobre $\mathfrak{S}$}, ou um \emph{$\mathfrak{S}$-pré-esquema formal}; $\vphi$ é chamado de \emph{morfismo estrutural} do $\mathfrak{S}$-pré-esquema $\mathfrak{X}$.
Se $\mathfrak{S}=\Spf(A)$, onde $A$ é um anel admissível, diz-se também que o $\mathfrak{S}$-pré-esquema formal $\mathfrak{X}$ é um \emph{$A$-pré-esquema formal}, ou um pré-esquema formal \emph{sobre $A$}.
Todo pré-esquema formal pode ser considerado como um pré-esquema formal sobre $\bb{Z}$ (munido da topologia discreta).

Se $\mathfrak{X}$ e $\mathfrak{Y}$ são dois $\mathfrak{S}$-pré-esquemas formais, diz-se que um morfismo $u:\mathfrak{X}\to\mathfrak{Y}$ é um \emph{$\mathfrak{S}$-morfismo} se o diagrama
\[
  \xymatrix{
    \mathfrak{X}\ar[rr]^u\ar[rd] & &
    \mathfrak{Y}\ar[ld]\\
    & \mathfrak{S}
  }
\]
(no qual as setas oblíquas são os morfismos estruturais) é comutativo.
Com essa definição, os $\mathfrak{S}$-pré-esquemas formais (para $\mathfrak{S}$ fixo) formam uma \emph{categoria}.
Denota-se por $\Hom_\mathfrak{S}(\mathfrak{X},\mathfrak{Y})$ o conjunto dos $\mathfrak{S}$-morfismos do $\mathfrak{S}$-pré-esquema formal $\mathfrak{X}$ no $\mathfrak{S}$-pré-esquema formal $\mathfrak{Y}$.
Quando $\mathfrak{S}=\Spf(A)$, diz-se também \emph{$A$-morfismo} em lugar de \emph{$\mathfrak{S}$-morfismo}.
\end{env}

\begin{env}[10.4.8]
\label{I.10.4.8}
Como todo esquema afim pode ser considerado como um esquema formal afim \sref{I.10.1.2}, todo pré-esquema (usual) pode ser considerado como um pré-esquema formal.
Além disso, de \sref{I.10.4.5} decorre que, para os pré-esquemas \emph{usuais}, os morfismos (resp. $S$-morfismos) de pré-esquemas \emph{formais} coincidem com os morfismos (resp. $S$-morfismos) definidos em \textsection\hyperref[section:I.2]{2}.
\end{env}

\subsection{Ideais de definição dos pré-esquemas formais}
\label{subsection:I.10.5}

\begin{env}[10.5.1]
\label{I.10.5.1}
Seja $\mathfrak{X}$ um pré-esquema formal; diz-se que um $\sh{O}_\mathfrak{X}$-ideal $\sh{J}$ é um \emph{feixe de ideais de definição} (ou um \emph{ideal de definição}) de $\mathfrak{X}$ se todo $x\in\mathfrak{X}$ possui uma vizinhança aberta formal afim $U$ tal que $\sh{J}|U$ seja um ideal de definição do esquema formal afim induzido por $\mathfrak{X}$ sobre $U$ \sref{I.10.3.3}; em virtude de \sref{I.10.3.1.1} e \sref{I.10.4.3}, para todo aberto $V\subset\mathfrak{X}$, $\sh{J}|V$ é então um ideal de definição do pré-esquema formal induzido por $\mathfrak{X}$ sobre $V$.

Diz-se que uma família $(\sh{J}_\lambda)$ de ideais de definição de $\mathfrak{X}$ é um \emph{sistema fundamental}
\oldpage[I]{187}
\emph{de ideais de definição} se existe um recobrimento $(U_\alpha)$ de $\mathfrak{X}$ por abertos formais afins tal que, para todo $\alpha$, a família dos $\sh{J}_\lambda|U_\alpha$ seja um sistema fundamental de ideais de definição \sref{I.10.3.7}\footnote{\emph{[Trad.] A fonte francesa impressa remete aqui a 10.3.6; a definição citada é 10.3.7.}} do esquema formal afim induzido por $\mathfrak{X}$ sobre $U_\alpha$.
Da observação final de \sref{I.10.3.7} decorre que, quando $\mathfrak{X}$ é um esquema formal afim, esta definição coincide com a dada em \sref{I.10.3.7}.
Para todo aberto $V$ de $\mathfrak{X}$, as restrições $\sh{J}_\lambda|V$ formam então um sistema fundamental de ideais de definição do pré-esquema formal induzido sobre $V$, em virtude de \sref{I.10.3.1.1}.
Se $\mathfrak{X}$ é um pré-esquema formal \emph{localmente noetheriano} e $\sh{J}$ é um ideal de definição de $\mathfrak{X}$, de \sref{I.10.3.6} decorre que as potências $\sh{J}^n$ formam um sistema fundamental de ideais de definição de $\mathfrak{X}$.
\end{env}

\begin{env}[10.5.2]
\label{I.10.5.2}
Sejam $\mathfrak{X}$ um pré-esquema formal e $\sh{J}$ um ideal de definição de $\mathfrak{X}$.
Então o espaço anelado $(\mathfrak{X},\sh{O}_\mathfrak{X}/\sh{J})$ é um \emph{pré-esquema} (usual), que é afim (resp. localmente noetheriano, resp. noetheriano) quando $\mathfrak{X}$ é um esquema formal afim (resp. um pré-esquema formal localmente noetheriano, resp. noetheriano); com efeito, reduzimo-nos imediatamente ao caso afim, e então a proposição já foi demonstrada em \sref{I.10.3.2}.
Além disso, se $\theta:\sh{O}_\mathfrak{X}\to\sh{O}_\mathfrak{X}/\sh{J}$ é o homomorfismo canônico, então $u=(1_\mathfrak{X},\theta)$ é um \emph{morfismo} (dito \emph{canônico}) de pré-esquemas formais $(\mathfrak{X},\sh{O}_\mathfrak{X}/\sh{J})\to(\mathfrak{X},\sh{O}_\mathfrak{X})$, pois, também aqui, isso foi visto no caso afim \sref{I.10.3.2}, ao qual se reduz imediatamente.
\end{env}

\begin{proposition}[10.5.3]
\label{I.10.5.3}
Sejam $\mathfrak{X}$ um pré-esquema formal e $(\sh{J}_\lambda)$ um sistema fundamental de ideais de definição de $\mathfrak{X}$.
Então o feixe de anéis topológicos $\sh{O}_\mathfrak{X}$ é o limite projetivo dos feixes pseudodiscretos de anéis \sref[0]{0.3.8.1} $\sh{O}_\mathfrak{X}/\sh{J}_\lambda$.
\end{proposition}

\begin{proof}
Como a topologia de $\mathfrak{X}$ admite uma base de abertos formais afins quase-compactos \sref{I.10.4.3}, reduzimo-nos ao caso afim, em que a proposição é consequência da Proposição \sref{I.10.3.5}, de \sref{I.10.3.2} e da Definição \sref{I.10.1.1}.
\end{proof}

Não é certo que todo pré-esquema formal admita ideais de definição.
Contudo:
\begin{proposition}[10.5.4]
\label{I.10.5.4}
Seja $\mathfrak{X}$ um pré-esquema formal localmente noetheriano.
Existe um maior ideal de definição $\sh{T}$ de $\mathfrak{X}$; é o único ideal de definição $\sh{J}$ tal que o pré-esquema $(\mathfrak{X},\sh{O}_\mathfrak{X}/\sh{J})$ seja reduzido.
Se $\sh{J}$ é um ideal de definição de $\mathfrak{X}$, então $\sh{T}$ é a imagem inversa, por $\sh{O}_\mathfrak{X}\to\sh{O}_\mathfrak{X}/\sh{J}$, do nilradical de $\sh{O}_\mathfrak{X}/\sh{J}$.
\end{proposition}

\begin{proof}
Suponhamos primeiro que $\mathfrak{X}=\Spf(A)$, onde $A$ é um anel ádico noetheriano.
A existência e as propriedades de $\sh{T}$ decorrem imediatamente de \sref{I.10.3.5}\footnote{\emph{[Trad.] A fonte francesa impressa remete a 10.3.4; a correção publicada em EGA~II remete a 10.3.5.}} e de \sref{I.5.1.1}, levando em conta a existência e as propriedades do maior ideal de definição de $A$ (\sref[0]{0.7.1.6} e \sref[0]{0.7.1.7}).

Para provar a existência e as propriedades de $\sh{T}$ no caso geral, basta mostrar que, se $U\supset V$ são dois abertos formais afins noetherianos de $\mathfrak{X}$, o maior ideal de definição $\sh{T}_U$ de $U$ induz o maior ideal de definição $\sh{T}_V$ de $V$; mas, como $(V,(\sh{O}_\mathfrak{X}|V)/(\sh{T}_U|V))$ é reduzido, isso decorre do que precede.
\end{proof}

Denota-se por $\mathfrak{X}_\text{red}$ o pré-esquema (usual) reduzido $(\mathfrak{X},\sh{O}_\mathfrak{X}/\sh{T}$).

\begin{corollary}[10.5.5]
\label{I.10.5.5}
Sejam $\mathfrak{X}$ um pré-esquema formal localmente noetheriano e $\sh{T}$ o maior ideal de definição de $\mathfrak{X}$; para todo aberto $V$ de $\mathfrak{X}$, $\sh{T}|V$ é o maior ideal de definição do pré-esquema formal induzido por $\mathfrak{X}$ sobre $V$.
\end{corollary}

\oldpage[I]{188}

\begin{proposition}[10.5.6]
\label{I.10.5.6}
Sejam $\mathfrak{X}$ e $\mathfrak{Y}$ dois pré-esquemas formais, $\sh{J}$ (resp. $\sh{K}$) um ideal de definição de $\mathfrak{X}$ (resp. $\mathfrak{Y}$), e $f:\mathfrak{X}\to\mathfrak{Y}$ um morfismo de pré-esquemas formais.
\begin{enumerate}
  \item[\rm{(i)}] Se $f^*(\sh{K})\sh{O}_\mathfrak{X}\subset\sh{J}$, existe um único morfismo $f':(\mathfrak{X},\sh{O}_\mathfrak{X}/\sh{J})\to(\mathfrak{Y},\sh{O}_\mathfrak{Y}/\sh{K})$ de pré-esquemas usuais que torna comutativo o diagrama
    \[
    \label{I.10.5.6.1}
      \xymatrix{
        (\mathfrak{X},\sh{O}_\mathfrak{X})\ar[r]^f &
        (\mathfrak{Y},\sh{O}_\mathfrak{Y})\\
        (\mathfrak{X},\sh{O}_\mathfrak{X}/\sh{J})\ar[r]_{f'}\ar[u] &
        (\mathfrak{Y},\sh{O}_\mathfrak{Y}/\sh{K})\ar[u]
      }
      \tag{10.5.6.1}
    \]
    em que as setas verticais são os morfismos canônicos.
  \item[\rm{(ii)}] Suponhamos que $\mathfrak{X}=\Spf(A)$ e $\mathfrak{Y}=\Spf(B)$ sejam esquemas formais afins, $\sh{J}=\mathfrak{J}^\Delta$ e $\sh{K}=\mathfrak{K}^\Delta$, onde $\mathfrak{J}$ (resp. $\mathfrak{K}$) é um ideal de definição de $A$ (resp. $B$), e $f=({}^a\vphi,\widetilde{\vphi})$, onde $\vphi:B\to A$ é um homomorfismo contínuo;
    para que se tenha $f^*(\sh{K})\sh{O}_\mathfrak{X}\subset\sh{J}$, é necessário e suficiente que $\vphi(\mathfrak{K})\subset\mathfrak{J}$, e então $f'$ é o morfismo $({}^a\vphi',\widetilde{\vphi'})$, onde $\vphi':B/\mathfrak{K}\to A/\mathfrak{J}$ é o homomorfismo deduzido de $\vphi$ pela passagem aos quocientes.
\end{enumerate}
\end{proposition}

\begin{proof}
\medskip\noindent
\begin{enumerate}
  \item[\rm{(i)}] Se $f=(\psi,\theta)$, a hipótese implica que a imagem por $\theta^\sharp:\psi^*(\sh{O}_\mathfrak{Y})\to\sh{O}_\mathfrak{X}$ do feixe de ideais $\psi^*(\sh{K})$ de $\psi^*(\sh{O}_\mathfrak{Y})$ está contida em $\sh{J}$ \sref[0]{0.4.3.5}.
    Passando aos quocientes, deduz-se portanto de $\theta^\sharp$ um homomorfismo de feixes de anéis
    \[
      \omega:\psi^*(\sh{O}_\mathfrak{Y}/\sh{K})=\psi^*(\sh{O}_\mathfrak{Y})/\psi^*(\sh{K})\to\sh{O}_\mathfrak{X}/\sh{J};
    \]
    além disso, como, para todo $x\in\mathfrak{X}$, $\theta_x^\sharp$ é um homomorfismo \emph{local}, o mesmo ocorre com $\omega_x$.
    O morfismo de espaços anelados $(\psi,\omega^\flat)$ é portanto \sref{I.2.2.1} o único morfismo $f'$ de espaços anelados que responde à questão.
  \item[\rm{(ii)}] A correspondência canônica e funtorial entre os morfismos de esquemas formais afins e os homomorfismos contínuos de anéis \sref{I.10.2.2} mostra que, no caso considerado, a relação $f^*(\sh{K})\sh{O}_\mathfrak{X}\subset\sh{J}$ implica que $f'=({}^a\vphi',\widetilde{\vphi'})$, onde $\vphi':B/\mathfrak{K}\to A/\mathfrak{J}$ é o único homomorfismo que torna comutativo o diagrama
    \[
    \label{I.10.5.6.2}
      \xymatrix{
        B\ar[r]^\vphi\ar[d] &
        A\ar[d]\\
        B/\mathfrak{K}\ar[r]_{\vphi'} &
        A/\mathfrak{J}
      }
      \tag{10.5.6.2}
    \]
    A existência de $\vphi'$ implica portanto que $\vphi(\mathfrak{K})\subset\mathfrak{J}$.
    Reciprocamente, se essa condição é satisfeita, denotando por $\vphi'$ o único homomorfismo que torna comutativo o diagrama \sref{I.10.5.6.2} e pondo $f'=({}^a\vphi',\widetilde{\vphi'})$, é claro que o diagrama \sref{I.10.5.6.1} é comutativo; a consideração dos homomorfismos ${}^a\vphi^*(\sh{O}_\mathfrak{Y})\to\sh{O}_\mathfrak{X}$ e ${}^a{\vphi'}^*(\sh{O}_\mathfrak{Y}/\sh{K})\to\sh{O}_\mathfrak{X}/\sh{J}$ correspondentes respectivamente a $f$ e $f'$ mostra então que isso implica a relação $f^*(\sh{K})\sh{O}_\mathfrak{X}\subset\sh{J}$.
\end{enumerate}
\end{proof}

É claro que a correspondência $f\mapsto f'$ definida acima é \emph{funtorial}.

\subsection{Pré-esquemas formais como limites indutivos de pré-esquemas}
\label{subsection:I.10.6}

\begin{env}[10.6.1]
\label{I.10.6.1}
Sejam $\mathfrak{X}$ um pré-esquema formal e $(\sh{J}_\lambda)$ um sistema fundamental de ideais de definição de $\mathfrak{X}$; para todo $\lambda$, seja $f_\lambda$ o morfismo canônico $(\mathfrak{X},\sh{O}_\mathfrak{X}/\sh{J}_\lambda)\to\mathfrak{X}$ \sref{I.10.5.2}; para $\sh{J}_\mu\subset\sh{J}_\lambda$, o homomorfismo canônico $\sh{O}_\mathfrak{X}/\sh{J}_\mu\to\sh{O}_\mathfrak{X}/\sh{J}_\lambda$ define um morfismo canônico
\oldpage[I]{189}
$f_{\mu\lambda}:(\mathfrak{X},\sh{O}_\mathfrak{X}/\sh{J}_\lambda)\to(\mathfrak{X},\sh{O}_\mathfrak{X}/\sh{J}_\mu)$ de pré-esquemas (usuais) tal que $f_\lambda=f_\mu\circ f_{\mu\lambda}$.
Os pré-esquemas $X_\lambda=(\mathfrak{X},\sh{O}_\mathfrak{X}/\sh{J}_\lambda)$ e os morfismos $f_{\mu\lambda}$ constituem portanto (em virtude de \sref{I.10.4.8}) um \emph{sistema indutivo} na categoria dos pré-esquemas formais.
\end{env}

\begin{proposition}[10.6.2]
\label{I.10.6.2}
Com as notações de \sref{I.10.6.1}, o pré-esquema formal $\mathfrak{X}$ e os morfismos $f_\lambda$ constituem um limite indutivo (T,~I,~1.8) do sistema $(X_\lambda,f_{\mu\lambda})$ na categoria dos pré-esquemas formais.
\end{proposition}

\begin{proof}
Seja $\mathfrak{Y}$ um pré-esquema formal e, para cada índice $\lambda$, seja
\[
  g_\lambda=(\psi_\lambda,\theta_\lambda):X_\lambda\to\mathfrak{Y}
\]
um morfismo tal que $g_\lambda=g_\mu\circ f_{\mu\lambda}$ para $\sh{J}_\mu\subset\sh{J}_\lambda$.
Essa última condição e a definição dos $X_\lambda$ implicam primeiro que todos os $\psi_\lambda$ são idênticos a uma mesma aplicação contínua $\psi:\mathfrak{X}\to\mathfrak{Y}$ dos espaços subjacentes; além disso, os homomorfismos $\theta_\lambda^\sharp:\psi^*(\sh{O}_\mathfrak{Y})\to\sh{O}_{X_\lambda}=\sh{O}_\mathfrak{X}/\sh{J}_\lambda$ formam um \emph{sistema projetivo} de homomorfismos de feixes de anéis.
Passando ao limite projetivo, deduz-se portanto um homomorfismo $\omega:\psi^*(\sh{O}_\mathfrak{Y})\to\varprojlim\sh{O}_\mathfrak{X}/\sh{J}_\lambda=\sh{O}_\mathfrak{X}$, e é claro que o morfismo $g=(\psi,\omega^\flat)$ de \emph{espaços anelados} é o \emph{único} morfismo que torna comutativos os diagramas
\[
\label{I.10.6.2.1}
  \xymatrix{
    X_\lambda\ar[rr]^{g_\lambda}\ar[rd]_{f_\lambda} & &
    \mathfrak{Y}\\
    & \mathfrak{X}\ar[ru]_g
  }
  \tag{10.6.2.1}
\]
Resta provar que $g$ é um morfismo de \emph{pré-esquemas formais}; como a questão é local sobre $\mathfrak{X}$ e $\mathfrak{Y}$, pode-se supor que $\mathfrak{X}=\Spf(A)$ e $\mathfrak{Y}=\Spf(B)$, sendo $A$ e $B$ anéis admissíveis, e que $\sh{J}_\lambda=\mathfrak{J}_\lambda^\Delta$, onde $(\mathfrak{J}_\lambda)$ é um sistema fundamental de ideais de definição de $A$ \sref{I.10.3.5}; como $A=\varprojlim A/\mathfrak{J}_\lambda$, a existência de um morfismo $g$ de esquemas formais afins que torna comutativos os diagramas \sref{I.10.6.2.1} decorre então da correspondência bijetiva \sref{I.10.2.2} entre morfismos de esquemas formais afins e homomorfismos contínuos de anéis, e da definição do limite projetivo.
Mas a unicidade de $g$ como morfismo de espaços anelados mostra que ele coincide com o morfismo denotado do mesmo modo no início da demonstração.
\end{proof}

A proposição seguinte estabelece, sob certas condições suplementares, a existência do limite indutivo de um sistema indutivo dado de pré-esquemas (usuais) na categoria dos pré-esquemas formais:
\begin{proposition}[10.6.3]
\label{I.10.6.3}
Sejam $\mathfrak{X}$ um espaço topológico e $(\sh{O}_i,u_{ji})$ um sistema projetivo de feixes de anéis sobre $\mathfrak{X}$, cujo conjunto de índices é $\bb{N}$.
Seja $\sh{J}_i$ o núcleo de $u_{0i}:\sh{O}_i\to\sh{O}_0$.
Suponha-se que:
\begin{enumerate}
  \item[\rm{(a)}] o espaço anelado $(\mathfrak{X},\sh{O}_i)$ é um pré-esquema $X_i$;
  \item[\rm{(b)}] para todo $x\in\mathfrak{X}$\footnote{\emph{[Trad.] A fonte francesa impressa diz $x\in X$; a correção publicada em EGA~II substitui $X$ por $\mathfrak{X}$.}} e todo $i$, existe uma vizinhança aberta $U_i$ de $x$ em $\mathfrak{X}$ tal que a restrição $\sh{J}_i|U_i$ é nilpotente; e
  \item[\rm{(c)}] os homomorfismos $u_{ji}$ são sobrejetivos.
\end{enumerate}

\oldpage[I]{190}
Seja $\sh{O}_\mathfrak{X}$ o feixe de anéis topológicos dado pelo limite projetivo dos feixes de anéis pseudodiscretos $\sh{O}_i$, e seja $u_i:\sh{O}_\mathfrak{X}\to\sh{O}_i$ o homomorfismo canônico.
Então o espaço anelado topológico $(\mathfrak{X},\sh{O}_\mathfrak{X})$ é um pré-esquema formal; os homomorfismos $u_i$ são sobrejetivos; seus núcleos $\sh{J}^{(i)}$ formam um sistema fundamental de feixes de ideais de definição de $\mathfrak{X}$, e $\sh{J}^{(0)}$ é o limite projetivo dos feixes de ideais $\sh{J}_i$.
\end{proposition}

\begin{proof}
Notemos primeiro que, em cada fibra, $u_{ji}$ é um homomorfismo sobrejetivo e \emph{a fortiori} um homomorfismo local; portanto, $v_{ij}=(1_\mathfrak{X},u_{ji})$ é um morfismo de pré-esquemas $X_j\to X_i$ ($i\geq j$) \sref{I.2.2.1}.
Suponhamos primeiro que cada $X_i$ seja um esquema afim de anel $A_i$.
Existe um homomorfismo \emph{de anéis} $\vphi_{ji}:A_i\to A_j$ tal que $u_{ji}=\widetilde{\vphi_{ji}}$ \sref{I.1.7.3}; por conseguinte \sref{I.1.6.3}, o feixe $\sh{O}_j$ é um $\sh{O}_i$-módulo quase-coerente sobre $X_i$ (para a lei externa definida por $u_{ji}$), associado a $A_j$ considerado como $A_i$-módulo por meio de $\vphi_{ji}$.
Para todo $f\in A_i$, seja $f'=\vphi_{ji}(f)$; por hipótese, os abertos $D(f)$ e $D(f')$ são idênticos em $\mathfrak{X}$, e o homomorfismo de $\Gamma(D(f),\sh{O}_i)=(A_i)_f$ em $\Gamma(D(f),\sh{O}_j)=(A_j)_{f'}$ correspondente a $u_{ji}$ não é outro senão $(\vphi_{ji})_f$ \sref{I.1.6.1}.
Mas, quando se considera $A_j$ como $A_i$-módulo, $(A_j)_{f'}$ é o $(A_i)_f$-módulo $(A_j)_f$, de modo que também se tem $u_{ji}=\widetilde{\vphi_{ji}}$, considerando desta vez $\vphi_{ji}$ como homomorfismo de \emph{$A_i$-módulos}.
Então, como $u_{ji}$ é sobrejetivo, conclui-se que $\vphi_{ji}$ também o é \sref{I.1.3.9}; e, se $\mathfrak{J}_{ji}$ é o núcleo de $\vphi_{ji}$, o núcleo de $u_{ji}$ é um $\sh{O}_i$-módulo quase-coerente igual a $\widetilde{\mathfrak{J}_{ji}}$.
Em particular, tem-se $\sh{J}_i=\widetilde{\mathfrak{J}_i}$, onde $\mathfrak{J}_i$ é o núcleo de $\vphi_{0i}:A_i\to A_0$.
A hipótese~(b) implica que $\sh{J}_i$ é \emph{nilpotente}: com efeito, como $\mathfrak{X}$ é quase-compacto, pode-se recobrir $\mathfrak{X}$ por um número finito de abertos $U_k$ tais que $(\sh{J}_i|U_k)^{n_k}=0$ e, tomando $n$ como o maior dos $n_k$, tem-se $\sh{J}_i^n=0$.
Conclui-se assim que $\mathfrak{J}_i$ é nilpotente \sref{I.1.3.13}.
Então o anel $A=\varprojlim A_i$ é admissível \sref[0]{0.7.2.2}, o homomorfismo canônico $\vphi_i:A\to A_i$ é sobrejetivo e seu núcleo $\mathfrak{J}^{(i)}$ é igual ao limite projetivo dos $\mathfrak{J}_{ik}$ para $k\geq i$; os $\mathfrak{J}^{(i)}$ formam um sistema fundamental de vizinhanças de $0$ em $A$.
As afirmações da Proposição~\sref{I.10.6.3} decorrem, neste caso, de \sref{I.10.1.1} e \sref{I.10.3.2}, pois $(\mathfrak{X},\sh{O}_\mathfrak{X})$ não é outro senão $\Spf(A)$.

Ainda nesse mesmo caso particular, notemos que, se $f=(f_i)$ é um elemento do limite projetivo $A=\varprojlim A_i$, todos os abertos $D(f_i)$ (abertos afins em $X_i$) se identificam com o aberto $\mathfrak{D}(f)$ de $\mathfrak{X}$, e o pré-esquema induzido por $X_i$ sobre $\mathfrak{D}(f)$ identifica-se, portanto, com o esquema afim $\Spec((A_i)_{f_i})$.

No caso geral, notemos primeiro que, para todo aberto quase-compacto $U$ de $\mathfrak{X}$, cada um dos $\sh{J}_i|U$ é nilpotente, como mostra o raciocínio precedente.
Veremos que, para todo $x\in\mathfrak{X}$, existe uma vizinhança aberta $U$ de $x$ em $\mathfrak{X}$ que é um \emph{aberto afim} para \emph{todos} os $X_i$.
Com efeito, tomemos $U$ aberto afim para $X_0$ e observemos que $\sh{O}_{X_0}=\sh{O}_{X_i}/\sh{J}_i$.
Como, pelo que precede, $\sh{J}_i|U$ é nilpotente, $U$ é também aberto afim para cada $X_i$, em virtude da Proposição~\sref{I.5.1.9}.
Nessas condições, para todo $U$ que satisfaça as condições precedentes, o estudo acima do caso afim mostra que $(U,\sh{O}_\mathfrak{X}|U)$ é um pré-esquema formal cujos $\sh{J}^{(i)}|U$ formam um sistema fundamental de feixes de ideais de definição, e que $\sh{J}^{(0)}|U$ é o limite projetivo dos $\sh{J}_i|U$; daí a conclusão.
\end{proof}

\begin{corollary}[10.6.4]
\label{I.10.6.4}
Suponhamos que, para $i\geq j$, o núcleo de $u_{ji}$ seja $\sh{J}_i^{j+1}$ e que $\sh{J}_1/\sh{J}_1^2$
\oldpage[I]{191}
seja de tipo finito sobre $\sh{O}_0=\sh{O}_1/\sh{J}_1$.
Então $\mathfrak{X}$ é um pré-esquema formal ádico e, se $\sh{J}^{(n)}$ é o núcleo de $\sh{O}_\mathfrak{X}\to\sh{O}_n$, tem-se $\sh{J}^{(n)}=\sh{J}^{n+1}$ e $\sh{J}/\sh{J}^2$ é isomorfo a $\sh{J}_1$.
Se, além disso, $X_0$ é localmente noetheriano (resp. noetheriano), então $\mathfrak{X}$ é localmente noetheriano (resp. noetheriano).
\end{corollary}

\begin{proof}
Como os espaços subjacentes a $\mathfrak{X}$ e $X_0$ são os mesmos, a questão é local e pode-se supor que todos os $X_i$ sejam afins; levando em conta as relações $\sh{J}_{ij}=\widetilde{\mathfrak{J}_{ji}}$ (com as notações da Proposição~\sref{I.10.6.3}), o problema reduz-se imediatamente às afirmações correspondentes da Proposição~\sref[0]{0.7.2.7} e do Corolário~\sref[0]{0.7.2.8}, observando que $\mathfrak{J}_1/\mathfrak{J}_1^2$ é então um $A_0$-módulo de tipo finito \sref{I.1.3.9}.
\end{proof}

Em particular, \emph{todo pré-esquema formal localmente noetheriano $\mathfrak{X}$} é limite indutivo de uma sequência $(X_n)$ de pré-esquemas (usuais) localmente noetherianos que satisfazem as condições da Proposição~\sref{I.10.6.3} e do Corolário~\sref{I.10.6.4}: basta considerar um feixe de ideais de definição $\sh{J}$ de $\mathfrak{X}$ \sref{I.10.5.4} e tomar $X_n=(\mathfrak{X},\sh{O}_\mathfrak{X}/\sh{J}^{n+1})$ (\sref{I.10.5.1} e Proposição~\sref{I.10.6.2}).

\begin{corollary}[10.6.5]
\label{I.10.6.5}
Seja $A$ um anel admissível.
Para que o esquema formal afim $\mathfrak{X}=\Spf(A)$ seja noetheriano, é necessário e suficiente que $A$ seja ádico e noetheriano.
\end{corollary}

\begin{proof}
A condição é evidentemente suficiente.
Reciprocamente, suponhamos que $\mathfrak{X}$ seja noetheriano, e sejam $\mathfrak{J}$ um ideal de definição de $A$ e $\sh{J}=\mathfrak{J}^\Delta$ o feixe de ideais de definição correspondente de $\mathfrak{X}$.
Os pré-esquemas (usuais) $X_n=(\mathfrak{X},\sh{O}_\mathfrak{X}/\sh{J}^{n+1})$ são então afins e noetherianos, logo os anéis $A_n=A/\mathfrak{J}^{n+1}$ são noetherianos \sref{I.6.1.3}; daí se conclui que $\mathfrak{J}/\mathfrak{J}^2$ é um $A/\mathfrak{J}$-módulo de tipo finito.
Como os $\sh{J}^n$ formam um sistema fundamental de feixes de ideais de definição de $\mathfrak{X}$ \sref{I.10.5.1}, tem-se $\sh{O}_\mathfrak{X}=\varprojlim\sh{O}_\mathfrak{X}/\sh{J}^n$ \sref{I.10.5.3}; conclui-se assim \sref{I.10.1.3} que $A$ é topologicamente isomorfo a $\varprojlim A/\mathfrak{J}^n$ e, portanto, é ádico e noetheriano \sref[0]{0.7.2.8}.
\end{proof}

\begin{remark}[10.6.6]
\label{I.10.6.6}
Com as notações da Proposição~\sref{I.10.6.3}, seja $\sh{F}_i$ um $\sh{O}_i$-módulo, e suponhamos dado, para $i\geq j$, um $v_{ij}$-morfismo $\theta_{ji}:\sh{F}_i\to\sh{F}_j$, de modo que $\theta_{kj}\circ\theta_{ji}=\theta_{ki}$ para $k\leq j\leq i$.
Como a aplicação contínua subjacente a $v_{ij}$ é a identidade, $\theta_{ji}$ é um homomorfismo de feixes de grupos abelianos sobre o espaço $\mathfrak{X}$; além disso, se $\sh{F}$ é o limite projetivo do sistema projetivo $(\sh{F}_i)$ de feixes de grupos abelianos, o fato de os $\theta_{ji}$ serem $v_{ij}$-morfismos permite definir sobre $\sh{F}$ uma estrutura de $\sh{O}_\mathfrak{X}$-módulo por passagem ao limite projetivo; munido dessa estrutura, diremos que $\sh{F}$ é o \emph{limite projetivo} (para os $\theta_{ji}$) do sistema de $\sh{O}_i$-módulos $(\sh{F}_i)$.
No caso particular em que $v_{ij}^*(\sh{F}_i)=\sh{F}_j$ e $\theta_{ji}$ é a \emph{identidade}, diremos, abreviadamente, que $\sh{F}$ é o limite projetivo de um sistema $(\sh{F}_i)$ tal que $v_{ij}^*(\sh{F}_i)=\sh{F}_j$ para $j\leq i$ (sem mencionar os $\theta_{ji}$).
\end{remark}

\begin{env}[10.6.7]
\label{I.10.6.7}
Sejam $\mathfrak{X}$ e $\mathfrak{Y}$ dois pré-esquemas formais, $\sh{J}$ (resp. $\sh{K}$) um ideal de definição de $\mathfrak{X}$ (resp. $\mathfrak{Y}$), e $f:\mathfrak{X}\to\mathfrak{Y}$ um morfismo tal que $f^*(\sh{K})\sh{O}_\mathfrak{X}\subset\sh{J}$.
Tem-se então, para todo inteiro $n>0$, $f^*(\sh{K}^n)\sh{O}_\mathfrak{X}=(f^*(\sh{K})\sh{O}_\mathfrak{X})^n\subset\sh{J}^n$; portanto, $f$ induz \sref{I.10.5.6} um morfismo de pré-esquemas (usuais) $f_n:X_n\to Y_n$, pondo $X_n=(\mathfrak{X},\sh{O}_\mathfrak{X}/\sh{J}^{n+1})$ e $Y_n=(\mathfrak{Y},\sh{O}_\mathfrak{Y}/\sh{K}^{n+1})$, e segue imediatamente das definições que os diagramas
\[
\label{I.10.6.7.1}
  \xymatrix{
    X_m\ar[r]^{f_m}\ar[d] &
    Y_m\ar[d]\\
    X_n\ar[r]_{f_n} &
    Y_n
  }
  \tag{10.6.7.1}
\]
\oldpage[I]{192}
são comutativos para $m\leq n$; em outras palavras, $(f_n)$ é um \emph{sistema indutivo} de morfismos.
\end{env}

\begin{env}[10.6.8]
\label{I.10.6.8}
Reciprocamente, seja $(X_n)$ (resp. $(Y_n)$) um sistema indutivo de pré-esquemas (usuais) que satisfaça as condições~(b) e (c) da Proposição~\sref{I.10.6.3}, e seja $\mathfrak{X}$ (resp. $\mathfrak{Y}$) seu limite indutivo.
Por definição dos limites indutivos, toda sequência $(f_n)$ de morfismos $X_n\to Y_n$ que forme um sistema indutivo admite um \emph{limite indutivo $f:\mathfrak{X}\to\mathfrak{Y}$}, que é o único morfismo de pré-esquemas formais para o qual os diagramas
\[
  \xymatrix{
    X_n\ar[r]^{f_n}\ar[d] &
    Y_n\ar[d]\\
    \mathfrak{X}\ar[r]_f &
    \mathfrak{Y}
  }
\]
são comutativos.
\end{env}

\begin{proposition}[10.6.9]
\label{I.10.6.9}
Sejam $\mathfrak{X}$ e $\mathfrak{Y}$ dois pré-esquemas formais localmente noetherianos, e $\sh{J}$ (resp. $\sh{K}$) um ideal de definição de $\mathfrak{X}$ (resp. $\mathfrak{Y}$); a aplicação $f\mapsto(f_n)$ definida em \sref{I.10.6.7} é uma bijeção do conjunto dos morfismos $f:\mathfrak{X}\to\mathfrak{Y}$ tais que $f^*(\sh{K})\sh{O}_\mathfrak{X}\subset\sh{J}$ sobre o conjunto das sequências $(f_n)$ de morfismos que tornam comutativos os diagramas \sref{I.10.6.7.1}.
\end{proposition}

\begin{proof}
Se $f$ é o limite indutivo de uma sequência desse tipo, resta mostrar que $f^*(\sh{K})\sh{O}_\mathfrak{X}\subset\sh{J}$.
Como a questão é local sobre $\mathfrak{X}$ e $\mathfrak{Y}$, pode-se reduzir ao caso em que $\mathfrak{X}=\Spf(A)$ e $\mathfrak{Y}=\Spf(B)$ são afins, sendo $A$ e $B$ anéis ádicos noetherianos, e $\sh{J}=\mathfrak{J}^\Delta$, $\sh{K}=\mathfrak{K}^\Delta$, onde $\mathfrak{J}$ (resp. $\mathfrak{K}$) é um ideal de definição de $A$ (resp. $B$).
Tem-se então $X_n=\Spec(A_n)$, $Y_n=\Spec(B_n)$, com $A_n=A/\mathfrak{J}^{n+1}$ e $B_n=B/\mathfrak{K}^{n+1}$, em virtude de \sref{I.10.3.6} e \sref{I.10.3.2}; além disso, $f_n=({}^a\vphi_n,\widetilde{\vphi_n})$, onde os homomorfismos $\vphi_n:B_n\to A_n$ formam um sistema projetivo, e portanto $f=({}^a\vphi,\widetilde{\vphi})$, onde $\vphi=\varprojlim\vphi_n$.
A comutatividade do diagrama \sref{I.10.6.7.1} para $m=0$ fornece então a condição $\vphi_n(\mathfrak{K}/\mathfrak{K}^{n+1})\subset\mathfrak{J}/\mathfrak{J}^{n+1}$ para todo $n$; passando ao limite projetivo, obtém-se $\vphi(\mathfrak{K})\subset\mathfrak{J}$, o que implica $f^*(\sh{K})\sh{O}_\mathfrak{X}\subset\sh{J}$ \sref{I.10.5.6}[ii].
\end{proof}

\begin{corollary}[10.6.10]
\label{I.10.6.10}
Sejam $\mathfrak{X}$ e $\mathfrak{Y}$ dois pré-esquemas formais localmente noetherianos, e seja $\sh{T}$ o maior ideal de definição de $\mathfrak{X}$ \sref{I.10.5.4}.
\begin{enumerate}
  \item[\rm{(i)}] Para todo ideal de definição $\sh{K}$ de $\mathfrak{Y}$ e todo morfismo $f:\mathfrak{X}\to\mathfrak{Y}$, tem-se $f^*(\sh{K})\sh{O}_\mathfrak{X}\subset\sh{T}$.
  \item[\rm{(ii)}] Há uma correspondência biunívoca canônica entre $\Hom(\mathfrak{X},\mathfrak{Y})$ e o conjunto das sequências $(f_n)$ de morfismos que tornam comutativos os diagramas \sref{I.10.6.7.1}, onde $X_n=(\mathfrak{X},\sh{O}_\mathfrak{X}/\sh{T}^{n+1})$ e $Y_n=(\mathfrak{Y},\sh{O}_\mathfrak{Y}/\sh{K}^{n+1})$.
\end{enumerate}
\end{corollary}

\begin{proof}
(ii) decorre imediatamente de (i) e da Proposição~\sref{I.10.6.9}.
Para demonstrar (i), pode-se reduzir ao caso em que $\mathfrak{X}=\Spf(A)$ e $\mathfrak{Y}=\Spf(B)$, sendo $A$ e $B$ noetherianos, e $\sh{T}=\mathfrak{T}^\Delta$, $\sh{K}=\mathfrak{K}^\Delta$, onde $\mathfrak{T}$ é o maior ideal de definição de $A$ e $\mathfrak{K}$ um ideal de definição de $B$.
Seja $f=({}^a\vphi,\widetilde{\vphi})$, onde $\vphi:B\to A$ é um homomorfismo contínuo; como os elementos de $\mathfrak{K}$ são topologicamente nilpotentes \sref[0]{0.7.1.4}[ii], também o são os de $\vphi(\mathfrak{K})$, de modo que $\vphi(\mathfrak{K})\subset\mathfrak{T}$, pois $\mathfrak{T}$ é o conjunto dos elementos topologicamente nilpotentes de $A$ \sref[0]{0.7.1.6}; daí a conclusão em virtude de \sref{I.10.5.6}[ii].
\end{proof}

\begin{corollary}[10.6.11]
\label{I.10.6.11}
Sejam $\mathfrak{S}$, $\mathfrak{X}$ e $\mathfrak{Y}$ três pré-esquemas formais localmente noetherianos, e sejam $f:\mathfrak{X}\to\mathfrak{S}$ e $g:\mathfrak{Y}\to\mathfrak{S}$ os morfismos que fazem de $\mathfrak{X}$ e $\mathfrak{Y}$ pré-esquemas formais sobre $\mathfrak{S}$.
Seja $\sh{J}$ (resp. $\sh{K}$, $\sh{L}$) um ideal de definição de $\mathfrak{S}$ (resp. $\mathfrak{X}$, $\mathfrak{Y}$), e suponhamos que $f^*(\sh{J})\sh{O}_\mathfrak{X}\subset\sh{K}$ e $g^*(\sh{J})\sh{O}_\mathfrak{Y}=\sh{L}$; ponhamos $S_n=(\mathfrak{S},\sh{O}_\mathfrak{S}/\sh{J}^{n+1})$, $X_n=(\mathfrak{X},\sh{O}_\mathfrak{X}/\sh{K}^{n+1})$ e $Y_n=(\mathfrak{Y},\sh{O}_\mathfrak{Y}/\sh{L}^{n+1})$.
Há então uma correspondência biunívoca canônica
\oldpage[I]{193}
entre $\Hom_\mathfrak{S}(\mathfrak{X},\mathfrak{Y})$ e o conjunto das sequências $(u_n)$ de $S_n$-morfismos $u_n:X_n\to Y_n$ que tornam comutativos os diagramas \sref{I.10.6.7.1}.
\end{corollary}

\begin{proof}
Para todo $\mathfrak{S}$-morfismo $u:\mathfrak{X}\to\mathfrak{Y}$, tem-se por definição $f=g\circ u$, donde
\[
  u^*(\sh{L})\sh{O}_\mathfrak{X}=u^*(g^*(\sh{J})\sh{O}_\mathfrak{Y})\sh{O}_\mathfrak{X}=f^*(\sh{J})\sh{O}_\mathfrak{X}\subset\sh{K}
\]
e o corolário decorre então da Proposição~\sref{I.10.6.9}.
\end{proof}

Note-se que, para $m\leq n$, o dado de um morfismo $f_n:X_n\to Y_n$ determina um único morfismo $f_m:X_m\to Y_m$ que torna comutativo o diagrama \sref{I.10.6.7.1}, como se vê imediatamente reduzindo ao caso afim; definiu-se assim uma aplicação $\vphi_{mn}:\Hom_{S_n}(X_n,Y_n)\to\Hom_{S_m}(X_m,Y_m)$, e os $\Hom_{S_n}(X_n,Y_n)$ formam, para os $\vphi_{mn}$, um \emph{sistema projetivo de conjuntos}; o Corolário~\sref{I.10.6.11} também se enuncia dizendo que existe uma bijeção canônica
\[
  \Hom_\mathfrak{S}(\mathfrak{X},\mathfrak{Y})\isoto\varprojlim_n\Hom_{S_n}(X_n,Y_n).
\]

\subsection{Produto de pré-esquemas formais}
\label{subsection:I.10.7}

\begin{env}[10.7.1]
\label{I.10.7.1}
Seja $\mathfrak{S}$ um pré-esquema formal; como os pré-esquemas formais sobre $\mathfrak{S}$ formam uma categoria, pode-se definir a noção de \emph{produto} de pré-esquemas formais sobre $\mathfrak{S}$.
\end{env}

\begin{proposition}[10.7.2]
\label{I.10.7.2}
Sejam $\mathfrak{X}=\Spf(B)$ e $\mathfrak{Y}=\Spf(C)$ dois esquemas formais afins sobre um esquema formal afim $\mathfrak{S}=\Spf(A)$.
Sejam $\mathfrak{Z}=\Spf(B\widehat{\otimes}_A C)$, $p_1$ e $p_2$ os $\mathfrak{S}$-morfismos que correspondem, por \sref{I.10.2.2}, aos $A$-homomorfismos canônicos (contínuos) $\rho:B\to B\widehat{\otimes}_A C$ e $\sigma:C\to B\widehat{\otimes}_A C$; então $(\mathfrak{Z},p_1,p_2)$ é um produto dos esquemas formais afins sobre $\mathfrak{S}$, $\mathfrak{X}$ e $\mathfrak{Y}$.
\end{proposition}

\begin{proof}
Por \sref{I.10.4.6}, basta verificar que a aplicação que associa a cada $A$-homomorfismo contínuo $\vphi:B\widehat{\otimes}_A C\to D$ (onde $D$ é um anel admissível que é uma $A$-álgebra topológica) o par $(\vphi\circ\rho,\vphi\circ\sigma)$ é uma bijeção
\[
  \Hom_A(B\widehat{\otimes}_A C,D)\isoto\Hom_A(B,D)\times\Hom_A(C,D),
\]
que é precisamente a propriedade universal do produto tensorial completado \sref[0]{0.7.7.6}.
\end{proof}

\begin{proposition}[10.7.3]
\label{I.10.7.3}
Dados dois $\mathfrak{S}$-pré-esquemas formais $\mathfrak{X}$ e $\mathfrak{Y}$, existe seu produto $\mathfrak{X}\times_\mathfrak{S}\mathfrak{Y}$.
\end{proposition}

\begin{proof}
A demonstração é idêntica à de \sref{I.3.2.6}, substituindo os esquemas afins (resp. os abertos afins) pelos esquemas formais afins (resp. os abertos formais afins), e a Proposição \sref{I.3.2.2} pela Proposição \sref{I.10.7.2}.
\end{proof}

Todas as propriedades formais do produto de pré-esquemas (\sref{I.3.2.7} e \sref{I.3.2.8}, de \sref{I.3.3.1} a \sref{I.3.3.12}) são válidas sem modificação alguma para o produto de pré-esquemas formais.

\begin{env}[10.7.4]
\label{I.10.7.4}
Sejam $\mathfrak{S}$, $\mathfrak{X}$ e $\mathfrak{Y}$ pré-esquemas formais, e sejam $f:\mathfrak{X}\to\mathfrak{S}$ e $g:\mathfrak{Y}\to\mathfrak{S}$ dois morfismos.
Suponhamos que existam, respectivamente em $\mathfrak{S}$, $\mathfrak{X}$ e $\mathfrak{Y}$, três sistemas fundamentais de feixes de ideais de definição $(\sh{J}_\lambda)$, $(\sh{K}_\lambda)$ e $(\sh{L}_\lambda)$, todos com o mesmo conjunto de índices $I$, tais que $f^*(\sh{J}_\lambda)\sh{O}_\mathfrak{X}\subset\sh{K}_\lambda$ e $g^*(\sh{J}_\lambda)\sh{O}_\mathfrak{Y}\subset\sh{L}_\lambda$ para todo $\lambda$.
Ponhamos $S_\lambda=(\mathfrak{S},\sh{O}_\mathfrak{S}/\sh{J}_\lambda)$, $X_\lambda=(\mathfrak{X},\sh{O}_\mathfrak{X}/\sh{K}_\lambda)$ e $Y_\lambda=(\mathfrak{Y},\sh{O}_\mathfrak{Y}/\sh{L}_\lambda)$; se $\sh{J}_\mu\subset\sh{J}_\lambda$, $\sh{K}_\mu\subset\sh{K}_\lambda$ e $\sh{L}_\mu\subset\sh{L}_\lambda$, observemos que $S_\lambda$ (resp. $X_\lambda$, $Y_\lambda$) é um subpré-esquema fechado de $S_\mu$ (resp. $X_\mu$, $Y_\mu$) com o \emph{mesmo}
\oldpage[I]{194}
espaço subjacente \sref{I.10.6.1}.
Como $S_\lambda\to S_\mu$ é um monomorfismo de pré-esquemas, vemos primeiro que os produtos $X_\lambda\times_{S_\lambda}Y_\lambda$ e $X_\lambda\times_{S_\mu}Y_\lambda$ são idênticos \sref{I.3.2.4}, e depois que $X_\lambda\times_{S_\mu}Y_\lambda$ se identifica com um subpré-esquema fechado de $X_\mu\times_{S_\mu}Y_\mu$ com o \emph{mesmo} espaço subjacente \sref{I.4.3.1}.
Assim, o produto $\mathfrak{X}\times_\mathfrak{S}\mathfrak{Y}$ é o \emph{limite indutivo} dos pré-esquemas usuais $X_\lambda\times_{S_\lambda}Y_\lambda$: com efeito, como em \sref{I.10.6.2}, podemos reduzir ao caso em que $\mathfrak{S}$, $\mathfrak{X}$ e $\mathfrak{Y}$ são esquemas formais afins.
Levando em conta \sref{I.10.5.6}[ii] e as hipóteses sobre os sistemas fundamentais de feixes de ideais de definição de $\mathfrak{S}$, $\mathfrak{X}$ e $\mathfrak{Y}$, vemos imediatamente que nossa afirmação decorre da definição do produto tensorial completado de duas álgebras \sref[0]{0.7.7.1}.

Além disso, seja $\mathfrak{Z}$ um $\mathfrak{S}$-pré-esquema formal, seja $(\sh{M}_\lambda)$ um sistema fundamental de ideais de definição de $\mathfrak{Z}$ cujo conjunto de índices é $I$, e sejam $u:\mathfrak{Z}\to\mathfrak{X}$ e $v:\mathfrak{Z}\to\mathfrak{Y}$ dois $\mathfrak{S}$-morfismos tais que $u^*(\sh{K}_\lambda)\sh{O}_\mathfrak{Z}\subset\sh{M}_\lambda$ e $v^*(\sh{L}_\lambda)\sh{O}_\mathfrak{Z}\subset\sh{M}_\lambda$ para todo $\lambda$.
Se pusermos $Z_\lambda=(\mathfrak{Z},\sh{O}_\mathfrak{Z}/\sh{M}_\lambda)$, e se $u_\lambda:Z_\lambda\to X_\lambda$ e $v_\lambda:Z_\lambda\to Y_\lambda$ forem os $S_\lambda$-morfismos correspondentes a $u$ e $v$ \sref{I.10.5.6}, verificamos imediatamente que $(u,v)_\mathfrak{S}$ é o limite indutivo dos $S_\lambda$-morfismos $(u_\lambda,v_\lambda)_{S_\lambda}$.

As considerações deste número aplicam-se, em particular, quando $\mathfrak{S}$, $\mathfrak{X}$ e $\mathfrak{Y}$ são localmente noetherianos, tomando como sistemas fundamentais de feixes de ideais de definição os sistemas formados pelas potências de um feixe de ideais de definição \sref{I.10.5.1}.
No entanto, observemos que $\mathfrak{X}\times_\mathfrak{S}\mathfrak{Y}$ não é necessariamente localmente noetheriano (veja-se, contudo, \sref{I.10.13.5}).
\end{env}

\subsection{Completamento formal de um pré-esquema ao longo de uma parte fechada}
\label{subsection:I.10.8}

\begin{env}[10.8.1]
\label{I.10.8.1}
Seja $X$ um pré-esquema (usual) \emph{localmente noetheriano}, e seja $X'$ uma parte fechada do espaço subjacente a $X$; designemos por $\Phi$ o conjunto dos feixes \emph{coerentes} de ideais $\sh{J}$ de $\sh{O}_X$ tais que o suporte de $\sh{O}_X/\sh{J}$ seja $X'$.
O conjunto $\Phi$ não é vazio (\sref{I.5.2.1}, \sref{I.4.1.4}, \sref{I.6.1.1}); nós o ordenaremos pela relação $\supset$.
\end{env}

\begin{lemma}[10.8.2]
\label{I.10.8.2}
O conjunto ordenado $\Phi$ é filtrante; se $X$ é noetheriano, para todo $\sh{J}_0\in\Phi$, o conjunto das potências $\sh{J}_0^n$ ($n>0$) é cofinal em $\Phi$.
\end{lemma}

\begin{proof}
Com efeito, se $\sh{J}_1$ e $\sh{J}_2$ pertencem a $\Phi$, e se pusermos $\sh{J}=\sh{J}_1\cap\sh{J}_2$, então $\sh{J}$ é coerente porque $\sh{O}_X$ é coerente (\sref{I.6.1.1} e \sref[0]{0.5.3.4}), e temos $\sh{J}_x=(\sh{J}_1)_x\cap(\sh{J}_2)_x$ para todo $x\in X$; portanto, $\sh{J}_x=\sh{O}_x$ para $x\not\in X'$, e $\sh{J}_x\neq\sh{O}_x$ para $x\in X'$, o que prova que $\sh{J}\in\Phi$.
Por outro lado, se $X$ é noetheriano e $\sh{J}_0$ e $\sh{J}$ pertencem a $\Phi$, existe um inteiro $n>0$ tal que $\sh{J}_0^n(\sh{O}_X/\sh{J})=0$ \sref{I.9.3.4}, o que implica $\sh{J}_0^n\subset\sh{J}$.
\end{proof}

\begin{env}[10.8.3]
\label{I.10.8.3}
Seja agora $\sh{F}$ um $\sh{O}_X$-módulo \emph{coerente}; para todo $\sh{J}\in\Phi$, $\sh{F}\otimes_{\sh{O}_X}(\sh{O}_X/\sh{J})$ é um $\sh{O}_X$-módulo coerente \sref{I.9.1.1}, cujo suporte está contido em $X'$, e que identificaremos quase sempre com sua restrição a $X'$.
Quando $\sh{J}$ percorre $\Phi$, esses feixes formam um \emph{sistema projetivo} de feixes de grupos abelianos.
\end{env}

\begin{definition}[10.8.4]
\label{I.10.8.4}
Dada uma parte fechada $X'$ de um pré-esquema localmente noetheriano $X$ e um $\sh{O}_X$-módulo coerente $\sh{F}$, chama-se \emph{completamento de $\sh{F}$ ao longo de $X'$}, e denota-se por $\sh{F}_{/X'}$ (ou por $\widehat{\sh{F}}$ quando não houver risco de confusão), a restrição a $X'$ do feixe
\oldpage[I]{195}
$\varprojlim_\Phi(\sh{F}\otimes_{\sh{O}_X}(\sh{O}_X/\sh{J}))$; dizemos que suas seções sobre $X'$ são as \emph{seções formais de $\sh{F}$ ao longo de $X'$}.
\end{definition}

É imediato que, para todo aberto $U\subset X$, temos $(\sh{F}|U)_{/(U\cap X')}=(\sh{F}_{/X'})|(U\cap X')$.

Por passagem ao limite projetivo, é claro que $(\sh{O}_X)_{/X'}$ é um feixe de anéis, e que $\sh{F}_{/X'}$ pode ser considerado como um $(\sh{O}_X)_{/X'}$-módulo.
Além disso, como existe uma base da topologia de $X'$ formada por abertos quase-compactos, podemos considerar $(\sh{O}_X)_{/X'}$ (resp. $\sh{F}_{/X'}$) como um \emph{feixe de anéis topológicos} (resp. de \emph{grupos topológicos}), limite projetivo dos feixes de anéis (resp. grupos) \emph{pseudodiscretos} $\sh{O}_X/\sh{J}$ (resp. $\sh{F}\otimes_{\sh{O}_X}(\sh{O}_X/\sh{J})=\sh{F}/\sh{J}\sh{F}$), e, por passagem ao limite projetivo, $\sh{F}_{/X'}$ torna-se então um \emph{$(\sh{O}_X)_{/X'}$-módulo topológico} (\sref[0]{0.3.8.1} e \sref[0]{0.3.8.2}); recordemos que, para todo aberto \emph{quase-compacto} $U\subset X$, $\Gamma(U\cap X',(\sh{O}_X)_{/X'})$ (resp. $\Gamma(U\cap X',\sh{F}_{/X'})$) é então o limite projetivo dos anéis (resp. grupos) discretos $\Gamma(U,\sh{O}_X/\sh{J})$ (resp. $\Gamma(U,\sh{F}/\sh{J}\sh{F})$).

Se agora $u:\sh{F}\to\sh{G}$ é um homomorfismo de $\sh{O}_X$-módulos, deduzem-se canonicamente homomorfismos $u_\sh{J}:\sh{F}\otimes_{\sh{O}_X}(\sh{O}_X/\sh{J})\to\sh{G}\otimes_{\sh{O}_X}(\sh{O}_X/\sh{J})$ para todo $\sh{J}\in\Phi$, e esses homomorfismos formam um sistema projetivo.
Por passagem ao limite projetivo e restrição a $X'$, eles dão um $(\sh{O}_X)_{/X'}$-homomorfismo contínuo $\sh{F}_{/X'}\to\sh{G}_{/X'}$, denotado por $u_{/X'}$ ou $\widehat{u}$, e chamado o \emph{completamento do homomorfismo $u$ ao longo de $X'$}.
É claro que, se $v:\sh{G}\to\sh{H}$ é um segundo homomorfismo de $\sh{O}_X$-módulos, temos $(v\circ u)_{/X'}=(v_{/X'})\circ(u_{/X'})$; portanto, $\sh{F}_{/X'}$ é um \emph{funtor aditivo covariante} em $\sh{F}$, da categoria dos $\sh{O}_X$-módulos coerentes para a categoria dos $(\sh{O}_X)_{/X'}$-módulos topológicos.

\begin{proposition}[10.8.5]
\label{I.10.8.5}
O suporte de $(\sh{O}_X)_{/X'}$ é $X'$; o espaço anelado topológico $(X',(\sh{O}_X)_{/X'})$ é um pré-esquema formal localmente noetheriano e, se $\sh{J}\in\Phi$, então $\sh{J}_{/X'}$ é um feixe de ideais de definição desse pré-esquema formal.
Se $X=\Spec(A)$ é um esquema afim de anel noetheriano $A$, $\sh{J}=\widetilde{\mathfrak{J}}$, onde $\mathfrak{J}$ é um ideal de $A$, e $X'=V(\mathfrak{J})$, então $(X',(\sh{O}_X)_{/X'})$ se identifica canonicamente com $\Spf(\widehat{A})$, onde $\widehat{A}$ é o completamento separado de $A$ para a topologia $\mathfrak{J}$-pré-ádica.
\end{proposition}

\begin{proof}
Evidentemente, basta provar a última afirmação.
Sabe-se por \sref[0]{0.7.3.3} que o completamento separado $\widehat{\mathfrak{J}}$ de $\mathfrak{J}$ para a topologia $\mathfrak{J}$-pré-ádica se identifica com o ideal $\mathfrak{J}\widehat{A}$ de $\widehat{A}$, e que $\widehat{A}$ é um anel noetheriano $\widehat{\mathfrak{J}}$-ádico tal que $\widehat{A}/\widehat{\mathfrak{J}}^n=A/\mathfrak{J}^n$ \sref[0]{0.7.2.6}.
Essa última igualdade mostra que os ideais primos abertos de $\widehat{A}$ são os ideais $\widehat{\mathfrak{p}}=\mathfrak{p}\widehat{A}$, onde $\mathfrak{p}$ é um ideal primo de $A$ que contém $\mathfrak{J}$, e que $\widehat{\mathfrak{p}}\cap A=\mathfrak{p}$; portanto, $\Spf(\widehat{A})=X'$.
Como $\sh{O}_X/\sh{J}^n=(A/\mathfrak{J}^n)\supertilde$, a proposição decorre imediatamente das definições.
\end{proof}

Diz-se que o pré-esquema formal assim definido é o \emph{completamento de $X$ ao longo de $X'$}, e denota-se por $X_{/X'}$ ou por $\widehat{X}$ quando não houver risco de confusão.
Quando se toma $X'=X$, pode-se tomar $\sh{J}=0$, e temos portanto $X_{/X}=X$.

É claro que, se $U$ é um subpré-esquema induzido sobre um aberto de $X$, então $U_{/(U\cap X')}$ se identifica canonicamente com o subpré-esquema formal induzido por $X_{/X'}$ sobre o aberto $U\cap X'$ de $X'$.

\begin{corollary}[10.8.6]
\label{I.10.8.6}
O pré-esquema (usual) $\widehat{X}_\red$ é o único subpré-esquema reduzido de $X$ cujo espaço subjacente é $X'$ \sref{I.5.2.1}.
Para que $\widehat{X}$ seja noetheriano, é necessário e suficiente que $\widehat{X}_\red$ o seja, e basta que $X$ o seja.
\end{corollary}

\oldpage[I]{196}

\begin{proof}
Como a determinação de $\widehat{X}_\red$ é local \sref{I.10.5.4}, podemos supor que $X$ seja um esquema afim de anel noetheriano $A$; com as notações de \sref{I.10.8.5}, o ideal $\mathfrak{T}$ dos elementos topologicamente nilpotentes de $\widehat{A}$ é a imagem inversa, pela aplicação canônica $\widehat{A}\to\widehat{A}/\widehat{\mathfrak{J}}=A/\mathfrak{J}$, do nilradical de $A/\mathfrak{J}$ \sref[0]{0.7.1.3}; portanto, $\widehat{A}/\mathfrak{T}$ é isomorfo ao quociente de $A/\mathfrak{J}$ por seu nilradical.
A primeira afirmação resulta então de \sref{I.10.5.4} e \sref{I.5.1.1}.
Se $\widehat{X}_\red$ é noetheriano, seu espaço subjacente $X'$ também o é; portanto, os $X_n'=\Spec(\sh{O}_X/\sh{J}^n)$ são noetherianos \sref{I.6.1.2}, e o mesmo vale para $\widehat{X}$ \sref{I.10.6.4}; a recíproca é imediata por \sref{I.6.1.2}.
\end{proof}

\begin{env}[10.8.7]
\label{I.10.8.7}
Os homomorfismos canônicos $\sh{O}_X\to\sh{O}_X/\sh{J}$ (para $\sh{J}\in\Phi$) formam um sistema projetivo e fornecem, passando ao limite projetivo, um homomorfismo de feixes de anéis $\theta:\sh{O}_X\to\psi_*((\sh{O}_X)_{/X'})=\varprojlim_\Phi(\sh{O}_X/\sh{J})$, onde $\psi$ denota a injeção canônica $X'\to X$ dos espaços subjacentes.
Denotaremos por $i$ (ou $i_X$) o morfismo (dito \emph{canônico})
\[
  (\psi,\theta):X_{/X'}\to X
\]
de espaços anelados.

Por tensorização, para todo $\sh{O}_X$-módulo coerente $\sh{F}$, os homomorfismos canônicos $\sh{O}_X\to\sh{O}_X/\sh{J}$ fornecem homomorfismos $\sh{F}\to\sh{F}\otimes_{\sh{O}_X}(\sh{O}_X/\sh{J})$ de $\sh{O}_X$-módulos que formam também um sistema projetivo e fornecem, portanto, passando ao limite projetivo, um homomorfismo canônico funtorial $\gamma:\sh{F}\to\psi_*(\sh{F}_{/X'})$ de $\sh{O}_X$-módulos.
\end{env}

\begin{proposition}[10.8.8]
\label{I.10.8.8}
\medskip\noindent
\begin{enumerate}
  \item[\rm{(i)}] O funtor $\sh{F}_{/X'}$ (em $\sh{F}$) é exato.
  \item[\rm{(ii)}] O homomorfismo funtorial $\gamma^\sharp:i^*(\sh{F})\to\sh{F}_{/X'}$ de $(\sh{O}_X)_{/X'}$-módulos é um isomorfismo.
\end{enumerate}
\end{proposition}

\begin{proof}
\medskip\noindent
\begin{enumerate}
  \item[\rm{(i)}] Basta provar que, se $0\to\sh{F}'\to\sh{F}\to\sh{F}''\to 0$ é uma sequência exata de $\sh{O}_X$-módulos coerentes e se $U$ é um aberto afim de $X$ de anel noetheriano $A$, então a sequência
\[
  0\to\Gamma(U\cap X',\sh{F}'_{/X'})\to\Gamma(U\cap X',\sh{F}_{/X'})\to\Gamma(U\cap X',\sh{F}''_{/X'})\to 0
\]
é exata.
Tem-se $\sh{F}|U=\widetilde{M}$, $\sh{F}'|U=\widetilde{M'}$ e $\sh{F}''|U=\widetilde{M''}$, onde $M$, $M'$ e $M''$ são três $A$-módulos de tipo finito tais que a sequência $0\to M'\to M\to M''\to 0$ é exata (\sref{I.1.5.1} e \sref{I.1.3.11}); seja $\sh{J}\in\Phi$ e seja $\mathfrak{J}$ um ideal de $A$ tal que $\sh{J}|U=\widetilde{\mathfrak{J}}$.
Tem-se então
\[
  \Gamma(U\cap X',\sh{F}\otimes_{\sh{O}_X}\sh{O}_X/\sh{J}^n)=M\otimes_A(A/\mathfrak{J}^n)
\]
\sref{I.1.3.12}; portanto, por definição do limite projetivo, tem-se
\[
  \Gamma(U\cap X',\sh{F}_{/X'})=\varprojlim_n(M\otimes_A(A/\mathfrak{J}^n))=\widehat{M},
\]
o completamento separado de $M$ para a topologia $\mathfrak{J}$-pré-ádica, e analogamente
\[
  \Gamma(U\cap X',\sh{F}'_{/X'})=\widehat{M'},\ \Gamma(U\cap X',\sh{F}''_{/X'})=\widehat{M''};
\]
nossa afirmação decorre então do fato de que, sendo $A$ noetheriano, o funtor $\widehat{M}$ em $M$ é exato na categoria dos $A$-módulos de tipo finito \sref[0]{0.7.3.3}.
\oldpage[I]{197}
  \item[\rm{(ii)}] A questão é local, de modo que podemos supor que se tem uma sequência exata $\sh{O}_X^m\to\sh{O}_X^n\to\sh{F}\to 0$ \sref[0]{0.5.3.2}; como $\gamma^\sharp$ é funtorial e os funtores $i^*(\sh{F})$ e $\sh{F}_{/X'}$ são exatos à direita (por (i) e \sref[0]{0.4.3.1}), tem-se o diagrama comutativo
\[
\label{I.10.8.8.1}
  \xymatrix{
    i^*(\sh{O}_X^m)\ar[r]\ar[d]_{\gamma^\sharp} &
    i^*(\sh{O}_X^n)\ar[r]\ar[d]_{\gamma^\sharp} &
    i^*(\sh{F})\ar[r]\ar[d]_{\gamma^\sharp} &
    0\\
    (\sh{O}_X^m)_{/X'}\ar[r] &
    (\sh{O}_X^n)_{/X'}\ar[r] &
    \sh{F}_{/X'}\ar[r] &
    0
  }
  \tag{10.8.8.1}
\]
cujas linhas são exatas.
Além disso, os funtores $i^*(\sh{F})$ e $\sh{F}_{/X'}$ comutam com as somas diretas finitas (\sref[0]{0.3.2.6} e \sref[0]{0.4.3.2}), de modo que nos reduzimos a provar nossa afirmação para $\sh{F}=\sh{O}_X$.
Tem-se $i^*(\sh{O}_X)=(\sh{O}_X)_{/X'}=\sh{O}_{\widehat{X}}$ \sref[0]{0.4.3.4}, e $\gamma^\sharp$ é um homomorfismo de \emph{$\sh{O}_{\widehat{X}}$-módulos}; basta, portanto, verificar que $\gamma^\sharp$ leva em si mesma a seção unidade de $\sh{O}_{\widehat{X}}$ sobre um aberto de $X'$, o que é imediato e mostra que, nesse caso, $\gamma^\sharp$ é a identidade.
\end{enumerate}
\end{proof}

\begin{corollary}[10.8.9]
\label{I.10.8.9}
O morfismo $i:X_{/X'}\to X$ de espaços anelados é plano.
\end{corollary}

\begin{proof}
Isso decorre de \sref[0]{0.6.7.3} e de \sref{I.10.8.8}[i].
\end{proof}

\begin{corollary}[10.8.10]
\label{I.10.8.10}
Se $\sh{F}$ e $\sh{G}$ são $\sh{O}_X$-módulos coerentes, existem isomorfismos canônicos funtoriais (em $\sh{F}$ e $\sh{G}$)
\[
\label{I.10.8.10.1}
  (\sh{F}_{/X'})\otimes_{(\sh{O}_X)_{/X'}}(\sh{G}_{/X'})\isoto(\sh{F}\otimes_{\sh{O}_X}\sh{G})_{/X'},
  \tag{10.8.10.1}
\]
\[
\label{I.10.8.10.2}
  (\shHom_{\sh{O}_X}(\sh{F},\sh{G}))_{/X'}\isoto\shHom_{(\sh{O}_X)_{/X'}}(\sh{F}_{/X'},\sh{G}_{/X'}).
  \tag{10.8.10.2}
\]
\end{corollary}

\begin{proof}
Isso decorre da identificação canônica de $i^*(\sh{F})$ com $\sh{F}_{/X'}$; a existência do primeiro isomorfismo é então um resultado válido para todos os morfismos de espaços anelados \sref[0]{0.4.3.3.1}, e a do segundo é um resultado válido para todos os morfismos planos \sref[0]{0.6.7.6}, por \sref{I.10.8.9}.
\end{proof}

\begin{proposition}[10.8.11]
\label{I.10.8.11}
Para todo $\sh{O}_X$-módulo coerente $\sh{F}$, o núcleo do homomorfismo canônico $\Gamma(X,\sh{F})\to\Gamma(X',\sh{F}_{/X'})$ induzido por $\sh{F}\to\sh{F}_{/X'}$ é formado pelas seções nulas em alguma vizinhança de $X'$.
\end{proposition}

\begin{proof}
Da definição de $\sh{F}_{/X'}$ resulta que a imagem canônica de uma tal seção é nula.
Reciprocamente, se $s\in\Gamma(X,\sh{F})$ tem imagem nula em $\Gamma(X',\sh{F}_{/X'})$, basta ver que todo $x\in X'$ admite uma vizinhança em $X$ na qual $s$ é nula, e podemos assim nos reduzir ao caso em que $X=\Spec(A)$ é afim, $A$ é noetheriano, $X'=V(\mathfrak{J})$ para algum ideal $\mathfrak{J}$ de $A$, e $\sh{F}=\widetilde{M}$ para algum $A$-módulo $M$ de tipo finito.
Então $\Gamma(X',\sh{F}_{/X'})$ é o completamento separado $\widehat{M}$ de $M$ para a topologia $\mathfrak{J}$-pré-ádica, e o homomorfismo $\Gamma(X,\sh{F})\to\Gamma(X',\sh{F}_{/X'})$ é o homomorfismo canônico $M\to\widehat{M}$.
Sabe-se por \sref[0]{0.7.3.7} que o núcleo desse homomorfismo é o conjunto dos $z\in M$ anulados por um elemento de $1+\mathfrak{J}$.
Portanto, tem-se $(1+f)s=0$ para algum $f\in\mathfrak{J}$; para todo $x\in X'$ deduz-se $(1_x+f_x)s_x=0$, e, como $1_x+f_x$ é invertível em $\sh{O}_x$ (pois $\mathfrak{J}_x\sh{O}_x$ está contido no ideal maximal de $\sh{O}_x$), tem-se $s_x=0$, o que demonstra a proposição.
\end{proof}

\begin{corollary}[10.8.12]
\label{I.10.8.12}
O suporte de $\sh{F}_{/X'}$ é igual a $\Supp(\sh{F})\cap X'$.
\end{corollary}

\begin{proof}
É claro que $\sh{F}_{/X'}$ é um $(\sh{O}_X)_{/X'}$-módulo de tipo finito (\sref{I.10.8.8}[ii] e \sref[0]{0.5.2.4}),
\oldpage[I]{198}
portanto seu suporte é fechado \sref[0]{0.5.2.2} e está evidentemente contido em $\Supp(\sh{F})\cap X'$.
Para mostrar que ele é igual a este último conjunto, reduzimo-nos imediatamente a provar que a igualdade $\Gamma(X',\sh{F}_{/X'})=0$ implica $\Supp(\sh{F})\cap X'=\emp$; isso decorre de \sref{I.10.8.11} e de \sref{I.1.4.1}.
\end{proof}

\begin{corollary}[10.8.13]
\label{I.10.8.13}
Seja $u:\sh{F}\to\sh{G}$ um homomorfismo de $\sh{O}_X$-módulos coerentes.
Para que $u_{/X'}:\sh{F}_{/X'}\to\sh{G}_{/X'}$ seja nulo, é necessário e suficiente que $u$ seja nulo em uma vizinhança de $X'$.
\end{corollary}

\begin{proof}
Com efeito, por \sref{I.10.8.8}[ii], $u_{/X'}$ se identifica com $i^*(u)$; portanto, se considerarmos $u$ como uma seção sobre $X$ do feixe $\sh{H}=\shHom_{\sh{O}_X}(\sh{F},\sh{G})$, $u_{/X'}$ é a seção sobre $X'$ de $i^*(\sh{H})=\sh{H}_{/X'}$ que lhe corresponde canonicamente (\sref{I.10.8.10.2} e \sref[0]{0.4.4.6}).
Basta, pois, aplicar \sref{I.10.8.11} ao $\sh{O}_X$-módulo coerente $\sh{H}$.
\end{proof}

\begin{corollary}[10.8.14]
\label{I.10.8.14}
Seja $u:\sh{F}\to\sh{G}$ um homomorfismo de $\sh{O}_X$-módulos coerentes.
Para que $u_{/X'}$ seja um monomorfismo (resp. um epimorfismo), é necessário e suficiente que $u$ seja um monomorfismo (resp. um epimorfismo) em uma vizinhança de $X'$.
\end{corollary}

\begin{proof}
Sejam $\sh{P}$ e $\sh{N}$ o conúcleo e o núcleo, respectivamente, de $u$, de modo que se tem a sequência exata $0\to\sh{N}\xrightarrow{v}\sh{F}\xrightarrow{u}\sh{G}\xrightarrow{w}\sh{P}\to 0$, e portanto, em virtude de \sref{I.10.8.8}[i], a sequência exata
\[
  0\to\sh{N}_{/X'}\xrightarrow{v_{/X'}}\sh{F}_{/X'}\xrightarrow{u_{/X'}}\sh{G}_{/X'}\xrightarrow{w_{/X'}}\sh{P}_{/X'}\to 0.
\]
Se $u_{/X'}$ é um monomorfismo (resp. um epimorfismo), tem-se $v_{/X'}=0$ (resp. $w_{/X'}=0$); portanto, existe uma vizinhança de $X'$ na qual $v=0$ (resp. $w=0$), em virtude de \sref{I.10.8.13}.
\end{proof}

\subsection{Prolongamento de um morfismo aos completamentos}
\label{subsection:I.10.9}

\begin{env}[10.9.1]
\label{I.10.9.1}
Sejam $X$ e $Y$ dois pré-esquemas (usuais) localmente Noetherianos, $f:X\to Y$ um morfismo, e $X'$ (resp. $Y'$) uma parte fechada do espaço subjacente $X$ (resp. $Y$), tais que $f(X')\subset Y'$.
Seja $\sh{J}$ (resp. $\sh{K}$) um feixe de ideais de $\sh{O}_X$ (resp. $\sh{O}_Y$) tal que o suporte de $\sh{O}_X/\sh{J}$ (resp. $\sh{O}_Y/\sh{K}$) seja $X'$ (resp. $Y'$) e que $f^*(\sh{K})\sh{O}_X\subset\sh{J}$; observemos que sempre existem tais feixes de ideais, pois, por exemplo, pode-se tomar como $\sh{J}$ o maior feixe de ideais de $\sh{O}_X$ que define um subpré-esquema de $X$ cujo espaço subjacente é $X'$ \sref{I.5.2.1}, e a hipótese $f(X')\subset Y'$ implica então $f^*(\sh{K})\sh{O}_X\subset\sh{J}$ \sref{I.5.2.4}.
Para todo inteiro $n>0$ tem-se $f^*(\sh{K}^n)\sh{O}_X\subset\sh{J}^n$ \sref[0]{0.4.3.5}; por conseguinte \sref{I.4.4.6}, se pusermos $X_n'=(X',\sh{O}_X/\sh{J}^{n+1})$ e $Y_n'=(Y',\sh{O}_Y/\sh{K}^{n+1})$, $f$ induz um morfismo $f_n:X_n'\to Y_n'$, e é imediato que os $f_n$ formam um sistema indutivo.
Denotaremos seu limite indutivo \sref{I.10.6.8} por $\widehat{f}:X_{/X'}\to Y_{/Y'}$, e diremos (por abuso de linguagem) que $\widehat{f}$ é o \emph{prolongamento de $f$ aos completamentos de $X$ e $Y$ ao longo de $X'$ e $Y'$}.
Verifica-se imediatamente que esse morfismo não depende da escolha dos feixes de ideais $\sh{J}$ e $\sh{K}$ que satisfazem as condições acima.
Basta, com efeito, vê-lo quando $X$ e $Y$ são esquemas afins Noetherianos de anéis $A$ e $B$, respectivamente; então $\sh{J}=\widetilde{\mathfrak{J}}$ e $\sh{K}=\widetilde{\mathfrak{K}}$, onde $\mathfrak{J}$ (resp. $\mathfrak{K}$) é um ideal de $A$ (resp. $B$), e $f$ corresponde a um homomorfismo de anéis $\vphi:B\to A$ tal que $\vphi(\mathfrak{K})\subset\mathfrak{J}$ (\sref{I.4.4.6} e \sref{I.1.7.4}); \emph{$\widehat{f}$ é então o morfismo que corresponde \sref{I.10.2.2} ao homomorfismo contínuo $\widehat{\vphi}:\widehat{B}\to\widehat{A}$}, onde $\widehat{A}$ (resp. $\widehat{B}$) é o completamento separado de $A$ (resp. $B$) para a topologia $\mathfrak{J}$-pré-ádica (resp. $\mathfrak{K}$-pré-ádica) \sref{I.10.6.8}; sabe-se que, se substituirmos $\sh{J}$ por outro
\oldpage[I]{199}
feixe de ideais $\sh{J}'=\widetilde{\mathfrak{J}'}$ tal que o suporte de $\sh{O}_X/\sh{J}'$ seja $X'$, então as topologias $\mathfrak{J}$-pré-ádica e $\mathfrak{J}'$-pré-ádica sobre $A$ são as mesmas \sref{I.10.8.2}.

Observemos que, por definição, a aplicação contínua $X'\to Y'$ dos espaços subjacentes a $X_{/X'}$ e $Y_{/Y'}$ que corresponde a $\widehat{f}$ não é senão a restrição de $f$ a $X'$.
\end{env}

\begin{env}[10.9.2]
\label{I.10.9.2}
Resulta imediatamente da definição anterior que o diagrama de morfismos de espaços anelados
\[
  \label{I.10.9.2.1}
  \xymatrix{
    \widehat{X}\ar[r]^{\widehat{f}}\ar[d]_{i_X} &
    \widehat{Y}\ar[d]^{i_Y}\\
    X\ar[r]_f &
    Y
  }
\]
é comutativo, sendo as setas verticais os morfismos canônicos \sref{I.10.8.7}.
\end{env}

\begin{env}[10.9.3]
\label{I.10.9.3}
Sejam $Z$ um terceiro pré-esquema, $g:Y\to Z$ um morfismo e $Z'$ uma parte fechada de $Z$ tal que $g(Y')\subset Z'$.
Se $\widehat{g}$ designa o completamento do morfismo $g$ ao longo de $Y'$ e $Z'$, resulta imediatamente de \sref{I.10.9.1} que
$\widehat{(g\circ f)}=\widehat{g}\circ\widehat{f}$.
\end{env}

\begin{proposition}[10.9.4]
\label{I.10.9.4}
Sejam $X$ e $Y$ dois $S$-pré-esquemas localmente Noetherianos, sendo $Y$ de tipo finito sobre $S$.
Sejam $f$ e $g$ dois $S$-morfismos de $X$ em $Y$ tais que $f(X')\subset Y'$ e $g(X')\subset Y'$.
Para que $\widehat{f}=\widehat{g}$, é necessário e suficiente que $f$ e $g$ coincidam em uma vizinhança de $X'$.
\end{proposition}

\begin{proof}
A condição é evidentemente suficiente (sem hipótese de finitude sobre $Y$).
Para ver que ela é necessária, observemos primeiro que a hipótese $\widehat{f}=\widehat{g}$ implica $f(x)=g(x)$ para todo $x\in X'$.
Por outro lado, sendo a questão local, podemos supor que $X$ e $Y$ são, respectivamente, vizinhanças abertas afins de $x$ e de $y=f(x)=g(x)$, com anéis Noetherianos, que $S$ é afim e que $\Gamma(Y,\sh{O}_Y)$ é uma $\Gamma(S,\sh{O}_S)$-álgebra de tipo finito \sref{I.6.3.3}.
Então $f$ e $g$ correspondem a dois $\Gamma(S,\sh{O}_S)$-homomorfismos $\rho$ e $\sigma$ de $\Gamma(Y,\sh{O}_Y)$ em $\Gamma(X,\sh{O}_X)$ \sref{I.1.7.3}, e, por hipótese, os prolongamentos por continuidade desses homomorfismos ao completamento separado de $\Gamma(Y,\sh{O}_Y)$ são iguais.
Conclui-se de \sref{I.10.8.11} que, para toda seção $s\in\Gamma(Y,\sh{O}_Y)$, as seções $\rho(s)$ e $\sigma(s)$ coincidem em uma vizinhança de $X'$ (dependente de $s$); como $\Gamma(Y,\sh{O}_Y)$ é uma álgebra de tipo finito sobre $\Gamma(S,\sh{O}_S)$, deduz-se imediatamente que existe uma vizinhança $V$ de $X'$ tal que $\rho(s)$ e $\sigma(s)$ coincidam em $V$ para \emph{toda} seção $s\in\Gamma(Y,\sh{O}_Y)$.
Se $h\in\Gamma(X,\sh{O}_X)$\footnote{\emph{[Trad.] A fonte francesa impressa tem aqui $\Gamma(Y,\sh{O}_X)$; o domínio deve ser $X$ para que $D(h)\subset X$ faça sentido.}} é tal que $D(h)$ é uma vizinhança de $x$ contida em $V$, conclui-se do que precede e de \sref{I.1.4.1}[d] que $f$ e $g$ coincidem em $D(h)$.
\end{proof}

\begin{proposition}[10.9.5]
\label{I.10.9.5}
Sob as hipóteses de \sref{I.10.9.1}, para todo $\sh{O}_Y$-módulo coerente $\sh{G}$, existe um isomorfismo canônico funtorial de $(\sh{O}_X)_{/X'}$-módulos
\[
  (f^*(\sh{G}))_{/X'}\isoto\widehat{f}^*(\sh{G}_{/Y'}).
\]
\end{proposition}

\begin{proof}
Se identificarmos canonicamente $(f^*(\sh{G}))_{/X'}$ com $i_X^*(f^*(\sh{G}))$, e $\widehat{f}^*(\sh{G}_{/Y'})$ com $\widehat{f}^*(i_Y^*(\sh{G}))$ \sref{I.10.8.8}, a proposição resulta imediatamente da comutatividade do diagrama de \sref{I.10.9.2}.
\end{proof}

\begin{env}[10.9.6]
\label{I.10.9.6}
Sejam agora $\sh{F}$ um $\sh{O}_X$-módulo coerente e $\sh{G}$ um $\sh{O}_Y$-módulo coerente.
Se $u:\sh{G}\to\sh{F}$ é um $f$-morfismo de $\sh{G}$ em $\sh{F}$, corresponde-lhe um $\sh{O}_X$-homomorfismo $u^\sharp:f^*(\sh{G})\to\sh{F}$, e por completamento um $(\sh{O}_X)_{/X'}$-homomorfismo contínuo $(u^\sharp)_{/X'}:(f^*(\sh{G}))_{/X'}\to\sh{F}_{/X'}$; em virtude de \sref{I.10.9.5}, existe um $\widehat{f}$-morfismo $v:\sh{G}_{/Y'}\to\sh{F}_{/X'}$,
\oldpage[I]{200}
e um só, tal que $v^\sharp=(u^\sharp)_{/X'}$.
Se considerarmos as triplas $(\sh{F},X,X')$ (sendo $\sh{F}$ um $\sh{O}_X$-módulo coerente e $X'$ uma parte fechada de $X$) como uma \emph{categoria}, cujos morfismos $(\sh{F},X,X')\to(\sh{G},Y,Y')$ consistem em um morfismo de pré-esquemas $f:X\to Y$ tal que $f(X')\subset Y'$ e um $f$-morfismo $u:\sh{G}\to\sh{F}$, pode-se dizer que $(X_{/X'},\sh{F}_{/X'})$ é um \emph{funtor} em $(\sh{F},X,X')$ com valores na categoria dos pares $(\mathfrak{Z},\sh{H})$ formados por um pré-esquema formal localmente Noetheriano $\mathfrak{Z}$ e um $\sh{O}_\mathfrak{Z}$-módulo $\sh{H}$, sendo os morfismos desta última categoria os pares formados por um morfismo $g$ de pré-esquemas formais e um $g$-morfismo.
\end{env}

\begin{proposition}[10.9.7]
\label{I.10.9.7}
Sejam $S$, $X$ e $Y$ três pré-esquemas localmente Noetherianos, $g:X\to S$ e $h:Y\to S$ dois morfismos, $S'$ uma parte fechada de $S$, e $X'$ (resp. $Y'$) uma parte fechada de $X$ (resp. $Y$) tal que $g(X')\subset S'$ (resp. $h(Y')\subset S'$); seja $Z=X\times_S Y$; suponhamos que $Z$ seja localmente Noetheriano, e seja $Z'=p^{-1}(X')\cap q^{-1}(Y')$, onde $p$ e $q$ são as projeções de $X\times_S Y$.
Nessas condições, o completamento $Z_{/Z'}$ identifica-se com o produto $(X_{/X'})\times_{S_{/S'}}(Y_{/Y'})$ dos $S_{/S'}$-pré-esquemas formais, identificando-se os morfismos estruturais com $\widehat{g}$ e $\widehat{h}$, e as projeções com $\widehat{p}$ e $\widehat{q}$.
\end{proposition}

\begin{proof}
É imediato que a questão é local para $S$, $X$ e $Y$, e reduzimo-nos assim ao caso em que $S=\Spec(A)$, $X=\Spec(B)$, $Y=\Spec(C)$, $S'=V(\mathfrak{J})$, $X'=V(\mathfrak{K})$ e $Y'=V(\mathfrak{L})$, onde $\mathfrak{J}$, $\mathfrak{K}$ e $\mathfrak{L}$ são ideais tais que $\vphi(\mathfrak{J})\subset\mathfrak{K}$ e $\psi(\mathfrak{J})\subset\mathfrak{L}$, denotando por $\vphi$ e $\psi$ os homomorfismos $A\to B$ e $A\to C$ que correspondem, respectivamente, a $g$ e $h$.
Sabe-se então que $Z=\Spec(B\otimes_A C)$ e que $Z'=V(\mathfrak{M})$, onde $\mathfrak{M}$ é o ideal $\Im(\mathfrak{K}\otimes_A C)+\Im(B\otimes_A\mathfrak{L})$.
A conclusão decorre de \sref{I.10.7.2} e do fato de que o produto tensorial completado $(\widehat{B}\otimes_{\widehat{A}}\widehat{C})^\wedge$ (onde $\widehat{A}$, $\widehat{B}$ e $\widehat{C}$ são, respectivamente, os completamentos separados de $A$, $B$ e $C$ para as topologias $\mathfrak{J}$-, $\mathfrak{K}$- e $\mathfrak{L}$-pré-ádicas) é o completamento separado do produto tensorial $B\otimes_A C$ para a topologia $\mathfrak{M}$-pré-ádica \sref[0]{0.7.7.2}.
\end{proof}

Observe-se, além disso, que, se $T$ é um $S$-pré-esquema localmente Noetheriano, $u:T\to X$ e $v:T\to Y$ são dois $S$-morfismos, e $T'$ é uma parte fechada de $T$ tal que $u(T')\subset X'$ e $v(T')\subset Y'$, então o prolongamento aos completamentos $((u,v)_S)^\wedge$ identifica-se com $(\widehat{u},\widehat{v})_{S_{/S'}}$.

\begin{corollary}[10.9.8]
\label{I.10.9.8}
Sejam $X$ e $Y$ dois $S$-pré-esquemas localmente Noetherianos tais que $X\times_S Y$ seja localmente Noetheriano; sejam $S'$ uma parte fechada de $S$, e $X'$ (resp. $Y'$) uma parte fechada de $X$ (resp. $Y$) cuja imagem em $S$ esteja contida em $S'$.
Para todo $S$-morfismo $f:X\to Y$ tal que $f(X')\subset Y'$, o morfismo gráfico $\Gamma_{\widehat f}$ identifica-se com o prolongamento $(\Gamma_f)^\wedge$ do morfismo gráfico de $f$.
\end{corollary}

\begin{corollary}[10.9.9]
\label{I.10.9.9}
Sejam $X$ e $Y$ dois pré-esquemas localmente Noetherianos, $f:X\to Y$ um morfismo, $Y'$ uma parte fechada de $Y$, e $X'=f^{-1}(Y')$.
Então o pré-esquema formal $X_{/X'}$ identifica-se, mediante o diagrama comutativo
\[
  \label{I.10.9.9.1}
  \xymatrix{
    X\ar[d]_f &
    X_{/X'}\ar[l]\ar[d]^{\widehat f}\\
    Y &
    Y_{/Y'}\ar[l]
  }
\]
com o produto $X\times_Y(Y_{/Y'})$ de pré-esquemas formais.
\end{corollary}

\begin{proof}
Basta aplicar \sref{I.10.9.7}, substituindo $S$ e $S'$ por $Y$, e $X$ e $X'$ por $X$.
\end{proof}

\begin{remark}[10.9.10]
\label{I.10.9.10}
\oldpage[I]{201}
Se $X$ é a soma disjunta $X_1\amalg X_2$ \hyperref[subsection:I.3.1]{(3.1)}, e $X'$ é a união $X'_1\cup X'_2$, onde $X'_i$ é uma parte fechada de $X_i$ ($i=1,2$), vê-se imediatamente que
\[
  X_{/X'}={X_1}_{/X'_1}\amalg {X_2}_{/X'_2}.
\]
\end{remark}

\subsection{Aplicação aos feixes coerentes sobre os esquemas formais afins}
\label{subsection:I.10.10}

\begin{env}[10.10.1]
\label{I.10.10.1}
Em todo este parágrafo, $A$ designará um \emph{anel ádico Noetheriano}, e $\mathfrak{J}$ um ideal de definição de $A$.
Sejam $X=\Spec(A)$ e $\mathfrak{X}=\Spf(A)$, identificando-se este último com a parte fechada $V(\mathfrak{J})$ de $X$ \sref{I.10.1.2}.
Além disso, as Definições \sref{I.10.1.2} e \sref{I.10.8.4} mostram que o \emph{esquema formal afim $\mathfrak{X}$} é idêntico ao \emph{completamento $X_{/\mathfrak{X}}$} do esquema afim $X$ ao longo da parte fechada $\mathfrak{X}$ de seu espaço subjacente.
A todo $\sh{O}_X$-módulo coerente $\sh{F}$ corresponde, portanto, um $\sh{O}_\mathfrak{X}$-módulo de tipo finito $\sh{F}_{/\mathfrak{X}}$, que é, além disso, um feixe de módulos topológicos sobre o feixe de anéis topológicos $\sh{O}_\mathfrak{X}$.
Mas todo $\sh{O}_X$-módulo coerente $\sh{F}$ é da forma $\widetilde{M}$, onde $M$ é um $A$-módulo de tipo finito \sref{I.1.5.1}; poremos $(\widetilde{M})_{/\mathfrak{X}}=M^\Delta$.
Além disso, se $u:M\to N$ é um $A$-homomorfismo de $A$-módulos de tipo finito, corresponde-lhe um homomorfismo $\widetilde{u}:\widetilde{M}\to\widetilde{N}$ e, por conseguinte, também um homomorfismo contínuo $\widetilde{u}_{/\mathfrak{X}}:(\widetilde{M})_{/\mathfrak{X}}\to(\widetilde{N})_{/\mathfrak{X}}$, que denotaremos por $u^\Delta$.
É imediato que $(v\circ u)^\Delta=v^\Delta\circ u^\Delta$; definiu-se assim um \emph{funtor aditivo covariante $M^\Delta$} da categoria dos $A$-módulos de tipo finito para a dos $\sh{O}_\mathfrak{X}$-módulos de tipo finito.
Quando $A$ é um anel \emph{discreto}, tem-se $M^\Delta=\widetilde{M}$.
\end{env}

\begin{proposition}[10.10.2]
\label{I.10.10.2}
\medskip\noindent
\begin{enumerate}
  \item[\rm{(i)}] $M^\Delta$ é um funtor exato em $M$, e existe um isomorfismo canônico funtorial de $A$-módulos $\Gamma(\mathfrak{X},M^\Delta)\isoto M$.
  \item[\rm{(ii)}] Se $M$ e $N$ são dois $A$-módulos de tipo finito, existem isomorfismos canônicos funtoriais
    \[
    \label{I.10.10.2.1}
      (M\otimes_A N)^\Delta\isoto M^\Delta\otimes_{\sh{O}_\mathfrak{X}}N^\Delta,
      \tag{10.10.2.1}
    \]
    \[
    \label{I.10.10.2.2}
      (\Hom_A(M,N))^\Delta\isoto\shHom_{\sh{O}_\mathfrak{X}}(M^\Delta,N^\Delta).
      \tag{10.10.2.2}
    \]
  \item[\rm{(iii)}] A aplicação $u\mapsto u^\Delta$ é um isomorfismo funtorial
    \[
    \label{I.10.10.2.3}
      \Hom_A(M,N)\isoto\Hom_{\sh{O}_\mathfrak{X}}(M^\Delta,N^\Delta).
      \tag{10.10.2.3}
    \]
\end{enumerate}
\end{proposition}

\begin{proof}
A exatidão de $M^\Delta$ decorre da exatidão dos funtores $\widetilde{M}$ \sref{I.1.3.5} e $\sh{F}_{/\mathfrak{X}}$ \sref{I.10.8.8}.
Por definição, $\Gamma(\mathfrak{X},M^\Delta)$ é o completamento separado do $A$-módulo $\Gamma(X,\widetilde{M})=M$ para a topologia $\mathfrak{J}$-pré-ádica; mas, como $A$ é completo e $M$ é de tipo finito, sabe-se \sref[0]{0.7.3.6} que $M$ é separado e completo, o que conclui a demonstração de (i).
O isomorfismo \sref{I.10.10.2.1} (resp. \sref{I.10.10.2.2}) provém da composição dos isomorfismos \sref{I.1.3.12}[i] e \sref{I.10.8.10.1} (resp. \sref{I.1.3.12}[ii] e \sref{I.10.8.10.2}).
Por fim, como $\Hom_A(M,N)$ é um $A$-módulo de tipo finito, podemos aplicar-lhe (i), que identifica $\Gamma(\mathfrak{X},(\Hom_A(M,N))^\Delta)$ com $\Hom_A(M,N)$, e utilizar \sref{I.10.10.2.2}, o que prova que o homomorfismo \sref{I.10.10.2.3} é um isomorfismo.
\end{proof}

Deduz-se da Proposição~\sref{I.10.10.2} toda uma série de consequências análogas às deduzidas do Teorema~\sref{I.1.3.7} e do Corolário~\sref{I.1.3.12}, cuja formulação deixamos ao leitor.

\oldpage[I]{202}
Observemos que a propriedade de exatidão de $M^\Delta$, aplicada à sequência exata $0\to\mathfrak{J}\to A\to A/\mathfrak{J}\to 0$, mostra que o feixe de ideais de $\sh{O}_\mathfrak{X}$ designado aqui por $\mathfrak{J}^\Delta$ coincide com aquele que havia sido denotado do mesmo modo em \sref{I.10.3.1}, em virtude de \sref{I.10.3.2}.

\begin{proposition}[10.10.3]
\label{I.10.10.3}
Sob as hipóteses de \sref{I.10.10.1}, $\sh{O}_\mathfrak{X}$ é um feixe coerente de anéis.
\end{proposition}

\begin{proof}
Se $f\in A$, sabe-se que $A_{\{f\}}$ é um anel ádico Noetheriano \sref[0]{0.7.6.11} e, como a questão é local, reduzimo-nos \sref{I.10.1.4} a provar que o núcleo de um homomorfismo $v:\sh{O}_\mathfrak{X}^n\to\sh{O}_\mathfrak{X}$ é um $\sh{O}_\mathfrak{X}$-módulo de tipo finito.
Temos então $v=u^\Delta$, onde $u$ é um $A$-homomorfismo $A^n\to A$ \sref{I.10.10.2}; como $A$ é Noetheriano, o núcleo de $u$ é de tipo finito, ou, equivalentemente, existe um homomorfismo $A^m\xrightarrow{w}A^n$ tal que a sequência $A^m\xrightarrow{w}A^n\xrightarrow{u}A$ seja exata.
Conclui-se \sref{I.10.10.2} que a sequência $\sh{O}_\mathfrak{X}^m\xrightarrow{w^\Delta}\sh{O}_\mathfrak{X}^n\xrightarrow{v}\sh{O}_\mathfrak{X}$ é exata, o que prova que o núcleo de $v$ é de tipo finito.
\end{proof}

\begin{env}[10.10.4]
\label{I.10.10.4}
Com as notações anteriores, ponha-se $A_n=A/\mathfrak{J}^{n+1}$, e seja $X_n$ o esquema afim $\Spec(A_n)=(\mathfrak{X},\sh{O}_\mathfrak{X}/\sh{J}^{n+1})$, sendo $\sh{J}=\mathfrak{J}^\Delta$ o feixe de ideais de definição de $\sh{O}_\mathfrak{X}$ correspondente ao ideal $\mathfrak{J}$.
Seja $u_{mn}$ o morfismo de pré-esquemas $X_m\to X_n$ correspondente ao homomorfismo canônico $A_n\to A_m$ para $m\leq n$; o esquema formal $\mathfrak{X}$ é o limite indutivo dos $X_n$ para os $u_{mn}$ \sref{I.10.6.3}.
\end{env}

\begin{proposition}[10.10.5]
\label{I.10.10.5}
Sob as hipóteses de \sref{I.10.10.1}, seja $\sh{F}$ um $\sh{O}_\mathfrak{X}$-módulo.
As condições seguintes são equivalentes:
\begin{enumerate}
  \item[\rm{(a)}] $\sh{F}$ é um $\sh{O}_\mathfrak{X}$-módulo coerente;
  \item[\rm{(b)}] $\sh{F}$ é isomorfo ao limite projetivo \sref{I.10.6.6} de uma sequência $(\sh{F}_n)$ de $\sh{O}_{X_n}$-módulos coerentes tais que $u_{mn}^*(\sh{F}_n)=\sh{F}_m$; e
  \item[\rm{(c)}] existe um $A$-módulo de tipo finito $M$ (determinado a menos de isomorfismo canônico pela Proposição~\sref{I.10.10.2}[i]) tal que $\sh{F}$ é isomorfo a $M^\Delta$.
\end{enumerate}
\end{proposition}

\begin{proof}
Mostremos primeiro que (b) implica (c).
Temos $\sh{F}_n=\widetilde{M_n}$, onde $M_n$ é um $A_n$-módulo de tipo finito, e a hipótese implica que $M_m=M_n\otimes_{A_n}A_m$ para $m\leq n$ \sref{I.1.6.5}; os $M_n$ formam, portanto, um sistema projetivo para os di-homomorfismos canônicos $M_n\to M_m$ ($m\leq n$), e resulta imediatamente da definição dos $A_n$ que esse sistema projetivo satisfaz as condições de \sref[0]{0.7.2.9}; por conseguinte, seu limite projetivo $M$ é um $A$-módulo de tipo finito tal que $M_n=M\otimes_A A_n$ para todo $n$.
Deduz-se que $\sh{F}_n$ é induzido sobre $X_n$ por $\widetilde{M}\otimes_{\sh{O}_X}(\sh{O}_X/\widetilde{\mathfrak{J}}^{n+1})$ e, portanto, $\sh{F}=M^\Delta$ por definição \sref{I.10.8.4}.
\end{proof}

Reciprocamente, (c) implica (b); com efeito, se $u_n$ é o morfismo de imersão $X_n\to X$, então $u_n^*(\widetilde{M})=(M\otimes_A A_n)\supertilde$ é induzido sobre $X_n$ por $\widetilde{M}\otimes_{\sh{O}_X}(\sh{O}_X/\widetilde{\mathfrak{J}}^{n+1})$, e $M^\Delta=\varprojlim u_n^*(\widetilde{M})$ por definição \sref{I.10.8.4}; como $u_m=u_n\circ u_{mn}$ para $m\leq n$, os $\sh{F}_n=u_n^*(\widetilde{M})$ satisfazem as condições de (b), donde resulta a afirmação.

Mostremos agora que (c) implica (a): com efeito, por definição, temos $\sh{O}_\mathfrak{X}=A^\Delta$; como $M$ é o conúcleo de um homomorfismo $A^m\to A^n$, deduz-se da Proposição~\sref{I.10.10.2} que $M^\Delta$ é o conúcleo de um homomorfismo $\sh{O}_\mathfrak{X}^m\to\sh{O}_\mathfrak{X}^n$ e, como o feixe de anéis $\sh{O}_\mathfrak{X}$ é coerente \sref{I.10.10.3}, o mesmo vale para $M^\Delta$ \sref[0]{0.5.3.4}.
\oldpage[I]{203}

Finalmente, (a) implica (b).
Considerado como $\sh{O}_\mathfrak{X}$-módulo, temos $\sh{O}_{X_n}=\sh{O}_\mathfrak{X}/\sh{J}^{n+1}=A_n^\Delta$; ora, $\sh{F}_n=\sh{F}\otimes_{\sh{O}_\mathfrak{X}}\sh{O}_{X_n}$ é um $\sh{O}_\mathfrak{X}$-módulo coerente \sref[0]{0.5.3.5} e, como também é um $\sh{O}_{X_n}$-módulo e $\sh{J}^{n+1}$ é coerente, conclui-se que $\sh{F}_n$ é um $\sh{O}_{X_n}$-módulo coerente \sref[0]{0.5.3.10}; além disso, é imediato que $u_{mn}^*(\sh{F}_n)=\sh{F}_m$ para $m\leq n$ (recordando que a aplicação contínua $X_m\to X_n$ dos espaços subjacentes é a identidade de $\mathfrak{X}$).
O feixe $\sh{G}=\varprojlim\sh{F}_n$ é, portanto, um $\sh{O}_\mathfrak{X}$-módulo coerente, pois vimos que (b) implica (a). Os homomorfismos canônicos $\sh{F}\to\sh{F}_n$ formam um sistema projetivo que, ao passar ao limite, dá um homomorfismo canônico $w:\sh{F}\to\sh{G}$, e tudo se reduz a provar que $w$ é bijetivo.
Como a questão agora é \emph{local}, basta considerar o caso em que $\sh{F}$ é o conúcleo de um homomorfismo $\sh{O}_\mathfrak{X}^p\to\sh{O}_\mathfrak{X}^q$; como esse homomorfismo é da forma $v^\Delta$, onde $v$ é um homomorfismo $A^p\to A^q$\footnote{\emph{[Trad.] A fonte francesa impressa tem aqui $A^m\to A^n$; os postos do homomorfismo precedente exigem $A^p\to A^q$.}}, temos que $\sh{F}$ é isomorfo a $M^\Delta$, onde $M=\Coker{v}$ \sref{I.10.10.2}.
Temos então, em virtude da Proposição~\sref{I.10.10.2}, $\sh{F}_n=M^\Delta\otimes_{\sh{O}_\mathfrak{X}}A_n^\Delta=(M\otimes_A A_n)^\Delta$ e, como a topologia $\mathfrak{J}$-ádica sobre $M\otimes_A A_n$ é discreta, temos $(M\otimes_A A_n)^\Delta=(M\otimes_A A_n)\supertilde$ (como $\sh{O}_{X_n}$-módulo); vimos acima que $M^\Delta=\varprojlim\sh{F}_n$, e $w$ é, portanto, a identidade neste caso.

\begin{corollary}[10.10.6]
\label{I.10.10.6}
Se $\sh{F}$ satisfaz a condição \emph{(b)} da Proposição~\sref{I.10.10.5}, o sistema projetivo $(\sh{F}_n)$ é isomorfo ao sistema dos $\sh{F}\otimes_{\sh{O}_\mathfrak{X}}\sh{O}_{X_n}$.
\end{corollary}

\begin{env}[10.10.7]
\label{I.10.10.7}
Sejam agora $A$ e $B$ dois anéis ádicos Noetherianos, e $\vphi:B\to A$ um homomorfismo contínuo; designaremos por $\mathfrak{J}$ (resp. $\mathfrak{K}$) um ideal de definição de $A$ (resp. de $B$) tal que $\vphi(\mathfrak{K})\subset\mathfrak{J}$, e poremos $X=\Spec(A)$, $Y=\Spec(B)$, $\mathfrak{X}=\Spf(A)$ e $\mathfrak{Y}=\Spf(B)$.
Seja $f:X\to Y$ o morfismo de pré-esquemas correspondente a $\vphi$ \sref{I.1.6.1}, e $\widehat{f}:\mathfrak{X}\to\mathfrak{Y}$ seu prolongamento aos completamentos \sref{I.10.9.1}, que também é o morfismo de pré-esquemas formais correspondente a $\vphi$ \sref{I.10.2.2}.
\end{env}

\begin{proposition}[10.10.8]
\label{I.10.10.8}
Para todo $B$-módulo $N$ de tipo finito, existe um isomorfismo canônico funtorial de $\sh{O}_\mathfrak{X}$-módulos
\[
  \widehat{f}^*(N^\Delta)\isoto(N\otimes_B A)^\Delta.
\]
\end{proposition}

\begin{proof}
Designando por $i_X:\mathfrak{X}\to X$ e $i_Y:\mathfrak{Y}\to Y$ os morfismos canônicos, temos \sref{I.10.8.8}, a menos de isomorfismos canônicos funtoriais, $N^\Delta=i_Y^*(\widetilde{N})$ e
\[
  (N\otimes_B A)^\Delta=i_X^*((N\otimes_B A)\supertilde)=i_X^*(f^*(\widetilde{N}))
\]
\sref{I.1.6.5}; a proposição decorre então da comutatividade do diagrama de \sref{I.10.9.2}.
\end{proof}

\begin{corollary}[10.10.9]
\label{I.10.10.9}
Para todo ideal $\mathfrak{b}$ de $B$, temos $\widehat{f}^*(\mathfrak{b}^\Delta)\sh{O}_\mathfrak{X}=(\mathfrak{b}A)^\Delta$.
\end{corollary}

\begin{proof}
Seja $j$ a injeção canônica $\mathfrak{b}\to B$, à qual corresponde a injeção canônica $j^\Delta:\mathfrak{b}^\Delta\to\sh{O}_\mathfrak{Y}$ de feixes de $\sh{O}_\mathfrak{Y}$-módulos; por definição, $\widehat{f}^*(\mathfrak{b}^\Delta)\sh{O}_\mathfrak{X}$ é a imagem do homomorfismo $\widehat{f}^*(j^\Delta):\widehat{f}^*(\mathfrak{b}^\Delta)\to\sh{O}_\mathfrak{X}=\widehat{f}^*(\sh{O}_\mathfrak{Y})$; mas esse homomorfismo se identifica com $(j\otimes 1)^\Delta:(\mathfrak{b}\otimes_B A)^\Delta\to\sh{O}_\mathfrak{X}=(B\otimes_B A)^\Delta$ pela Proposição~\sref{I.10.10.8}.
Como a imagem de $j\otimes 1$ é o ideal $\mathfrak{b}A$ de $A$, a imagem de $(j\otimes 1)^\Delta$ é, portanto, $(\mathfrak{b}A)^\Delta$, em virtude da Proposição~\sref{I.10.10.2}, donde a conclusão.
\end{proof}

\subsection{Feixes coerentes sobre os pré-esquemas formais}
\label{subsection:I.10.11}

\begin{proposition}[10.11.1]
\label{I.10.11.1}
\oldpage[I]{204}
Se $\mathfrak{X}$ é um pré-esquema formal localmente Noetheriano, o feixe de anéis $\sh{O}_\mathfrak{X}$ é coerente, e todo feixe de ideais de definição de $\mathfrak{X}$ é coerente.
\end{proposition}

\begin{proof}
A questão é local, de modo que se reduz ao caso de um esquema formal afim Noetheriano, e a proposição decorre então das Proposições~\sref{I.10.10.3} e~\sref{I.10.10.5}.
\end{proof}

\begin{env}[10.11.2]
\label{I.10.11.2}
Seja $\mathfrak{X}$ um pré-esquema formal localmente Noetheriano, $\sh{J}$ um feixe de ideais de definição de $\mathfrak{X}$, e $X_n$ o pré-esquema (usual) localmente Noetheriano $(\mathfrak{X},\sh{O}_\mathfrak{X}/\sh{J}^{n+1})$, de modo que $\mathfrak{X}$ é o \emph{limite indutivo} da sequência $(X_n)$ para os morfismos canônicos $u_{mn}:X_m\to X_n$ \sref{I.10.6.3}.
Com essas notações:
\end{env}

\begin{theorem}[10.11.3]
\label{I.10.11.3}
Para que um $\sh{O}_\mathfrak{X}$-módulo $\sh{F}$ seja coerente, é necessário e suficiente que ele seja isomorfo ao limite projetivo de uma sequência $(\sh{F}_n)$, onde os $\sh{F}_n$ são $\sh{O}_{X_n}$-módulos coerentes tais que $u_{mn}^*(\sh{F}_n)=\sh{F}_m$ para $m\leq n$ \sref{I.10.6.6}.
O sistema projetivo $(\sh{F}_n)$ é então isomorfo ao sistema dos $u_n^*(\sh{F})=\sh{F}\otimes_{\sh{O}_\mathfrak{X}}\sh{O}_{X_n}$, onde $u_n$ é o morfismo canônico $X_n\to\mathfrak{X}$.
\end{theorem}

\begin{proof}
A questão é local, de modo que se reduz ao caso em que $\mathfrak{X}$ é um esquema formal afim Noetheriano, e o teorema é então consequência da Proposição~\sref{I.10.10.5} e do Corolário~\sref{I.10.10.6}.
\end{proof}

Pode-se, portanto, dizer que \emph{o dado de um $\sh{O}_\mathfrak{X}$-módulo coerente equivale ao de um sistema projetivo $(\sh{F}_n)$ de $\sh{O}_{X_n}$-módulos coerentes tais que $u_{mn}^*(\sh{F}_n)=\sh{F}_m$ para $m\leq n$}.

\begin{corollary}[10.11.4]
\label{I.10.11.4}
Se $\sh{F}$ e $\sh{G}$ são dois $\sh{O}_\mathfrak{X}$-módulos coerentes, pode-se definir (com as notações do Teorema~\sref{I.10.11.3}) um isomorfismo canônico funtorial
\[
\label{I.10.11.4.1}
  \Hom_{\sh{O}_\mathfrak{X}}(\sh{F},\sh{G})\isoto\varprojlim_n\Hom_{\sh{O}_{X_n}}(\sh{F}_n,\sh{G}_n).
  \tag{10.11.4.1}
\]
\end{corollary}

\begin{proof}
O limite projetivo do segundo membro entende-se com respeito às aplicações $\theta_n\mapsto u_{mn}^*(\theta_n)$ ($m\leq n$) de $\Hom_{\sh{O}_{X_n}}(\sh{F}_n,\sh{G}_n)$ em $\Hom_{\sh{O}_{X_m}}(\sh{F}_m,\sh{G}_m)$.
O homomorfismo \sref{I.10.11.4.1} faz corresponder a um elemento $\theta\in\Hom_{\sh{O}_\mathfrak{X}}(\sh{F},\sh{G})$ a sequência $(u_n^*(\theta))$; define-se um homomorfismo inverso ao precedente fazendo corresponder ao sistema projetivo $(\theta_n)\in\varprojlim_n\Hom_{\sh{O}_{X_n}}(\sh{F}_n,\sh{G}_n)$ seu limite projetivo em $\Hom_{\sh{O}_\mathfrak{X}}(\sh{F},\sh{G})$, levando em conta o Teorema~\sref{I.10.11.3}.
\end{proof}

\begin{corollary}[10.11.5]
\label{I.10.11.5}
Para que um homomorfismo $\theta:\sh{F}\to\sh{G}$ seja sobrejetivo, é necessário e suficiente que o homomorfismo correspondente $\theta_0=u_0^*(\theta):\sh{F}_0\to\sh{G}_0$ o seja.
\end{corollary}

\begin{proof}
A questão é local, de modo que se reduz ao caso em que $\mathfrak{X}=\Spf(A)$, sendo $A$ um anel ádico Noetheriano, $\sh{F}=M^\Delta$, $\sh{G}=N^\Delta$ e $\theta=u^\Delta$, onde $M$ e $N$ são $A$-módulos de tipo finito e $u$ um homomorfismo $M\to N$; temos ainda $\theta_0=\widetilde{u_0}$, onde $u_0$ é o homomorfismo $u\otimes 1:M\otimes_A A/\mathfrak{J}\to N\otimes_A A/\mathfrak{J}$; a conclusão decorre de que $\theta$ e $u$ (resp. $\theta_0$ e $u_0$) são simultaneamente sobrejetivos (\sref{I.1.3.9} e \sref{I.10.10.2}) e de que $u$ e $u_0$ são simultaneamente sobrejetivos \sref[0]{0.7.1.14}.
\end{proof}

\begin{env}[10.11.6]
\label{I.10.11.6}
O Teorema~\sref{I.10.11.3} mostra que se pode considerar todo $\sh{O}_\mathfrak{X}$-módulo coerente $\sh{F}$ como um \emph{$\sh{O}_\mathfrak{X}$-módulo topológico}, considerando-o como limite projetivo dos feixes de grupos \emph{pseudodiscretos} $\sh{F}_n$ \sref[0]{0.3.8.1}.
Resulta então do Corolário~\sref{I.10.11.4} que todo homomorfismo $u:\sh{F}\to\sh{G}$ de $\sh{O}_\mathfrak{X}$-módulos coerentes é automaticamente \emph{contínuo}
\oldpage[I]{205}
\sref[0]{0.3.8.2}.
Além disso, se $\sh{H}$ é um $\sh{O}_\mathfrak{X}$-submódulo coerente de um $\sh{O}_\mathfrak{X}$-módulo coerente $\sh{F}$, então, para todo aberto $U\subset\mathfrak{X}$, $\Gamma(U,\sh{H})$ é um subgrupo \emph{fechado} do grupo topológico $\Gamma(U,\sh{F})$, pois, como o funtor $\Gamma$ é exato à esquerda, $\Gamma(U,\sh{H})$ é o núcleo do homomorfismo $\Gamma(U,\sh{F})\to\Gamma(U,\sh{F}/\sh{H})$, que é \emph{contínuo} pelo que precede, já que $\sh{F}/\sh{H}$ é coerente \sref[0]{0.5.3.4}; a afirmação resulta do fato de $\Gamma(U,\sh{F}/\sh{H})$ ser um grupo topológico separado.
\end{env}

\begin{proposition}[10.11.7]
\label{I.10.11.7}
Sejam $\sh{F}$ e $\sh{G}$ dois $\sh{O}_\mathfrak{X}$-módulos coerentes.
Podem-se definir (com as notações do Teorema~\sref{I.10.11.3}) isomorfismos canônicos funtoriais de $\sh{O}_\mathfrak{X}$-módulos topológicos \sref{I.10.11.6}
\[
\label{I.10.11.7.1}
  \sh{F}\otimes_{\sh{O}_\mathfrak{X}}\sh{G}\isoto\varprojlim_n(\sh{F}_n\otimes_{\sh{O}_{X_n}}\sh{G}_n),
  \tag{10.11.7.1}
\]
\[
\label{I.10.11.7.2}
  \shHom_{\sh{O}_\mathfrak{X}}(\sh{F},\sh{G})\isoto\varprojlim_n\shHom_{\sh{O}_{X_n}}(\sh{F}_n,\sh{G}_n).
  \tag{10.11.7.2}
\]
\end{proposition}

\begin{proof}
A existência do isomorfismo \sref{I.10.11.7.1} resulta da fórmula
\[
  \sh{F}_n\otimes_{\sh{O}_{X_n}}\sh{G}_n=(\sh{F}\otimes_{\sh{O}_\mathfrak{X}}\sh{O}_{X_n})\otimes_{\sh{O}_{X_n}}(\sh{G}\otimes_{\sh{O}_\mathfrak{X}}\sh{O}_{X_n})=(\sh{F}\otimes_{\sh{O}_\mathfrak{X}}\sh{G})\otimes_{\sh{O}_\mathfrak{X}}\sh{O}_{X_n}
\]
e do Teorema~\sref{I.10.11.3}.
O isomorfismo \sref{I.10.11.7.2}, no qual ambos os membros são considerados como feixes de módulos sem topologia, resulta da definição das seções de $\shHom_{\sh{O}_\mathfrak{X}}(\sh{F},\sh{G})$ e $\shHom_{\sh{O}_{X_n}}(\sh{F}_n,\sh{G}_n)$, e da existência do isomorfismo \sref{I.10.11.4.1}, aplicado ao pré-esquema induzido sobre um aberto formal afim Noetheriano arbitrário de $\mathfrak{X}$.
Resta provar que o isomorfismo \sref{I.10.11.7.2} é bicontínuo sobre um conjunto quase-compacto, e assim se reduz ao caso em que $\mathfrak{X}=\Spf(A)$, sendo $A$ um anel ádico Noetheriano, e portanto, em virtude de \sref{I.10.10.5}, àquele em que $\sh{F}=M^\Delta$ e $\sh{G}=N^\Delta$, sendo $M$ e $N$ dois $A$-módulos de tipo finito; levando em conta \sref{I.10.10.2.1}, \sref{I.10.10.2.3} e o Corolário~\sref{I.1.3.12}[ii], reduz-se a mostrar que o isomorfismo canônico $\Hom_A(M,N)\isoto\varprojlim_n\Hom_{A_n}(M_n,N_n)$ (com $M_n=M\otimes_A A_n$ e $N_n=N\otimes_A A_n$) é contínuo, o que já foi provado em \sref[0]{0.7.8.2}.
\end{proof}

\begin{env}[10.11.8]
\label{I.10.11.8}
Como $\Hom_{\sh{O}_\mathfrak{X}}(\sh{F},\sh{G})$ é o grupo das seções do feixe de grupos topológicos $\shHom_{\sh{O}_\mathfrak{X}}(\sh{F},\sh{G})$, ele é munido de uma topologia de grupo.
Se $\mathfrak{X}$ é \emph{Noetheriano}, resulta de \sref{I.10.11.7.2} que os subgrupos $\Hom_{\sh{O}_\mathfrak{X}}(\sh{F},\sh{J}^n\sh{G})$ (para $n$ arbitrário) formam um sistema fundamental de vizinhanças de $0$ nesse grupo.
\end{env}

\begin{proposition}[10.11.9]
\label{I.10.11.9}
Sejam $\mathfrak{X}$ um pré-esquema formal Noetheriano, e $\sh{F}$ e $\sh{G}$ dois $\sh{O}_\mathfrak{X}$-módulos coerentes.
No grupo topológico $\Hom_{\sh{O}_\mathfrak{X}}(\sh{F},\sh{G})$, os homomorfismos sobrejetivos (resp. injetivos, bijetivos) formam uma parte aberta.
\end{proposition}

\begin{proof}
Em virtude do Corolário~\sref{I.10.11.5}, o conjunto dos homomorfismos sobrejetivos de $\Hom_{\sh{O}_\mathfrak{X}}(\sh{F},\sh{G})$ é a imagem inversa, pela aplicação contínua $\Hom_{\sh{O}_\mathfrak{X}}(\sh{F},\sh{G})\to\Hom_{\sh{O}_{X_0}}(\sh{F}_0,\sh{G}_0)$, de uma parte do grupo discreto $\Hom_{\sh{O}_{X_0}}(\sh{F}_0,\sh{G}_0)$, donde resulta a primeira afirmação.
Para demonstrar a segunda, cubramos $\mathfrak{X}$ por um número finito de abertos formais afins Noetherianos $U_i$.
Para que $\theta\in\Hom_{\sh{O}_\mathfrak{X}}(\sh{F},\sh{G})$ seja injetivo, é necessário e suficiente que todas as suas imagens pelas aplicações (contínuas) de restrição $\Hom_{\sh{O}_\mathfrak{X}}(\sh{F},\sh{G})\to\Hom_{\sh{O}_\mathfrak{X}|U_i}(\sh{F}|U_i,\sh{G}|U_i)$ sejam injetivas; reduz-se assim ao caso afim, e isso já foi provado em \sref[0]{0.7.8.3}.
\end{proof}

\subsection{Morfismos ádicos de pré-esquemas formais}
\label{subsection:I.10.12}

\begin{env}[10.12.1]
\label{I.10.12.1}
\oldpage[I]{206}
Sejam $\mathfrak{X}$ e $\mathfrak{S}$ dois pré-esquemas formais \emph{localmente noetherianos};
diremos que um morfismo $f:\mathfrak{X}\to\mathfrak{S}$ é \emph{ádico} se existe um ideal de definição $\sh{J}$ de $\mathfrak{S}$ tal que $\sh{K}=f^*(\sh{J})\sh{O}_\mathfrak{X}$ seja um ideal de definição de $\mathfrak{X}$;
diz-se também então que $\mathfrak{X}$ é um \emph{$\mathfrak{S}$-pré-esquema ádico} (para $f$).
Quando isso ocorre, para \emph{todo} ideal de definição $\sh{J}_1$ de $\mathfrak{S}$, $\sh{K}_1=f^*(\sh{J}_1)\sh{O}_\mathfrak{X}$ é um ideal de definição de $\mathfrak{X}$.
Com efeito, como a questão é local, pode-se supor que $\mathfrak{X}$ e $\mathfrak{S}$ são afins e noetherianos;
existe então um inteiro $n$ tal que $\sh{J}^n\subset\sh{J}_1$ e $\sh{J}_1^n\subset\sh{J}$ (\sref{I.10.3.6} e \sref[0]{0.7.1.4}), donde $\sh{K}^n\subset\sh{K}_1$ e $\sh{K}_1^n\subset\sh{K}$.
A primeira dessas relações mostra que $\sh{K}_1=\mathfrak{K}_1^\Delta$, onde $\mathfrak{K}_1$ é um ideal aberto de $A=\Gamma(\mathfrak{X},\sh{O}_\mathfrak{X})$, e a segunda mostra que $\mathfrak{K}_1$ é um ideal de definição de $A$ \sref[0]{0.7.1.4}, donde a nossa afirmação.
\end{env}

Resulta imediatamente do que precede que, se $\mathfrak{X}$ e $\mathfrak{Y}$ são dois $\mathfrak{S}$-pré-esquemas ádicos, \emph{todo $\mathfrak{S}$-morfismo $u:\mathfrak{X}\to\mathfrak{Y}$ é ádico}:
com efeito, se $f:\mathfrak{X}\to\mathfrak{S}$ e $g:\mathfrak{Y}\to\mathfrak{S}$ são os morfismos estruturais, e $\sh{J}$ é um ideal de definição de $\mathfrak{S}$, tem-se $f=g\circ u$, e portanto $u^*(g^*(\sh{J})\sh{O}_\mathfrak{Y})\sh{O}_\mathfrak{X}=f^*(\sh{J})\sh{O}_\mathfrak{X}$ é um ideal de definição de $\mathfrak{X}$ e, por hipótese, $g^*(\sh{J})\sh{O}_\mathfrak{Y}$ é um ideal de definição de $\mathfrak{Y}$.

\begin{env}[10.12.2]
\label{I.10.12.2}
No que segue, suporemos fixados um pré-esquema formal localmente noetheriano $\mathfrak{S}$ e um ideal de definição $\sh{J}$ de $\mathfrak{S}$;
poremos $S_n=(\mathfrak{S},\sh{O}_\mathfrak{S}/\sh{J}^{n+1})$.
Os $\mathfrak{S}$-pré-esquemas ádicos (localmente noetherianos) formam evidentemente uma \emph{categoria}.
Diremos que um sistema indutivo $(X_n)$ de $S_n$-pré-esquemas (usuais) localmente noetherianos é um \emph{$(S_n)$-sistema indutivo ádico} se os morfismos estruturais $f_n:X_n\to S_n$ são tais que, para $m\leq n$, os diagramas
\[
\label{I.10.12.2.1}
  \xymatrix{
    X_n \ar[d]_{f_n}
    & X_m \ar[l] \ar[d]^{f_m}\\
    S_n
    & S_m \ar[l]
  }
  \tag{10.12.2.1}
\]
sejam comutativos e \emph{identifiquem $X_m$ com o produto $X_n\times_{S_n}S_m=(X_n)_{(S_m)}$}.
Os sistemas indutivos ádicos formam uma \emph{categoria}:
basta, com efeito, definir um morfismo $(X_n)\to(Y_n)$ de tais sistemas como um \emph{sistema indutivo de $S_n$-morfismos $u_n:X_n\to Y_n$} tal que $u_m$ se identifique com $(u_n)_{(S_m)}$ para $m\leq n$.
Posto isso:
\end{env}

\begin{theorem}[10.12.3]
\label{I.10.12.3}
Há uma equivalência canônica entre a categoria dos $\mathfrak{S}$-pré-esquemas ádicos e a categoria dos $(S_n)$-sistemas indutivos ádicos.
\end{theorem}

A equivalência em questão é obtida da seguinte maneira:
se $\mathfrak{X}$ é um $\mathfrak{S}$-pré-esquema ádico e $f:\mathfrak{X}\to\mathfrak{S}$ é o morfismo estrutural, então $\sh{K}=f^*(\sh{J})\sh{O}_\mathfrak{X}$ é um ideal de definição de $\mathfrak{X}$, e faz-se corresponder a $\mathfrak{X}$ o sistema indutivo dos $X_n=(\mathfrak{X},\sh{O}_\mathfrak{X}/\sh{K}^{n+1})$, sendo $f_n:X_n\to S_n$ o morfismo estrutural correspondente a $f$ \sref{I.10.5.6}.
Mostremos primeiro que $(X_n)$ é um \emph{sistema indutivo ádico}:
se $f=(\psi,\theta)$, tem-se $\psi^*(\sh{J})\sh{O}_\mathfrak{X}=\sh{K}$, logo $\psi^*(\sh{J}^n)\sh{O}_\mathfrak{X}=\sh{K}^n$ para todo $n$, e (pela exatidão do funtor $\psi^*$) $\sh{K}^{m+1}/\sh{K}^{n+1}=\psi^*(\sh{J}^{m+1}/\sh{J}^{n+1})(\sh{O}_\mathfrak{X}/\sh{K}^{n+1})$ para $m\leq n$;
nossa conclusão decorre, portanto, de \sref{I.4.4.5}.
Além disso, verifica-se imediatamente que a um $\mathfrak{S}$-morfismo $u:\mathfrak{X}\to\mathfrak{Y}$ de $\mathfrak{S}$-pré-esquemas ádicos corresponde (com notações evidentes)
\oldpage[I]{207}
um sistema indutivo de $S_n$-morfismos $u_n:X_n\to Y_n$ tal que $u_m$ se identifica com $(u_n)_{(S_m)}$ para $m\leq n$.

O fato de se ter definido assim uma equivalência resultará da seguinte proposição mais precisa:

\begin{proposition}[10.12.3.1]
\label{I.10.12.3.1}
Seja $(X_n)$ um sistema indutivo de $S_n$-pré-esquemas;
suponhamos que os morfismos estruturais $f_n:X_n\to S_n$ sejam tais que os diagramas de \sref{I.10.12.2.1} sejam comutativos e identifiquem $X_m$ com $X_n\times_{S_n}S_m$ para $m\leq n$.
Então o sistema indutivo $(X_n)$ satisfaz as condições \emph{(b)} e \emph{(c)} de \sref{I.10.6.3};
seja $\mathfrak{X}$ seu limite indutivo e seja $f:\mathfrak{X}\to\mathfrak{S}$ o morfismo limite indutivo do sistema indutivo $(f_n)$.
Então, se $X_0$ é localmente noetheriano, $\mathfrak{X}$ é localmente noetheriano e $f$ é um morfismo ádico.
\end{proposition}

\begin{proof}
Como o feixe de ideais de $\sh{O}_{S_n}$ que define o subpré-esquema $S_m$ de $S_n$ é nilpotente, o mesmo ocorre, em virtude de \sref{I.4.4.5}, com o feixe de ideais de $\sh{O}_{X_n}$ que define o subpré-esquema $X_m$ de $X_n$; portanto, as condições de \sref{I.10.6.3} são satisfeitas.
A questão é, pois, local sobre $\mathfrak{X}$ e $\mathfrak{S}$, e pode-se supor que $\mathfrak{S}=\Spf(A)$, $\sh{J}=\mathfrak{J}^\Delta$ (sendo $A$ um anel $\mathfrak{J}$-ádico noetheriano) e $X_n=\Spec(B_n)$;
se $A_n=A/\mathfrak{J}^{n+1}$, a hipótese implica que $B_0$ é noetheriano e, se pusermos $\mathfrak{J}_n=\mathfrak{J}/\mathfrak{J}^{n+1}$, então $B_m=B_n/\mathfrak{J}_n^{m+1}B_n$.
O núcleo de $B_n\to B_0$ é, portanto, $\mathfrak{K}_n=\mathfrak{J}_nB_n$, e o núcleo de $B_n\to B_m$ é $\mathfrak{K}_n^{m+1}$ para $m\leq n$;
além disso, como $A_1$ é noetheriano, $\mathfrak{J}_1$ é de tipo finito sobre $A_1$, logo $\mathfrak{K}_1/\mathfrak{K}_1^2$ é de tipo finito sobre $B_1$ e, \emph{a fortiori}, de tipo finito sobre $B_0=B_1/\mathfrak{K}_1$;
o fato de $\mathfrak{X}$ ser noetheriano decorre então de \sref{I.10.6.4};
se $B=\varprojlim B_n$, tem-se $\mathfrak{X}=\Spf(B)$ e, se $\mathfrak{K}$ é o núcleo de $B\to B_0$, então $B_n=B/\mathfrak{K}^{n+1}$.
Se $\rho_n:A/\mathfrak{J}^{n+1}\to B/\mathfrak{K}^{n+1}$ é o homomorfismo correspondente a $f_n$, tem-se
\[
    \mathfrak{K}/\mathfrak{K}^{n+1} = (B/\mathfrak{K}^{n+1})\rho_n(\mathfrak{J}/\mathfrak{J}^{n+1})
\]
Como o homomorfismo $\rho:A\to B$ correspondente a $f$ é igual a $\varprojlim\rho_n$, o ideal $\mathfrak{J}B$ de $B$ é denso em $\mathfrak{K}$; como todo ideal de $B$ é fechado \sref[0]{0.7.3.5}, tem-se também $\mathfrak{K}=\mathfrak{J}B$.
Se $\sh{K}=\mathfrak{K}^\Delta$, a igualdade $f^*(\sh{J})\sh{O}_\mathfrak{X}=\sh{K}$ decorre então de \sref{I.10.10.9} e conclui a demonstração.
\end{proof}

\begin{env}[10.12.3.2]
\label{I.10.12.3.2}
A equivalência precedente fornece, para dois $\mathfrak{S}$-pré-esquemas ádicos $\mathfrak{X}$ e $\mathfrak{Y}$, uma \emph{bijeção canônica}
\[
  \Hom_\mathfrak{S}(\mathfrak{X},\mathfrak{Y})\isoto\varprojlim_n\Hom_{S_n}(X_n,Y_n)
\]
em que o limite projetivo é relativo às aplicações $u_n\to(u_n)_{(S_m)}$ para $m\leq n$.
\end{env}

\subsection{Morfismos de tipo finito}
\label{subsection:I.10.13}

\begin{proposition}[10.13.1]
\label{I.10.13.1}
Seja $\mathfrak{Y}$ um pré-esquema formal localmente noetheriano, $\sh{K}$ um ideal de definição de $\mathfrak{Y}$, e seja $f:\mathfrak{X}\to\mathfrak{Y}$ um morfismo de pré-esquemas formais.
As condições seguintes são equivalentes.
\begin{enumerate}
  \item[\rm{(a)}] $\mathfrak{X}$ é localmente noetheriano, $f$ é um morfismo ádico \sref{I.10.12.1} e, se pusermos $\sh{J}=f^*(\sh{K})\sh{O}_\mathfrak{X}$, o morfismo $f_0:(\mathfrak{X},\sh{O}_\mathfrak{X}/\sh{J})\to(\mathfrak{Y},\sh{O}_\mathfrak{Y}/\sh{K})$ deduzido de $f$ é de tipo finito.
  \item[\rm{(b)}] $\mathfrak{X}$ é localmente noetheriano e é o limite indutivo de um $(Y_n)$-sistema indutivo ádico $(X_n)$ tal que o morfismo $X_0\to Y_0$ seja de tipo finito.
\oldpage[I]{208}
  \item[\rm{(c)}] Todo ponto de $\mathfrak{Y}$ possui uma vizinhança aberta formal afim noetheriana $V$ que satisfaz a seguinte propriedade:
\begin{enumerate}
  \item[(\textbf{Q})] $f^{-1}(V)$ é a união de uma família finita de abertos formais afins noetherianos $U_i$ tais que o anel ádico noetheriano $\Gamma(U_i,\sh{O}_\mathfrak{X})$ seja topologicamente isomorfo ao quociente de uma álgebra de séries formais restritas \sref[0]{0.7.5.1} sobre $\Gamma(V,\sh{O}_\mathfrak{Y})$ por um ideal (necessariamente fechado).
\end{enumerate}
\end{enumerate}
\end{proposition}

\begin{proof}
É imediato que (a) implica (b), em virtude de \sref{I.10.12.3}.
Para demonstrar que (b) implica (c), pode-se, como a questão é local sobre $\mathfrak{Y}$, supor que $\mathfrak{Y}=\Spf(B)$, onde $B$ é noetheriano e ádico;
seja $\sh{K}=\mathfrak{K}^\Delta$, sendo $\mathfrak{K}$ um ideal de definição de $B$.
Como, por hipótese, $X_0$ é de tipo finito sobre $Y_0$, $X_0$ é a união de um número finito de abertos afins $U_i$ tais que o anel $A_{i0}$ do esquema afim induzido por $X_0$ sobre $U_i$ seja uma álgebra de tipo finito sobre o anel $B/\mathfrak{K}$ de $Y_0$ \sref{I.6.3.2}.
Em virtude de \sref{I.5.1.9}, $U_i$ é também um aberto afim de cada um dos pré-esquemas noetherianos $X_n$ e, se $A_{in}$ é o anel do esquema afim induzido por $X_n$ sobre $U_i$, a hipótese (b) implica, para $m\leq n$, que $A_{im}$ é isomorfo a $A_{in}/\mathfrak{K}^{m+1}A_{in}$.
Consequentemente, o pré-esquema formal induzido por $\mathfrak{X}$ sobre $U_i$ é isomorfo a $\Spf(A_i)$, onde $A_i=\varprojlim_n A_{in}$ \sref{I.10.6.4};
$A_i$ é um anel $\mathfrak{K}A_i$-ádico, e $A_i/\mathfrak{K}A_i$, que é isomorfo a $A_{i0}$, é uma álgebra de tipo finito sobre $B/\mathfrak{K}$.
Conclui-se, pois, de \sref[0]{0.7.5.5} que $A_i$ é topologicamente isomorfo a um quociente de uma álgebra de séries formais restritas sobre $B$ (por um ideal necessariamente fechado, pois tal álgebra é noetheriana \sref[0]{0.7.5.4}).

Para demonstrar que (c) implica (a), pode-se restringir ao caso em que $\mathfrak{X}=\Spf(A)$ é também afim, sendo $A$ um anel ádico noetheriano isomorfo a um quociente de uma álgebra de séries formais restritas sobre $B$ por um ideal fechado.
Então, por \sref[0]{0.7.5.5}, $A/\mathfrak{K}A$ é uma álgebra de tipo finito sobre $B/\mathfrak{K}$ e $\mathfrak{K}A=\mathfrak{J}$ é um ideal de definição de $A$; portanto, em virtude de \sref{I.10.10.9}, as condições de (a) são satisfeitas.
\end{proof}

Observe-se que, se as condições da Proposição~\sref{I.10.13.1} são satisfeitas, a propriedade~(a) é válida para \emph{todo} ideal de definição $\sh{K}$ de $\mathfrak{Y}$ (em virtude de \emph{(c)}) e, por conseguinte, na propriedade~(b), \emph{todos} os $f_n$ são morfismos de tipo finito.

\begin{corollary}[10.13.2]
\label{I.10.13.2}
Se as condições de \sref{I.10.13.1} são satisfeitas, todo aberto formal afim noetheriano $V$ de $\mathfrak{Y}$ possui a propriedade~(\textbf{Q}) e, se $\mathfrak{Y}$ é noetheriano, $\mathfrak{X}$ também o é.
\end{corollary}

\begin{proof}
Isso resulta imediatamente de \sref{I.10.13.1} e \sref{I.6.3.2}.
\end{proof}

\begin{definition}[10.13.3]
\label{I.10.13.3}
Quando as propriedades equivalentes (a), (b) e (c) de \sref{I.10.13.1} são satisfeitas, diz-se que o morfismo $f$ é \emph{de tipo finito}, ou que $\mathfrak{X}$ é um \emph{$\mathfrak{Y}$-pré-esquema formal de tipo finito}, ou um \emph{pré-esquema formal de tipo finito sobre $\mathfrak{Y}$}.
\end{definition}

\begin{corollary}[10.13.4]
\label{I.10.13.4}
Sejam $\mathfrak{X}=\Spf(A)$ e $\mathfrak{Y}=\Spf(B)$ dois esquemas formais afins noetherianos;
para que $\mathfrak{X}$ seja de tipo finito sobre $\mathfrak{Y}$, é necessário e suficiente que o anel ádico noetheriano $A$ seja isomorfo ao quociente de uma álgebra de séries formais restritas sobre $B$ por um ideal fechado.
\end{corollary}

\begin{proof}
Com as notações de \sref{I.10.13.1}, se $\mathfrak{X}$ é de tipo finito sobre $\mathfrak{Y}$, então $A/\mathfrak{K}A$ é uma $(B/\mathfrak{K})$-álgebra de tipo finito em virtude de \sref{I.6.3.3}, e $\mathfrak{K}A$ é um ideal de definição de $A$ \sref{I.10.10.9}.
Conclui-se então por \sref[0]{0.7.5.5}.
\end{proof}

\begin{proposition}[10.13.5]
\label{I.10.13.5}
\medskip\noindent
\begin{enumerate}
  \item[\rm{(i)}] O composto de dois morfismos de pré-esquemas formais de tipo finito é de tipo finito.
\oldpage[I]{209}
  \item[\rm{(ii)}] Sejam $\mathfrak{X}$, $\mathfrak{S}$ e $\mathfrak{S}'$ três pré-esquemas formais localmente noetherianos (resp. noetherianos), e sejam $f:\mathfrak{X}\to\mathfrak{S}$ e $g:\mathfrak{S}'\to\mathfrak{S}$ dois morfismos.
    Se $f$ é de tipo finito, $\mathfrak{X}\times_\mathfrak{S}\mathfrak{S}'$ é localmente noetheriano (resp. noetheriano) e é de tipo finito sobre $\mathfrak{S}'$.
  \item[\rm{(iii)}] Seja $\mathfrak{S}$ um pré-esquema formal localmente noetheriano, e sejam $\mathfrak{X}'$ e $\mathfrak{Y}'$ dois $\mathfrak{S}$-pré-esquemas formais localmente noetherianos tais que $\mathfrak{X}'\times_\mathfrak{S}\mathfrak{Y}'$ seja localmente noetheriano.
    Se $\mathfrak{X}$ e $\mathfrak{Y}$ são $\mathfrak{S}$-pré-esquemas formais localmente noetherianos, e $f:\mathfrak{X}\to\mathfrak{X}'$ e $g:\mathfrak{Y}\to\mathfrak{Y}'$ são dois $\mathfrak{S}$-morfismos de tipo finito, $\mathfrak{X}\times_\mathfrak{S}\mathfrak{Y}$ é localmente noetheriano e $f\times_\mathfrak{S}g$ é um $\mathfrak{S}$-morfismo de tipo finito.
\end{enumerate}
\end{proposition}

\begin{proof}
Pelo raciocínio formal de \sref{I.3.5.1}, (iii) decorre de (i) e (ii); basta, portanto, provar (i) e (ii).

Sejam $\mathfrak{X}$, $\mathfrak{Y}$ e $\mathfrak{Z}$ três pré-esquemas formais localmente noetherianos, e sejam $f:\mathfrak{X}\to\mathfrak{Y}$ e $g:\mathfrak{Y}\to\mathfrak{Z}$ dois morfismos de tipo finito.
Se $\sh{L}$ é um ideal de definição de $\mathfrak{Z}$, $\sh{K}=g^*(\sh{L})\sh{O}_\mathfrak{Y}$ é um ideal de definição de $\mathfrak{Y}$ e $\sh{J}=f^*(g^*(\sh{L}))\sh{O}_\mathfrak{X}$ é um ideal de definição de $\mathfrak{X}$.
Ponhamos $X_0=(\mathfrak{X},\sh{O}_\mathfrak{X}/\sh{J})$, $Y_0=(\mathfrak{Y},\sh{O}_\mathfrak{Y}/\sh{K})$ e $Z_0=(\mathfrak{Z},\sh{O}_\mathfrak{Z}/\sh{L})$, e sejam $f_0:X_0\to Y_0$ e $g_0:Y_0\to Z_0$ os morfismos correspondentes a $f$ e $g$, respectivamente.
Como, por hipótese, $f_0$ e $g_0$ são de tipo finito, $g_0\circ f_0$ também o é \sref{I.6.3.4}, e corresponde a $g\circ f$;
portanto, $g\circ f$ é de tipo finito em virtude de \sref{I.10.13.1}.

Nas condições de (ii), $\mathfrak{S}$ (resp. $\mathfrak{X}$, $\mathfrak{S}'$) é o limite indutivo de uma sequência $(S_n)$ (resp. $(X_n)$, $(S'_n)$) de pré-esquemas localmente noetherianos, e pode-se supor \sref{I.10.13.1} que $X_m=X_n\times_{S_n}S_m$ para $m\leq n$.
O pré-esquema formal $\mathfrak{X}\times_\mathfrak{S}\mathfrak{S}'$ é então o limite indutivo dos pré-esquemas $X_n\times_{S_n}S'_n$ \sref{I.10.7.4}, e tem-se
\[
    X_m\times_{S_m}S'_m = (X_n\times_{S_n}S_m)\times_{S_m}S'_m = (X_n\times_{S_n}S'_n)\times_{S'_n}S'_m.
\]
Além disso, $X_0\times_{S_0}S'_0$ é localmente noetheriano, pois $X_0$ é de tipo finito sobre $S_0$ \sref{I.6.3.8}.
Conclui-se primeiro, por \sref{I.10.12.3.1}, que $\mathfrak{X}\times_\mathfrak{S}\mathfrak{S}'$ é localmente noetheriano;
além disso, como $X_0\times_{S_0}S'_0$ é de tipo finito sobre $S'_0$ \sref{I.6.3.8}, resulta de \sref{I.10.12.3.1} e \sref{I.10.13.1} que $\mathfrak{X}\times_\mathfrak{S}\mathfrak{S}'$ é de tipo finito sobre $\mathfrak{S}'$, o que conclui a prova de (ii) (a afirmação relativa aos pré-esquemas noetherianos é consequência imediata de \sref{I.6.3.8}).
\end{proof}

\begin{corollary}[10.13.6]
\label{I.10.13.6}
Sob as hipóteses de \sref{I.10.9.9}, se $f$ é um morfismo de tipo finito, sua extensão $\widehat{f}$ aos completamentos também o é.
\end{corollary}

\subsection{Subpré-esquemas fechados de pré-esquemas formais}
\label{subsection:I.10.14}

\begin{proposition}[10.14.1]
\label{I.10.14.1}
Seja $\mathfrak{X}$ um pré-esquema formal localmente noetheriano e seja $\sh{A}$ um feixe coerente de ideais de $\sh{O}_\mathfrak{X}$.
Se $\mathfrak{Y}$ é o suporte (fechado) de $\sh{O}_\mathfrak{X}/\sh{A}$, o espaço topologicamente anelado $(\mathfrak{Y},(\sh{O}_\mathfrak{X}/\sh{A})|\mathfrak{Y})$ é um pré-esquema formal localmente noetheriano, que é noetheriano se $\mathfrak{X}$ o é.
\end{proposition}

\begin{proof}
Observe-se que $\sh{O}_\mathfrak{X}/\sh{A}$ é coerente em virtude de \sref{I.10.10.3} e \sref[0]{0.5.3.4}; portanto, seu suporte $\mathfrak{Y}$ é fechado \sref[0]{0.5.2.2}.
Seja $\sh{J}$ um ideal de definição de $\mathfrak{X}$ e seja $X_n=(\mathfrak{X},\sh{O}_\mathfrak{X}/\sh{J}^{n+1})$;
o feixe de anéis $\sh{O}_\mathfrak{X}/\sh{A}$ é o limite projetivo dos feixes $\sh{O}_\mathfrak{X}/(\sh{A}+\sh{J}^{n+1})=(\sh{O}_\mathfrak{X}/\sh{A})\otimes_{\sh{O}_\mathfrak{X}}(\sh{O}_\mathfrak{X}/\sh{J}^{n+1})$ \sref{I.10.11.3}, todos os quais têm suporte $\mathfrak{Y}$.
O feixe $(\sh{A}+\sh{J}^{n+1})/\sh{J}^{n+1}$ é um $\sh{O}_\mathfrak{X}$-módulo coerente, pois $\sh{J}^{n+1}$ é coerente; portanto, $(\sh{A}+\sh{J}^{n+1})/\sh{J}^{n+1}$ também é um $(\sh{O}_\mathfrak{X}/\sh{J}^{n+1})$-módulo coerente \sref[0]{0.5.3.10};
se $Y_n$ é o subpré-esquema fechado de $X_n$ definido por esse feixe de ideais, é imediato que $(\mathfrak{Y},(\sh{O}_\mathfrak{X}/\sh{A})|\mathfrak{Y})$
\oldpage[I]{210}
é o pré-esquema formal limite indutivo dos $Y_n$ e, como as condições de \sref{I.10.6.4} são satisfeitas, isso prova que esse pré-esquema formal é localmente noetheriano, e noetheriano se $\mathfrak{X}$ o é (pois então $Y_0$ o é em virtude de \sref{I.6.1.4}).
\end{proof}

\begin{definition}[10.14.2]
\label{I.10.14.2}
Chama-se \emph{subpré-esquema fechado} de um pré-esquema formal $\mathfrak{X}$ todo pré-esquema formal da forma $(\mathfrak{Y},(\sh{O}_\mathfrak{X}/\sh{A})|\mathfrak{Y})$, onde $\sh{A}$ é um \erratum[II]{ideal coerente de $\sh{O}_\mathfrak{X}$};
diz-se que esse pré-esquema é o subpré-esquema definido por $\sh{A}$.
\end{definition}

É claro que a correspondência assim definida entre os \erratum[II]{ideais coerentes de $\sh{O}_\mathfrak{X}$} e os subpré-esquemas fechados de $\mathfrak{X}$ é bijetiva.

O morfismo de espaços topologicamente anelados $j=(\psi,\theta):\mathfrak{Y}\to\mathfrak{X}$, onde $\psi$ é a injeção $\mathfrak{Y}\to\mathfrak{X}$ e $\theta^\sharp$ o homomorfismo canônico $\sh{O}_\mathfrak{X}\to\sh{O}_\mathfrak{X}/\sh{A}$, é evidentemente \sref{I.10.4.5} um morfismo de pré-esquemas formais, chamado de \emph{injeção canônica} de $\mathfrak{Y}$ em $\mathfrak{X}$.
Observe-se que, se $\mathfrak{X}=\Spf(A)$, onde $A$ é noetheriano e ádico, então $\sh{A}=\mathfrak{a}^\Delta$, onde $\mathfrak{a}$ é um ideal de $A$ \sref{I.10.10.5}; segue imediatamente do que precede que $\mathfrak{Y}=\Spf(A/\mathfrak{a})$, a menos de isomorfismo, e que $j$ corresponde \sref{I.10.2.2} ao homomorfismo canônico $A\to A/\mathfrak{a}$.

Diz-se que um morfismo $f:\mathfrak{Z}\to\mathfrak{X}$ de pré-esquemas formais localmente noetherianos é uma \emph{imersão fechada} se ele se fatora como $\mathfrak{Z}\xrightarrow{g}\mathfrak{Y}\xrightarrow{j}\mathfrak{X}$, onde $g$ é um isomorfismo de $\mathfrak{Z}$ sobre um subpré-esquema fechado $\mathfrak{Y}$ de $\mathfrak{X}$ e $j$ é a injeção canônica.
Como $j$ é um monomorfismo de espaços anelados, $g$ e $\mathfrak{Y}$ são necessariamente \emph{únicos}.

\begin{proposition}[10.14.3]
\label{I.10.14.3}
Uma imersão fechada é um morfismo de tipo finito.
\end{proposition}

\begin{proof}
Reduz-se imediatamente ao caso em que $\mathfrak{X}$ é um esquema formal afim $\Spf(A)$ e $\mathfrak{Y}=\Spf(A/\mathfrak{a})$;
a proposição decorre então da Proposição~\sref{I.10.13.1}[\emph{c}].
\end{proof}

\begin{lemma}[10.14.4]
\label{I.10.14.4}
Seja $f:\mathfrak{Y}\to\mathfrak{X}$ um morfismo de pré-esquemas formais localmente noetherianos, e seja $(U_\alpha)$ um recobrimento de $f(\mathfrak{Y})$ por abertos formais afins noetherianos de $\mathfrak{X}$ tais que os $f^{-1}(U_\alpha)$ sejam abertos formais afins noetherianos de $\mathfrak{Y}$.
Para que $f$ seja uma imersão fechada, é necessário e suficiente que $f(\mathfrak{Y})$ seja um subconjunto fechado de $\mathfrak{X}$ e que, para todo $\alpha$, a restrição de $f$ a $f^{-1}(U_\alpha)$ corresponda \sref{I.10.4.6} a um homomorfismo sobrejetivo $\Gamma(U_\alpha,\sh{O}_\mathfrak{X})\to\Gamma(f^{-1}(U_\alpha),\sh{O}_\mathfrak{Y})$.
\end{lemma}

\begin{proof}
As condições são evidentemente necessárias.
Reciprocamente, se elas são satisfeitas, e se denotamos por $\mathfrak{a}_\alpha$ o núcleo de $\Gamma(U_\alpha,\sh{O}_\mathfrak{X})\to\Gamma(f^{-1}(U_\alpha),\sh{O}_\mathfrak{Y})$, define-se um feixe coerente de ideais $\sh{A}$ de $\sh{O}_\mathfrak{X}$ tomando $\sh{A}|U_\alpha=\mathfrak{a}_\alpha^\Delta$ e tomando $\sh{A}$ igual a $\sh{O}_\mathfrak{X}$ no complementar da união dos $U_\alpha$.
\footnote{\emph{[Trad.] A fonte francesa imprime que $\sh{A}$ é nulo nesse complementar e, na frase seguinte, que o próprio $\mathfrak{a}_\alpha^\Delta$ tem suporte $U_\alpha\cap f(\mathfrak{Y})$. O ideal definidor deve ser o ideal unitário fora de $f(\mathfrak{Y})$, enquanto seu quociente tem esse suporte; ambas as correções são necessárias para que os feixes se colem no subpré-esquema fechado cujo espaço subjacente é $f(\mathfrak{Y})$.}}
Com efeito, como $f(\mathfrak{Y})$ é fechado e o suporte do quociente $\bigl(\sh{O}_\mathfrak{X}|U_\alpha\bigr)/\mathfrak{a}_\alpha^\Delta$ é $U_\alpha\cap f(\mathfrak{Y})$, tudo se reduz a provar que $\mathfrak{a}_\alpha^\Delta$ e $\mathfrak{a}_\beta^\Delta$ induzem o mesmo feixe sobre todo aberto formal afim noetheriano $V\subset U_\alpha\cap U_\beta$.
Mas a restrição a $f^{-1}(U_\alpha)$ de $f$ é uma imersão fechada desse pré-esquema formal em $U_\alpha$, $f^{-1}(V)$ é um aberto formal afim noetheriano de $f^{-1}(U_\alpha)$, e a restrição de $f$ a $f^{-1}(V)$ é uma imersão fechada;
se $\mathfrak{b}$ é o núcleo do homomorfismo sobrejetivo $\Gamma(V,\sh{O}_\mathfrak{X})\to\Gamma(f^{-1}(V),\sh{O}_\mathfrak{Y})$ correspondente a essa restrição, é imediato \sref{I.10.10.2} que $\mathfrak{a}_\alpha^\Delta$ induz $\mathfrak{b}^\Delta$ sobre $V$.
Definido assim o feixe de ideais $\sh{A}$, é claro que $f=j\circ g$\footnote{\emph{[Trad.] A fonte francesa imprime $f=g\circ j$; com $g:\mathfrak{Y}\to\mathfrak{Z}$ e $j:\mathfrak{Z}\to\mathfrak{X}$, o composto bem tipado é $j\circ g$.}}, onde $j:\mathfrak{Z}\to\mathfrak{X}$ é a injeção canônica do subpré-esquema fechado $\mathfrak{Z}$ de $\mathfrak{X}$ definido por $\sh{A}$, e $g$ é um isomorfismo de $\mathfrak{Y}$ sobre $\mathfrak{Z}$.
\end{proof}

\begin{proposition}[10.14.5]
\label{I.10.14.5}
\medskip\noindent
\begin{enumerate}
  \item[\rm{(i)}] Se $f:\mathfrak{Z}\to\mathfrak{Y}$ e $g:\mathfrak{Y}\to\mathfrak{X}$ são imersões fechadas de pré-esquemas formais localmente noetherianos, então $g\circ f$ é uma imersão fechada.
\oldpage[I]{211}
  \item[\rm{(ii)}] Sejam $\mathfrak{X}$, $\mathfrak{Y}$ e $\mathfrak{S}$ três pré-esquemas formais localmente noetherianos, seja $f:\mathfrak{X}\to\mathfrak{S}$ uma imersão fechada e seja $g:\mathfrak{Y}\to\mathfrak{S}$ um morfismo.
    Então o morfismo $\mathfrak{X}\times_\mathfrak{S}\mathfrak{Y}\to\mathfrak{Y}$ é uma imersão fechada.
  \item[\rm{(iii)}] Seja $\mathfrak{S}$ um pré-esquema formal localmente noetheriano, e sejam $\mathfrak{X}'$ e $\mathfrak{Y}'$ dois $\mathfrak{S}$-pré-esquemas formais localmente noetherianos tais que $\mathfrak{X}'\times_\mathfrak{S}\mathfrak{Y}'$ seja localmente noetheriano.
    Se $\mathfrak{X}$ e $\mathfrak{Y}$ são $\mathfrak{S}$-pré-esquemas formais localmente noetherianos, e $f:\mathfrak{X}\to\mathfrak{X}'$ e $g:\mathfrak{Y}\to\mathfrak{Y}'$ são $\mathfrak{S}$-morfismos que são imersões fechadas, então $f\times_\mathfrak{S}g$ é uma imersão fechada.
\end{enumerate}
\end{proposition}

\begin{proof}
Em virtude de \sref{I.3.5.1}, basta novamente provar (i) e (ii).

Para provar (i), pode-se supor que $\mathfrak{Y}$ (resp. $\mathfrak{Z}$) seja um subpré-esquema fechado de $\mathfrak{X}$ (resp. $\mathfrak{Y}$), definido por um feixe coerente $\sh{J}$ (resp. $\sh{K}$) de ideais de $\sh{O}_\mathfrak{X}$ (resp. $\sh{O}_\mathfrak{Y}$);
se $\psi$ é a injeção $\mathfrak{Y}\to\mathfrak{X}$ dos espaços subjacentes, então $\psi_*(\sh{K})$ é um feixe coerente de ideais de $\psi_*(\sh{O}_\mathfrak{Y})=\sh{O}_\mathfrak{X}/\sh{J}$ \sref[0]{0.5.3.12}, e portanto também um $\sh{O}_\mathfrak{X}$-módulo coerente \sref[0]{0.5.3.10};
o núcleo $\sh{K}_1$ de $\sh{O}_\mathfrak{X}\to(\sh{O}_\mathfrak{X}/\sh{J})/\psi_*(\sh{K})$ é, portanto, um feixe coerente de ideais de $\sh{O}_\mathfrak{X}$ \sref[0]{0.5.3.4}, e $\sh{O}_\mathfrak{X}/\sh{K}_1$ é isomorfo a $\psi_*(\sh{O}_\mathfrak{Y}/\sh{K})$, o que prova que $\mathfrak{Z}$ é isomorfo a um subpré-esquema fechado de $\mathfrak{X}$.

Para provar (ii), pode-se reduzir imediatamente ao caso em que $\mathfrak{S}=\Spf(A)$, $\mathfrak{X}=\Spf(B)$ e $\mathfrak{Y}=\Spf(C)$, onde $A$ é um anel noetheriano $\mathfrak{J}$-ádico, $B=A/\mathfrak{a}$ (sendo $\mathfrak{a}$ um ideal de $A$) e $C$ é uma $A$-álgebra topológica ádica e noetheriana.
Tudo se reduz então a provar que o homomorfismo $C\to C\widehat{\otimes}_A(A/\mathfrak{a})$ é \emph{sobrejetivo}:
mas $A/\mathfrak{a}$ é um $A$-módulo de tipo finito, e sua topologia é a topologia $\mathfrak{J}$-ádica;
resulta então de \sref[0]{0.7.7.8} que $C\widehat{\otimes}_A(A/\mathfrak{a})$ se identifica com $C\otimes_A(A/\mathfrak{a})=C/\mathfrak{a}C$, de onde segue a afirmação.
\end{proof}

\begin{corollary}[10.14.6]
\label{I.10.14.6}
Sob as hipóteses de \sref{I.10.14.5}[ii], sejam $p:\mathfrak{X}\times_\mathfrak{S}\mathfrak{Y}\to\mathfrak{X}$ e $q:\mathfrak{X}\times_\mathfrak{S}\mathfrak{Y}\to\mathfrak{Y}$ as projeções, de modo que o diagrama
\[
  \xymatrix{
    \mathfrak{X}\ar[d]_f &
    \mathfrak{X}\times_\mathfrak{S}\mathfrak{Y}\ar[l]_p\ar[d]^q\\
    \mathfrak{S} &
    \mathfrak{Y}\ar[l]^g
  }
\]
é comutativo.
Para todo $\sh{O}_\mathfrak{X}$-módulo coerente $\sh{F}$, tem-se então um isomorfismo canônico de $\sh{O}_\mathfrak{Y}$-módulos
\[
\label{I.10.14.6.1}
  u:g^*f_*(\sh{F})\isoto q_*p^*(\sh{F}).
  \tag{10.14.6.1}
\]
\end{corollary}

\begin{proof}
Sabe-se que definir um homomorfismo $g^*f_*(\sh{F}) \to q_*p^*(\sh{F})$ equivale a definir um homomorfismo $f_*(\sh{F}) \to g_*q_*p^*(\sh{F}) = f_*p_*p^*(\sh{F})$ \sref[0]{0.4.4.3}:
toma-se como $u$ o homomorfismo que corresponde por essa adjunção a $f_*(\rho)$,\footnote{\emph{[Trad.] A fonte francesa escreve $u=f_*(\rho)$; esses morfismos vivem sobre espaços distintos e têm feixes fonte e alvo distintos, sendo $f_*(\rho)$ o adjunto de $u$.}} onde $\rho$ é o homomorfismo canônico $\sh{F}\to p_*p^*(\sh{F})$ \sref[0]{0.4.4.3}.
Para ver que $u$ é um isomorfismo, pode-se reduzir imediatamente ao caso em que $\mathfrak{S}$, $\mathfrak{X}$ e $\mathfrak{Y}$ são os espectros formais de anéis ádicos noetherianos $A$, $B$ e $C$, respectivamente, que satisfazem as condições de \sref{I.10.14.5}[ii];
tem-se então $\sh{F}=M^\Delta$, onde $M$ é um $(A/\mathfrak{a})$-módulo de tipo finito \sref{I.10.10.5}, e os dois membros de \sref{I.10.14.6.1} se identificam, respectivamente, em virtude de \sref{I.10.10.8}, com $(C\otimes_A M)^\Delta$ e $((C/\mathfrak{a}C)\otimes_{A/\mathfrak{a}}M)^\Delta$; daí o corolário, pois $(C/\mathfrak{a}C)\otimes_{A/\mathfrak{a}}M = (C\otimes_A(A/\mathfrak{a}))\otimes_{A/\mathfrak{a}}M$ se identifica canonicamente com $C\otimes_A M$.
\end{proof}

\oldpage[I]{212}
\begin{corollary}[10.14.7]
\label{I.10.14.7}
Sejam $X$ um pré-esquema usual localmente noetheriano, $Y$ um subpré-esquema fechado de $X$, $j$ a injeção canônica $Y\to X$, $X'$ um subconjunto fechado de $X$ e $Y'=Y\cap X'$;
então $\widehat{j}:Y_{/Y'}\to X_{/X'}$ é uma imersão fechada e, para todo $\sh{O}_Y$-módulo coerente $\sh{F}$, tem-se
\[
    \widehat{j}_*(\sh{F}_{/Y'}) = (j_*(\sh{F}))_{/X'}.
\]
\end{corollary}

\begin{proof}
Como $Y'=j^{-1}(X')$, basta utilizar \sref{I.10.9.9} e aplicar \sref{I.10.14.5} e \sref{I.10.14.6}.
\end{proof}

\subsection{Pré-esquemas formais separados}
\label{subsection:I.10.15}

\begin{definition}[10.15.1]
\label{I.10.15.1}
Sejam $\mathfrak{S}$ um pré-esquema formal, $\mathfrak{X}$ um $\mathfrak{S}$-pré-esquema formal e $f:\mathfrak{X}\to\mathfrak{S}$ o morfismo estrutural.
Chama-se morfismo diagonal $\Delta_{\mathfrak{X}|\mathfrak{S}}:\mathfrak{X}\to\mathfrak{X}\times_\mathfrak{S}\mathfrak{X}$ (também denotado por $\Delta_\mathfrak{X}$) o morfismo $(1_\mathfrak{X},1_\mathfrak{X})_\mathfrak{S}$.
Diz-se que $\mathfrak{X}$ é separado sobre $\mathfrak{S}$, ou que é um $\mathfrak{S}$-esquema formal, ou que $f$ é um morfismo separado, se a imagem do espaço subjacente a $\mathfrak{X}$ por $\Delta_\mathfrak{X}$ é um subconjunto fechado do espaço subjacente a $\mathfrak{X}\times_\mathfrak{S}\mathfrak{X}$.
Diz-se que um pré-esquema formal $\mathfrak{X}$ é separado, ou que é um esquema formal, se é separado sobre $\bb{Z}$.
\end{definition}

\begin{proposition}[10.15.2]
\label{I.10.15.2}
Suponha-se que os pré-esquemas formais $\mathfrak{S}$ e $\mathfrak{X}$ sejam, respectivamente, limites indutivos de sequências $(S_n)$ e $(X_n)$ de pré-esquemas usuais, e que o morfismo $f:\mathfrak{X}\to\mathfrak{S}$ seja o limite indutivo de uma sequência de morfismos $f_n:X_n\to S_n$.
Para que $f$ seja separado, é necessário e suficiente que o morfismo $f_0:X_0\to S_0$ o seja.
\end{proposition}

\begin{proof}
Com efeito, $\Delta_{\mathfrak{X}|\mathfrak{S}}$ é então o limite indutivo da sequência de morfismos $\Delta_{X_n|S_n}$ \sref{I.10.7.4}; os espaços subjacentes a $\mathfrak{X}$ e $\mathfrak{X}\times_\mathfrak{S}\mathfrak{X}$ são respectivamente idênticos aos de $X_0$ e $X_0\times_{S_0}X_0$, e, sob essas identificações, a imagem de $\mathfrak{X}$ por $\Delta_{\mathfrak{X}|\mathfrak{S}}$ é idêntica à imagem de $X_0$ por $\Delta_{X_0|S_0}$;
daí a conclusão.
\end{proof}

\begin{proposition}[10.15.3]
\label{I.10.15.3}
Suponha-se, no que segue, que todos os pré-esquemas formais (resp. morfismos de pré-esquemas formais) considerados sejam limites indutivos de sequências de pré-esquemas usuais (resp. de morfismos de pré-esquemas usuais).
\begin{enumerate}
    \item[\rm{(i)}] O composto de dois morfismos separados é separado.
    \item[\rm{(ii)}] Se $f:\mathfrak{X}\to\mathfrak{X}'$ e $g:\mathfrak{Y}\to\mathfrak{Y}'$ são dois $\mathfrak{S}$-morfismos separados, então $f\times_\mathfrak{S}g$ é separado.
    \item[\rm{(iii)}] Se $f:\mathfrak{X}\to\mathfrak{Y}$ é um $\mathfrak{S}$-morfismo separado, então o $\mathfrak{S}'$-morfismo $f_{(\mathfrak{S}')}$ é separado para toda extensão $\mathfrak{S}'\to\mathfrak{S}$ do pré-esquema formal de base.
    \item[\rm{(iv)}] Se o composto $g\circ f$ de dois morfismos é separado, então $f$ é separado.
\end{enumerate}
\emph{(Fica subentendido neste enunciado que, se um mesmo pré-esquema formal $\mathfrak{Z}$ intervém várias vezes em uma mesma proposição, ele é considerado como limite indutivo da \emph{mesma} sequência $(Z_n)$ de pré-esquemas usuais em todos os lugares onde figura, e os morfismos de $\mathfrak{Z}$ em outro pré-esquema formal (resp. de um pré-esquema formal em $\mathfrak{Z}$) como limites indutivos de morfismos dos $Z_n$ em pré-esquemas usuais (resp. de pré-esquemas usuais nos $Z_n$)).}
\end{proposition}

\begin{proof}
Com as notações de \sref{I.10.15.2}, tem-se, com efeito, $(g\circ f)_0=g_0\circ f_0$ e $(f\times_\mathfrak{S}g)_0=f_0\times_{S_0}g_0$;
as afirmações de \sref{I.10.15.3} são então consequências imediatas de \sref{I.10.15.2} e das afirmações correspondentes de \sref{I.5.5.1} para os pré-esquemas usuais.
\end{proof}

Deixa-se ao leitor o cuidado de enunciar, para a mesma classe de pré-esquemas formais e de morfismos que em \sref{I.10.15.3}, as proposições correspondentes a \sref{I.5.5.5},
\oldpage[I]{213}
\sref{I.5.5.9} e \sref{I.5.5.10} (substituindo «aberto afim» por «aberto formal afim que satisfaz a condição (b) de \sref{I.10.6.3}»).

Um raciocínio análogo mostra também que todo esquema formal afim \emph{noetheriano} é separado, o que justifica a terminologia.

\begin{proposition}[10.15.4]
\label{I.10.15.4}
Sejam $\mathfrak{S}$ um pré-esquema formal localmente noetheriano, e $\mathfrak{X}$ e $\mathfrak{Y}$ dois $\mathfrak{S}$-pré-esquemas formais localmente noetherianos tais que $\mathfrak{X}$ ou $\mathfrak{Y}$ seja de tipo finito sobre $\mathfrak{S}$ (de modo que $\mathfrak{X}\times_\mathfrak{S}\mathfrak{Y}$ seja localmente noetheriano) e que $\mathfrak{Y}$ seja separado sobre $\mathfrak{S}$.
Seja $f:\mathfrak{X}\to\mathfrak{Y}$ um $\mathfrak{S}$-morfismo; então o morfismo gráfico $\Gamma_f=(1_\mathfrak{X},f)_\mathfrak{S}:\mathfrak{X}\to\mathfrak{X}\times_\mathfrak{S}\mathfrak{Y}$ é uma imersão fechada.
\end{proposition}

\begin{proof}
Pode-se supor que $\mathfrak{S}$ seja o limite indutivo de uma sequência $(S_n)$ de pré-esquemas localmente noetherianos, que $\mathfrak{X}$ (resp. $\mathfrak{Y}$) seja o limite indutivo de uma sequência $(X_n)$ (resp. $(Y_n)$) de $S_n$-pré-esquemas, e que $f$ seja o limite indutivo de uma sequência $(f_n:X_n\to Y_n)$ de $S_n$-morfismos;
então $\mathfrak{X}\times_\mathfrak{S}\mathfrak{Y}$ é o limite indutivo da sequência $(X_n\times_{S_n}Y_n)$, e $\Gamma_f$ o limite indutivo da sequência $(\Gamma_{f_n})$ \sref{I.10.7.4};
por hipótese, $Y_0$ é separado sobre $S_0$ \sref{I.10.15.2}, de modo que o espaço $\Gamma_{f_0}(X_0)$ é um subespaço fechado de $X_0\times_{S_0}Y_0$;
como os espaços subjacentes a $\mathfrak{X}\times_\mathfrak{S}\mathfrak{Y}$ (resp. $\Gamma_f(\mathfrak{X})$) e a $X_0\times_{S_0}Y_0$ (resp. $\Gamma_{f_0}(X_0)$) são idênticos, vê-se desde já que $\Gamma_f(\mathfrak{X})$ é um subespaço \emph{fechado} de $\mathfrak{X}\times_\mathfrak{S}\mathfrak{Y}$.
Observe-se agora que, quando $(U,V)$ percorre o conjunto dos pares formados por um aberto formal afim noetheriano $U$ (resp. $V$) de $\mathfrak{X}$ (resp. $\mathfrak{Y}$) tais que $f(U)\subset V$, os abertos $U\times_\mathfrak{S} V$ formam um recobrimento de $\Gamma_f(\mathfrak{X})$ em $\mathfrak{X}\times_\mathfrak{S}\mathfrak{Y}$ e, se $f_U:U\to V$ é a restrição de $f$ a $U$, então $\Gamma_{f_U}:U\to U\times_\mathfrak{S}V$ é a restrição de $\Gamma_f$ a $U$.
Se mostrarmos que $\Gamma_{f_U}$ é uma imersão fechada, então $\Gamma_f$ será uma imersão fechada \sref{I.10.14.4}; em outras palavras, somos reduzidos ao caso em que $\mathfrak{S}=\Spf(A)$, $\mathfrak{X}=\Spf(B)$ e $\mathfrak{Y}=\Spf(C)$ são afins (sendo $A$, $B$ e $C$ anéis ádicos noetherianos), e $f$ corresponde a um $A$-homomorfismo contínuo $\vphi:C\to B$;
então $\Gamma_f$ corresponde ao único homomorfismo contínuo $\omega: B\widehat{\otimes}_A C\to B$ que, composto com os homomorfismos canônicos $B\to B\widehat{\otimes}_A C$ e $C\to B\widehat{\otimes}_A C$, dá respectivamente a identidade e $\vphi$.
Mas é claro que $\omega$ é \emph{sobrejetivo}, de onde segue a afirmação.
\end{proof}

\begin{corollary}[10.15.5]
\label{I.10.15.5}
Sejam $\mathfrak{S}$ um pré-esquema formal localmente noetheriano e $\mathfrak{X}$ um $\mathfrak{S}$-pré-esquema de tipo finito;
para que $\mathfrak{X}$ seja separado sobre $\mathfrak{S}$, é necessário e suficiente que o morfismo diagonal $\mathfrak{X}\to\mathfrak{X}\times_\mathfrak{S}\mathfrak{X}$ seja uma imersão fechada.
\end{corollary}

\begin{proposition}[10.15.6]
\label{I.10.15.6}
Uma imersão fechada $j:\mathfrak{Y}\to\mathfrak{X}$ de pré-esquemas formais localmente noetherianos é um morfismo separado.
\end{proposition}

\begin{proof}
Com as notações de \sref{I.10.14.2}, $j_0:Y_0\to X_0$ é uma imersão fechada, portanto um morfismo separado, e basta aplicar \sref{I.10.15.2}.
\end{proof}

\begin{proposition}[10.15.7]
\label{I.10.15.7}
Sejam $X$ um pré-esquema (usual) localmente noetheriano, $X'$ um subconjunto fechado de $X$ e $\widehat{X}=X_{/X'}$.
Para que $\widehat{X}$ seja separado, é necessário e suficiente que $\widehat{X}_\red$ o seja, e basta que $X$ seja separado.
\end{proposition}

\begin{proof}
Com as notações de \sref{I.10.8.5}, para que $\widehat{X}$ seja separado, é necessário e suficiente que $X'_0$ o seja \sref{I.10.15.2}; como $\widehat{X}_\red=(X'_0)_\red$, isso equivale a que $\widehat{X}_\red$ seja separado \sref{I.5.5.1}[vi].
\end{proof}

% Projeção TeX, baseada primeiramente na fonte, do índice de notações de EGA I.
% Linhagem do rascunho: Markdown comparativo de Jacob C. Reinhold; rege a digitalização francesa do NUMDAM.
\clearpage
\section*{Índice de notações}
\phantomsection
\label{section:index-notation}
\begingroup
\small
\setlength{\columnsep}{2em}
\begin{multicols}{2}
\raggedcolumns
\newcommand{\EGAIndexLine}[2]{\par\noindent\hangindent=1em\hangafter=1 #1\ \dotfill\ \textup{#2}\par}

\phantomsection\label{section:index-notation-chapter0}
\medskip\noindent\textbf{Capítulo 0}\par\nopagebreak
\EGAIndexLine{$M_{[\phi]}$}{(0.1.0.2)}
\EGAIndexLine{$B\mathfrak{J}$, $\mathfrak{J}B$}{(0.1.0.3)}
\EGAIndexLine{$r(\mathfrak{a})$}{(0.1.1.1)}
\EGAIndexLine{$\mathfrak{R}(A)$}{(0.1.1.2)}
\EGAIndexLine{$S_{f}$}{(0.1.2.1)}
\EGAIndexLine{$S^{-1}A$, $S^{-1}M$, $m/s$, $i_A^S$, $i_M^S$, $i^S$}{(0.1.2.2)}
\EGAIndexLine{$A_{f}$, $M_{f}$, $A_{\mathfrak{p}}$, $M_{\mathfrak{p}}$}{(0.1.2.3)}
\EGAIndexLine{$S^{-1}u$}{(0.1.3.1)}
\EGAIndexLine{$\rho^{T,S}_{A}$, $\rho^{T,S}_{M}$, $\rho^{T,S}$}{(0.1.4.1)}
\EGAIndexLine{$\operatorname{Supp}(M)$}{(0.1.7.1)}
\EGAIndexLine{$\mathcal{F}|U$, $u|U$}{(0.3.1.5)}
\EGAIndexLine{$\mathcal{F}_{x}$, $s_{x}$, $\Gamma(U, \mathcal{F})$, $u(s)$, $\operatorname{Supp}(\mathcal{F})$}{(0.3.1.6)}
\EGAIndexLine{$\psi_{*}(\mathcal{F})$}{(0.3.4.1)}
\EGAIndexLine{$\psi_{*}(u)$}{(0.3.4.2)}
\EGAIndexLine{$\psi_{x}$}{(0.3.4.4)}
\EGAIndexLine{$\mathcal{G} \to \mathcal{F}$}{(0.3.5.1)}
\EGAIndexLine{$u^{\sharp}$, $v^{\flat}$, $\rho_{\mathcal{G}}$}{(0.3.5.3)}
\EGAIndexLine{$\psi^*(\mathcal{G})$, $\psi^*(v)$, $\sigma_{\mathcal{F}}$}{(0.3.5.5)}
\EGAIndexLine{$\mathcal{O}_{X}$, $\mathcal{O}_{X,x}$, $\mathcal{O}_{x}$, $1$, $e$}{(0.4.1.1)}
\EGAIndexLine{$\check{\mathcal{F}}$, $\bigwedge^{p} \mathcal{F}$}{(0.4.1.4)}
\EGAIndexLine{$\mathcal{J}\mathcal{F}$}{(0.4.1.5)}
\EGAIndexLine{$\Psi_{*}(\mathcal{F})$, $\Psi_{*}(u)$}{(0.4.2.1)}
\EGAIndexLine{$\Psi_{*}(\mathcal{C})$}{(0.4.2.4)}
\EGAIndexLine{$\Psi^*(\mathcal{G})$, $\Psi^*(v)$}{(0.4.3.1)}
\EGAIndexLine{$\Psi^*(\mathcal{C})$}{(0.4.3.4)}
\EGAIndexLine{$\Psi^*(\mathcal{J})\mathcal{O}_{X}$, $\mathcal{J}^{c}_{X}$}{(0.4.3.5)}
\EGAIndexLine{$\mathcal{G} \to \mathcal{F}$}{(0.4.4.1)}
\EGAIndexLine{$u_{\theta}^{\sharp}$, $u^{\sharp}$, $v_{\theta}^{\flat}$, $v^{\flat}$, $\rho_{\mathcal{G}}$, $\sigma_{\mathcal{F}}$}{(0.4.4.3)}
\EGAIndexLine{$u_{1} \otimes u_{2}$}{(0.4.4.4)}
\EGAIndexLine{$\mathcal{L}^{-1}$}{(0.5.4.3)}
\EGAIndexLine{$\mathcal{L}^{\otimes n}$}{(0.5.4.4)}
\EGAIndexLine{$\Gamma_{*}(\mathcal{L})$, $\Gamma_{*}(\mathcal{L}, \mathcal{F})$}{(0.5.4.6)}
\EGAIndexLine{$\mathcal{O}_{X}^{*}$}{(0.5.4.7)}
\EGAIndexLine{$\mathfrak{m}_{x}$, $\kappa(x)$, $f(x)$}{(0.5.5.1)}
\EGAIndexLine{$X_{f}$}{(0.5.5.2)}
\EGAIndexLine{$\operatorname{Im}(M' \otimes_{A} N')$}{(0.6.0)}
\EGAIndexLine{$\widehat A$, $\widehat M$}{(0.7.2.3) e (0.7.3.1)}
\EGAIndexLine{$A\{T_{1}, \cdots, T_{r}\}$}{(0.7.5.1)}
\EGAIndexLine{$A\{S^{-1}\}$}{(0.7.6.1)}
\EGAIndexLine{$\mathfrak{a}\{S^{-1}\}$}{(0.7.6.9)}
\EGAIndexLine{$A_{\{f\}}$, $\mathfrak{a}_{\{f\}}$}{(0.7.6.14)}
\EGAIndexLine{$A_{\{S\}}$}{(0.7.6.14)}
\EGAIndexLine{$(M \otimes_{A} N)^{\wedge}$, $M \widehat{\otimes}_{A} N$}{(0.7.7.1)}
\EGAIndexLine{$u \widehat{\otimes} v$}{(0.7.7.3)}

\phantomsection\label{section:index-notation-chapterI}
\medskip\noindent\textbf{Capítulo I}\par\nopagebreak
\EGAIndexLine{$\operatorname{Spec}(A)$, $\mathfrak{j}_{x}$, $\mathfrak{m}_{x}$, $\kappa(x)$, $f(x)$, $M_{x}$, $r(E)$, $V(E)$, $V(f)$, $D(f)$}{(I, \hyperref[I.1.1.1]{1.1.1})}
\EGAIndexLine{$\mathfrak{j}(Y)$}{(I, \hyperref[I.1.1.3]{1.1.3})}
\EGAIndexLine{${}^{a}\phi$}{(I, \hyperref[I.1.2.1]{1.2.1})}
\EGAIndexLine{$S_{f}'$}{(I, \hyperref[I.1.3.1]{1.3.1})}
\EGAIndexLine{$\rho_{g,f}$}{(I, \hyperref[I.1.3.3]{1.3.3})}
\EGAIndexLine{$\widetilde A$, $\widetilde{M}$, $\theta_{f}$}{(I, \hyperref[I.1.3.4]{1.3.4})}
\EGAIndexLine{$\widetilde u$}{(I, \hyperref[I.1.3.5]{1.3.5})}
\EGAIndexLine{$\widetilde{\phi}$}{(I, \hyperref[I.1.6.1]{1.6.1})}
\EGAIndexLine{$A(X)$}{(I, \hyperref[I.1.7.1]{1.7.1})}
\EGAIndexLine{$\mathcal{O}_{X/Y}$}{(I, \hyperref[I.2.1.6]{2.1.6})}
\EGAIndexLine{$\operatorname{Hom}(X, Y)$}{(I, \hyperref[I.2.2.1]{2.2.1})}
\EGAIndexLine{$\operatorname{Hom}_{S}(X, Y)$, $1_X$}{(I, \hyperref[I.2.5.2]{2.5.2})}
\EGAIndexLine{$\Gamma(X/S)$}{(I, \hyperref[I.2.5.5]{2.5.5})}
\EGAIndexLine{$X \amalg Y$}{(I, \hyperref[I.3.1.1]{3.1.1})}
\EGAIndexLine{$X \times_{S} Y$, $X \times Y$, $(g, h)_{S}$, $u \times_{S} v$, $u \times v$}{(I, \hyperref[I.3.2.1]{3.2.1})}
\EGAIndexLine{$X \times_{A} Y$, $X \otimes_{A} B$, $(g, h)_{A}$, $u \times_{A} v$}{(I, \hyperref[I.3.2.1]{3.2.1})}
\EGAIndexLine{$X_{(S')}$}{(I, \hyperref[I.3.3.6]{3.3.6})}
\EGAIndexLine{$f_{(S')}$}{(I, \hyperref[I.3.3.7]{3.3.7})}
\EGAIndexLine{$\Gamma_{f}$}{(I, \hyperref[I.3.3.14]{3.3.14})}
\EGAIndexLine{$X(T)$}{(I, \hyperref[I.3.4.1]{3.4.1})}
\EGAIndexLine{$P \times_{R} Q$}{(I, \hyperref[I.3.4.2]{3.4.2})}
\EGAIndexLine{$X(T)_{S}$}{(I, \hyperref[I.3.4.3]{3.4.3})}
\EGAIndexLine{$X(B)$, $X(B)_{A}$}{(I, \hyperref[I.3.4.4]{3.4.4})}
\EGAIndexLine{$X \otimes_{y} B$, $X \otimes_{\mathcal{O}_y} B$}{(I, \hyperref[I.3.6.3]{3.6.3})}
\EGAIndexLine{$Z \leq Y$}{(I, \hyperref[I.4.1.10]{4.1.10})}
\EGAIndexLine{$f^{-1}(Y)$}{(I, \hyperref[I.4.4.1]{4.4.1})}
\EGAIndexLine{$\mathcal{N}_{X}$}{(I, \hyperref[I.5.1.1]{5.1.1})}
\EGAIndexLine{$X_{\mathrm{red}}$}{(I, \hyperref[I.5.1.3]{5.1.3})}
\EGAIndexLine{$f_{\mathrm{red}}$}{(I, \hyperref[I.5.1.5]{5.1.5})}
\EGAIndexLine{$\Delta_{X/S}$, $\Delta_{X}$, $\Delta$}{(I, \hyperref[I.5.3.1]{5.3.1})}
\EGAIndexLine{$\operatorname{rg}_{K}(X)$}{(I, \hyperref[I.6.4.5]{6.4.5})}
\EGAIndexLine{$n(X)$}{(I, \hyperref[I.6.4.8]{6.4.8})}
\EGAIndexLine{$\Gamma_{\mathrm{rat}}(X/Y)$}{(I, \hyperref[I.7.1.2]{7.1.2})}
\EGAIndexLine{$R(X)$}{(I, \hyperref[I.7.1.3]{7.1.3})}
\EGAIndexLine{$\mathcal{R}(X)$}{(I, \hyperref[I.7.3.1]{7.3.1})}
\EGAIndexLine{$L(A)$}{(I, \hyperref[I.8.1.2]{8.1.2})}
\EGAIndexLine{$\delta(f)$}{(I, \hyperref[I.8.2.1]{8.2.1})}
\EGAIndexLine{$\mathcal{F} \otimes_{\mathcal{O}_{S}} \mathcal{G}$, $\mathcal{F} \otimes_{S} \mathcal{G}$}{(I, \hyperref[I.9.1.2]{9.1.2})}
\EGAIndexLine{$\overline{\mathcal{G}}$}{(I, \hyperref[I.9.4.1]{9.4.1})}
\EGAIndexLine{$\overline Y$}{(I, \hyperref[I.9.5.11]{9.5.11})}
\EGAIndexLine{$\operatorname{Spf}(A)$, $\mathcal{O}_{\mathfrak{X}}$ ($\mathfrak{X}=\operatorname{Spf}(A)$)}{(I, \hyperref[I.10.1.2]{10.1.2})}
\EGAIndexLine{$\mathfrak{D}(f)$}{(I, \hyperref[I.10.1.4]{10.1.4})}
\EGAIndexLine{${}^{a}\phi$, $\widetilde{\phi}$}{(I, \hyperref[I.10.2.1]{10.2.1})}
\EGAIndexLine{$\mathfrak{J}^{\Delta}$}{(I, \hyperref[I.10.3.1]{10.3.1})}
\EGAIndexLine{$\mathfrak{X} \times_{\mathfrak{S}} \mathfrak{Y}$}{(I, \hyperref[I.10.7.3]{10.7.3})}
\EGAIndexLine{$\mathcal{F}_{/X'}$, $\widehat{\mathcal{F}}$, $u_{/X'}$, $\widehat{u}$}{(I, \hyperref[I.10.8.4]{10.8.4})}
\EGAIndexLine{$X_{/X'}$, $\widehat{X}$}{(I, \hyperref[I.10.8.5]{10.8.5})}
\EGAIndexLine{$\widehat{f}$}{(I, \hyperref[I.10.9.1]{10.9.1})}
\EGAIndexLine{$M^{\Delta}$, $u^{\Delta}$}{(I, \hyperref[I.10.10.1]{10.10.1})}
\EGAIndexLine{$\Delta_{\mathfrak{S}|\mathfrak{X}}$, $\Delta_{\mathfrak{X}}$}{(I, \hyperref[I.10.15.1]{10.15.1})}
\end{multicols}
\endgroup

% Projeção brasileira, baseada primeiramente na fonte, do índice terminológico de EGA I.
% Rege a digitalização francesa do NUMDAM. As entradas conservam a ordem do índice
% da fonte para que cada linha possa ser cotejada diretamente com as pp. 219--224 impressas.
\clearpage
\section*{Índice terminológico}
\phantomsection
\label{section:index-terminology}
\begingroup
\footnotesize
\setlength{\columnsep}{2em}
\begin{multicols}{2}
\raggedcolumns
\newcommand{\EGATermIndexLine}[2]{\par\noindent\hangindent=1em\hangafter=1 #1\ \dotfill\ \textup{#2}\par}

\noindent\emph{A ordem das entradas segue o índice da fonte francesa;
por isso, as traduções em português nem sempre aparecem em ordem alfabética portuguesa.}\par\medskip

\EGATermIndexLine{Fecho de um subpré-esquema}{(I, \hyperref[I.9.5.11]{9.5.11})}
\EGATermIndexLine{$\mathcal O_X$-álgebra}{(0.4.1.3)}
\EGATermIndexLine{Anel ádico; anel $\mathfrak J$-ádico}{(0.7.1.9)}
\EGATermIndexLine{Anel admissível}{(0.7.1.2)}
\EGATermIndexLine{Anel completo de frações}{(0.7.6.5)}
\EGATermIndexLine{Anel de frações}{(0.1.2.2)}
\EGATermIndexLine{Anel de funções racionais}{(I, \hyperref[I.2.1.6]{2.1.6} e \hyperref[I.7.1.3]{7.1.3})}
\EGATermIndexLine{Anel de um esquema afim}{(I, \hyperref[I.1.7.1]{1.7.1})}
\EGATermIndexLine{Anel integral}{(0.1.0.6)}
\EGATermIndexLine{Anel topologizado linearmente}{(0.7.1.1)}
\EGATermIndexLine{Anel local}{(0.1.0.7)}
\EGATermIndexLine{Anel local dominante}{(I, \hyperref[I.8.1.1]{8.1.1})}
\EGATermIndexLine{Anel local de $X$ ao longo de $Y$; anel local de $Y$ em $X$}{(I, \hyperref[I.2.1.6]{2.1.6})}
\EGATermIndexLine{Anel pré-ádico; anel $\mathfrak J$-pré-ádico}{(0.7.1.9)}
\EGATermIndexLine{Anel pré-admissível}{(0.7.1.2)}
\EGATermIndexLine{Anel reduzido}{(0.1.1.1)}
\EGATermIndexLine{Anel regular}{(0.4.1.3)}
\EGATermIndexLine{Anéis locais relacionados}{(I, \hyperref[I.8.1.4]{8.1.4})}
\EGATermIndexLine{Anulador de um $\mathcal O_X$-módulo}{(0.5.3.7)}
\EGATermIndexLine{Aplicação de espectros associada a um homomorfismo de anéis}{(I, \hyperref[I.1.2.1]{1.2.1})}
\EGATermIndexLine{Aplicação racional; $S$-aplicação racional}{(I, \hyperref[I.7.1.2]{7.1.2})}
\EGATermIndexLine{Aplicação racional definida em um ponto}{(I, \hyperref[I.7.2.1]{7.2.1})}
\EGATermIndexLine{Aplicação racional induzida em um subconjunto aberto}{(I, \hyperref[I.7.1.2]{7.1.2})}
\EGATermIndexLine{Aplicação racional induzida em $\operatorname{Spec}(\mathcal O_x)$}{(I, \hyperref[I.7.2.8]{7.2.8})}

\EGATermIndexLine{$\mathcal O_X$-módulo coerente}{(0.5.3.1)}
\EGATermIndexLine{$\mathcal O_X$-álgebra coerente}{(0.5.3.6)}
\EGATermIndexLine{Completamento de um $\mathcal O_X$-módulo e de um homomorfismo de $\mathcal O_X$-módulos ao longo de um subconjunto fechado}{(I, \hyperref[I.10.8.4]{10.8.4})}
\EGATermIndexLine{Completamento de um pré-esquema ao longo de um subconjunto fechado}{(I, \hyperref[I.10.8.5]{10.8.5})}
\EGATermIndexLine{Componente irredutível}{(0.2.1.5)}
\EGATermIndexLine{Composto de um $\psi$-morfismo e um $\psi'$-morfismo}{(0.3.5.2)}
\EGATermIndexLine{Composto de um $\Psi$-morfismo e um $\Psi'$-morfismo}{(0.4.4.2)}
\EGATermIndexLine{Condição de colagem}{(0.3.3.1 e 0.4.1.6)}
\EGATermIndexLine{Corpo de base de um pré-esquema algébrico}{(I, \hyperref[I.6.4.1]{6.4.1})}
\EGATermIndexLine{Corpo residual de um ponto geométrico}{(I, \hyperref[I.3.4.5]{3.4.5})}

\EGATermIndexLine{Diagonal de $X\times_S X$}{(I, \hyperref[I.5.3.9]{5.3.9})}
\EGATermIndexLine{Di-homomorfismo}{(0.1.0.2)}
\EGATermIndexLine{Domínio de definição de uma aplicação racional}{(I, \hyperref[I.7.2.1]{7.2.1})}
\EGATermIndexLine{Dual de um $\mathcal O_X$-módulo}{(0.4.1.4)}

\EGATermIndexLine{Elemento topologicamente nilpotente}{(0.7.1.1)}
\EGATermIndexLine{$\mathcal O_X$-módulo gerado por uma família de seções}{(0.5.1.2)}
\EGATermIndexLine{Lugar dos zeros de uma seção}{(0.5.5.1)}
\EGATermIndexLine{Álgebra integral; álgebra integral finita}{(0.1.0.5)}
\EGATermIndexLine{Espaço anelado}{(0.4.1.1)}
\EGATermIndexLine{Espaço localmente anelado}{(0.5.5.1)}
\EGATermIndexLine{Espaço anelado induzido em um subconjunto aberto}{(0.4.1.2)}
\EGATermIndexLine{Espaço anelado normal, reduzido ou regular}{(0.4.1.3)}
\EGATermIndexLine{Espaço anelado obtido por colagem}{(0.4.1.6)}
\EGATermIndexLine{Espaço de Kolmogorov}{(0.2.1.2)}
\EGATermIndexLine{Espaço irredutível}{(0.2.1.1)}
\EGATermIndexLine{Espaço noetheriano}{(0.2.2.1)}
\EGATermIndexLine{Espaço quase-compacto}{(0.2.2.4)}
\EGATermIndexLine{Espaço subjacente de um espaço anelado}{(0.4.1.1)}
\EGATermIndexLine{Espaço anelado topológico}{(0.4.1.1)}

\EGATermIndexLine{Feixe algébrico sobre um espaço anelado}{(0.4.1.3)}
\EGATermIndexLine{Feixe associado a um pré-feixe}{(0.3.5.7)}
\EGATermIndexLine{Feixe associado a um $A$-módulo sobre $\operatorname{Spec}(A)$}{(I, \hyperref[I.1.3.4]{1.3.4})}
\EGATermIndexLine{Feixe com valores em uma categoria}{(0.3.1.1)}
\EGATermIndexLine{Feixe coerente de anéis}{(0.5.3.5)}
\EGATermIndexLine{Feixe de anéis normal em um ponto, normal, reduzido em um ponto, reduzido, regular em um ponto ou regular}{(0.4.1.3)}
\EGATermIndexLine{Feixe graduado de anéis}{(0.4.1.3)}
\EGATermIndexLine{Feixe de funções racionais}{(I, \hyperref[I.7.3.1]{7.3.1})}
\EGATermIndexLine{Feixe de torção}{(I, \hyperref[I.7.4.1]{7.4.1})}
\EGATermIndexLine{Feixe de ideais}{(0.4.1.3)}
\EGATermIndexLine{Feixe de ideais de definição}{(I, \hyperref[I.10.3.3]{10.3.3} e \hyperref[I.10.5.1]{10.5.1})}
\EGATermIndexLine{Feixe induzido}{(0.3.7.1)}
\EGATermIndexLine{Feixe localmente simples}{(0.3.6.1)}
\EGATermIndexLine{Feixe obtido por colagem}{(0.3.3.1)}
\EGATermIndexLine{Feixe pseudodiscreto}{(0.3.8.1)}
\EGATermIndexLine{Feixe simples}{(0.3.6.1)}
\EGATermIndexLine{Feixe estrutural de um espaço anelado}{(0.4.1.1)}
\EGATermIndexLine{Feixe estrutural de um esquema afim}{(I, \hyperref[I.1.3.4]{1.3.4})}
\EGATermIndexLine{Feixe sobre uma base de subconjuntos abertos}{(0.3.2.2)}
\EGATermIndexLine{Talo de um feixe}{(0.3.1.6)}
\EGATermIndexLine{Fibra de um morfismo de pré-esquemas}{(I, \hyperref[I.3.6.3]{3.6.3})}
\EGATermIndexLine{Álgebra finita}{(0.1.0.5)}
\EGATermIndexLine{Função racional}{(I, \hyperref[I.7.1.2]{7.1.2})}

\EGATermIndexLine{Generalização de um ponto}{(0.2.1.2)}
\EGATermIndexLine{$\mathcal O_X$-módulo graduado}{(0.4.1.3)}
\EGATermIndexLine{Gráfico de um morfismo}{(I, \hyperref[I.5.3.11]{5.3.11})}

\EGATermIndexLine{Homomorfismo contínuo de feixes de anéis topológicos}{(0.3.1.4)}
\EGATermIndexLine{Homomorfismo definido por uma seção}{(0.5.1.1)}
\EGATermIndexLine{Homomorfismo local de anéis locais}{(0.1.0.7)}
\EGATermIndexLine{$\phi$-homomorfismo de módulos}{(0.1.0.2)}

\EGATermIndexLine{Ideal de definição de um anel admissível}{(0.7.1.2)}
\EGATermIndexLine{$\mathcal O_X$-ideal}{(0.4.1.3)}
\EGATermIndexLine{$\mathcal O_{\mathfrak X}$-ideal de definição de um pré-esquema formal $\mathfrak X$}{(I, \hyperref[I.10.3.3]{10.3.3} e \hyperref[I.10.5.1]{10.5.1})}
\EGATermIndexLine{Ideal primo}{(0.1.0.6)}
\EGATermIndexLine{Imagem direta de um $\mathcal O_X$-módulo}{(0.4.2.1)}
\EGATermIndexLine{Imagem direta de um pré-feixe}{(0.3.4.1)}
\EGATermIndexLine{Imagem de uma $S$-seção}{(I, \hyperref[I.5.3.11]{5.3.11})}
\EGATermIndexLine{Imagem fechada de um pré-esquema por um morfismo}{(I, \hyperref[I.9.5.3]{9.5.3})}
\EGATermIndexLine{Imagem inversa de um $\mathcal O_Y$-módulo}{(0.4.3.1)}
\EGATermIndexLine{Imagem inversa de um $S$-morfismo}{(I, \hyperref[I.3.3.7]{3.3.7})}
\EGATermIndexLine{Imagem inversa de um pré-feixe}{(0.3.5.3)}
\EGATermIndexLine{Imagem inversa de um $S$-pré-esquema}{(I, \hyperref[I.3.3.6]{3.3.6})}
\EGATermIndexLine{Imagem inversa de um subpré-esquema}{(I, \hyperref[I.4.4.1]{4.4.1})}
\EGATermIndexLine{Imersão, imersão fechada ou imersão aberta de pré-esquemas}{(I, \hyperref[I.4.2.1]{4.2.1})}
\EGATermIndexLine{Imersão fechada de pré-esquemas formais}{(I, \hyperref[I.10.14.2]{10.14.2})}
\EGATermIndexLine{Imersão local de pré-esquemas}{(I, \hyperref[I.4.5.1]{4.5.1})}
\EGATermIndexLine{Injeção canônica de um espaço anelado induzido em um subconjunto aberto}{(0.4.1.2)}
\EGATermIndexLine{Injeção canônica de um subpré-esquema}{(I, \hyperref[I.4.1.7]{4.1.7})}
\EGATermIndexLine{Injeção canônica de um subpré-esquema formal}{(I, \hyperref[I.10.14.2]{10.14.2})}
\EGATermIndexLine{Injeção geométrica}{(I, \hyperref[I.3.5.4]{3.5.4})}
\EGATermIndexLine{Inverso de um $\mathcal O_X$-módulo invertível}{(0.5.4.3)}
\EGATermIndexLine{$\mathcal O_X$-módulo invertível}{(0.5.4.1)}
\EGATermIndexLine{Isomorfismo associado a uma imersão}{(I, \hyperref[I.4.2.1]{4.2.1})}
\EGATermIndexLine{Isomorfismo local de pré-esquemas}{(I, \hyperref[I.4.5.1]{4.5.1})}

\EGATermIndexLine{Limite projetivo de $\mathcal O_{X_n}$-módulos}{(I, \hyperref[I.10.6.6]{10.6.6})}
\EGATermIndexLine{$\mathcal O_X$-módulo localmente livre}{(0.5.4.1)}
\EGATermIndexLine{Caráter local de um ponto com valores em um anel local}{(I, \hyperref[I.3.4.5]{3.4.5})}

\EGATermIndexLine{Módulo de frações}{(0.1.2.2)}
\EGATermIndexLine{Módulo fielmente plano}{(0.6.4.1)}
\EGATermIndexLine{Módulo plano}{(0.6.1.1)}
\EGATermIndexLine{Módulo $A$-plano}{(0.6.2)}
\EGATermIndexLine{Módulo quase-finito}{(0.7.4.1)}
\EGATermIndexLine{$\mathcal O_X$-módulo}{(0.4.1.3)}
\EGATermIndexLine{Morfismo de espaços anelados ou de espaços anelados topológicos}{(0.4.1.1)}
\EGATermIndexLine{Morfismo fielmente plano}{(0.6.7.8)}
\EGATermIndexLine{Morfismo plano}{(0.6.7.1)}
\EGATermIndexLine{Morfismo de pré-feixes definidos sobre uma base de subconjuntos abertos}{(0.3.2.3)}
\EGATermIndexLine{Morfismo de pré-esquemas}{(I, \hyperref[I.2.2.1]{2.2.1})}
\EGATermIndexLine{Morfismo de $S$-pré-esquemas ou de $A$-pré-esquemas}{(I, \hyperref[I.2.5.2]{2.5.2})}
\EGATermIndexLine{Morfismo birracional}{(I, 6.5.6)}
\EGATermIndexLine{Morfismo de tipo finito}{(I, \hyperref[I.6.3.1]{6.3.1})}
\EGATermIndexLine{Morfismo diagonal}{(I, \hyperref[I.5.3.1]{5.3.1})}
\EGATermIndexLine{Morfismo dominante}{(I, \hyperref[I.2.2.6]{2.2.6})}
\EGATermIndexLine{Morfismo fechado}{(I, \hyperref[I.2.2.6]{2.2.6})}
\EGATermIndexLine{Morfismo gráfico de um morfismo}{(I, \hyperref[I.3.3.14]{3.3.14})}
\EGATermIndexLine{Morfismo localmente de tipo finito}{(I, \hyperref[I.6.6.2]{6.6.2})}
\EGATermIndexLine{Morfismo dominado por outro}{(I, \hyperref[I.4.1.8]{4.1.8})}
\EGATermIndexLine{Morfismo aberto}{(I, \hyperref[I.2.2.6]{2.2.6})}
\EGATermIndexLine{Morfismo quase-compacto}{(I, \hyperref[I.6.6.1]{6.6.1})}
\EGATermIndexLine{Morfismo radicial}{(I, \hyperref[I.3.5.4]{3.5.4})}
\EGATermIndexLine{Morfismo reduzido}{(I, \hyperref[I.5.1.5]{5.1.5})}
\EGATermIndexLine{Morfismo separado}{(I, \hyperref[I.5.4.1]{5.4.1})}
\EGATermIndexLine{Morfismo estrutural de um $S$-pré-esquema}{(I, \hyperref[I.2.5.1]{2.5.1})}
\EGATermIndexLine{Morfismo sobrejetivo}{(I, \hyperref[I.2.2.6]{2.2.6})}
\EGATermIndexLine{Morfismo universalmente injetivo}{(I, \hyperref[I.3.5.4]{3.5.4})}
\EGATermIndexLine{Morfismos equivalentes}{(I, \hyperref[I.7.1.1]{7.1.1})}
\EGATermIndexLine{$A$-morfismo ou $S$-morfismo de pré-esquemas}{(I, \hyperref[I.2.5.2]{2.5.2})}
\EGATermIndexLine{Morfismo de pré-esquemas formais}{(I, \hyperref[I.10.4.5]{10.4.5})}
\EGATermIndexLine{Morfismo ádico de pré-esquemas formais}{(I, \hyperref[I.10.12.1]{10.12.1})}
\EGATermIndexLine{Morfismo diagonal de pré-esquemas formais}{(I, \hyperref[I.10.15.1]{10.15.1})}
\EGATermIndexLine{Morfismo de tipo finito de pré-esquemas formais}{(I, \hyperref[I.10.13.3]{10.13.3})}
\EGATermIndexLine{Morfismo separado de pré-esquemas formais}{(I, \hyperref[I.10.15.1]{10.15.1})}
\EGATermIndexLine{Morfismo estrutural de um $\mathfrak S$-pré-esquema formal}{(I, \hyperref[I.10.4.7]{10.4.7})}
\EGATermIndexLine{$A$-morfismo ou $\mathfrak S$-morfismo de pré-esquemas formais}{(I, \hyperref[I.10.4.7]{10.4.7})}
\EGATermIndexLine{$\psi$-morfismo de pré-feixes}{(0.3.5.1)}
\EGATermIndexLine{$\Psi$-morfismo de um $\mathcal O_Y$-módulo a um $\mathcal O_X$-módulo}{(0.4.4.1)}

\EGATermIndexLine{Nilradical de um anel}{(0.1.1.1)}
\EGATermIndexLine{Nilradical de uma $\mathcal O_X$-álgebra}{(I, \hyperref[I.5.1.1]{5.1.1})}
\EGATermIndexLine{Número geométrico de pontos de um pré-esquema}{(I, \hyperref[I.6.4.8]{6.4.8})}

\EGATermIndexLine{Subconjunto aberto afim}{(I, \hyperref[I.2.1.1]{2.1.1})}
\EGATermIndexLine{Subconjunto aberto formal afim; subconjunto aberto formal afim ádico; subconjunto aberto formal afim noetheriano}{(I, \hyperref[I.10.4.1]{10.4.1})}

\EGATermIndexLine{Subconjunto multiplicativo de um anel}{(0.1.2.1)}
\EGATermIndexLine{Subconjunto multiplicativo saturado}{(0.1.4.3)}
\EGATermIndexLine{Módulo $f$-plano sobre $\mathcal O_X$}{(0.6.7.1)}
\EGATermIndexLine{Ponto de um pré-esquema com valores em um anel}{(I, \hyperref[I.3.4.4]{3.4.4})}
\EGATermIndexLine{Ponto de um pré-esquema com valores em um pré-esquema}{(I, \hyperref[I.3.4.1]{3.4.1})}
\EGATermIndexLine{Ponto de um $A$-pré-esquema com valores em uma $A$-álgebra}{(I, \hyperref[I.3.4.4]{3.4.4})}
\EGATermIndexLine{Ponto de um $S$-pré-esquema situado sobre um ponto $s\in S$}{(I, \hyperref[I.2.5.1]{2.5.1})}
\EGATermIndexLine{Ponto de um $S$-pré-esquema com valores em um $S$-pré-esquema}{(I, \hyperref[I.3.4.3]{3.4.3})}
\EGATermIndexLine{Ponto fechado}{(0.2.2.6)}
\EGATermIndexLine{Ponto genérico}{(0.2.1.2)}
\EGATermIndexLine{Ponto geométrico de um pré-esquema}{(I, \hyperref[I.3.4.5]{3.4.5})}
\EGATermIndexLine{Ponto geométrico situado sobre $s$}{(I, \hyperref[I.3.4.5]{3.4.5})}
\EGATermIndexLine{Ponto geométrico localizado em $x$}{(I, \hyperref[I.3.4.5]{3.4.5})}
\EGATermIndexLine{Ponto $K$-racional}{(I, \hyperref[I.3.4.5]{3.4.5})}
\EGATermIndexLine{Pré-feixe constante}{(0.3.6.1)}
\EGATermIndexLine{Pré-feixe sobre uma base de subconjuntos abertos}{(0.3.2.1)}
\EGATermIndexLine{Pré-esquema}{(I, \hyperref[I.2.1.2]{2.1.2})}
\EGATermIndexLine{Pré-esquema artiniano}{(I, \hyperref[I.6.2.1]{6.2.1})}
\EGATermIndexLine{Pré-esquema conexo}{(I, \hyperref[I.2.1.7]{2.1.7})}
\EGATermIndexLine{Pré-esquema de base}{(I, \hyperref[I.2.5.1]{2.5.1})}
\EGATermIndexLine{Pré-esquema obtido por redução módulo $\mathfrak J$}{(I, \hyperref[I.3.7.1]{3.7.1})}
\EGATermIndexLine{Pré-esquema induzido em um subconjunto aberto}{(I, \hyperref[I.2.1.8]{2.1.8})}
\EGATermIndexLine{Pré-esquema integral}{(I, \hyperref[I.2.1.7]{2.1.7})}
\EGATermIndexLine{Pré-esquema irredutível}{(I, \hyperref[I.2.1.7]{2.1.7})}
\EGATermIndexLine{Pré-esquema localmente integral}{(I, \hyperref[I.2.1.7]{2.1.7})}
\EGATermIndexLine{Pré-esquema localmente noetheriano}{(I, \hyperref[I.6.1.1]{6.1.1})}
\EGATermIndexLine{Pré-esquema reduzido associado a um pré-esquema}{(I, \hyperref[I.5.1.3]{5.1.3})}
\EGATermIndexLine{$A$-pré-esquema; pré-esquema sobre $A$}{(I, \hyperref[I.2.5.1]{2.5.1})}
\EGATermIndexLine{$K$-pré-esquema algébrico}{(I, \hyperref[I.6.4.5]{6.4.5})}
\EGATermIndexLine{$S$-pré-esquema; pré-esquema sobre $S$}{(I, \hyperref[I.2.5.1]{2.5.1})}
\EGATermIndexLine{$S$-pré-esquema dominante}{(I, \hyperref[I.2.5.1]{2.5.1})}
\EGATermIndexLine{Pré-esquema de tipo finito sobre $S$; $S$-pré-esquema de tipo finito}{(I, \hyperref[I.6.3.1]{6.3.1})}
\EGATermIndexLine{Pré-esquema obtido por extensão do pré-esquema de base}{(I, \hyperref[I.3.3.6]{3.3.6})}
\EGATermIndexLine{Pré-esquema separado sobre $S$}{(I, \hyperref[I.5.4.1]{5.4.1})}
\EGATermIndexLine{Pré-esquema formal}{(I, \hyperref[I.10.4.2]{10.4.2})}
\EGATermIndexLine{Pré-esquema formal ádico}{(I, \hyperref[I.10.4.2]{10.4.2})}
\EGATermIndexLine{Pré-esquema formal localmente noetheriano}{(I, \hyperref[I.10.4.2]{10.4.2})}
\EGATermIndexLine{Pré-esquema formal noetheriano}{(I, \hyperref[I.10.4.2]{10.4.2})}
\EGATermIndexLine{Pré-esquema formal sobre $A$; $A$-pré-esquema formal}{(I, \hyperref[I.10.4.7]{10.4.7})}
\EGATermIndexLine{Pré-esquema formal sobre $\mathfrak S$; $\mathfrak S$-pré-esquema formal}{(I, \hyperref[I.10.4.7]{10.4.7})}
\EGATermIndexLine{Pré-esquema formal de tipo finito sobre $\mathfrak S$; $\mathfrak S$-pré-esquema formal de tipo finito}{(I, \hyperref[I.10.13.3]{10.13.3})}
\EGATermIndexLine{Pré-esquema formal separado sobre $S$}{(I, \hyperref[I.10.15.1]{10.15.1})}
\EGATermIndexLine{Apresentação finita (módulo que admite uma)}{(0.1.0.5)}
\EGATermIndexLine{Apresentação finita ($\mathcal O_X$-módulo que admite uma)}{(0.5.2.5)}
\EGATermIndexLine{Princípio de indução noetheriana}{(0.2.2.2)}
\EGATermIndexLine{Produto de $S$-pré-esquemas}{(I, \hyperref[I.3.2.1]{3.2.1})}
\EGATermIndexLine{Produto de $\mathfrak S$-pré-esquemas formais}{(I, \hyperref[I.10.7.1]{10.7.1})}
\EGATermIndexLine{Produto fibrado de conjuntos}{(I, \hyperref[I.3.4.2]{3.4.2})}
\EGATermIndexLine{Produto tensorial de feixes sobre pré-esquemas distintos}{(I, \hyperref[I.9.1.2]{9.1.2})}
\EGATermIndexLine{Produto tensorial completado de álgebras}{(0.7.7.5)}
\EGATermIndexLine{Produto tensorial completado de módulos}{(0.7.7.2)}
\EGATermIndexLine{Projeções canônicas de um produto}{(I, \hyperref[I.3.2.1]{3.2.1})}
\EGATermIndexLine{Extensão canônica de um submódulo}{(I, \hyperref[I.9.4.1]{9.4.1})}
\EGATermIndexLine{Extensão de um morfismo aos completamentos}{(I, \hyperref[I.10.9.1]{10.9.1})}
\EGATermIndexLine{$p$-ésima potência exterior de um $\mathcal O_X$-módulo}{(0.4.1.4)}

\EGATermIndexLine{$\mathcal O_X$-módulo quase-coerente}{(0.5.1.3)}
\EGATermIndexLine{$\mathcal O_X$-álgebra quase-coerente}{(0.5.1.3)}

\EGATermIndexLine{Radical de um ideal}{(0.1.1.1)}
\EGATermIndexLine{Radical de um anel}{(0.1.1.2)}
\EGATermIndexLine{Posto de um $\mathcal O_X$-módulo localmente livre}{(0.5.4.1)}
\EGATermIndexLine{Posto de um $\mathcal O_X$-módulo sem torção}{(I, \hyperref[I.7.4.2]{7.4.2})}
\EGATermIndexLine{Posto de um $K$-esquema finito}{(I, \hyperref[I.6.4.5]{6.4.5})}
\EGATermIndexLine{Posto separável de um $K$-esquema finito}{(I, \hyperref[I.6.4.8]{6.4.8})}
\EGATermIndexLine{Restrição de um espaço anelado a um subconjunto aberto}{(0.4.1.2)}
\EGATermIndexLine{Restrição de um morfismo a um subconjunto aberto}{(0.4.1.2)}
\EGATermIndexLine{Restrição de um morfismo de pré-esquemas a um subpré-esquema}{(I, \hyperref[I.4.1.7]{4.1.7})}
\EGATermIndexLine{Restrição de um pré-esquema a um subconjunto aberto}{(I, \hyperref[I.2.1.8]{2.1.8})}
\EGATermIndexLine{Restrição de uma aplicação racional a um subconjunto aberto}{(I, \hyperref[I.7.1.2]{7.1.2})}

\EGATermIndexLine{Esquema}{(I, \hyperref[I.5.4.1]{5.4.1})}
\EGATermIndexLine{Esquema afim}{(I, \hyperref[I.1.7.1]{1.7.1})}
\EGATermIndexLine{Esquema local; esquema local em um ponto de um pré-esquema}{(I, \hyperref[I.2.4.1]{2.4.1})}
\EGATermIndexLine{$K$-esquema algébrico}{(I, \hyperref[I.6.4.1]{6.4.1})}
\EGATermIndexLine{Esquema finito sobre $K$; $K$-esquema finito}{(I, \hyperref[I.6.4.5]{6.4.5})}
\EGATermIndexLine{$S$-esquema}{(I, \hyperref[I.5.4.1]{5.4.1})}
\EGATermIndexLine{Esquema formal afim}{(I, \hyperref[I.10.1.2]{10.1.2})}
\EGATermIndexLine{$\mathfrak S$-esquema formal}{(I, \hyperref[I.10.15.1]{10.15.1})}
\EGATermIndexLine{Seções de um feixe}{(0.3.1.6)}
\EGATermIndexLine{$S$-seção de um $S$-pré-esquema}{(I, \hyperref[I.2.5.5]{2.5.5} e \hyperref[I.5.3.11]{5.3.11})}
\EGATermIndexLine{$S$-seção racional de um $S$-pré-esquema}{(I, \hyperref[I.7.1.2]{7.1.2})}
\EGATermIndexLine{Seção unidade de $\mathcal O_X$}{(0.4.1.1)}
\EGATermIndexLine{Séries formais restritas}{(0.7.5.1)}
\EGATermIndexLine{Soma de pré-esquemas}{(I, \hyperref[I.3.1.1]{3.1.1}; o índice da fonte imprime 3.6.1)}
\EGATermIndexLine{$\mathcal O_X$-subálgebra gerada por um submódulo}{(0.4.1.3)}
\EGATermIndexLine{Subpré-esquema}{(I, \hyperref[I.4.1.3]{4.1.3})}
\EGATermIndexLine{Subpré-esquema associado a uma imersão}{(I, \hyperref[I.4.2.1]{4.2.1})}
\EGATermIndexLine{Subpré-esquema fechado}{(I, \hyperref[I.4.1.3]{4.1.3})}
\EGATermIndexLine{Subpré-esquema fechado definido por um feixe de ideais}{(I, \hyperref[I.4.1.2]{4.1.2})}
\EGATermIndexLine{Subpré-esquema formal fechado}{(I, \hyperref[I.10.14.2]{10.14.2})}
\EGATermIndexLine{Especialização de um ponto}{(0.2.1.2)}
\EGATermIndexLine{Espectro de um anel}{(I, \hyperref[I.1.1.1]{1.1.1})}
\EGATermIndexLine{Espectro formal de um anel admissível}{(I, \hyperref[I.10.1.2]{10.1.2})}
\EGATermIndexLine{Suporte de um módulo}{(0.1.7.1)}
\EGATermIndexLine{Suporte de um feixe de grupos}{(0.3.1.6)}
\EGATermIndexLine{Sistema fundamental de ideais de definição}{(I, \hyperref[I.10.3.6]{10.3.6} e \hyperref[I.10.5.1]{10.5.1})}
\EGATermIndexLine{Sistema indutivo ádico de pré-esquemas}{(I, \hyperref[I.10.12.2]{10.12.2})}

\EGATermIndexLine{Topologia $\mathfrak J$-ádica; topologia $\mathfrak J$-pré-ádica}{(0.7.1.9 e 0.7.2.3)}
\EGATermIndexLine{Topologia espectral}{(I, \hyperref[I.1.1.2]{1.1.2})}
\EGATermIndexLine{$\mathcal O_X$-módulo invertível trivial}{(I, \hyperref[I.2.4.8]{2.4.8})}
\EGATermIndexLine{$\mathcal O_X$-módulo de tipo finito}{(0.5.2.1)}

\EGATermIndexLine{Valor de uma seção em um ponto}{(0.5.5.1)}

\medskip\noindent\textit{Nota terminológica.} A fonte emprega «ádico» somente
para anéis ou módulos separados e completos; esse uso difere um pouco do habitual.
\end{multicols}
\endgroup

\clearpage
\nocite{*}
\bibliography{the}
\bibliographystyle{amsalpha}
\end{document}
